% =====================================================================
%  Harald Høffding, ETIK.  En Fremstilling af de etiske Principer og deres
%  Anvendelse paa de vigtigste Livsforhold.  Fjerde, paany gennemsete og
%  delvis ændrede Udgave.  København og Kristiania: Gyldendalske Boghandel,
%  Nordisk Forlag, 1913.  vii + 606 pp.
%
%  TRANSCRIPTION -- GENERATED; DO NOT EDIT BY HAND.  Made in three steps, all rerunnable:
%    python3 ebook-method/etik/draft.py --emit --force
%      this template + the scan's text layer (the National Library of Norway's OCR, laid
%      into its scan by ebook-method/etik/nb_scan.py; ~/bibliotek/Høffding, Harald/
%      1913-etik-4udg.pdf, see SCANS.tsv), then patches.PATCHES;
%    python3 ebook-method/collate/check.py etik --italics --apply
%      the decisions in ebook-method/etik/decisions.tsv, taken against an independent
%      tesseract reading of the images (rules.py writes the ones settled by rule, among them
%      by the SAGA e-book of the 3rd edition, 1905, as a third reading: pagemap.EBOOK);
%    python3 ebook-method/etik/draft.py --post
%      patches.POST: corrections written against the collated text.
%  Run them chained (&&): each step exits without writing if a patch fails to match.
%  STATE (2026-10-04): COLLATED, NOT YET COMPLETE.  Pages 1--599 (the text) collated
%  word by word against a second OCR (tesseract) and, as a third reading, the e-book
%  of the 3rd edition (1905); every disagreement settled by rule or read at the
%  image.  Still to do: the footnotes (NB's OCR often ran a note into the body or lost
%  its mark), the italics, letterspaced names (\emph), about 110 readings the
%  collation could not apply itself, and the title page, the two prefaces, the
%  Indhold (pp. I--VII) and the Register (pp. 601--606).
%
%  CONVENTIONS.  "% --- p. N ---" + \opage{N} at the first word of printed page N;
%  a word broken over the turn is split with %.  Footnotes as printed, 1), 2), ...
%  numbered per page (\footnote[k]).  Names in small capitals \textsc; italics
%  \textit; letterspacing \emph.  Running heads, folios and the signature line at a
%  sheet's foot ("H. Høffding: Etik." and the sheet number) are not transcribed;
%  pp. 2, 180, 182, 262 and 600 are blank.  Print governs; misprints are kept and
%  marked PRINTED AS IS.
% =====================================================================
\documentclass[12pt]{article}
\usepackage[utf8]{inputenc}
\usepackage[T1]{fontenc}
\usepackage{libertinus}
\usepackage{libertinust1math}
\usepackage[danish]{babel}
\usepackage{textalpha}            % Greek: the subdivisions α, β, ... and quotations
\usepackage[protrusion=true,expansion=true]{microtype}
\usepackage{setspace}
\usepackage[margin=1.3in]{geometry}
\usepackage{fancyhdr}
\usepackage{hyperref}
\hypersetup{pdftitle={Etik (4. Udgave, 1913)}, pdfauthor={Harald Høffding}, hidelinks}
\pagestyle{fancy}
\fancyhf{}
\fancyhead[LE,RO]{\thepage}
\fancyhead[C]{\textit{Harald Høffding: Etik (1913)}}
\setstretch{1.3}
\usepackage{marginnote}
\usepackage{graphicx}
\DeclareUnicodeCharacter{0254}{\reflectbox{c}}  % ɔ: the old Danish id-est mark (a reversed c)
\renewcommand{\thefootnote}{\arabic{footnote})}
\newcommand{\tocline}[2]{\noindent #1\dotfill\ #2\par}
\newcommand{\opage}[1]{\marginnote{\footnotesize\textit{[#1]}}}
% the headings (etik): a part title on its own page, a group of chapters (Social Etik: A. FAMILIEN., 1. ÆGTESKABET.),
% a chapter (roman numeral over the title, as printed)
\newcommand{\partitle}[1]{\clearpage\vspace*{6em}\begin{center}{\Large #1}\end{center}\clearpage}
\newcommand{\grouphead}[1]{\bigskip\begin{center}{\large #1}\end{center}}
\newcommand{\chapterhead}[2]{\bigskip\begin{center}#1\\[3pt]{\large #2}\end{center}\medskip}

\begin{document}

% The title page, the prefaces to the first and the fourth edition and the Indhold
% (PDF 7-19): to be set.

% --- p. 1 ---
\setcounter{footnote}{0}%

\opage{1}\partitle{ETIKENS FORUDSÆTNINGER}

% --- p. 3 ---
\setcounter{footnote}{0}%

\opage{3}\chapterhead{I.}{POSITIV MORALITET OG VIDENSKABELIG ETIK.}

1. Etiske Domme indeholde Vurderinger af menneskelige
Handlinger. Naar vi i det virkelige Liv kalde en Handling god
eller ond, forklare vi derved ikke, hvorledes Handlingen er
bleven til, men vi udtale, hvilken Værdi den har i vore Øjne.
Enhver saadan Vurdering forudsætter dels, at der er en Trang, en
Følelse, som driver os til at bedømme Handlingen (hvad enten
nu Dommen blot fældes i vore egne Tanker eller tillige i Ord,
som rettes til Andre), --- dels, at vi have en Maalestok, et
Ideal, med hvilket Handlingen kan sammenholdes, og efter hvilket
den kan prøves. Motivet til etisk Vurdering kalder jeg vor
Etiks Grundlag eller Vurderingsmotiv. Dets Karakter bestemmes
væsentligt ved den Grundværdi, af hvilken vort Livs højeste
Formaal og dermed Maalestokken for vor etiske Vurdering afhænge.
Maalestokken for den etiske Vurdering bliver bestemmende
for vor Etiks Indhold, idet den afgør, hvilke Handlinger,
Livsretninger og Livsformer der kunne kaldes etisk gode.

En etisk Anskuelses Karakter beror paa, hvilket Grundlag
den forudsætter, hvilken Maalestok den anlægger, og hvilket Forhold
der finder Sted mellem dens Grundlag og dens Maalestok.

Men etiske Domme fældes for det Meste ikke med udtrykkelig
og klar Bevidsthed om de Principer, de forudsætte. En saa%
% --- p. 4 ---
\setcounter{footnote}{0}%
\opage{4}dan Bevidsthed fører ved videre Udvikling til videnskabelig Etik.
Men den vaagner kun, naar Reflexion og Tvivl gøre sig gældende.
Det Spørgsmaal fremtræder da her naturligt, hvilken
Berettigelse og hvilken Betydning den bevidste Fremdragen og
Diskussion af de etiske Principer da i det Hele har.

2. Ikke blot fældes der etiske Domme, Domme om Godt
og Ondt, før egentlig Tænken og videnskabelig Forsken vaagne;
men etiske Domme, som skulle have personlig Sandhed og
praktisk Betydning, maa stedse udspringe af en levende Følelse
og Drift, der ikke lader os i Ro, før vi faa udtalt os. Det er
Betingelsen for et sundt og dygtigt Liv, at vi ikke hente Principerne
for vore vigtigste Afgørelser langvejs fra og ved møjsommelig
Eftertænken, men at Afgørelsen udspringer af, hvad
der er blevet Kød og Blod i os. I de etiske Domme lægge vi
vor Personlighed for Dagen; de maa da være bestemte ved vor
hele Natur og ikke blot ved Ræsonnementer, som vi kunne anstille
i vore ledige Timer, om vi have saadanne. Livet giver os
desuden ikke altid Stunder til at tænke efter, men kræver, at
Dommen afsiges øjeblikkeligt.

Og selv om vi have Tid og Evne til at anstille Eftertænkning,
ville saa ikke vore Følelser og vore Drifter være saa
stærke, at de helt igennem --- uden at vi selv mærke det ---
bestemme vor Tænkning og lede den i Stedet for at ledes af
den? Er det ikke en Fordom, som Psykologien forlængst har
omstyrtet, at Fornuften skulde være suveræn Myndighed i vort
Indre, og har det ikke ofte nok vist sig, at hvad der fremtraadte
som Fornuftens Udsagn, i Virkeligheden var et Udslag
af Hjertets Trang? I Logik og Matematik kan maaske den rene
Fornuft tale gennem vore Domme; men dens Røst er for svag,
hvor det gælder det mest Konkrete, det mest Personlige af Alt,
nemlig menneskelige Handlinger.

Vil Reflexion og Kritik her ikke endog være betænkelig,
om ikke absolut skadelig? Reflexionen har altid en vis opløsende
Indflydelse, berøver os den instinktmæssige Sikkerhed
og Tryghed, hvormed vi begynde vort Liv. Den svækker vor
Kraft selv der, hvor den ikke lammer os ved de Tvivl, den
vækker. Vi handle ikke længer med samlet Energi og udtale
% --- p. 5 ---
\setcounter{footnote}{0}%
\opage{5}ikke længer vore Domme med Fynd og Frejdighed, ja blive
maaske tilsidst tilbøjelige til at holde vore Domme tilbage,
fordi det synes os umuligt at komme til nogen sikker Afgørelse.
Men saa tørres Livet ud og staar tilsidst helt stille. Ti etiske
Domme ere ikke blotte teoretiske Kuriositeter. De ere ikke blot
Udtryk for Følelse, men øve ogsaa bestemmende Indflydelse
baade paa den Dømmendes egne og paa Andres Følelse og
Handlen. De ere ikke blot motiverede, men ere ogsaa selv igen
Motiver, praktiske Kræfter af den største Betydning.

Betænkeligheden bliver endnu større, naar vi udvide Betragtningen
ud over det enkelte Individ. Den Enkeltes Følelser
og Drifter ere bestemte ved hele Slægtens Natur, Livsvilkaar og
Traditioner. Det enkelte Individ modtager nogle af sine mest
indgribende Evner og Instinkter som Arv fra Slægten. Det føres
dernæst, ved at udvikle sig og opdrages som Led af Familien,
Stammen og Staten, ind i en vis aandelig Atmosfære, idet Livsvaner,
Forestillinger, Tilskyndelser og Opgaver fremtræde for
det uden at kunne gøres til Genstand for dets Overvejelse og
Valg, men som Noget, det uvilkaarligt optager i sig. Som det
er udsprunget af Slægtens Moderskød, saaledes indsuger det
med Modermælken Slægtens Overleveringer. Dets Maade at
tænke, føle og handle paa er en ubevidst Arv fra tidligere Slægtled.
Gennem Raceinstinkter og Traditioner, gennem uvilkaarlig
Efterligning og Indøvelse grundlægges den Enkeltes Etik for
ham, førend hans bevidste Eftertanke kan gribe ind. Villien
ledes uvilkaarligt ind paa bestemte Baner, og af den uvilkaarligt
dannede Villiesretning bestemmes Livsinteresser og Formaal,
af disse igen Dommene om Godt og Ondt. Dyderne blive ---
som Jacobi et Sted siger --- til, før de faa Navne, og før de
gøres til Fordringer: „Livets Bog maa skrives, før man kan udarbejde
Registret til den; vore Moralteorier ere saadanne bagefter
udarbejdede Registre.“ Og --- føjer Jacobi til —de udarbejdes
i Regelen af Folk, der ikke forstaa Noget af selve Bogen!
I hvert Tilfælde er Bogen vigtigere end Registret, og Registret
kan ikke erstatte Bogen. Ti i de etiske Domme, de stærke
Følelsesudbrud, i hvilke der udtales Vurdering af menneskelige
Handlinger, have vi ikke blot det enkelte Individs Sindelag ud%
% --- p. 6 ---
\setcounter{footnote}{0}%
\opage{6}trykt, men det er Resultatet af Slægtens Erfaringer, som lægger
sig for Dagen i dem. Individet laver ikke selv sin Etik, „hitter“
ikke paa den, konstruerer den ikke absolut fra først af. Og kun
derved er det, at den faar sin rette Styrke. Den i Slægten
levende Etik er en Betingelse for det menneskelige Livs Sundhed
og Kraft. Den, som ved sin Reflexion og Kritik vil opløse,
hvad Naturen har sammenføjet, har ikke blot en uhyre Modstand
at overvinde, men maa ogsaa vel vide, hvad han gør, at han
ikke skal ende i et Uføre og drage Andre med sig derned.

Denne Slægtens virkelige og i Livet virkende Etik har man
kaldet den positive Moralitet. Den fremtræder i gængse Domme
og Grundsætninger, som ofte iklædes Ordsprogets Form, og som
enten kunne være blivende Udtryk for en Nations, en Stammes,
et religiøst Samfunds Livsvisdom, eller have en mere kortvarig
Tilværelse som et Aarhundredes eller en Tidsalders „offentlige
Mening“. Den fremtræder ogsaa i de personlige Forbilleder
(Religionsstiftere, Helte, Lovgivere o. s. v.), til hvilke den enkelte
Slægt, eller Slægt efter Slægt, ser hen som det højeste
Udtryk for det sande Menneskelige. Den fremtræder endeligt i
den Maade, paa hvilken de menneskelige Livsforhold og Livsformer
(Familie, borgerligt Samfund, Stat, Kirke) ere ordnede.
Den positive Lovgivning indeholder stedse en vis Del af den
positive Moralitet. Fælles for alle Former af positiv Moralitet
er det, at en Handling, som strider mod dem, bliver Genstand
for en Tilbagevirken fra Samfundets Side. En Krænkelse af de
herskende Grundsætninger, Forbilleder eller Livsformer medfører
et af Forargelse eller Harme fremkaldt Udbrud, den offentlige
Menings Dom, af hvilken den Handlendes fremtidige Ære afhænger.
Eller det bliver ikke ved denne moralske Censur, men
den Handlende maa yde Erstatning og maaske gennem Smerte
eller endog Død sone, hvad han har gjort, og berolige den
Harme, han har vakt. Denne Hævdelse af den positive Moralitet
gennem social Paavirkning betegnes ved Udtrykket Sanktion.
Den behøver ikke at føles af Individet som noget blot Ydre;
formedelst den oprindelige, uvilkaarlige Indleven i den positive
Moralitet vil Individet kunne føle dens Magt i sit Inderste som
Et med sin egen Samvittigheds Magt, --- hvis man idethele her
% --- p. 7 ---
\setcounter{footnote}{0}%
\opage{7}vil tale om Samvittighed; Samvittigheden existerer egentligt her
slet ikke som rent individuel Faktor.

Den positive Moralitet er et Leje, som Menneskelivets Strøm
er bleven ledet ind paa og stedse mere har uddybet, og som
Strømmen vedbliver at følge, saalænge nye Kanaler eller Opdæmninger
ikke hindre det. Men vil det være muligt at gribe
ind her med den bevidste Tænkning? Og hvis det er muligt,
vil det saa ikke medføre en skadelig Svækkelse af den Kraft,
med hvilken Strømmen iler frem?

3. Overfor saadanne velbegrundede og meget betydningsfulde
Indvendinger vil jeg forsøge at vise, at videnskabelig Etik
er saa langt fra at være enten umulig, eller skadelig eller unødvendig,
at den netop selv udvikler sig af den positive Moralitet,
virker klarende og frigørende tilbage paa denne, og løser Opgaver,
som ellers ikke kunne løses.

Allerførst er det klart, at man ikke kan forsvare den positive
Moralitet og hævde dens Betydning som Betingelse for det
menneskelige Livs Sundhed og Kraft, uden at man har en
Maalestok, efter hvilken man bedømmer den. Denne Maalestok
kan man ikke, naar man ikke vil dreje sig i en Kreds, hente
fra selve den positive Moralitet. Man maa da søge at klare sig,
hvilket Princip man lægger til Grund for sin Vurdering, --- men
saa er man allerede inde paa den Vej, der fører til videnskabelig
Etik. Saa snart man slaar fast, at Sundhed og Kraft er
bedre end Sygelighed og Mathed, har man allerede fastslaaet et
Vurderingsprincip. Det Spørgsmaal vil da naturligt rejse sig,
om den positive Moralitet virkeligt overalt fyldestgør dette Princips
Fordringer, --- om der ikke er andre Principer, der maa
lægges til Grund ved Vurderingen, --- og Diskussionen er da i
fuld Gang!

Det er selvmodsigende at optræde som Forsvarer for den
positive Moralitet, dersom man i al Reflexion og Kritik ser en
Fare. Overfor det, der skal virke med Ubevidsthedens Magt,
er Tavshed det Bedste. En alvorlig Tvivl kan kun mødes med
afgørende Modgrunde eller med Paavisning af, at der er Spørgsmaal,
som kunne henstaa ubesvarede, uden at Livet lider derunder.
Det hører til Menneskets ædleste Fornødenheder at spørge
% --- p. 8 ---
\setcounter{footnote}{0}%
\opage{8}„hvorfor?“, og Undertrykkelse af saadan Spørgen bliver let en
Undertrykkelse af det aandelige Liv idethele. Ogsaa i Spørgen
og Tvivl kan Livets Sundhed og Kraft vise sig.

Livet fører ad sin egen Vej til at spørge. Kun hvor Synskredsen
er snever og Opgaverne begrænsede, er der fuld Tryghed.
Men med voxende Erfaring begynder Sammenligningen
mellem de forskellige Love, de forskellige Idealer, de forskellige
Institutioner, som man lærer at kende hos forskellige Folkeslag
til forskellige Tider. Eller de nye Erfaringer og Situationer stille
Opgaver, som ikke kunne løses efter den overleverede Etik.
Eller den store Mangfoldighed af etiske Domme, som man fælder
og hører Andre fælde, søges ordnet, for at der kan skelnes mellem
det mere og det mindre Betydningsfulde. De etiske Domme
ere ganske vist fra først af uvilkaarlige Følelsesudbrud, og „om
Følelser kan man ikke disputere“; men man kan disputere om
de Forestillingers Gyldighed, til hvilke Følelserne ere knyttede, og
om de Handlingers Værdi, til hvilke Følelserne føre.

Det er sikkert et vanskeligt Punkt i den aandelige Udvikling,
naar Reflexionen og Kritiken vaagne hos det enkelte Individ
eller hos Folket. Instinktet og Driften miste deres umiddelbare
Kraft og Sikkerhed, og det er det store Spørgsmaal, om
der kan vindes Erstatning for dette Tab\footnote[1]{Smlgn. min Psykologi VII B, 3.}. Men hvor det er
Livet selv, der lægger os Spørgsmaalene i Munden, maa de
enten besvares, eller det maa afgøres, hvorfor de ikke kunne
besvares. Og det maa vel bemærkes, at Sikkerhed og Kraft
ikke ere absolute Goder. Søvngængeren gaar sikrere, end den
Vaagne vilde kunne; men vi vække eller stanse ham dog for
at hindre ham i at nærme sig Afgrunden. Den største Kraft
kan lægge sig for Dagen i en højst fordærvelig Retning. Naar
man da søger at forminske denne Kraft, forminsker man ogsaa
dens fordærvelige Virkninger. Og man maa da forsøge at skaffe
den bedre Indsigt, som vindes, al den Energi, der tidligere stod
i det umiddelbare Instinkts Tjeneste. Al sund aandelig Udvikling
bestaar jo i, at der overføres Energi fra snevrere til større
Formaal. Intet frembringes af Intet.

% --- p. 9 ---
\setcounter{footnote}{0}%
\opage{9}Den videnskabelige Etik vil ikke og kan ikke sætte sig i
den positive Moralitets Sted. Den vil kun begrunde, udvikle og
udvide den. Vi søge i den videnskabelige Etik kun at forstaa
os selv, at blive klare over, efter hvilke Principer vi føre vort
Liv, og at bringe disse Principer til større Klarhed og indbyrdes
Harmoni. Der finder i det menneskelige Aandsliv en stadig
Vexelvirkning Sted mellem det Bevidste og det Ubevidste,
saavel som mellem Erkendelse, Følelse og Villie indbyrdes.
Hvad der erhverves paa ét Omraade, kan komme alle andre
aandelige Omraader til Gode.

4. Der maa skelnes mellem to Opgaver for den videnskabelige
Etik. Den er dels en historisk, dels en filosofisk Videnskab.
Den historiske eller sammenlignende Etik søger at fremstille
den positive Moralitet, saaledes som den fremtræder hos
et givet Folk til en given Tid, at paavise, hvilken Udvikling
den under forskellige Forhold undergaar, og at sammenligne de
forskellige Former, den kan antage hos forskellige Folk og til
forskellige Tider. Den søger at finde Aarsagerne til disse forskellige
Udviklingstrin og Former i bestemte fysiske, psykologiske
og historiske Forhold. Den filosofiske Etik har ikke Beskrivelse
og Forklaring af givne etiske Fænomener, men Vurdering
af saadanne til Opgave. Ved Vurdering menes en Undersøgelse
af, hvorvidt en Handling eller Ordning af menneskelige
Livsforhold stemmer med det Formaal, som vi forudsætte for
al menneskelig Stræben, altsaa ogsaa med Grundværdien, den
Værdi, til hvilken vi føre alle andre Værdier tilbage. Den videnskabelige
Etik kan med Rette kaldes en praktisk Videnskab, da
den forudsætter et Formaal, som skal virkeliggøres gennem menneskelige
Handlinger. Enhver etisk Dom forudsætter et saadant
Formaal, ti Følelsen sættes kun i Bevægelse ved Synet af eller
Tanken om en Handling, naar denne enten fremmer eller hæmmer
Noget, hvis Bestaaen og Fremgang ligger os paa Hjerte.
En etisk Dom forudsætter en Livsinteresse, der enten begunstiges
eller krænkes ved en menneskelig Handling eller ved en
Ordning af menneskelige Forhold.

Det historisk Udviklede er ikke uden videre etisk berettiget.
Vi kunne forstaa, hvorledes den positive Moralitet har udviklet
% --- p. 10 ---
\setcounter{footnote}{0}%
\opage{10}sig til den Form, den har antaget i vort Land og til vor Tid,
og dog kunne vi bryde Staven over den i større eller mindre
Omfang. Den historiske Etik kan kun lære os, hvad der faktisk
bestaar, og hvorledes det er blevet et Bestaaende; men naar
man ikke antager, at det Virkelige som saadant er Et med det
Fornuftige, bliver der endnu en Vurdering af det at anstille.
Den filosofiske Etik er en systematisk gennemført Vurdering paa
et bestemt, i den menneskelige Natur givet Grundlag og efter
den til dette Grundlag svarende Maalestok. En saadan gennemført
Undersøgelse kan føre til Fordømmelse af visse Handlinger
og Livsformer og til Billigelse af andre. Det kan være, at der
paavises Inkonsekvenser indenfor den positive Moralitet, og da
den Forudsætning, at alle etiske Domme maa stemme overens
indbyrdes, nødvendigvis maa ligge til Grund, bliver det Opgaven
at udvikle den positive Moralitet saaledes, at Modsigelsen forsvinder.
Det kan fremdeles ved den sammenlignende Betragtning
af forskellige Former for positiv Moralitet vise sig, at den ene
Form har Fortrinnet for den anden; dette forudsætter en bestemt
Maalestok, hvorved Sammenligningen begrundes, og det
kan føre til, at man søger Midler til at udvikle den Type af
positiv Moralitet, der erkendes at staa lavere, i Retning af den,
der erkendes at staa højere. Det kan endeligt ofte ske, at de
oprindelige Motiver til visse Handlinger og Institutioner ikke
længere kunne anerkendes, uden at dog derfor selve disse Handlinger
og Institutioner miste deres Værdi og Betydning. De
kunne under de forandrede Forhold og ved at ses fra en anden
Side end tidligere medføre Virkninger, som man forhen ikke
kendte. Karakteregenskaber som ridderligt Mod, Kyskhed, Selvbeherskelse,
der ere udviklede under raa og barbariske Forhold,
kunne tages i Tjeneste for helt andre Opgaver end de oprindelige.
Privatejendommen er historisk opstaaet under Indflydelse
af Magtlysten og Erhvervsdriften og er bleven anerkendt af Samfundsmagten
i Fredens Interesse. Men ved den etiske Drøftelse
af dens blivende Betydning og Berettigelse er der helt andre
Hensyn, der blive at tage i Betragtning. Det gaar her, som naar
man i en By anvender de gamle Volde, der oprindeligt vare
anlagte til Forsvar, til Spadseregange og Lystanlæg. Der fore%
% --- p. 11 ---
\setcounter{footnote}{0}%
\opage{11}gaar i saadanne Tilfælde en Forskydning paa det etiske Omraade,
hvorved de gamle Motiver og Værdier falde bort, og nye
Motiver og Værdier træde i deres Sted, men Handlinger, Vaner
og Institutioner vedblive at være de samme. Det bliver herved
indlysende, at den historiske Forklaring af et etisk Fænomen
ikke uden Videre er afgørende for dets Vurdering, da Muligheden
af en saadan Forskydning stedse staar aaben. Den filosofiske
Etik agter intet Bestaaende, blot fordi det bestaar; men
idet den søger at føre udover det Bestaaende, søger den at vise,
hvorledes det, der historisk har udviklet sig, kan anvendes og
omsættes i nye Former.

Medens den historiske og den filosofiske Etik fra denne Side
set have en meget forskellig Karakter, vil det fra en anden Side
set vise sig vanskeligt at holde dem skarpt ude fra hinanden. ---
Den historiske Etiker vil ikke kunne frigøre sig for visse etiske
Forudsætninger, som kunne faa Indflydelse paa det Lys, i hvilket
han ser Fortiden. Han maa endog have Blik for og Øvelse
i at anstille en systematisk Vurdering, forat han skal kunne opfatte
den indre Sammenhæng (eller Mangel paa Sammenhæng)
i andre Tiders og Folkeslags etiske Anskuelser og Sædvaner.
Og --- som det allerede er antydet --- saa snart han anstiller en
vurderende Sammenligning mellem forskellige Former for positiv
Moralitet og kalder nogle for lavere, andre for højere, forudsætter
han en bestemt Maalestok (og en bestemt Grundværdi)
og er altsaa allerede inde paa den filosofiske Etiks Omraade. ---
Den filosofiske Etiker vil naturligvis stedse være Barn af sin
Tid og sit Folk, og den i disse herskende positive Moralitet vil
paa mange Maader faa Indflydelse paa hans Resultater. Men
dette medfører kun, at den filosofiske Etiks Arbejde stedse maa
optages paany. Vi maa udskille og prøve det i vore Forgængeres
Arbejde, som kun skyldes forbigaaende Forhold og
Indflydelser. I Kants Etik f. Ex. mærker man tydeligt, at han
er en Tysker, at han er opdragen under pietistisk Indflydelse
og som moden Mand levede i Rationalismens Tid. Nogle af
Fortrinene og nogle af Fejlene i hans Etik ere at forklare
af disse Grunde. Men en saadan Afhængighed af Folk og
% --- p. 12 ---
\setcounter{footnote}{0}%
\opage{12}Tid\footnote[1]{Et interessant Exempel paa, hvorledes den nationale Aandsretning bestemmer de enkelte Tænkeres Arbejde og Resultater, frembyder den italienske Filosofis Historie, særligt hvad Etiken angaar. Smlgn. Luigi Ferri: National Character and Classicism in Italian Ethics. (International Journal of Ethics 1895). --- Smlgn. for det gamle Grækenlands Vedkommende: Leopold Schmidt: The Unity of the Ethics of ancient Greece. (ibid. 1894).} hindrer ikke, at den store filosofiske Etiker kan udtale
Tanker, som have Gyldighed for mange Slægter. Filosofisk Etik
bliver kun da umuliggjort ved den historiske Begrænsning, dersom
den vil være en ren Fornuftvidenskab, en Fremstilling af
evige Sandheder. Den er stedse og kan aldrig blive Andet end
den systematiske Gennemførelse af et historisk givet Standpunkt,
og den kan kun konstruere Idealer og Normer for visse bestemte
Livsbetingelser; den er paa sin Vis --- ligesom den positive
Moralitet --- et Middel i Kampen for Tilværelsen: et Middel
til Hævdelse af et opnaaet Livstrin og til Opnaaelse af et
nyt Livstrin, som kan fortjene Navn af et højere. Dens videnskabelige
Karakter beror paa den Klarhed, hvormed dens Grundlag
og dens Maalestok fastsættes, og paa den Konsekvens, med
hvilken dens Principer gennemføres i det Enkelte. Dens praktiske
Betydning beror paa, i hvilken Grad de saaledes vundne
Tanker og Resultater kunne gaa over i Kød og Blod hos større
eller mindre Kredse af Menneskeslægten, vinde Tilknytningspunkter
i stærke Følelser og Vaner og blive ledende ved Ordningen
af de forskellige Livsforhold. Erfaring viser, at der ofte
kan være en lang og indviklet Vej fra den almindelige etiske
Tanke til dens praktiske Gennemførelse. ---

(1913). I den nyeste Tid har den franske sociologiske Skole,
hvis vigtigste Repræsentanter ere Durkheim og Le vy Bruhl, med
stor Energi fremhævet Betydningen af det historiske Studium af
den positive Moralitet til forskellige Tider. De vente, at der engang
paa Grundlag af en saadan sammenlignende Etik vil kunne udvikles
en „moralsk Kunst“, der staar i samme Forhold til „Videnskaben
om Sæderne“ (la science des mœurs), som Lægekunsten
staar til Lægevidenskaben. Naar Videnskaben har afgjort,
hvorledes enhver Samvittighedsforpligtelse er bleven til,
% --- p. 13 ---
\setcounter{footnote}{0}%
% CHECK p. 13: notes paired with the marks in order (the layer misread a number)
\opage{13}er bleven bestyrket og anerkendt, --- hvilke Virkninger den har
frembragt, og hvilken Funktion den udøver i det sociale Liv,
--- saa ville vi ogsaa vide, i hvilken Grad det er muligt og
gavnligt at ændre den. Da vil der kunne udøves en moralsk
Kunst. I Mellemtiden kan og skal man dog ikke lægge Hænderne
i Skødet, men følge den bedste Erkendelse, vor Viden
til given Tid gør mulig.\footnote[1]{Levy Bruhl: La Morale et la Science des Mæurs. p. 150. 223.} Desuden er det kun tilsyneladende,
at der er saa stor Strid mellem Filosofernes etiske Teorier. De
ere i Virkeligheden forskellige Synspunkter for den moralske
Grundkendsgerning (le fait moral): Samfundslivet, der paa én
Gang staar som Autoritet overfor de enkelte Individer og som
Kilde til de højeste Goder, disse have Adgang til.\footnote[2]{Durkheim i Bulletin de la Société Française de Philosophie. 1906. p. 113—139.}

Ogsaa denne betydningsfulde Retning indenfor etisk Forsken
kommer, tildels imod sin Villie, til at anerkende Forskellen
mellem historisk og filosofisk Etik og Vexelvirkningen mellem
dem. Den lader sig ikke nøje med at paavise „Moral“ som
Kendsgerning (le fait moral), men tillægger umiddelbart denne
Kendsgerning Værdi. Og den antyder, at al Forpligtelse skal
vurderes efter dens Virkninger og dens Betydning for det sociale
Liv; dermed er da strax forudsat en fra selve Kendsgerningerne
forskellig Maalestok. Og det etiske Problem angaar netop Muligheden
og Begrundelsen af en saadan Maalestok.

\chapterhead{II.}{TEOLOGISK OG FILOSOFISK ETIK.}

1. Den teologiske Etik staar i bestemt Modsætning baade
til den historiske Etik og til den filosofiske Etik. --- Den bygger
paa Overlevering, paa Sandheden som given, paa Idealet
% --- p. 14 ---
\setcounter{footnote}{0}%
\opage{14}som historisk aabenbaret. Forsaavidt falder den ganske vist ind
under den historiske Etik: ligesom den positive Retsvidenskab
fremstiller Systemet af de Regler og Bud, der ere fastslaaede i
et vist Lands juridiske Love, saaledes fremstiller enhver teologisk
Etik Systemet af de etiske Domme, der følge af et vist
kirkeligt Samfunds Trosbekendelse. Men den teologiske Etik
anerkender ikke den videnskabelige Historieforsknings Metode,
da den Aabenbaring, hvorpaa den bygger, skal skyldes et Gennembrud
af overnaturlige Kræfter, som ikke kan forklares ved
de fysiske, psykologiske og sociale Love, til hvilke den historiske
Videnskab støtter sig. Den fordrer en Særstilling for sit historiske
Grundlag og paastaar, at dette maa betragtes paa en helt
anden Maade end den øvrige Historie. --- Den nærmer sig den
filosofiske Etik, forsaavidt den ligesom denne anstiller en Vurdering
af de historisk givne Handlinger og Livsformer. Men den
anstiller ikke denne Vurdering efter Principer, som kunne vises
at have deres Udspring i den menneskelige Natur. Den tager
sit Udgangspunkt ikke i Menneskenaturen, men over denne, i
den overnaturlige Aabenbaring af et Ideal. Den bygger paa det
absolute Autoritetsprincip.\footnote[1]{Smlgn. om Autoritetens Begreb mit Skrift Om Grundlaget for den humane Etik. København 1876. Kap. 111.}

Autoritet betyder Myndighed, og Myndighed kan lægge sig
for Dagen paa meget forskellig Maade. Den kan grunde sig paa
fysisk Magt, saa at den virker ved den Frygt, den indgyder.
Den kan grunde sig paa Overlegenhed i Indsigt og i Dygtighed
og virke ved at afgive et Forbillede for Andre, som de beundrende
søge at efterligne. Den kan grunde sig paa Evne og
Villie til at hjælpe og skærme, saa at den ikke blot virker ved
at fremkalde Frygt og Beundring, men ved at vække Taknemmelighed,
Tillid og Kærlighed. Det menneskelige Liv kan ikke
undvære saadanne Autoriteter. Den ældre Slægt er den yngres
Opdrager og Forbillede. Forældrene ere Børnenes Autoriteter,
Lærerne Disciplenes; overalt er den videre Fremskredne Autoritet
for den, der ikke er saa langt fremme paa Vejen. Autoritetsforholdet
er et vigtigt Element i den positive Moralitet.

% --- p. 15 ---
\setcounter{footnote}{0}%
\opage{15}Det skyldes dog altid en vis Reflexion, en begyndende
Tænkning, naar Autoritetsprincipet slaas udtrykkeligt fast og opfattes
som ubetinget. Menneskene følge uvilkaarligt det, der
forekommer dem stort og fremragende, og behøve ikke nogen
særlig Befaling dertil. De følge endog af Efterligningsinstinkt
og Vane det Ubetydelige og Ligegyldige, naar det blot virker
ved Gentagelsens Magt, og naar der foreligger Præcedenser.
Først i Tider, hvor Kritiken er vaagnet, og hvor Tvivlen afløser
den trygge Hvilen i det halvt Ubevidste, opstiller man med Bevidsthed
Autoritetsprincipet i dets absolute Form som sidste
Grundlag for al Etik. Naar det ene Spørgsmaal kalder det andet
frem, saa det synes, som om Intet staar fast, vækkes Trangen
til et absolut Holdepunkt, der ikke drages med ind i Hvirvelen.
Menneskelig Følelse og Tænkning synes at være et usikkert
Grundlag for det Etiske. Menneskelivets største Spørgsmaal:
„hvad er godt og hvad er ondt?“ synes at behøve det stærkeste
Grundlag, og hvilket Grundlag kan være stærkere end et almægtigt
Væsens Villie? Naar man kan appellere til en saadan
Villie, maa al Diskussion høre op, og de larmende Tvivlsmaal
forstumme overfor den ubetingede Myndighed.

Godt bliver da det, som stemmer med Guds Villie, og det
er kun godt, fordi det stemmer med Guds Villie (bonum est,
qvia deus vult). Det var den Lære, som allerede \textsc{Platon} havde
haft Anledning til at bekæmpe (i Dialogen Eutyfron), men som
først ret formuleredes henimod Middelalderens Slutning af \textsc{Duns}
Scotus, der endog ligefrem udtalte, at hvis Gud havde foreskrevet
Mord, saa var Mord ingen Synd. I saa kras en Form
ville vel kun Faa vedkende sig denne Lære, men ethvert teologisk
Standpunkt indeholder en Tilnærmelse dertil. I Dogmet
om Pavens Ufejlbarhed og i den protestantiske Ortodoxis snevre
Bogstavtro indeholdes det ubetingede Autoritetsprincip. Gennem
Kirkens Overhoveds levende Ord eller gennem de af Kirken
overleverede Bøgers skrevne Ord kundgøres, hvad der ifølge
Guds Villie er godt eller ondt, og de, til hvem Ordet lyder,
have kun at modtage og følge det i ubetinget Lydighed.

Lad os forsøge at stille os paa dette Standpunkt for at se,
% --- p. 16 ---
\setcounter{footnote}{0}%
\opage{16}om dets Forudsætninger ikke medføre indre Modsigelser og
uovervindelige Vanskeligheder.

2. Vi spørge først, hvoraf man véd, at Noget er Guds
Villie, og hvilke Garantier man har derfor. Det kan ikke nytte,
at jeg véd, at hvis Noget er Guds Villie, er det godt, naar jeg
ikke med ubetinget Sikkerhed véd, hvad der i det enkelte Tilfælde
er Guds Villie. Denne Sikkerhed maa ifølge Sagens Natur
paa dette Standpunkt være mig det Vigtigste. Den Tro, den
teologiske Etik bygger paa, bliver derfor konsekvent en Tro
paa Kirken, den religiøse Aabenbarings og Traditions Bærer.
Kirken garanterer mig, at dens Overhoved er den rette Efterfølger
af Kirkens Stifter, eller at de bibelske Bøger indeholde den
sande Aabenbaring. I det katolske Ufejlbarhedsdogme finder
derfor den teologiske Etik sin konsekvente Afslutning. --- Hvis
jeg derimod umiddelbart vil erfare Guds Villie, maa jeg henholde
mig til Aandens Vidnesbyrd i mit eget Indre; men hvorledes
kan jeg bestemt skelne mellem den guddommelige Stemme
og mine egne naturlige Følelser og Tanker? Hvilken Maalestok
har jeg til at drage Grænsen mellem det Guddommelige og det
Menneskelige? Vi miste her den absolute Fasthed, som netop
er det, der skulde være Autoritetsprincipets Fortrin. Vi glide
fra det teologiske Standpunkt over til det psykologiske.

,Ja“, siger man, „vi klamre os netop til Autoritetsprincipet,
fordi vi føle, at vi trænge til det. Det er vor egen Afmagt
og Tvivl, der kaster os i Armene paa den ubetingede
Myndighed“. --- Men ved at betone denne personlige Trang
rejser man et psykologisk Problem i Stedet for det rent teologiske
Problem. Ikke Autoriteten selv, men min Trang til den
bliver da det egentlige Grundlag. Jeg retfærdiggør min Underkastelse
under Autoriteten ved min Trang; men kan da selve
denne Trang retfærdiggøres? Er det godt at tilfredsstille denne
Trang? Hvorledes skelner jeg, hvilke af mine mange Ønsker
og Fornødenheder det er godt at tilfredsstille, og hvilke der
skulle undertrykkes? Hvis man til Afgørelsen heraf beraaber
sig paa Autoriteten, drejer man sig i en Kreds.

En lignende Kredsgang i sin Tænken kommer man til,
naar man, som nogle Teologer forsøge, vil forbinde de to Sæt%
% --- p. 17 ---
\setcounter{footnote}{0}%
\opage{17}ninger: „det Gode er godt, fordi Gud vil det I', og: „Gud vil
det, fordi det er godt“. --- Dersom det Gode er det, der stemmer
med Guds Villie, hvad vil saa det sige, at Gud vil det,
fordi det er godt? Det kan aabenbart kun betyde, at han vil
det, fordi det er hans Villie. Derved komme vi altsaa ikke af
Stedet. --- Dersom paa den anden Side Gud først erkender Noget
som godt og derefter vil det, saa er der en Lov og Regel, for
hvilken hans Villie bøjer sig, og hans Villie er altsaa ikke i og
for sig selv Aarsag til, at det Gode er godt. Selv de, der tillægge
sig det fortroligste Kendskab til Guddommens Psykologi,
ville dog blive noget betænkelige ved den Delthed, der herved
statueres i Guddommens Væsen. Konsekvente teologiske Tænkere
have ogsaa indset, at der paa Autoritetsprincipets Standpunkt
ikke kan være Tale om Godt eller Ondt for selve Guddommens
Vedkommende. Gud kan hverken have Samvittighed
eller Pligt. Enhver Overførelse af etiske Forestillinger fra Mennesker
til Gud fører til Modsigelser, og den eneste Udvej er at
blive staaende ved Paaberaabelsen af den guddommelige Villie
og bøje sig for den i Lydighed. Man opgiver ellers selve det
teologiske Princip. Saa snart man anerkender, at der gives andre
Kilder til Kundskab om det Gode end selve Visheden om, hvad
Guds Villie er, er man egentlig fra den teologiske Etik gaaet
over til den filosofiske. Det Gode erkendes --- det bliver Hovedsætningen,
og at det saaledes erkendte Gode er Guds Villie,
bliver en Slutning, --- ganske vist en Slutning af stor praktisk
Betydning, idet derved en overnaturlig Sanktion bliver mulig,
men en Slutning, der grunder sig paa helt andre Forudsætninger
end Hovedsætningen og derfor meget godt kunde falde bort,
uden at denne faldt bort. ---

Den ubetingede Autoritet udtaler en Mangfoldighed af Bud
og Forbud. Saasnart Eftertanken vaagner, forsøge vi naturligt at
bringe Sammenhæng og Overensstemmelse mellem de enkelte
Autoritetskundgørelser indbyrdes. Men dette kan kun ske ved,
at vi anvende vor Fornufts Grundlov paa dem og gaa ud fra,
at den ene Kundgørelse ikke kan stride mod den anden. Men
naar vi have Ret til at anvende denne Fornuftlov, hvorfor skulde
da andre af vor naturlige Erkendelses Love være udelukkede?
% --- p. 18 ---
\setcounter{footnote}{0}%
\opage{18}Og have vi ikke allerede ved hin Anvendelse brudt med det
ubetingede Autoritetsprincip?

3. Det er dog ikke blot disse indre Modsigelser, som afholde
den filosofiske Etik fra at lægge teologiske Forudsætninger
til Grund.

Hvad der fra først af har fremkaldt filosofisk Etik og stedse
paany vækker Interessen for den, er den Overbevisning, at den
sidste Grund til det Etiske maa ligge i selve den menneskelige
Natur. For Alt, hvad Mennesket skal kunne anerkende som
sandt og godt, som skønt og stort, maa Maalestokken tilsidst
ligge i det selv. Principerne for al Forstaaelse og al Vurdering,
for al teoretisk og praktisk Virksomhed, maa findes i dets eget
Indre. Hvor højt Idealerne end staa over det, og hvor mægtigt
end Lovens Majestæt kundgør sig for dets Samvittighed, saa
ere de dog kun Ideal og Lov for det ved dets frie Anerkendelse,
ved den Tilslutning, det i Kraft af sin Natur yder dem. Sokrates
blev derfor Etikens Grundlægger ved sit Bud: Kend dig selv!
Han opstillede derved Personlighedsprincipet eller Subjektivitetens
Princip, paa én Gang Principet for den frie Forsken og for den
frie Samvittighed. Blind Lydighed som blivende Tilstand strider
mod dette Princip; man affører sig derved sin frie Personlighed
og gør sig til en upersonlig Maskine.

Dernæst medvirker ved Bestræbelsen for at udvikle en filosofisk
Etik ogsaa et andet Motiv, nemlig Ønsket om at gøre det
Etiske saa uafhængigt som muligt af omtvistelige Forudsætninger.
Netop fordi Dommene om Godt og Ondt gribe saa stærkt
ind i det menneskelige Liv og betinge de skarpeste Modsætninger
indenfor dettes Omraade, er det af stor Vigtighed, at ingen
andre Stridsspørgsmaal blandes sammen med de etiske. Derfor
maa etiske Spørgsmaal saavidt muligt gøres uafhængige af religiøse
og metafysiske Problemer. Dersom Etiken skulde vente,
indtil der opstod Enighed om dogmatiske Spørgsmaal, vilde den
faa Lov at vente længe. Det er da Umagen værdt at prøve, om
Enigheden paa det etiske Omraade ikke skulde være større end
paa det religiøse og metafysiske Omraade. Og da det netop er
den etiske Betydning af de religiøse og metafysiske Problemer,
som skal give disse den store Interesse, ligger det i Sagens
% --- p. 19 ---
\setcounter{footnote}{0}%
\opage{19}Natur, at man maa have et Grundlag og en Maalestok for det
Etiske, som --- til en vis Grad idetmindste --- maa være uafhængige
af de religiøse og metafysiske Antagelser\footnote[1]{En teologisk Kritiker (i Tidsskriftet Vidar 1888) har fundet „en Svaghed“ i mit Standpunkt udtrykt i de Vendinger: „saa uafhængigt som muligt af omtvistelige Forudsætninger“ og „til en vis Grad idetmindste“, der forekomme i dette Stykke. Jeg har valgt disse Udtryk med velberaad Hu, idet jeg er overbevist om, at der her kun kan være Tale om en Stræben efter at reducere de omtvistelige Forudsætninger til det mindst Mulige, en Stræben, der ikke af nogen Enkelt kan fuldt gennemføres. Enhver Tænker har sin Begrænsning og maa vide, at han ubevidst opererer med Forudsætninger, som maaske først senere kunne fremdrages og prøves. Jeg har altsaa brugt hine Vendinger af videnskabelig Forsigtighed. De rokke ikke ved Principet: at det Etiske skal undersøges for sig alene og ikke udledes af religiøse eller metafysiske Antagelser. Derved udelukkes naturligvis ikke, at den etiske Tænker kan opstille Ideer, man maaske vil kunne kalde religiøse eller metafysiske (smlgn. IV, 5 og XXXI, 3), naar han kun ikke gør Brug af dem ved Begrundelsen af det Etiske.}. Ti idet man
tillægger saadanne Antagelser en etisk Betydning, vurderer man
dem, og der maa da spørges, paa hvilket Grundlag man stiller
sig ved denne Vurdering, og hvilken Maalestok man anlægger.

Den almindelige videnskabelige Grundsætning, at man ikke
uden Nødvendighed maa forøge Principerne, har altsaa her ogsaa
en praktisk Betydning. Idet den filosofiske Etik søger at
stille sig udenfor Striden mellem de forskellige religiøse Bekendelser
og mellem de forskellige metafysiske Teorier (Spiritualisme,
Materialisme o. s. v.), vil den ikke blot opnaa et sikrere
Grundlag for det Etiske, men den vil ogsaa virke til Bedste for
Tolerance, Religionsfrihed og Lærefrihed. Den vil modvirke den
hadefulde Diskussionsmaade, som bestaar i at udlede umoralske
Konsekvenser af Modstanderens Anskuelser.

Det etiske Problem staar derfor indenfor Filosofien meget
selvstændigt og uafhængigt i Forhold til de andre filosofiske
Hovedproblemer. Der er naturligvis en Sammenhæng mellem
Erkendelsesteorien, som undersøger vor Erkendelses Grundlag,
Forudsætninger og Grænser, Metafysiken, som søger at forme
en Verdensanskuelse, en almindelig Teori om Tilværelsens Væsen,
% --- p. 20 ---
\setcounter{footnote}{0}%
% CHECK p. 20: notes paired with the marks in order (the layer misread a number)
\opage{20}og Etiken\footnote[1]{Smlgn. om den indre Sammenhæng mellem de store Problemer Den menneskelige Tanke, dens Former og Opgaver (1910). p. 246—254.}. Men Sammenhængen er ikke saa direkte og umiddelbar,
som man ofte har antaget. Og det vil i hvert Tilfælde
have sin Betydning at holde de forskellige Problemer saa meget
som muligt ude fra hverandre; derved udelukkes ikke, at de
kunne belyse hverandre paa mange Maader.

Den filosofiske Etik, som skal være en videnskabelig Lære
om den etiske Vurdering, maa stille sig paa samme Grundlag
som al anden Videnskab. De Principer, Resultater og Hypoteser,
som andre Videnskaber ere naaede til, søger den ikke at
rokke ved. Den forlanger ingen Særstilling. Dersom den skulde
opstille Forudsætninger, som stred mod dem, paa hvilke andre
Videnskaber bygge, saa vilde den komme ind paa det Omtvisteliges
Omraade. Sikrest vil den gaa, naar den bygger paa de
fælles Forudsætninger. Det er dog endnu en hyppig Antagelse,
at Etiken idetmindste paa ét Punkt, nemlig ved Spørgsmaalet
om Aarsagslovens Anvendelse paa Villieslivet, skulde være nødsaget
til at opstille Forudsætninger, som ikke blot afvige fra,
men ligefrem stride mod de almindelige videnskabelige Forudsætninger.
Dette Spørgsmaal vil senere blive særligt behandlet.

4. Naar der i det Foregaaende har været Tale om teologisk
Etik, saa maa det vel mærkes, at den kristelige Etik ikke
i alle sine Former og paa alle sine Udviklingstrin er en teologisk
Etik. Kristendommen begyndte ikke med et teologisk System,
ligesaalidt som den begyndte med en kirkelig Organisation.
Forud for System og Organisation gaar her som overalt,
hvor der finder Udvikling Sted, en mere ubestemt Livsform.
Den Etik, som forkyndes i de tre første Evangelier, er for en
stor Del en Menneskekærlighedens Lære. Den støttes ved teologiske
Forestillinger, og der henvises til Løn og Straf ien
hinsidig Verden. Men den Gud, til hvis gengældende Magt der
henvises, fremtræder dog hovedsageligt som Fuldkommenhedens,
det vil her efter Sammenhængen (Matth. V. 43---48)\footnote[2]{I Parallelstedet, Lucas VI. 36, staar Barmhjertighed i Stedet for Fuldkommenhed.} sige:
% --- p. 21 ---
\setcounter{footnote}{0}%
\opage{21}Menneskekærlighedens Ideal. Jesus stifter det messianske Rige,
til hvilket Adgang staar aaben for dem, som ere rene af Hjertet,
som hungre og tørste efter Retfærdighed, som ere fredsommelige
og elske Andre som sig selv. Det nye Riges Aand
og Vilkaar vare helt andre end de stolte Verdensrigers, som
Historien hidtil havde kendt, og end de, som det jødiske Folk
havde tænkt sig i sine nationale Fantasier.

Det Betydningsfulde ved den ældste kristelige Etik er den
Udvidelse og Inderliggørelse af Menneskekærligheden, som den
indleder. Først ved den især af Paulus grundlagte dogmatiske
Udvikling af Kristendommen bleve teologiske Lærdomme af afgørende
Betydning for den kristelige Etik. Da fik Kristendommen
to Principer, Troen og Kærligheden, som det ikke er
lykkedes at bringe i fuldstændig Harmoni med hinanden. Trosforskellighederne
bleve Skranker for Menneskekærligheden. Den
mest afgørende Maalestok for et Menneskes Værd blev nu ikke,
om det havde et inderligt og kærligt Sind, var rent af Hjertet
og søgte Sandheden, men om det havde den rette Tro. Først
hvor Trosforskellighederne behandles saaledes, som Kristendommen
selv havde behandlet Nationalforskellighederne og andre
ydre Forskelle, først dér har den store etiske Ide, der udtaltes
i den ældste Kristendom, faaet sin fulde Udvikling\footnote[1]{Smlgn. Om Grundlaget for den humane Etik. p. 70 f. --- I dette Skrift tog jeg ved min Omtale af dette Punkt ikke tilstrækkeligt Hensyn til den ældste, forud for den bestemtere teologiske Udformning liggende Kristendom.}.

Dog maa det til rigtig Forstaaelse af Urkristendommens
Etik vel mærkes, at dens dybeste Motiv alligevel ikke var selve
Menneskekærligheden, men den levende, begejstrede Forvisning
om, at Himmeriges Rige var nær, og at alle jordiske Tings
Ende nærmede sig. Overfor dette store Haab tabte alle menneskelige
Formaal og Opgaver deres Vægt. Der kunde ikke blive
Tale om sammenhængende etisk Stræben, og selv Menneskekærlighedens
Ide fik ikke Lov at udfolde sine Konsekvenser.
Først da Kirken skød hin Forventning i Baggrunden, blev Kristendommen
en Kulturmagt. Og Menneskekærlighedens Ide (som
% --- p. 22 ---
\setcounter{footnote}{0}%
\opage{22}iøvrigt ikke først er undfangen af Urkristendommen) har vist
sig at have Livskraft baade ud over hin ekstatiske Forventnings
og ud over de kirkelige Dogmers Grænser\footnote[1]{Smlgn. min Afhandling: Hedenske Sandhedssøgere (i „Mindre Arbejder'' I) og Slutningsafsnittet i min Bog Søren Kierkegaard som Filosof (1892).}.

\chapterhead{III.}{ETIKENS PRINCIPER OG METODE.}

1. Etiske Domme, Domme om Godt og Ondt, opstaa fra
først af uden Begrundelse. Ligesom vi overfor et Kunstværk
eller overfor Naturen kunne lette vort Hjerte ved at udbryde:
„Det er skønt!“ saaledes kunne vi overfor en menneskelig Handling
lette vort Hjerte ved at udbryde: „Det er godt!“ Vi ledes
af en umiddelbar Følelse, der vækkes ved Synet af eller Erindringen
om Handlingen. Hvilke Handlinger vi kalde gode (eller
onde), vil for den Enkelte fra først af bero paa, hvilken «positiv
Moralitet“ (smlgn. I, 2) han har levet sig ind i. Et etisk
Problem opstaar først, naar denne positive Moralitet ikke taler
tydeligt, eller naar den viser sig at indeholde Selvmodsigelser.
Da trænges der til Principer, til ledende Tanker, som kunne bestemme
Valget mellem de foreliggende Muligheder.

Man kunde tvivle om Muligheden af at opstille begrundede
etiske Domme, altsaa om Muligheden af en filosofisk Etik. Etiske
Domme indeholde jo en Fordring, udsige, hvad der bør være,
ikke blot, hvad der er. Selv naar de udsige Ros eller Dadel om
en virkeligt udført Handling, betyde de egentligt en Fordring
om, at saadanne Handlinger skulle udføres eller undlades for
Fremtiden. Men gaar al Videnskab ikke ud paa at forklare,
finde Aarsagerne til det, som er? Kan der gives en Videnskab
om det, som ikke er, men bør være? Den filosofiske Etik gaar
% --- p. 23 ---
\setcounter{footnote}{0}%
\opage{23}jo et Skridt videre: den vil være en Lære om, hvilke etiske
Domme vi bør fælde! Hvorledes er en saadan Videnskab mulig?
Dertil kommer, at Domme om Godt og Ondt udspringe af Følelsen,
--- men „om Følelser kan man ikke disputere.“ Og Erfaring
viser jo ogsaa, at Diskussioner om etiske Spørgsmaal hyppigt
ende med, at Enhver beraaber sig paa sin Følelse.

Til Belysning af dette Problem maa først bemærkes, at
etiske Domme gaa ud paa at angive Midler og Veje til at opnaa
Noget, der har Værdi for os, og som vi derfor mere eller
mindre bevidst gøre til Formaal. Den etiske Vurdering betragter
den enkelte Handling, den enkelte Egenskab eller den enkelte
sociale Indretning i Forhold til et Formaal. Jo mere uvilkaarligt
Bedømmelsen og Vurderingen foregaar, desto mindre
Bevidsthed have vi om Formaalet; vi gribes umiddelbart og
hvile ganske i Tanken om Genstanden, der vurderes. Vurderin
gen kan foregaa ganske instinktmæssigt. Men naar Vurderingen
i Tvivlstilfælde skal grundes paa bestemte Principer, saa maa
det først afgøres, hvilket Formaal vi sætte Handlingen i Forhold
til. Dommen om Handlingen vil afhænge af, om denne ved nærmere
Undersøgelse viser sig at fremme eller hæmme det forudsatte
Formaal. Enhver Handling vil bevirke Noget; ellers
vilde den være umulig og absurd. Denne Bemærkning er ikke
unødvendig, ti store etiske Tænkere, især \textsc{Immanuel} \textsc{Kant}, have
ment, at det var en Nedværdigelse af Villien at lægge særlig
Vægt paa Virkningen og tilkende Formaalet en bestemmende
Indflydelse paa Vurderingen. Herimod gælder \textsc{Schleiermachers}
træffende Bemærkning\footnote[1]{Ueber den Begriff des höchsten Gutes. Werke zur Philosophie. 2 Band. p. 452.}: „Hvorfor handler jeg dog, naar jeg
Intet vil bevirke?“

Vi faa allerede fast Grund under Fødderne, naar det staar
fast, at enhver Handling maa have et Formaal. Ti da er det
muligt at undersøge, om Handlingen virkeligt tjener dette Formaal
eller ikke. Forholdet mellem Vejen og Endemaalet, mellem
Midlet og Formaalet kan videnskabeligt fastslaas. Staar Formaalet
fast, maa visse bestemte Midler anvendes. Forholdet
% --- p. 24 ---
\setcounter{footnote}{0}%
\opage{24}mellem Middel og Formaal er jo det samme som mellem Aarsag
og Virkning: hvad vi kalde Middel, er den Aarsag, der bevirker
Opnaaelsen af Formaalet. Her se vi da Muligheden af
en Videnskab om det, der bør være: Ordet bør udtrykker Fordringen
om Midler til et givet Formaal. I enhver Trang eller
Drift ligger en saadan Fordring, og hvad vi kalde Samvittighed,
er en dyb indre Trang til at hævde eller frembringe noget
Værdifuldt i Verden. --- Ved det nødvendige Forhold mellem
Formaal og Midler bliver en videnskabelig Etik mulig. Men
Formaalet forudsættes givet --- og, som det vil vise sig, Problemet
opstaar igen, naar der spørges, hvilket Formaal vi forudsætte
som det højeste, --- hvilken Værdi der for os er Grundværdien.

Til samme Resultat, som Betragtningen af Handlingen i
dens Forhold til Formaalet har ført os, komme vi ogsaa ved en
nærmere Betragtning af den Sætning, at Domme om Godt og
Ondt ere Følelsesytringer. Følelsesytringen forudsætter her, at
Anskuelsen af eller Forestillingen om en Handling har bevæget
os i vort Indre, har fremmet eller hæmmet en uvilkaarlig Stræben,
en bevidst eller ubevidst Interesse i os. Lyst- og Ulystfølelser
hænge --- som Følelseslivets Biologi viser --- sammen
med, at vort Liv eller dog en Del af vort Liv fremmes eller
hæmmes. Derved bestemmes, hvad der faar Værdi for os og
altsaa kan gøres til Formaal. Al Værdi af hænger af Livsforholdene
og af Livets Trin. Det menneskelige Følelsesliv er derfor kun
forstaaeligt i Sammenhæng med de menneskelige Livsforhold. Og
kun i sine simpleste Former og i de voldsomste Affekter er
Følelsen ganske blind. Sin bestemte, fuldt udviklede Form faar
den kun ved, at den forbindes med Forestillinger. Glæden forudsætter
Forestillingen om en lystvækkende Genstand, Sorgen
Forestillingen om Tabet af en saadan. Haab og Frygt forudsætte
Forestillinger om Muligheden af visse lyst- eller ulystvækkende
Begivenheder.

Derved faa vi ogsaa her fastere Grund under Fødderne.
„Om Følelser kan man ikke disputere“; men man kan disputere
om de Forestillinger, til hvilke Følelserne ere knyttede,
idet man kan undersøge disse Forestillingers Forhold til Erfa%
% --- p. 25 ---
\setcounter{footnote}{0}%
\opage{25}ringen. Og man kan undersøge selve de Tilstande, som vise sig
værdifulde, og søge at finde, under hvilke Betingelser de indtræde;
man vil da deraf kunne drage Slutninger angaaende
Maaden, hvorpaa saadanne værdifulde Tilstande kunne tilvejebringes i Fremtiden.

Altsaa ikke blot i Forholdet mellem Middel og Formaal,
der spiller en Rolle ved alle menneskelige Handlinger, men
ogsaa i Forholdet mellem Erfaringen og de Forestillinger, ved
hvilke de Følelser, der faa Luft i Vurderingen, bestemmes, se
vi en Mulighed for Begrundelse af etiske Domme og altsaa for
Opstilling af en filosofisk Etik. Og de to Forhold kunne ikke
skilles; ti et Formaal forudsætter stedse en Følelse. Kun hvad
der sætter vor Følelse i Bevægelse, kunne vi gøre til Formaal.

Etiske Domme angaa menneskelige Handlinger og menneskelige
Livsordninger (Institutioner), disse sidste dog kun, forsaavidt
de ere fremgaaede af og atter kunne ændres ved menneskelige
Handlinger. For det Meste er man nu i Etiken gaaet ud
fra Handlingen, der skal vurderes, og har ikke undersøgt selve
Vurderingen. Det vil vise sig, at man maa begynde med at
drøfte selve Vurderingens Betingelse.

2. Handlingens Udspring er i Individets Indre, i dets Instinkter
og Drifter, dets Tanker og Følelser. Men efter at have
udviklet sig her, træder den som Virkning over i den ydre Verden.
Her opstaar da det Spørgsmaal, til hvilken Del af hele
dette Forløb den etiske Vurdering især skal holde sig.

Til en fuldkommen Vurdering vilde udfordres, at Handlingens
hele Forløb droges med ind. Den etiske Karakteristik af
en Handling som god eller ond vilde da forudsætte, at man
kunde efterspore den lige fra dens første Spirer i Individets
indre Liv og gennem alle dens forgrenede og fjerne Virkninger
i Omverdenen. Etikerne have ikke lagt lige stor Vægt paa alle
Dele af Handlingens hele Historie. Nogle have taget deres Udgangspunkt
fra Handlingens eget Udspring, altsaa fra de Motiver,
det Sindelag, som afføder den. Handlingens ydre Virkninger,
mene de, ere uoverskuelige og staa ikke fuldstændigt under
Villiens Herredømme; det er den indre Handling, i hvilken
Villien lægger sig for Dagen, og derfor er det kun den, der kan
% --- p. 26 ---
\setcounter{footnote}{0}%
\opage{26}blive Genstand for etisk Vurdering. Andre gaa derimod ud fra
Handlingens Følger og Virkninger i den ydre Verden. Det er
kun disse, mene de, som give Handlingen praktisk Betydning.
Til dem maa Vurderingen først og fremmest holde sig, især da
Erfaring viser, at Handlinger, som udspringe af samme Motiver,
kunne have meget forskellige Virkninger, og Handlinger, som
udspringe af forskellige Motiver, kunne have samme Virkning.
Kun sekundært udstrækker Vurderingen sig til Motiverne og
Sindelaget, forsaavidt det kan godtgøres, at visse bestemte Motiver
gennemgaaende føre til en vis bestemt Handlemaade. Der
skelnes altsaa her mellem Vurderingen af Handlingerne og
Vurderingen af den handlende Person; hin er primær, denne
sekundær.

Denne Modsætning har spillet en stor Rolle i Etikens Historie.
Man kunde betegne den som en Modsætning mellem
subjektiv Etik og objektiv Etik. I Etikens nyere Historie repræsenteres
den subjektive Etik af Kant, den objektive Etik af
\textsc{Bentham}, og Debatten paa den filosofiske Etiks Omraade staar
væsentligt mellem de fra disse to Tænkere udgaaende Retninger.
En Afgørelse af Striden maa bero paa en nøjere Bestemmelse
af Forholdet mellem subjektiv og objektiv Etik. En saadan
nøjere Bestemmelse søger jeg i dette Kapitel at give. Jeg
vil især støtte mig til den Betragtning, at selv om man er enig
med Bentham i, at en objektiv Vurdering af Handlinger kun er
mulig, hvis man tager sit Udgangspunkt fra deres Virkninger
og Følger, saa er Vurdering dog altid en subjektiv Virksomhed,
og ethvert Vurderingsprincip maa forudsætte visse subjektive
Betingelser. Striden mellem Kants og Benthams Skole har altfor
udelukkende drejet sig om Forholdet mellem Handlingens
Motiver og dens Virkninger. Maaske vil Problemet stille sig
klarere og blive lettere at løse, naar man gaar længere tilbage
og spørger om selve Vurderingens Motiver, eller om, hvilke
psykologiske Forudsætninger selve den Omstændighed, at der tilstræbes
en almindelig og objektiv Vurdering af menneskelige Handlinger,
hviler paa. Ligesom den teoretiske Erkendelse ved Analyse
fører tilbage til visse Principer, som det er Erkendelsesteoriens
Sag at paavise, og som staa i nøje Sammenhæng med
% --- p. 27 ---
\setcounter{footnote}{0}%
\opage{27}den menneskelige Erkendelsesevnes Natur, saaledes maa ogsaa
den praktiske Vurdering vise sig at bygge paa Principer, der
have deres Grund i den menneskelige Aand, selv om Spørgsmaalet
her skulde være noget mere indviklet end paa det teoretiske
Omraade. Baade Erkendelsesteori og Etik føre os tilbage
til Subjektiviteten som sidste Grundlag. Opgaven er netop
at vise, hvorledes der trods det uundgaaelige Subjektive i Grundlaget
kan naas en objektiv Erkendelse og en objektiv Vurdering.
Her er altsaa en Parallel mellem Erkendelsesproblemet og det
etiske Problem, som man ikke altid har været opmærksom paa.

Ligesom Grundlaget (Vurderingsmotivet) er Etikens subjektive
Princip, saaledes er den Maalestok, efter hvilken Handlingerne
vurderes ud fra dette Grundlag, Etikens objektive Princip
og bestemmer Etikens Indhold (kfr. I, 1). Vurdering forudsætter
nemlig to Ting: et Motiv, som driver til at udøve den
vurderende Aandsvirksomhed, og et Kriterium, efter hvilket
Vurderingen udføres\footnote[1]{I mit Skrift Om Grundlaget for den humane Etik (1876) har jeg endnu ikke tilstrækkeligt skelnet mellem Begreberne Grundlag og Motiv, eller tydeligere: mellem Vurderingsmotiv og Handlingsmotiv, hvorimod Adskillelsen mellem Grundlag og Indhold allerede er gjort i det nævnte Skrift. --- I min Afhandling Filosofiske Problemer (Universitetsprogram. Novbr. 1902) har jeg indført Begrebet Grundværdi som Forudsætning for Vurderingsmotivet. Ethvert Motiv forudsætter Erfaring af en Værdi, og Vurderingsmotivet forudsætter Erfaring af en Værdi, der for os betinger alle andre Værdier.}.

For nærmere at godtgøre Betydningen af det her fremdragne
Synspunkt vil jeg søge at give en Fremstilling af, hvorledes
baade Etikens subjektive og objektive Princip (Grundlaget
og Indholdet) udspringe af den etiske Vurderings Natur og Udviklingsmaade.
Derved vil da det nærmere Forhold imellem
dem vinde i Klarhed. --- Men da Vurderingen af Handlinger og
Livsforhold ikke blot hos forskellige Individer og Folkeslag foregaar
paa Grundlag af forskellige Motiver og efter forskellige
Maalestokke, men tillige hos hvert enkelt Individ betinges ved
sammensatte Motiver, vil det være nødvendigt, at vi tænke os
Forholdene simplere, end de i Virkeligheden ere. Intet Menneske
% --- p. 28 ---
\setcounter{footnote}{0}%
\opage{28}forfølger blot et eneste Formaal; Livet er broget og mangfoldigt
i os som udenfor os, og de forskellige Formaal kæmpe om
Eneherredømmet indenfor det enkelte Individs Bevidsthed saavel
som indenfor Samfundet. De enkelte Øjeblikke saavel som
de enkelte Interesser kræve deres Ret; Individet føler sig som
en lille Verden og dog som Led af en større Verden, eller af
flere Verdener, der hver kræve deres Ret. Derfor er den menneskelige
Samvittighed, den Følelse, der bestemmer Vurderingen,
i Regelen et meget sammensat Produkt. Det vil da være
hensigtsmæssigt at undersøge enkelte Elementer af dette sammensatte
Hele hvert for sig. Den følgende Fremstilling skrider
frem fra de simpleste til de mest omfattende Vurderinger; den
er ikke en historisk Fremstilling, men ordner Standpunkterne
efter Omfanget af de Formaal, der blive bestemmende for Vurderingen.

3. Det er en Kendsgerning, at Menneskene vurdere deres
egne og Andres Handlinger og kalde dem gode eller onde efter
Udfaldet af denne Vurdering. Hvorledes er nu en saadan Vurdering
mulig?

Vi tænke os først det simpleste Tilfælde, nemlig hvor det
er det handlende Subjekt, der dømmer sin egen Handling, uden
at det har nogen Bevidsthed om andre Væsener, til hvilke der
kunde være at tage Hensyn. Det handlende og vurderende
Subjekt tænkes altsaa som en lille Verden for sig selv. Vi
foretage en Abstraktion for at faa et saa simpelt Exempel som
muligt.

Den første Forudsætning for, at en Handling kan vurderes,
er, at den erindres; den maa altsaa ikke være forsvunden for
Bevidstheden, naar den er gaaet over fra Beslutning til Virkelighed.
Billedet af Handlingen maa kunne kaldes frem igen som
Genstand for Betragtning. Men dette er ikke tilstrækkeligt. Et
saadant Erindringsbillede kunde i og for sig staa som noget
rent Ligegyldigt, der ikke satte Sindet i Bevægelse. Kun hvis
Handlingen paa en eller anden Maade, paa et eller andet Stadium
af sin Udvikling, har grebet ind i Individets hele Tilstand
og derved enten vakt Lyst eller Ulyst, vil Billedet af Handlingen
kunne vække Lyst eller Ulyst.

% --- p. 29 ---
\setcounter{footnote}{0}%
\opage{29}Her fremtræder i sin allersimpleste Form den vigtige, allerede
berørte Sandhed, at al Vurdering af Handlinger forudsætter
et Subjekt med Evne til at føle Lyst eller Ulyst. Vurderingen
forudsætter, at der stilles en Fordring til Handlingerne, som
disse i større eller mindre Grad kunne tilfredsstille. Men en
saadan Fordring staar aldeles umotiveret, naar Handlingen ikke
formaar at vække Lyst eller Ulyst. Dette er, som allerede bemærket,
kun et andet Udtryk for, at Vurdering forudsætter et
Formaal, efter hvilket Handlingen kan maales. Formaal kan kun
det blive, der har Værdi for os, og som ikke umiddelbart falder
i vor Lod, men kun kan opnaas gennem et Arbejde. Den umiddelbare
Værdi kundgør sig gennem Lystfølelse, medens der vil
opstaa en Ulystfølelse, naar Handlingen hindrer Virkeliggørelsen
af en umiddelbar Værdi.

I det simple Tilfælde, vi tænkte os, kan den Følelse, der
betinger Formaalet, og som altsaa skal tilfredsstilles ved Handlingen,
kun være Individets egen. Individet vil da stemple Handlingen
som god eller ond, alt efter den Maade, paa hvilken den
har grebet ind i dets Liv. Hvilken Karakter og Betydning, Vurderingen
nærmere faar, vil bero paa, om Følelserne af Lyst
og Ulyst kun ere bestemte ved og svare til de øjeblikkelige
Tilstande hver for sig, eller om de bestemmes ved Hensynet
til Individets Liv som Helhed og til de Vilkaar, under hvilke
det udvikler sig.

4. Jo lavere Bevidsthedslivet er, des mere isolerede og
selvstændige ere de enkelte Øjeblikke i Forhold til hverandre,
og des ringere Rolle spiller Erindringen og Tanken om Selvet
som en Helhed, der omfatter de enkelte Livsøjeblikke med alt
deres Indhold. Det er da kun et halvt ubevidst Instinkt, som
hindrer Individet i at gaa fuldstændigt op i det enkelte Øjeblik.
Selvopholdelsestrangen fører Individet til i det enkelte nærværende
Øjeblik ogsaa at tage Hensyn til Fremtiden og til at benytte
Erfaringerne fra Fortiden. Men jo mere Individet gaar op
i de enkelte Øjeblikke, des mindre Mulighed bliver der for en
Vurdering, da der ingen Sammenligning og Vexelvirkning mellem
de forskellige Tilstande kan finde Sted. Handlingen selv er
maaske glemt i det Øjeblik, da dens Virkning gør sig gældende
% --- p. 30 ---
\setcounter{footnote}{0}%
\opage{30}i Bevidstheden. Hvert Øjeblik udfyldes paa sin Maade, med sin
Følelsestilstand, som ikke faar nogen Indflydelse paa de andre
Øjeblikke. Det enkelte Øjeblik staar som absolut Egoist overfor
de andre Øjeblikke, vil ikke til Fordel for dem opgive sit
Krav paa Tilfredsstillelse.

Der viser sig her Muligheden af et Standpunkt, hvor al
Vurdering falder bort, fordi der vel rører sig Lyst- og Ulystfølelser,
men disse blot svare til Øjeblikkets Tilstand, ikke til
Livet som Helhed. Intet Formaal opstilles; selv Instinkterne
med deres ubevidste Formaal trænges tilbage; og naar ethvert
Hensyn til et Formaal falder bort, bliver Vurdering umulig.
Det er et barnligt Standpunkt; men det er muligt at forsøge
paa at gøre det til det definitive Standpunkt. Det er i Etikens
Historie forsøgt af Aristippos fra Kyrene\footnote[1]{Om Forholdet mellem Aristips Standpunkt og Søren Kierkegaards „æstetiske Stadium“ se min Bog S. Kierkegaard som Filosof. p. 84. Note.}. Det er Principet
om Øjeblikkets Suverænitet, han hævder. Dette er det mest
radikale etiske Standpunkt, som kan tænkes. Det har de færrest
mulige Forudsætninger, saa faa Forudsætninger, at enhver
Vurdering falder bort. Ti det Gode for Aristippos var den rent
momentane Lystfølelse. Hvorfor skal --- saaledes er Tankegangen
--- det ene Øjeblik ofres for eller underordnes det andet? Det
ene har i og for sig ligesaa megen Ret til at være som det
andet. --- Paa dette Standpunkt --- og kun paa dette --- falder
det Gode ganske sammen med Lystfølelsen, det Onde sammen
med Ulystfølelsen. --- Det følger af sig selv, at et saadant Standpunkt
ikke kan være et aldeles primitivt og naivt Standpunkt.
Det vil være sorgløst og tankeløst, som Barnet er det, --- men
netop den udtrykkelige Villen viser, at det kun er en Efterligning
af det barnlige Standpunkt. Der maa en vis Kunst og
Selvbeherskelse til for at hindre den uvilkaarlige Higen ud over
det nærværende Øjeblik. Allerede Instinkterne tilvejebringe en
Forbindelse mellem de enkelte Øjeblikke, som den principielle
Gennemførelse af Øjebliksstandpunktet ikke kan taale. Heri
ligger netop en stor Vanskelighed for den praktiske Hævden af
% --- p. 31 ---
\setcounter{footnote}{0}%
\opage{31}dette Standpunkt. Men det har sin store Betydning for os at
se, at dets Berettigelse teoretisk er bleven hævdet. Det kunde
efter Sagens Natur kun ske indenfor en fremskreden Kultur.
De simpleste Standpunkter ere ingenlunde altid de, der tidligst
fremtræde i Historien. Men der er en naturlig Tendens til at
opsøge Grænsetilfælde, til at se, hvorledes vore Begreber formes,
naar vi tage Livets Yderpunkter for os; og de Forsøg paa
Øjebliksetik, der foreligge, ere meget lærerige. Maaske kunde
der gives endnu fuldkomnere Fremstillinger af dette Standpunkt,
end der historisk foreligger.

Et saadant Standpunkt mangler ikke Berettigelse. Man
kunde indvende, at det ophæver al Etik, da det udelukker Vurdering,
og da al etisk Bestræbelse forudsætter, at et mere begrænset
Hensyn underordnes et mere omfattende Hensyn. En
saadan Underordning bliver her ikke Tale om, da Livet tænkes
at bestaa af absolut suveræne Øjeblikke! --- Men dertil maa
svares, at Etiken selv maa godtgøre sin Berettigelse til at opstille
Fordringen om Opgivelse af en Tilfredsstillelse i det ene
Øjeblik til Fordel for andre Øjeblikke. Den Mulighed frembyder
sig, at hvert Øjeblik kunde udfyldes af en umiddelbar Værdi
af en og samme Grad, eller at der i hvert Tilfælde var en Harmoni
mellem de Værdier, der opfylde de enkelte Øjeblikke. Bevisbyrden
paahviler den, der fordrer Opofrelse og Resignation.
Dette kan kun negtes af en absolut asketisk Anskuelse, d. v. s.
en Anskuelse, for hvem Askesen er Formaal, ikke blot Middel.
Hvert Øjeblik har en naturlig Ret til at være og har tillige,
om man saa kan sige, sit Selvopholdelsesinstinkt, idet Trangen
til fuld Tilfredsstillelse i Øjeblikket kun under en vis Modstand
giver efter for andre Tilskyndelser. Selv paa et rent instinktmæssigt
Livstrin mærkes denne Modstand.

5. Hvis Principet om Øjeblikkets Suverænitet praktisk
kunde gennemføres, vilde intet Ræsonnement kunne omstyrte
det. Men der gives vel neppe noget bevidst Individ, hos hvem
der ikke rører sig Instinkter og Drifter, der føre ud over Øjeblikket.
Og hos menneskelige Individer vil Erindring og Forventning
stedse gøre sig gældende og blive Tilknytningspunkter
for Følelser af Lyst og Ulyst, der bestemmes ved Individets
% --- p. 32 ---
\setcounter{footnote}{0}%
\opage{32}varige eller stadigt genkommende Vilkaar. Ved enhver Besindelse
over sig selv og sin Handlen vil Individet hæve sig over de
enkelte Øjeblikke i deres Forskellighed og Isolerethed, og dets
Følelse vil (i Besindelsens Øjeblikke idetmindste) svare til den
Maade, hvorpaa Livshelheden, og ikke det enkelte Øjeblik, staar
for Bevidstheden. Der danner sig først et realt Selv, naar der
er en Kreds af Følelser og Forestillinger, som danne en fast
Kærne i Bevidstheden, selv om de ikke i hvert Øjeblik gøre
sig gældende. Denne faste Kærne udgøres fremfor Alt af de
Værdier, Individet kender, de Formaal, det sætter sig. Disse
Værdier og Formaal beherske dets hele Liv, udgøre den Interesse,
som samler Livet til en Helhed og lader det staa
som Mere end en Række af isolerede Øjeblikke. Enhver stærk
Følelse har en Tendens til at brede sig og beherske hele Bevidsthedslivet,
til at farve Alt, hvad der udgør Bevidsthedens
Indhold. Ved denne Følelsens Expansion kan ikke blot en øjeblikkeligt
vakt Stemning strække sine Virkninger ud over flere
Øjeblikke, men der kan ogsaa danne sig et realt Selv, en herskende
Grundstemning, der betegner det individuelle Følelseslivs
Niveau i Modsætning til de øjeblikkelige Svingninger.

Hvor nu det enkelte Øjebliks Følelsestilstand, betragtet som
Virkning af Individets egen Handling, i Bevidstheden støder sammen
med den ved Forestillingen om Livshelheden bestemte Følelsestilstand,
vil der opstaa en ny Følelse, som er bestemt ved
Forholdet mellem dem, et Forhold, som enten kan være harmonisk
eller disharmonisk. I denne Følelse ved Forholdet mellem
den momentane Tilstand og den ved Hensynet til Livets
Helhed bestemte Tilstand bestaar Vurderingen. Den udtrykker
Forholdet mellem to Værdier. Evnen til saadan Vurdering er
Samvittigheden, saaledes som den kunde tænkes at ytre sig hos
et aldeles isoleret Individ. I videste Betydning er Samvittighed
en Forholdsfølelse og forudsætter kun, at der i Bevidstheden gør
sig et Forhold gældende mellem et Centralt og et Periferisk,
et mere og et mindre Omfattende. Det enkelte Øjeblik og den
enkelte Handling, der har givet det dets Beskaffenhed, vurderes
--- paa det Standpunkt, som her er forudsat --- efter den Maade,
paa hvilken de kunne føjes ind som Led i det individuelle Livs
% --- p. 33 ---
\setcounter{footnote}{0}%
\opage{33}Helhed. Herved virker ikke blot det reale Selvs og de herskende
Formaals Tendens til at bestemme Alt i Bevidstheden, men ogsaa
en naturlig, i Bevidsthedens almindelige Væsen dybt grundet
Trang til Enhed og Sammenhæng i de sjælelige Tilstande. Vi
ønske, at vort Liv skal være en Helhed, ikke blot fordi kun
derved vor højeste Livsinteresse kan raade, men ogsaa fordi
Enhed og Sammenhæng i og for sig ere Goder. Det er Personlighedens
formale Side, der ytrer sig i Trangen til at ville Et,
saa at alt Andet, som vi attraa foruden dette Ene, kun er Middel
for det. Det er en aandelig Selvopholdelsestrang, som aldrig
ganske kan mangle, og saa sjeldent man i Virkeligheden finder
fuld Konsekvens, saa umuligt er det, at en Bestræbelse efter
at være konsekvent nogensinde ganske kan mangle\footnote[1]{Smlgn. med Hensyn til de her anvendte psykologiske Ideer og Love min Psykologi V B, 5; VI E, 1 og F, 4.}.

Der stiller sig her den Opgave for Individet at bringe Harmoni
tilveje mellem de enkelte Dele af sit Liv. Det er en Opgave,
som neppe hos noget Menneske løses uvilkaarligt, uden
nogen bevidst Stræben. Her faar derfor ogsaa Vurderingen af
de tidligere Handlinger efter den Maade, hvorpaa de have bidraget
til Løsningen af denne Opgave, Betydning for Individet.
Vurderingen muliggøres altsaa ikke blot ved den centrale, til
Livshelheden svarende Følelse, men motiveres ogsaa ved den.
En fin og udviklet Sans for, hvad der tjener det individuelle
Liv, hvis enkelte Led Øjeblikkene ere, er en Betingelse for
dets Bestaaen og Udvikling. Den er en Art højere Selvopholdelsesinstinkt
og behøver ikke at indskrænke sig til blot at angaa,
hvad det fysiske Livs Bestaaen udkræver, men kan omfatte
de ideelle Fornødenheder ikke mindre end de fysiske. --- I
\textsc{Platons} og \textsc{Aristoteles}'s Etik har denne harmoniske Individualisme
faaet sit mest karakteristiske Udtryk. Den etiske Dyd
betragtes af dem som en aandelig Sundhed og Harmoni. Især
er Aristoteles's Definition af Dyden (som den ved ethvert enkelt
Individs Natur bestemte Midte) af Interesse i denne Henseende.
Dog er de nævnte Tænkeres hele Etik ikke udtømt
med denne individualistiske Lære.

% --- p. 34 ---
\setcounter{footnote}{0}%
\opage{34}Jeg bruger til Betegnelse af dette Standpunkt Ordet „Individualisme“
og ikke „Egoisme“, da der ved dette sidste Ord
helst maa tænkes paa en bevidst Tilsidesættelse og Underordnelse
af Andres Vel for Ens eget. Individualisten behøver ikke at
være Egoist, men han kan blive det.

6. Paa et saadant individualistisk Standpunkt vil Opgaven
ikke blot være den at bestemme, hvor megen Energi der maa
forbruges i de enkelte Øjeblikke, men ogsaa den at anvende
den Energi, der raades over, paa saa afvexlende og forskelligartet
en Maade, som foreneligt er med Livshelhedens Interesse.
Ifølge en psykologisk Naturlov betinges Følelsernes Liv og Friskhed
ved, at de forskellige Tilstande træde i et vist Modsætningsforhold
til hverandre, medens Ensformighed og Gentagelse virke
dæmpende eller sløvende. --- Desuden vil ogsaa Livets naturlige
Gang føre med sig, at Individet efterhaanden ikke blot gaar op
i den fysiske Selvopholdelse, men faar Fornødenheder af mere
ideel og mere sammensat Beskaffenhed. Individets Liv vil naa
des større Fylde, jo flere forskellige Retninger det kan brede
sig i, jo rigere og mangfoldigere et Indhold det kan omfatte,
uden at Enheden og den samlede Kraft derved svækkes.

Da Individet her tænkes at have sat sig Formaal, der
pege ud over Øjeblikket og ud over Tilfredsstillelsen af en enkelt
Side i dets Natur, kan der blive Tale om en Fordring, en
Skullen, en Lov, og de enkelte Øjeblikke og de enkelte Handlinger
vurderes efter, hvorvidt de tjene det opstillede Formaal.
Det kunde f. Ex. være intellektuel eller æstetisk Nydelse og
Udvikling, der sættes som Formaal, og dermed vil da den Fordring
være given, der bliver at stille til de enkelte Øjeblikke
og Handlinger.

Den etiske Lov paa Individualismens Standpunkt findes ved
Formulering af, hvad det harmoniske Forhold mellem Livshelhedens
Interesse og de enkelte Øjeblikkes Trang kræver. Den
vil bestaa af to Hovedbud, et negativt og et positivt: 1) det
enkelte Øjeblik maa ikke have større Selvstændighed, end der
svarer til dets Betydning indenfor Livshelheden; 2) men paa
den anden Side skal der i det enkelte Øjeblik leves saa rigt
% --- p. 35 ---
\setcounter{footnote}{0}%
\opage{35}og fyldigt, som foreneligt er med Livshelhedens Bevarelse og
med Opnaaelsen af det Formaal, der staar som det Højeste.

God bliver altsaa her den Handling, som bevarer Livshelheden
og giver Livsindholdet Fylde og Liv, ond bliver den, der
har en mere eller mindre udpræget Tendens til at sprænge eller
indsnevre Livshelheden og dens Indhold. Ondt er saaledes det
enkelte Øjeblik og den enkelte Drift i deres oprørske Isolation
fra det øvrige Liv; og Ondet vil ligge saa meget des dybere
og blive Genstand for en saa meget des stærkere Forkastelse i
den vurderende Følelse, jo mere det er villet, d. v. s. fremgaaet
som Frugt af Overvejelse og Valg, ikke blot af øjeblikkelig Tilskyndelse.
Den Styrke, hvormed Vurderingen gør sig gældende,
vil bero paa den Styrke, hvormed Livshelhedens Interesser gøre
sig gældende i Individets centrale Følelser.

7. Det gælder om Individualismen eller Principet om Individets
Suverænitet, hvad der gjaldt om Øjeblikkets Suverænitet,
at den ikke kan omstyrtes ved Ræsonnement. En absolut
Individualist eller Egoist vilde være absolut uangribelig. Naar
Hævdelsen af hans eget Liv, dets Bestaaen, dets Fylde, er det
eneste Formaal, han erkender, saa gives der ingen logisk
Overgang fra dette Standpunkt til et andet. Med tilstrækkelig
Konsekvens og Selvbeherskelse og med tilstrækkelig Magt over
Forholdene vil da et Individ kunne gennemføre den Grundsætning:
Du skal handle saaledes, at du selv stedse staar som
eneste Formaal, som absolut Midtpunkt i Verden. --- Skal en
Forandring ske, maa den Følelse, som bestemmer Grundværdien
og dermed Vurderingen, ændres ved at knyttes til en mere omfattende
Kreds af Forestillinger end den, der blot angaar Individets
eget Liv. Før det er sket, kan det heller ikke nytte at appellere
til Samvittigheden; ti denne er, som vi have set, en Forholdsfølelse,
Udtryk for Forholdet mellem det Centrale og det Periferiske
i Individets Følelsesliv, og bestemmes altsaa selv ved, hvad det
Centrale er.

Den filosofiske Etik har ofte, især i tidligere Tid, ment at
maatte gøre Fordring paa at være en ren Fornuftvidenskab, der
ikke appellerede til andre Forudsætninger end saadanne, som
ligge i den menneskelige Fornufts Væsen. Men dette strider
% --- p. 36 ---
\setcounter{footnote}{0}%
\opage{36}imod Etikens Karakter som praktisk Videnskab. Handling kan
kun vurderes efter Formaal, og Fastsættelsen af et Formaal
forudsætter Lyst- og Ulystfølelse hos det fastsættende Subjekt.
Afset fra en Bevidsthed med Evne til at føle Lyst eller Ulyst
har derfor etisk Vurdering ingen Betydning. Men paa den anden
Side ligger der i den blotte Evne til at føle Lyst og Ulyst endnu
ikke Noget om, hvor omfattende den Kreds af Forestillinger er,
til hvilken Lyst- og Ulystfølelsen knyttes.

Fra Øjebliksstandpunktet til det individualistiske Standpunkt
skete Overgangen ved, at der i Bevidstheden dannede sig Følelser,
som bestemtes ved Livshelhedens Interesser, og som
derfor kunde blive stærke Bindeled mellem de skiftende Øjeblikke.
Støttet til dem føler Individet sin Enhed trods den stadige
Skiften. Dersom der gives et videre Standpunkt end
Individualismens, maa der gives Følelser, der ved at røre sig
hos det enkelte Individ kunne knytte ham til en mere omfattende
Helhed paa lignende Maade, som de enkelte Øjeblikke
og Tilskyndelser i ham ere knyttede til den individuelle Livshelhed.
Der maa være en Magt, som forbinder de enkelte Individer
indbyrdes og ophæver Isolationen imellem dem.

8. Kun tilnærmelsesvis kan i Praxis Individualismen gennemføres.
Individets Suverænitet, Opfattelsen af Individet som
en afsluttet og aldeles selvstændig Helhed, viser sig at bero
paa en voldsom og kunstig Abstraktion. Individet bliver til ud
af Slægten og lever hele sit Liv som en Del af Slægtens Liv,
med en Organisation, i hvilken det arver Følger af tidligere
Generationers Handlen og Liden, under Livsvilkaar og i en
aandelig Atmosfære, som Slægtens Udvikling har bragt til Veje.
Og ligesom Selvopholdelsesinstinktet ophæver Isolationen af de
enkelte Øjeblikke i Livet og derved bliver Grundlaget for Følelser,
der bestemmes ved Livshelhedens Interesser, saaledes
rører der sig i de sympatiske Instinkter Kræfter, som ophæve
Isolationen af de enkelte Individer og hævde Slægtens Livsvilkaar
i disses Indre. I den mest primitive Form fremtræder
de sympatiske Instinkter ved Familieforholdets Grundlæggelse.
Saa meget end Familieforholdets Former og Indretninger variere
hos forskellige Folkeslag og til forskellige Tider, og saa løst og
% --- p. 37 ---
\setcounter{footnote}{0}%
\opage{37}vilkaarligt særligt Forholdet mellem de to Køn ofte viser sig,
saa er der dog et Forhold, som ifølge Sagens Natur ikke kan
ophæves eller væsentligt forandres, nemlig Forholdet mellem
Moder og Barn. I dette Forhold voxer den sympatiske Følelse
umiddelbart ud af Naturinstinktet. Det danner den faste Kærne,
ved hvilken Familielivets højere Former blive mulige. Der
lægges her Grunden til et Livsfællesskab, indenfor hvilket Sympatifølelserne
kunne plejes og naa en saadan Styrke, at de efterhaanden
kunne omfatte større Kredse. Moderkærligheden vedbliver
dog stedse at staa som Forbillede og Maalestok for al
Sympati, baade hvad Styrke og Renhed angaar. Et Minde om,
at Familieforholdet er de sympatiske Følelsers stadige Kilde, er
det, naar den almindelige Menneskekærlighed finder sit mest
slaaende Udtryk i den Sætning, at alle Mennesker ere Brødre.

Det er her ikke Stedet at gaa ind paa en nærmere psykologisk
Undersøgelse af de sympatiske Følelsers Natur og af
deres forskellige Karakter efter de Elementer, hvoraf de bestaa,
og efter det Omfang, i hvilket de ytre sig (smlgn. herom min
Psykologi VI C)\footnote[1]{Jeg bruger, som det ikke blot ses af Texten her, men ogsaa af min Psykologi (VI C), Ordet Sympati i Betydning af en Medfølelse, i hvilken man umiddelbart gør Andres Lyst eller Ulyst til sin egen, eller, som man ogsaa kan udtrykke det, sætter sig paa Andres Stade og føler Livet paa den Maade, de maa føle det. I denne Betydning er Ordet brugt ikke blot i den ældre danske Psykologi (Treschow: Om den menneskelige Natur. København 1812. p. 355. F. C. Sibbern: Psykologisk Pathologie. København 1828. p. 27), men ogsaa i senere Værker over Følelsens Psykologi (Domrich: Die psychischen Zustånde. Jena 1849. p. 218. Bain: Emotions and Will. 3. ed. London 1875. p. 111) og i etiske Værker (Hume, A. Smith og i vore Dage Wundt: Ethik\textsuperscript{1} I. p. 208). --- Afset fra denne Betydning bruges Ordet Sympati i tre andre Betydninger. --- Det kan betegne en passiv,\textbackslash{} reflexagtig Genspejling af Andres Tilstande. Synet af Andres Nydelse eller Lidelse lader os selv uvilkaarligt stemmes paa lignende Maade, saa vi maaske endog f. Ex. føle Smerte i vort eget Ben, naar vi se en Andens Ben lemlæstes. Lidelsen virker her i Regelen stærkere end Nydelsen. --- Det kan dernæst betegne en Evne til at forestille os Andres Følelser, hvad enten vi dele dem eller ikke. Her spille Takt og Slutning en større Rolle end umiddelbar Enhedsfølelse. Den, vi paa denne Maade sympatisere med, vedbliver dog at staa}. Hvor Sympatien har naat sin fulde Renhed,
% --- p. 38 ---
\setcounter{footnote}{0}%
\opage{38}er den en Følelse af Lyst eller Ulyst, som blot er bestemt ved,
at andre Væsener føle Lyst eller Ulyst. Hvad dens Omfang angaar,
var det det betydningsfuldeste Punkt i dens Udviklingshistorie,
da den udvidede sig fra blot at omfatte Familie, Nation
og Race til at gælde hele Slægten. Allerede den græske Filosofi
(den peripatetiske og den stoiske Skole) førte til Ideen om
en almindelig Menneskekærlighed, grundet paa alle Menneskers
naturlige Sammenhøren iet stort Samfund\footnotemark[1]. Men større historisk
Betydning fik denne Ide først ved, som et Hovedbud i
Kristendommen, at blive en af de ledende Tanker i en stor
Verdensreligion.

9. Etisk set er den store Betydning heraf, at Horizonten
som Objekt. --- Endeligt kan det betyde en Følelse af Velbehag ved at
genfinde hos Andre Følelser af Lyst eller Ulyst, vi selv i lignende Tilfælde
vilde have. Det er en Styrkelse, en Selvbekræftelse, vi derved erfare.
--- Det er ørkesløst at disputere om, hvilken af alle disse Betydninger
der er den „rigtige“. --- Sympati i de tre sidste Betydninger af
Ordet kan træde i Modsætning til Sympati i den første Betydning af Ordet.
Men de forskellige Arter af Sympati ere beslægtede, og navnligt ville de
tre sidste Arter (den reflektoriske, den intellektuelle og den interesserede
Sympati) kunne bane Vejen for den første gennem psykiske Overgange
og Forskydninger.

\footnotemark[1] Cicero: De finibus V, 65 kfr. De officiis I, 16—17; III, 6. ---
Smlgn. min Af handling: Hedenske Sandhedssøgere (Mindre Arbejder I. p.
159 f.; 167—174). --- Flere teologiske Kritikere have klaget over, at det
ikke vises, hvorfra de sympatiske Følelser opstaa. Af min Psykologi vilde
de have kunnet se, at jeg antager et fysiologisk Udgangspunkt og en
psykologisk-sociologisk Udviklingsgang af disse Følelser. Med Hensyn til
Grad, Art og Omfang ere de underkastede bestemte Betingelser, som jeg
i Korthed har angivet. Det er mærkværdigt, hvor hyppig den Mening er,
at Egenkærligheden er den eneste naturlige Følelse. Denne Mening findes
især paa to modsatte Standpunkter: paa den ene Side hos Skeptikere og
blaserede Verdensmænd, der herved beraabe sig paa deres formentlige
Erfaringer, og paa den anden Side hos ortodoxe Teologer, efter hvem al
sand Kærlighed kommer fra en overnaturlig Kilde. Ja, selv en oplyst og
kritisk Teolog som Adolf Harnack negter, at Menneskekærligheden kan
være «Naturprodukt“, og mener, at kun Egenkærligheden er naturlig.
„Den til sig selv overladte Oplysning“ fører efter hans Mening til Egoisme
og „almindelig Opløsning“. Se hans Foredrag paa den evangelisk sociale
Kongres i Frankfurt a. M. 17. Maj 1894.
% CHECK p. 38: note mark without a note

% --- p. 39 ---
\setcounter{footnote}{0}%
\opage{39}er bleven udvidet, saa at Følelsen af Lyst og Ulyst nu ikke
blot bestemmes ved Individets egen Skæbne, men ved Livsvilkaarene
for et Samfund, af hvilket Individet kun er et enkelt
Led. Naar nu saadanne Følelser blive raadende, ville Handlingerne
vurderes efter det Forhold, i hvilket de ved dem frembragte
Tilstande staa til dette Samfunds, den større Helheds
Interesser. Maalestokken hentes nu ikke blot fra Individets
isolerede Liv; dette staar som hørende til en større Verdensorden,
og det kommer an paa, om Handlingerne hæmme eller
fremme denne.

Det vilde ikke være rigtigt at sige, at Sympatien paa dette
Standpunkt bliver Et med den etiske Følelse eller Samvittigheden.
Denne fremtræder ogsaa her som en Forholdsfølelse,
bestemt ved Forholdet mellem, hvad der er den raadende eller
centrale Følelse hos Individet, og Handlingernes Resultater. Forskellen
mellem dette Standpunkt og det individualistiske er den,
at Grundlaget er mere omfattende; Vurderingen vil derfor kunne
falde anderledes ud: en og samme Handling vil fra det individualistiske
Standpunkt kunne kaldes god, fra Slægtens ond.
Naar Individet umiddelbart føler sine egne personlige Interesser
som underordnede Hensynet til den store Helheds Tarv, af
hvilken det i Sympatien føler sig som et enkelt Led, fremtræder
den etiske Følelse som Pligtfølelse. Allerede paa Individualismens
Standpunkt vilde der kunne tales om Pligt. Ti i Begrebet Pligt
behøver der kun at ligge det formelle Forhold mellem et mere
begrænset, og et mere omfattende Hensyn, at dette sidste skal
have Forrang fremfor hint. Individet kan i det enkelte Øjeblik
føle sig forpligtet ved Hensynet til sin egen Livshelhed. Pligtfølelsen
udspringer --- hvor Pligten er Mere end noget udefra
Paatvunget --- af Personlighedens Indre og betinges ved Fordringen
om og Trangen til at bevare Sammenhængen i sit Liv
gennem Hævdelse af de højeste Formaal og Livsinteresser i de
enkelte Øjeblikke og under de specielle Forhold, --- til at genfinde
sig selv efter sit centrale Væsen i den enkelte Tilstand
eller Handling. I Pligten ytrer sig da en Grundvillie: vi skulle,
fordi vi i vor inderste Grund ville. Vor centrale Villen bliver
en Skullen, fordi der kan rejse sig periferiske Villiestendenser
% --- p. 40 ---
\setcounter{footnote}{0}%
\opage{40}i Modsætning til den, og fordi vi dog ville forblive „tro mod
os selv og mod vor bedste Stræben“. Men paa det Standpunkt,
som vi nu beskrive, vil selve Hensynet til Individets
Livshelhed igen være underordnet Hensynet til Slægtens Livshelhed,
og ligeledes vil den Sympati, som har et snævrere
Samfund til Genstand, være underordnet den, der er rettet mod
et mere omfattende Samfund. Eller rettere: Individet har her
gjort Samfundets eller Slægtens Livsinteresse til sin egen højeste
Livsinteresse, Genstand for sin egen centrale Villen.

Fra en anden Side set fremtræder den etiske Følelse under
sin videre Udvikling som Retfærdighedsfølelse. Den Sympati,
som ligger til Grund, vil nemlig, naar den er Mere end et
blindt Instinkt, i sin Ytringsmaade (baade hvad Art og hvad
Grad angaar) ledes af Hensynet til de Væseners Ejendommelighed,
hvilke den gælder. Og naar dens Omfang udvides til
alle Væsener, der have Evne til at føle Lyst og Ulyst, maa
den i hvert enkelt Tilfælde ytre sig saaledes, at det enkelte
Væsen, den er knyttet til, fyldestgøres efter sin Ejendommelighed,
uden at derved andre Væsener krænkes i deres ligesaa
udprægede Ejendommeligheder. Enhver Forskel og Ulighed, som
gøres gældende, maa være begrundet ved, hvad det Riges Tarv,
til hvilket baade Individet selv og de andre Væsener høre,
fordrer. Sympati er, fra den aktive Side betragtet, Trang til
Meddelelse; denne Meddelelse maa, naar den ikke skal ske i
Blinde, bestemmes ved Principer, og disse maa hentes fra selve
Sympatiens Natur. Naar Sympatien er universel, kan den Forskel,
den gør i Fordelingen af sine Goder, kun skyldes den
Grund, at disse Goder ved at fordeles paa anden Maade ikke
virkeligt eller ikke i saa høj Grad vilde være Goder for dem,
de tilfaldt, eller ikke vilde blive til saa stor Fremgang for alle
Livsinteresser. Paa Sympatiens Grund udvikler den etiske Følelse
sig saaledes til en Følelse for fordelende Retfærdighed.
Ligesom Individet føler sig som En blandt Mange, saaledes betragter
det ogsaa ethvert andet Individ som En blandt Mange,
forsaavidt ikke særegne Grunde motivere særegne Hensyn. ---
Paa Individualismens Standpunkt vilde der kunne findes en Analogi
til den fordelende Retfærdighed, idet Forholdet mellem de
% --- p. 41 ---
\setcounter{footnote}{0}%
\opage{41}forskellige Øjeblikke og Tilskyndelser svarer til Forholdet mellem
de forskellige individuelle Væsener paa det Standpunkt, vi her
tale om. \textsc{Platon} har netop brugt Udtrykket Retfærdighed om
det harmoniske Forhold mellem de forskellige Sider eller Dele
af det enkelte Individs Bevidsthedsliv.

Det er denne logiske Karakter ved den etiske Følelse, den
Strenghed, hvormed den fordrer Begrundelse for enhver Forskel
og Ulighed, som statueres, der især har gjort, at man har ment,
at den etiske Lov er et Udtryk for den rene Fornuft. Hvad
\textsc{Kant} formulerede som det kategoriske Imperativs Indhold, udtrykker
i Virkeligheden kun den Upartiskhed, den Bortseen fra
uvedkommende og tilfældige Hensyn, som den paa Sympatien
byggede etiske Følelse fordrer. Kant fordrede, at man ved en
Handlings Vurdering skulde stille sig paa et universelt Standpunkt,
og han mente, at ethvert andet Standpunkt indeholdt en
Modsigelse. Men for at en saadan Modsigelse skal komme frem,
maa der være aldeles bestemte psykologiske Betingelser tilstede;
der maa mærkes et Sammenstød mellem to Interesser. Det
gælder atter her, at Samvittigheden forudsætter et Forhold, i
hvilket intet af Leddene er forsvindende lille. Hin Fordring er
altsaa ingenlunde selvfølgelig. Den er jo ogsaa først kommen
frem paa et bestemt Trin af den historiske Udvikling. Hos Kant
selv er Ideen om denne Fordring bleven til under Indflydelse
af en Undersøgelse af det menneskelige Samfundslivs Udvikling\footnote[1]{Smlgn. mit Værk: Den nyere Filosofis Historie'\textsuperscript{2}, II p. 72 —76 og min Afhandling: Rousseaus Indflydelse paa den definitive Form for Kants Etik (Oversigt over det kgl. danske Videnskabernes Selskabs Forhandlinger 1896).},
en Undersøgelse, ved hvilken han dog netop oversaa de sympatiske
Følelsers Udvikling under Samfundslivets Indflydelse, hvorfor
ogsaa det kategoriske Imperativ kom til at staa som en mystisk
Magt i Menneskenaturen. Kant betragtede al Følelse af Lyst
og Ulyst (ogsaa sympatisk Følelse) som egoistisk og afskar sig
derved Muligheden for at forstaa, hvorledes der hos Mennesket
kan udvikle sig Interesser og Formaal, som ligge ud over dets
egen Bestaaen og Nyden. Den etiske Lov kommer derfor hos
% --- p. 42 ---
\setcounter{footnote}{0}%
\opage{42}Kant paa hemmelighedsfuld Maade, som en Kundgørelse fra en
anden Verden, ind i Menneskets erfaringsmæssigt givne Natur.

10. En etisk Lov bliver til, naar en mere omfattende
Helheds Livsbetingelser formuleres i bestemte Tanker. Paa det
Standpunkt, vi her have for os, Humanitetsetikens Standpunkt,
kan dens Indhold ikke være andet end den Grundsætning, at
Handlingerne skulle føre til saa stor Velfærd og Fremgang for
saa mange bevidste Væsener som muligt. Og heri indbefattes igen
to Hovedbud, et negativt og et positivt: 1) intet enkelt Individ
bør der tildeles Mere, end der tilkommer det efter den Plads,
det ifølge sin Ejendommelighed indtager indenfor Slægten; 2)
men paa den anden Side skal ethvert Individs Evner og Drifter
udvikles og tilfredsstilles saa fuldt og rigt, som foreneligt er
med, hvad Slægtens Liv som Helhed kræver. Disse to Bud
følge med logisk Nødvendighed af Begrebet om Samfundet som
en til Enhed forbunden Mangfoldighed af bevidste Væsener.
Mod Samfundets Enhed strider det, at et enkelt Individ eller
enkelte Individer paa vilkaarlig Vis foretrækkes eller tilsidesættes
for Andre: enhver Særstilling maa begrundes ved, hvad
de fælles Livsvilkaar kræve; men paa den anden Side er et
Samfund des fuldkomnere, jo friere og selvstændigere de enkelte
Led røre sig, jo flere forskellige Muligheder de virkeliggøre,
samtidigt med, at Enheden bevares og stedse faar en
inderligere Karakter og en højere Gyldighed.

Naar den etiske Følelse paa Sympatiens Grundlag udvikler
sig til Pligt- og Retfærdighedsfølelse, vil det være det i denne
Lov udtalte Princip, som tilsidst er Maalestokken for de etiske
Domme, der fældes\footnote[1]{Medens jeg her er gaaet fra Begrebet Velfærd til Begrebet Retfærdighed, gaar jeg den omvendte Vej i et Forelæsningsgrundrids fra 1901, der under Navn af „Human Etik“ er optrykt i Mindre Arbejder, Tredie Række. Jeg opstiller der den Hypotese, at Retfærdighed er den ypperste etiske Karakteregenskab („Dyd“), og undersøger --- efterat have givet en Beskrivelse af den --- derefter, hvorvidt og hvorledes en saadan Antagelse kan begrundes. (1913).}. Paa dette Standpunkt følger det tillige
af sig selv, at de etiske Domme ikke blot angaa Individets
egne Handlinger, men ogsaa andre Individers, og at Maale%
% --- p. 43 ---
\setcounter{footnote}{0}%
\opage{43}stokken maa være den samme. Kun vil Individet overfor sine
egne Handlinger bedre være i Stand til at følge Handlingen
tilbage til dens Udspring og vil altsaa kunne vurdere større
Strækninger af dens Historie, end naar det staar overfor andre
Individers Handlinger.

God vil altsaa en Handling være, naar den bevarer og udvikler
bevidste Væseners Velfærd, ond, naar den gaar ud paa
det Modsatte. Det Onde vil her, ligesom paa Individualismens
Standpunkt, være det Opløsende og Isolerende. Naar et enkelt
Individ enten selv gør sig eller af Andre gøres til absolut Formaal,
opløses de bevidste Væseners Samfund. Det Onde er
derfor Egoismen i dens forskellige Grader og under dens forskellige
Former. Og Dommen over den bliver des strengere,
jo mere bevidst den er; ti des dybere og fastere Rod har Handlingen
i hele Sindelaget, og des vanskeligere kan der raades
Bod paa Isolationen og Opløsningen. Selve den strenge Vurdering
finder sin Motivering og sin Begrænsning ved Ønsket om en
saadan Genoprettelse. Den er derfor Mere end en æstetisk Dom,
den har et bestemt Formaal, --- det samme Formaal som det,
efter hvilket Handlingen vurderes.

11. Det objektive Princip, Principet for Bestemmelsen af
Etikens Indhold og for Vurderingen af de menneskelige Handlinger,
bliver altsaa her Principet om den almindelige Velfærd.
Ifølge det har ingen Handling, og ingen ved Handling grundlagt
Institution eller Livsordning, Værdi, uden forsaavidt den
er til Fremme for bevidste Væseners Liv og Lykke. Der er
mange andre Ting, der virke til Fremme eller til Skade herfor,
end menneskelige Handlinger. Allerede ved det ubevidste Naturlivs
Gang virkes der for eller mod bevidste Væseners Velfærd;
men Vurderingen af saadanne Virkninger har ikke nogen etisk
Karakter, fordi de ikke udspringe af nogen Bevidsthed og altsaa
ikke kunne paavirkes ved nogen Dom. Dommen over, hvad
der sker udenfor de menneskelige Handlingers Omraade, faar
en æstetisk eller religiøs Karakter, det vil sige, har Karakteren
af Stemningsudbrud overfor det, der væsentligt maa tages, som
det er. Den etiske Dom kan derimod selv motiveres ved det
Princip, ifølge hvilket den fældes, og maa tjene det bestemte
% --- p. 44 ---
\setcounter{footnote}{0}%
\opage{44}Formaal at be virke større Velfærd. Derfor kan den kun angaa
saadanne Begivenheder, der kunne motiveres ved en Dom, altsaa
Handlinger. Dette træder tydeligst frem, hvor det vurderende
Individ er det samme som det handlende, eller hvor den etiske
Dom anerkendes af det handlende Individ; i andre Tilfælde vil
det blive en særlig Opgave, hvorledes det handlende Individ
kan bringes til at anerkende Dommen, en Opgave af psykologisk-pædagogisk
Natur. Men tillige følger heraf, at den etiske
Vurdering selv kan vurderes --- og det efter sit eget Princip.
Naar etiske Domme kun kunne motiveres ved deres praktiske
Indflydelse paa de handlende Væsener, de angaa, saa er det jo
klart, at de tilsidst maa betragtes som pædagogiske Midler. Af
de to Begreber, Vurdering og Opdragelse, er Opdragelsen det
højeste, det mest omfattende. En Vurdering, der ikke selv kan
tjene som Middel i det Formaals Tjeneste, der danner Grundlaget
for selve Vurderingen, vilde være en haandgribelig Inkonsekvens.
Paa denne Maade viser, som ogsaa de følgende
Undersøgelser ofte i det Enkelte ville godtgøre, Etiken ifølge
sit eget Princip ud over sig selv.

Jeg bruger Ordet Velfærd hellere end Ord som Lykke eller
Nytte, fordi disse let kunne hidføre og ogsaa have hidført Misforstaaelser.
Jeg tænker ved Ordet Velfærd paa Alt, hvad der
fører til Tilfredsstillelse af den menneskelige Naturs Trang efter
dens hele Omfang. Etiken skal tage Hensyn til alle Livets Trin
og kan derfor ikke gaa ud fra en Distinktion mellem ydre og
indre, lavere og højere Velfærd. En saadan Distinktion er jo
allerede en Vurdering og kan altsaa først finde Sted, naar Vurderingsprincipet
er givet. De Misforstaaelser, Velfærdsprincipet
(eller som dets vigtigste Forkæmper, \textsc{Bentham}, kalder det: Principet
om den størst mulige Lykke for det størst mulige Antal
Mennesker)\footnote[1]{Bentham er ikke Ophavsmand til denne Formel. Han har selv erklæret, at han har den enten fra Priestley eller fra Beccaria. Men begge disse Forfattere ere igen paavirkede af Hutcheson, hos hvem} har været Genstand for, skrive sig for en stor
Del fra, at man ikke har holdt sig dette for Øje. --- Dog er
der ogsaa et andet Punkt, som tydeligere betegnes ved Ordet
% --- p. 45 ---
\setcounter{footnote}{0}%
\opage{45}Velfærd end ved andre Udtryk, man kunde bruge. Det peger
nemlig hen paa en total Tilstand. Øjeblikkelige Lyst- og Ulystfølelser
ere intet sikkert Kriterium paa Livet som Helhed. Efter
en psykologisk Hypotese ere Lyst og Ulyst ganske vist som
Regel Udtryk for Livets Fremgang eller Tilbagegang: Smerte
er Tegn paa en begyndende Opløsning af Livet, Lyst paa dettes
normale og harmoniske Udfoldelse. (Smlgn. min Psykologi VI D,
2. 3). Men af enkelte isolerede Lyst- og Ulystfølelser kan intet
Sikkert sluttes; heller ikke en blot Opsummering af dem vilde
føre til Maalet. Der maa derimod søges en Sammenhæng mellem
Følelserne og den hele Karakter, den reale Enhed i Bevidstheden,
som de høre til, saavel som med Organismens hele
Tilstand. Ligeledes maa de enkelte Individers Lyst- eller Ulystfølelse
betragtes i Sammenhæng med den hele sociale Tilstand.
Den saakaldte Utilitarianisme, den især af Bentham grundlagte
etiske Skole, der har den store Fortjeneste med Energi at have
gjort Velfærdsprincipet gældende, har skadet sin Sag ved at
lægge en psykologisk Teori til Grund, der opløser Bevidstheden
i en Mængde af Fornemmelser og Følelser og Samfundet i en
Mængde af Individer. Den Betydning, Følelsen af Lyst og Ulyst
har for den varige og omfattende Velfærd, kan ikke udtømmes
ved et simpelt Regnestykke. --- Det følger af sig selv, at Bevisbyrden
paahviler den, der fordrer Opgivelse af en Lystfølelse i
det enkelte Øjeblik eller hos det enkelte Individ: der maa paavises
de Interesser, til Fordel for hvilke Opofrelsen skal finde
Sted (smlgn. § 4).

12. Velfærdsprincipet giver os ingenlunde en Tryllenøgle,
der kan lukke op overalt uden videre. Det er en Misforstaaelse
af, hvad et Princip er, naar man mener, at det med ét Slag
klarer alle de enkelte Tilfælde for os. Vi søge Principer, netop
fordi de enkelte Tilfælde kunne være saa indviklede og 
% CHECK p. 45: note whose mark was not found
\footnote[1]{Formelen ofte forekommer. (Se Den ny ere Filosofis Historie\textsuperscript{2} I. p. 400). Hutcheson er igen paavirket af sit ivrige Studium af de senere Stoikere, hos hvem det egentlige Udspring af det berømte Princip maa søges. Dette er paavist af W. R. Scott: Francis Hutcheson. Cambridge 1900. p. 272—279. Smlgn. ogsaa ovenfor p. 38 Note.}sammen%
% --- p. 46 ---
\setcounter{footnote}{0}%
\opage{46}såtte, at vi kun faa nogenlunde Overblik ved at gaa ud fra
visse bestemte Forudsætninger. Og, som allerede antydet, et
Princip har væsentligt den Betydning at fordele Bevisbyrden.
Velfærdsprincipet forlanger et Bevis af den, der vil hæmme
Livet i dets uvilkaarlige Udfoldelse, og som vil paaføre Smerte,
medens der intet særligt Bevis behøves for Berettigelsen i saadanne
Tilfælde, hvor man fremmer Livsudviklingen og skaffer
forøget Lykke til Veje. Og uagtet den filosofiske Etik er en
praktisk Videnskab, tilfredsstiller den dog ogsaa en teoretisk
Interesse: nemlig den at begrunde vor Vurdering af Handlingerne.
En saadan Begrundelse bliver umulig, naar vi ikke ---
gennem mange Mellemled maaske --- gaa tilbage fra den indviklede
historiske Sammenhæng, i hvilken hver enkelt Handling
foreligger, til et almindeligt Princip, der i og for sig ikke indeholder
Andet end den Aand og Retning, i hvilken den etiske
Vurdering skal gaa. Velfærdsprincipet staar i Etiken som Aarsagsprincipet
i Erkendelsesteorien. Hverken teoretisk eller praktisk
give de almindelige Principer os fuldstændige Løsninger af
de specielle Problemer. Læren om Principerne stilles vel forrest
ved den syntetiske eller systematiske Fremstilling af en
Videnskab, men deri ligger ikke, at de ere fundne først; tvertimod
findes de fra først af ved Analyse af visse fremtrædende
Fænomener og opstilles saa hypotetisk som Forudsætninger for
de specielle Forklaringer.

Den filosofiske Etiks Ræsonnement maa ikke forvexles med
den Maade, paa hvilken vi under den praktiske Overvejelse
drøfte, om en Handling skal foretages. Under den praktiske
Overvejelse ledes vi af Instinkter og Drifter, af Motiver, der
for en stor Del ikke komme fuldstændigt frem i vor Bevidsthed,
af Tanker og Følelser, hvis første Udspring vi ikke drøfte. Vi
ledes af den „positive Moralitet“, i hvilken vi have levet os
ind, og som for en stor Del er et Arvegods i Slægten. Den
etiske Kunst gaar forud for den etiske Videnskab; denne søger
dels at paavise, dels at berigtige de Principer, af hvilke hin
ubevidst ledes. Den teoretiske Reflexion vil naturligvis ofte virke
berigtigende og lutrende tilbage paa den praktiske Overvejelse;
derfor er fuldstændig Diskussionsfrihed om etiske Spørgsmaal
% --- p. 47 ---
\setcounter{footnote}{0}%
\opage{47}af saa overordentligt stor Betydning. Men kun i Tvivlstilfælde
kommer bevidst Reflexion til at spille en direkte Rolle ved
etiske Afgørelser\footnote[1]{Smlgn. mine Etiske Undersøgelser. København 1891. p. 26—29.}.

13. Naar Hensynet til den største og mest omfattende Velfærd
skal være Principet for al etisk Begrundelse, saa kan det
naturligvis ikke selv etisk begrundes. Hvis det virkeligt er det
rette etiske Princip, det vil sige: hvis det virkeligt betegner
den afgørende Maalestok for den Vurdering, vi ud fra vort Standpunkt
anstille, --- saa træder først med dets --- bevidste eller
ubevidste --- Anerkendelse den etiske Betragtningsmaade i Kraft,
og der kan ikke føres nogen etisk Diskussion med den, der
negter det. Kun ad indirekte Vej kan der være Tale om en
Begrundelse af selve Principet. Forsaavidt nemlig den, der negter
vort Princips Gyldighed, selv opstiller andre Principer, kan man
søge at godtgøre, at disse forudsætte hint og logisk set ere afledede
i Forhold til det. Dette vilde være den eneste frugtbare
Maade, paa hvilken Debatten kunde føres mellem Kants og Benthams
Skole. Denne Vej er \textsc{Henry} Sidgwick med stor Dygtighed
slaaet ind paa i sit Værk The Methods of Ethics. Han
gaar dér endnu videre og søger at vise, hvorledes den i det
moderne Europa herskende moralske Tradition kun finder sin
fulde Forklaring, naar man tænker sig den opstaaet under Indflydelse
af en ubevidst Utilitarianisme. Spredte Bud og Afgørelser
kunne bringes i indbyrdes Sammenhæng, og den Tvivl
og Ubestemthed, der kan herske med Hensyn til Afgørelsen af
enkelte Spørgsmaal, kan overvindes, naar man lægger Velfærdsprincipet
til Grund. Ligeledes vinde vi derved en Forklaring
af, hvorfor visse Dyder til visse Tider og hos visse Folkeslag
have spillet en særligt fremtrædende Rolle.

Men overfor den, der resolut stiller sig paa Individualismens
(endsige paa Øjeblikssuverænitetens) Standpunkt, hjælper dette
som sagt ikke. Disse Standpunkter ere, saalænge de ere konsekvente,
logisk uangribelige. Det Samme gælder overfor dem,
der kun udstrække Vurderingsprincipet til at gælde for en snevrere
social Gruppe, idet de gøre Familien, Kasten, Nationen,
% --- p. 48 ---
\setcounter{footnote}{0}%
\opage{48}Racen eller Sekten til Et og Alt. Det snevrere Samfund kan
derved komme til at staa som Individualist eller Egoist i Forhold
til det mere omfattende. Vi kunne paa denne Maade
tænke os ligesaa mange etiske Systemer, som der er større
eller mindre Helheder. Den etiske Verden strækker sig fra
Øjeblikket til Menneskeslægten, men mellem disse Yderpunkter
gives der mange Punkter, hvor man kan stanse, og udover
hvilke kun den psykologisk-historiske Udvikling fører ved de
Forandringer, den volder i det menneskelige Følelsesliv. Et stort
Exempel herpaa er den Maade, hvorpaa der som Følge af den
kulturhistoriske Udvikling efter Alexander den Stores Erobringer
opstod en fælles Menneskelighedsfølelse, der blev Grundlag for
en ny Etik. Det bliver Psykologiens og Historiens Sag at vise,
hvorledes de forskellige etiske Standpunkter og Vurderingsmotiver
blive til. Et Standpunkt afløses af et andet gennem en Udviklingsgang,
der underligger psykologiske og historiske Love. Der
er en kontinuerlig Sammenhæng mellem Standpunkterne eller
Vurderingsmotiverne, forsaavidt som de alle udspringe af den
menneskelige Natur og blive til under den menneskelige Histories
Gang\footnote[1]{Smlgn. Etiske Undersøgelser. p. 7 —14.}. Den historiske eller sammenlignende Etik (smlgn. I, 4),
som netop støtter sig til Historie og Psykologi, tager fat dér,
hvor den filosofiske Etik, der søger Begrundelsen for de etiske
Domme, stanser. Grænsen for den filosofiske Etiks Videnskabelighed
er altsaa ikke Grænsen for al videnskabelig Behandling
af etiske Domme. I Historien kæmpe de forskellige etiske Systemer
og søge ad Opdragelsens og Agitationens Vej at fortrænge
eller rettere absorbere hverandre. Den Maade, paa hvilken
denne Opdragelse og Agitation øves, maa naturligvis konsekvent
være bestemt ved det enkelte Systems egne Principer: man opdrager
Andre, ikke efter deres, men efter sine egne Synspunkter.
Hvad det gælder om, er at tilvejebringe et hidtil manglende
psykologisk Grundlag, vække et vist bestemt Vurderingsmotiv.
Man vil forhøje de Andres Niveau, ikke sænke sit eget. Men
før Opdragelsen er fuldendt, kunne de Andre ikke anerkende
de Principer, efter hvilke de ere blevne opdragne. ---

% --- p. 49 ---
\setcounter{footnote}{0}%
\opage{49}Kort sagt: enhver Maalestok for Handlingers Vurdering
støtter sig til bestemte psykologisk-historiske Forudsætninger. Den,
der skal anerkende og for Alvor anvende Principet om den
størst mulige Velfærd for det størst mulige Antal af bevidste
Væsener, maa ikke være Egoist eller Individualist, fanatisk Patriot
eller Sekterer, men maa formaa at følge de menneskelige
Handlinger med uinteresseret og universel Sympati. Dette er
det objektive Princips subjektive Forudsætning. Afset fra denne
bliver det kun intellektuel Kuriositet, naar man anvender Principet
til Bedømmelse og drager Konsekvenser af det.

Det var \textsc{Benthams} Hovedfejl, at han ikke tydeligt saa, at
et subjektivt Princip dannede Forudsætningen for det objektive
Princip. Han fordrede et objektivt Princip for at kontrollere og
regulere de forskellige subjektive Synsmaader og Paastande paa
det etiske Omraade, og et saadant fandt han i Principet om
størst mulig Lykke for saa Mange som muligt. Men den Subjektivitet
(den Samvittighed), der skal reguleres ved det objektive
Princip, ligger stedse selv til Grund for dette Princips Anerkendelse.
En Etik, som ikke tager dette med i Betragtning,
faar stedse en dogmatisk Karakter, saa fortrinligt end dens Princip
maatte egne sig til Leder under den specielle Diskussion.
Den filosofiske Etik maa udtrykkeligt konstatere, hvilket Standpunkt
den vurderende Subjektivitet staar paa. Vejen til dette
Standpunkt kan, historisk set, være lang og bugtet\footnote[1]{Autoriteten. Der søges her paavist, hvorledes Autoritetsforholdet er den vigtigste Kilde til den moralske Følelses Udvikling i Historien. Jeg vilde dog nu noget mere fremhæve Medfølelsens og Fællesskabsfølelsens Betydning i denne Henseende.}. ---

Nogle Ord om, hvorledes hele den Opfattelse af Etikens Betydning
og Begrænsning, som udvikles i dette Kapitel, er bleven
til, ville her være paa deres Plads. --- Jeg gik ud fra det utvivlsomme
Faktum, at der i Livet fældes etiske Domme, Domme,
der billige eller misbillige menneskelige Handlinger eller Livsordninger,
--- Domme om Godt og Ondt. Den første Forudsætning,
vi møde med overfor saadanne Domme, er, at de ikke ind-

') Smlgn. Om Grundlaget for den humane Etik (1876). Kap. III:
% --- p. 50 ---
\setcounter{footnote}{0}%
\opage{50}byrdes maa stride mod hverandre. Den Maalestok, som anvendes
ved Billigelsen eller Misbilligelsen, maa stedse være den
samme; Vurderingen maa gennemføres konsekvent ud fra det
givne Standpunkt. Det er en Maalestok, vi anlægge paa alle
etiske Systemer: der maa være nødvendig Sammenhæng mellem
det Formaal, der lægges til Grund, og de anvendte Midler og
Veje. Da Sokrates lagde saa stor Vægt paa det Bud „Kend dig
selv!“ opstillede han et Princip, hvis Gyldighed og Anvendelighed
ere uafhængige af den reale Forskel mellem etiske Systemer.
Man maa kunne blive enig om, om det, hvor der er
sat et Formaal af bestemt Indhold og Omfang, er konsekvent
at udtale en vis Billigelse eller Misbilligelse af en foreliggende
Handling. Men med den formale Konsekvens er ikke Alt afgjort.
Hvilke Domme der fældes, beror ikke blot paa Tænkningens
Konsekvens, men ogsaa paa, hvilket Formaal der forudsættes.
Naar Formaalet (og med det Grundlaget, Vurderingsmotivet)
ændres, ændres ogsaa hele Systemet af etiske Domme.
To Individer, der lægge helt forskellige højeste Formaal til Grund
for deres Vurdering, kunne derfor kun fortsætte Diskussionen
med hinanden til et vist Punkt. Det forskellige Formaal medfører
forskellig Vurdering. Og man faar Menneskene til at
sætte sig nye Formaal gennem praktisk, personlig Paavirkning
(Opdragelse og Agitation), ikke gennem blotte teoretiske Slutninger.
Mennesket voxer med sine større Formaal, siger
\textsc{Schiller}. Men Mennesket maa allerede være voxet for at
kunne stille sig større Formaal, og af en saadan Væxt afhænger
det derfor ogsaa, hvilke etiske Sætninger der kunne
komme til at staa som begrundede\footnote[1]{Denne hele Tankegang har jeg først udviklet i en Afhandling om Den filosofiske Etiks Principer (Det kgl. danske Videnskabernes Selskabs Oversigter 1886), som optoges i nærværende Værks første Udgave (1887), hvor det udgjorde Størstedelen af Kap. III. Senere har jeg i første Kapitel af Etiske Undersøgelser (1891) og i mine Ziiricherforelæsninger (Etische Principienlehre. Bern 1896) behandlet det samme Emne, og i Den menneskelige Tanke, dens Former og Opgaver (1910) kom jeg endnu engang tilbage til det (p. 346—359). Den ledende Tanke er: hvor langt er Be-}.

% --- p. 51 ---
\setcounter{footnote}{0}%
\opage{51}Godt og Ondt ere Forholdsbegreber, --- som alle Begreber
ere det\footnote[1]{Se herom Den menneskelige Tankep. 237 —245; 346—349.}. De udtrykke et Forhold til et bevidst eller ubevidst
Formaal og derigennem til en Villie. Det lyder idealistisk, men
det er i Grunden meningsløst, naar man siger, at der gives
noget i og for sig — d. v. s. af ethvert Formaal uafhængigt
--- Godt eller Ondt. Gaar man blot langt nok tilbage, vil man
finde et Formaal og en Villie bag enhver etisk Dom.

(1913). Ganske kort kan Brodden i det etiske Problem
udtrykkes ved, at Værdi kun kan maales ved Værdi. Der
maa tilsidst forudsættes en Grundværdi, der bestemmer alle
andre Værdiers Gyldighed og deres indbyrdes Plads. Enhver
Vurdering (eller «Omvurdering“) forudsætter, bevidst eller
ubevidst, en saadan Grundværdi. En nærmere Undersøgelse
vil vise, at enhver saadan Grundværdi betinges ved Hensynet
til en individuel eller social Helhed og Betingelserne for deres
Bestaaen og Udvikling. Etikens rationelle Karakter beror paa
Muligheden af at finde og formulere disse Betingelser. Men
stedse paany vil Modsætningen og Brydningen mellem forskellige
individuelle og sociale Helheder stille Problemet. Værdibegreberne
ere de mest specielle, de mest indholdsrige og de
mest betingede af alle vore Grundbegreber (Kategorier)*).

14. De etiske Domme vise altsaa stedse --- om man end
nok saa meget begrunder dem ved et objektivt Vurderingsprincip
--- tilbage til et subjektivt Grundlag. De ere tilsidst Udtryk
for en Følelse. --- Men om Følelser, siger man, lader der
sig ikke disputere. Dette er forsaavidt sandt, som en Følelse
er et psykologisk Faktum og maa tages som saadant. Men i
enhver Følelse indgaar der visse bestemte Forestillinger, af
hvilke dens Karakter og Retning afhænger, og disse Forestillingers
Sammenhæng og Gyldighed kan diskuteres. En saadan
Diskussion vil, omend langsomt, virke tilbage paa Grundlaget.
Man vil saaledes kunne vise, med hvor liden Ret Individualisten
grundelse mulig paa dette Omraade? --- I Grundlaget for den humane
Etik (1876) saa jeg endnu ikke dette Problem skarpt.

Smlgn. Psykologi V D, 5.
% --- p. 52 ---
\setcounter{footnote}{0}%
\opage{52}betragter sig som et isoleret og enestaaende Væsen. Man vil
kunne paavise det Ulogiske i at indskrænke sin Menneskekærlighed
til kun at omfatte menneskelige Væsener af en vis
Farve, en vis Afstamning eller en vis Tro. Den fuldstændige
Flytning af Tyngdepunktet vil dog først tilvejebringes, naar
personlige og historiske Erfaringer virke med.

Saaledes er enhver praktisk Videnskab stillet. Betydningen
af dens Sætninger beror tilsidst paa en Følelsesinteresse, der
dikterer det Formaal, hvortil Midlerne søges. Nationaløkonomien
forudsætter Erhvervstrangen som given, og dens Læresætninger
staa og falde med denne Trangs Virken. Den undersøger de
Veje og Midler, ved hvilke det af denne Trang satte Formaal
kan naas. Og forsaavidt der hos Menneskene røre sig andre
Fornødenheder, som faa Indflydelse paa deres Adfærd, saa faa
de nationaløkonomiske Sætninger en abstrakt og hypotetisk Karakter,
gælde kun under en Forudsætning, der i Virkeligheden
ikke fuldstændigt bekræfter sig. Ja, der maa gaas et Skridt
videre. Nationaløkonomien er ikke en af Etiken aldeles uafhængig
Videnskab. Den er ikke blot en Lære om Produktionen,
men ogsaa en Lære om Fordelingen. Den betragter ikke det
økonomiske Omraade som aldeles isoleret, men ser det i dets
Sammenhæng med den almindelige Kultur, og tilsidst med det
menneskelige Samfundslivs etiske Formaal. Den nationaløkonomiske
Opgave bliver tilsidst den at ordne og lede Produktion og
Fordeling saaledes, at en menneskelig Tilværelse bliver opnaaelig
for saa Mange som muligt i Samfundet. Nationaløkonomien
er altsaa en etisk Videnskab og hviler tilsidst paa samme Grundlag
som den humane Etik. Det samme gælder om Retslæren.
Naar man skelner mellem Moral og Ret, er den væsentlige Forskel
den, at Rettens Indhold kan gennemføres og virkeliggøres
ved ydre Magt, medens det egentligt Moralske kræver indre
Anerkendelse og Villiestilslutning. Men kun den bøjer sig for
Magten, der enten anerkender Magtanvendelsen som berettiget
eller indser sin egen Afmagt; og paa den anden Side anvendes
Magten til Rettens Haandhævelse kun, hvor der er Evne og
Villie dertil hos de Personer, hvem den tilhører. Retslærens
Sætninger faa derfor ogsaa en abstrakt og hypotetisk Karakter.
% --- p. 53 ---
\setcounter{footnote}{0}%
\opage{53}Der ses endog i den almindelige Retslære bort fra, om Midlerne
til Rettens Haandhævelse faktisk ere til Stede. I videste Forstand
forstaas nemlig ved Retsforhold saadanne Livsforhold, som
egne sig til retslig Ordning. Man bygger altsaa paa Ideen om
et Samfund, hvor Alt, hvad der egner sig dertil, er retsligt ordnet,
ligesom Nationaløkonomien i mange af sine Sætninger bygger
paa Ideen om et Samfund, hvor Erhvervstrangen virker, og
den humane Etik paa Ideen om et Samfund, hvor Menneskekærlighed
(uinteresseret og universel Sympati) hersker, og almindelig
Velfærd tilstræbes. Og i sidste Instans falder ogsaa
for Retslærens, ligesom for Nationaløkonomiens Vedkommende,
den bestemte Forskel fra Etiken bort, ti Begrundelsen af Magtanvendelsen
er i Retslæren tilsidst af etisk Natur, idet den søges
i Nødvendigheden af, at der dannes en fast, ydre Livsorden med
bestemte Grænser for Handlefriheden, for at det etiske Liv og
de højere menneskelige Bestræbelser kunne udfolde sig i Ly
deraf. Ogsaa Retslæren er altsaa en etisk Videnskab og hviler
paa samme Grundlag som den humane Etik. --- I ethvert Tilfælde
kommer man ved at drøfte det sidste Grundlag for Nationaløkonomi
og Retslære tilbage til Problemer, der ere analoge
med dem, Spørgsmaalet om Begrundelsen af etiske Domme
fører til.

Baade Nationaløkonomi og Retslære have, ligesom Etiken,
ofte lidt under den Misforstaaelse, at de skulde være rene Fornuftvidenskaber.
De have da faaet en dogmatisk og uhistorisk
Karakter, der strider mod deres Natur som praktiske Videnskaber.
De forudsætte en Subjektivitet med visse bestemte Interesser,
og den Mulighed er ikke udelukket, at disse Interesser
kunne ændres.

15. Dersom man i Etiken kun vilde gaa ud fra det subjektive
Princip, Grundlaget, vilde Etiken kun blive en Lære om
den etiske Følelse. Da nu enhver Etik dog maa lære Noget om,
hvad der bør gøres, vilde man blive nødt til at udlede Indholdet
af Grundlaget. Men naar dette ikke sker saaledes, at der
først søges et objektivt Princip, bliver Etiken kun en Række
Postulater, som det endog overfor dem, der staa paa samme
Grundlag, bliver umuligt at begrunde.

% --- p. 54 ---
\setcounter{footnote}{0}%
\opage{54}Jeg har søgt at vise, hvorledes det subjektive og det objektive
Princip udvikle sig med hinanden og svare til hinanden.
En uinteresseret og universel Sympati kan ikke --- uden at opgive
sig selv --- lægge andet Princip til Grund for Handlingens
Vurdering end Principet om den almindelige Velfærd. Af Kortsynethed
og Utaalmodighed kan den komme til at handle mod
dette Princip; men den handler da ogsaa mod sin egen Grundværdi,
sit Formaal, og jo mere den bliver sig dette bevidst, des
mere vil den ogsaa i det Enkelte følge Velfærdsprincipet. Samvittigheden
er ikke ufejlbar: derfor behøver den at ledes af et
objektivt Princip. I ethvert bestemt Øjeblik maa Vurderingen
naturligvis ske efter den Indsigt, som er forhaanden; der gives
ingen anden Domstol end den saa godt som muligt oplyste Samvittighed.
Men derfor kan Afgørelsen alligevel være objektivt
urigtig; en konsekventere, paa mere udstrakte Erfaringer bygget
Anvendelse af Velfærdsprincipet vil kunne godtgøre det. Maaske
blev dog den større Konsekvens og den mere udstrakte Erfaring
selv kun mulig ved, at Afgørelsen i det givne Tilfælde
skete efter den Indsigt, som havdes. Samvittigheden er den
højeste Autoritet, men en Autoritet, der stedse kan gøre sig
selv fuldkomnere. Det objektive Princip gør ikke blot Forhandling
mellem forskellige Samvittigheder (paa samme Grundlag)
mulig, men er ogsaa Betingelsen for, at det enkelte Individs
Samvittighed kan dømme sig selv.

Paa denne Maade virker altsaa Indholdet tilbage paa Grundlaget,
det objektive Princip paa det subjektive. Subjektiviteten
kontrollerer sig selv ved det Princip, der indeholdes i det Formaal,
den sætter sig. Paa tilsvarende Maade kontrollerer den
sig selv paa det teoretiske Omraade ved Identitetsprincipet og
Kausalitetsprincipet, uagtet ogsaa disse Principer i sidste Instans
ere satte ved en subjektiv Virksomhed\footnote[1]{Smlgn. min Psykologi VB, 11; D, 3.}.

16. Naar man nærmere betragter Forholdet mellem Grundlag
og Indhold, vil det let kunne synes, som om Slutningen fra
Grundlag til Indhold har større Sikkerhed end den fra Indhold
til Grundlag. Historien viser os, at et og samme etiske Indhold,
% --- p. 55 ---
\setcounter{footnote}{0}%
\opage{55}samme Grundsætninger, Bud og Love, ere opstillede paa meget
forskellige Grundlag. Allerede den Omstændighed, at visse Pligtbud
komme igen indenfor meget forskellige Religioner, er nok
til at vise dette. Etiske Systemer paa forskelligt Grundlag kunne
meget godt frembyde fælles Elementer, skønt det vil vise sig,
at disse Elementer indtage en noget forskellig Plads i de forskellige
Systemer. Egoistens System udelukker f. Ex. ikke Hensynet
til Familie og Stat, skønt disse kun staa der som Midler.
Den nationale og den humane Etik kunne lægge en overordentlig
stor Vægt paa det enkelte Individs Selvhævdelse, skønt
denne ikke, som i Egoistens System, er det absolute Formaal.
--- Holde vi os til Velfærdsprincipet som det betydningsfuldeste
Exempel, vilde det saa ikke kunne motiveres paa anden Maade,
end det her er sket? --- Hertil er at svare, at historisk set
synes det ganske vist saa, men at der stedse vil kunne paavises
Inkonsekvenser og Ufuldkommenheder, naar Indhold og
Grundlag ikke fuldstændigt og reciprokt svare til hinanden.

Der er især to Maader, paa hvilke man har søgt et Grundlag
for Velfærdsprincipet, foruden den i denne Undersøgelse
anvendte.

Man har ment, at Bestræbelsen for at fremme almindelig
Velfærd er den Vej, ad hvilken det enkelte Individ bedst fremmer
sin egen Velfærd, idet han for sin egen Person derved
naar at sætte Formaal igennem, han ellers ikke, i hvert Tilfælde
ikke saa let og saa sikkert, kunde naa. Bestræbelsen for almindelig
Velfærd bliver altsaa blot en Omvej, der maa slaas
ind paa, for at den Enkelte kan naa sin egen Tilfredsstillelse.
Ethvert etisk Bud angiver en saadan Omvej, og Etiken bliver
den systematiske Lære om saadanne Omveje. Dette var den
Opfattelse af Etikens Grundlag, som den moderne Utilitarianismes
Stifter, \textsc{Bentham}, nærmest hyldede, og som han især
paa en kras Maade har udtalt i et efter hans Død udgivet
Skrift (Deontology)\footnote[1]{I Principles of Morals and Legislation, som er Benthams Hovedskrift, træder denne Lære ikke saa bestemt frem. Han udtaler sig der ikke klart om Etikens Grundlag. Et enkelt Sted hedder det: „The die-}. I min Afhandling „Om Grundlaget for
% --- p. 56 ---
\setcounter{footnote}{0}%
\opage{56}den humane Etik“ har jeg kritiseret Bentham fra denne Side,
og kun fra denne Side. Her vil jeg tilføje, at det er en stor
og meget prekær Hypotese, Bentham benytter som Grundlag
for sin Etik. Det er nemlig intet mindre end den dristige Forudsætning,
at der gives en Harmoni mellem alle Individers
egoistiske Interesser, naar de blot hvert for sig ere klare over,
hvilke disse egentligt ere. Den, som arbejder for sin egen velforstaaede
Interesse, arbejder altsaa, efter denne Antagelse, ogsaa
for alle Andres. Derfor behøves der egentligt ikke andre
Motiver end egoistiske, og selve Velfærdsprincipets Opstilling
kan ogsaa motiveres egoistisk. Enhver vurderende Subjektivitet
bedømmer Handlingerne efter den Maade, hvorpaa de fremme
almindelig Velfærd, fordi selve det vurderende Individ finder
sig selv personligt bedst tjent med saadanne Handlinger.

Man vil let se, at det er en Hypotese, der staar paa meget
svage Fødder og stadigt udæsker Tvivlen. Den fører ind i meget
indviklede Undersøgelser, som aldrig vil kunne afsluttes. Det
er derfor et meget betænkeligt Grundlag at opføre Etiken paa.
Og dertil kommer, at selv om man antager dette Grundlag, vil
det Indhold, som bygges derpaa, kun være et ydre Middel, ikke
umiddelbart i og for sig svare til, hvad der ligger i Grundlaget.
Etikens Indhold bliver et System af Indrømmelser, som Individet,
der egentligt staar paa den absolute Individualismes Standpunkt,
maa lade sig afpresse. Ti den lige Vej er dog altid den
bedste og Omvej kun Nødhjælp.

Et andet Grundlag, paa hvilket Velfærdsprincipet er fremtraadt,
er det teologiske. Ogsaa her støtter man sig til vidtgaaende
og disputable Forudsætninger. Den etiske Vurdering
gøres herved afhængig af de religiøse Problemers Løsning.
Men Bestræbelserne for at opstille en filosofisk Etik ere netop
fremgaaede af Ønsket om at gøre det etiske Problem 
% CHECK p. 56: note whose mark was not found
\footnote[1]{tates of utility are neither more nor less than the dictates of the most extensive and enlightened benevolence“ (X, 36). Men af denne Sætning gøres der aldeles ikke Brug ved Opstillingen af selve det etiske Vurderingsprincip (ch. I— II). --- Smlgn. min Karakteristik af Bentham i Den nyere Filosofis Historie' Il. p. 363—368.}uafhæn%
% --- p. 57 ---
\setcounter{footnote}{0}%
% CHECK p. 57: notes paired with the marks in order (the layer misread a number)
\opage{57}gigt af de religiøse Problemer. Dersom Etiken bygges paa
teologisk Grundlag, maa disse Bestræbelser siges at være forfejlede.
Dertil kommer, at den teologisk grundlagte Etik rammes
af en lignende Indvending som den Etik, der bygger paa den
forudsatte Harmoni mellem de velforstaaede egoistiske Interesser.
Hele Etiken bliver nemlig Midlet for Tilfredsstillelse af en
Interesse, der i sig selv er forskellig fra --- i hvert Tilfælde
ikke falder ganske sammen med — selve den etiske Interesse.
Individets Forhold til det, der antages at være Guddommens
Villie, er et Forhold for sig, og det skal først bevises, at der
er Overensstemmelse mellem Guddommens og Velfærdsprincipets
Krav. I Virkeligheden har ogsaa enhver teologisk Etik, selv om
den stiller Velfærdsprincipet meget højt, en Klasse af særegne
Pligter, som skyldes det teologiske Grundlag, men ikke begrundes
ved Velfærdsprincipet, ja endog kunne komme i Strid
med dette\footnote[1]{Se Om Grundlaget for den humane Etik. p. 68 —72.}.

Ved Opstillingen af en Videnskabs Principer bør Sparsommelighedens
Lov raade. I den af uinteresseret og universel
Sympati betingede etiske Følelse have vi et Grundlag, hvis
Realitet er utvivlsom, skønt der kan være Strid om, hvor stor
en Rolle det spiller i Verden. Menneskekærligheden har, som
det i det Foregaaende blev antydet, langsomt arbejdet sig frem
og er endnu en kæmpende Magt, --- maa kæmpe om Herredømmet
baade i det Ydre og i selve den Enkeltes Indre. Om
den vil blive det Vurderingsmotiv, der en Gang vil sejre hos
Alle, vide vi ikke, og vi behøve ikke at vide det\footnote[2]{Axel Hågerstrom (Undersokning af den empiristiska Etikens möjlighet med särskild hänsyn till dess moderna hufvudformer. Upsala 1895. p. 20 f.) mener, at jeg maa antage den universelle Sympati for lige stor hos alle Mennesker, da Etiken ellers intet almengyldigt Grundlag faar og derved mister Karakteren af Videnskab. Min Mening er nu dog, at den universelle Sympati udvikler sig baade sporadisk og successivt, og at det Samme gælder hvilketsomhelst andet Vurderingsmotiv, man vil nævne, --- og netop deri ser jeg en Begrænsning af Etikens (ikke blot min Etiks, men enhver Etiks) Videnskabelighed. --- En teologisk Kritiker (N. Tejsen i „Vidar“ 1888) spørger, hvad den filosofiske Etik vil gøre,}. Den spiller
% --- p. 58 ---
\setcounter{footnote}{0}%
\opage{58}allerede nu en saadan Rolle i Verden, at det er af meget stor
Interesse at undersøge, hvilke etiske Domme der konsekvent
udspringe af den. Og det vil vise sig at være et overordentligt
hensigtsmæssigt Valg at lægge netop dette Vurderingsmotiv til
Grund. Ti der er neppe noget andet Motiv, der i den Grad
som dette afføder Interesse for at forfølge Handlingernes Virkninger
i deres hele Forløb, saa langt som det overhovedet er
muligt. Individualisten stanser sin Undersøgelse, naar han har
sikret sig, at en Handlings Virkninger ikke berøre hans Velfærd;
den, der staar paa Nationaletikens Grund, er ubekymret,
naar en Handlings Virkninger blot ikke berøre hans Folk. End
ikke Hypotesen om en Harmoni mellem de forskellige individuelle
Interesser (hos Enkelte eller hos Nationerne) bliver for
disse Standpunkter af afgørende Betydning. Hvad angaar det
Individualisten, om Andres Interesser stemme med hans, naar
han blot har Midler til at sikre sin egen Interesse? I det
Højeste kan han føle en vis Tryghed ved Tanken om hin
Harmoni. Universel Sympati som Vurderingsmotiv fører derimod
til at efterspore den Maade, paa hvilken Handlingernes Virkninger
gribe ind i alle personlige Væseners Liv, og afføder derved
en Interesse for Undersøgelser af naturvidenskabelig,
psykologisk, historisk, social og statistisk Art, der ganske kan
gaa sammen med den videnskabelige Interesse for saadanne
Undersøgelser. I Kraft af det valgte Vurderingsmotiv føres vi
altsaa ud over Etiken, eller føres Etiken til at støtte sig til
Naturvidenskab, Historie, Psykologi, Statistik og Nationaløkonomi,
saa at Vurderingen kan blive saa objektiv og saa realistisk som
muligt. Og denne Overgang sker med indre Konsekvens og paa
en mere direkte Maade end paa andre Standpunkter, der forsøge
at optage og begrunde Velfærdsprincipet. Til den universelle
Sympati som Vurderingsmotiv svarer Velfærdsprincipet som
et objektivt Vurderingsprincip for menneskelige Handlinger; det
udfolder kun nærmere de logiske Konsekvenser af det 
% CHECK p. 58: note whose mark was not found
\footnote[1]{hvis den som Vurderingsmotiv forudsatte Følelse mangler. Her kan der svares med et Kontraspørgsmaal: hvad gør den teologiske Etik overfor den, der hverken tror paa Gud eller Djævelen? Il}psykolo%
% --- p. 59 ---
\setcounter{footnote}{0}%
\opage{59}Giske Grundlag og gør en konsekvent Begrundelse af dettes
enkelte Ytringer mulig. Etikens Grundlag og Indhold ere her
rykkede hinanden saa nær, som det idethele er muligt.

Enhver etisk Opfattelse, der gaar ud fra en Autoritet som
Grundlag --- denne være nu overnaturlig eller naturlig, bestaa
i Magthaveres Villie, i Forbilleders Dragen eller i den offentlige
Menings Tvang --- maa dog indrømme, at saasnart Individet
umiddelbart kan gøre Autoritetens Formaal til sine egne Formaal,
vil dermed det eneste sikre og simple Grundlag for det
Etiske være naat, og tillige Muligheden sikrest være given for
en stadig Berigtigelse af det Etiske. Det er en Omvej, naar jeg
kun gennem en Autoritets Villie træder i Forhold til de højeste
Formaal, min Villie anerkender. I den energiske og intellektuelt
udviklede Menneskekærlighed haves det eneste Grundlag, hvor
de individuelle og universelle Formaal umiddelbart kunne identificeres.
Overfor en Kraft som denne kan al Autoritet kun have
foreløbig og pædagogisk Betydning. Al Autoritet kan forsvinde,
uden at dermed det blivende Grundlag for det Etiske forsvinder.
Et er, at den etiske Lov (naar vi se bort fra Øjebliksstandpunktet
og den absolute Individualismes Standpunkt) er en
omfattende Helheds Livsbetingelser, formulerede i bestemte Tanker,
et Andet, hvorledes disse Tanker anerkendes af det enkelte
Individ. Selv hvor hine Livsbetingelser indskærpes af ydre Love
og af den offentlige Mening, kan jo Individet ad sin egen Overbevisnings
Vej være kommen til dem. Maaske har den Enkelte
endog en langt klarere og mere vidtgaaende Indsigt i Samfundets
Livsbetingelser, end Lovene og den offentlige Mening hidtil
have formaaet at udtrykke, --- og skulde saa det Sindelag,
med hvilket han følger denne Indsigt i sin Handlen, ikke fortjene
Navn af etisk, fordi ingen Autoritet sanktionerer den? Det
Vurderingsmotiv, som det vil have størst Interesse at prøve,
maa dog være det, som mest direkte formaar at vurdere enhver
Autoritet efter den Maalestok, al Autoritet maa fyldestgøre for
at kunne anses som berettiget. Bygger man al Etik paa Autoritet,
hvor vil man saa finde et Standpunkt, ud fra hvilket den
etiske Vurdering af Autoriteter kan foregaa?

17. Forskellen mellem Grundlag og Indhold falder ikke
% --- p. 60 ---
\setcounter{footnote}{0}%
\opage{60}sammen med Forskellen mellem individuel og social Etik. Til
Etikens Indhold høre baade de individuelle Karakteregenskaber
og de sociale Livsformer, som stemme med Vurderingsprincipet;
det indbefatter altsaa baade individuel og social Etik. Det Spørgsmaal,
om den individuelle og den sociale Etik ere uafhængige
af hinanden, eller om den ene af dem er bestemmende for den
anden, maa afgøres efter Velfærdsprincipet. Der kan indenfor
den specielle Etik gøre sig en Individualisme gældende af anden
Art end den forhen omtalte, idet den ikke bygger paa Individets
Suverænitet, men begrundes ved Hensynet til den almindelige
Velfærd, der kræver saa mange selvstændige og ejendommelige
Udgangspunkter for Handlen som muligt. Paa tilsvarende
Maade gaar det med Begrundelsen af de mindre Samfunds Bestaaen
og Betydning indenfor de større. Det almindelige Velfærdsprincip
fordrer Hensyn til alle bevidste Væsener, hvis Lyst- og
Ulystfølelser kunne sættes i Bevægelse ved Handlingerne.
Men netop ifølge dette Hensyn kan det blive af største Betydning,
at Interessen koncentrerer sig om en snevrere Kreds.
Gennem Livet i Familien og i den bestemte Kaldsvirksomhed
arbejdes der maaske paa bedste Maade for Nationens Fremgang,
og gennem Deltagelsen i det nationale Liv for Menneskehedens.
Alle disse Spørgsmaal ere af mere speciel Art end det,
som det er dette Kapitels Opgave at drøfte.

18. Som vi saa i § 2, er Modsætningen mellem ..subjektiv“
og „objektiv“ Etik fra først af opstaaet ved den forskellige Vægt,
der lagdes paa Handlingens indre og dens ydre Side. Jeg har
søgt at vise, at denne Modsætning gaar længere tilbage og
hænger nøje sammen med Etikens Karakter som praktisk Videnskab.
Jeg vender nu til Slutning tilbage til hine to Sider ved
Handlingen for at faa Lejlighed til at bestemme Forholdet mellem,
hvad jeg her har kaldt det etiske Grundlag eller Vurderingsmotivet,
og de enkelte Handlingers Motiver.

Etikens Historie viser os, at Vurderingen først angaar den
ydre Handling og dens Virkninger, men at den efterhaanden
udvider sig til ogsaa at gælde Motiverne, Sindelaget, Karakteren
hos den Handlende. Det er jo naturligt, at man først holder sig
til det, som synes at ligge klart for Dagen, og som kan blive
% --- p. 61 ---
\setcounter{footnote}{0}%
\opage{61}Genstand for sanselig Iagttagelse. Det aandelige Syn er rettet
mod det Ydre, før det formaar at opfatte det Indre. Handlingerne
have tillige paa de mere primitive Trin væsentligt Karakteren
af Reflexbevægelser og Instinktytringer; Motiverne træde
ikke særligt frem, og Interessen dvæler ikke særligt ved dem.
Der skelnes heller ikke mellem Moral og Ret, mellem indre og
ydre Overensstemmelse med den gældende Lov eller Vedtægt.
Etisk Lov, Sædvane og juridisk Lov have endnu ikke udsondret
sig fra hverandre. En betydningsfuld Side ved den etiske Udvikling
bestaar netop i, at disse Forskelligheder gøre sig gældende.
Dette forudsætter en Evne til at forstaa Handlingens
Afhængighed af bestemte Motiver og disses Tendens til at virke
i en vis bestemt Retning. Den etiske Vurdering udstrækkes da
ogsaa til det Indre. De store Revolutioner paa det etiske Omraade
fremtræde væsentligt som Fremskridt i Inderliggørelse af
den etiske Vurdering. Et saadant Fremskridt have vi, hvor
Etiken skiller sig fra Retslæren ved at betone Handlingens
indre Kilde som det Væsentlige. Det gælder ikke blot om at
udføre eller afholde sig fra den ydre Handling, men om Overensstemmelse
mellem den hele indre Tilstand og den etiske
Lovs Fordring. En saadan Inderliggørelse er tillige en Generalisation.
Ti naar et Motiv (f. Ex. Egoisme, Had) forkastes,
forkastes dermed ogsaa alle enkelte Handlinger, som udspringe
af det, og disse behøve ikke særligt at opregnes. Ligeledes er
Billigelsen af et Motiv Billigelse af alle Handlinger, som udspringe
deraf. Overgangen fra den ydre Vurdering til den indre
er derfor endeligt en stor Simplificering af den etiske Lov. Det
er ikke Buddenes Mangfoldighed, men de Karakteregenskaber,
som skulle herske, og den Retning, i hvilken Samfundsformerne
skulle udvikles, det gælder at fastsætte.

Exempler paa en saadan Inderliggørelse, Generalisation og
Simplificering have vi ved Kristendommens Brud med Jødedommen
og ved Protestantismens Brud med Katolicismen, og
idethele overalt, hvor en rent etisk Betragtningsmaade udskiller
sig fra en Betragtning, der støtter sig til Sædvane, Autoritet og
ydre Magt.

Den etiske Lov fordrer Tilstedeværelsen af det Sindelag,
% --- p. 62 ---
\setcounter{footnote}{0}%
\opage{62}der bærer den selv i Slægten. Det er dette, som Kant har udtrykt
i den Sætning, at det er en Pligt at have Samvittighed.
Da Anerkendelse af Pligter forudsætter Samvittighed, kunde
det synes, som om man her drejede sig i en Kreds. Men det
er kun tilsyneladende; ti Grundlag og Handlingsmotiv behøve
ikke at falde sammen, og det behøver ikke heller at være nogen
Ufuldkommenhed, naar de ikke falde sammen. Ifølge Velfærdsprincipet
kan det nemlig meget vel være nødvendigt, at andre
Motiver end selve Pligtfølelsen gøre sig gældende. Det kan f.
Ex. være nødvendigt eller sundt, at Selvopholdelsesinstinktet
eller ogsaa den umiddelbare Medfølelse leder Mennesket til at
foretage Handlinger til egen og Andres Velfærd, og at Samvittigheden
ikke rører sig ved enhver enkelt Handling. Det
kan endog være et Tegn paa Fuldkommenhed, naar Handlinger,
der kræve Anstrengelse og Opofrelse, foretages, uden at det
sker af Pligtfølelse. Aandelig Øvelse fører jo til, at hvad der
fra først af kun kunde udføres gennem flere psykologiske Mellemled,
gennem udtrykkelig Besindelse og Villiesanspændelse, tilsidst
kan udføres direkte og uden særlig Bevidsthed om, hvorfor
man egentligt handler saaledes.

19. Al Etik er praktisk Idealisme. Den forudsætter, at vi
stille os Formaal, men et Formaal er ikke noget Værende, men
Noget, som skal være. Al Etik forudsætter derfor en Drift, en
levende Følelse og Higen forbunden med en Forestilling om
det, hvorefter der higes. Dersom Virkeligheden fyldestgjorde
enhver Trang, vilde der intet Ideal og altsaa ingen Etik gives.
Men et Ideal, som Etiken skal kunne bruge, maa ikke blot
staa over Virkeligheden; det maa tillige have Tilknytningspunkter
i den givne Virkelighed, saa der udfra denne idetmindste bliver
en Tilnærmelse til det mulig. Det maa — subjektivt set ---
kunne finde Anerkendelse, og det maa --- objektivt set ---
kunne gennemføres mere og mere i Erfaringens Verden. Det
er virkelige Mennesker, som skulle ville og handle, og en
virkelig Verden, som er Skuepladsen for deres Villen og
Handlen. Hvad Etiken fordrer, maa være fysisk, psykologisk
og historisk muligt, ikke staa i Strid med den virkelige Verdens
Love. Den, som forkaster enhver Etik, fordi den opstiller Idealer,
% --- p. 63 ---
\setcounter{footnote}{0}%
\opage{63}kan dermed kun mene, at dens Idealer mangle Tilknytningspunkter
i Virkeligheden. Han vil ikke kunne undgaa selv at
opstille Idealer, saa at han omstyrter Etiken ved en ny Etik.
Naar han f. Ex. fordrer, at man skal være saa ærlig at indrømme,
at Etik er en Umulighed, eller naar han maaske endog
forkaster Etiken, fordi alle Menneskelivets Ulykker skrive sig
fra „Moralen“, saa taler han i Ærlighedens eller Menneskekærlighedens
Navn, optræder altsaa som etisk Idealist. Diskussionen
kan kun dreje sig om, hvilke Idealer man skal stille op.
Ingen slipper for Idealer, hvor meget han end vrider og vender
sig. Men Idealerne kunne være af meget forskellig Værdi, og
om Værdien strides der.

Ogsaa teoretisk Videnskab har stedse en idealistisk Karakter.
Den opererer med Principer og Forudsætninger, som ere
saa simple, at den virkelige Erfaring aldrig fuldstændigt dækker
dem. Al sammenhængende, konstruerende Tænken bliver
kun mulig ved, at vi fastholde en enkelt Side, et enkelt Element
ved det Givne og drage alle Konsekvenserne deraf. Saaledes
betragter Geometrien Tingene blot som udstrakte i Rummet
og udleder Udstrækningens almindelige Love uden at tage
Hensyn til Tingenes andre Egenskaber. Paa alle Omraader maa
vi for at kunne tænke skarpt tænke os Tingene simplere, end
de i Virkeligheden ere. Derved udelukkes det ikke, at Tænkningens
Resultater kunne have Gyldighed for Tingene, naar det
blot er væsentlige og betydningsfulde Sider ved disse, vi behandle
i vor konstruktive Tænkning. Ligesom Geometrien konstruerer
et idealt Rum, konstruerer Etiken paa Psykologiens og
Historiens Grund en ideal Samvittighed. I Virkeligheden er ethvert
Menneskes Samvittighed sammensat af forskellige, undertiden
endog stridende Elementer. Efterligning og Overlevering,
Egoisme og Æretrang, Hensyn til Andres Mening og Frygt for
Straf, Pietet og Menneskekærlighed, Religiøsitet og Samfundsfølelse
--- disse og endnu flere Motiver røre sig indenfor den
sædvanlige Samvittighed, der ikke er saa simpel en Følelse,
som man ofte antager. De forskellige mulige etiske Systemer,
som i det Foregaaende (4—13) ere karakteriserede, kæmpe ikke
blot i Historien og i Slægten, men ogsaa indenfor det enkelte
% --- p. 64 ---
\setcounter{footnote}{0}%
\opage{64}Menneskes Bevidsthed. Enhver har --- eller kan have --- flere
Interesser og Formaal, der kæmpe med hverandre om, hvilken
af dem der tilsidst skal bestemme den Vurdering, Mennesket
anstiller af sine og Andres Handlinger. Enhver, der vil forsøge
paa at udvikle en sammenhængende Etik, maa lægge et enkelt
Vurderingsmotiv til Grund som herskende og bestemmende. Og
idet et enkelt Motiv lægges til Grund, er der allerede foretaget
en Idealisering. Den filosofiske Etiks Historie viser os, hvorledes
man har prøvet at gennemføre de forskellige mulige Vurderingsmotiver.
Saaledes gik \textsc{Hobbes} og \textsc{Bentham} ud fra Egoismen
eller den velforstaaede Egeninteresse, \textsc{Hutcheson}, \textsc{Hume}
og \textsc{Adam} \textsc{Smith} fra Menneskenes paa Fællesinteresser grundede
Sympati, \textsc{Hegel} fra Statens og Kultursamfundets Ide. Den
filosofiske Etik har her den Betydning at bringe de forskellige
Formaal, Mennesker kunne sætte sig, til fuld Bevidsthed og
drage alle Konsekvenser af dem. De forskellige Idealer belyses
derved saaledes, at Kampen mellem dem kan føres med større
Klarhed; det er et af de vigtigste Bidrag, den teoretiske Etik
kan yde til den praktiske etiske Stræben. Bidragets Værdi vil
afhænge af, hvor betydningsfuldt det til Undersøgelse valgte
Vurderingsmotiv er. --- Udfra det valgte Grundlag vurderes alle
Motiver og Tendenser efter, hvorvidt de direkte eller indirekte
føre i Retning af det Ideal, Grundlaget fører til at opstille.
Forskellen mellem teoretisk og praktisk Videnskab er her den,
at i Teorien ere vore Idealer Tilnærmelser til Virkeligheden,
som er det Maal, Tanken søger at naa, medens det i Praxis er
Virkeligheden, der skal ændres og nærmes til Idealet.

20. I tre Henseender fjerne de etiske Idealer sig fra det
virkeligt Givne.

For det Første er der i den virkelige Villen og Handlen
Meget, som ligefrem strider mod den etiske Fordring. Overfor
dette stiller Etiken sig hæmmende og negerende, fordrer det udelukket
fra den etiske Verden. I det praktiske Villiesliv svarer
hertil den hæmmende Virksomhed, hvorved uvilkaarlige, oprindelige
eller erhvervede Tilskyndelser og Tilbøjeligheder kunne
trænges tilbage\footnote[1]{Psykologi Il, 4 e; 6 e. --- IV, 4.}.

% --- p. 65 ---
\setcounter{footnote}{0}%
\opage{65}Dernæst frembyder den virkelige Villen og Handlen ofte
kun svage og ufuldstændige Virkeliggørelser af det, den etiske
Fordring indeholder. Her gælder det da Forøgelse baade i Grad
og Omfang; det gælder Potensering og Udvidelse. Etiken udtaler
her et Excelsior! Retningen er god, men Bevægelsen
mangler Kraft og Fylde. I det praktiske Villiesliv svarer hertil
Anspændelsen af Opmærksomheden, Opbuddet af Energien, idethele
Villiens Evne til gennem sin Indflydelse paa Forestillinger
og Følelser at virke tilbage paa sig seiv\footnote[1]{Psykologi VII B, 1 c: 3; 5 a}. --- Eller Retningen
er god, men Motivet kan ikke anerkendes, saa at en helt anden
Begrundelse er nødvendig, for at Villien eller Handlingen kan
blive fuldkommen. En Villiesretning, der er opstaaet af et Motiv,
kan senere (ved Forskydning, se I, 4) maaske kun anerkendes
formedelst et helt andet Motiv\footnotemark[1].

Endeligt kan den virkelige Villen og Handlen mangle Sammenhæng,
Enhed og Harmoni. Forskellige og indbyrdes stridende
Tendenser og Drifter kunne gøre sig gældende. Det gælder da
om Kombination og Koncentration, om at sammenarbejde og forsone
de stridende og spredte Kræfter. Hertil svarer paa det praktiske
Villieslivs Omraade Indøvelsen af sammensatte Bevægelser
og Forestillingsrækker, og Evnen til at fatte en Beslutning
og derved gøre Ende paa Overvejelsen\footnotemark[1].

Baade ved den hæmmende, ved den forhøjende eller motivændrende
og ved den forenende Virksomhed er Velfærdsprincipet
ledende. Kun hvor det sker Fyldest, faa disse Virksomheder
Værdi. Man kan ligesaa vel hæmme noget Godt som
noget Ondt, ligesaa godt forhøje noget Ondt som noget Godt;
den etiske Værdi ligger derfor ikke i disse Virksomheder i og
for sig, men i den Retning, i hvilken de udøves.

De omtalte Processer finde Anvendelse ikke blot i det enkelte
Individs, men ogsaa i Samfundets og Slægtens etiske Udvikling.
Der kan være individuelle Livsytringer, som stride mod
Slægtens Velfærd, skønt de ikke stride mod Individets egen,

') Psykologi VII A, 5; B, 2.

 Psykologi VI C, 2.
% CHECK p. 65: note mark without a note

% --- p. 66 ---
\setcounter{footnote}{0}%
\opage{66}isolerede Velfærd; de maa hæmmes. En Bestræbelse for social
Velfærd kan mangle Styrke eller Omfang eller ret Motivering;
her bliver da Anvendelse for Potensering eller Udvidelse eller
Forskydning. Endeligt kan der mangle den Sammenhæng og
Harmoni mellem de forskellige Livsytringer og Bestræbelser
indenfor Samfundet og Slægten, som er nødvendig for, at almindelig
Velfærd kan naas; det gælder da at sammenarbejde
og organisere de spredte Virksomheder. ---

Ogsaa fra denne Side viser det sig, at den filosofiske Etik
baade er konservativ og radikal. (Smlgn. I, 4). Den tager sine
Udgangspunkter i det Givne. Paa bar Bund kan den intet udrette.
Men dette Givne søger den saa at udvikle videre paa de
tre anførte Maader.

21. I mere eller mindre stærk Modsætning til den i dette
Kapitel fremsatte Teori om de etiske Principer og om Muligheden
af at begrunde dem have flere danske Forfattere søgt at
finde et efter deres Mening mere objektivt og rationelt Grundlag
for Etiken. Deres Bestræbelser gaar ud paa at hævde Etikens
Videnskabelighed; Enkelte af dem mene endog for første
Gang at have paavist et videnskabeligt Grundlag for Etiken. Særligt
gaa deres Bestræbelser ud paa at frigøre den fra den empiriske
og subjektive Begrænsning, den efter min Opfattelse nødvendigvis
underligger. Det er nu ikke sandsynligt, at enhver
saadan Begrænsning skulde kunne falde bort. Ti det viser sig
overalt, hvor man gaar tilbage til det principielle Grundlag for
vor Erkendelse, at Begrundelse og Begrænsning følges ad. Og
det kan nu engang ikke være anderledes, da enhver Begrundelse
hviler paa visse Forudsætninger. At overse dette fører til
Dogmatisme. Det har været min Bestræbelse klart og bestemt
at pege paa de Punkter, hvor en streng, rationel Etik hviler
paa psykologiske og historiske Forudsætninger, der ikke overalt
og til alle Tider gøre sig gældende. Imidlertid indrømmer jeg,
at de omtalte Arbejder kunne tjene til at supplere min Fremstilling
ved at belyse Sider ved Problemerne, som denne maaske
ikke har ladet komme til deres fulde Ret. Og da de tillige
vidne om den Interesse for etiske Problemer, som i de senere
% --- p. 67 ---
\setcounter{footnote}{0}%
\opage{67}Aar har rørt sig herhjemme, skal jeg gaa noget nærmere ind
paa disse Arbejder, hvad det afgørende Hovedpunkt angaar.

De Vanskeligheder, som gøre sig gældende ved ethvert
Forsøg paa videnskabelig Etik, have deres Grund i, at alle
etiske Begreber og Sætninger angaa Værdiforhold. Etiske Domme
udsige Forholdet mellem menneskelige Handlinger og de Formaal,
de skulle tjene, og Begrebet Formaal forudsætter igen Begrebet
Værdi, naar man ikke tager Ordet „Formaal“ i billedlig
Betydning. Men Værdi kan igen kun det have, der tænkes i
Forhold til et levende Væsens Følelse. Der kommer derved et
subjektivt Element ind. Og desuden kende alle levende Væsener
ikke nødvendigvis de samme Værdier, hvorfor de heller ikke
stille sig de samme Formaal. Her kan da Værdi staa imod
Værdi, i hvert Tilfælde, naar man gaar tilbage til Grundværdierne,
hvoraf de specielle Værdier afhænge.

I det Foregaaende har jeg vist, hvorledes jeg selv stiller
mig til disse Vanskeligheder. Enhver Etik maa, fra hvilken
Grundværdi den saa gaar ud, søge Midler og Veje til at frembringe
det Værdifulde i Virkelighedens Verden. Den maa stille
sig paa Erfaringens Grund, baade naar det gælder om at paavise,
hvori det Værdifulde bestaar, og naar det gælder om at
finde Midler og Veje til at frembringe det. I begge Henseender
grunder Etiken sig efter min Opfattelse paa Biologi, Psykologi
og Sociologi. Alle menneskelige Værdier ere knyttede til Livet,
det personlige Livs og Samfundslivets Bestaaen og Udfoldelse,
og Livets Betingelser maa læres af Livet selv. Dette er Grundtanken
i min Etik, og ved den bliver det muligt, at Etiken
trods sit subjektive Udgangspunkt kan blive en objektiv Videnskab
(se ovenfor § 11 og 14—16). Man vil forhaabentligt ikke
i min Fremstilling finde et eneste Forsøg paa Begrundelse, der
ikke støtter sig til Biologi, Psykologi eller Sociologi. Ligesaa lidt
som Sansekvaliteternes subjektive Karakter gør Fysiken som objektiv
Videnskab umulig, ligesaa lidt gør Værdiernes Subjektivitet
Etiken som objektiv Videnskab umulig. Det gælder begge
Steder at finde de objektive Forhold, til hvilke de subjektive
Bestemmelser svare. Men en væsentlig Forskel mellem Etikens
og Fysikens Stilling fremkommer ved, at medens Sansekvalite%
% --- p. 68 ---
\setcounter{footnote}{0}%
\opage{68}terne nogenlunde ere de samme for Alle, variere Værdierfaringerne
hos de forskellige Individer og Folkeslag og til de forskellige
Tider. Jeg er for meget Empiriker og Historiker til at
kunne overse denne Kendsgerning eller til at mene at kunne
slippe fra den ved abstrakte Konstruktioner. Den videnskabelige
Etiks Opgave bliver da for mig at undersøge, hvilke specielle
Værdier der konsekvent kunne afledes af de forskellige Grundværdier
eller Værdityper. (Se § 3—13). Derved fremtræder Muligheden
af forskellige etiske Systemer. Disse ville dog ikke
falde ganske udenfor hverandre; de kunne mødes i enkelte Sætninger,
som da opstilles med forskellig Begrundelse. (Smlgn.
§ 13 og 19). Selvom de indeholde de samme Elementer, ville
disse dog være forbundne paa forskellig Maade, og Klangfarven
vil derfor blive forskellig. Eller maaske vil det, der staar som
Grundværdi i det ene System, staa som afledet Værdi i det
andet System. En saadan Undersøgelse vil i hvert Tilfælde have
sin Betydning, selvom den kun betragtes som en foreløbig Drøftelse
(en „Phånomenologie des Geistes“). Der sker en stadig
Udvikling paa det etiske Omraade, baade hvad de enkelte Individer
og hvad Slægten angaar. Men vi stikke uheldigvis midt
i denne Udvikling og kunne ikke hæve os over de Ufuldkommenheder,
som følge deraf. Endnu staar Striden, og jeg kan
ikke se Andet end, at den Vej, jeg er slaaet ind paa, i hvert
Tilfælde maa bringe større Klarhed angaaende Stridens Genstand.
I Længden vil naturligvis Klarhed over Konsekvenserne af de
forskellige Standpunkter indvirke paa Udviklingens Gang\footnote[1]{Det har været af meget stor Interesse for mig at se, at den franske Matematiker og Filosof Henri Poincaré indtog en ganske lignende Stilling til det etiske Problem som den, jeg i dette Kapitel har gjort gældende. Se hans Derniéres Pensées. p. 236—238 (ogsaa p. 228 f.). (1913).}.

Naar man nu ikke finder den Maade, paa hvilken jeg stiller
mig overfor de nævnte Vanskeligheder, fyldestgørende, vil der
især være tre Veje, man kan slaa ind paa.

a. Man kan forsøge at reducere Værdibegrebet, som jo er
det, der volder Vanskelighederne, til dets mest elementære Former.
De mest elementære Værdier vil der være mindst Strid
om, og de ville lettest kunne bringes i bestemt Sammenhæng
% --- p. 69 ---
\setcounter{footnote}{0}%
\opage{69}med objektive Forhold. Denne Vej gaar i vor Literatur Dr.
\textsc{Starcke} (Samvittighedslivet. 1894—97) og Dr. N. H. \textsc{Bang}
(Begrebet Moral. 1897).

b. Man kan gaa et Skridt videre og forsøge at eliminere
hele Værdibegrebet, stryge alle psykiske Elementer og Forudsætninger
og grunde Etiken paa Biologien som rent objektiv
Videnskab. De etiske Sætninger blive da biologiske Deduktioner.
En saadan „biologisk“ Etik er hævdet af Dr. S. \textsc{Hansen}
(Etikens Begrundelse. 1903).

c. Man kan endeligt søge at reducere Modsætningerne mellem
de forskellige Grundværdier til rent logiske Modsigelser,
som falde bort, naar der gøres Alvor af at opfatte Mennesket
som tænkende Væsen. Denne Vej gaar Professor \textsc{Kroman} (Begrebet
„Det Etiske“. Universitetsprogram Novbr. 1903).

Det er interessant, at alle tre Veje ere forsøgte af danske
Forfattere. Jeg skal nu forsøge at vise, at ingen af disse tre
Veje føre til Rom.

a. Dr. Starckes Opfattelse af Etiken har øjensynligt udvik
let sig under Indflydelse af hans Studier over primitive Samfundsformer,
Studier, der bl. A. have sat Frugt i hans fortrinlige
kritiske Arbejde over den primitive Familie. For en Forsker,
der med stor Energi har fordybet sig i saadanne Studier,
ligger det nær at anvende Forhold fra de mere primitive Samfund
paa Samfundslivets idethele. Klanen staar da ogsaa i hans
Fremstilling af „Samvittighedslivet“ som den Ramme, der omslutter
den Enkeltes Liv derved, at de Andres Agt og Ære
bliver den vigtigste Betingelse for hans trygge Livsførelse. „Moralens
Fordringer\dots{} komme til Mennesket udefra\dots{} De
sædelige Krav ere snart af den Beskaffenhed, at de fra Samfundsmagtens
Side søges gennemførte med Magt, saa den, der
overtræder dem, straffes; snart ere de henviste til den Sanktion,
som den offentlige Mening, andre Menneskers gode Omdømme
og venlige Sindelag kan give dem\dots{} I Andres Ret til at
kræve Individet til Ansvar for hans Handlinger ligger Moralens
Tyngdepunkt“. (Samvittighedslivet. I. p. 120). Denne Opfattelse
af det Etiske fører konsekvent til, at man begrænser Etikens
Indhold. Alt, hvad der ikke behøves til at tilvejebringe den
% --- p. 70 ---
\setcounter{footnote}{0}%
\opage{70}Tryghedsfølelse, som opstaar hos den Enkelte ved, at han véd
sig og sin Færd i Overensstemmelse med den straffende Myndighed
og den offentlige Mening, falder udenfor Etiken. Den
Enkelte kan, naar denne Grænse er naaet, gaa sine egne Veje.
Hele Samfundslivet staar som blot Middel for, at den Enkelte
kan føle sig tryg. Og naar dette Formaal er opnaaet, kan den
Enkelte søge sin Lykke, som han vil. --- Spørgsmaalet er, om
Individ og Samfund kunne stilles overfor hinanden paa denne
udvortes Maade. Selvom Individet mener at have affundet sig
med Samfundet, vil dog den Maade, paa hvilken det søger sin
Lykke, kunne være af den største sociale Betydning.“ Det kan
ikke være ligegyldigt, hvorledes de Enkelte føre deres indre
Liv, eftersom det er i det indre, personlige Liv, at de Værdier
blive til, som senere kunne ændre og udvikle Samfundet. Den
Enkelte har derfor selv i sine inderligste Forhold et socialt Ansvar,
som kun han selv kan gøre gældende. Han er her selv
Vogter og Dommer. Det er lagt i hans Haand, om der paa
dette Centralsted i Verden skal opstaa nye Værdier, eller om
Spirerne til saadanne skulle gaa til Grunde. Dr. Starcke lærer
nu ogsaa, at det er selve Samfundslivets Betingelser, den Enkelte
skal føle sig i Overensstemmelse med, for at han kan føle
sig tryg. Men naar dette er saa, er det klart, at Løn og Straf,
Agt og Foragt --- de eneste Sanktioner, som anerkendes i Dr.
Starckes Fremstilling --- kun kunne være opdragende Midler.
Hvorfor skal nu netop Tryghedsfølelsen være Vurderingsmotiv?
Og kan det ikke for den Enkelte, der har tilstrækkelig Magt,
være ligegyldigt, om de Andre føle sig trygge, naar blot han
selv gør det? Det elementære Vurderingsliv, der ligger til Grund
i Dr. Starckes Fremstilling, fører altsaa ikke med Nødvendighed
til de Resultater, han vil begrunde.

Dr. Bang mener, at Sprogbrugen allerede har afgjort, hvad
Moral er. „En Dom, som erklærer en Handling moralsk eller
umoralsk, udsiger, at Handlingen som almindelig vil fremme
eller hæmme Samfundets Vel, og at det er i det almene Vels
Interesse, at Samfundet gennem sine Organer, Staten eller den
offentlige Mening, fremmer eller hæmmer Handlinger af den
Art.“ (Begrebet Moral. p. 113). Hertil bemærker '1 jeg,) at, hvor%
% --- p. 71 ---
\setcounter{footnote}{0}%
\opage{71}ledes det saa end forholder sig med Sprogbrugen, saa gælder
den Begrundelse, der gives i denne Udtalelse, kun for den, der
interesserer sig for det almene Vel; for enhver Anden bliver
den kun et Tankeexperiment. Erfaring viser, at Mennesker med
større eller mindre Tilnærmelse kunne gøre deres eget Vel, ikke
Samfundets, til Formaal; for Saadanne vil en Deduktion ud fra
Samfundets Vel intet virkeligt Overbevisende have. De kunne
indrømme, at visse Fordringer konsekvent maa stilles, naar man
gaar ud fra Samfundets Vel; men de gaa selv ikke ud fra dette,
naar de ikke tvinges dertil, og tvinges kunne de ikke, naar de
have Mod og Magt nok til at trodse Staten og den offentlige
Mening. Og ifølge Dr. Bang strækker Etikens Indhold sig ikke
ud over, hvad der kan hævdes gennem Statens Love og den
offentlige Menings Kendelser. Udenfor den forholdsvis snevre
Kreds af Pligter, som disse Myndigheder med Rette fordre fyldestgjorte,
kan den Enkelte føre sit eget Liv og søge Lykken
ad sine egne Veje.

Saaledes fører en Anlæggelse af de mest elementære Synspunkter
nødvendigvis til en Dualisme eller et rent udvortes Forhold
mellem Individ og Samfund. Og hos Dr. Bang fremtræder
denne Dualisme paa en særligt paafaldende Maade, idet han
udtaler, at „de beundringsværdigste Sjælsegenskaber, det store
Menneskes Egenskaber“ maaske netop komme frem hos den,
der bryder med „Moralen“. „Vi have aldrig paastaaet, at moralsk
Handlen er den fuldkomneste Handlen“ (p. 145). „Moralitet
er en Fuldkommenhed, en Dyd, blandt andre; det er den
fundamentale sociale Dyd, men ikke den eneste, ikke den altomfattende,
og heller ikke den højeste“ (p. 190). Men hvor er
da den Maalestok, i Kraft af hvilken „Moraliteten“ sammenlignes
med andre Dyder? Der maa dog være en saadan Maalestok,
ti ellers vilde det være meningsløst at tale om højere og
lavere Egenskaber. Desværre angives denne Maalestok ikke ---
skønt det netop maa være dens Beskaffenhed, en Behandling af
det etiske Problem skulde angaa. Hvorfor er det „Moralske“
godt, og naar hører det op at være godt for at give Plads for
„højere“ og „højeste“ Egenskaber? Det er Problemets Kærne,
og det er den, jeg har forsøgt at drage frem

% --- p. 72 ---
\setcounter{footnote}{0}%
\opage{72}Hvad Beraabelsen paa den sædvanlige Sprogbrug angaar,
tager Dr. Bang den egentligt selv tilbage, da han med Beklagelse
indrømmer, at Sprogbrugen er ubestemt. Ordet «moralsk“
bruges jo ogsaa om rent individuelle Klogskabsregler, og ædel
eller fortjenstlig Handlen, der overskrider det „Pligtmæssiges“
Minimum, kaldes ogsaa moralsk Handlen (p. 30. 35). Det forekommer
mig, at naar man betragter Sprogbrugen som Autoritet,
maa man ogsaa følge dens Svingninger. Den levende Sprogbrug
er stedse paa Vandring, og den kan ledes af et ligesaa rigtigt
Instinkt, naar den bevæger sig, som naar den staar stille. Hvor
rigtigt dette Instinkt er, der leder den i dens Varieren, vil man
opdage, naar man lader Ordene fare og holder sig til Iagttagelse
af, hvilke forskellige Formaal Menneskene kunne stille sig, og
hvilken Indflydelse dette konsekvent faar paa deres vurderende
Domme. De Spørgsmaal, som den videnskabelige Undersøgelse
af de vurderende Domme faar at behandle, ere da, hvorledes
Menneskene komme til at stille sig Formaal, og hvorledes disse
Formaal bestemme deres Vurdering af deres egen og Andres
Færd. Det, at Sprogbrugen varierer paa den omtalte Maade,
viser, at allerede den almindelige Bevidsthed er opmærksom
paa Problemet i dets større Omfang\footnote[1]{I sin Bog Naturlig Ret (1907) har Severin Christensen givet et i flere Henseender interessant Udkast til en Etik, der begrænses til at formulere de elementære Betingelser for menneskeligt Samliv. Der er her gjort et Forsøg paa at isolere den mest elementære Del af Etiken og behandle den for sig. Det minder om Fichtes Forsøg paa at opstille en af alle etiske Forudsætninger uafhængig Retslære. (Se herom nedenfor Kap. XXXVII, 4). Ligesaa lidt som de i Texten omtalte Forsøg undgaar denne Opfattelse de af Værdibegrebet følgende Vanskeligheder. Der maa atter her spørges, om det ikke for den Enkelte, der har tilstrækkelig Magt, kan være ligegyldigt, hvilke Samlivets Betingelser ere. Han overtræder dem, naar han kan, og han overholder dem kun, naar han kan bruge dem til sine egne Formaal. (1913). f}.

b. Den «biologiske“ Etik, saaledes som Dr. S. Hansen har
søgt at grundlægge den, vil undvære Værdibegrebet og mener
derved at slippe for de Vanskeligheder, det medfører. Vel indrømmes
det, at al etisk Undersøgelse gaar ud paa, hvorvidt en
vis Handling er Middel til Opnaaelsen af et Formaal; men om
% --- p. 73 ---
\setcounter{footnote}{0}%
\opage{73}det til Grund liggende Formaal skal der her ikke kunne være
Tvivl; det er et Formaal, som er givet paa Forhaand, og som
Etiken ikke selv opstiller (p. 107). Det af Naturen selv satte
Formaal er Slægtens Bestaaen og Fremgang eller den størst
mulige Livsudfoldelse (p. 24).

Er dette nu virkeligt et Formaal, som er absolut nødvendigt?
Dr. Hansen indrømmer selv, at andre Formaal kunne
stilles. Hint Formaal er, siger han, „saa lidt vilkaarligt valgt
som vel muligt“ (p. 28), opstillet „med mindst Vilkaarlighed“
(p. 30); et andet Maal kan „ikke med større logisk Ret“ opstilles,
og der kan i hvert Tilfælde ikke opstilles noget Maal,
som strider mod dette (p. 234). Hvor bliver da det logisk Tvingende
af? Der skal jo kunne opstilles forskellige Formaal, og
intet enkelt har paa Forhaand den absolute „logiske“ Ret. Det
er i Virkeligheden kun dette, jeg har paastaaet, og det er deri,
at for mig Vanskelighederne ved Etikens Forudsætninger ligge.
Dr. Hansen indrømmer mig (med hvem han paa dette Punkt
mener at være „saa uenig som vel muligt“) Alt, hvad jeg ønsker.
Der maa ogsaa for ham være flere konsekvente etiske Systemer
mulige. Til denne Erkendelse er jeg kommen ad Erfaringens
Vej, støttet til, hvad Biologi, Psykologi og Sociologi have lært
mig. Og er der Noget, det sidste halve Aarhundredes Biologi
har indskærpet, saa er det jo dog især Kampen for Livet ---
en Kamp mellem levende Væsener om, hvis Liv der skal bestaa
og udfolde sig. Det er en Kamp mellem Individer indbyrdes
og mellem Grupper af Individer. Vi staa endnu midt i
denne Kamp; den føres i den menneskelige Verden saavelsom
i den øvrige levende Natur. Det er i den, de etiske Ideer
skulle godtgøre deres Gyldighed, ikke i Abstraktionernes Verden,
ikke for „Livet overhovedet“, men for aldeles bestemte,
individuelt og historisk udformede Livstrin.

Efter Dr. Hansen er det tvingende Formaal nu egentligt
givet med det bestemte Udfald af Kampen. Han indrømmer, at
han bruger Ordet „Formaal“ i overført Betydning (p. 24). Det
betyder egentligt kun en Bevægelsestendens i en vis Retning,
saa man altsaa med samme Ret kan sige, at den rullende Lavine
har sit „Maal“ i Dalen, eller Floden sit „Maal“ i Havet.
% --- p. 74 ---
\setcounter{footnote}{0}%
\opage{74}Men kan man da overalt sætte det faktiske Resultat som Maal?
Her er det netop, Værdibegrebet melder sig som uundgaaeligt.
Ti om vi kunde, vilde vi stanse Lavinen, naar den truer med
at ødelægge Byen nede i Dalen, og aflede Flodens Vand, naar
det kan bruges til at frugtbargøre vore Marker. Det Uundgaaelige
maa vi finde os i; men det staar ikke altid som Noget, der
kan blive Maal for os. Derved bliver Etik noget Andet end
Anerkendelse af den Stærkeres Ret og Underkastelse under
blinde, brutale Fakta. Her gælder det stolte Ord, som blev sagt
om Cato den Yngre, at Guderne holdt med den sejrende Sag,
Cato med den besejrede. Det Sejrende har ikke altid Værdien
paa sin Side. Erfaringen viser os baade Udvikling og Opløsning.
Hvilken af dem er egentligt Naturens „Maal“? Mefistofeles
mener, at Alt, hvad der opstaar, er værdigt til at gaa
under, og at det derfor var bedst, at Intet opstod. Hvorledes
vil man kunne gendrive ham? En konsekvent advocatus diaboli
vilde i hvert Tilfælde neppe give tabt overfor lyriske Vendinger
som f. Ex. „Livets evige Ret til at hævde sig selv“ (p. 111).
Hvorledes kan der være Tale om Ret, før Etiken er traadt i
Kraft? Den «biologiske“ Etik hævder jo udtrykkeligt, at „Maalet“ selv ikke kan begrundes etisk.

Naar Erfaringen viser os forskellige Bevægelsestendenser,
saa er der ikke Andet at gøre end at fastholde og begunstige
den Tendens, der for os har størst Værdi, og Værdi erfares gennem
Lystfølelse. Her komme vi da ind i den subjektive Verden,
man ikke kan komme udenom, hverken i Etiken eller i
Logiken, naar man ikke vil dogmatisere. Men der existerer
ingen Kløft mellem det Subjektive og det Objektive. Enhver
Tænkning over Værdier fører til at undersøge Betingelserne for
deres Opstaaen og derved Muligheden for at frembringe dem
paany og udvikle dem videre. Ingen Retning i Etiken har indskærpet
det med større Energi end Nytte- eller Velfærdsmoralen.
Ifølge denne gaar det etiske Arbejde ud paa at omsætte
de naturlige Energier i saadanne Former, i hvilke de kunne
virke til værdifulde Resultater. Menneskelig Drift og Lidenskab
virke ikke --- ligesaalidt som Lavinen og Floden --- altid i Retning
af værdifulde „Maal“; derfor kræver Etiken en Metamor%
% --- p. 75 ---
\setcounter{footnote}{0}%
\opage{75}fose af dem, hvorved de --- med Bevarelse af Energien ---
kunne virke i andre Retninger. Det er en Omsætning fra det
Reale til det Ideale, der kræves. Første Gang har Platon i stor
Stil stillet denne Opgave (i „Symposion“). Denne Omsætning
skal ske, ikke i en Fantasiverden, men i selve den virkelige
Verden. Derfor har jeg ogsaa stedse hævdet en nøje Sammenhæng
mellem Vurderingsproblemet og de andre Hovedproblemer.
„Hvad der for mig er det højeste Formaal, der bestemmer den
Værdi, jeg tillægger Alt andet, det maa dog være saaledes beskaffent,
at det kan hævdes og tilstræbes indenfor den Verden,
hvis Virkelighed kan godtgøres ved det mest konsekvente Kriterium,
jeg formaar at anvende [nemlig den videnskabeligt paaviste
Aarsagssammenhæng]. Noget kan kun godtgøre sig som
værdifuldt gennem den Maade, paa hvilken det griber ind i den
virkelige Verden. Og paa den anden Side faar selve denne
virkelige Verden større eller mindre Værdi for mig, alt eftersom
den viser sig at fremme eller hæmme en Stræben hen efter
det Formaal, der bestemmer mine Værdidomme“:).

Paa alle Omraader har Opstillingen af et Princip først og
fremmest den Betydning at fordele Bevisbyrden. Aarsagssætningens
Betydning er, at Bevisbyrden paahviler den, der vil hævde,
at en Begivenhed ingen Aarsag har. Og paa analog Maade betyder
Velfærdsprincipet, at Bevisbyrden paahviler den, der hævder,
at en Smerte er bedre end en Lystfølelse. Naar jeg slaar
Bægeret af Haanden paa den Tørstige, maa jeg godtgøre, at
det er skadeligt for hans Sundhed at drikke, eller at Vandet
var giftigt. En Smerte eller en Skuffelse kan ifølge Velfærdsprincipet
være værdifuldere end en Lystfølelse, hvis den er en
nødvendig Skærsild, eller hvis den er et medvirkende Element
i en værdifuldere Tilstand. Den etiske Betragtning støtter sig
her tilsidst til Følelsens Biologi. Derfor kan man meget vel betone
det subjektive Element i de etiske Domme-) uden at opgive
Muligheden af en Begrundelse.
% CHECK p. 75: note whose mark was not found
\footnote[1]{Det psykologiske Grundlag for logiske Domme. (Vid. Selsk. Skr. 6. Række) p. 65.}
% CHECK p. 75: note whose mark was not found
\footnote[2]{Det er undgaaet Dr. Hansens Opmærksomhed, at jeg ved „etiske Domme“ fra først af (se I, 1) mener de Domme om Godt og Ondt, Men- neskene fælde i det virkelige Liv. Det etiske Problem stilles, nåar der spørges, om saadanne Domme kunne begrundes. Videnskabelig Etik er et System af begrundede etiske Domme. Med Begrundelsen vil en nærmere Bestemmelse og videre Udvikling følge.}

% --- p. 76 ---
\setcounter{footnote}{0}%
\opage{76}c. Professor Kroman er enig med mig i, at man ikke kan
komme udenom Værdibegrebet, og at der ved alle etiske Domme,
som fældes i det virkelige Liv, medvirker et Følelseselement af
afgørende Betydning: „Udtrykket „for—Skyld“ kan med Mening
kun bruges overfor et følende Væsen“ (p. 52). Men han er
overbevist om, at Mennesket „til Syvende og Sidst“ maa forbinde
de forskellige Følelseselementer paa modsigelsesfri Maade.
Den rette Handlen, den, som den ideale Samvittighed kræver
og anerkender, vil være „den med Menneskets inderste blivende
Væsen overensstemmende, den til Syvende og Sidst modsigelsesfrie“
(p. 37). Det Etiske er ingen historisk Tilfældighed,
men „noget med Menneskets inderste blivende Væsen nødvendigvis
Sammenhængende, et i alle væsentlige Træk konstant
Udslag af det for samtlige Individer fælles Grundlag“ (p. 44).

Det er ubestrideligt, at naar der gør sig en Modsigelse
gældende mellem forskellige Interesser eller Tendenser, maa
der opstaa en Trang til at skaffe den af Vejen, enten ved et
Valg imellem Modsætningerne eller ved en Harmonisering af
dem. Saaledes virker Modsigelsen fremaddrivende. Men det
store Spørgsmaal er netop, om Modsigelsen er til Stede i det
enkelte Tilfælde, og i hvilken Grad den mærkes. Og det er
her, Erfaringen, den psykologiske og historiske Erfaring, viser
os overordentlig store Forskelligheder. Disse Forskelligheder ere
ingenlunde historiske Tilfældigheder. De maa antages at have
deres bestemte Aarsager, som det er Psykologiens og Historiens
Opgave at finde. Man har ingen Ret til at gaa ud fra, at
alle Følelsestendenser findes hos ethvert Menneske i en saadan
Grad af Udvikling, at det nu kun skulde gælde at forene dem
paa en modsigelsesfri Maade. Derfor kan jeg ikke komme saa
let over de forskellige Stadier og Standpunkter, som min Kollega.
Jeg har ganske vist søgt at godtgøre, at Lombrosos bekendte
Teori om Forbrydermennesket ikke kan begrundes ved Erfaring\footnote[1]{Smlgn. Etiske Undersøgelser (1891). p. 73—81.};
% --- p. 77 ---
\setcounter{footnote}{0}%
\opage{77}men den modsatte Antagelse, at alle de Elementer, som en ideal
Samvittighed fordrer, skulde være til Stede hos Enhver, i enhver
Race og til enhver Tid, kan heller ikke begrundes.

Selv om alle Elementer vare til Stede, vil det endnu være
af Betydning, i hvilket Forhold de staa til hverandre. Forat der
skal være en Modsigelse til Stede, --- en Modsigelse af en saadan
Art, at den kan virke som Motiv (baade Vurderingsmotiv
og Handlingsmotiv), --- maa de forskellige Elementer hvert for
sig fremtræde med stor Energi. Der kommer ingen Modsigelse
frem, naar et Element kun lyder som en svag Tone, en fjern
Klang, medens det modsatte døver Sjælens Øre med sin Kimen.
Der kræves altsaa en vis Udvikling, forat Modsigelsen kan gøre
sig gældende, og især forat den kan blive en fremaddrivende
og motiverende Kraft. Saalænge Udviklingen endnu foregaar,
hos det enkelte Individ og i Slægten, staar derfor „den ideale
Samvittighed“ som en Abstraktion, der ikke dækkes af noget
virkeligt Standpunkt. Enhver af os, ethvert Slægtled og ethvert
Folk staar paa et vist bestemt Trin af Udvikling, og af dette
bestemte Livstrin bestemmes alle Værdier, Formaal, Idealer og
Normer. Vil man bygge Etik paa Grundlag af Erfaring, kan
man ikke forbigaa de Vanskeligheder, heri ligge. Naar ere vi
ved „det Syvende og Sidste“? Og hvorlænge skal en Tendens
ytre sig og være raadende, forat den skal kunne fortjene Prædikat
af „blivende“? Her er stedse kun Mulighed for Tilnærmelse,
i Livet og i Tanken.

løvrigt tror jeg, at man ikke tør betragte Modsigelsen mellem
Elementerne i Menneskenaturen som den eneste Udviklingsbetingelse.
Ofte foregaar den biologisk-psykologiske Udvikling
som en positiv Væxt, idet et Element ubevidst og uvilkaarligt
naar en saadan Udviklingsgrad, at det faar Overhaand i
Sjælen og bestemmer den Klangfarve, hvormed de sjælelige
Elementers hele Sammenspil optræder. Hvor saaledes et enkelt
Element ved Expansion bliver raadende, spiller Modsigelsen ingen
Rolle. Og det er maaske de største Ting paa det sjælelige Omraade,
der foregaa paa denne Maade. Det er ingenlunde saaledes,
at vi altid først have en Række enkelte Værdier og Formaal,
og saa bortskaffe Modsigelsen paa en eller anden Maade.
% --- p. 78 ---
\setcounter{footnote}{0}%
\opage{78}Der er Værdier, som uvilkaarligt komme til at raade, og der er
Formaal, vi stille os, fordi vi ere voxede, uden at vi vide det.
Ikke blot kan det gaa, som i Parablen, at vi saa Sæden, og
den saa voxer, bliver til Straa og sætter Ax, medens vi sove;
men selv Sæden kan være saaet, uden at vi ane det. Som jeg
et andet Sted har sagt, gælder det ikke blot, at Mennesket
voxer med sine større Formaal, men ogsaa, at det maa være
voxet forat kunne sætte sig større Formaal. Og foruden denne
ubevidste Væxt kan der foregaa Værdi- og Motivforskydninger,
hvis Resultater ikke altid ere til at forudse. (Smlgn. Psykologi
VI C og nærværende Værk XIII, 4). Heller ikke ved saadanne
Overgange er Modsigelsen det Afgørende. En psykologisk Proces
er nu engang noget Andet end en formel logisk Tankerække,
og Mennesket er ganske vist „et tænkende Væsen“,
men dets Tanke udfolder sig paa hvert givet Trin under Indflydelse
af aldeles bestemte historiske Forhold.

Det psykologisk-historiske Grundlag, hvis Konsekvenser og
specielle Anvendelser jeg undersøger i min Etik, fremtraadte i
Stoicismen i de sidste førkristelige Aarhundreder; da opstilledes
første Gang Ideen om en almindelig Humanitetsfølelse (caritas
generis humani) som det mest omfattende Vurderingsmotiv. Uddybet
og inderliggjort ved Kristendommen og klaret gennem
den nyere Tids intellektuelle og sociale Udvikling, stiller det
for den etiske Tænkning de mest omfattende Opgaver, noget
System kan ønske sig at behandle. Til disse Opgaver hører ogsaa
den at afgøre, hvorledes man i Tanken og i Livet skal
stille sig overfor dem, hos hvem hint Grundlag ikke er udviklet,
eller som ikke anerkende det. En Paavisning af Selvmodsigelser
vil herved kunne have sin Betydning. Men Menneskets
Væsen er for sammensat og for skiftende til, at den blotte
Stræben efter at befries fra Selvmodsigelse skulde kunne afgive
et tilstrækkeligt Grundlag for det Etiske.

% --- p. 79 ---
\setcounter{footnote}{0}%

\opage{79}\chapterhead{IV.}{SAMVITTIGHEDENS TEORI.}

1. I det foregaaende Kapitel (§ 5 og 9) er Samvittigheden
nærmest beskreven som en Følelse. Men ingen Følelse er
ganske uden Sammenhæng med Bevidsthedslivet idethele, og
det er af Vigtighed særlig at fremhæve dette for at forstaa
Samvittighedens Natur.

Det er interessant, at Begrebet Samvittighed fra først af
hænger nøje sammen med Begrebet Bevidsthed. Det tilsvarende
Ord (συνείδησις) blev hos Grækerne først brugt i den almindelige
Betydning af Selvbevidsthed, Evne til Opfattelse af sine
egne indre Tilstande, før det blev brugt i den speciellere Betydning
af etisk Selvvurdering. Og det tyske Gewissen (= Mitwissen)
har ligeledes først havt en mere almindelig teoretisk
Betydning, før det fik den specielle etiske Betydning\footnote[1]{Rudolf Eucken: Geschichte der philosophischen Terminologie. Leipzig 1879. p. 175. --- Det danske „Samvittighed“ er overført fra Nedertysk (vistnok i det 14. eller 15. Aarhundrede). Det nedertyske samwitticheit havde (som Oversættelse af conscientia) Betydningen „Bevidsthed“ og Betydningen „Samvittighed“. Ældre end det er oldhøjtysk gewissen eller gewizzida, der betød Viden, Erkendelse og først senere (efter det latinske conscientia) antog Betydningerne Bevidsthed og Samvittighed. Jeg skylder sprogvidenskabelige Kolleger disse Oplysninger.}. Der
ligger i Udtrykket Samvittighed, at Mennesket har Evne til
ligesom at dele sig i to Personer, af hvilke den ene er Objekt
for den anden. Vi vises herved hen til Bevidsthedens almindelige
Karaktermærke: en sammenfattende Virksomhed, ved
hvilken de enkelte Erfaringer blive stillede i Forhold til hverandre.
Alt Bevidsthedsliv er Udtryk for en Trang til Enhed
og Sammenhæng, for en Syntese. Og i Samvittigheden ytrer
denne Ejendommelighed sig paa en særligt inderlig og koncentreret
Maade. Saa stor Modsætning der end kan være mellem
det, Mennesket erkender som Idealet, og dets egen virkelige
% --- p. 80 ---
\setcounter{footnote}{0}%
\opage{80}Villen og Handlen, saa fører Samvittigheden dog til, at Mennesket
vedkender sig begge: „Det er mit Ideal, hvor meget det
end strider mod min Færd!“ eller: „Det har jeg villet og gjort,
hvor meget det end strider mod mit Ideal!“ Samvittighedens
Opstaaen forudsætter, at der indenfor Bevidstheden gør sig en
Forskel gældende mellem et centralt Indhold, det Raadende og
Herskende, og et periferisk Indhold, de enkelte forbigaaende
Tanker, Stemninger og Villiesytringer. Hvor der har udviklet
sig en Forestilling om et Hovedformaal, et Ideal, et Forbillede,
maa den --- og de til den knyttede Følelser og Impulser ---
høre til det Centrale i Bevidstheden. Samvittigheden er da en
Form for eller en Del af det Centrales Reaktion mod det Periferiske,
en Forholdsfølelse, der faar en forskellig Karakter, alt
eftersom der er et harmonisk eller et disharmonisk Forhold
mellem det Centrale og det Periferiske i Personligheden. Der
kan være stor Modstand at overvinde, før den umiddelbare
Sammenholden og Sammenlignen af det Højeste og det Laveste,
Bevidstheden rummer, kan komme i Stand. Selverkendelsens
Helvedfart, for at bruge Kants Udtryk, kræver Opbydelsen af
stor aandelig Energi. Men den kommer i Stand efter almindelige
psykologiske Love og udtrykker kun Bevidsthedens almindelige
Natur paa en særligt udpræget Maade. Der er derfor
ingen Grund til at antage Samvittigheden for en speciel Evne
eller Sans, der skulde være indplantet én Gang for alle og hos
Alle\footnote[1]{Friburgi 1895. p. 60) støtter sig, naar han tillægger mig Antagelsen af en særegen „sensus moralis“.}. Muligheden af Samvittighed ligger i Bevidsthedens almindelige
Væsen; men om denne Mulighed udvikles, og i
hvilken Retning den udvikles, --- hvilket specielt Indhold altsaa
Samvittigheden faar, det beror paa de sociale Forhold, under
hvilke Individet bliver til og udvikler sig (se I, 2). Da Kontrastforholdet
mellem de forskellige Elementer spiller en saa
stor Rolle ved Samvittighedens Opstaaen, er det psykologisk
forstaaeligt, hvad Erfaringen synes at godtgøre, at den „onde“
Samvittighed opstaar før den „gode“: Kontrasten er naturligvis

') Jeg ved ikke, hvorpaa Victor Cathrein (Philosophia moralis,
% --- p. 81 ---
\setcounter{footnote}{0}%
\opage{81}størst, hvor der ikke blot er en Forskel, men en skarp Modstrid
til Stede. Det personlige etiske Liv bliver derfor ofte først
egentligt til derved, at man føler sig i en faktisk Strid med
det, man anser for det Rette.

Samvittighedens specielle Karakter vil bero paa, hvilke Formaal
eller Livsinteresser der for den Enkelte ere de højeste og
udgøre hans reale Selv. Herpaa beror jo ogsaa, som vi have
set, Forskellen mellem de forskellige mulige etiske Systemer
lige fra det egoistiske til det universelt-humane\footnote[1]{Samvittigheden har, naar der ses bort fra de forskellige Grundværdier, en rent formel Karakter. Men selv det rent Formelle har sin Værdi. Det er derfor, jeg i Den menneskelige Tanke (p. 241—245) har skelnet mellem formale og reale Værdiforskelligheder. Til de første høre Forskellene mellem elementære og ideelle Værdier, mellem umiddelbare og middelbare Værdier og mellem potentielle og aktuelle Værdier. I sin rent formelle Karakter er Samvittigheden en Vogter over et harmonisk Forhold mellem disse Modsætninger. Til en anden Gruppe af Værdiforskelligheder hører dels Forskellen mellem biologiske (økonomiske), intellektuelle, æstetiske, etiske og religiøse Værdier, dels Forskellen mellem individuelle, sociale og kosmiske Værdier. --- Smlgn. ogsaa min Psykologi VI C, 8 a. (1913).}. Men saa omfattende
Livsinteresser den Enkelte end har knyttet sig til, saa
føler han dog i Samvittigheden stedse sig selv som et harmonisk
eller disharmonisk Hele. Hans egentlige Jeg er det, der lever
i hine Interesser; hans ydre Jeg er det, der fremtræder i hans
virkelige Villen; og i Samvittigheden mærker han Forholdet
mellem sit egentlige og sit ydre Jeg, som jo dog begge høre
med til hans Personlighed. Fornegter han et af dem, bliver
Samvittigheden ikke til; men saa bedrager han ogsaa sig selv.
Der ytrer sig i Samvittigheden, som F. C. \textsc{Sibbern} har vist,
en dyb Selvomsorg, en aandelig Selvopholdelsestrang. Og den
ytrer sig først og fremmest deri, at man vil være ærlig mod
sig selv, --- selv om man ikke har været tro mod sig selv.

Det er ikke altid og ikke fra først af, at Samvittigheden
har denne frie, personlige Karakter. Den kan være autoritetsbunden,
og dens Opgave er da at sammenholde Autoritetens
Kundgørelser med Individets egen Villen og Handlen. Det reale
% --- p. 82 ---
\setcounter{footnote}{0}%
\opage{82}Seiv bestemmes da af Lydigheden og Pieteten, og det kommer
an paa, om disse Følelser krænkes eller fyldestgøres ved Menneskets
Færd. Den katolske Kirkes aandelige Magtstilling beror
paa, at den ligesom har forgrenet sig gennem Skriftestolene ud
i sine Tilhængeres personlige Centra. Derved tilvejebringes stor
Fasthed og Bundethed i visse Ting, men tillige stor Frihed i
andre Ting. Ti hvor Autoriteten ikke taler tydeligt, har jeg
efter den katolske Etik\footnote[1]{denne noget strengere Lære ikke den herskende i de katolske Skriftestoles Teori og Praxis. Se om dette Spørgsmaal \}. P. Gury, S. J.: Compendium theologiæ moralis\textsuperscript{6}. Lugd. 1899. I. p. 133—140.} --- i Kraft af det Princip, at en usikker
Lov ingen Lov er (lex non certo promulgata non obligat), ---
Ret til at følge, hvilken Mening jeg vil, selv den mindst sandsynlige
--- naar den blot endnu er sandsynlig! --- Man slutter
her: om jeg drukner paa 10 eller 3 Favne Vand er ligegyldigt!
Naar jeg én Gang er udenfor den tilforladelige Autoritetskundgørelse,
behøver jeg ikke at tage det saa nøje! Med Autoritetsbundetheden
følger konsekvent Probabilismen.

Anderledes hvor Individet ikke har Maalestokken udenfor
sig, men hvor Afgørelsen er lagt i dets egen Haand. Den frie
Samvittighed har ikke saa let ved at føle sig tryg; den kender
ikke den bratte Grænse mellem Sikkerhed og Usikkerhed. Den
har baade en inderligere Karakter --- idet den Enkelte her er
sin egen Dommer, og en mere omfattende Opgave --- idet den
stadigt skal søge at lære af Livet.

Det er ikke altid, at det Formaal eller den Livsinteresse,
med hvilken den faktiske Villen og Handlen sammenholdes, selv
træder frem i Bevidstheden. Der kan føles Billigelse eller Misbilligelse
overfor en Handling, uden at man kan gøre sig Rede
for, hvad det er, der er blevet fyldestgjort eller krænket. Vi

*) Victor Cathrein: Philosophia moralis. p. 152. --- Alfonso de'
Liguori, en af den katolske Moralteologis Autoriteter, indskærpede en
vis Forsigtighed ved Overvejelsen af Sandsynligheden for, at en Lov kunde
være den rette. Han fordrede, at de Muligheder, mellem hvilke Individet
havde frit Valg, ikke blot overhovedet hver for sig skulde have en alvorlig
Grund, men ogsaa at de ved udtrykkelig Sammenholden og Sammenlignen
skulde staa som lige eller omtrent (sic) lige sandsynlige. Dog er
% --- p. 83 ---
\setcounter{footnote}{0}%
% CHECK p. 83: notes paired with the marks in order (the layer misread a number)
\opage{83}have her en Analogi til saadanne Tilfælde, hvor der indtræder
en Kontrastvirkning til en Fornemmelse, der ikke selv kan godtgøres
at have været i Bevidstheden, eller som i hvert Tilfælde
har været der saa kort og flygtigt, at den ikke er bleven bemærket
og erindret\footnote[1]{Psykologi III, 6. --- Det psykol. Grundlag for log. Domme. p. 35—39; 47.}. I saadanne Tilfælde faar Samvittigheden
en hemmelighedsfuld Karakter, der især betones af dem, som
negte Samvittighedens naturlige Oprindelse og betragte ethvert
Forsøg paa at forklare dens Opstaaen ad psykologisk-historisk
Vej som «Materialisme“. Forklaringen ligger i mange Tilfælde
i, at en stadig Interesse formedelst Gentagelsen vil kunne synke
ned under Bevidsthedens Tærskel, men at den derfor ikke hører
op at virke. Der er et stadigt Vexelvirkningsforhold mellem
det Bevidste og det Ubevidste.

I saadanne Tilfælde antager Samvittigheden Karakteren af
et Instinkt. Den instinktmæssige Virken er overalt karakteriseret
ved, at den tjener et Formaal, som ikke selv er Genstand for
Bevidsthed. Den dunkle og ubetingede Maade, Samvittigheden
ofte virker paa, kan maaske kun forklares ved den Antagelse,
at nedarvede Elementer virke med. Man kan med Rette tale
om en Racesamvittighed, og den fælles Tradition vil ikke altid
være tilstrækkelig Forklaring. Men det maa dog --- selvom
man gaar ind paa Antagelsen af arvelige Elementer i Samvittigheden
--- vel fastholdes, at nedarvede Dispositioner stedse
ere ubestemte, og at den Maade, paa hvilken de udvikles, vil
bero paa de sociale Forhold og Overleveringer, i hvilke Individet
voxer op. Det enkelte Individ tilegner sig uvilkaarligt den positive
Moralitet og faar derigennem de Formaal og Livsinteresser,
der bestemme dets Samvittighedsliv. Den i Samfundet herskende
Sæd og Skik føles umiddelbart som forpligtende. Ligesom Følelseslivet
idethele taget udvikler sig under Livsforholdenes stadige,
men stille Magt\footnote[2]{Smlgn. Psykologi III, 7 og VI F, 4 c.}, saaledes kan heller ikke Samvittighedslivet
hos den Enkelte altid føres tilbage til aldeles bestemte og
enkelte Erfaringer. De „uskrevne Loves“ hemmelighedsfulde
% --- p. 84 ---
\setcounter{footnote}{0}%
\opage{84}Karakter beror tildels herpaa. Barnet hører, at Omgivelserne
billige eller misbillige dets Adfærd, at det er Genstand for
etiske Domme, og i sin store Efterligningstrang gentager det
ved Lejlighed disse Domme og optræder saaledes uvilkaarligt
som sin egen Dommer, betragter sig fra Omgivelsernes Standpunkt.
Samtidigt ledes det --- ligeledes uvilkaarligt --- til i sin
Færd at efterligne Omgivelserne; det indøver saaledes de Dyder,
der særligt forlanges i det Samfund, det lever i. Det handler
--- som Aristoteles\footnote[1]{Eth. Nic. II, 1. f|S\textasciicircum{}\textasciitilde{}} siger --- godt, før det bliver godt, ligesom
det kun ved at spille paa Cithar kan blive Citharspiller.
Handlingen gaar forsaavidt forud for Evnen. Baade gennem
passiv og aktiv Efterligning lever Barnet sig saaledes ind i
Samfundets praktiske Etik og faar derved den første Form for
Samvittighedsliv givet. Samvittighed existerer her endnu ikke
som fri, individuel Vurderingsevne, men er et Ekko, der ---
maaske dybt og inderligt, men dog udefra --- genlyder i den
Enkeltes Indre. Som det stadigt gaar med Barnet, gaar det paa
tidligere Kulturtrin --- og overalt, hvor den positive Moralitet
er eneraadende --- med Menneskene idethele. Deres etiske Ideer
og deres Livsførelse bestemmes først og fremmest gennem det
sociale Milieu. Frygt og Ærefrygt for Stammens Guder, der jo
repræsentere Fortidens Erfaringer, Agt for Sæd og Skik og for
den offentlige Mening φάτις) bestemme --- som særligt
det græske Folks ældste Historie lærer os --- paa dette Trin
Menneskenes Domme og Handlinger. Ærefrygt for Guderne afholder
f. Ex. fra at bruge forgiftede Pile (Odyss. I, 261—363),
og en Søn, der i Vrede vilde dræbe sin Fader, afholdes derfra
ved Tanken om «Folkets Tale og Menneskenes mange Haansord“:
han frygtede at blive benævnet Fadermorder! (Iliad. IX, 460 f.).
I sin ædleste Form fremtræder den etiske Færd paa dette Trin
som bestemt ved en Hensynsfuldhed, der kan stige til Undseelse
og Ærefrygt (αἴδως), og som føles ikke blot overfor
Oldinge og Mægtige, men ogsaa overfor Ligestillede, ja særligt
overfor Fattige og Lidende, hvor den bliver Et med Medlidenhed.
Paa karakteristisk Maade se vi her Sympatien som Vurderings%
% --- p. 85 ---
\setcounter{footnote}{0}%
% CHECK p. 85: notes paired with the marks in order (the layer misread a number)
\opage{85}og Handlingsmotiv faa Indflydelse paa den positive Moralitet og
forædle denne, hindre den i at blive blot Hensyn til, hvad de
Andre sige. Grækerne skelnede mellem hin ædle Undseelse
og den Skamfølelse (cuo\%6vrj), der kun er Sky for at paadrage
sig Dadel. Overgangen til egentlig Samvittighed i fri, individuel
Form betegnes ved den Fordring, som \textsc{Demokrit} opstillede:
„Hav Undseelse for dig selv, og du vil ikke komme til at
skamme dig for Nogen!“\footnote[1]{\textsuperscript{1} P. Na torp: Die Ethica des Demokritos. Marburg 1893. p. 10. 102. Smlgn. ogsaa den interessante Undersøgelse af Begreberne Undseelse og Skam hos L. Schmidt: Die Ethik der alten Griechen. I. p. 168—186.}. Sokrates's Fordring om Selverkendelse
som Betingelse for al ret Handlen var det vigtigste Skridt
til at forlægge Tyngdepunktet fra det Ydre til det Indre, og
Samvittighedens Emancipation fuldbyrdes af Stoikerne, der ogsaa
have slaaet Udtrykket Samvittighed i etisk Betydning fast. ---
Ved den antike Kulturs Opløsning og Undergang blev denne
Udvikling afbrudt, og med Middelalderens Kirkeherredømme kom
Autoritetsprincipet frem i den skarpeste og mest bevidste Form,
det nogensinde har haft. Protestantismen og den nyere Filosofi
hævdede igen i Praxis og Teori den frie Samvittigheds Sag\footnote[2]{Se mit Værk: Den nyere Filosofis Historie\textsuperscript{2}. I. p. 8—74.}.
Saalænge Samvittigheden kun er „Autoriteternes Ekko“, saa
længe bestaar der en Modsætning mellem Frihed og Samvittighed:
ti saa længe vil Beraabelse paa Samvittigheden være Et med
Opgivelse af ethvert Forsøg paa at begrunde sin Dømmen og
Handlen. Samvittighedens Fred bliver da let en Sovepude, medens
Samvittighedsfuldheden --- „Lydigheden mod det indre Du
skal!“ --- bliver en finere Art af Trældom. Forstaar man ved
Samvittighed kun et socialt Ekko, er det konsekvent, naar det
proklameres, at „det frie Fornuftmenneske“ maa være samvittighedsløst\footnote[3]{Bruno Wille: Die Philosophie der Befreiung durch das reine Mittel. Berlin 1894. p. 243 ff.}.
Men som vi have set, hænger Samvittigheden
nøje sammen med Bevidsthedens almindelige Natur. Og først
naar Selverkendelsen er vaagnet, saa at Mennesket kan blive
sig sine højeste Formaal bevidst, vil en klar og grundig Under%
% --- p. 86 ---
\setcounter{footnote}{0}%
\opage{86}søgelse af Dommes og Handlingers Berettigelse være mulig,
og ikke blot være mulig, men være en aandelig Fornødenhed.
Særligt vil, som det er vist ovenfor (III, 16), den Samvittighed,
hvis Vurderingsmotiv er den universelle Sympati, nødvendigvis
blive ført til med alle forhaandenværende Midler at undersøge
Handlingens eller Dommens Indflydelse paa Livet i os og udenfor
os. Hvor langt Undersøgelsen kan føres i det enkelte Tilfælde,
derom kan der naturligvis Intet siges paa Forhaand. Det vil
gaa her som ved al Erkendelse: der hvor den Ene stanser, vil
maaske en Anden kunne fortsætte. Begrænsede („bornerede“)
ere vi Alle mere eller mindre.

Den frie Samvittighed vil have en Tendens til at udvikle
sig til en højere Grad af udtrykkelig Bevidsthed, end den instinktive,
til positiv Moralitet bundne Samvittighed besidder. Der
vil uvilkaarligt danne sig Idealbilleder, som mere eller mindre
afvige fra de overleverede Forbilleder. I Stedet for det blotte
Instinkt, der virker uvilkaarligt, træder da en Drift med Forestilling
om et Formaal. Erfaringer og Reflexioner bestemme
Idealbilledernes Udvikling. Individet prøver at hævde sit Ideal
overfor Modstand i sit eget Indre eller udefra; derved træder
det klarere frem for ham, men der bliver tillige Mulighed for
at ændre og omforme det: Modstanden, som erfares, kan jo
undertiden erkendes som berettiget eller dog føre til fornyet
Selvprøvelse. Og selvom Mennesket af Svaghed ikke altid formaar
at hævde sit Ideal overfor Modstanden, bliver Forsøget
ikke virkningsløst; Kraften øves og kan maaske voxe sig stærk
til den næste Kamp. Det er gennem saadanne Prøvelser, at
Samvittigheden gør sine Erfaringer. I det Enkelte udfolder den
sig i Kraft af almindelige psykologiske Love. Snart er det formedelst
Lovene for Forestillingernes Association, at Tanker og
Betænkeligheder komme op; snart er det en Kontrastvirkning,
det skyldes, naar Samvittigheden vækkes. Gennem Forestillingssammensmeltninger
og Følelsesblandinger, gennem Forskydninger
og Udskydninger saavelsom gennem udtrykkelig Koncentration
af Villien i bestemte Retninger danner der sig en Resultant af
Individets Erfaringer, der bestemmer den Maade, paa hvilken
det stiller sig overfor nye Opfordringer til Dømmen og Handlen.

% --- p. 87 ---
\setcounter{footnote}{0}%
\opage{87}Uden ganske at miste sin instinkt- eller driftagtige Karakter
vil Samvittigheden kunne udvikle sig til, hvad man træffende
har kaldet praktisk Fornuft. Der vil opstaa en Trang til
Konsekvens og Totalitet i de etiske Forestillinger. Handlingerne
--- ikke blot de virkelige, men ogsaa de blot tænkte --- forfølges
i deres Konsekvenser; der danner sig Forestillinger om
alle de Væsener med Evne til at føle Lyst og Smerte, som berøres
af Handlingerne og disses Virkninger. Her vil særligt en
udviklet historisk Erfaring være af Betydning. Ved dens Hjælp
førtes den senere græske Filosofi ud over Modsætningen mellem
Hellener og Barbarer, og ved dens Hjælp førtes den nyere
Filosofi til en Humanitetsbevidsthed, der hævede sig over Klasse- og
Religionsforskelligheder. Baade historisk Erfaring og logisk
Konsekvens bidrager til Samvittighedens stadige Selvkorrigeren.
Vurderingen sker efterhaanden mere bevidst og efter bestemte
Principer. En filosofisk Etik er kun en metodisk Fremstilling
af Samvittighedens Indhold, saaledes som det former sig paa et
saadant Trin. Hvad Kant skildrede i det kategoriske Imperativ
og i den formale og universale moralske Fornuftlov, var et Resultat
af en lang historisk og psykologisk Udviklingsgang. ---

Hvor Samvittigheden virker som Instinkt, véd Individet ikke,
hvad det gør. Hvor den virker som Drift, har det en dæmrende
Forestilling derom. Men hvor den virker som praktisk
Fornuft, opstaar en klar Bevidsthed om Idealer og Regler.

2. Samvittigheden fremtræder under højst forskellige Former
og i højst forskellige Grader hos de enkelte Individer. Ikke
blot kan den snart træde frem som et halvt ubevidst Instinkt,
snart som Drift, snart som praktisk Fornuft. Men snart kan
umiddelbar Sympati, snart Pligtfølelse, snart Retfærdighedsfølelse
være Hovedelementet i den. Den kan snart ytre sig overvejende
negativt og hæmmende, snart mere positivt, dels forhøjende,
dels forbindende. (Smlgn. III, 20). Den kan snart ytre sig som
entusiastisk Hengivelse, snart som en mere roligt og stadigt
virkende Trang. Den kan snart med Inderlighed hævde et lille
Samfunds Livsinteresser, snart bestemmes ved omfattende Blik
paa store Kredses Tarv. Det vilde være umuligt at skildre
endog blot Hovedformerne for dens Fremtræden hos de enkelte
% --- p. 88 ---
\setcounter{footnote}{0}%
\opage{88}Individer. Men det har sin store Betydning at være opmærksom
paa disse Forskelligheder, da vi endnu lide af en Dogmatisme,
der skærer Alle over én Kam. Det vil netop følge af Velfærdsprincipet,
at ethvert Individ saa meget som muligt bør leve og
virke efter sin Ejendommelighed, og intet Sted maa dette da
gælde i højere Grad, end naar det gælder selve det etiske
Grundlag. Netop hvor det største Ansvar hviler paa Individet,
idet det skal hævde de højeste Interesser, det kender, vil Samvittighedens
Individualisation være en nødvendig Fordring: ti
kun formedelst den kan Individet virke med fuld Kraft. Hvor
Individet ikke har handlet efter sin fulde Ejendommelighed,
har det (som \textsc{Schleiermacher} har sagt) heller ikke været helt
aktivt. Og jo mere Individet i sin betydningsfuldeste Handlen
kan virke ud af sin inderste Ejendommelighed, desmere staar
det ikke blot som Middel, men ogsaa som Formaal; dets egen
dybeste Tilfredsstillelse ledsager da dets Handlen. ---

Man maa sikkert endog gaa et Skridt videre. Der kan
gives Mennesker, som ingen egentlig etisk Følelse besidde, og
som heller ikke behøve nogen. Det, som de kunne yde, yde
de af ganske Hjerte, uden at de anstille nogen Vurdering af
deres egne eller Andres Handlinger. De gaa fuldstændigt op i
en Virksomhed, som ganske stemmer med deres Evner og
Drifter, uden nogensinde at tvivle om dens Berettigelse og Betydning,
og uden at deres Iver kølnes. Det kan være Dyrkning
af Kunst eller Videnskab, Virksomhed i Samfundets Tjeneste
eller for deres Families Underhold. Eller de høre til de lykkelige
Naturer, der sprede Lys og Glæde omkring sig, blot ved
at være til. Om der gives mange af den Art Mennesker, er
ikke let at sige; men der er Intet til Hinder for, at der i det
Hele kan gives saadanne. Af Hjertet gøre de Lovens Gerninger
uden at have Loven, og hvad skulde Etiken have at indvende
derimod? Etiken er til for Livets Skyld, og Livet ikke for
Etikens.

Der kan saaledes være overordentligt store Forskelligheder
mellem Individerne i Henseende til Samvittighed, og det vilde
være dogmatisk her at fastslaa en bestemt Type som den eneste
rette. Disse Forskelligheder træde naturligvis især frem dér,
% --- p. 89 ---
\setcounter{footnote}{0}%
\opage{89}hvor man fjerner sig fra den positive Moralitet; ti denne har i
alle sine Former en Tendens til at betragte alle Individer som
ens. Forsøgene paa at udvikle en filosofisk Etik vidne om de
Forskelligheder, der faktisk ere til Stede, idet de enkelte Filosofer
hver for sig beskrive og fremstille den Art af Samvittighed,
som de bedst kende. En «personlig Ligning“ gør sig her gældende\footnote[1]{Smlgn. Frank Chapman Sharp: The personal equation in Ethics. (Transactions of the Wisconsin Academy. 1894—1895). --- Sharp anvender Begrebet personlig Ligning paa et andet Punkt, end jeg gør det i „Etiske Undersøgelser“.},
og den filosofiske Etiks Historie har blandt Andet den
store Betydning at gøre opmærksom paa de typiske Forskelligheder
paa det etiske Omraade. Og dog hylde de fleste Filosofer
den Antagelse, at de etiske Fordringer, den etiske Lovs Indhold
maa være ens for Alle, saa at Samvittigheden kun kan blive
en formel Evne til at anerkende og fyldestgøre Loven. Det
turde dog være, at Forskellighederne ikke blot angaa Samvittighedens
Form, Maade eller Bevidsthedsgrad, men ogsaa blive
bestemmende for den etiske Lovs Indhold baade i kvalitativ og
kvantitativ Henseende.

Hvis man ved en etisk Lov forstaar et Indbegreb af Fordringer,
der fastsættes af en ydre Autoritet eller af den absolute
Fornuft uden Hensyn til de Individer, for hvem de skulle
gælde, da er det klart, at individuelle Forskelligheder ikke kunne
faa Indflydelse paa Lovens Indhold. Det „objektivt Rigtige“
fordres --- og fordres i lige Grad af Alle. Men naar man anerkender
Velfærdsprincipet, maa man ogsaa fastholde, at den rette
etiske Lov ikke kan være funden, saalænge Individets personlige
Ejendommelighed ikke kommer til sin Ret. Hverken i Henseende
til Evners og Motivers Art eller i Henseende til deres
Grad ere de forskellige Individer ens, og en Lov, der i kvalitativ
og kvantitativ Henseende vilde fordre „det Samme“ af Alle,
vilde altsaa faktisk netop fordre et højst forskelligt Arbejde,
lægge højst forskellige Byrder paa de Enkeltes Skuldre. Skal
der være Lighed for den etiske Lov, maa der altsaa ikke fordres
„det Samme“ af Alle, og den Enkeltes Samvittighed maa
% --- p. 90 ---
\setcounter{footnote}{0}%
\opage{90}derfor stille Fordringer, som i kvalitativ og i kvantitativ Henseende
vilde være forskellige fra dem, Andres Samvittighed stiller.
Snart fordres der Andet og Mere, snart Andet og Mindre.

I kvalitativ Henseende vil man snarest indrømme dette.
Hvert Individ har --- paa Grund af sin særegne Begavelse ---
sit Kald, sin særegne Retning, i hvilken det skal arbejde. Kun
hvor Kastevæsenet i en eller anden Form bestaar, fordres der
samme Virksomhedsretning af højst forskellige Individer. Under
friere Samfundsformer erkendes det, at Menneskets Opgaver
ikke bestemmes blot ved almindelige Love og ydre Forhold,
men først og fremmest ved dets egen ejendommelige Natur.
Det var Noget, man allerede havde Blik for i den græske Oldtid.
\textsc{Sokrates} indskærpede Nødvendigheden af at kende sin Natur
og sine Evner for at kunne finde sine Opgaver. \textsc{Stoikerne}
førte denne Tanke videre: den, som ikke følger sin egen Natur,
vil ikke kunne bevare den personlige Sammenhæng og Enhed
i sit Liv; derfor skal man gøre som de gode Skuespillere, der
ikke vælge sig de største Roller, men vælge de Roller, der
passe bedst til deres Begavelse!\footnote[1]{Cicero: De officiis I c. 31. --- Cicero følger her Stoikeren Panaitios (fra det 2. Aarh. f. Kr.). Se Schmekel: Die Philosophie der mittleren Stoa. p. 41; 219.} --- Samvittigheden har her
at værne om den individuelle Ejendommelighed, men tillige at
holde Selvforblindelse ude.

Vanskeligere vil man gaa ind paa, at den etiske Lov ogsaa
i Kvantum varierer for de forskellige Individer, og dog er det
ikke muligt at se, hvorledes denne Konsekvens kan undgaas,
naar man tilstrækkeligt fastholder Forskellen mellem en etisk
og en juridisk Lov. Etisk set kan der være ydet et langt større
Arbejde og være vist en langt større Selvbeherskelse og Selvopofrelse
af et Menneske, der dog ikke naar op til det Gennemsnitsniveau,
der udefra forlanges, end af Mangen, som med
Lethed fyldestgør den „sociale“ Fordring, men som til Gengæld
intet Arbejde har gjort for at udvikle sine Evner udover Gennemsnitsniveauet,
skønt der var Mulighed for det. Og skulde
saa alligevel hin etisk set staa under denne? En saadan Absurditet
% --- p. 91 ---
\setcounter{footnote}{0}%
\opage{91}hilder enhver Etik sig i, der indtager et rent objektivt Standpunkt.
Velfærdsprincipet fordrer, som det senere vil blive vist,
som nødvendig Konsekvens, at et personligt Væsen ikke paa
noget Punkt maa behandles som blot Middel. Men det vilde
være Tilfældet, hvis Etiken havde sagt sit sidste Ord med
Formuleringen af en universel Lov, der slog en Streg over
alle personlige Forskelligheder. Den sande etiske Lov maa
individualisere, maa for den Enkelte angive det Arbejde, han
kan gøre. De Dyder og Pligter, der fordres af den Enkelte,
maa være Midler og Veje til Udvikling af hans personlige Ejendommelighed.
--- Derfor ere etiske Love heller ikke saa lette
at finde, som man ofte mener. Opgaven maa være den at udfinde,
hvorledes den Enkelte ved at følge sin Lov fyldestgør
den almindelige Lov paa bedste Maade. Og denne Opgave løses,
naar den almindelige, for Alle gyldige Lov kun angiver den
Retning, i hvilken der skal stræbes, den Type, der skal fastholdes.
Kun med dogmatisk Vilkaarlighed og ved at overse de
dybest liggende psykologiske og etiske Problemer kan man fastslaa
„Mandevilliens qvantum satis“ for Alle paa én Gang. Til
enhver given Tid vil et Samfund, gennem sine Ledere eller
gennem den offentlige Mening, fastslaa visse Egenskaber i et
vist Kvantum som dem, der kunne kræves, og som staaende i
ethvert «normalt“ Menneskes Magt. Det, der i de enkelte Tilfælde
forstaas ved Ord som Tapperhed, Selvbeherskelse o. s. v.,
skyldes de Gennemsnitserfaringer, man mener at have gjort angaaende
den Modstand, som kan overvindes. Man anvender her
et Skøn, som man stoler paa, gælder for alle Individer. Man
ser ikke, at der her gælder en personlig Ligning, der sætter
en Grænse for det almengyldige etiske Raisonnement. Hvis
Menneskene vare mere ens, vilde det være lettere at være
Moralfilosof. Men at gøre sig Sagen lettere, end den er, kan
dog umuligt være god Filosofi\footnote[1]{Smlgn. Etiske Undersøgelser. p. 51—83, hvor jeg udførligt har fremstillet og belyst den individuelle Relativitet eller den personlige Ligning i Etiken. --- I Oldtiden antydedes Læren om den personlige Ligning af Aristoteles, i den nyere Tid af Hutcheson, Jacobi, Schleier-}. Det er ingen Løsning af Pro%
% --- p. 92 ---
\setcounter{footnote}{0}%
\opage{92}blemet at skelne mellem Vurdering og Tilregnelse og overlade
til denne sidste at tage Hensyn til individuelle Forskelle i Anlæg
og Evner. Denne Udvej har Alexius \textsc{Meinong} forsøgt: „Den
moralske Vurdering i sin Almindelighed kan ikke indlade sig
paa personlige Enkeltheder, medens den mod den enkelte Personlighed
rettede Tilregnelse ikke let kan være personlig nok“\footnote[1]{1894. p. 201. (Smlgn. min Anmeldelse af dette Værk i Gottinger gelehrte Anzeigen. 1896. Nr. 4). --- Smlgn. om den juridiske og den teologiske Opfattelse mine Etiske Undersøgelser. p. 63.}.
Dette er dog en altfor juridisk Betragtningsmaade. Juristerne
have plejet at skelne skarpt mellem den almindelige Lov og
den individualiserende Straffebehandling, og mellem Dom og Benaadning.
Ogsaa den almindelige Opfattelse og den offentlige
Mening anvende denne Distinktion; ligeledes Teologerne, naar
de skelne mellem Lov og Naade. Men Loven og Dommen ere
uretfærdige, naar de ikke tage den Handlendes individuelle Forudsætninger
med i Betragtning, --- og ikke mindre ere Straffuldbyrdelse,
Benaadning og Naade uforsvarlige, naar de ikke
ere Udtryk for en højere Retfærdighed.

3. Da alle etiske Domme forudsætte Samvittigheden som
psykologisk Grundlag, er den den højeste Autoritet og Lovgiver,
og enhver anden Myndighed, af hvad Art nævnes kan, er underordnet
og afledet i Forhold til den. At ville gaa ud over sin
Samvittighed er at ville gaa ud over sig selv. Hvor jeg bøjer
mig for et andet Menneske, hvis Dom jeg har større Tillid til
end til min egen, er det kun berettiget, dersom det sker i Kraft
af min Samvittighed. Samvittigheden er ufejlbar, hvis man med
dens Ufejlbarhed mener dette, at den i hvert enkelt Øjeblik er
macher, Beneke, og i vor egen Literatur af A. S. Ørsted (Om Grænserne
mellem Teori og Praxis i Sædelæren. 1811. Optrykt i Eunomia I)
og af S. Kierkegaard (Uvidenskabelig Efterskrift. 1846. p. 376.). (Samlede
Værker VII. p. 426). („Hvor er Grænsen for det enkelte Individ i
hans konkrete Existens mellem, hvad der er Mangel paa Villie, og hvad
der er Mangel paa Evne, hvad der er Dvaskhed og jordisk Selviskhed,
og hvad der er Endelighedens Indskrænkning?“) Smlgn. allerede Holberg:
Moralske Tanker. 1744. p. 463.

') Psychologisch-etische Untersuchungen zur Werth-Theorie, Graz
% --- p. 93 ---
\setcounter{footnote}{0}%
\opage{93}den højeste Dommer. Vel maa Samvittigheden, hvor den rører
sig for Alvor, netop føre til at søge en objektiv Afgørelse ved
Hjælp af de bedste Grunde og de nøjagtigste Erfaringer, som
kunne tilvejebringes; men Anvendelsen af det objektive Kriterium
kan ikke altid gennemføres, før Øjeblikket til Handling
er kommet, og vi nødes da til at handle efter den bedste Overbevisning,
vi lige til dette Øjeblik have vundet. Der er saaledes
ikke, som man har indvendt imod mig\footnote[1]{Giacomo Laviosa: La filosofia del diritto in Inghilterra. I Torino 1897. p. 629—631.}, en Modsigelse
i at erklære Samvittigheden for ufejlbar og dog fastholde, at
den skal ledes af et objektivt Kriterium.

Men denne Ufejlbarhed er ikke af objektiv Art. Den højeste
etiske Autoritet kan selv vise sig at være paa Vildspor. Det
vilde være dogmatisk at paastaa, at Samvittigheden aldrig skuffer
os. Ved en saadan Paastand lukker man paa vilkaarlig Vis
Øjnene for nogle af Livets mest tragiske Konflikter. Den reneste
og alvorligste Overbevisning kan bero paa en fuldstændig Vildfarelse.
Et aldeles umiddelbart og ubedrageligt Kendemærke paa
Sandheden have vi ligesaa lidt paa det etiske Omraade som
paa andre Omraader-). I Samvittighedens Form træder enhver
alvorligere Overbevisning op; men den blotte Form borger ikke
for den absolute Sandhed. \textsc{Fichte} begrundede Paastanden om
Samvittighedens Ufejlbarhed ved, at det højeste Maal for os er
at handle selvstændigt, ud af vor egen Personlighed; gøre vi
virkeligt det, saa er det Fuldkomne naat, saaledes som det kan
virkeliggøres paa dette bestemte Punkt i Tilværelsen\footnotemark[1]. Men
det Tragiske er, at den inderligste personlige Overbevisning
kan vise sig at have været i Modstrid med, hvad en objektiv
Vurdering af Forholdene indskærper. Kunne de, der korsfæstede
Jesus, ikke have handlet efter bedste Overbevisning? Skulde
Kant have Ret i, at en Kætterdommer ikke kan have en god
Samvittighed? Var det ikke i god Tro, at Aristoteles forsvarede
% CHECK p. 93: note whose mark was not found
\footnote[2]{Smlgn. om Sandhedskriteriet paa andre Omraader: Psykologi V D, 3; VII B, 4. --- Den menneskelige Tanke. p. 99—108.}
% CHECK p. 93: note whose mark was not found
\footnote[3]{Smlgn. f. Ex. Appellation. 1799. p. 32.}
% CHECK p. 93: note mark without a note

% --- p. 94 ---
\setcounter{footnote}{0}%
\opage{94}Slaveriets Berettigelse, at Calvin, med Melanchtons Bifald, lod
Servet bestige Baalet, og at Sand myrdede „Landsforræderen“
Kotzebue? ---

Ikke mindre dogmatisk overhugger man Knuden fra modsat
Side. Hvor man betragter Etiken blot fra den objektive Side,
som Læren om Samfundsformer og ydre Handlinger, erklærer
man uden videre den subjektive Følelse for uberettiget overfor
de objektive Forhold og deres Krav. Etikere af saa forskelligt
Standpunkt som \textsc{Hegel} og \textsc{Bentham} mødes i Mistro til Samvittigheden.
Denne staar endog efter Hegel lige paa Springet
til at være det onde Princip i Verden, da den subjektive Overbevisning
ligesaa vel kan gaa ud paa noget objektivt Forkasteligt
som paa noget objektivt Rigtigt\footnote[1]{\textsuperscript{1} Smlgn. H egel: Philosophie des Rechts. § 139---140. --- Bentham: Principles of Morals and Legislation. ch. II. § 11---19. --- James Mill dadlede en slet Handling lige strengt, hvilket dens Motiv saa end havde været. Ros og Dadel vare for ham motiverende Kræfter; derfor maatte de ikke paavirkes af Handlingens Motiver. Stuart Mills Levned. Dansk Overs. p. 52!. Der gaves efter hans Mening intet Motiv, som under alle Omstændigheder virker paa rette Maade; end ikke de sociale Følelser; derfor er det af stadig Vigtighed at beregne Handlingens Virkninger efter Nytteprincipet. James Mill: Fragment on Mackintosh. London 1835. p. 176 f).}. --- Det er let nok, som
disse Etikere gøre, at advare mod Selvtillid og Selvraadighed
og forlange, at man skal bøje sig for en objektiv Lov. Den
Lov, vi følge, maa jo dog stedse kundgøre sig for os i Samvittighedens
Form. Det Lys, som gør os alt Andet klart, maa
tilsidst findes i os selv, og hvem borger os i det enkelte Tilfælde
for, at det ikke er en Lygtemand, vi lade os lede af?

Der foreligger her Muligheden af en Konflikt mellem de
to Principer, paa hvilke Etiken bygger, Grundlaget og Indholdet
(Vurderingsmotivet og Vurderingsprincipet). Der kan neppe gives
nogen anden Løsning deraf end den allerede antydede, at i Øjeblikket
maa Samvittigheden lydes, da der ingen højere Autoritet
kan gives. Det forudsættes naturligvis, at den taler tydeligt og
efter fornødent Overlæg. Men dertil maa føjes, at Samvittigheden
kan kontrollere og korrigere sig selv; den senere, mere
% --- p. 95 ---
\setcounter{footnote}{0}%
\opage{95}øvede og erfarne Samvittighed dømmer den tidligere. Netop
ved at følge sin bedste Overbevisning i Øjeblikket kan Mennesket
bane sig Vej til en rigtigere Overbevisning, som lærer ham, at
hvad der i Beslutningens Øjeblik stod for ham som det Bedste,
dog var urigtigt og fordærveligt\footnote[1]{Beraabelsen paa den gode Samvittighed sker ofte saaledes, at man netop spærrer sig Muligheden for en saadan Korrigeren. Kfr. Albrecht Ritschl's træffende Bemærkning: „A\textbackslash{}an kann beobachten, dass wer im Falle des Streites sein Gewissen als die höchste sittliche Instanz für sich ausspielt, und sich auf keine weitere Ueberlegung sittlicher Regeln einlässt, hiebei meistens eine Gereiztheit hervorkehrt, welche kein ganz gutes Gewissen verråth“. Ueber das Gewissen. p. 17.}. Saalænge der handles efter
bedste Overbevisning, er Menneskets Indre sundt, hvorledes
det end forholder sig med Handlingens objektive Karakter; og
denne indre Sundhed var den Spire, af hvilken den rigtigere
Indsigt kunde udvikle sig. Derfor kan ogsaa fra et etisk Synspunkt
en fordærvelig Handling, der udførtes i den Overbevisning,
at den var god, sættes over en god Handling, der udførtes
i den Overbevisning, at den var slet. I første Tilfælde
var selve Kilden ren, i sidste Tilfælde var den fordærvet. I
den personlige Inderlighed, med hvilken det, der betragtes som
sandt og godt, fastholdes og udøves, ligger det Punkt, paa hvilket
det ene Individ bedst forstaar det andet, den ene Slægt bedst
kan sympatisere med den anden. Paa de ydre Handlingers og
Livsformers Omraade ere Forskellighederne og Modsætningerne
ofte saa store, at Forstaaelse er umulig.

Kun den, der har Mod til at fejle, kan udrette noget Stort.
Villigheden til at udsætte sig for en saadan Vildfarelse, er Villighed
til at lide for Sandheden. Det er da heller ikke de
Kolde og Sneverhjertede, men de, der have Begejstring for det
Sande og Gode, der fejle paa denne Maade. Hvo Intet vover,
Intet vinder.

Men Samvittighedens Evne til at korrigere sig selv bliver
kun udviklet, naar der kan findes et bestemt Princip eller Kriterium
for Vurderingen. Et saadant kan Samvittigheden naturligvis
ikke lade sig foreskrive udefra; det maa udspringe af
Samvittighedens egen Natur. Jeg har allerede søgt at vise (se
% --- p. 96 ---
\setcounter{footnote}{0}%
\opage{96}111, 10 —15), at naar Samvittigheden bestemmes ved en uinteresseret
og universel Sympati, er Velfærdsprincipet det eneste
Kriterium, der kan være Tale om. Det udtrykker ikke Andet,
end hvad der ligger i det Formaal, som en Personlighed, der
er besjælet af Sympati, maa stille sig.

Den etiske Udvikling vil dog neppe nogensinde gaa for sig
uden atter og atter gentagne Konflikter mellem den subjektive
Overbevisning og den objektive Sammenhæng. Ligesom det er
en uendelig Opgave, der er stillet vor teoretiske Videnskab,
nemlig at forklare os Virkeligheden ved Aarsagsprincipets Hjælp,
saaledes er det ogsaa en uendelig Opgave, der er stillet os
paa det etiske Omraade, nemlig at vurdere alle Handlinger og
Livsforhold efter Velfærdsprincipets Krav.

4. I nøje Sammenhæng med Begrebet Autoritet staar Begrebet
Sanktion. Autoriteten byder eller forbyder, men det er
Sanktionen, der lader Buddet eller Forbuddet staa ved Magt.
Sanktionen bestaar i den Lyst eller Smerte, som knyttes til
Befalingens Overholdelse eller Overtrædelse, den Løn eller Straf,
man paadrager sig ved sin Handling, det Himmerige eller det
Helvede, man ved Handlingen nærmer sig til.

Kun hvor Autoriteten selv er en ydre, staar Sanktionen i
dette udvortes Forhold til Handlingen. Og i denne udvortes
Form har den ikke nogen umiddelbar etisk Betydning. Handlingens
etiske Karakter beror subjektivt set paa dens Udspring
af det inderste Sindelag, objektivt set paa dens Overensstemmelse
med Velfærdsprincipet. Hvilken etisk Betydning kunde
det have, om der dertil føjedes en Lyst eller Smerte, som
ikke udsprang af selve Handlingen? Fordi jeg handler godt nu,
hvorfor skal jeg saa i et følgende Øjeblik opnaa en Lystfølelse?
Fordi jeg handler ondt nu, hvorfor skal jeg saa i et senere
Øjeblik lide en Smerte? Der er etisk set intet Selvindlysende
heri. Hele Forestillingen om Løn og Straf hører hjemme i
Pædagogiken, ikke i Etiken. Det kan være fornødent til Opdragelse,
at Løn og Straf forøge den Magt, som Følelsen af
Handlingens Værd og Betydning ikke altid besidder i tilstrækkelig
Grad. Men det er ingen umiddelbar etisk Nødvendighed.
Naar man paastaar, at Tanken om Gengældelsen er en i sig
% --- p. 97 ---
\setcounter{footnote}{0}%
\opage{97}selv klar Tanke, saa er det, fordi man ledes af et Instinkt (til
Hævn eller til Tak), hvis etiske Berettigelse og Nødvendighed
ikke fra først af er given. Det „Selvindlysende“ er meget ofte
det Blinde. Det Uvilkaarlige i Tilskyndelsen gør, at man ikke
spørger efter Grunde, og saa mener man, at der heller ikke
behøves Grunde.

Den ydre Sanktion, som bestaar i Løn og Straf, kan som
sagt kun være opdragende. Den etiske Sanktion maa være en
indre. Den kan kun bestaa i den Handlendes Følelse af Harmoni
med sin egen højeste Overbevisning, af Overensstemmelse
mellem sit Ideal (sin Grundvillie) og sin virkelige Villen. Der
opstaar derved en indre Fred, som kan være stærkere end al
Modsigelse og Modstand udefra. Individet kan have Følelsen af,
at der raader en Kraft i det, der vilde kunne omskabe Verden,
hvis den raadede i Alle. Det føler sit Væsen udvidet og forhøjet.
De Mennesker, det beundrer højest, vilde det trygt kunne
tænke sig som Tilskuere ved sin Handling, selvom alle Motiverne
laa utildækkede for dem. Det er ikke Hovmod, som her
rører sig, men en stille Følelse af i al Ringhed at være beslægtet
med det Store. Modstykket hertil er den Følelse af Forladthed,
der kan gribe den, som indser, at hans Handlen maa
isolere ham fra Alle. Da Richard III vaagner efter sin Drøm
før Slaget, er han fortvivlet, fordi han véd, at Ingen kan elske
ham. End ikke Drømmens og Erindringens Verden kan for ham
indeslutte Billedet af Nogen, som billiger hans Færd.

En saadan indre Sanktion er ikke blot Handlingens Virkning,
men bestaar i en Følelse, som allerede kan være tilstede,
før Handlingen udføres. Den er da kun Fortsættelsen og den
fulde Udvikling af det, der førte til Beslutningen og gjorde
denne mulig. „Saligheden“, siger Spinoza, „er ikke Dydens
Løn, men er selve Dyden“. Den er af samme Art som den
Tilfredsstillelse, der følger med Fyldestgørelsen af en dybtliggende
Trang. Handlingens Kilde og Handlingens Virkning
hænge her nøje sammen.

Denne Tilfredsstillelse kan være saa stor og stærk, at alt
Andet taber sin Værdi i Forhold til den. Et stort og skønt Øje%
% --- p. 98 ---
\setcounter{footnote}{0}%
% CHECK p. 98: notes paired with the marks in order (the layer misread a number)
\opage{98}blik kan staa over et langt, men betydningsløst Liv. Herved
bliver Selvopofrelsen psykologisk forstaaelig. Der kan imellem
den Tilfredsstillelse, som føles ved at udføre en opofrende Handling,
og al anden mulig Lystfølelse være en saadan Afstand, at
denne sidste saa godt som forsvinder for Bevidstheden. Billedet
af Livet, som det former sig, naar den opofrende Handling
hører med til det, kan være saa straalende, at Livet, som det
vilde være uden denne Handling, mister al Værdi. Dette er kun
et enkelt Exempel paa en bekendt psykologisk Lov. Allerede
i det gamle buddhistiske Læredigt Dhammapada hedder det
(v. 111 —112): „Bedre end at leve hundrede Aar i Blindhed
og Spredthed er det at leve en eneste Dag i Forstaaelse og
Begejstring. Bedre end at leve hundrede Aar i Sløvhed og
Villiesmathed er det at leve en eneste Dag i Udfoldelse af
Villieskraft.“ Ogsaa Aristoteles har set dette. „Den brave
Mand“, siger han\footnote[1]{Eth. Nic. IX, 9.}, „vil gøre Meget for sine Venner og for sit
Fædreland; om det er nødvendigt, vil han endog dø for dem.
Han vil ofre Gods og Æresbevisninger, og idethele alle de
omtvistelige Goder, idet han erhverver sig det Skønne; ti han
vil hellere føle en stor Glæde i en kort Tid, end en ringe
Glæde i en lang Tid, hellere leve et Aar skønt\footnote[2]{Ved det Skønne forstaar Aristoteles det, som vi ville for dets egen Skyld, ikke som Middel til noget Andet.} end mange
Aar paa ligegyldig Vis, hellere udføre én stor og skøn Handling
end mange smaa. Saaledes gaar det jo vel dem, der dø
for Andre: de vinde for sig selv noget meget Skønt“. ---
\textsc{Immanuel} \textsc{Kant} udtrykker en beslægtet Tanke i følgende Ord\footnote[3]{Kritik der praktischen Vernunft. Kehrbachs Udg. p. 106 f.}:
„Holdes en retskaffen Mand under den største Ulykke i sit
Liv, som han kunde have undgaaet, naar han blot havde kunnet
sætte sig ud over sin Pligt, ikke oppe af den Bevidsthed, at
han har hævdet Menneskeheden i sin Person i dens Værdighed
og har æret den, saa han ikke har Aarsag til at skamme
sig for sig selv og sky Selvprøvelsens indre Skue?“

5. Hvor stor en Rolle denne indre Sanktion spiller i det
% --- p. 99 ---
\setcounter{footnote}{0}%
\opage{99}virkelige Liv, gør her Intet til Sagen. Der er maaske meget
faa menneskelige Handlinger, som ikke behøve anden Støtte end
den. Men det er af Betydning at holde fast, at den kan være
tilstrækkelig til at bære det etiske Liv. Ved dette Faktum betinges
Etikens Selvstændighed overfor Dogmatik og Metafysik,
og bliver den uafhængig af Trospostulater og Hypoteser.

Fra et etisk Standpunkt maa man nære stor Betænkelighed
ved den Maade, paa hvilken det Etiske saa ofte gøres afhængigt
af visse bestemte religiøse eller spekulative Antagelser. For det
Første ligger da den Betragtning nær, at den, der ikke længere
hylder disse Dogmer, dermed ogsaa har emanciperet sig fra det
Etiske og maaske endog vil handle konsekvent, naar han følger
det Princip: „Lader os æde og drikke, ti i Morgen skulle vi
dø!“ Men dernæst berøver man Handlingen dens egentlige
etiske Karakter ved at vende Opmærksomheden mod det, som
ligger udenfor Handlingens Væsen og Kilde, og erklære Hensynet
til Løn eller Straf for et nødvendigt Motiv.

Der gives dem, som ikke kunne fastholde det Etiskes
Gyldighed uden ved Hjælp af Troen paa en højere Tingenes
Orden, hvor alt det er fuldkomment og fuldendt, der i den
Virkelighed, vi kende, er ufuldkomment og brudstykkeagtigt.
Men Trangen til at leve i en saadan Tro maa ikke betragtes
som en almenmenneskelig Fornødenhed, selv om det er en
Trang, der har rørt sig hos nogle af de ypperste Mennesker,
der have levet.

End ikke Troen paa et Fremskridt indenfor Erfaringens
Verden bør man tillægge en absolut Værdi for Etiken. En saadan
Tro kan det, teoretisk set, have sin Vanskelighed at hævde.
Men selvom den ikke skulde staa sin Prøve, selvom den sejrende
Tendens i Verdensløbet gik Etiken imod, vilde der derfor
ikke rokkes ved de etiske Principer. Opgaverne vilde blive
andre. Der vilde blive større Anvendelse for Medlidenhed og
Resignation; men hvor det etiske Sindelag var til Stede, vilde
man holde med den besejrede Sag, skønt Guderne holdt med
Sejrherrerne. Den etiske Værdi beror ikke paa den blotte Magt.
Den, som følger sin Samvittighed, kan naturligvis kun gøre det,
fordi den er det Stærkeste i ham selv; men han gør det ikke,
% --- p. 100 ---
\setcounter{footnote}{0}%
\opage{100}fordi den er det Stærkeste i Verden. Legenden om Kristoforus,
der vilde tjene den Stærkeste, er derfor uetisk, som saa mange
Legender. Om Verdensudviklingen ender (hvis den overhovedet
ender) som en Tragedie eller som en Komedie, kunne vi ikke
vide Noget om, og i hvert Tilfælde forandrer den, der lyder
Samvittighedens Lov, ikke derfor Noget i sin Rolle. Vor Etik
er de Vejfarendes Etik (Ethica viatorum). Vi se Vejen aaben
for os, men dens Ende kende vi ikke. Derimod vide vi, at vi
kunne komme videre, selvom vi ikke skulde kunne gaa den
til Ende; og vi vide, at Stansning er Tilbagegang og Død.

I selve den Kendsgerning, at Samvittigheden er opstaaet
og har udviklet sig indenfor den menneskelige Verden, have
vi dog et Vidnesbyrd om, at der virker værdifulde Kræfter i
Tilværelsen. Der er her, selvom det er i en forsvindende Krog
af Universet, bleven en Magt til, som kan bære Livet og ikke
er bange for at udtale sin Dom om det. Vi leve derfor ikke
blot „paa Muligheder“, men paa en virkelig Grundvold, som vi
søge at gøre stedse bredere og stærkere. Det etiske Liv er en
Kamp for et Rige, som er i sin Vorden, som vil kunne voxe,
og hvis Skæbne staar i nøjeste Sammenhæng med Alt, hvad
der giver Livet Værdi og Betydning. Til denne Betragtning
trækker Etiken sig tilbage og lader de dogmatiske og spekulative
Diskussioner gaa deres egen Gang.

6. Samvittighedens Fødselsstund er det Øjeblik, da der opstaar
en Følelse, der er bestemt ved Forskellen mellem Ideal
og Virkelighed. Den opstaar, som ethvert Instinkt og enhver
Drift opstaar. Dens Dødsstund vilde være det Øjeblik, da denne
Forskel for bestandig faldt bort. Dette kan ifølge Sagens Natur
ske paa to Maader, enten ved at Virkeligheden faar Bugt med
Idealet, eller ved at Idealet faar Bugt med Virkeligheden.

Holde vi os foreløbigt til den første Mulighed, saa vil der
her igen kunne være forskellige Maader, paa hvilke Virkeligheden
saa at sige kan kvæle Idealet. Dette kan ske ved aandelig
Svækkelse, ved Blaserethed eller ved Hengivelse til rent
dyriske Drifter. Men slige Tilfælde have en individuel og speciel
Karakter og høre ind under den specielle Etik. Her hvor vi
drøfte Etikens almindelige Forudsætninger, spørge vi derimod,
% --- p. 101 ---
\setcounter{footnote}{0}%
\opage{101}om der ikke skulde kunne gives en saadan videnskabelig Opfattelse
af Virkeligheden, at der ikke længere blev nogen Plads
for Idealet.

En saadan Opfattelse foreligger efter Manges Mening i den
moderne Udviklingshypotese. I Følge denne hersker der jo i
hele Naturen, den menneskelige Verden indbefattet, en ubønhørlig
Udvælgelsens Lov, i Kraft af hvilken det, der ikke kan
lempe sig efter de givne Betingelser, gaar under, og kun det
bestaar, der passer til Forholdene. Saa synes jo den fysiske
Kraft at være sat i Højsædet, og Herredømmet i Verden at
være overgivet til Brutaliteten. Kunne vore Idealer da blive
Mere end fromme Ønsker og Suk, som ikke ændre Noget i
Sagernes Gang? Ved den naturlige Udvælgelse og i Kampen
for Livet trampes jo selv det Ædleste og Bedste ned, naar det
ikke kan eller vil sætte Haardt imod Haardt!

Det er ikke Etikens Sag at lære os Virkeligheden at kende.
Det overlader den til andre Videnskaber. Den formaar hverken
at bekræfte eller at afkræfte nogen af de Hypoteser, som rent
teoretisk have vist sig berettigede og begrundede. Men den har
naturligvis en stor Interesse af at undersøge, om en saadan
Hypotese kan forenes med de etiske Principer.

Se vi nu nærmere paa Udviklingshypotesen, da hviler
denne, naar den anvendes paa det menneskelige Liv, paa en
Grundtanke, der er saa langt fra at stride mod Etikens Forudsætninger,
at den netop maa siges at være en af disse, den
Tanke nemlig, at menneskelig Udvikling, saasnart den er mere
end blot fysisk Væxt, kun bliver mulig ved, at der vækkes
bestemte Ønsker og Drifter, der kunne virke som bevægende
Kræfter. Den største Hindring for Udvikling, ja endog for den
blotte Bestaaen, er Sløvhed og Ligegyldighed. Ved de rent
elementære Former for menneskelig Udvikling er det Nøden,
som driver fremad: men under Udviklingens Gang opstaar der
Fornødenheder, som ere rettede mod Mere end den blotte Selvopholdelse.
Kun hvor ingen Trang, ingen Drift rører sig, og
hvor ingen kan vækkes, er Haabet ude. Kun den Stræbende
kan reddes. Denne Sandhed, som idealistiske Opfattelser saa
ofte have proklameret i mere eller mindre vage Former, er ved
% --- p. 102 ---
\setcounter{footnote}{0}%
\opage{102}den moderne Udviklingshypotese kun lagt os nærmere og mere
indtrængende, om man vil: mere brutalt, paa Sinde.

Det er ikke ligegyldigt, om vi med Bevidsthed og Villie
arbejde med paa Udviklingen eller ikke. Selve Driften til etisk
Vurderen og etisk Handlen er en af de Kræfter, der virke med
til at bestemme Udviklingens Gang og Retning. Netop naar vi
antage, at denne Drift har udviklet sig efter bestemte Naturlove,
maa det være os klart, at den ikke er i Strid med Udviklingshypotesen.
Den er ganske vist i sine friere og højere Former
vaagnet sent og har endnu ikke øvet sin fulde Indflydelse. Den
er i Vorden, men netop derfor har den Fremtiden for sig.

Skønt den moderne Udviklingshypotese saa stærkt har betonet
de ydre Omstændigheders Indflydelse, saa hævder den dog
ogsaa, at denne Indflydelse i hvert enkelt Tilfælde nærmere
bestemmes og begrænses ved de indre Betingelser, med hvilke
ethvert Væsen træder ud i Kampen for Tilværelsen. Jo højere
et levende Væsen staar, des mere formaar det at gribe aktivt
ind i denne Kamp. Og Udviklingshypotesen fører endog ligefrem
over i Etiken ved at vise os, hvorledes Tilværelseskampen
paa højere Trin bliver en fælles Kamp for Bevarelsen og Udviklingen
af det menneskelige Liv.

Naar Talen er om Kampen for Livet, er man tilbøjelig til
kun at tænke paa de mest elementære og brutale Former. Men
Livet har mange Former og Trin, og Kampen for Livet faar
derfor ogsaa i det Enkelte en meget forskellig Karakter. Den,
som arbejder i Videnskabens og Kunstens Tjeneste, eller den,
som kæmper for at holde sig oppe og ikke fornegte, hvad han
anser for Ret, kæmper ligesaavel for Livet som den, der klamrer
sig til den nøgne Existens og følger det fysiske Selvopholdelsesinstinkt\footnote[1]{Smlgn. Om Grundlaget for den humane Etik. p. 30 —32. --- Etiske Undersøgelser, p. 17 —23. j}.
Det er den specielle Etiks Sag at vise, hvilken Værdi
de forskellige Livstrin og Livsformer have i Forhold til hverandre.

7. Udviklingshypotesen fører ikke blot over i Etiken; den
fører hos nogle af sine mest optimistiske Repræsentanter endog
% --- p. 103 ---
\setcounter{footnote}{0}%
\opage{103}udover Etiken. I Følge \textsc{Spencer} har den etiske Følelse kun sin
Plads i en Mellemperiode af Menneskehedens Udvikling, en Mellemperiode,
som vi naturligvis ere langt fra at have gennemløbet.
Al Øvelse bevirker, at Handlingerne kunne foretages med
mindre og mindre Modstand og Vanskelighed og med mindre
og mindre udtrykkelig Opmærksomhed og Villiesanspændelse.
Anvendes dette paa det etiske Liv, nødes vi til at antage, at
der i den menneskelige Natur efterhaanden opsamles saamegen
„organisk Moralitet“ eller uvilkaarlig Evne og Trang til rigtig
Handlen, at en særlig etisk Følelse ikke længere gør sig gældende
og heller ikke vil være nødvendig, fordi der vil være
bragt Harmoni tilveje mellem Menneskets Instinkter og det, som
Hensynet til Slægtens Velfærd fordrer. Hvad der forhen maatte
gøres til Genstand for et udtrykkeligt Pligtbud, udøves nu instinktmæssigt
og uden, at Overlæg behøves\footnote[1]{Data of Ethics. p. 128. 275.}. Den virkelige
Menneskenatur vilde altsaa da have faaet en saa ideal Karakter,
at Samvittighedens Tid var forbi.

Etiken kan i og for sig ikke have Noget at indvende imod
den Antagelse, at dens Tid en Gang vil være forbi. Hvis den
dør en lykkelig Død, det vil sige, afløses af en fuldkommen og
harmonisk Tilstand, saa er det jo et Tegn paa, at den har gjort
sin Gerning. Fuldkomne Væsener have ingen Etik og kunne
ingen have, da Ideal og Virkelighed falde sammen. Ligesom
Autoritetsforholdet (i dets forskellige Skikkelser) afløses af den
frie Samvittighed, saaledes vil denne igen kunne tænkes afløst
af en umiddelbar Leven i det Gode og Rette. Der gives (som
tidligere omtalt) maaske allerede nu enkelte Naturer, for hvis
Vedkommende dette kan siges at være sket.

Men selvom denne Mulighed maa indrømmes, ere vi dog
i Virkeligheden saa langt fjernede fra en saadan Tilstand, at det
bliver en Trossag eller et rent spekulativt Anliggende, om man
vil antage, at den nogensinde ganske indtræder, og Etiken, der
nødigt indlader sig paa vidtgaaende Spekulationer, kan ikke
lægge stor Vægt derpaa. Saalænge vi endnu ere i „Lømmelalderen“,
om ikke i „Fostertilstanden“ (for at bruge Sibberns
% --- p. 104 ---
\setcounter{footnote}{0}%
\opage{104}Udtryk), niaa man ikke ved saadanne Fremtidsfantasier svække
Opmærksomheden for de store Hindringer, der hæmme Fremskridtet.
Vi kunne styrke vort Haab og Mod ved Tanken om,
at Menneskenaturen undergaar en langsom, men stadig Forandring,
og at der kan være Grund til at tro, at denne Forandring
er en Fremgang. Men en væsentlig Indflydelse paa Behandlingsmaaden
af Etiken vil denne Antagelse ikke kunne faa.

Vi maa vel endog antage, at selve Fremskridtet vil medføre
nye Idealer og nye Opgaver. Ligesom vi ved Sammenligning
mellem det civiliserede Menneskes Etik og den Vildes
Etik finde, at der er store Omraader, som hist ere dragne ind
under det etiske Synspunkt, medens de her endnu ligge aldeles
upaaagtede, saaledes vil det sandsynligvis gaa med vor nuværende
Etik, naar den engang kan sammenlignes med et
højere Udviklingstrin. Der vil sikkert frembyde sig mange Forhold,
hvis etiske Betydning vi paa Grund af vor Haardhudethed,
vor Uvidenhed og vor Egoisme endnu ikke have faaet Blik for.
Idealet vil højne sig paany, hver Gang vi mene at være naaede
til Enden.

\textsc{Spencer} gaar (ligesom \textsc{Kant}) ud fra, at Pligtfølelsen nødvendigvis
er forbunden med en Følelse af Tvang og derfor af
Ulyst. Men i Begrebet Pligt behøver der, som tidligere bemærket
(III, 9), kun at ligge, at et mere begrænset Hensyn
underordnes et mere omfattende, uden at denne Forskel eller
Modsætning mellem et „Lavere“ og et „Højere“ behøver at
føles som Tvang. Tvangsfølelsen kan falde bort, uden at Pligtfølelsens
Tid dermed er forbi.

8. Kan man gøre Mere end sin Pligt? Kan man indlægge
sig Fortjeneste ved at overtræffe den Fordring, som etisk kan
stilles? --- Dette Spørgsmaal kan kun besvares med ja fra et
Standpunkt, der tænker sig den etiske Fordring stillet udefra
til Mennesket, enten fra en overnaturlig Autoritet eller fra andre
Mennesker; en saadan ydre Fordring eller Forventning kan man
overtræffe, men ikke Samvittighedens indre Fordring. Katolicismen
svarer ja. Den skelner mellem Bud og Raad. Ved
ikke blot at lyde Buddene, men ogsaa følge Raadene indlægger
% --- p. 105 ---
\setcounter{footnote}{0}%
\opage{105}Mennesket sig en særegen, overskydende Fortjeneste\footnote[1]{Smlgn. mit Skrift: Den engelske Filosofi i vor Tid. København 1874. p. 82 f.}. Den
populære Etik, som er tilbøjelig til at tænke sig den etiske Lov
i større eller mindre Lighed med en juridisk Lov, fastholder
ligeledes Forskellen mellem det Pligtmæssige og det Fortjenstlige.
Og det Samme er Tilfældet med saadanne Etikere, der
ere tilbøjelige til at gøre den etiske Fordring til Et med det,
som Samfundet eller den offentlige Mening kan fordre (Mill og
Bain)\footnotemark[1].

Det er ikke let at se, hvorledes man, naar man ikke bliver
staaende ved en saadan udvortes Opfattelse af den etiske Lov,
vil kunne fastholde Forskellen mellem Pligt og Fortjeneste.
Den, hos hvem Samvittigheden er tilstrækkeligt udviklet, maa
føle det som sin Pligt at gøre Alt, hvad han formaar, for at
fremme Velfærd i Verden, og naar han gaar ud over det
Maal, som Andre eller endog han selv havde stillet, saa er det
kun Tegn paa, at Maalet ikke var stillet tilstrækkeligt højt. Man maa
ikke forvexle det, som visse Mennesker vente eller fordre, med
det, som virkeligt kan naas. Selv den højeste Opofrelse, et

') Begrundelsen af Forskellen mellem præceptum og consilium sker
dog undertiden paa en Maade, der netop fører til det modsatte Resultat
af det tilsigtede. At det er ædlest (nobilius) at ofre den timelige Lykke,
begrunder Cathrein (Philos. mor. p. 208) saaledes: „Opgivelsen af ydre
Goder fjerner ikke blot mange Hindringer og Farer for Frelsen, men
yder ogsaa stadig Lejlighed til fuldkommen Øvelse i Dyd.“ --- Og det
skulde ikke være Pligt at foretage et saadant Offer, naar man er i Stand
dertil? Det maa være Pligt at virkeliggøre alt det Værdifulde, man formaar
at tænke og ville. Det maa være Pligt at indlægge sig Fortjeneste,
naar man er i Stand dertil. Naar man kan afværge „Hindringer og Farer''
af væsentlig Betydning for sig selv eller for Andre, vilde det stride mod
en udviklet Pligtfølelse at undlade det. En af mine Tilhørere, en Katolik,
som har studeret paa en teologisk Læreanstalt i Udlandet, sagde mig, at
hans Lærer i Moralteologi til sin Fremstilling af Forskellen mellem Bud
og Raad føjede en Opfordring til ikke at gøre Brug af denne Distinktion
i Praxis. --- Den hellige Teresa fandt ingen Ro i sin Sjæl, før hun
havde indlagt sig al den hun formaaede: det var hendes
„Pligt“ at indlægge sig --- Smlgn. mine Bemærkninger i
Göttinger gelehrte Anzeigen. 1896. p. 305.
% CHECK p. 105: note mark without a note

% --- p. 106 ---
\setcounter{footnote}{0}%
\opage{106}Menneske kan øve, er hans simple Pligt, naar den under de
foreliggende Forhold virkeligt er gavnlig og mulig. Den populære
Etik holder sig til et vist sædvanemæssigt Jevnmaal, er
glad, naar det naas, og beundrer, hvad der gaar ud derover.
Det kan være pædagogisk rigtigt at knytte særligt Bifald („Fortjeneste“)
til Handlinger, der gaa ud over det Sædvanlige; men
den, der udøver en saadan Handling, har dog kun det rette
Sindelag, naar han føler, at han blot har gjort sin Pligt, og at
han stod som den, der ikke kunde Andet. For den ydre Tilskuer,
der møder med sit Jevnmaal, kan det være beundringsværdigt
at se en Handling, der skyder langt ud over det. Men
en saadan Handling kan derfor meget godt hos den Handlende
selv være fremgaaet af Pligtfølelse. Det Niveau, fra hvilket de
forskellige Individer gaa ud i deres Dømmen og Handlen, ligger
ikke altid i samme Højde; og af en højt udviklet Pligtfølelse
og af en med store Muligheder begavet Natur kan der fremgaa
Handlinger, som for andre Naturer staa som Mere end
Pligt. --- Ogsaa paa dette Punkt er det Anerkendelsen af de
store personlige Forskelligheder, der afgør Striden.

Selvom man af pædagogiske Grunde vil fastholde hin
Forskel mellem Bud og Raad, vil Forskellen stedse være
relativ: hvad der paa et lavere Udviklingstrin staar som Noget,
der tilraades, men ikke fordres, er maaske paa et højere Udviklingstrin
en af de mest elementære etiske Fordringer. For
den Vilde kan det være en fortjenstlig Handling at skaane den
overvundne og afvæbnede Fjende; men det hører til de elementære
Bud i civiliserede Nationers Etik.

\chapterhead{V.}{VILLIENS FRIHED.}

1. Den etiske Vurdering opstaar fra først af som et Følelsesudbrud.
Men den har sin blivende Betydning ved sin moti%
% --- p. 107 ---
\setcounter{footnote}{0}%
\opage{107}verende, villiesbestemmende Kraft. Derved bliver den selv en
medvirkende Aarsag til Udvikling henimod den højeste Velfærd.
Det er af stor Vigtighed at fastholde dette, at Samvittigheden
selv er en Aarsag. Ti derved indtager man det rette Standpunkt
overfor det saa ofte drøftede Spørgsmaal, hvorvidt Etiken kan
gaa ind paa at anerkende Aarsagslovens Gyldighed for det
menneskelige Villiesliv. Man vil da kunne indse, at Etiken
ikke blot ikke behøver at fordre nogen Indskrænkning i Aarsagslovens
Gyldighed, men at en saadan Indskrænkning vilde
være til Skade for Etiken selv, ja endog gøre den umulig.

Det er jo en hyppig Paastand, at al sand Etik staar og
falder med den Antagelse, at den menneskelige Villie ikke, i
hvert Tilfælde ikke helt igennem, er Aarsagsloven underlagt.
Der kan, mener man, kun være Tale om moralsk Ansvar, naar
Mennesket kan begynde en Aarsagsrække absolut forfra, ɔ: kan
være Aarsag uden at være Virkning. Det er en saadan absolut
Begyndelse, man mener med det lidet heldige Udtryk „Villiens
Frihed“. Lidet heldigt er dette Udtryk allerede af den
Grund, at det kan bruges og bruges i en hel Række forskellige
Betydninger, som ofte forvexles med hverandre.

2. Anvendt paa den menneskelige Villie kan Ordet „Frihed“\footnote[1]{Dette Ord bruges nu overvejende i negativ Betydning, saa at man, saasnart man hører det, maa spørge: „Frihed for hvad?“ Denne negative Betydning er dog neppe den oprindelige. Steinthal (Allgemeine Ethik. Berlin 1885. p. 356 f.) bemærker om Ordet „frei“: „Det kommer af en Rod, som betyder pleje, skaane, vise Kærlighed, elske og gøre sig elsket. Fri er altsaa kær.“ Ordet er beslægtet med Ord som Freund, Freude, Friede.}
bruges idetmindste i 6 (sex) forskellige Betydninger.

a. Den Betydning af Ordet Frihed, som Debatten om „Villiens
Frihed“ ene har med at gøre, er den, ifølge hvilken en
„fri“ Villie ikke er Aarsagsloven underlagt, ikke er Led i Aarsagsrækken
som andre Fænomener, men er Aarsag uden at
være Virkning. Man kunde kalde Frihed i denne Betydning af
Ordet „Aarsagsfrihed“. Det er om den, at Striden staar, idet
Determinismen negter den, Indeterminismen antager den. Det,
% --- p. 108 ---
\setcounter{footnote}{0}%
\opage{108}som man er fri for, er altsaa her: Alt. At ville frit er at ville
uden Aarsag, uafhængigt af Alt, hvad der er gaaet forud.

b. Frihed kan dernæst blot betyde Fraværelse af ydre
Tvang. Her udelukkes da ikke alle Aarsager, men kun de, som
ligge udenfor den villende Personlighed. Den er fri, hvis Beslutning
ikke ved ydre Vold hindres i at gaa over til Handling.
Det er altsaa snarere Handlingsfrihed end Villiesfrihed, der her
er Tale om. Jeg har Frihed til at gaa ud af mit Værelse, naar
jeg har Nøglen i Lommen, og naar Døren ikke er spærret; i
modsat Tilfælde er jeg ufri og maa blive, hvor jeg er.

c. Frihed kan ogsaa være en Frihed for indre Tvang. Man
kalder ofte den Villen ufri, som udspringer af Smerte eller af
Frygt i Modsætning til den Villen, som udspringer af Lyst eller
Haab. Selvom Døren staar aaben, gaar jeg maaske ikke ud,
fordi jeg frygter for at blive overfaldet af en lurende Fjende.
Jeg gør i saadanne Tilfælde ikke, hvad jeg har Lyst til; jeg
føler en Hindring, som bevirker, at jeg ikke handler med Sindets
hele og fulde Tilslutning. „Du bygger dig et Hus“, siger
Augustinus, „fordi du ellers ikke vilde have nogen Bolig. Det
er da Nødvendigheden, ikke din frie Villie, som faar dig til at
ville.“ Den „frie“ Villie betegne vi i denne Betydning ofte i
daglig Tale som vor „gode“ Villie. At jeg ikke gør Noget med
min „gode“ Villie, vil sige, at jeg gør det nødigt, gør det, men
med en stærk Ulystfølelse, gør det, fordi jeg foretrækker et
mindre Onde for et større Onde. Den, som udleverer sine
Penge til Røveren, der truer ham paa Livet, gør det med velberaad
Hu og ved en bevidst Villiesakt; men han kan sige:
„jeg har villet det --- men af Tvang“ (coactus volui, som en
gammel romersk Jurist allerede træffende har udtrykt det)\footnote[1]{Smlgn. Jhering: Der Zweck im Recht. I. 2. Aufl. p. 17.}.

d. For det Fjerde kan Villiens „Frihed“ betyde Villiens
Evne, Kraft og Dygtighed. Det drejer sig her om, hvor meget
Villien kan udrette, ikke om, hvorvidt den er afhængig eller
uafhængig af, hvad der er gaaet forud. Man kan være Indeterminist,
altsaa mene, at Villien ikke er betinget ved noget Forudgaaende,
og dog antage, at denne aarsagsfrie Villie spiller
% --- p. 109 ---
\setcounter{footnote}{0}%
\opage{109}en saare lille Rolle i Verden. Og man kan være Determinist,
altsaa mene, at Villien helt igennem er betinget ved, hvad der
gaar forud, og dog antage, at denne aarsagsbestemte Villie spiller
en overordentligt stor Rolle i Verden. I den berømte Strid
mellem \textsc{Augustinus} og \textsc{Pelagius} drejede det sig snarere om
Villiens (den naturlige Menneskevillies) Kraft og Dygtighed\footnote[1]{Augustinus hævdede, at den menneskelige Villie er splidagtig med sig selv og ikke formaar at lade sig løfte af Sandheden. „Sjælen byder sig selv at ville, og den vilde ikke byde det, hvis den ikke allerede vilde! Og dog sker ikke det, den byder; men den vil heller ikke af sit ganske Væsen (ex toto), og derfor byder den heller ikke af sit ganske Væsen\dots{}. Dersom Villien var fuldstændig (plena), vilde den ikke byde, at den skulde være, ti da vilde den allerede være! . . . . Det er en Sjælens Sygdom; ti Sjælen rejser sig ikke i sin Helhed: vel støttes den af Sandheden, men den tynges af Vanen. Der er da to Villier, siden ingen af dem er fuldstændig; hvad den ene har, mangler den anden\dots{}. Det var mig selv, der vilde, og det var mig selv, der ikke vilde, og jeg sønderreves i mit Inderste.“ Confessiones. VIII, 9—lo.}
end om Villiens Forhold til den Sætning, at enhver Ting har
sin Aarsag. Det gjaldt jo om at afgøre, om Mennesket trængte
til overnaturlig Hjælp for at handle godt, --- altsaa, om den
menneskelige Natur raadede over tilstrækkelige Motiver, ikke
om der overhovedet behøvedes Motiver. Det samme Stridspunkt
genoptoges som bekendt paa Reformationstiden. I den augsburgske
Konfession (I, 18) bruges da ogsaa Ordene Kraft eller
Evne (vis) og Frihed (libertas) enstydige. --- Hvor let man kan
forvexle denne Betydning af Ordet med den førstnævnte (Aarsagsfriheden),
ses af den Brug, Indeterministerne hyppig gøre
af Ordet „Evne“. De tale da om Friheden som Evne til at
begynde absolut forfra. Men dersom de skelne mellem selve
den virkelige Begyndelse og Evnen til at begynde, saa bliver
jo Villien, den virkelige Villen, alligevel afhængig, nemlig af
sin „Evne“. Efter den strenge indeterministiske Tankegang maa
der Intet gaa forud for den „frie“ Villen, ikke en Gang Evnen.
Skal der nemlig forbindes nogen Mening med Ordet Evne, saa
betyder det de Betingelser i vor Natur, som maa være til Stede,
forat vi kunne udføre en vis Virksomhed. Evne til at ville kan
% --- p. 110 ---
\setcounter{footnote}{0}%
\opage{110}jo dog fornuftigvis kun betyde de indre Betingelser for, at en
Villen kan indtræde.

e. Meget hyppigt forstaar man ved Villiens „Frihed“ Valgfrihed,
Evne til at vælge. Et Valg forudsætter ingenlunde Aarsagslovens
Ugyldighed. Det forudsætter kun, at man har Forestillinger
om forskellige mulige Handlinger, som man overvejer,
det vil sige, sammenligner indbyrdes. Afgørelsen af, hvilken af
Handlingerne der skal udføres, beror da ikke paa øjeblikkelig
Tilskyndelse eller isolerede Affekter, men bestemmes gennem
en indre Debat mellem en Mangfoldighed af Forestillinger og
Følelser. „Frihed“ betyder her ikke Modsætningen til Nødvendighed,
men Modsætningen til Blindhed. Gennem Valgfriheden
lægger en dybere, mere sammensat Nødvendighed sig for Dagen
end i den af øjeblikkelige Affekter og Drifter fremkaldte Handling.
Den „frie“ Villie betyder her den modne, selvbevidste
Villie, som dog i Valgets Øjeblik siger: „Jeg kan ikke Andet!“
For den ydre Tilskuer kan det se ud, som om Mennesket i et
saadant Øjeblik „ligesaa godt“ kunde have villet det Modsatte,
--- at f. Ex. Luther i Worms „ligesaa godt“ kunde have udtalt
en Tilbagekaldelse. Men for den, der kendte Luthers Karakter
og Alt, hvad der bevægede sig i hans Indre, maatte det staa
som en Umulighed, at han kunde have villet anderledes, end
han gjorde. Den ydre Tilskuer kender ikke de indre Betingelser,
der gøre Udslaget. Og paa en rent udvortes Maade stiller
den populære Opfattelse sig i Regelen til Villien.

Den, som kun kender én Mulighed, kan ikke vælge, „har
intetValg“. Aarsagen til, at han kun kender én Mulighed, kan
være meget forskellig. Det kan være begrænset Erfaring, Forsømmelse
i at bruge de Erfaringer, han har gjort, øjeblikkelig
Uoplagthed, Lidenskab, Svækkelse eller ligefrem Sindssygdom.
Af saadanne Aarsager kan det ogsaa forklares, naar de forskellige
Muligheder under Overvejelsen ikke fremstille sig med den
Grad af Tydelighed, de ellers vilde have. Men alt dette gaar
forud for Valget. Om et Valg kan finde Sted, beror altsaa paa,
hvad der forud er foregaaet. Valgfriheden strider ikke mod
Determinismen, og det er ikke om den, at Striden staar.

f. Endelig kan Ordet Frihed betegne Villien som den, der
% --- p. 111 ---
\setcounter{footnote}{0}%
\opage{111}ledes af etiske Motiver. I denne Betydning af Ordet er kun
den Gode fri. Der forudsættes her en saadan aandelig Udvikling
og Øvelse, at Samvittigheden kan komme til at spille en
afgørende Rolle ved enhver Overvejelse og Beslutning. Men
dette forudsætter aabenbart igen, at der bestaar en psykologisk
Aarsagssammenhæng. Ved Opfordring eller Anledning til Handling
maa da — i Kraft af Lovene for Forestillingernes Forbindelse
med hverandre og med Følelserne --- Samvittigheden
kunne vækkes. Frihed betyder her, at visse Tanker og Følelser
ere blevne herskende og have fortrængt andre. Frihed staar
her i Modsætning til Trældom under sanselige og egoistiske
Drifters og Lidenskabers Herredømme. Undertiden kalder man
dette for den sande eller den højere Frihed. Denne Betydning
af Ordet er den ældste; den stammer fra Sokrates (eller Antisthenes?)
og den anvendes af Augustinus, Spinoza o. Fl. Den har
aldeles Intet med Striden mellem Determinismen og Indeterminismen
at gøre.

Det er kun den førstnævnte Betydning, som vi i denne
Sammenhæng have Brug for, saa betydningsfulde de øvrige
ogsaa kunne være i andre Henseender.

3. Det indrømmes vel nu af de Fleste, at det ikke er teoretiske,
men praktisk-etiske Motiver, som føre til at paastaa, at
vor Villie ikke er Aarsagsloven underlagt. Hvis man betragtede
Villien blot som Psykolog eller Historiker, vilde man neppe
falde paa at opstille en saadan Paastand. Der vilde stadigt være
mange Villiesytringer, som man ikke kunde finde nogen Forklaring
af; men deri vilde ikke ligge nogen Grund til at paastaa,
at de ingen Aarsag havde. Derimod mener man, at en
saadan Paastand er en nødvendig Forudsætning for Etiken.
Skulde det nu virkeligt forholde sig saaledes? Skulde vi være
nødte til at anerkende en saadan Disharmoni mellem vor intellektuelle
og vor etiske Natur, at vi for ikke at fornegte det
Etiskes Gyldighed maatte fornegte den Grundsætning, i Kraft
af hvilken Tilværelsen ene kan blive forstaaelig for os? --- Vi
maa idetmindste se os vel for, inden vi gaa ind paa en saa
fortvivlet Opfattelse. Det er ikke Alle, hvem det falder saa
ganske let at skyde Aarsagspostulatet bort og opstille andre
% --- p. 112 ---
\setcounter{footnote}{0}%
\opage{112}Postulater i Stedet, omtrent som man affører sig sin daglige
Frakke for at iføre sig en Selskabskjole.

Jeg vil da forsøge at godtgøre Indeterminismens\footnote[1]{Etiske“. (Universitetsprogram. Novbr. 1903.) p. 87. Men Problemet kommer her igen: Er det Centrale eller „Væsenet“ altid ens? Har „jeg selv“ frembragt det fra først af? Kan det lige let udløses til Virksomhed i hvert Øjeblik? Og er det altid godt, om det udløses?} Uholdbarhed
og at vise, at Determinismen endog er en uundværlig
Forudsætning for Etiken.

a. Indeterminismen sønderriver Baandet mellem Individet
og Slægten, ja, mellem Individet og hele den øvrige Tilværelse.
Individet staar ikke længere som et ejendommeligt Led i Tilværelsens
store Sammenhæng, men rives ud af denne netop
paa de mest afgørende Punkter. Det bliver derfor umuligt for
Indeterminismen at opfatte Tilværelsen som en Helhed. Enhver
dyberegaaende filosofisk eller religiøs Anskuelse bliver umulig.
Den eneste religiøse Anskuelse, som er forenelig med Indeterminismen,
er Polyteismen; ti ethvert Væsen, som absolut kan
begynde en Aarsagsrække, er en lille Gud (eller en lille Djævel),
et absolut Væsen, og vi faa altsaa ligesaa mange Guder
som „frie“ Mennesker. Maaske bryder man sig heller ikke saa
meget om en saadan Helhedsopfattelse. Men den her fremførte
Betragtning har dog sin Betydning, især naar der polemiseres
mod Determinismen som en ugudelig eller antireligiøs Lære.

I vor egen Literatur er Indeterminismen bleven forsvaret af Heegaard
\{Om Intolerance. 1878), Wilkens (Sociologi. 1882), Kroman
\{Vor Naturerkendelse. 1883), H. Scharling (Kristelig Sædelære. 1884).
--- Naar man (som den sidste Forfatter) mener, at der gives en „højere
Formidling“ af Determinisme og Indeterminisme, da er jeg ude af Stand
til at diskutere dette Spørgsmaal, simpelthen, fordi jeg ikke kan forbinde
nogen Mening med det. Enten gælder Aarsagsloven for Villien, eller den
gælder ikke. Hvilket Tredie kan der gives? --- Om Prof. Kromans
Standpunkt overfor dette Problem i 2. Udgave af hans Tænke- og Sjælelære
maa jeg henvise til mine Psykologiske Undersøgelser. (Vidensk.
Selsk. Skrifter 6. Række) p. 99—103. I sin seneste Udtalelse om dette
Spørgsmaal finder Prof. Kroman Muligheden for en tildels aarsagsfri
Villen ved at skelne mellem det Centrale og det Periferiske i Mennesket
(eller mellem Menneskets „Væsen“ og dets „Karakter“). Begrebet „det
% --- p. 113 ---
\setcounter{footnote}{0}%
\opage{113}Dersom man opfatter Guddommen som et absolut og almægtigt
Væsen, bliver Antagelsen af en aarsagsfri Villie hos endelige
Væsener ligefrem selvmodsigende. Og hvis man protesterer
herimod og hævder, at vi her staa overfor et „Mysterium“, saa
maa man se til, hvorledes man kan skelne mellem Mysterium
og Selvmodsigelse\footnote[1]{Arminianerne overfor de ortodoxe Kalvinister). Naar Kroman selv udtaler, at en matematisk Verdensformel er umulig, hvis Mennesket har fri Villie (p. 257), saa maa han ogsaa indrømme, at det er selvmodsigende paa en Gang at tro paa en alvidende Gud og at antage Villiens Aarsagsfrihed. Ti den alvidende Gud maa være summus mathematicus og kunne forudsige Alt med absolut Nøjagtighed.}.

b. En betydningsfuld Forvexling af to forskellige Betydninger
af Ordet Frihed begaa Indeterministerne, naar de snart udtrykke
sig saaledes, at Villien ingen Aarsag har, snart saaledes, at vi
selv ere Aarsag til vore Villiesakter. Den anden, fjerde og femte
Betydning sammenblandes her med den første. Ti naar „vi
selv“ ere Aarsag til vore Villiesakter, saa er vor Villie ikke
aarsagsfri, men har sin Aarsag i vor Natur. „Vi selv“ ere i
hvert Øjeblik ikke noget „Abstrakt“, men noget aldeles Bestemt,
nemlig de bestemte Tanker og Følelser, Anlæg, Instinkter og
Drifter, der røre sig hos os; i dem maa da Villiesakternes Udspring
søges. De to Begreber: Selvbestemmelse og Aarsagsfrihed,
der ofte betragtes som identiske, ophæve i Virkeligheden
hinanden, saasnart man med Ordet „Selv“ forbinder en bestemt
Mening.

Indeterministen svarer maaske hertil: „Jeg mener netop,
at vi i Aarsagsrækken maa stanse ved Mennesket selv. Gaa vi
videre, og lade ogsaa hans Villie [og Væsen?] have en Aarsag,
saa kommer jo den sidste Aarsag til hans Beslutning til at ligge
langt ud over ham selv; altsaa er han selv ikke Aarsag“. Hertil
er kun at svare, at paa denne Maade vilde der slet ingen
Aarsager gives i Verden. Det vilde da ikke være Lynet, der

!) Det er neppe historisk rigtigt, naar Kroman (Vor Naturerkendelse.
p. 256) mener, at den antireligiøse Følelse i Regelen har taget
Parti for Determinismen. Man kunde snarere paastaa det Modsatte; idetmindste
have de rationalistiske Retninger gerne taget Parti for Indeterminismen
(smlgn. Pelagius overfor Augustinus, Erasmus overfor Luther,
% --- p. 114 ---
\setcounter{footnote}{0}%
\opage{114}slog Manden ihjel, men de Aarsager (Luftens Elektricitet o. s. v.),
der fremkaldte Lynet; men disse vilde igen ikke egentlig være
Aarsager til Lynet --- o. s. v. i en uendelig Række. Denne Betragtning
har sin Interesse; men hvorfor anvender man den
kun overfor Villien? --- Enhver Ting og ethvert Væsen i Verden
er baade Virkning og Aarsag. Jo ejendommeligere og
rigere udstyret et Væsen er, des flere Betingelser forudsætter
det, og des flere Virkninger kan der fremgaa af det. Det vil
vise sig, at Etiken meget godt kan, ja at den maa antage denne
Betragtningsmaade.

Naar Selvbestemmelse særligt stilles i Modsætning til Inerti,
maa det vel mærkes, at netop de allerbetydningsfuldeste Villiesakter
øves i Begejstring, i Hengivelse, i Optagethed af store
Ideer og under Fordybelse i Tanker, som Individet ikke selv
med Bevidsthed har frembragt, i hvert Tilfælde ikke i selve
det Øjeblik, da Villiesakten foregaar. Det er nu engang aldrig
det nøgne og abstrakte Selv, der vil Noget. Det Selv, der vil,
er revet ud af sin Ligevægtstilstand ved Noget, som vel finder
Tilknytningspunkter i det, men hvis Virken Selvet ikke er
eneste Aarsag til.

c. Dersom Villiesakten eller en Del af den ikke er Aarsagsloven
underlagt, staar den som noget Isoleret og Tilfældigt
i Forhold til den hele Personlighed. Den bliver ikke Kød af
dennes Kød, Blod af dennes Blod. Det er ikke blot Tilværelsen,
Indeterminismen ikke kan opfatte som et Hele; heller ikke
i den enkelte Personlighed kan Sammenhængen bevares. Det er
en underlig Skæbne, Indeterminismen her frister. Den mener
at hævde Menneskets Værdighed overfor Determinismen, som
efter dens Paastand gør Mennesket til en blot Maskine, —og
saa kommer den selv til at gøre Mennesket til Noget, som er
langt ringere end en Maskine, til noget Sammenhængsløst og
Tilfældigt. Naar ét og samme Motiv, uden at der er sket nogen
Forandring i de indre eller ydre Forhold, under hvilke der
handles, snart efterfølges af én Beslutning, snart af en anden,
hvorved er saa egentligt denne „Frihed“ forskellig fra Lunefuldhed
og Tilfældighed? Hvilken Værdi vil man idethele taget
kunne tillægge den? Og hvorledes kan man under saadanne
% --- p. 115 ---
\setcounter{footnote}{0}%
\opage{115}Forudsætninger vedblive at tale om, at et Menneske har en vis
Karakter?

d. Den etiske Dom om min Handlen bygger paa, at Handlingen
virkeligt er min Handling. Derfor er den etiske Dom
kun klar og bestemt, hvor den psykologiske Sammenhæng mellem
Motiver og Beslutning ligger klart for Dagen. Jo mindre
min Handling kan forstaas ud fra Kendskabet til min Karakter
og mine Forhold, des snarere vil jeg blive betragtet som utilregnelig,
og des mindre vil jeg kunne drage mig selv til Ansvar.
Ved at opgive Aarsagssammenhængen i Villieslivet opgiver
jeg netop Kendemærket paa Tilregneligheden!

Den juridiske Betragtningsmaade stemmer her med den
etiske. Straffeloven betragter oprørt Sindsstemning og Mangel
af Overlæg som formildende Omstændigheder, Gentagelse, Tilbagefald
og klart Overlæg som skærpende Omstændigheder. Til
Grund for denne Betragtningsmaade ligger den Opfattelse, at en
forbryderisk Villiesakt maa straffes des strengere, jo større Sammenhæng
den viser sig at staa i med Menneskets hele Karakter.
Uagtet den juridiske Dom væsentligt har med den ydre
Handling, ikke med det indre Sindelag at gøre, er det dog nødvendigt
for Dommeren at skaffe sig Indsigt i Handlingens Motiver
og Udspring. Og det ikke blot for at kunne bestemme Straffens
Grad, men ogsaa for at kunne have en fuldstændigt sikker
Forvisning om, at Handlingen virkeligt er udøvet. Jo bedre
Dommeren forstaar, hvorledes Handlingen nødvendigt udsprang
af de foreliggende indre og ydre Forhold, des mere kan han
være sikker paa, at Individet virkeligt har udøvet Handlingen\footnote[1]{Smlgn. min Psykologi VII B, 4. --- Anselm v. Feuerbach (Aktenmässige Darstellung merkwiirdiger Verbrechen. Giessen 1829. II. p. 502 f.) hævder, at en Dommer ikke bør fælde en Dødsdom, naar han ikke har forstaaet, hvorledes Handlingen er bleven til.}.

e. Naar flere nyere Forfattere (Heegaard, Kroman) mene,
at Postulatet om Aarsagsfrihed kan indskrænkes til det mindst
Mulige, til „en meget lille Størrelse, en Ubetydelighed“, saa er det
ikke godt at forstaa, hvilken Betydning det kan have, at en
saa ringe Del af Villiesakten er aarsagsfri, da det jo dog er
Villiesakten som Helhed, den etiske Dom gælder. Man ser, at
% --- p. 116 ---
\setcounter{footnote}{0}%
\opage{116}Indeterminismen, der først opløser Tilværelsens Sammenhæng
og dernæst Personlighedens Sammenhæng, ender med ikke engang
at opfatte den enkelte Villiesakt som noget Helt. Visse
p. Ct. af dens Elementer skulle være aarsagsbestemte, andre
aarsagsfrie. Vi faa ikke at vide, om den aarsagsfrie Bestanddel
er lige stor ved alle Villiesakter af samme Individ og ved Villiesakter
af forskellige Individer, ikke heller, hvorledes den, der
drager sig selv til Ansvar, endsige den, der skal dømme Andre,
kan opdage, hvor stor Bestanddelen er, og om den er der i
det enkelte Tilfælde. Men ogsaa afset herfra staar denne Antagelse
som noget højst Forunderligt.

Dersom det, som man har sagt\footnote[1]{Kroman: Vor Naturerkendelse. p. 246.}, er „Determinismens
Akilleshæl“, at Individet efter enhver lav Handling skulde have
Ret til at sige: „Det var nødvendigt, altsaa er det ikke min
Skyld!“ --- saa er det ikke godt at indse, hvorledes den Indeterminisme,
der kun antager en ubetydelig Del af Handlingen
for aarsagsfri (og forsaavidt ogsaa kan kaldes, og kalder
sig selv, en relativ Determinisme), kan undgaa at blive ramt
af den samme Vanskelighed, saafremt denne idethele er tilstede.
Individet vil jo da med Rette sige: „Kun en Tusindedel af min
Handling kan jeg gøre for, da kun saa Meget af den er aarsagsfrit!
Hvorfor drages jeg da til Ansvar for hele Handlingen?
Jeg har ikke begaaet noget Mord, men kun Tusindedelen af et
Mord!“ --- Mest konsekvent vilde det dog være, at man indrømmede
Umuligheden af at bestemme, hvorvidt et Menneske
var tilregneligt for en enkelt bestemt Handling. Denne Konsekvens
drager \textsc{Kant} et enkelt Sted, idet han udtaler, at „hvor
stor en Del af vor Karakter der er ren Virkning af Friheden,
og hvor stor en Del der kan tilskrives den blotte Natur og
Temperamentets uforskyldte Fejl eller lykkelige Beskaffenhed,
kan Ingen udgrunde“; derfor „forbliver Handlingernes egentlige
Moralitet ganske skjult for os“. (Kritik der reinen Vernunft\footnotemark[1] p.
551). --- Hin --- beskedne Indeterminisme hindrer sig netop ved sin
Beskedenhed i at opnaa, hvad den vil opnaa, og den prøver
% CHECK p. 116: note mark without a note
% --- p. 117 ---
\setcounter{footnote}{0}%
\opage{117}paa den umulige Ting at faa en centnertung Vægt til at hænge
i et Spindelvæv. Dersom Ansvaret virkeligt staar og falder med
Aarsagsfriheden, saa maa vi antage, at den Handlende ligesaa
godt kunde have undladt hele Villiesakten, hele Handlingen.

Den Art Beskedenhed indtræder i Teoriernes Historie gerne
paa det Punkt, hvor deres Livskraft og Selvtillid er forsvunden,
og hvor man allerede har faaet Fjenden ind i sin egen Lejr
uden at være sig det bevidst. De gamle Paastande holde sig
da en Tid endnu, ligesom rudimentære Organer; men det er
en Illusion, naar man mener, at de virkeligt udøve nogen
Funktion.

f. Ordene Ansvar, Skyld og Tilregnelighed ere, som saa
mange andre etiske Udtryk, overførte fra det juridiske til det
etiske Omraade, eller rettere, skrive sig fra en Tid, da Forskellen
mellem disse to Omraader endnu ikke ret gjorde sig gældende.
At jeg drages til Ansvar eller betragtes som tilregnelig
for en Handling, vil sige, at jeg staar som den, hvem Løn eller
Straf for Handlingen tildeles. Hvorfor Handlingen belønnes eller
straffes, er et Spørgsmaal for sig. Det er, som allerede vist
(IV, 4) aldeles ikke selvindlysende, at en Handling skal belønnes
eller straffes. Men selv om vi gaa ind paa Gengældelsesteorien
og betragte det som selvindlysende, er det aldeles ikke
nødvendigt, at den Villie, der drages til Ansvar, skal være absolut
første Aarsag, begynde en Aarsagsrække absolut forfra.
Gengældelsen fordrer kun, at der er Nogen, man kan holde sig
til, og det er der, naar Handlingens Udspring af en menneskelig
Villie er konstateret. Det primitive Gengældelsesinstinkt gaar
endda ikke saa langt tilbage; det tilfredsstilles, naar blot et
Menneske af samme Familie eller Stamme som Gerningsmanden
lider for Handlingen, og det fordrer Tilfredsstillelse, hvad enten
der foreligger en virkelig villet Handling eller ikke. Forestillingen
om individuel Skyld og individuel Tilregnelighed har langsomt
arbejdet sig frem. Men Fordringen om, at den individuelle
Villie ingen Aarsag skal have, skriver sig fra en metafysisk
Spekulation, som hverken det primitive Instinkt eller den etiske
Skyldfølelse indlader sig paa.

Hvorfor gaa vi da i vor etiske Skyldfølelse ikke længere
% --- p. 118 ---
\setcounter{footnote}{0}%
% CHECK p. 118: notes paired with the marks in order (the layer misread a number)
\opage{118}tilbage end til Villien? Hvorfor holder Etiken op her og overlader
Alt, hvad der gaar forud, til Psykologien? --- Dette Spørgsmaal
skal her søges besvaret fra et rent etisk Standpunkt. Det
juridiske Ansvar kan først behandles, naar vi i Læren om Staten
komme ind paa Begrundelsen af Statens Straffemyndighed.

Etisk set kan Følelsen af Skyld, Ansvar eller Tilregnelighed
kun betyde, at jeg føler min Handling underlagt Samvittighedens
Dom. Den indre Sanktion (IV, 4) træder i Kraft. Forudsætningen
herfor er, at jeg er mig bevidst virkeligt at have
villet Handlingen. Jeg tilregner mig den, fordi den ikke er mig
noget Fremmed, men uden min Beslutning ikke vilde være i
Verden. Jeg har Skylden, fordi jeg har villet. Naar jeg nu fornemmer
Modstriden mellem den Villiesakt, jeg vedkender mig,
og det, som jeg erkender for det Rette, opstaar Angerfølelsen,
en Følelse af indre Disharmoni, Uværdighed og Selvforagt, der
kan stige til den højeste aandelige Smerte. Den opstaar som en
Kontrastfølelse, fremkaldt ved Modsætningen mellem min virkelige
Villen og min Anerkendelse af Idealet. Denne Følelse opstaar
under de antydede Forhold med psykologisk Nødvendighed.

Ogsaa fra et deterministisk Standpunkt kan man meget
godt definere Anger som „Misfornøjelse med sig selv over ikke
at have handlet anderledes“\footnote[1]{Kro man: Vor Naturerkendelse. p. 243.}. En saadan Misfornøjelse er kun
en anden Form for det levende Ønske, der i Angerens Øjeblik
opstaar om, at man dog havde handlet anderledes, et Ønske,
der, som jeg strax skal vise, har sin store praktiske Betydning.
Et saadant Ønske opstaar naturligt, naar man i Besindelsens Øjeblik
underkaster sin Villen og Handlen en alvorlig Prøvelse og
finder, at den ikke kan bestaa for denne. Men af dette Ønske
og den derved vakte Misfornøjelse kan man ikke slutte, at man
i Handlingens Øjeblik ligesaa godt havde kunnet handle helt
anderledes, end man virkeligt gjorde. Hos Mange opstaar der
maaske den Illusion; men de datere da blot den dyrekøbte Erfaring,
de have vundet gennem Fejl og Anger, tilbage til et
tidligere Øjeblik\footnote[2]{Psykologi VII B, 5 b (Slutning). --- En udførligere Paavisning af de Maader, paa hvilke Forestillingen om en aarsagsløs Villen kan opstaa, har jeg givet i mine Psykologiske Undersøgelser. p. 92—99.}.

% --- p. 119 ---
\setcounter{footnote}{0}%
\opage{119}Men med den psykologiske Forklaring af Angeren er endnu
ikke dens etiske Værdi og Begrundelse given. Angeren er en
Smertefølelse, og det er kun fra et absolut asketisk Standpunkt,
at Smerte i og for sig er noget Godt. Hvorfor skal jeg ikke
søge saa hurtigt som muligt at glemme min forkastelige Handling,
ligesom jeg naturligt søger at holde ubehagelige Minder
borte? Da Hamlet har vakt sin Moders Samvittighed, siger
hun: „O Hamlet, du har søndersprængt mit Hjerte!“ Og han
svarer:

\begin{quote}
Saa kast den slette Halvdel bort, og lev\\
des renere da med den anden Halvdel!
\end{quote}

Hvorfor mener Etiken, at det ikke kan gaa saa rask, som Hamlet
mener, men at Angerens Smerte har sin store Betydning?

Det er ikke, fordi den hylder Principet „Lige for Lige“,
eller „Øje for Øje, Tand for Tand“. Gengældelseslæren er ikke
nogen etisk Lære. Etiken kan ikke i Angeren se „Guds frygtelige
og uudsigelige Vrede“, som \textsc{Melanchton} kalder den, ikke
heller, med Katolikerne, se en Bod eller Soning i den. --- En
mere indgaaende Kritik af Gengældelseslæren, end der indeholdes
i de allerede fremførte Bemærkninger, vil der blive en bedre
Lejlighed til ved Undersøgelsen af Statens Straffemyndighed.
Her vil jeg indskrænke mig til et enkelt Punkt, som har særlig
Interesse i denne Sammenhæng. --- I Følge Gengældelseslæren
maatte Angeren (hvor den er den eneste Straf) være
stærkest hos den, der har øvet den største Forbrydelse. Men
det er den faktisk ikke, og dette finder sin psykologiske Forklaring
deri, at Angeren beror paa et Kontrastforhold mellem
Idealet og den virkelige Villen, et Kontrastforhold, der naturligvis
ikke kan blive til, uden hvor Forestillingen om Idealet bliver
levende. Jo mere levende og udviklet denne Forestilling er, des
flere Handlinger ville blive betragtede i dens Lys, og des skarpere
ville de blive bedømte. Derfor er Angerfølelsen ofte stærkest
hos de reneste og bedste Karakterer. Der er en psykologisk
Sandhed i Fortællingen om Dronning Dagmar, som ikke fik Ro
i Døden, fordi hun havde snørt sine Ærmer en Søndag. Mod
en snehvid Baggrund ses Pletter, der ellers ikke vilde mærkes.

Psykologisk set fremtræder Angeren --- ligesom fra en anden
% --- p. 120 ---
\setcounter{footnote}{0}%
\opage{120}Side set Pligten (III, 9) --- som Udtryk for Personlighedens
Stræben efter at bevare sin Enhed og Sammenhæng. Angeren
forudsætter, at Bevidstheden om det Gode er vakt, hvad den
maaske ikke var i Handlingens Øjeblik, eller ogsaa, at den er
bleven forhøjet i Art eller Omfang eller begge Dele, idet et
højere etisk Stade er naaet. Der føles Medlidenhed med Handlingens
Offer, og man føler Værdien af de Interesser, som ere
krænkede. Angeren er da allerede et Tegn paa, at det Gode er
ved at sejre i Mennesket. Men Angerens Smerte voldes ved,
at Bevidstheden om den virkelige Villen og Handlen staar i
den skarpeste Strid med Bevidstheden om Idealet. Derfor siger
\textsc{Dante} (Skærsilden XXXI), at Beatrice, d. v. s. den ideale Tanke,
var Aarsag til den Smerte, „Angerens Nelder“ voldte ham.
Trods denne indre smertelige Modsætning, der viser Mennesket
dets egen Sjæl i et helt andet Lys end før, stræber det at bevare
Forbindelsen med sin Fortid. Det kan hverken slippe det
nye Jeg eller fornegte det gamle Jeg: de udgøre begge dets
eget Jeg. Det genkender --- trods det nye Jeg --- ærligt og
uforbeholdent sig selv som Aarsag til den tidligere Handlen.
Gennem denne stadige Stræben efter at identificere sig med sig
selv og faa det gamle og det nye Jeg til at dække hinanden
ved at prøve paa at tænke dem sammen, føles Kontrasten saameget
des stærkere. Kun i Form af Smerte kan Personligheden
her hævde sin Enhed. Men Dækningsforsøgene kunne føre til
en Lutring, en Sammensmeltning, hvorved det nye Jeg sejrer,
men tillige beriges med den Erfaring, Mennesket har gjort om
sin Grænse. Bevidstheden om Idealet bliver saameget des inderligere,
som Livets Modsætninger ere prøvede i en Dybde, til
hvilken de, hvis Udvikling mere forløber i lige og ubrudt Linie,
ikke naa ned\footnote[1]{Denne Udvikling belyses ved de i min Psykologi VI B, 2 d og e udviklede Love for Følelsers Forbindelse og Sammensmeltning og kunde selv dér være anført som Exempel.}. Hvad der er erhvervet gennem en saadan Skærsild,
har en Fasthed, som ikke let vindes paa anden Vis. Den
Energi, som Forblindelsen forhen tog i sin Tjeneste, kan nu
gennem en Metamorfose komme til at virke i en helt anden
% --- p. 121 ---
\setcounter{footnote}{0}%
% CHECK p. 121: notes paired with the marks in order (the layer misread a number)
\opage{121}Retning. Paa denne Maade vidner Angeren om Personlighedens
Kausalitet og om den psykiske Energis Bestaaen. Allerede
\textsc{Brochner} har (i sit Skrift Det Religiøse i dets Enhed med det
Humane. 1869. p. 157. 172) udtalt denne Opfattelse af Angeren:
„Det forudgaaende Liv slutter sig gennem Angeren sammen
med det nuværende til en kontinuerlig, intensiv Livshelhed“.
„Under Kraftens Forvandling er her, ligesom paa det
Fysiskes Gebet, Kraftsummen blivende“\footnote[1]{Smlgn. allerede Giordano Bruno (Den nyere Filosofis Historie\textsuperscript{2} I. p. 141 f.).}.

Kun i den udtrykkelige Bevidsthed om den tidligere Tilstands
Beskaffenhed har man Tegnet paa, at man er ude over
den. Slappes Villien til at holde Modsætningen mellem det
Gamle og det Nye ud, indtræder Metamorfosen og den faste
Grundlæggelse af den nye Karakter ikke. Den, der ikke kan
taale at tænke paa sin tidligere Adfærd, --- den, der ikke kan
gennemgaa Selverkendelsens Skærsild, --- er endnu hildet i sin
Fortid, ligesom den Sindssyge ikke er fuldt helbredet, naar han
ikke vil indrømme, at han har været sindssyg. Der er naturligvis
her kun Tale om den Dom, den Enkelte holder i sit eget
Indre; det er ikke sagt, at nogen Anden har Ret til at gøre
hans Regnskab op for ham. --- Hvor Grænsen ligger mellem
den nødvendige Selverkendelse og den sygelige Rugen over,
hvad der ikke kan være anderledes, hører til de allervanskeligste
Spørgsmaal i Livet; ingen almindelig Formel kan belære os
derom. Melankolske Naturer komme let --- i Kraft af Følelsens
Expansion'\footnote[2]{Psykologi VI F, 4.} --- til at anse sig for mere skyldige, end de ere.
Derved lammes Virksomheden, og Driften til Fremgang kues\footnote[3]{Der ligger ofte en egoistisk Sentimentalitet i den rugende Anger og den smaalige Selvanklage. Goethe fortæller (i en af sine Samtaler med Eckermann) om en lille Dreng, der ikke kunde beroliges over en lille Fejl, han havde begaaet. „Det var mig,“ siger Goethe, „ukært at bemærke dette; ti det vidner om en altfor øm Samvittighed, der vurderer det egne Selv saa højt, at det ikke vil tilgive det Noget. En saadan Samvittighed fremkalder hypokondre Mennesker, naar den ikke opvejes af en stor Virksomhed.“ I Wilhelm Meisters Lehrjahre (IV, 12) hedder det:}.

% --- p. 122 ---
\setcounter{footnote}{0}%
\opage{122}Ti den etiske Betydning af den indtrængende Bevidsthed
om Handlingens Forkastelighed, som Angerens Smerte dels forudsætter,
dels vækker, er den at virke ansporende og fremaddrivende.
Der fremkaldes en Opmærksomhed for Ting, man
forhen gik blind og haard forbi. Der rejser sig en Drift efter
at naa et højere Stade, en Higen, som naturligt udspringer af
den pinlige Modsigelse, som føles. Kun ved at være Motiv er
Angeren af etisk Natur. Den er Middel for fremtidig Udvikling.
Etiken ser fremad, ikke tilbage, eller rettere, den ser kun tilbage
for saameget des bedre at kunne se fremad. Etiske Domme
ere, som flere Gange bemærket i det Foregaaende, fra først af
Følelsesudbrud overfor Synet af eller Erindringen om en Handling.
Dette vedblive de at være, hvor vel de end begrundes.
Men en saadan Doms Udtalelse er kun berettiget, hvor den kan
være Aarsag, Motiv til fortsat eller ændret Villen. Hverken
overfor os selv eller overfor Andre have vi Ret til at dømme,
naar Dommen ikke kan virke motiverende. En stor Del af de
Vanskeligheder, Mange finde ved at forene Etik og Determinisme,
skriver sig fra den gammeltestamentlige Maade, paa hvilken
de opfatte Begreber som Ansvar og Tilregnelighed. Kun
den bør drages til Ansvar, hos hvem Regnskabsopgørelsen kan
blive Motiv --- afskrækkende eller opmuntrende Motiv. Hvis
den populære Opfattelse mener noget Andet med Ordene Anger,
Ansvar, Tilregnelighed o. s. v., saa er det uholdbare Begreber,
den har dannet.

I Ønsket om at have handlet anderledes og i Smerten ved
ikke at have gjort det anerkender man maaske Noget, som i
sin Tid gjorde sig gældende under Overvejelsen, men den Gang
ikke kunde faa Magt. Vi se her den store Betydning af den
indre Debat før Beslutningen. Selvom de Drifter og Tilskyndelser,
der havde Retten paa deres Side, maatte bukke under, har
denne Kamp dog maaske ikke været forgæves. De kunne have
gjort Ligevægten mindre stadig, selvom de ikke ganske kunde
% CHECK p. 122: note whose mark was not found
\footnote[1]{„Die Eigenliebe låsst uns sowohl unsere Tugenden als unsere Fehler viel bedeutender, als sie sind, erscheinen“. --- Smlgn. ogsaa mine Bemærkninger i Rousseau og hans Filosofi. p. 13 f. ---}
% --- p. 123 ---
\setcounter{footnote}{0}%
\opage{123}overvinde Modstanden mod Bevægelse i den Retning, i hvilken
de pege. Men det Arbejde, de have gjort, behøver ikke at gaa
til Spilde. De stige op igen, naar Sindet er blevet roligere, og
ved Kontrasten kunne de nu maaske opnaa den Magt, de forhen
manglede, eller de faa Hjælp og Supplering ved nye Tilskyndelser
og Drifter, der paa deres Side ikke alene vilde være
i Stand til at bestemme Villien.

g. Dersom Aarsagsloven ikke gjaldt paa det aandelige Omraade,
vilde den etiske Stræben være haabløs. Kun hvor der
hersker Lovmæssighed, kan min Villie gribe saaledes ind, at
der udrettes Noget. At vi kunne gribe ind i den ydre Natur og
faa den til at adlyde vore Formaal, muliggøres ved, at vi kende
Naturens Love og vide, hvilke Betingelser vi maa skaffe til
Veje for at opnaa, hvad vi ville. Paa lignende Maade ere vi
stillede overfor den menneskelige Natur. Kun hvis der hersker
Lovmæssighed, kunne vi ændre den paa de Punkter, hvor den
strider mod vore Idealer. Hvad det gælder om, er at skaffe
Motiver af den rette Art og Styrke, og dette maa ske i Kraft
af de psykologiske Naturlove. Men hvad vilde al mulig Anstrengelse
nytte, dersom et og samme Motiv under samme Forhold
snart efterfulgtes af én Beslutning, snart af en ganske anden?
Hvad der ingen Aarsag har, kan jeg ikke forberede; overfor det
staar jeg som overfor det Tilfældige og Lunefulde. Min fremtidige
Villen raader jeg kun for, dersom der er et Aarsagsforhold
mellem min nuværende og min fremtidige Villen. Da kan
det, jeg nu saar som et Frø, voxe til en kraftig Plante.

Vi forstaa nu, hvorfor Ansvaret ikke gaar længere tilbage
i Aarsagsrækken end til Villien. Det er, fordi det er den, det
gælder at forandre. Hvad der ligger forud for Villiesakten, interesserer
os etisk set kun, forsaavidt det har Indflydelse paa
Villien. Kun indirekte drages vi derfor til Ansvar for vore Tanker
og Følelser, nemlig forsaavidt de ere villiesbestemmende
Kræfter. Det Omfang, i hvilket de skulle drages ind, vil naturligvis
være meget forskelligt i de forskellige Tilfælde.

Af Intet kommer Intet. Denne Sætning maa Etiken for sin
egen Skyld lade gælde. Det kommer an paa at vække og nære
Ønsker og Drifter af rette Art hos mig selv og Andre. Hvor
% --- p. 124 ---
\setcounter{footnote}{0}%
\opage{124}langt den Enkelte i det enkelte Tilfælde kan naa, derom kan
kun Erfaring og Forsøg overbevise ham.

h. Endnu bør med nogle faa Ord den underlige Paastand
belyses, at Deterministen konsekvent maa staa som blot Tilskuer
ved Livet, baade ved sit eget og ved Andres Liv. Som om
man ikke kan føle Glæde eller Sorg ved Livet og Lyst til at
gribe ind i det, fordi man antager, at det underligger bestemte
Love! Dersom mit Ansvar slet ikke giver mig Tid til at tænke,
kan jeg naturligvis ikke være Determinist. Men saa kan jeg
heller ikke være Indeterminist. Ti det bliver man ogsaa kun
ved at tænke. Indeterminismen er ligesaavel en Teori som Determinismen,
og der er dem, for hvem det kræver den største
intellektuelle Anstrengelse at sætte sig ind i, hvad Indeterminismen
egentligt mener. Gaar jeg i hvert Øjeblik op i Praxis,
kan jeg ikke teoretisere. Agerdyrkeren faar ingen Tid til at
studere Kemi, naar han uafbrudt skal pløje og harve. Det er
naturligvis ikke Alle, for hvem det ligger at studere Psykologi
og Historie. Der kan være Naturer, hos hvem den stadige Interesse
for at finde Aarsagerne til sine egne og Andres Handlinger
kan svække den praktiske Interesse. Men det er da kun
den almindelige Modsætning mellem teoretiske og praktiske Naturer,
vi staa overfor. Enhver maa her se til, hvor megen Plads
han giver Teorien i sin aandelige Husholdning. Det er naturligvis
uetisk at optræde som blot Forsker, hvor det gælder om
Handling. Hvem vilde, om han end er nok saa ivrig Fysiker,
studere Lysvirkningernes Natur, naar han stod ved en af sine
Kæres Dødsleje? Men vil man derfor sige, at Fysik og Kærlighed
udelukke hinanden? Jeg kan (hvad enten jeg er Indeterminist
eller Determinist) ikke studere Villiens Psykologi i samme
Øjeblik, som jeg skal handle; men fordi jeg ikke paa samme
Tid kan gøre to forskellige Ting, behøve disse to Ting dog derfor
ikke logisk at stride mod hinanden. Jeg kan ikke paa én
Gang staa paa Benene og paa Hovedet; men jeg kan (maaske)
gøre det Ene først, det Andet senere.

En Forudsætning ligger der ganske vist til Grund for hele
den etiske Betragtning, hvad enten den er deterministisk eller
indeterministisk, nemlig den, at Tilværelsen ikke er fuldendt,
% --- p. 125 ---
\setcounter{footnote}{0}%
\opage{125}færdig eller fuldkommen, men at der er et Arbejde at gøre for
menneskelig Villen. Den menneskelige Villie maa være en af
de Aarsager, paa hvilke Verdensudviklingen beror. Hvilken
Vægt den har i Forhold til andre Aarsager, kan kun Erfaring
lære; men den maa ikke være Nul, hvis en Etik skal være
mulig. Enhver Etik, der staar paa Virkelighedens Grund, er en
Lære om, hvorledes Udviklingstendenser, der i mere elementære
Former ytre sig i det ydre Naturliv, kunne føres videre i menneskelig
Aand. Verdensløbet maa for en Del foregaa gennem
menneskelig Villie, og ethvert Forsøg paa en afsluttende Verdensanskuelse
maa ogsaa tage dette Moment med i Betragtning\footnote[1]{Smlgn. nedenfor VII, 1 og 4 og XXXI, 3, saavel som mine Filosofiske Problemer (Universitetsprogram Novbr. 1902). III, 5 og IV, 2.}.

\chapterhead{VI.}{DET ETISKE ONDE.}

1. Hvad det etiske Onde er, følger allerede af, hvad der
tidligere er udviklet (III, 6; 10; 20). Det bestaar i en mere
eller mindre bevidst Isolation af det enkelte Øjeblik fra Individets
Livshelhed, og af det enkelte Individ fra Slægtens Livshelhed.
Deraf følger ikke blot en Modstand mod det, Individets
eller Slægtens Velfærd kræver, men ogsaa en Slappelse af Energien
og en Opløsning af Sammenhængen i Individet eller i
Slægten. --- Vi gaa her lidt nærmere ind paa de psykologiske
Aarsager til denne Isolation.

To Aarsager forklare det tragiske Misforhold, som betegnes
ved Begrebet det Onde. For det Første Udviklingens sporadiske
Karakter, som bestaar i, at de enkelte Evner i den enkelte
Personlighed og de enkelte Individer indenfor Slægten ikke udvikle
sig samtidigt, ejheller lige stærkt og hurtigt, men saaledes,
at snart det ene Element, snart det andet optræder med den
% --- p. 126 ---
\setcounter{footnote}{0}%
\opage{126}største Energi; derved bliver der Mulighed for, at de ikke
komme til harmonisk at indordne sig i den større Helhed, men
endog gøre bestemt Modstand derimod. Under to Former kan
denne Disharmoni opstaa. Det kan være, at det enkelte Element
(den enkelte Evne hos Individet, det enkelte Individ i Slægten)
med Energi fortsætter en Tendens, som forud har udviklet sig.
I Kraft af en Art Inerti vedbliver det at bestaa og virke som
forhen, uagtet der nu fra Helhedens Side maa stilles andre
Fordringer til det end før. Eller det kan være, at det enkelte
Element uvilkaarligt varierer, slaar ind paa nye Retninger, netop
fordi Livet rører sig saa kraftigt i det, at det maa have rigeligere
Afløb end forhen, --- men at der nu ikke viser sig nogen
Mulighed for at finde en Sammenhæng, i hvilken det Nye
kunde indgaa og vinde Værdi. Hvad der fra først af syntes at
gøre et helt nyt Livstrin muligt, bliver nu et Brud paa Livets
Sammenhæng, noget Forvredet og Forvrænget. --- Hverken den
konserverende eller den varierende Tendens kan undværes; men
begge indeholde Mulighed for Brud med Livets Sammenhæng
og Krænkelse af den Grundvillie, som Vurderingen forudsætter.
De ere Kræfter, der uvilkaarligt ville det Gode, men ofte virke
det Onde. De kunne føre til Isolation, saa at Delen sætter sig
i Helhedens Sted, fordi den ikke vil ændre sin Tilstand ---
hvad enten denne nu væsentligt er Fortsættelse eller stadig
Varieren.

Den anden Aarsag ligger i, at det Uvilkaarlige gaar forud
for det Vilkaarlige, det Ubevidste eller dunkelt Bevidste forud
for det klart Bevidste. Den Tanke, der kan lede, skal først
vindes, og vindes ofte kun gennem Vildfarelse. Man bliver ofte
først klog af Skade. Den Lov, siger Aischylos, har Zeus sat,
at Visdom kun vindes ved Smerte (jid-\&ei Individet
væver sig ind i Livet uden at vide, i hvilken Sammenhæng det
kommer ind. Og selve den første Tanke opstaar uvilkaarligt:
Overgangen fra det Uvilkaarlige til det Vilkaarlige maa naturligvis
ske uvilkaarligt. De første Formaal, Mennesket sætter
sig, sætter det uden Overlæg:

% CHECK p. 126: centred line: heading or verse
\begin{center}/ffc\textasciicircum{}J'\end{center}

% --- p. 127 ---
\setcounter{footnote}{0}%
\opage{127}Hvordan den første Kundskab vakt kan blive,

Véd ingen Mand, vèd ej, hvorfra den vælder,
Og hvordan først hans Lyster kom til Live.

Som Honningdriften hos en Bi sig melder,

\begin{quote}
Kom de til ham; hin Villiens Førstetanke\\
Fortjener ingen Ros, og Last ei heller.\\
(Dante: Skærsilden XVIII).
\end{quote}

Og dog kan denne „Førstetanke“ blive afgørende for hele
Livets Retning. Den bliver Udgangspunktet for Tankelivet, det
første Associationscentrum, der bestemmer de følgende Tankers
Gang og Art. Men der er ingen Sikkerhed for, at den Retning,
som hermed gives, stemmer med de højeste Livsinteresser.

Isolation og Blindhed ere da de to Træk, det Onde, naar
vi se det efter dets Mulighed, frembyder. Det Onde er Trægheden,
der ikke vil ændre Tilstand, Tungheden, der ikke vil
lade sig løfte, Sneverheden, der ikke vil lade sig udvide. Det
er den fulde Modsætning til det Vurderingsmotiv, som i det Foregaaende
blev forudsat: ti Sympatien er en Trang til Meddelelse,
til at gøre sig til Et med Andre, til at aabne sig for Verden. ---
Den græske Etik lagde mest Vægt paa Forblindelsen, der vækker
Overmod og Frækhed (ὕβρις). \textsc{Homer} og \textsc{Tragikerne}
opfattede Syndens Udspring som Bedaarelse (ἄτη), \textsc{Sokrates}
ligefrem som Uvidenhed. Den nyere Tids Etik fremhæver
mest den af Træghed fremgaaende Isolation. Saaledes \textsc{Spinoza}
og J. G. \textsc{Fichte}. Denne sidste Tænker viser i sin interessante
psykologiske Udvikling\footnote[1]{nyere Filosofis Historie\textsuperscript{2}. I. p. 325---327. --- F. C. Sibbern har i sin Psychologie (1856) p. 119 og i sin Moralfilosofi (1878) p. 7---11 anvendt sin Lære om al Udviklings sporadiske Gang paa Problemet om det Onde.}, hvorledes der af Trægheden som Modstand
mod at gaa ud over den givne Tilstand og Retning kan
opstaa Fejghed og Falskhed, idet man hellere opgiver sin Frihed
end optager Kampen, og idet man, naar der alligevel skal

*) J. G. Fichte: Sittenlehre. 1798. p. 262 ff. Se min Bog Den nyere
Filosofis Historie\textsuperscript{2}. II. p. 150 f. --- For Spinoza var Synden en aandelig
Afmagt (Tract. polit. II, 7. 11. 20), der dannede Modsætningen til den
Maade, paa hvilken Selvhævdelsestrangen formaar at føre Menneskets
Udvikling frem fra lavere til højere Former (Ethica. III, 6 ff.). Se Den
% --- p. 128 ---
\setcounter{footnote}{0}%
\opage{128}kæmpes, hellere bruger List end ærlige Vaaben. For begge
Tænkere er det etiske Onde Modsætningen til den aandelige
Kraft og Frihed. Denne aandelige Frihed hang for dem nøje
sammen med Samfundslivet, og derved blev det Onde indirekte
ogsaa Brud paa den sociale Harmoni, Krænkelse af den fælles
Frihedsstræben. Men paa det Standpunkt, paa hvilket vi her
have stillet os (III, 8---11), maa det etiske Onde direkte karakteriseres
ved sin antisociale Karakter, ved Fraværelsen af de
Motiver og Egenskaber, der gøre aandeligt Fællesskab muligt,
ikke blot af dem, der gøre aandelig Frihed mulig.

Naar Bevidstheden om Modsætningen til et højere Livstrin
(hvilket man dog naturligvis ikke anerkender som et højere)
opstaar, og det nu en Gang indledede Stadium dog fastholdes,
stiger Trægheden til Trods. Kontrasten virker her fæstnende og
æggende og forøger Modstanden. Standpunktet udformes til et
System, der kan fastholdes med streng Konsekvens. Naar et
saadant isolerende Stade kaldes ondt, sker denne Vurdering ud
fra et andet etisk System end det, som den, der staar paa hint
Stade, selv anerkender. Ethvert Standpunkt hævder sin Ret til
at være, sætter sin Vurdering imod al anden Vurdering. Her
have vi da i den skarpeste Form den Kamp mellem etiske Systemer,
som vi tidligere (III, 13) betragtede fra et rent teoretisk
Synspunkt. Det er den store Kamp i Verden, hvor det, der fra
den ene Side stemples som Træghed, Isolation, Forblindelse og
Trods, fra den anden Side erklæres at være Udslag af naturlig
og berettiget Selvhævdelse eller maaske endog Spirer til nye
Fremskridt. I denne Kamp er det ikke blot Tanke, der staar
mod Tanke, men Lidenskab, der staar mod Lidenskab, og Villie,
der staar mod Villie. Til Grund for al etisk Vurdering ligger
jo en Grundvillie, en formaalsættende Virksomhed. Hvor Vurdering
staar mod Vurdering, er det det aandelige Livs mest
fundamentale Kræfter, der have rejst sig mod hverandre.

Den filosofiske Etik staar anderledes overfor denne Kamp
end den teologiske Etik. Den lærer af Psykologien og Historien,
hvorledes Standpunkter og Systemer blive til, og hvorledes de
successivt udfolde sig af Elementer, der have deres berettigede
Plads i Livets Sammenhæng. Denne successive Udfoldelse af
% --- p. 129 ---
\setcounter{footnote}{0}%
% CHECK p. 129: notes paired with the marks in order (the layer misread a number)
\opage{129}de Standpunkter, der forkastes og skulle overvindes, kan den
gennemførte teologiske Etik ikke antage. Den anerkender ikke
den simple Sandhed, at det Uvilkaarlige gaar forud for det
Vilkaarlige, og at Delene ikke altid slutte sig sammen til et
Hele. Synden er for den Hovmod, --- ligesom Lydighed er
Indbegrebet af al Dyd. Naar Mennesket isolerer sig, løsriver
sig fra Kilden til alt Godt, saa er Aarsagen at søge i dets Stolthed.
Hvoraf Stoltheden og Hovmodet igen opstaa, faar man
ikke at vide; men fra Augustinus\footnote[1]{De civitate dei. XIV, 12---13. (Malæ voluntatis initium quæ potest esse nisi superbia?) De moribus eccl. cath. c. 12. --- Se Om Grundlaget for den humane Etik. p. 137—141. --- Ved Siden af den stærke Betoning af Hovmodet som den onde Villies Udspring gaar hos Augustinus den stærke Fremhæven af det sanselige Begær (concupiscentia), der især aabenbarer sig i Kønstrangen. Da denne rører sig uafhængigt af Villien, er det klart, at Naturen er fordærvet, og da Slægtens Forplantning nu netop sker gennem den med sanseligt Begær fuldbyrdede Avlingsakt, er det intet Under, at Menneskeheden er bleven til et Syndens Kaos (massa peccati). Smlgn. A. Harnack: Lehrbuch der Dogmengeschichte.\textsuperscript{3} 11. p. 29. Men det er dog Augustins udtrykkelige Lære, at det er den hovmodige Villie, der har fordærvet Naturen, saa at det Onde derefter opstaar uvilkaarligt og behersker Viliien. De civ. dei XIV, 3—6.} af staa de i Spidsen for
alle Laster. Medens Syndens Udvikling efter den filosofiske Opfattelse
foregaar fra det Ubevidste og Uvilkaarlige til det Bevidste
og Vilkaarlige, saaledes som Erfaringen viser os det ved
al sjælelig Udvikling, bliver det Opgaven for den teologiske Etik
at vise, hvorledes de mindre bevidste og vilkaarlige Former for
det Onde udvikle sig af de mest bevidste og vilkaarlige. I den
middelalderlige Etik er der gjort Forsøg paa at forklare de syv
Kapitalsynder genetisk, saaledes at der af Hovmodet successivt
udvikler sig Misundelse, Had, Sløvhed (acedia), Gærrighed, Umaadelighed
og Vellyst\footnote[2]{Hugo a St. Victore udleder acedia af de tre første Synder, der føre Mennesket til Væmmelse ved sig selv; af acedia udspringe igen de tre sidste Synder, idet Mennesket, der har mistet den indre Glæde, vender sig udad for at finde Trøst (Liebner: Hugo von St. Victor. Leipzig 1832. p. 278---2801. --- Hos Dan te antydes den samme Udviklingsgang: Skærsilden. XVII. Smlgn. Ozanam: Dante et la philosophie catholique au 13e siècle. Paris 1843. p. 94. --- Andreas Sunesøn: Hexaemeron.}. Selvom der i denne Skildring indflettes
% --- p. 130 ---
\setcounter{footnote}{0}%
\opage{130}interessante Iagttagelser, saa er den dog bygget paa en psykologisk
Umulighed. F. C. \textsc{Sibbern} har med Rette sagt, at den
teologiske Etik udleder det Onde af det højeste Onde. Det er
at vende op og ned paa Sagen. Først gennem en Række af
Udviklingstrin stiger den ubevidste Træghed, der kan være gavnlig,
til bevidst Afvisning af ethvert Forsøg paa at drage Individet
ud af dets Isolation. Først paa Højdepunktet af den ved
Trægheden indledede Modstand gør Individet alt Andet til Middel
for sin Nydelse eller sin Magtfølelse. I Stedet for at være
ét blandt flere har det nu forsøgt at gøre sig selv til absolut
Formaal. I det Enkelte er det saa langt fra, at der her kun
kan gennemløbes en eneste Udviklingsrække, at der for den
psykologiske Undersøgelse fremstiller sig en stor Mangfoldighed
af Veje.

De speciellere Former for det Onde ville bero paa de almindelige
Følelsesdispositioner eller Temperamenter, med hvilke
de enkelte Individer fra først af møde. Det mørke Temperament
fører hos passive Naturer til Forstemthed, Mismod og
Uimodtagelighed. Individet unddrager sig Alt, hvad der vil lægge
Beslag paa det, særligt ogsaa de oplivende og vækkende Paavirkninger.
Hos aktive Naturer opstaar under Forudsætning af
dette Temperament Skadefryd og Misundelse, Haardhed og
Grusomhed. Man føler Andres Lyst og Lykke som skærende
Disharmoni eller betragter den som Udtryk for Dumhed og
Kortsynethed. Der er Tilbøjelighed til at berøres af Alt paa en
smertelig Maade, og Ulystfølelsernes Hyppighed frembringer let
ed. Gertz. p. 153 har samme Rækkefølge (som synes at hidrøre fra St.
Gregor). --- Ifølge Thomas Aqvinas kunne alle syv Hovedsynder være
Udgangspunktet (radix); men da der til Grund for dem alle ligger Higen
efter egen Fremragen og Højhed, kan Hovmodet (superbia), den ubændige
Trang til selvisk Højhed (inordinatus amor propriæ excellentiæ) med
Rette kaldes Begyndelse til al Synd, ligesom den er Dronningen for alle
Synder. Thomas opstiller følgende Rækkefølge, som han dog ikke anser
for den eneste mulige: Hovmod (superbia), Gerrighed (avaritia), Vellyst
(luxuria), Misundelse (invidia), Fraadseri (gula), Vredagtighed (ira). Sløvhed
% CHECK p. 130: note whose mark was not found
\footnote[1]{(acedia). En Parallelisme mellem Kapitalsynder og Kardinaldyder anser Thomas ikke for mulig. Summa theologica. Pars 11. Qvæstio}
% --- p. 131 ---
\setcounter{footnote}{0}%
\opage{131}Egoisme ved den Opmærksomhed, Individet nødes til at henvende
paa sig selv. Smerten hindrer det i Hengivelse, hindrer
det i at glemme sig selv. --- Det lyse Temperament fører i
Passivitetens Form til Nydelsessyge, til at søge Tilfredsstillelse
for de i Øjeblikket opstigende Ønsker og Drifter eller til at
svælge i straalende Fantasier og derved tabe Evnen til at arbejde
paa de Opgaver, Livet stiller. Hos aktive Naturer fører
det til hensynsløs Virkedrift og Higen efter Selvstændighed, til
Herskesyge eller Fanatisme. Ogsaa her kan Haardhed og Grusomhed
udvikle sig, især naar den egne Evne til at lide Smerte
er ringe. Der kan her ofte vises stor Opofrelse og Anstrengelse
for at tilfredsstille Trangen til at fornemme sine Kræfter,
den egoistiske Magtsyge. Det bliver den heroiske Form for det
Onde. --- Under saadanne forskellige Former sætter Individet
sig fast i sig selv, gør sig selv til Centrum og umuliggør den
fulde Hengivelse. Endnu flere Former vil man finde, naar man
tager Hensyn til de mangfoldige Motivforskydninger, der blive
mulige udfra den oprindelige Disposition, og tillige til den Maade,
paa hvilken de almindelige psykologiske Elementer optræde under
de forskellige historiske og sociale Forhold. Det Onde er ikke
blot et psykologisk, men ogsaa et sociologisk Fænomen.

At det Onde saaledes kan fremtræde i mange forskellige
Former, er af stor praktisk Betydning. Naar den Ene ikke har
sin Ufuldkommenhed der, hvor den Anden har sin, kan gensidig
Kritik og etisk Paavirkning blive mulig. Den Enkelte kan
blive Forbillede i visse Henseender, selvom han ikke kan blive
det i alle Henseender. Hvis en saadan Forskelligartethed ---
Differentiation eller „Arbejdsdeling“ --- ikke fandt Sted, vilde
Mulighederne for Udvikling og Fremskridt være langt ringere,
end de ere.

2. Jeg gaar lidt nærmere ind paa det Spørgsmaal: hvor
høj en Grad af Bevidsthed kan der være til Stede i de nu skildrede
Tilstande? --- --- Prædikatet „Ondt“ udtales fra et Standpunkt,
som den, paa hvem det anvendes, ikke deler. Øjebliksmennesket
deler jo ikke Individualistens Standpunkt, og Individualisten
deler ikke dens Anskuelse, der stiller sig paa Familiens,
Stammens, Statens eller hele Slægtens Standpunkt. 9» Der%
% --- p. 132 ---
\setcounter{footnote}{0}%
\opage{132}som den, der staar paa det snevrere eller lavere Standpunkt,
havde fuldstændig og klar Bevidsthed om Prædikatets Gyldighed
og Anvendelighed paa ham, saa maatte ogsaa de tilsvarende
Følelser og Drifter sættes i Bevægelse, og hans Handlen maatte.
blive en anden. Der maa altid i det Onde være Uvidenhed og
Forblindelse: Uvidenhed, forsaavidt som Forestillingerne ikke
ere klare og fuldstændige; Forblindelse, forsaavidt som det er
visse herskende Følelser og Lidenskaber, der hindre Forestillingernes
rette Udvikling. Der kan være en Forestilling om et
andet Livstrin, men denne Forestilling høres kun som en fjern
Røst, vinder ikke Magt, vækker ingen stærk Følelse, gaar ikke
ind i det reale Selv. \textsc{Sokrates} havde i Hovedsagen Ret i, at
Synd er Uvidenhed, naar man blot ved Uvidenhed tænker paa
Mere end manglende Forstandsoplysning, eller naar man til
Uvidenheden føjer Forblindelsen for at udtrykke, at der ikke
blot mangler den rette Erkendelse, men at der ogsaa er positive
Hindringer for, at den kan udvikle sig og faa Magt over Sindet.
De, som handle ondt, „vide ikke, hvad de gøre“. Aarsagen
til, at de ikke vide det, ligger ganske vist dybere, end Sokrates
tænkte sig, eller i hvert Tilfælde dybere, end han forstod at
udtrykke. Men den Selverkendelse, Sokrates fordrede, vilde,
naar den gennemførtes, ogsaa bringe det skjulte Grundlag for
det hele Standpunkt for Dagen. Al Etik maa derfor tilsidst
være enig med Sokrates.

Til en egentlig Villiesakt hører nu jo dog, at man skal
være sig klart bevidst, hvad man gør! --- Her maa derfor skelnes
mellem Bevidstheden om de faktiske Forhold, som Handlingen
angaar, og Bevidstheden om Handlingens etiske Karakter.
Uvidenheden og Forblindelsen behøve ikke at angaa de
faktiske Forhold. Jeg kan vide bestemt, hvad jeg gør, som
f. Ex. Kætterdommeren véd, hvad han gør, naar han fælder Dødsdommen
over et Menneske paa Grund af hans Tro, uden dog
at være klar over, at det fra et højere etisk Standpunkt set er
noget Forkasteligt. En saadan Klarhed vilde gøre Handlingen
umulig. Ti det er psykologisk umuligt at handle mod sin bestemte
og fulde Overbevisning, naar denne ikke fordunkles og
fortrænges af andre Tilskyndelser. Man kan „have“ en Overbevis%
% --- p. 133 ---
\setcounter{footnote}{0}%
\opage{133}ning i den Forstand, at der er Mulighed for, at den kan fremkaldes
og blive levende i Bevidstheden, naar der bliver tilstrækkelig
Opfordring eller Anledning dertil. Men Spørgsmaalet er,
om den spontant (eller dog ved et Minimum af ydre Anledning)
trænger sig frem i Sindet i Overvejelsens og Handlingens Øjeblik;
ti kun da er det egentligt en bestemt og fuld Overbevisning.
Her er uendeligt mange Grader og Nuancer af Klarhed,
og det er dogmatisk uden videre at fastslaa den højeste Grad
som given.

Man kunde mene, at der alligevel findes en saadan Villen
mod bedre Viden, naar Mennesket søger at fordunkle den etiske
Bevidsthed hos sig for at slippe fra den Splid og Delthed, som
opstaar, hvor den ikke strax kan blive herskende, eller naar Mennesket
søger at hindre Opstaaen af Tanker, han frygter vilde
ændre hele hans nuværende Vurdering, hans nuværende System\footnote[1]{Smlgn. Psykologi VI F, 2.}.
Men i saadanne Tilfælde har jo den højere etiske Bevidsthed
ikke faaet fast Fod; ellers kunde Pinligheden ved den
delte Tilstand ikke blive det raadende Motiv. Der øves ikke
Ondt med fuld og klar Bevidsthed om, at det er Ondt.

Det er først i den nyeste Tid, at store Digtere have skildret
det Onde med en finere og dybere Psykologi end den, der
ligger til Grund for den sædvanlige Djævletype, hvilken endog
Shakespeare altfor meget har fulgt i sin Richard III og sin Jago.
I første Linie staar her Begyndelsen af Dostojewskys Raskolnikow.
Men jeg maa ogsaa henvise til en enkelt Replik i \textsc{Ibsens}
„Rosmersholm“. Skildringen af Rebekkas Karakter kan jeg ikke
paatage mig at forsvare; men en enkelt Replik er mesterlig.
„Der er da vel to Slags Viilie i et Menneske, skulde jeg mene!
Jeg vilde ha' Beate væk. Paa den ene Maaden eller paa den
anden. Men jeg troede aldrig, at det skulde komme alligevel.
Ved hvert Skridt, som jeg fristed og voved fremad, syntes jeg
ligesom noget skreg indeni mig: Nu ikke længer! Ikke et Skridt
længer! --- Og saa kunde jeg ikke la' det være endda. Jeg
maatte friste et bitte lidet Gran til. Bare et eneste et. Og saa
et til --- og altid et til. --- Og saa kom det. --- Det er paa
% --- p. 134 ---
\setcounter{footnote}{0}%
\opage{134}den Vis, sligt noget gaar for sig“. En mindre Digter vilde have
skildret hende som en Djævel, der med fuld Bevidsthed lagde
sin Plan og klogt beregnede de mange smaa Knappenaalsstik,
der skulde føre til Maalet.

Den teologiske Forestilling om „Synd mod den hellige Aand“
er ligeledes en psykologisk Umulighed. \textsc{Luther} betegner den i
en Prædiken over dette Emne som „Djævlesynd, der, skønt
aabenbar og overbevist, ikke vil overbevises“. Han siger: „Denne
Synd har jeg i forrige Tider, da jeg endnu var en lærd Doktor,
aldrig villet tro, var til i Verden“. Senere kom han til at
frygte for, at Papisterne, der ikke kunde modbevise den rette
Lære, vilde blive stikkende i denne Synd! Som lærd Doktor
var Luther en bedre Psykolog, end da han var bleven den store
Agitator. Naar man ikke vil overbevises, hvorledes kan man
saa være overbevist? Saalænge man ikke vil overbevises, maa
der jo dog være en Del af Sindet tilbage, hvor Overbevisningen
ikke er trængt hen med sin fulde Virkning; hvorledes kunde
ellers en saadan Villie opstaa? --- Luther siger: „Naar Nogen
kommer saa vidt, at han Intet vil høre eller se, ja tilmed vil
forsvare sin Gudsbespottelse, saa staar han ikke længer til at
raade og hjælpe“. Men den, der forsvarer sig, beraaber sig dog
paa Grunde; han anerkender altsaa endnu Sandhedens Magt,
kun at han ikke anerkender som Sandhed, hvad hans Modstandere
kalde saa. Her er altsaa endnu noget Godt tilbage. Det aldeles
Djævelske vilde være, slet ingen Sandhed at anerkende; men
det vilde tilsidst føre til fuldstændig Meningsløshed og Tankeløshed
og til fuldstændig Afmagt. Den absolut konsekvente Ondskab
vilde altsaa være absolut Dumhed, og det er et rigtigt
Instinkt, der lod Folketroen forvandle den onde Djævel til den
dumme Djævel\footnote[1]{Allerede i den gammelpersiske Teologi fremtræder dette Træk: Ahriman aner i sin Vankundighed ikke, at Ormuzd, den gode Magt, er til, og da han har opdaget denne Magt, undervurderer han den, ligeledes af Vankundighed. Ed v. Lehmann: Zarathustra. En Bog om Persernes gamle Tro. København 1899—1902. II. p. 161.}.

3. Et Skridt videre fremme end Forblindelsen ligger For%
% --- p. 135 ---
\setcounter{footnote}{0}%
\opage{135}hærdelsen. Men heller ikke den er Et med bevidst Ondskab.
Forhærdelse er Forblindelse, der har sat sig fast, er gaaet over
i Kød og Blod, saaledes at en Ændring af Karakteren er overordentligt
vanskelig. En saadan Forhærdelse vilde det dog være
uberettiget at gøre til Et med Uforbederlighed. Vor Opdragelseskunst,
baade den individuelle og den sociale, er neppe saa vidt
fremskreden, at vi tør erklære et Menneske for absolut uforbederligt.
Vor Kunst og vore Midler slaa maaske ikke til. Men
der er hidtil paa dette Omraade hverken sat den Kraft eller
den Kløgt ind, som Menneskeslægten har anvendt paa saamange
andre Omraader. Man har ladet sig nøje med at hensætte slige
Individer i et „Tugthus“ og give dem Anvisning paa et Helvede
i den anden Verden. Den Kunst og den store, taalmodige
Menneskelighedsfølelse, hvilke ene kunne overvinde de Hindringer,
der her kunne foreligge, ere endnu kun i deres Vorden\footnotemark[1].

\chapterhead{VII.}{VELFÆRDENS TEORI.}

1. I de tre nærmest foregaaende Kapitler have vi beskæftiget
os med Etikens Grundlag. Vi gaa nu over til en nøjere
Fastsættelse og Belysning af Principet for Etikens Indhold. Til
det Grundlag, hvorpaa vi her have stillet os, svarer som tidligere
vist (III, 10—15), kun ét saadant, nemlig Velfærdsprincipet.
Det giver det almindelige Svar paa det Spørgsmaal: hvad er
det Gode?

Man har af misforstaaet idealistisk Interesse ofte overset,
at det Godes Begreb er et rent formelt Begreb, og har drevet
en vis Kultus med det, uden at give nogen nærmere Bestemmelse
deraf. Det sande Gode, mente man, maa være det i og

) Smlgn. udførligere herom mine Etiske Undersøgelser p. 73—81
% CHECK p. 135: note mark without a note

% --- p. 136 ---
\setcounter{footnote}{0}%
\opage{136}for sig Gode, hvis Værdi ikke afhænger af noget Andet. Men
hvad er det, der har sin Værdi i sig selv? Hvilket Indhold
tilfredsstiller denne Betingelse?

Som tidligere udviklet maa dette Spørgsmaal besvares forskelligt
paa de forskellige etiske Standpunkter. Naar en uinteresseret
og universel Sympati er bestemmende for den etiske
Vurdering, kan kun det være godt, der bevarer og forhøjer bevidste
Væseners Velfærd, forøger deres Lyst eller forminsker
deres Smerte. Enhver Handling, som virker i denne Retning
uden at medføre fjernere Virkninger i modsat Retning, er dermed
retfærdiggjort, enhver Handling, om hvilken det Modsatte
gælder, er forkastelig.

Etiken arbejder her kun videre paa, hvad der fra Naturens
Haand er indledet. Ti i det Store og Hele er Lystfølelse
knyttet til sund og naturlig Brug af Kræfterne og til det, der
bevarer og fremmer Livet, Ulystfølelse til det Modsatte. Smerte
er Tegn paa begyndende Opløsning af Livet\footnote[1]{Psykologi VI D.}. Naar Etiken
vurderer vor Villen og Handlen efter den Grad, i hvilken Lyst
derved kommer til at sejre over Smerte, arbejder den --- formedelst
den nøje Sammenhæng mellem Lyst og Liv, Smerte og
Død --- paa Udviklingen af det menneskelige Liv, de menneskelige
Evner og Kræfter i saa rigt Maal og paa saa harmonisk
Vis som muligt. De etiske Problemer angaa derfor ogsaa Vurdering
af de enkelte Øjeblikkes Lyst og Smerte i Forhold til
Individets Livshelhed og af de enkelte Individers Lyst og Smerte
i Forhold til Slægtens Livshelhed. Kun saavidt denne Vurdering
er mulig, kan Vurderingen af Villiesakter og Handlinger begrundes,
og har den etiske Handlen et bestemt Indhold. Enhver
Villiesakt og enhver Handling er som en Sten, der kastes
i Vandet; Bevægelsen forplanter sig i større eller mindre Kredse.
Og Vurderingen beror paa, hvorledes Handlingen griber ind i
bevidste Væseners Liv og Følelse. Ligesom den teoretiske Videnskab
forklarer et Naturfænomen ved et andet, saaledes vurderer
Etiken den ene Følelse ved andre Følelser: den Tilfredsstillelse,
som den Handlende føler ved eller opnaar ved sin
% --- p. 137 ---
\setcounter{footnote}{0}%
\opage{137}Villen og Handlen, er kun at kalde ubetinget god, naar Handlingens
Virkninger ikke gribe forstyrrende ind i andre Væseners
Lystfølelse. Ethvert saadant Indgreb kræver Retfærdiggørelse
ved at paavises som nødvendigt Middel til en saameget
des større Velfærd for den Lidende selv eller for Andre. Saaledes
maa ifølge Velfærdsprincipet Bevisbyrden fordeles. Det
er Tilføjelsen af Smerte, der skal begrundes, medens Lystfølelsen
har sin foreløbige Retfærdiggørelse i sig selv. Denne Grundsætning
er først opstillet af \textsc{Platon} (Protagoras. Heises Overs.
p. 198 f.), som lærte, at naar sanselig Nydelse er at laste, er
det paa Grund af dens Følger, ikke paa Grund af den Lyst,
den skaffer i Øjeblikket\footnote[1]{Af nyere Tænkere har Spinoza hævdet den (Ethica IV, 43. 60. App. c. 30); ligeledes Malebranche (Recherche de la vérité IV, 10). Derimod nærede Arnauld, i sin Polemik med Malebranche, stor Betænkelighed ved den. Smlgn. Bay les Dictionnaire. Art. Épicure.}.

Velfærd ytrer sig som en varig Tilstand af Lystfølelse. Den
kan maaske kun vindes ved en Udviklingsgang, der fører gennem
Smerte og Savn; men den er det Maal, for hvis Skyld
Smerte og Savn udholdes. Man behøver ikke at være Optimist
for at anerkende Velfærd som Etikens objektive Princip. Ti
dermed siges aldeles Intet om, hvor let det er at naa Maalet,
eller i hvilken Grad det idethele taget er muligt at naa det.
Dersom Pessimismen skulde have Ret, dersom al Lystfølelse, al
Lykke var en Illusion, og smertelig Brynde og Trang det Inderste
i os Alle, --- saa vilde dog dermed Velfærdsprincipet
ikke omstyrtes. Det vilde da blive Opgaven at bringe saa
megen Lindring som muligt, og saadan Lindring vilde da være
al den Velfærd, der kunde opnaas.

Kun fra et absolut asketisk Standpunkt, det vil sige et
saadant, som vilde hævde, at Smerte i sig selv og i al Evighed
var at foretrække for Lyst, kan Velfærdsprincipet her bestrides.
Men et saadant Standpunkt kan heller ikke gennemføres;
selv Asketen plager kun sig selv og Andre for at
frembringe større Lystfølelse derved. Enhver etisk Tankegang
fører tilsidst ud til et Punkt, hvor den sidste Grund til en
% --- p. 138 ---
\setcounter{footnote}{0}%
\opage{138}Handlings Billigelse viser sig at være, at der paa et eller andet
Punkt i Verden frembringes Lyst eller tilintetgøres Smerte.

Som allerede tidligere bemærket (III, 11), skyldes Modstanden
mod dette Princip for en Del Misforstaaelser, som Ordene
have voldet, og jeg bruger hellere Ordet „Velfærd“ i Stedet for
„Lykke“ eller „Nytte“ for at undgaa saadanne Misforstaaelser.
--- Ved „Lykke“ tænker man ofte paa noget Tilfældigt og Udvortes.
„Lykke kaldes det“, siger P. E. \textsc{Muller} i sin „Synonymik“,
„naar et Gode vederfares os, som vi ikke selv kunne
bevirke“. Det er i denne Betydning, at Digteren taler om „Lykkens
lunefulde Spil“. Ofte bruges Ordet Lykke ogsaa om mæt
Tilfredshed, Hvile i, hvad der er opnaaet. Hvilken af disse Betydninger
man tager, saa er det klart, at der kan være Noget,
som er bedre end „Lykke“. Men hvorfor skulde Ordet Lykke
ikke kunne bruges om den Lystfølelse, der er grundet paa og
knyttet til Selvvirksomhed og kun kan fastholdes ved stadigt
fortsat Udvikling? --- Her skal da foreløbigt ingen Art eller
Grad af Lystfølelse udelukkes; Rangforordningen mellem de forskellige
Lystfølelser kan først fastsættes ved de følgende speciellere
Undersøgelser. Men i og for sig maa enhver Etik, der
anerkender Sympatien som Vurderingsmotiv, ogsaa finde det
værdifuldt, at man bibringer Andre Goder, som de selv ikke
kunne bevirke; kunde de selv bevirke disse Goder, behøvede
de jo ingen Hjælp. Og Mættelse er jo i og for sig ikke noget
Ringe, hvad enten den følger paa legemlig eller aandelig Hunger
og Tørst. --- Ved „Nytte“ tænker man let kun paa, hvad
der er blot Middel for noget Andet, altsaa ikke har Værdi i sig
selv. Det er et Forretningsudtryk og skurrer i Manges Øren,
blot fordi det ikke er højtideligt nok\footnote[1]{vilde ikke kalde Agtelse for en Følelse; den skulde hverken være Lyst}. Men idet man med

!) Til højtideligt Brug have vi det gamle danske Ord „Gavn“. Ordet
Nytte kom fra Tyskland. „Idet Gavn sjeldnere hørtes i daglig Tale, fik
det en desto almindeligere og hæderligere Betydning, og Præsten brugte
det paa Prædikestolen“ (P. E. Müller). Naar Jesus siger: „Hvad gavner
det et Menneske, om han vinder den hele Verden, men tager Skade paa
sin Sjæl“, saa vilde det, i Følge P. E. Müller, ikke have været rigtigt at
bruge Ordet „nytter“ i Stedet for „gavner“! --- (Tilføjelse 1913). Kant
% --- p. 139 ---
\setcounter{footnote}{0}%
\opage{139}overspændt Indignation vender sig bort fra denne jævne Prosa,
kommer man let til at overse, at det „at have Værdi i sig selv“
ikke kan betyde Andet end at være umiddelbar Genstand for
en Følelsestrang. Værdi „i sig selv“ er det samme som umiddelbar
Værdi. Som sidste Maalestok genfinde vi stedse et eller
flere bevidste Væseners Lyst og Smerte. --- Med Velfærd mener
jeg den „sande“ Lykke eller den „sande“ Nytte, overladende
til de følgende Afsnit nærmere at paavise dens Betingelser.

2. Men Velfærdens almindelige Begreb indeholder nogle
Vanskeligheder, som vi allerede her maa fremdrage\footnote[1]{I mine Etiske Undersøgelser, p. 24—42 har jeg drøftet de forskellige Indvendinger, der kunde rejses eller ere blevne rejste mod Velfærdsprincipet, til Dels noget udførligere og fra en noget anden Side end i det Følgende.}.

Kan en varig Tilstand af Lystfølelse ikke naas paa meget
forskellige Udviklingstrin og ad meget forskellige Veje? „Sand“
eller rettere virkelig Lykke kan jo ikke betyde Andet end fuld
Tilfredsstillelse af al den Trang, der kan føles Og den Trang,
man føler, retter sig efter den Natur, man nu engang har. Maa
da den Enes Velfærd eller Lykke ikke være ligesaa god som
den Andens, da de jo hver for sig ikke kunne naa videre end
at faa Tilfredsstillelse for deres Trang, og den Ene naturligvis
ikke behøver Tilfredsstillelse for en Trang, som ikke han selv,
men en Anden har? Hvorledes kan da Velfærdsprincipet afgive
Maalestok for Vurdering af Lystfølelserne paa de forskellige
Trin og hos de forskellige Væsener?

Og for det Andet. Fører Udviklingen af de menneskelige
Evner og Kræfter, fører den hele Kulturudvikling virkeligt til
eller Ulyst. Fichte skelnede skarpt mellem „Genuss“ og „Würde“.
Ifølge He gel ere de lykkelige Tider Historiens tomme Blade, fordi der
mangler Modsætninger. Carlyle skelner skarpt mellem „happiness“ og
„blessedness“. --- Paa den anden Side hævdede Rousseau, at der ikke
gives nogen Lykke uden Mod. Grønbech lærer (i sin Bog Lykkemand
og Nidding), at det germanske Ord Lykke indeholder en potentiel moralsk
Vurdering; Lykken er Villiessiden i Mennesket, det sidste og dybeste
Udtryk for dets Væsen; den samler Ættens Livsbetingelser i sig. Og
J her ing udtaler: „Wohlsein ist aufgespeicherte Arbeitskraft“. --- Man ser.
hvorledes Ordceremoniellet kan faa Indflydelse paa Teoriernes Skæbne.
% --- p. 140 ---
\setcounter{footnote}{0}%
\opage{140}Velfærd? Ved Kulturen forøges ganske vist de Midler, Slægten
raader over, og der uddannes Evner og Færdigheder, som ikke
før vare til Stede. Men forøges samtidigt ikke ogsaa Fornødenhederne,
baade de materielle og de ideelle, og dermed Mulighederne
for Savn, Smerte og Disharmoni? Selv Udvikling og
Udvidelse af Sympatien og Pligtfølelsen truer jo Velfærden, da
man derved bliver saarlig paa langt flere Punkter, stiller større
Fordringer til sig selv og Andre og ikke saa let finder Fred og
Ro i sit Sind!

Disse to Spørgsmaal vil jeg undersøge hvert for sig.

3. \textsc{Stuart} Mill har sagt: „Det er bedre at være et utilfredsstillet
Menneske end et veltilfreds Svin, bedre at være en
utilfredsstillet Sokrates end en veltilfreds Taabe.“ Han begrunder
denne Paastand ved, at selvom Svinet og Taaben vilde
mene det Modsatte, maatte deres Dom afvises, fordi de ikke
kende det højere Standpunkt, hvorfra Mennesket og Sokrates
betragte Livet, hvorimod Mennesket kender Svinets Fornødenheder,
og Sokrates gennemskuer Taaben. Man maa holde sig
til deres Dom, der kende begge Slags Fornødenheder og derfor
kunne anstille en Vurdering af dem\footnote[1]{Moral grundet paa Nytte- og Lykkeprincipet. Dansk Overs. p. 10—12. --- En lignende Tankegang findes allerede hos Platon (Statens 9de Bog. Heises Overs. p. 86—92). Ogsaa Voltaire har drøftet Problemet. Se Histoire d'un bon brahmin.}.

Jeg maa dog tage Svinet og Taaben noget i Forsvar. Vanskeligheden
er større, end Mill antager. Mennesket kender ganske
vist Svinets Fornødenheder, og Sokrates vil ikke have Vanskelighed
ved at sætte sig ind i Taabens. Men Mennesket har ikke
Svinets, Sokrates ikke Taabens Fornødenheder som eneraadende,
og det er dog dette, hvorpaa det i denne Sammenhæng kommer
an. Mennesket kan ikke gøre sig til et Svin uden at ophøre
at være Menneske\footnotemark[1], og Sokrates vil neppe kunne gøre sig
% CHECK p. 140: note whose mark was not found
\footnote[2]{Hvis det lykkedes, vilde Nydelsen maaske blive saameget des større, saafremt nemlig Leipzigerstudenterne i Goethes Faust have Ret, naar de synge:}
% CHECK p. 140: note whose mark was not found
\footnote[3]{Uns ist so kannibaliseh wohl,}
% CHECK p. 140: note whose mark was not found
\footnote[4]{als wie fiinfhunderl Sauen.}
% CHECK p. 140: note mark without a note

% --- p. 141 ---
\setcounter{footnote}{0}%
\opage{141}saaledes til Et med Taaben, at de sokratiske Fornødenheder
ganske falde bort. Naar nu Svinet kan faa fuld Tilfredsstillelse
af al sin Trang, er dets Lykke da ikke større end Menneskets,
som vanskeligt faar alle sine Ønsker og al sin Higen opfyldt?
Og Taaben, der ikke har mange Tanker og ikke stiller store
Fordringer til Livet, er han ikke lykkeligere end Sokrates, der
tilbringer hele sit lange Liv med at lære sig selv at kende og
at vække Andre, for saa tilsidst at erklære, at Døden i Grunden
er bedre end Livet?

Det gælder jo her ikke om en bestemt Sum af ydre Goder,
men om Goder, der virkeligt ere Goder for dette bestemte Væsen.
En og samme Pengesum kan jo være en stor Vinding for
den Ene og en ren Ubetydelighed for den Anden. Det er en
psykologisk Hovedlov, at en Følelses Grad bestemmes ved den
Forandring, der frembringes i Individets hele Tilstand.

Endeligt har den, der har vundet fuld Tilfredsstillelse for
sin Trang, intet Motiv til at sammenligne sin Tilstand med Andres.
At Svinet ikke lader sig genere af Mennesket, Taaben ikke af
Sokrates, kommer naturligvis af Mangel paa Evne. Men selv,
hvor de intellektuelle Betingelser for at anstille en saadan Sammenligning
ere til Stede, vil man ikke kunne føle nogen Brod
ved Sammenligningen, ingen Trang til at komme ud over sin
Tilstand. Det vil gaa som i \textsc{Dantes} Paradis, hvor der er forskellige
Trin af Salighed, men hvor de, der ere stillede paa de
lavere Trin, ingen Uro føle ved Tanken om de højere, da deres
„Villie er bragt til Hvile, saa kun i, hvad den har, den finder
Gammen“. Digteren fortsætter:

\begin{quote}
Da blev det klart mig, at der alle Vegne\\
er Paradis i Himlen, hvor jeg vanker,\\
skønt Gud forskelligt lader Naaden regne\footnote[1]{Dante: Paradiset. IIl v. 88—90. --- Smlgn. Augustin us: De civitate dei XXII, 30.}.
\end{quote}

I Dantes Paradis er derfor al Stræben og Udvikling forbi.
Men vil det ikke gaa saaledes overalt? Og hvorledes kan Velfærdsprincipet
lægges til Grund for en Sammenligning af for%
% --- p. 142 ---
\setcounter{footnote}{0}%
\opage{142}skellige Væseners Tilstande, naar disse hver for sig have Karakteren
af fuldstændig Tilfredshed?

Efter nu saa skarpt som muligt at have stillet det Problem,
som her frembyder sig, skal jeg søge at besvare det.

Ligesom der i den ydre Natur ikke findes nogen absolut
Hvile, men overalt Bevægelse under forskellige Skikkelser, saaledes
er ogsaa Velfærd som absolut Tilfredsstillelse af enhver
Trang for ethvert Væsen kun tilnærmelsesvis mulig, og enhver
Tilnærmelse til den vil være Anfægtelse underlagt. Det ligger
deri, at ethvert Væsen er underlagt ydre og indre Forandringer.
Naturen og Livet staa ikke stille, men vedblive at virke
baade i Individet og udenom det, og disse Forandringer ville
vække ny Trang og stille nye Opgaver.

Hvad de ydre Forandringer angaar, behøves ingen nærmere
Paavisning. Vi søge stedse forgæves at trække os tilbage i en
idyllisk Tilværelse. Tidligere eller senere afbrydes Idyllen af
Forandringer i Omverdenen, der gøre nye Anstrengelser og Arbejder
nødvendige.

Og med de ydre Forandringer følge de indre. Forestillinger
om det Nye og Fremmede ville i Længden ikke kunne
holdes borte, og de ville sætte Følelse og Villie i Bevægelse.
Dersom vi tør anvende vore psykologiske Love paa Forholdene
i Dantes Paradis, saa vil der hos de Salige med Forestillingerne
om de højere Livstrin ogsaa opstaa en Higen og Længsel efter
at opnaa dem. --- Og selvom nye Forestillinger ikke opstaa og
virke, saa vil dog selve Følelsen behøve Modsætning og Afvexling
for at bevare sig frisk og levende. Ensformighed og
Gentagelse sløve og svække. Derfor behøver Idyllen som Kunstart
til Baggrund en Antydning af et større, mangfoldigt og bevæget
Liv for at kunne gøre sin Virkning. Saaledes danner i
\textsc{Goethes} „Hermann und Dorothea“ Revolutionskrigen og dens
Virkninger den urolige Baggrund\footnote[1]{havde skildret Overgangen fra den tunge Kamp med Virkeligheden til den ideale olympiske Harmoni, fattede han den Plan at forfatte en Idyl, en Skildring af en over al Strid hævet salig Tilværelse. Han ansaa dette for den højeste poetiske Opgave. Men han saa Nødvendigheden af, at der}. --- Endeligt vil der, naar

') Efterat Schiller i det storartede Digt „Das Ideal und das Leben“
% --- p. 143 ---
\setcounter{footnote}{0}%
\opage{143}Ensformigheden og Gentagelsen ikke fører til Tilbagegang og
Svækkelse, naturligt opstaa en Trang til Bevægelse og Virksomhed.
Vi bevæge os jo ikke blot for at opnaa bestemte Formaal,
men ogsaa for at bruge Kræfterne og faa Afløb for opsamlet
Energi. Denne Bevægelsestrang vil voxe, jo mere der samles
Kræfter, jo længer Hvilen varer. Den Virksomhed, den fører
til, vil føre til nye Erfaringer og dermed til nye Kampe.

Selvom disse Aarsager ikke virke hos den Enkelte, der har
naat en vis Ligevægtstilstand, ville de dog virke i Slægten efter
ikke mange Generationers Forløb. Det vil vise sig, at hvis der
ikke sker en Tilbagegang eller en Opløsning, vil enhver Opnaaelse
af et Formaal kun være Begyndelsen til, at et nyt Formaal
sættes. Velfærd er derfor en Illusion, dersom dermed menes
en passiv, en Gang for alle frembragt Tilstand. Den maa bestaa
i Virksomhed, i Arbejde, i Udvikling. Det ligger allerede
i Ordet: Velfærd er at fare vel --- altsaa at være i Bevægelse,
ikke i Hvile. Naar Ordet „fare“ ogsaa kan betyde „at være i
en vis Tilstand“, er det, ifølge Molbech, en „figurlig og uegentlig“
Betydning. Velfærd som Hvile kan kun betyde foreløbig
Afslutning, Opnaaelse af et nyt Niveau, hvorpaa en ny Udvikling
kan indledes. Vor Definition af Velfærd som en varig
Tilstand af Lystfølelse behøve vi dog derfor ikke at tage tilbage.
Hvad der maa afvises, er blot Forestillingen om en passiv Tilstand.
Virksomhed er jo dog ogsaa en Tilstand, ligesaavel som
Hvile er det. Ved en naturlig Illusion tænke vi os gerne vore
fjernere Formaal som hvilende Tilstande. Vi kunne paa Grund
af Fjernheden ikke se den Mangfoldighed, de indeholde, og de
Opgaver, de stille. Af denne Illusion river den virkelige Erfaring
os stedse ud. Vi lære at forstaa, at kun den Lykke,
der findes i Arbejde paa Livets Fremgang hos os selv og Andre,
er bygget paa fast Grund.

Der er her en betydningsfuld Parallel mellem Velfærdens
dog blev mindet om Menneskelivets Strid og Kamp: „Hovedpersonerne
skulle vel være Guder, men ved Herkules kan jeg endnu finde et Tilknytningspunkt
mellem dem og Menneskeheden og bringe en Bevægelse
ind i Billedet“. (Briefwechsel zwischen Schiller und Wilhelm v. Humboldt.
% CHECK p. 143: note whose mark was not found
\footnote[1]{Stuttgart und Tübingen 1830. p. 328). --- Planen blev ikke udført.}
% --- p. 144 ---
\setcounter{footnote}{0}%
\opage{144}Teori og Samvittighedens Teori. (Smlgn. IV, 3 —4). Ingen af
dem kan antage nogen absolut Afslutning, men begge føre os
til Tanken om en uoverskuelig Udviklingsgang. Og der er tillige
en Vexelvirkning mellem dem. Ti jo rigere en Forestilling
man har om, hvad der hører til Velfærd, des mere udvikles og
øves Samvittigheden, da den faar des flere Opgaver at behandle.
Og jo finere og skarpere udviklet Samvittigheden er, des større
Fordringer stiller den til det, der (hos Andre og hos os selv)
skal fortjene Navn af Velfærd. Den, hvis Øje er sløvet, og hvis
Synskreds er snever, er let at tilfredsstille.

Naar store Aander ofte dømme nedsættende om den Lykke,
der kan opnaas i Verden, da tænke de mere paa de passive Tilstande,
der udfylde Mellemrummene mellem den store Virksomheds
Tider, end paa selve den Virken, der har gjort dem til
de store Aander. Det er psykologisk forstaaeligt, at de ide
virksomme Tider, da Livet rørte sig stærkt i dem, ikke tænkte
paa deres egen Tilstand, men gik op i Arbejdet for deres Sag.
Det er altid misligt, naar Opmærksomheden fra Tanken om Formaalet
glider over til den Lystfølelse, dets Opnaaelse kan skænke\footnote[1]{Psykologi VII, B, 1 a; 3.}.
Derved kommer man til at sondre, hvad der ikke skal sondres,
og derved bliver Tilstanden forvandlet fra en aktiv til en passiv.
I samme Øjeblik, som man spørger sig selv, om man er
lykkelig, isolerer man let sig selv fra det Indhold, man lever
i og for. Hint Spørgsmaal kommer slet ikke frem, saalænge
man gaar op i dette Indhold; at det kommer frem, tyder paa
en begyndende Tvivl. Det gælder jo om al etisk Reflexion, om
al Erkendelse af Godt og Ondt, at den forudsætter (ikke som
i den gamle Fortælling: medfører) Ophør af den paradisiske
Tilstand. Reflexionen kan naturligvis vækkes udefra —af den,
der bringer Slangen ind i Paradiset. Men saa forudsættes, at
den er vakt hos Andre, og den vil ikke virke, hvor der ikke
var en begyndende Disharmoni forud. Lykken er Expansion,
ikke Reflexion, eller --- for at bruge Ludvig \textsc{Feilbergs} Udtryk
--- „Ligeløb“, ikke „Kredsning“. I en vis Forstand kan
man derfor sige, at vi kun ret kende Lykken, naar den er for%
% --- p. 145 ---
\setcounter{footnote}{0}%
\opage{145}svunden. Stræben efter Lykke bliver da en Stræben efter paa
et højere Trin at genfrembringe det, vi have kendt paa et tidligere
Trin. Og jo mere selve Virksomheden er forbunden med
Lystfølelse, jo mindre Middel og Maal --- Stræben efter Lykken
og selve Lykken --- falde udenfor hinanden, des fastere er Velfærden
grundet. Det er ingenlunde ligegyldigt, hvorledes Middel
og Maal, Aarsag og Virkning her forholde sig til hinanden,
--- om de staa i rent udvortes Forbindelse, eller om der er en
indre Sammenhæng imellem dem, saa at Formaalet allerede
rører sig i Midlet, og Virkningen kun er en Udfoldelse af, hvad
der ligger i Aarsagen. To Tilstande kunne tilsyneladende være
ganske ens, og dog kan den ene paa Grund af sin Tilblivelsesmaade
faa større Værdi end den anden ved de skjulte Elementer
(den potentielle Energi), den indeholder, nemlig dels Mulighed for
rig Erindring, dels Mulighed for fremtidig Udvikling. Værdien 2
kan her være større, naar den er naat ved 7 ÷ 5, end naar
den er naat ved 1 + 1. Den Velfærd, der er naat gennem
Smerte og Kamp, kan være større end den tilsyneladende
samme Velfærd, der naas ved rolig og sagte Udfoldelse, --- kan
indeholde flere Muligheder end denne. --- Et fuldstændigt og
en Gang for alle opnaaet Formaal formaar ikke længere at sætte
Virksomhedsdriften i Bevægelse, og den Følelse, som knytter
sig til Forestillingen om det, kan derfor heller ikke blive saa
stærk og levende, som den, der betinges ved Forestillingen om
et Formaal, der stadigt kan virkeliggøres i højere Grad og i
større Udstrækning. Det Højeste, vi kunne tænke os, er en
Fremgang, hvor hvert enkelt Skridt føles som et Gode, fordi
det sætter vore Kræfter i Bevægelse uden dog at anspænde
dem over Evne.

Herved ere vi førte over til det andet af de opkastede
Spørgsmaal, om Kulturudviklingens etiske Betydning.

4. Velfærd bestaar i Virksomhed, men derfor er der ikke
til enhver Virksomhed knyttet Velfærd. Ved denne simple Sætning
er Forholdet mellem Etik og Kultur foreløbigt betegnet.
Kultur er det videre Begreb, Etik det snevrere. Al etisk Stræben
er Kulturvirksomhed, men ikke al Kulturvirksomhed er
etisk.

% --- p. 146 ---
\setcounter{footnote}{0}%
\opage{146}Ved Kultur forstaa vi Bearbejdelse af det af Naturens
Haand givne Stof. Dog kan der ingen faste Grænser drages
mellem Natur og Kultur. Hvad der fra ét Synspunkt set er
Kultur, er fra et andet Synspunkt set Natur. Ti hvad den ene
Slægt møjsommeligt frembringer, er for de følgende et givet
Grundlag. End ikke i den dyriske Verden finde vi altid det
rene og bare Naturliv. Fuglen har sin Rede, som den selv har
bygget, og Ræven sin Hule, som den selv har uddybet. Og
desuden har enhver Generation i sin Natur en Arv fra tidligere
Generationers Erfaring og Virken.

Begrebet Kultur er i og for sig et rent formelt Begreb.
Bearbejdelse og Virksomhed siger os Intet om, hvorvidt der
arbejdes og virkes til Gavn og Fremgang. Den saakaldte Kulturbevidsthed
kan ofte være en hjerteløs Glæde over at se de
mange Hjul i Virksomhed, de mange hverandre krydsende
Traade, som knyttes sammen til et sindrigt Væv. Man tænker
da ikke paa, hvad der opnaas ved det Hele, eller paa dem,
der knuses under de rullende Hjul. Hvilken Værdi har egentligt
det hele rastløse, fra Slægt til Slægt fortsatte Arbejde?
Vore Færdigheder og vore Fornødenheder stige; men var det
ikke bedre at slippe for disse nye Færdigheder, og disse nye
Fornødenheder, dersom Velfærd kunde opnaas paa simplere og
lettere Vilkaar? --- Vi føres ved dette Spørgsmaal igen tilbage
til den første Vanskelighed, der laa i Velfærdens Begreb. De
to Vanskeligheder, dette Begreb fører os ud i, hænge saa nøje
sammen, at Overvindelsen af den anden danner Fortsættelsen
af Overvindelsen af den første. Fra Velfærd som passiv Tilstand
maatte vi gaa over til Velfærd som Virksomhed; men
her viser det sig igen, at en nærmere Begrænsning er nødvendig.

Denne Begrænsning er given ved den almindelige Teori
om Lystfølelsens Sammenhæng med den sunde, naturlige og livsudviklende
Brug af Kræfterne. Virksomheden hører op at være
Velfærd, naar den enten overspænder Kræfterne, eller naar
Kræfterne spredes og splittes, eller naar Kraftanvendelsen ensidigt
sker i en enkelt Retning paa andre vigtige Retningers
Bekostning. Hvor Modstanden er større, end nødvendigt er for
% --- p. 147 ---
\setcounter{footnote}{0}%
\opage{147}Kræfternes Udfoldelse, vil der opstaa Smerte og Lidelse; og en
saadan smertelig Virksomhed kan kun retfærdiggøres som Overgang
til en Tilstand, hvor der er tilvejebragt et mere harmonisk
Forhold mellem de ydre og de indre Betingelser. Hvor
Kræfterne ikke tilsidst samles mod et fælles Maal, opløses Livet
i Brudstykker, og ingen samlet Livsfølelse bliver mulig. Og
endeligt, hvor nogle Kræfter øves paa andres Bekostning, vil
der tidligt eller sent vise sig et Misforhold mellem det Indre
og det Ydre; de ensidigt udviklede indre Kræfter ville ikke
svare til de flersidige Fordringer, de ydre Livsforhold stille.
I disse tre Henseender kan Kulturudviklingen komme paa Vildspor.
Den har en Tendens til at slaa ind paa enhver ny Virksomhedsretning,
der frembyder sig, idet den ledes af den øjeblikkelige
Tilfredsstillelse, som derved kan vindes. Der er noget
Blindt og Hensynsløst i Kulturvirksomheden, som baade hænger
sammen med dens Fortrin og dens Fejl. Den lader nødigt
nogensomhelst Mulighed uprøvet; den har sine Følehorn ude i
alle Retninger; derved gøres de betydningsfulde Erfaringer;
men derved kan den ogsaa føre til Disharmoni og Lidelse for
de enkelte Individer og for Slægten.

Kulturen er ikke Genstand for Valg. Den er Fortsættelse
af Naturudviklingen og kan ligesaa lidt stanses som denne.
Trangen til rastløs Udvikling, til at prøve nye Veje og til at
uddanne Specialiteter og frembringe forskelligartede Former er
det store Middel i Naturens Haand til at muliggøre levende
Væseners fortsatte Bestaaen, naar disse tiltage i høj Grad paa
et forholdsvis lille Omraade. Differentiationen og Rastløsheden
udspringe af Befolkningstætheden\footnote[1]{Smlgn. Ch. Darwin: The Variation of Animals and Plants under Domestication. London 1868. I. p. 5---7. --- E. Durkheim: De la division du travail social. Paris 1893. p. 282—294.}. Men fordi Kulturudviklingen,
saaledes som den faktisk foregaar, er naturnødvendig,
er det ikke sagt, at den er beundringsværdig. Etiken føler
ingen særlig Beundring for en Verdensorden, under hvilken
Udvikling ofte ikke er mulig uden Overspænding af Kræfterne
og uden Ensidighed og Spredthed i Evnernes Brug. Den er
% --- p. 148 ---
\setcounter{footnote}{0}%
\opage{148}ikke saa haardhudet, at den over det tilsyneladende Glimrende
i de ydre Resultater kan glemme den Angst og Smerte, den
Sved og det Blod, som de koste. Derfor fordrer den, at de
tunge Byrder skulle lettes, at de splittede Kræfter skulle samles,
og at alle værdifulde Evner skulle udvikles. De Fordringer,
den saaledes stiller til Kulturudviklingen, ere beslægtede med
dem, den stiller til den menneskelige Villen og Handlen idethele
(III, 20). Gennem Indsigt i Betingelserne for Kulturudviklingen
søger den at finde Midler til at ændre, mildne og
humanisere, saa at Kultur og Velfærd kunne blive bragte i saa
nøje Forbindelse som muligt. Den er ikke sentimental og kortsynet,
saa at den skulde glemme, at Fremgangen kun kan ske
gennem Anstrengelse og Lidelse. Kun i Deltagelse i Kulturarbejdet
finder den en virkelig Opgave og et virkeligt Indhold
for Villieslivet. Og særligt indser den, at først gennem Kulturudviklingen
og den Differentiation, denne medfører, komme de
personlige Ejendommeligheder til Udfoldelse og Anvendelse.
Vel gøres hver Enkelt mere afhængig af de Andre; men med
denne gensidige Afhængighed følger en gensidig Supplering og
en inderligere Sammenslutning, end der er mulig, saalænge
Arbejdsdelingen ikke er traadt i Kraft. Baade Mangfoldigheden
og Sammenhængen kunne nu blive større og fastere. Dette
kan Kulturen bevirke --- men den har Sidevirkninger, som
hindre den i at udfolde alle sine lykkebringende Virkninger.

Den etiske Hovedopgave overfor Kulturen er den at indskærpe,
at Livet ikke skal gøres til et blot Middel for upersonlige
Opgavers Løsning. Kulturen skal naa sin Fuldendelse
i Personligheden, saa at denne selv kommer til at staa som
den højeste Kultur. Det personlige Liv maa hverken ødes i
frugtesløs Kamp mod uovervindelig Modstand eller spredes og
indsnevres ved for stor Mangfoldighed eller Ensidighed. Kulturen
indeholder en Mulighed for, at Menneskene kunne komme
til at sætte sig højere Formaal og til at voxe med deres højere
Formaal; det er denne Mulighed, Etiken fordrer benyttet.
Sværmeriet for Naturtilstanden er uetisk, naar det fører til at
opgive Arbejdet. Et karrigt Indhold er let nok at ordne; hvor
der ingen Forskelligheder og Modsætninger er, er det let nok
% --- p. 149 ---
\setcounter{footnote}{0}%
\opage{149}at holde aandelig Disharmoni borte. Ofte kan vel en Indskrænkning
af Fornødenhederne være den eneste Mulighed til at vinde
Sejr over Hindringerne. Men ofte kan netop det, at der føles
Savn og Mangel, vække nye Kræfter og derved muliggøre en
Tilfredsstillelse, som ellers ikke havde været mulig. Ingen af
de to Paastande, hverken den, at Fornødenhederne skulle indskrænkes,
eller den, at de skulle forøges, for at Velfærd kan
opnaas, er derfor ubetinget sand.

Dette faar sin Anvendelse, naar der bliver Spørgsmaal om,
med hvilken Ret man ophæver den Harmoni og Lykke, som er
til Stede, for at føre Menneskene ind paa nye Veje paa den
aandelige eller materielle Kulturs Omraade. Saaledes naar vilde
Nationer rives ud af deres „Naturtilstand“, --- naar man berøver
Barnet dets Illusioner eller vækker nye Fornødenheder
hos det, --- naar man vækker Tvivlen, hvor der før raadede
en blind, men tryg Tro. Afgørelsen kan i saadanne Tilfælde
kun ske ved Hjælp af Velfærdsprincipet. Dersom man er sikker
paa, at der ad den nye Vej vil kunne vindes en ny Harmoni,
en ny Tilfredsstillelse, som har fastere Grund og rigere Indhold
end den tidligere, saa er det berettiget at gribe ind. Og det
bliver Pligt, hvor man staar overfor en Søvngængertilstand, der
kan ende med Ulykke for det sovende Individ selv og Andre.
Et Menneskes Idyl kan medføre Skade og Smerte for Andre
eller dog forringe deres Lykke. Der kan være ubrugte Kræfter,
som Slægten trænger til, skønt Individet hidtil ikke har følt
Trang til at bruge dem. Selvom de skulle vækkes ved et
smerteligt Stød, kan det være Pligt at bibringe det. Det kan i
det enkelte Tilfælde være overordentligt vanskeligt at afgøre,
om et saadant smerteligt Stød skal bibringes. Den etiske Kalkule
vil aldrig kunne føre til fuld Sikkerhed: der maa \textsc{voves},
og der maa gøres et Spring ud i Mørket. Velfærdsprincipet
kan herved kun tjene som ledende Synspunkt. Det strider ikke
mod det, at der voves; men det strider mod det, at Springet
ud i Mørket voves, naar man ikke har et Haab om, at der
gives Lys paa den anden Side af Mørket. Forholdet kan oplyses
ved en Analogi, der dog er Mere end en Analogi. Det
strider ikke mod Erfaringsvidenskabens Metode at opstille
% --- p. 150 ---
\setcounter{footnote}{0}%
\opage{150}Hypoteser, naar man kun fastholder Bestræbelsen for at vinde
Erfaringens Bekræftelse for disse Hypoteser. Enhver Hypotese
er et Vovestykke, men Videnskabens Historie viser os, at uden
saadanne Vovestykker var vor Erkendelse ikke gaaet frem.
Saaledes vil det ogsaa være en Konsekvens af Velfærdsprincipet,
at der maa voves, at der ikke altid kan ventes, til
fuld Garanti for lykkebringende Virkninger er opnaaet. Velfærdsprincipet
maa udtrykkeligt føre til en Betænkelighed ved
at lade Reflexion og Overvejelse af Grunde for og imod spille
for stor en Rolle, da derved den energiske og resolute Handlen,
der faktisk viser sig at være af saa stor Betydning, udelukkes.
Paul Hensel\footnote[1]{Ethisches Wissen und ethisches Handeln. Ein Beitrag zur Methodenlehre der Ethik. Freiburg i. Br. 1889. p. 31.} siger i sin Kritik af Utilitarianismen: „Den
forsigtige Utilitarianer maa i hvert Tilfælde forholde sig afvisende
overfor et Forsøg paa at føre Samfundet ind paa nye
og uprøvede Veje; ti ethvert saadant Forsøg forstyrrer utvivlsomt
megen Lykke, og det er problematisk, om det formaar at
skabe ny Lykke“. Men netop Velfærdsprincipet fører til ikke
at være altfor forsigtig. Forsigtighed er en Borgemesterdyd,
men det vilde ikke stemme med Velfærdsprincipet, om vi alle
vare Borgemestre, og selv Borgemestre kunne være for forsigtige.
Utilitarianismens Tilhængere have især i en tidligere
Tid ofte --- ligesom Empiristerne idethele --- undervurderet
Vovestykkets, Hypotesens Betydning. Men det var en afgjort
Inkonsekvens fra deres Side, og derpaa kan ingen afgørende
Indvending mod Velfærdsprincipet bygges. Naar Hensel desuden
udtaler: „Hvor store Ideer træde i Kamp med hverandre,
paa Vendepunkterne i Historien, hvor en ny Tid kommer til
Verden under tusindfoldige Kvaler, dér maa Utilitarianismen
forstumme“, --- saa overser han, at de store Ideer i Historien
netop stedse pege hen til et højere Livstrin, et Personlighedernes
Rige af fuldkomnere Art, end de bestaaende Samfund
muliggjorde, og at det kun var i Tillid til den større Værdi af
de Goder, dette Rige vilde bringe, at man vovede Kampen.
Men han overser tillige, at Historien faktisk ikke naar sine
% --- p. 151 ---
\setcounter{footnote}{0}%
\opage{151}Resultater hverken ad de korteste eller ad de reneste Veje, og
at man ikke behøver at billige disse Veje, fordi man fastholder
Værdien af det, der gennem mange Forskydninger er opnaaet
ved at følge dem. Resultatet helliger ikke Metoden.

Dersom vi stedse vilde rette vor Færd efter den Trang og
de Fornødenheder, som Menneskene i Øjeblikket havde, saa
vilde vi derved sænke Niveauet for det menneskelige Liv og
den menneskelige Kultur. Velfærdsprincipet fordrer, at vi ikke
sky Striden mod Fordomme og Træghed. Det Bedste, man kan
gøre for Andre, er ofte simpelthen at lade dem føle, at de
endnu i deres Ønsker og Fornødenheder staa for lavt og ikke
stille tilstrækkeligt store Fordringer. Saaledes gaar --- for at
tage et enkelt Exempel --- den store Kunstner ofte sin ensomme
Gang, uforstaaet eller endog miskendt af den store
Mængde. Og dog følger han Velfærdsprincipet, uden at han
maaske tænker derpaa, ved strengt at hævde Kunstens Fordringer.
Han forøger Slægtens aandelige Kapital og skænker
den en Kraft, som i Fremtiden kan komme til at virke i store
Kredse. Kun en kortsynet Opfattelse og Anvendelse af Velfærdsprincipet
holder sig til den i Øjeblikket herskende Trang
i Stedet for at se hen til Slægtens varige Livsbetingelser og
til de stadige Kilder for nyt Liv og ny Virksomhed. \textsc{Goethe}
sagde (i en Samtale med Eckermann): „Jeg har i mit Kald
som Forfatter aldrig spurgt: hvad vil den store Masse, og hvorledes
kan jeg gavne det Hele? Jeg har stedse kun tragtet efter
at gøre mig selv indsigtsfuldere og bedre, at forhøje min egen
Personligheds Værdi og saa stedse kun udtale, hvad jeg havde
erkendt som Godt og Sandt“. Hvis han havde villet rette sig
efter Publikums Smag og Evner, vilde han kunne have gjort
større øjeblikkelig Virkning, men hans Værker vilde saa ogsaa
have havt en kortere Levetid, og han vilde ikke være bleven
den store Lærer for mange Slægter. --- Hvor stor en Horisont
der i det enkelte Tilfælde skal anlægges ved Drøftelsen af en
Handlings eller en Bestræbelses Værdi, derom kan der Intet
siges i Almindelighed. Med en kraftig Overbevisning om at
have Betingelser for at udrette Noget forbinder sig i Regelen
en mere eller mindre bevidst og tillidsfuld Appel til Fremtidens
% --- p. 152 ---
\setcounter{footnote}{0}%
\opage{152}Dom. Der maa voves, hver Gang Noget skal sættes ind i
Verdens store, for os uoverskuelige Sammenhæng. Og selv
hvad den senere følgende Erfaringsbekræftelse angaar, den, der
historisk afgør Vovestykkets Værdi, gælder det, at den ikke
kan blive helt afsluttet nogensinde, saalænge den historiske
Udvikling varer. Der vil stedse kunne blive Tale om en Omvurdering.
Nietzsche har endog forlangt en „Omvurdering af
alle Værdier“. Det er at tage Munden vel fuld; men en
Grænse for, hvorlangt den Omvurdering, der stadigt gaar for
sig og maa gaa for sig, skal udstrækkes, kan ikke paa Forhaand
sættes. Det vil derfor være umuligt at danne et absolut
afsluttet etisk System. Enhver Videnskab, der fastholder Nødvendigheden
af Erfaringsbekræftelse (Verifikation), maa til
enhver Tid staa med Ufuldendthedens Præg. Det er derfor
ingen afgørende Indvending mod den Opfattelse af Etiken, som
her gøres gældende, at den kun fører til „empirisk Vished“\footnote[1]{möjlighet. Upsala 1895. p. 16. 129. il.)}.
Ligesom de menneskelige Vurderingsmotiver ere i stadig Udvikling,
saaledes vil ogsaa en mere udstrakt Anvendelse af
Vurderingsprinciperne stedse være mulig. Vi ere stadigt undervejs
--- og ikke alene vi: Verden, Tilværelsen idethele er
maaske heller ikke færdig, men i stadig Udvikling. Selve den
etiske Stræben er jo et Tegn paa, at Verden ikke er færdig.
Der er en Del af Verdensarbejdet, som kun Mennesker kunne
gøre. En Del af Verdensudviklingen foregaar gennem selve den
menneskelige Handlen.

5. De Veje, den menneskelige Handlen følger for at naa
sine Formaal, ere ingenlunde ligegyldige. Som vi have set, er
det ifølge Velfærdsprincipet en Ufuldkommenhed, naar Vej og
Maal, Aarsag og Virkning staa i udvortes Forhold til hinanden.
Disse Veje ville være des fuldkomnere, jo mere de Midler, der
anvendes, ere i selve Formaalets Aand, eller, som \textsc{Bruno} \textsc{Wille}
har udtrykt det, ere rene Midler. Den bekendte Sætning, at
Hensigten helliger Midlet, gaar ud fra et rent udvortes Forhold
mellem Middel og Formaal. Jo mere Midlet er i Formaalets

*) Axel Hägerström: Undersökning af den empiristiska etikens
% --- p. 153 ---
\setcounter{footnote}{0}%
\opage{153}Aand, altsaa i sig selv er „helligt“, des mere er Formaalet allerede
opnaaet, skønt vi endnu ere undervejs: det er opnaaet i
det Sindelag, der fører til Virken i Formaalets Aand og altsaa
er en Antecipation af Formaalet. Selve den Sætning, at Hensigten
helliger Midlet, kan iøvrigt have sin Berettigelse. Der
kan være Tilfælde, hvor vi i et Formaals Tjeneste gøre og
maa gøre Ting, vi ellers ikke vilde gøre --- ikke fordi de ere
slette, men fordi de ere indifferente: hvad der under andre
Forhold ingen Værdi har for os, faar nu Værdi som Middel.
Den moderne jesuitiske Etik indskrænker Sætningen til at gælde
saadanne Tilfælde\footnote[1]{Cathrein: Philosophia moralis. p. 80.}. Men dertil kommer endnu saadanne Tilfælde,
hvor et mindre Onde maa paaføres, for at et større
Onde kan undgaas. Naar et Menneske, der er ved at dø af
Sult, slaar en Rude ind i en Bagerbutik og tager et Stykke
Brød, foretager han en Handling, der kun helliges ved Hensigten.
Herhen hører ogsaa al Opdragelse gennem Straflidelse.
Efter Velfærdsprincipet er jo Tilføjelse af Smerte berettiget og
nødvendig, naar den er eneste Middel til et uundværligt Gode;
men det er ogsaa den eneste Grund, der kan gøre Smertetilføjelsen
tilstedelig. Naturligvis vil da det Afgørende være, om
Formaalet selv virkeligt er „helligt“; er det ikke det, kan det
ikke hellige Midlet. Er det saaledes egoistiske Formaal eller
Partiformaal (kirkelige, politiske, literære o. s. v.), der ligge til
Grund, saa bliver Sætningen betænkelig. Og selvom Formaalet
virkeligt er værdifuldt, er det ikke sagt, at det formaar at hellige
Midlet. Anvendelsen af et Middel kan medføre Sidevirkninger,
der volde større Fordærv, end Opgivelsen af den gode
Hensigt vilde gøre. Den, der vilde myrde i ædel Hensigt, vilde
sikkert ved et saadant Brud paa de elementære Love for menneskeligt
Samliv frembringe et større Onde, end han kunde afhjælpe
ved at opnaa sin Hensigt. Attentater ramme i Regelen
kun de ydre Følger af eller Symptomer paa en slet social Tilstand,
ikke den egentlige Aarsag, og de bidrage ofte til at gøre
Ondt værre. Betydningsfulde Formaal kræve en lang Udviklingsgang
til deres Virkeliggørelse og muliggøres ofte netop kun,
% --- p. 154 ---
\setcounter{footnote}{0}%
\opage{154}nåar der ikke sker noget Brud paa de elementære etiske Love,
man mener sig berettiget til at bryde i Bevidsthed om sin gode
Hensigt. Bedre sociale Tilstande opnaas ved at virke paa Folkets
Karakter og Tænkemaade, ved at udvikle dets Kræfter og
Hjælpekilder. Det staar endeligt som et psykologisk Problem,
hvorledes den ædle Hensigt og det raa eller grusomme Middel
kunne bo Side om Side i samme Sjæl. Det kan vel kun forklares
ved den dybe Indignation, Modstanden mod det store
Formaal avler: men Handlinger, der blive til paa denne Maade,
ere snarere at opfatte som Symptomer, som Udslag af en
ulykkelig social Tilstand og af den forpinte Stemning, en saadan
har vakt hos idealistiske Naturer, end som Villiesytringer,
der kunne have forbilledlig Karakter og gøre Krav paa „Hellighed“.
Og netop her ligge Sidevirkninger saa overordentligt nær:
gennem Motivforskydning kan det, der først kun modstræbende
valgtes som Middel, brede sig i Sjælen og tilsidst blive Hovedformaalet.
Hvad der indlededes i rent ideal Form, kan gennem
stigende Bitterhed ende i blind Hadefuldhed og Ødelæggelseslyst.
Det er en farlig Vej, man betræder, naar man begynder
at vælge Midler, der først skulle „helliges“. Statskupenes og
Attentaternes Historie vidner herom. --- Foruden den uvilkaarlige
Motivforskydning gives der ogsaa en vilkaarlig Forskydning,
nemlig naar man „drejer Hensigten“ for at faa Lov til at udføre
Handlinger, man har Lyst til. Det var Tilladeligheden af
dette, Pascal beskyldte Jesuiterne for at lære: skønt Kristus
forbyder at gengælde Ondt med Ondt, skulde vi dog have Lov
at forfølge vore Fjender, naar vi blot gøre det --- ikke for at
tilføje dem Ondt, men --- for at hævde vor Ære!\footnotemark[1]

:) Lettres écrites å un provincial. VII.
% CHECK p. 154: note mark without a note

% --- p. 155 ---
\setcounter{footnote}{0}%

\opage{155}\chapterhead{VIII.}{INDIVIDUEL OG SOCIAL ETIK.}

1. Der er nogle Begreber, som uundgaaeligt ville blive
dannede, hvergang en etisk Teori udvikles i sit fulde Omfang,
idet de betegne forskellige Sider ved eller Synspunkter for al
etisk Handlen.

En Handling maa have et Formaal, som bestaar i det
Værdifulde, der tilstræbes. Begrebet om det er Begrebet om
det Gode. Ethvert Formaal staar som et Gode, og dette Begreb
er i sin Almindelighed fælles for alle praktiske Videnskaber,
for Nationaløkonomi og Retslære, saavel som for Etik. (Smlgn. III, 1---14). En Etiks hele Karakter og Retning bestemmes ved
de Formaal, som fra først af lægges til Grund, og derfor opstaa
de største Vanskeligheder i enhver etisk Diskussion, naar
Striden kommer til at angaa de fundamentale Formaal (Grundværdierne,
Vurderingsmotiverne). (Smlgn. III, 12---13). Hvad
der faar Værdi og derfor kan blive Formaal, afhænger af
Livets, særligt Sjælelivets Trang og Stræben og vil være forskelligt
paa Livets forskellige Trin. (Smlgn. I, 4; III, 13. 21).

I Begrebet Formaal ligger allerede Forudsætningen om en
vis Afstand mellem det Tilstræbte og den givne Tilstand. Det
umiddelbart givne Værdifulde behøver ikke at opstilles som Formaal,
uden forsaavidt det gælder at bevare det i Fremtiden.
Naar nu Afstandsmomentet særligt fremhæves, idet Vanskelighederne
ved Formaalets Opnaaelse skærpes, opstaar Begrebet
Pligt. Tilvejebringelsen af det, der er anerkendt som et Gode,
bliver da en bevidst Opgave; Mennesket føler sig bundet til at
arbejde for det og til at skyde modstaaende Tendenser tilbage.
Forudsætningen herfor er, at Mennesket vil fastholde Sammenhængen
og Overensstemmelsen med sig selv, forblive identisk
med sig selv, tro mod sig selv gennem de skiftende Øjeblikke
og Forhold. Kun da vil det hævde Overbevisningen om den
anerkendte Værdis Gyldighed ud over Anerkendelsens Øjeblik.
% --- p. 156 ---
\setcounter{footnote}{0}%
\opage{156}Fra Selvopholdelsens Synspunkt bliver det Pligt at underordne
Øjeblikkets Tilfredsstillelse under Hævdelsen af Livets mere omfattende
Værdier; fra Slægtens Synspunkt bliver det Pligt at
underordne de rent individuelle Interesser under de fælles.
(Smlgn. III, 5---7. 9—10).

Begrebet Pligt angaar særligt de enkelte Handlinger, der
for Selvbesindelsen staa som nødvendig Konsekvens af det anerkendte
Formaal. Lægges derimod Vægten paa de blivende
Karakteregenskaber, som efter deres Natur ville føre til Bestræbelser
for at virkeliggøre Formaalet, opstaar Begrebet Dyd.
Mellem Pligt og Dyd er der i Etiken en lignende Forskel som
mellem Sindsbevægelse og Sindelag i Psykologien. (Smlgn. Psykologi
VI E, 5). Pligt forudsætter en Koncentration af den etiske
Følelse paa et enkelt Punkt, med en bevidst Anstrengelse i det
enkelte Øjeblik; Dyd forudsætter en varig Stemning og Stræben,
uden at den Konsekvens, med hvilken den enkelte Handling
følger af det anerkendte Formaal, behøver at være Genstand
for fuld Bevidsthed. Ligesom Overgangen fra den umiddelbare
Værdifølelse til Pligtfølelsen sker ved stigende Betoning
af Afstandsmomentet, saaledes sker Overgangen fra Pligt til Dyd
ved stigende Indøvelse i de Stemninger og Bestræbelser, det
anerkendte Formaal fremkalder, indtil den enkelte Handling tilsidst
fremgaar af uvilkaarlig Trang.

Begreberne Pligt og Dyd udspringe med direkte Konsekvens
af Begrebet om det Gode. Begrebet Ret er derimod en
indirekte Konsekvens af det. Etisk berettiget er Alt, hvad der
ikke strider imod eller hæmmer den Grundværdi, mit højeste
Formaal forudsætter, og derfor heller ikke strider mod nogen
Pligt eller Dyd. Jeg er i min gode etiske Ret overalt, hvor
min Samvittighed ingen Indsigelse kan gøre, og hvor jeg derfor
ogsaa fordrer, at Andre skulle anerkende min Færd. Udebliver
denne Anerkendelse, appellerer jeg trygt fra den ufuldkomne
Samvittighed til en fuldkommen Samvittighed, som jeg
haaber vil udvikle sig, og som jeg stræber at faa udviklet hos
dem. Det etisk Berettigede er den frie Livsudfoldelse, den frie
Brug af Kræfterne. Livet er ikke til for Samvittighedens Skyld,
men Samvittigheden for Livets Skyld; derfor hævder Samvittig%
% --- p. 157 ---
\setcounter{footnote}{0}%
\opage{157}heden mig i min Færd, naar denne ikke krænker Livets Værdier.
Først og fremmest staar selve Grundværdien som berettiget,
kun underlagt den Betingelse, at den skal virkeliggøres
paa en med dens eget Væsen stemmende Maade. (Smlgn. VII,
5). --- Det etisk Berettigede er ikke blot af negativ Natur. Det
er urigtigt at kalde det det Tilladelige. Deri vilde ligge, at det
blot taales. Men hvor Livet har udfoldet sig frit, uden at Samvittigheden
behøver at gribe regulerende ind, der er det Maal,
hvorefter al etisk Udvikling stræber, allerede naat. Det er af
afgørende Betydning, at det samvittighedsfrie Omraade ikke
indskrænkes. Etiken peger her ud over sig selv, til en Tingenes
Orden, hvor den selv ikke behøver at existere. Da det
er Pligt at arbejde henimod en saadan Tingenes Orden, er det
ogsaa en Pligt at hævde sin Ret, hvor den falder sammen med
Livets egen Ret. Det er her Retten, der begrunder Pligten.
Det etisk Berettigede udspringer af samme Kilde som Begreberne
Pligt og Dyd, men nærmere ved Kildens Udspring
end disse. ---

Enhver Etiks Karakter bestemmes ved det Forhold, i hvilket
disse fire Grundbegreber træde til hverandre. De hænge saa
nøje sammen, at man stedse, paa en mere eller mindre naturlig
Maade, vil kunne finde sin Vej fra ethvert af dem som Udgangspunkt
til de andre tre. De angive de vigtigste formelle
Synspunkter, fra hvilke enhver Handling og ethvert Forhold
kan betragtes i etisk Henseende. Det vil naturligvis være af
afgørende Betydning til Karakteristik af en etisk Teori, hvilket
af disse Begreber den tager som Udgangspunkt. Hvis man vil
grunde Etiken paa Erfaring, kan der dog ikke være Tvivl om,
at Begrebet om det Gode maa være Etikens første Begreb. En
Handling, en Stræben og en Vurdering forudsætte nødvendigvis
et Formaal, og et Formaal forudsætter igen en umiddelbar Værdi.
Dette Udgangspunkt greb da ogsaa den græske Etik, og dens
Betydning for alle Tider beror ikke mindst derpaa. I skarp
Modsætning hertil vilde Kant lægge Pligtbegrebet til Grund og
bestemme de andre etiske Grundbegreber ud fra det. Kun paa
en kunstig Maade kunde han ad denne Vej udlede specielle
% --- p. 158 ---
\setcounter{footnote}{0}%
\opage{158}Formaal\footnote[1]{Smlgn. Den nyere Filosofis Historie.“ II. p. 254 f.}; Hovedformaalet blev for ham selve Pligtvillien, som
han udtalte det i sin berømte Sætning: „Der kan ikke tænkes
Noget i Verden eller udenfor Verden, som kunde kaldes godt
uden Begrænsning, uden ene og alene den gode Villie.“ Paa
Spørgsmaalet om, naar en Villie da er god, svarer Kant: naar
den er i Overensstemmelse med et Fornuftvæsens Natur og derfor
i sine Principer den samme hos ethvert Fornuftvæsen. Paa
denne Maade bliver Pligten selv det højeste Gode. Dyden bestaar
for Kant i den Energi, med hvilken Mennesket øver sin
Pligt (die Stårke des Menschen in Befolgung seiner Pflicht.
Tugendlehre § IX). Der er derfor mange Pligter, men kun én
Dyd og ét Gode. --- Overfor Kant fremhævede \textsc{Schleiermacher},
at man ikke kan ville og handle uden at have et Formaal. (Se
ovenfor III, 1). Han viste nærmere i en Række interessante
Afhandlinger (som findes i andet Bind af hans filosofiske Skrifter),
hvorledes Begreberne det Gode, Pligten og Dyden indbyrdes
føre over til hverandre. I sin egen Etik lagde han Begrebet
om Formaalet til Grund og afledede de to andre Grundbegreber
af det. (Philosophische Sittenlehre. § 113—121). Begrebet
Ret opstiller han ikke i denne Sammenhæng. Dog har
han indirekte forberedt det gennem sin Kritik af Begrebet det
Tilladelige, der --- som han træffende viser --- hører hjemme
i Retslæren, ikke i Etiken, og kun kommer ind i Etiken, naar
man forvexler denne med Retslæren\footnotemark[1]. I samme Retning som
Schleiermacher maa naturligvis Utilitarianismen, og overhovedet
al Velfærdsetik gaa. --- Det tredie Grundbegreb, Dyden eller
den etiske Karakteregenskab, lægges i nyere Tid til Grund i
\textsc{Herbarts} Etik\footnotemark[1]. Etiken bliver her først og fremmest en Beskrivelse
af Egenskaber, en Opstilling af Forbilleder. Foruden
Herbart gaar ogsaa den svenske Idealisme og hos os F. C.
81—84.:) Smlgn. nedenfor VIII, 6, og Den nyere Filosofis Historie.\textsuperscript{2}\footnotemark[2] Il. p.
443—445. Ueber den Begriff des Erlaubten. (Werke zur Philos. II.) p.
% CHECK p. 158: note mark without a note

% CHECK p. 158: note mark without a note

% --- p. 159 ---
\setcounter{footnote}{0}%
\opage{159}Sibberns Moralfilosofi i denne Retning\footnote[1]{Bi berg kritiserede Schleiermachers Lære om de tre Grundbegreber og fordrede Dydsbegrebet lagt til Grund, fordi det angiver den etiske Handlings Kilde (fons actionis), medens Pligtbegrebet angiver den enkelte Handling (actio), der stammer fra denne Kilde, og det Gode det Resultat, der opstaar ved Handlingen (opus morale). (Smlgn. Axel Nyblæus: Den filosofiska forskningen i Sverige. III, 2. p. 228). „Al Sædelighed“, siger Boström (Föreläsningar i Etiken. Udg. af Ribbing. p. 70), „beror paa Villien som Princip eller Udgangspunkt\dots{} Den kan betragtes 1) som hvilende: som Beskaffenhed, men tillige Grund for Virksomheden; 2) som vordende: hvorved den har Ideen til Lov; 3) som tilbleven eller opnaaet, eller som det højeste Gode.“ --- F. C. Sibbern lægger Læren om Sindelaget (Grunddyderne) til Grund, fordi det afhænger af den Maade, paa hvilken Mennesket er indordnet i Livets store Helhed. Derefter følger Læren om Formaalene og tilsidst Læren om Pligterne. Moralphilosophie som Retsindighedsog Tilbørlighedslære.}. Vanskeligheden ved
at gennemføre den beror paa, at en Karakteregenskabs Værdi
kun kan begrundes ved, at den indeholder en Tendens til at
virke i Retning af et værdifuldt Maal. Begrebet Dyd afhænger
ligesom Begrebet Pligt af Begrebet det Gode.

Da jeg i Aaret 1896 holdt Forelæsninger i Zürich over
etisk Principlære, holdt jeg mig væsentligt til Begrebet Formaal
og drøftede de Problemer, deri indeholdes. Da Pligtbegrebet
kun udspringer, hvor et Afstandsforhold indtræder mellem Værdi
og Virkelighed, og da Pligtens Værdi afhænger af Formaalets
Værdi, omtalte jeg ikke særligt det. I en Diskussion, der afholdtes
efter Foredragenes Slutning, blev det indvendt imod
mig, at Pligtbegrebet savnedes i min Fremstilling. Jeg svarede:
„Pligtbegrebet hører ikke til Principlæren. Ti Pligt opstaar
først, naar det principielt Anerkendte skal fastholdes og gennemføres
i Virkeligheden. Jeg beundrer Kants Skildring af Pligtbevidstheden,
der i store, fundamentale Træk fremstiller denne
betydningsfulde Side af det etiske Liv. Men Kant har grebet
fejl deri, at han opfatter Pligtbegrebet rent formelt og ikke som
betinget ved reale Formaal. Hans psykologiske Forklaring er
ikke saa god som hans psykologiske Beskrivelse. Men ethvert
etisk Standpunkt kan optage Kants Beskrivelse af Pligtbevidst%
% --- p. 160 ---
\setcounter{footnote}{0}%
\opage{160}heden, saa snart det drejer sig om, hvorledes de af Villien satte
Formaal og Principer skulle fastholdes og gennemføres i Livets
Kamp.“ ---

De omtalte fire Begreber ville atter og atter komme til
Anvendelse ved de enkelte Spørgsmaal indenfor den specielle
Etik. Men det vilde ikke være hensigtsmæssigt at lægge dem
til Grund for Inddelingen af Etiken. De ere rent formelle Begreber,
og det samme reale Indhold vilde komme igen under
alle fire Synspunkter. Desuden beror Forholdet mellem dem i
hvert enkelt Tilfælde paa individuelle og historiske Betingelser.
Hvad der for det ene Menneske har umiddelbar Værdi, har
maaske for det andet Menneske kun middelbar Værdi og kræver
derfor en særegen Anstrengelse, naar det skal holdes fast: hvad
der for hin er et Gode, bliver for denne en Pligt. Maaske er
hos en Tredie Karakteren saa udviklet, at der uden Anstrengelse
udløses en Stræben efter at hævde det Værdifulde.

I Stedet for den rent formelle Inddeling, der kunde bygges
paa de omtalte Grundbegreber, er det hensigtsmæssigere at
følge en mere real Inddeling\footnote[1]{Se om formale og reale Værdiforskelligheder: Den menneskelige Tanke. p. 241—248. Smlgn. ovenfor IV, 1. i}. Den etiske Verden strækker sig
fra Øjeblikket til Menneskeheden. Hvad der udfylder Øjeblikkene,
vurderes fra den individuelle Livshelheds Synspunkt, hvad der
udfylder Individets Liv, vurderes tillige fra Samfundets Synspunkt,
og ethvert større eller mindre Samfund vurderes tilsidst
fra Menneskehedens Synspunkt. Den mest indgribende Modsætning
vil her være den mellem det enkelte Individ og Samfundet.
De forskellige Arter af Samfund frembyde fælles Træk,
og de bestaa alle af Individer som deres sidste Elementer.
Spørgsmaalet bliver da om Forholdet mellem individuel og social
Etik. Omfatter den ene af disse den anden, eller staa de
sideordnede med lige Ret?

2. Se vi først historisk paa Forholdet mellem Individ og
Samfund, da træffe vi paa tidligere Udviklingstrin ikke Individet
som Indehaver af særlige Pligter og Rettigheder. Det betragtes
kun som socialt Element. Den Vægt, der til forskellige
% --- p. 161 ---
\setcounter{footnote}{0}%
\opage{161}Tider og hos forskellige Folk lægges paa individuelle Dyder
som Selvbeherskelse, Maadehold o. lign., finder sin Forklaring
ved, at det var saadanne Egenskaber, som Stammen til den givne
Tid behøvede. Den Beundring, som næres for de individuelle
Dyder, skyldes altsaa oprindeligt deres sociale Betydning. Den
primitive Etik er den sociale Etik, og den individuelle Etik udvikler
sig af denne. For Grækenlands Vedkommende repræsentere
de filosofiske Skoler Udviklingen af en individuel, til Dels
en individualistisk Etik. Der lægges stigende Vægt paa Harmonien
i Individets egen Natur, paa Individets Sindelag og Karakter,
og Samfundet betragtes tilsidst som blot ydre Middel for
Individets Selvudvikling. I den kristelige Ide om den Enkeltes
Sjæl, som skal reddes for enhver Pris, ligger den samme Henpegen
mod den enkelte Personlighed. De moderne Emancipationsbestræbelser
stille mere og mere Individet frit overfor de
sociale Former og se det, der foregaar i det og med det, som
det egentligt Afgørende. Og dog vil det vise sig, at den sociale
Etik ikke blot historisk, men ogsaa begrebsmæssigt stedse forudsættes
af den individuelle Etik.

3. Man kunde mene, at den individuelle Etik maatte være
den egentlige Etik, da det jo stedse er enkelte, bestemte Individer,
vi paa det etiske Omraade have med at gøre. Etiske
Domme angaa Villen og Handlen, og der villes og handles
stedse af enkelte, bestemte Individer. Et Samfund handler kun,
forsaavidt de enkelte Individer handle. Og ethvert af disse
handlende Individer synes tilsidst kun at have med sig selv at
gøre. Det skal handle efter sin egen Samvittighed; det kan
ikke direkte sætte sig ind i, hvad der foregaar i Andres Indre;
det har da kun at sørge for, at det ikke selv tager Skade paa
sin Sjæl; og naar Enhver gør saa, vil det staa vel til med
Samfundet!

Denne Synsmaade har øvet stor Indflydelse. Den kommer,
som allerede berørt, først frem i den græske Filosofi. Det var
især Kynikerne og Stoikerne, som hævdede et saadant Selvhensyn,
en højere Selvopholdelse som det Vigtigste i etisk Henseende.
De ere i denne Henseende Kristendommens Forgængere. I vor

H. Høffding: Btik. 
% --- p. 162 ---
\setcounter{footnote}{0}%
\opage{162}moderne Tid er S. \textsc{Kierkegaard} den store Repræsentant for
den Enkeltes Etik. Under Paavirkning fra ham findes den samme
Tankegang i \textsc{Heegaards} Bog „Om Intolerance“, hvor Trangen
til Samkvem med Andre endog karakteriseres som en Trang til
„at se, høre, nyde, eller at beundres, smigres og tilbedes“. I
\textsc{Henrik} \textsc{Ibsens} „Dukkehjem“ kommer den samme Tendens frem,
naar Nora i Kraft af det Princip, at vi have Pligter mod os
selv før imod Andre, drager sig ud af sine allernærmeste Forhold
--- for at blive Menneske!

Denne Opfattelse er trods al sin Ensidighed fremgaaet af
en betydningsfuld etisk Tankegang. De enkelte Personligheder
ere det levende Stof, i hvilket det Etiske har sin Virkelighed.
Men deraf følger ikke, at hele Tilværelsen, Samfundet og alle
andre Individer skulle gøres til Midler for den Enkeltes Selvudvikling.
Det vilde medføre en Isolation, en ensidig Optagethed
af sig selv, der let vilde gaa over til Egoisme. Der gives ligesaavel
en asketisk Egoisme, som der gives en nydende eller
en magtsyg Egoisme. Den Enkelte mister Evnen til Selvhengivelse.
Blikket fæstes mod den egne Velfærd, ikke mod den
store, fælles Velfærd, af hvilken den egne er en Del.

Der næres her ængstelige Forestillinger om, hvad Selvet
eller Personligheden taaler og behøver. Man tror kun at kunne
bevare sig selv ved uafbrudt at have Opmærksomheden henvendt
paa sig selv. Men et Selv er ikke saa skrøbelig en
Ting. Det bliver netop skrøbeligt ved den megen Selvomsorg.
Et sundt og kraftigt Selv naas ved rask og kraftigt Arbejde paa
virkelige Opgaver, og et stort og omfattende Selv naas kun,
naar der arbejdes paa fælles Opgaver, naar man gør sig til Et
med Noget, der er større end Selvet. Kun ved Selvforglemmelse
voxer og udvikles Selvet. Ved at leve i Andre og i Andet end
sig selv faar Individet nye og kraftige Impulser ogsaa for sit
eget individuelle Liv. Det hæves op over sin oprindelige snevre
Kreds og fyldes af et rigere Indhold. Alt hvad et Menneske
lever i og for gennem Sympatien, er en Berigelse og Udvidelse
af dets Personlighed. Det er her paa ethvert Punkt baade Formaal
og Middel. H. C. \textsc{Ørsteds} Grundsætning: „Glem dit Selv,
men tab det ikke!“ udtaler træffende, hvad det her gælder om.
% --- p. 163 ---
\setcounter{footnote}{0}%
\opage{163}For den individuelle Etiks egen Skyld maa den altsaa ikke
gøres til den eneste Etik.

Overvurderingen af den individuelle Etik hænger sammen
med en abstrakt Forestillingsmaade, som drager Mennesket ud
af de bestemte, historisk givne Forhold og først i det saaledes
isolerede Væsen mener at have det sande Menneske. Men jeg
er ikke først Menneske i al Almindelighed, og derefter Menneske
i disse specielle Forhold. Jeg er ikke først Menneske,
og saa Dansk, Borger, Familiefader o. s. v. Individualismen
kommer naturligt frem i Tider, hvor de historisk givne Livsforhold
tabe deres absolute Myndighed og ikke staa som de
eneste mulige og berettigede. Da har den den store Betydning
at nære en Frihedsbevidsthed, som kan blive en frugtbar Kilde
for ny Udvikling. Men ikke blot medfører den, som jeg har
søgt at vise, sine store Betænkeligheder fra selve den individuelle
Etiks Standpunkt; --- det vil ikke være muligt ud fra en saadan
Betragtningsmaade at begrunde en Opfattelse af det menneskelige
Samfundsliv, hvorved dette bliver Mere end blot Middel.
Dettes Interesser og Opgaver kunne i det Højeste blive Øvelsesexempler
for de enkelte Individer, Øvelsesexempler, som i og
for sig ligesaa godt kunde ombyttes med andre.

En helt anden Art af Individualisme fremtræder i den
nyeste Tid især i litterær og æstetisk Form. Frygt for Nivellering
og Uniformering og for Mængdens aandelige Tyranni har især
virket til at fremkalde den, og den fremtræder som en Kultus
af det Sære, det Ensomme, det Ørkenagtige, det Egensindige,
det Vovede. Paa det praktiske Omraade fører dette til Kultus
af «Overmennesket“, og Overmennesket tænkes især som den
overlegne Magt, der er hævet ikke blot over Fordomme og
Sædvaner, men ogsaa over de Modsætninger, indenfor hvilke
Massens Liv bevæger sig. Hvad der især synes at vække Beundring
overfor det tænkte Overmenneske, er den Foragt, dette
fra sit Standpunkt formaar at føle overfor den store Mængde.
Det er dette Afstandsmoment („das Pathos der Distanz“), der
stilles i Forgrunden. Dermed forbinder sig da ofte den Længsel
og Beundring, som de Svage og Syge naturligt føle overfor
den store Kraft og den store Sundhed. Disse forskellige Ten%
% --- p. 164 ---
\setcounter{footnote}{0}%
% CHECK p. 164: notes paired with the marks in order (the layer misread a number)
\opage{164}denser lade sig paavise hos \textsc{Friedrich} \textsc{Nietzsche}. Saaledes
fordrer han (Jenseits von Gut und Böse. Aphor. 258) af en herskende
Stand, at den ikke skal begrunde sin Værdi ved den
Funktion, den udøver i Samfundet, men at „dens Grundoverbevisning
skal være, at Samfundet ikke er til for Samfundets
Skyld, men kun som Underbygning og Grundlag, hvorpaa en
udsøgt Art af Væsener formaar at hæve sig til sin højere Opgave
og overhovedet til en højere Væren.“\footnote[1]{Jeg kan her ikke drøfte, hvorvidt vi i saadanne Udtalelser har Nietzsches egentlige Grundtanke. Herom maa jeg henvise til mit Skrift Moderne Filosofer (1904).} Dog kommer
Problemet om Forholdet mellem individuel og social Etik igen
indenfor hin udsøgte Art af Væsener, nemlig ved Spørgsmaalet
om deres Forhold til hverandre og til det Samfund, de danne.
Om et af de Overmennesker, Nietzsche mest beundrede, Napoleon,
er det blevet sagt, at han ikke kunde taale stolte Sjæle.
Det blev hans Nemesis. Mange af Nietzsches Tilhængere undgaa
af gode Grunde denne Nemesis. Der røber sig idethele
hos dem en ret uinteresseret Følelse, forsaavidt som de selv
synes meget langt fra at indtage det høje Stade, hvorfra den
store Foragt kunde motiveres.

Paa en mere filosofisk Maade fremtræder Individualismen
i den idealistiske Anarkisme, saaledes som den f. Ex. læres af
den tyske Forfatter \textsc{Bruno} \textsc{Wille}. Bruno Willes Bog Die Philosophie
der Befreiung durch das reine Mittel (Berlin 1894) giver
en interessant og lærerig Fremstilling af dette Standpunkt. Der
fordres her, at Individualiteten skal udvikle sig indefra, og at
alle Idealer skulle være selvskabte. Der protesteres mod al
aandelig og fysisk Tvang, og en saadan Tvang finder Wille ogsaa
i moralsk Paavirkning fra Andres Side. Idet han nu opfatter
Samvittigheden som et blot ydre, socialt Produkt, erklærer han,
at det frie Fornuftmenneske maa være samvittighedsløst.

Bruno Wille bygger paa en Tanke, der ogsaa har spillet
en betydningsfuld Rolle i den danske Tænknings Historie\footnote[2]{Se min Bog Søren Kierkegaard som Filosof. Kbhvn. 1892. p. 22—27; 98—107.}').
Personligheden maa selv frembringe --- eller paany frembringe
% --- p. 165 ---
\setcounter{footnote}{0}%
\opage{165}--- den Sandhed, der skal være Sandhed for den. De etiske
Domme maa udspringe af den Enkeltes egen Erfaring og Overbevisning;
ellers ere de kun et Ekko af Omgivelsernes Domme.
Enhver skal være sin egen indre Dommer, sin egen Autoritet.
Al Indvirkning udefra skal gaa ud paa at frembringe Evnen til
etisk Selvstændighed. Men hvad den idealistiske Anarkisme
overser, er den pædagogiske Betydning, som Efterligning og
Lydighed, overhovedet Autoritetsforholdet kan have. Den Enkelte
behøver Hjælp fra Andre til at blive fri og selvstændig.
Denne Hjælp kan meget ofte bedst ydes ad anden Vej end
ved moralsk Bedømmelse. Direkte Moraliseren hindrer ofte Vækkelsen
af Selvvirksomheden, idet den etiske Vurdering bibringes
Individet helt færdig og formuleret i Stedet for at frembringes
af det. Dog er direkte Moralisering virkeligt et „rent Middel“
(smlgn. VII, 5), naar den gaar ud paa at vække Selverkendelsen
og Udviklingstrangen. Denne Vækkelse kan ofte være smertelig;
men der gives Smerter, som ere nødvendige Gennemgangsled
til et højere Trin.

Forgæves vil man stræbe efter at naa det Højeste ad Isolationens
Vej. Ikke blot faar den Enkelte sin Opdragelse i Samfundet,
men ogsaa sit Indhold og sine Opgaver faar han kun
ved Hengivelse til Slægtens Liv. Kun her bliver han delagtig
i de højere Formaal, med hvilke han kan voxe. I den Kultus
af Foragten, som den æstetiske Individualisme i den nyeste Tid
er tilbøjelig til, kommer Umuligheden af Isolation for Dagen
ved en ejendommelig Nemesis. Naar det er, hvad Nietzsche
har kaldet Afstandens Pathos, der bestemmer det væsentlige
Indhold, som dette Standpunkt faar, saa fører det netop til Afhængighed
af den store Mængde; ti der maa være Nogen, man
kan foragte, Nogen, man kan føle sig i Afstand fra. Det ophøjede
Stade bruges væsentligt til derfra at se ned paa dem,
der ikke have naat det. Man glemmer, at den sande aandelige
Magt og Overlegenhed viser sig i, at man har Evne til at bære
ikke blot sit eget, men mange Andres Liv. Der gives, siger
\textsc{Spinoza}, ingen Maade, paa hvilken et Menneske bedre kan
ytre sin Aandskraft end ved at udvikle Andre til at leve efter
deres egen frie Erkendelse.

% --- p. 166 ---
\setcounter{footnote}{0}%
\opage{166}4. Som Modstykke til denne Opfattelse optræder den, der
fordrer fuldstændig Opgivelse af sig selv og betragter Selvforglemmelse
som den eneste Dyd. Indenfor den nyere filosofiske
Etik er et saadant Standpunkt gjort gældende af J. G. \textsc{Fichte}
og A. Comte. Det sande Liv bestaar efter Fichte i at ofre sig
for Slægten: „Der gives kun én Dyd, at glemme sig selv som
Person, og kun én Last, at tænke paa sig seiv\dots{} Den, som
idethele taget tænker paa sig selv som Person og begærer Liv,
Tilværelse og Nydelse uden i og for Slægten, han er --- ved
hvilke gode Gerninger han saa ellers paa andre Omraader søger
at skjule sin Vanførhed --- dog i Bund og Grund et simpelt,
lille, slet og derfor usaligt Menneske“\footnote[1]{Politique positive. IV. p. 45—49. --- Der er paa dette Punkt en vis Modsætning mellem Comtes tidligere og hans senere Standpunkt \{Den nyere Filosofis Historie.'\textsuperscript{2} Il. p. 346 f. og 357), som danner en Analogi til Modsætningen mellem Fichtes to Stadier (ib. 142 f. og 151 /' f.).}. Comte betegner ved
det af ham indførte Ord Altruisme\footnotemark[1] den fuldstændige Opgaaen
i Andre, den absolute Modsætning til Egoismen. Vel udtaler
Comte, at Altruismen ikke undertrykker Selvhævdelsen; men
dennes Berettigelse skal dog kun bero paa, at den helliges ved
at arbejde i hins Tjeneste: baade Pligt og Lykke bestaa i vivre
pour autrui\footnotemark[1].

Ogsaa denne Opfattelse er fremkaldt ved bestemte historiske
Forhold, for Fichtes og Comtes Vedkommende ved Modsætningen
til det 18. Aarhundredes Individualisme og Lykkemoral. Den
gaar paa sin Vis, ligesom den individualistiske, ud fra en skarpere
Modsætning mellem Selvhensyn og Hensyn til Andre, end
der ifølge de virkelige Forhold behøver at være. Af Velfærdsprincipet
følger, at den Enkelte kun er En blandt Mange; men
der følger tillige, at den Enkelte ogsaa i Virkeligheden er En
blandt de Mange. Enhver tæller med, og Giveren skal ikke
slette sig ud for Modtagerens Skyld, ligesaa lidt som Modtageren
skal gøres til blot Middel for Udviklingen af gode Egenskaber
hos Giveren. Kun den isolerede, i sig selv hildede Selvtilfredsstillelse
strider mod Velfærdsprincipet. Dersom Trangen

!) Die Grundziige des gegenwärtigen Zeitalters. Berlin 1806. p. 70 f.

 Af alter (den anden, d. v. s. Næsten), ligesom Egoisme af ego (jeg).
% CHECK p. 166: note mark without a note

% --- p. 167 ---
\setcounter{footnote}{0}%
\opage{167}til Selvopholdelse, Selvhævdelse og Selvudvikling var af det
Onde, vilde vor inderste Natur være ond, og al Etik vilde da
være selvmodsigende og umulig.

Vi vinde kun Kræfter og Midler til at virke for det, der
er forskelligt fra os selv og større end os selv, derved, at vi
bevare og udvikle vort eget Væsen. Hvis Etiken fordømmer
Selvopholdelsesinstinktet, fordømmer den sine egne Midler. Netop
naar man stiller sig paa Altruismens Grund, bliver det en særlig
Opgave for det enkelte Individ at styrke og udvikle sine Kræfter
og Evner. Selvhævdelse og Selvudvikling arbejde i samme Retning
som Altruismen. Ja, som SPENCER bemærker, den Lykke,
som fra et sundt og frisk Individ kan udstraale over Andre, er
ofte mere værd end den, man bevidst og med stor Opofrelse
kan berede dem. „Skønt vore etiske Ræsonnementer tage lidet
Hensyn dertil, saa er det dog aabenbart, at da Lykke og Sorg
ere smitsomme, er den Selvomsorg, der fører til Sundhed og
Oprømthed, en Velgerning mod Andre“\footnote[1]{Data of Ethics. London 1879. p. 194.}. Og dette gælder ikke
mindst om den højere, ideelle Udvikling af Personligheden ved
Arbejde for Interesser, der maaske ikke direkte komme Andre
til Gode, som f. Ex. ved kunstneriske og videnskabelige Bestræbelser.

Ogsaa i anden Henseende vilde Etiken ved at lægge den
absolute Altruisme til Grund komme i Modsigelse med sig selv.
Al Opofrelse for Andre gaar jo dog ud paa at støtte deres Selvopholdelse,
Selvhævdelse og Selvudvikling, og hvorledes kunne
disse være berettigede hos Modtagerne, naar de ikke ere det
hos Giverne? Den absolute Opofrelse for Andre vil endog kunne
have en skadelig Indflydelse paa disse. Den, der stadigt modtager
Opofrelser fra Andres Side, bliver let til Egoist, naar han
da ikke gaar helt i Søvne uden at vide, hvad der gaar for sig.
Han gøres til et passivt Væsen, der tager, hvad der gives, uden
at sætte sig ind i, hvad der ofres for ham. Hans Evne til Arbejde
svækkes, han bliver et Snyltedyr, ikke et selvstændigt,
selvvirkende Væsen. Den rette Opofrelse maa være den, der
hverken berøver Giver eller Modtager deres Selvstændighed,
% --- p. 168 ---
\setcounter{footnote}{0}%
% CHECK p. 168: notes paired with the marks in order (the layer misread a number)
\opage{168}men gør dem begge til selvstændige Medlemmer af Personlighedernes
Rige.

Saadanne Standpunkter som Fichtes og Comtes nærme sig
til at opfatte Samfundet eller Slægten som et mystisk Hele, der
har en Existens afset fra de enkelte Individer. En lignende
Tendens fremtræder hos Hegel\footnote[1]{Se Den nyere Filosofis Historie.'\textsuperscript{2} II. p. 181.} og i den nyeste Tid hos
\textsc{Wundt} og Paul Carus. Efter Wundts Opfattelse\footnote[2]{Wundt: Ethik. Eine Untersuchung der Thatsachen und Gesetze des sittlichen Lebens.\textsuperscript{3} II. p. 116.} bestaar
det «offentlige Vel“ ikke i Summen af Velfærd hos saa mange
Enkelte som muligt, og det almindelige Fremskridt bestaar ikke
i Fremskridtet hos saa mange Individer som muligt. Ti ethvert
Individ er et Døgnvæsen: „Hvor rigt velsignet og hvor fuldkommen
den individuelle Existens saa ogsaa maatte være, den
er dog blot en Draabe i Livets Ocean. Hvad kan individuel
Lykke og individuel Smerte betyde for Verden?“ Og Carus\footnote[3]{P. Carus: The Ethical Problem. Chicago 1890. p. III kfr. p. 33, 38, 40. --- Smlgn. Diskussionen mellem Carus og mig i The Monist. Juli 1891. Jeg kan ikke se, at Carus ved sine nærmere Forklaringer berøver sin Ide dens Uklarhed, saalænge han (ligesom Wundt) negter, at man i Etiken stedse maa tænke sig et Indbegreb af Individer (af større eller mindre Omfang, maaske uoverskueligt Omfang), naar man vil operere med Begrebet Samfund. L)}
betegner Samfundslivet som „overindividuelt“; Etikens Formaal
er efter ham „hverken Selvets egen eller andre Individers Velfærd,
men Velfærd for de overindividuelle Interesser.“

Hvis Begrebet Samfund skal kunne anvendes videnskabeligt,
maa det ske saaledes, at man paa ethvert Punkt kan afgøre,
hvilken Gruppe af Individer det repræsenterer. Dette Begrebs
store Betydning beror paa, at det udtrykker de samtidigt
og successivt existerende Individers fælles og blivende Interesser
i Modsætning til enkelte Individers Interesser eller til en lille
Gruppes eller en enkelt Generations Interesser. Den etiske Betragtning
fører (paa det Standpunkt, hvor vi her have stillet
os) til at betragte de menneskelige Handlinger ikke blot med
Hensyn til den Enkeltes eller en begrænset Kreds's eller Tidsalders
Velfærd, men med Hensyn til Samfundets eller hele
% --- p. 169 ---
\setcounter{footnote}{0}%
\opage{169}Slægtens Velfærd, saavidt vi kunne efterspore dennes Betingelser.
Men saasnart det viser sig umuligt at omsætte Begrebet
om et Samfund til Begrebet om en Gruppe af Individer, der
leve under visse bestemte Forhold, giver hint Begreb os ingen
etisk Belæring; ingen etiske Normer kunne udledes af det.
Mystik træder i Stedet for Tænkning.

En saadan Mystik kan have sin Betydning. Det kan vise
sig umuligt at forme et tydeligt Begreb om den store Mængde
af menneskelige Interesser, i hvilke en Handling eller en Livsordning
under visse Forhold kan gribe ind paa afgørende Maade.
Udtryk som Samfund og Slægt betegne meget godt det Uafsluttede
og Uoverskuelige i mange af disse Virkningsrækker.
Og de betegne tillige de Muligheder eller Dispositioner, den
potentielle Energi, som menneskeligt Arbejde kan tilvejebringe
og ophobe, og som maaske først i fjerne Tider kunne omsættes
til aktuelle Værdier, der føles af virkelige, personlige Væsener.
Men naar man erklærer en saadan Omsætning for umulig eller
unødvendig, bliver man stikkende i Mystiken. En Velfærd, der
ikke paa et eller andet Stadium er bestemte Individers Velfærd,
er en Selvmodsigelse. Og en Handling eller en Livsordning,
der ikke paa et eller andet Tidspunkt leder til Velfærd for bestemte
Individer, har ingen etisk Værdi. En Mulighed, der aldrig
kan blive til Virkelighed, er en Umulighed.

Det enkelte Individ er kun en Draabe i Oceanet. Og
Oceanet er ikke til for den enkelte Draabes Skyld alene. Men
hvad er et Ocean, der ikke bestaar af Draaber? Og vil ikke
hele Oceanet være klart, naar hver enkelt Draabe er klar? Og
kun da er det fuldkomment klart.

Ligesom der er dem, der ikke kunne se Skoven for lutter
Træer, saaledes er der dem, der ikke kunne se Træerne for
lutter Skov. I Etiken viser dette sig ved, at menneskelig Stræben
betragtes som Middel for overmenneskelige Formaal. Enhver
Etik, der ikke vil opgive Muligheden af fremskridende Verifikation,
maa fastholde, at Samfundet stedse repræsenteres af bestemte
Individer. Derfor behøver den ikke at overse, at etisk
Virken, ligesom al Kraftudfoldelse, hænger sammen med den
hele Verdensproces.

% --- p. 170 ---
\setcounter{footnote}{0}%
\opage{170}5. Man kunde forsøge at gaa en Middelvej og stille individuel
og social Etik ved Siden af hinanden som to selvstændige
Omraader. Saaledes antage \textsc{Bentham} og \textsc{Stuart} \textsc{Mill}, at der
gives et Omraade, hvor Individet er eneraadende, og et andet,
hvor sociale Hensyn raade. \textsc{Mill} mener endog, at saadanne
Fejl, der kun angaa os selv og ikke faa Indflydelse paa Andres
Velfærd, ikke egentligt ere „moralske“ Fejl. Taabelig Adfærd,
slet Smag, lave og sanselige Drifters Herredømme, Mangel paa
personlig Værdighed og paa Selvagtelse kunne efter Mill være
tilstede, uden at Andres Interesser derved i mindste Maade berøres.
Ingen er, mener Mill, sine Medmennesker i mindste
Maade ansvarlig for sin Selvudvikling, „da det ikke er af Hensyn
til Samfundets Vel, at man skal staa til Ansvar herfor“\footnote[1]{Om Friheden. Dansk Overs. p. 141 f. --- Benthams Udtalelser om den „private Etik“ (Principles of Morals and Legislation. XVII, 3, 20) ere ikke ganske tydelige, og der fremtræder hos ham en Tendens til helt at underordne den sociale Etik under den individuelle. Smlgn. ovenfor III, 16. --- Hos nogle danske Forfattere i den nyeste Tid (Starcke og Bang) er en lignende dualistisk Teori som Mill's kommen frem. Se ovenfor 111, 21.}.

Hvad Mill bekæmper, er Statsmagtens og den offentlige
Menings (det „moralske Politis“) Indblanding i Individets inderste
Anliggender. Han frygter for Aandstyranni og usund Tvang,
hvis man gaar ud fra, at der Intet gives i et Menneskes Adfærd,
som kun vedkommer ham selv. Men han sammenblander
her to Ting. Et er det, at Individet helt igennem, selv i sit
inderste Indre og med alle sine Egenskaber, maa betragtes og
betragte sig selv som Led af Slægten, og at hans Villen og
Handlen derfor ere underlagte en etisk Vurdering, hvis Maalestok
er given i Principet om almindelig Velfærd, --- et ganske
Andet er det, hvem der i det enkelte Tilfælde skal optræde
som den etiske Doms Forkynder og Haandhæver. Etisk set
staar og falder Enhver først og fremmest for sin egen indre
Domstol i Samvittigheden, og Etiken er en Lære om, hvilke
Domme denne indre Domstol konsekvent maa fælde. Den Grad,
i hvilken et Menneske udvikler sine Evner og Kræfter, maa
fremfor alt Andet falde ind under denne Domstol. A. S. \textsc{Ørsted}
% --- p. 171 ---
\setcounter{footnote}{0}%
% CHECK p. 171: notes paired with the marks in order (the layer misread a number)
\opage{171}har med Rette sagt, at værre end at begaa et Brud paa Ejendomsretten
er det at bortødsle den Kraft, med hvilken vi kunne
virke i Fremskridtets Tjeneste, ved Egennytte, Forfængelighed
eller dorsk Uvirksomhed\footnote[1]{Om Grænserne mellem Teori og Praxis i Sædelæren. (Eunomia. Første Del. København 1815). p. 105.}. Det er dog derfor ikke sagt, at det
er hensigtsmæssigt at tilstede andre Mennesker en Ret til at
blande sig i den Enkeltes indre etiske Husholdning. Man kan
indrømme det Ene og negte det Andet. Mill derimod negter
Selvudviklingens etiske Betydning, fordi han frygter for derved
at give Andre Ret til at trænge ind i Individets private Færd.
Men han kunde naa sit Maal, nemlig Hævdelsen af den individuelle
Frihed overfor socialt Tryk og moralsk Politi, uden at
udelukke Selvudviklingen fra det etiske Omraade. Intet er betænkeligere
end den store Misforstaaelse, at Etik og offentlig
Mening ere det Samme, og at et Forhold ophører at være etisk
Bedømmelse underlagt, fordi ingen ydre Myndighed har Ret til
at paatale det\footnote[2]{I sit fortrinlige Værk The Utilitarians (III, p. 286—299) har Leslie Stephen i lignende Retning som jeg kritiseret Mill's Lære om Grænsen mellem individuel og social Etik.}.

Der gives slet Intet i Individets frie Udvikling, som ikke
kan faa etisk Betydning. Hvad der optræder i den Enkeltes
indre Liv, kan maaske være den første Spire til store sociale
Processer. Enhver individuel Udvikling kan frembringe nye
Former og Forbilleder; de ensomme Stier, ad hvilke den Enkelte
vandrer, kunne senere blive alfare Veje. Den Energi, den
Troskab og den Visdom, som han anvender til at opsøge og
bane disse Veje, kunne blive af afgørende Betydning. Det er
ikke muligt at skille de Enkeltes Selvudvikling og indre Husholdning
fra Samfundets og hele Slægtens Liv. Selv i Legen,
i Kunsten, i Troen og i det frie Stemningsliv hos den Enkelte
er det Etiske potentielt tilstede. Ligesom der i den ydre Natur
kun tilsyneladende gives en absolut Hvile, men Energien stedse
er tilstede, skønt ofte i saakaldet potentiel Form, saaledes gives
der i det menneskelige Liv intet Punkt, som ikke muligvis kan
% --- p. 172 ---
\setcounter{footnote}{0}%
\opage{172}underligge en etisk Dom. Man har særligt villet hævde det
erotiske Forhold og Religionen som Privatsager. Det er vistnok
sket af taktiske Grunde, nemlig for i Æstetiken ikke at blive
forulempet af Moralen, og for i Socialpolitiken at kunne bevæge
sig uden at støde sammen med den kirkelige Tro. I Virkeligheden
ere baade Kønsforholdet og det religiøse Forhold af
meget stor etisk Betydning. Netop paa disse Omraader have
Egenskaber som Mod, Troskab, Ærlighed og Sandhedskærlighed
den største Værdi; netop her kunne lavt Sindelag, Vaklen og
Ængstelighed blive skæbnesvangre. Og den blotte Konstateren
af saadanne Egenskabers Tilstedeværelse er jo en etisk Dom.

6. Det indre Forhold mellem individuel og social Etik er
egentligt givet med Opstillingen af selve Velfærdsprincipet. Naar
Grundlaget for Vurderingen er den universelle og uinteresserede
Sympati, og Maalestokken den almindelige Velfærd, saa følger,
at den individuelle Etik maa være underordnet den sociale, uden
derfor at forsvinde i Forhold til denne. Den fulde Begrundelse
af de individuelle Dyder, Pligter og Rettigheder vil kun kunne
naas ved at betragte Individet i det sociale Milieu, som Led af
den kæmpende og stræbende Slægt, hvis Arv det har modtaget
i sin Natur og i sine ydre Livsforhold, og hvis Udvikling det
skal føre videre, saavidt det formaar. Gennem Deltagelsen i
Slægtens Liv og Arbejde vil det udvikle sit eget Væsen.

Den nærmere Bestemmelse af Forholdet mellem individuel
og social Etik vil kunne vindes ved Hjælp af nogle Sætninger,
der kunne udledes af det almindelige Velfærdsprincip.

Den simpleste Form, i hvilken dette Princip fremtræder,
er den, at Tilføjelsen af Smerte stedse skal begrundes, medens
Frembringelsen af Lyst og Glæde i sig selv er berettiget. Anvendelsen
af denne Sætning er kun i de allersimpleste Tilfælde
simpel og let, men man vil ved nærmere Betragtning finde
Sætningen som sidste Forudsætning ogsaa i sammensatte Tilfælde.
Des mærkeligere er det, at Sætningen er funden betænkelig;
den siger egentligt kun, hvad Alle --- med Undtagelse
af de absolute Asketer, om der gives saadanne --- maa indrømme.
Selv Asketerne have i Regelen betragtet Smerten som
Bods- og Øvelsesmiddel, ikke som Formaal i og for sig.

% --- p. 173 ---
\setcounter{footnote}{0}%
\opage{173}I denne simple Sætning kan nu ved en nærmere psykologisk
Betragtning en anden meget vigtig Sætning findes indesluttet.

I Smerten kundgør der sig en Hæmning eller en Opløsning
af Livet. Dette er tydeligt ved legemlig Smerte, f. Ex.
Smerten ved Saar, Tandpine, Stensmerter, Kræftlidelse. Smerten
svarer her til et Overgangsstadium mellem det uhæmmede, udelte
Liv og den fuldstændige Hæmning og Opløsning, der kommer
med Døden. En græsk Forfatter har sagt, at hvis Mennesket
var et absolut enkelt Væsen, vilde det ikke kunne føle Smerte:
ti da vilde ingen Opløsning kunne finde Sted. Men det Omvendte
gælder ogsaa: kun fordi Livets Elementer stedse ere i
den inderligste Enhed og Sammenhæng hos det levende Væsen
og ikke udgøre et blot Aggregat, kun derfor kan der føles
Smerte. Det gælder for det sjælelige Liv saavel som for det
legemlige. Ogsaa aandelig Smerte udtrykker en Hæmning eller
en Opløsning. Sorg over et Tab optræder snart som en Hæmning,
idet Tanken om det Tabtes Uigenkaldelighed fortrænger
alle andre Tanker, snart som en Opløsning, idet Modsætningen
mellem det Tabtes Værdi og Tabet truer med at sprænge Bevidstheden.
Tvivlen er en Ulysttilstand, idet Bevidstheden drages
i modsatte Retninger af stridende Muligheder og Tendenser.
Angerens Smerte betinges ved den voldsomme Modsætning mellem
vort anerkendte Ideal og Tanken om vor virkelige Villen og
Handlen. Paa den anden Side er Alt, hvad der fremmer Bevidsthedslivets
Enhed og Harmoni, forbundet med Lystfølelse.
Vort Bevidsthedsliv stræber uvilkaarligt efter at overvinde de
Lidelser og Smerter, der ikke kunne fjernes, ved Opbyden af
højere, sammenfattende Kræfter. Der danner sig en ny Følelse,
hvor Bitterheden er optagen som Element i en mere sammensat
eller blandet Følelse. De fleste af vore saakaldte højere Følelser
have denne sammensatte Karakter. Men paa alle Trin
staar vort Følelsesliv som Udtryk for vort Livs Fremgang eller
Tilbagegang. Forskellen mellem Følelserne beror i denne Henseende
paa den mere eller mindre centrale Stilling, de indtage
indenfor vort Livs Helhed\footnote[1]{Smlgn. om de her benyttede psykologiske Synspunkter min Psykologi VI B, C og D.}.

% --- p. 174 ---
\setcounter{footnote}{0}%
\opage{174}Formedelst denne nøje Sammenhæng mellem Følelsen og
Individets Væsen idethele fører Velfærdsprincipet umiddelbart
over til den frie Personligheds Princip, d. v. s. til den Grundsætning,
at intet personligt Væsen tør behandles og betragtes
kun som Middel, men stedse tillige skal være Formaal. I de
personlige Væsener have vi jo de Centralsteder, til hvilke vore
Handlingers Virkninger forplante sig, og hvor deres vigtigste
Betydning fremtræder. Paa disse Punkter i Verden føles Livets
Værdi, --- og Livets Værdi existerer (som aktuel Værdi) kun,
naar den føles. En Krænkelse eller en Hæmning her vil altsaa
nødvendigvis forringe Livets Værdi, og hvis der ikke føles nogen
Lidelse ved Hæmningen, fordi denne frembringer Sløvhed, saa
er det netop den største Krænkelse af Velfærdsprincipet, idet
selve Evnen til at føle Værdi ophæves. --- Velfærdsprincipet
forudsætter, at vi formaa at hensætte os paa de fremmede Centralsteder
og at overveje, hvilken Indflydelse de Handlinger,
der skulle bedømmes, her øve. Paa denne Evne beror jo ogsaa
den egentlige Sympatis Mulighed.

Dette Princip, der saaledes følger af Velfærdsprincipet, er
opstillet af \textsc{Immanuel} \textsc{Kant} ved Hjælp af en helt anden Begrundelse,
som det vil have sin Interesse at betragte. Da Kant
betragter den etiske Lov som rent formal og apriorisk, kan han
egentligt ikke af den udlede bestemte Formaal for den etiske
Handlen. Ethvert Formaal, der skulde sættes, vilde føre udover
Lovens rene Form. Og dog indser Kant, at enhver Handling
maa have et Formaal. Den Vanskelighed, som derved fremkommer,
mener han at løse paa følgende Maade. Da den menneskelige
Personlighed har Evne til med Frihed og Inderlighed
at underlægge sig selv en universel Lov, saa staar en saadan
Personlighed som et absolut Formaal og bør aldrig nedværdiges
ved at behandles som blot Middel. Det er det Etiske i dets
Inderlighed og Ophøjethed, der er Menneskets Adelsmærke, som
under ingen Omstændigheder tør berøves ham ved, at man behandler
ham som blot Middel. I Kritik der praktischen Vernunft
(Kehrbachs Udg. p. 106) udtaler Kant sig herom saaledes: „I
hele Verden kan Alt, hvad man vil, og hvorover man har Magt,
ogsaa bruges blot som Middel; kun Mennesket, og med det et%
% --- p. 175 ---
\setcounter{footnote}{0}%
\opage{175}hvert Fornuftvæsen, er Formaal i sig selv. Det er nemlig formedelst
sin Friheds Selvlovgivning Subjekt for den moralske
Lov, der er hellig. Netop paa Grund af denne Frihedens Selvlovgivning
er enhver Villie, endog den egne Villie for det enkelte
Individ selv, begrænset ved den Betingelse, der bestaar i
Overensstemmelse med det fornuftige Væsens Selvlovgivning,
--- nemlig saaledes begrænset, at dette Væsen ikke kan underlægges
nogen Hensigt, som ikke er mulig efter en Lov, der
kunde udspringe af selve dets egen Villie, --- altsaa saaledes,
at det aldrig behandles blot som Middel, men tillige som det,
der seiv er Formaal.“

Det er Kants Fortjeneste at have uddybet Personlighedsbegrebet,
idet han paa én Gang fremhæver det personlige Livs
Inderlighed og Universalitet. Personligheden er en indre Verden,
som følger sin egen Lov og dog kan staa i Sammenhæng med
den store Verden. Men Spørgsmaalet er, om Kant her har anvendt
denne betydningsfulde Tanke paa berettiget Maade. Personlige
Væsener lære vi kun at kende ved Erfaring, og kun
ved Erfaring --- gennem den Vexelvirkning, Samfundslivet fører
med sig --- lære vi at sætte os paa deres Stade og føle, hvorledes
Livet former sig i dem. Men af den rent formale Lov
kan hverken Evnen hertil eller Interessen herfor udledes; der
kan intetsomhelst bestemt Indhold udledes af den. Maaske staar
hint Personlighedsprincip som saameget betydningsfuldere, som
det har udviklet sig paa det virkelige Livs og Historiens Grund,
endnu før vi kunne anvise det dets systematiske Plads i Etiken.
At det direkte følger af Velfærdsprincipet, kan der ikke være
Tvivl om.

Det var fremdeles en Ensidighed hos Kant, at Personlighedsbegrebet
begrænsedes til det strengt Moralske. Paa alle
aandelige Omraader og i alle Livsforhold kan Personligheden
fremtræde, og kan Loven for dens Virken udspringe af den
selv. Overalt --- ikke blot paa det moralske Omraade --- er
Selvvirksomhed Personlighedens Kendetegn; den fører Kampen
for Livet paa sin egen Vis, udfra sit eget Centrum. I Rytmen
af Lyst og Smerte ytrer sig paa alle Livstrin Begunstigelsen
% --- p. 176 ---
\setcounter{footnote}{0}%
\opage{176}eller Hæmningen af denne maaske udefra vakte, men ved en
indre Lov bestemte Virksomhed.

Men skønt Kants Begrundelse ikke kan anerkendes, er det
dog en stor Fortjeneste, han har indlagt sig ved at opstille den
frie Personligheds Princip. Og selvom man indrømmer Gyldigheden
af den Begrundelse, han har forsøgt, er det dog af
ikke ringe Betydning, at den vigtige Sætning kan naas ad forskellige
Veje. Baade efter Kants Begrundelse og efter den, jeg
har forsøgt, staar et personligt Væsen baade som Formaal og
som Middel. Derpaa beror Betydningen af Personlighedens frie
Udvikling.

Friheden er Formaal: ti den uhindrede Udfoldelse af Evner
og Kræfter er et Gode, idet den er forbunden med den inderligste
Tilfredsstillelse, medens Følelsen af Tvang og Hindring
er en Ulystfølelse. Ved Friheden spørge vi: hvorfor ikke? Ved
Tvangen spørge vi: hvorfor? Bevisbyrden paahviler dem, der
ville anvende Tvang. Kun som Middel og Forberedelse til Kræfternes
frie Udfoldelse kunne Tvang og Skranker være berettigede,
saaledes f. Ex. naar det ene Individs grænseløse Udfoldelse
vilde gøre et eller flere andre Individers Udvikling umulig.

Men Friheden er ogsaa Middel: Individernes frie Udfoldelse
frembringer nye Muligheder og fører til Opdagelse af nye
Retninger for hele Slægtens Liv. Det Nye begynder stedse
paa enkelte Punkter og udbreder sig saa til større Kredse. De
personlige Væsener ere ikke blot de Centralsteder, hvor Livets
Værdi bliver følt, men tillige ogsaa de Steder, fra hvilke Livsbevægelserne
stedse paany gaa ud. Ved Friheden skabes der
nye Midtpunkter for selvstændig Virken, og da Slægten bestaar
af personlige Midtpunkter, bliver ogsaa Slægtens Liv derved
rigere og kraftigere. Heri ligger de moderne Emancipationsbevægelsers
store Betydning: at Slaven, Bonden, Arbejderen og
Kvinden komme til at staa som frie og selvstændige Led af
Slægten, kan kun gøre dennes Liv rigere og frugtbarere. Frigørelsen
kræves af Samfundets egen Interesse. Denne Tanke
indskærper \textsc{Stuart} \textsc{Mill} netop i sin Bog Om Friheden (p.
115—119) i mærkelig Modstrid med sin Paastand om den individuelle
Selvudvikling som indifferent i moralsk Henseende.
% --- p. 177 ---
\setcounter{footnote}{0}%
\opage{177}Han fordrer Fredning af de enkelte Personligheders Selvstændighed
og Originalitet, fordi Stødet til alt Klogt og Ædelt kommer
fra enkelte Individer, og fordi Massernes aandelige Herredømme
vilde føre til almindelig Middelmaadighed. Mærkeligt nok
kom samtidigt med det nævnte Skrift en analog Tanke frem i
\textsc{Charles} \textsc{Darwin}'s Lære om det organiske Livs Udvikling gennem
de nye Midler og Veje for Tilværelseskampen, som individuelle
Variationer gøre mulige. Kun hvor saadanne Variationer
fremkomme (af en eller anden Aarsag), kun der kan den naturlige
Udvælgelse, som Livsforholdene øve, komme til at virke.
Baade biologisk og sociologisk viser saaledes den frie Udfoldelse
af de enkelte Individers Natur sig som et betydningsfuldt
Middel for Fremskridtet.

Vi kunne her fra en ny Side belyse Forholdet mellem Motiv
og Handling. Som vi have set (III, 18), gaar Vurderingen
konsekvent tilbage fra Handlingen til Motivet, fordi det maa
være Hovedsagen at ramme den Kilde, fra hvilken Handlingen
udspringer. Vi gik der ud fra Vurderingens og de vurderende
Dommes praktiske Betydning. Her kunne vi supplere denne
Betragtning ved Hjælp af den frie Personligheds Princip. Et
personligt Væsens indre Liv bør nemlig --- i Kraft af dette
Princip --- ikke gøres til blot Middel til Frembringelse af ydre
Virkninger. Det maa ikke være blot Hjul i en Maskine. Selv
hvor Individet betragtes som Led i en større Helhed, skal dets
Handlen udspringe af dets ejendommelige Væsen, af dets indre
Sindelag. Kun da kan Handlingen kaldes fuldkommen, naar
den paa én Gang udspringer af personlig Ejendommelighed og
paa fremmende og lykkebringende Maade griber ind i det store
Hele. De Pligter og Dyder, som Samfundet fordrer af den
Enkelte, skulle tillige være Midler og Former for Individets
egen Udvikling.

Det er den store pædagogiske (eller maaske bedre: psykologiske)
Kunst, som der baade fra enkelte Individers og fra
Samfundsmagtens Side kan blive Lejlighed til at øve: at fremkalde
Motiver, der gøre Selvvirksomhed mulig. Hvis man ved
Tvang vilde forstaa enhver Indgriben fra Andres Side, vilde en
saadan motivvækkende Virksomhed være en Tvang, selvom den
% --- p. 178 ---
\setcounter{footnote}{0}%
\opage{178}øvedes nok saa indirekte og sokratisk. Det er dog naturligere
ved Tvang blot at forstaa en ydre Paavirken, der fremkalder
Smerte eller Frygt (smlgn. V, 2 c). Og den motivvækkende
Indgriben forudsætter ikke nødvendigvis Tvang i denne snevrere
Betydning af Ordet. Et Ansvar følger med al saadan Indgriben;
men naar denne bestaar i egentlig Tvang, behøver den en
ganske særlig Retfærdiggørelse. Der er her Meget at lære af
de teoretiske Anarkister, som forkaste al Autoritet og al aandelig
og fysisk Tvang. Naar man spørger, hvorledes en saadan
Anarkist kan opdrage sine Børn, vil han svare, at han søger
paa indirekte Maade at vække de uvilkaarlige Følelser, som
efter hans Mening udgøre det etiske Grundlag, saaledes at Børnene
uden Bud og uden Tvang komme til at gøre det Rette.
Han vil altsaa forsøge at gennemføre, hvad Rousseau kaldte
„negativ Opdragelse“. Spørgsmaalet (som Anarkisten afgør vel
hurtigt) er, hvorvidt al motivvækkende Indvirken kan foregaa
paa denne Maade, og i hvor høj Grad egentlig Tvang kan indskrænkes.
Men her have vi kun med det almindelige Princip,
det vil sige, med Fordelingen af Bevisbyrden at gøre. ---

Men herved opstaar igen det Spørgsmaal, om der hos hvert
enkelt Individ er den samme Mulighed for Opstaaen af Motiver
til de Handlinger, der maa kræves fra Samfundets Side. Naar
man nu ikke dogmatisk gaar ud fra Antagelsen af, at alle Mennesker
i den nævnte Henseende ere aldeles ens, fremkommer
her, som tidligere vist (se IV, 2), den uundgaaelige Fordring
om Individualisering af den etiske Lov. Den etiske Lov kan da
ikke være fuldstændigt given ved Bestemmelse af „Samfundslivets
Betingelser“, ligegyldigt, hvorledes de enkelte Individer
ere stillede med Hensyn til at fyldestgøre disse Betingelser.
Eller rettere: naar man ikke gaar ud fra et rent udvortes Forhold
mellem Individ og Samfund, kan man ikke fastsætte Samfundslivets
Betingelser uden derunder ogsaa at indbefatte de
enkelte Individers Ejendommeligheder i kvalitativ og kvantitativ
Henseende. Betingelse for et sundt Samfundsliv er det jo dog
først og fremmest, at der er et harmonisk Forhold mellem de
Enkeltes Udvikling og det Heles Fordringer. Og et saadant
Forhold forudsætter igen, at den Opgave, der stilles den En%
% --- p. 179 ---
\setcounter{footnote}{0}%
\opage{179}kelte, svarer til de kvalitative og kvantitative Forudsætninger,
han medbringer. Da Samfundet bestaar af alle de Enkelte, er
dette ikke en udvortes Fordring, der stilles til det. Kun der er
der et virkeligt Samfund tilstede, hvor den Enkelte paa ethvert
Punkt kan være baade Formaal og Middel. En Modstrid mellem
Samfundsfordring og individuel Natur tyder derfor paa en
Ufuldkommenhed hos selve Samfundet, ikke blot --- eller ikke
altid --- hos Individet.

Ikke blot den sædvanlige Moralprædiken overser dette.
Men enhver Opfattelse, der lader det Etiske komme „udefra“
til Mennesket --- hvad enten man bygger paa Autoritetsprincipet,
eller paa en Opfattelse af Slægten som Andet og Mere
end et organisk Hele af Individer, eller paa den offentlige Mening
som Indbegreb af al Etik --- enhver saadan Opfattelse maa
naturligvis forkaste denne Betragtning. Derved afskærer man
sig fra Behandling af nogle af de dybest liggende personlige
Problemer. En illusorisk og rent formel Afslutning af Teorien
købes med Tilsidesættelse af Opgaver, som ganske vist kun tilnærmelsesvis
kunne løses, men som det dog er af stor Betydning
at fastholde. Som ved sit første Grundlag (Vurderingsmotivet),
saaledes staar Etiken ogsaa ved den konsekvente Gennemførelse
af sine Principer overfor det Irrationelle (for at bruge
en fra Aritmetiken hentet Analogi). Der kan her kun være
Tale om at finde saa mange Decimaler som muligt. ---

Kun ved den her gennemførte Betragtning kan der blive
et harmonisk Forhold mellem individuel og social Etik. De
vise gensidigt hen til hinanden, idet Individet er at betragte
som organisk Led af Samfundet, og Samfundet som et organisk
Hele af Individer. Den individuelle Etik fremstille vi først,
fordi den frembyder de simpleste Forhold; men vi fastholde
stadigt det sociale Synspunkt som det til Grund liggende, ligesom
vi i den sociale Etik stedse ville lægge det Synspunkt til
Grund, at Samfundslivets højeste Formaal og bedste Middel er
den frie Udfoldelse af de enkelte Individers ejendommelige
Kræfter.
% --- p. 181 ---
\setcounter{footnote}{0}%

\opage{181}\partitle{INDIVIDUEL ETIK}

% --- p. 183 ---
\setcounter{footnote}{0}%

\opage{183}\chapterhead{IX.}{DEN INDIVIDUELLE ETIKS INDDELING.}

1. Det Forhold, som ovenfor er fastsat mellem individuel
og social Etik, fører naturligt til den Inddeling, som maa følges
indenfor den individuelle Etik.

Idet vi fra første Færd betragte Individet som Led af Samfundet,
og indenfor Samfundet baade som Formaal og som Middel,
viser der sig to Tendenser i det personlige Liv, som maa
anerkendes og fremmes: Selvhævdelsen og Hengivelsen.

Det enkelte Individ er Formaal, da det er selve Samfundets
Repræsentant; i dets Liv rører sig Samfundets Liv, og ved
dets Krænkelse krænkes Samfundets Liv. Det er et af de Endepunkter,
i hvilke Udviklingen løber ud. Jo mere det formaar
at hævde sig og udfolde sine Evner og Drifter, saa at Livet
rører sig i det med fuld Kraft og Harmoni, des højere et Trin
er der naat --- ikke blot for selve den Enkelte, men for Samfundet.
Samfundet bestaar af Individer, og dets Liv maa derfor
være des fyldigere og kraftigere, jo mere hver Enkelt formaar
at udvikle, hvad der ligger i hans Natur. Den Enkelte er en
lille Verden, hvis Bestaaen og Udvikling er et selvgyldigt Formaal.
Selvhævdelsen --- i dens forskellige Former: Selvopholdelse,
Selvbeherskelse og Selvstændighed --- træder derved frem
som en væsentlig Dyd. Den er en Pligt, den Enkelte har paa
én Gang mod sig selv og mod Samfundet. Selvhensynet og
Hensynet til Samfundet gaa her saa umiddelbart i Et, at det
% --- p. 184 ---
\setcounter{footnote}{0}%
\opage{184}kan blive ligegyldigt, hvilket af Synspunkterne man betoner.
Der behøves ingen Appel til bevidst Samfundsinteresse eller til
det abstrakte Velfærdsprincip for at vise, at Selvhævdelsen maa
betragtes som værdifuld ogsaa af Andre end af den, der selv
øver den. En kraftig og harmonisk Udfoldelse af Personligheden
er Genstand for umiddelbar Sympati og Beundring. Som allerede
\textsc{Hume}\footnote[1]{Enquiry concerning the Principles of Morals VI, 1. Kfr. Treatise III, 3, 1.} har bemærket, er det endnu vanskeligere at forklare
Beundringen hos Andre for de Dyder, der komme et Individ
selv til Gode, af egoistisk Interesse, end det er at forklare
Anerkendelsen af sociale Dyder paa denne Maade. Vi betragte
ikke nødvendigvis Andres Selvhævdelse som blot Middel. Det
er en unødvendig Omvej at gaa, naar Hobbes-) forklarer Afskyen
for Drukkenskab derved, at denne let fører til Brud paa
Freden og overhovedet gør, at man ikke kan stole paa Individet.
Ved Synet af manglende Selvbeherskelse hos Andre føle
vi umiddelbart et Skaar i den kraftige og harmoniske Selvudfoldelse,
vi ønske, at Andres Livsførelse skal give os et Billede
af. Vi erkende, at hvor et saadant Billede træder os i Møde,
dér er paa dette Punkt af Verden det Højeste naat. Den stiltiende
Forudsætning ligger naturligvis til Grund, at Kraften og
Harmonien hos den Enkelte ikke strider mod Betingelserne for,
at Andre kunne stræbe efter at naa et lignende Maal for deres
Vedkommende, men maaske endog er et Middel dertil. Der
gives jo, som tidligere udviklet (VIII, 5), neppe Noget indenfor
den individuelle Udvikling, selv i dennes frieste Ejendommelighed,
som ikke kan blive af social Betydning; i al Selvhævdelse
er der en potentiel social Etik tilstede. Men ved Anerkendelsen
af Selvhævdelsen som Formaal peger Etiken ud over sig selv.
Dersom Etiken er til for Livets Skyld, og Livet ikke for Etikens,
vil der være et Grænsepunkt, hvor Normernes Gyldighed
hører op, fordi den umiddelbare Livsudfoldelse ingen Retfærdiggørelse
behøver. Hvor dette Grænsepunkt ligger, og hvor nær
man i det enkelte Tilfælde er det, kan være vanskeligt at 
% CHECK p. 184: note whose mark was not found
\footnote[2]{De cive 111, 25.}af%
% --- p. 185 ---
\setcounter{footnote}{0}%
\opage{185}gøre, nåar Tvivlen rejser sig. Men den umiddelbare Selvhævdelse
har Værdi som Formaal, ikke blot som Middel; derom kan der
efter det hele Synspunkt, som her er anlagt, ikke være Tvivl.

Da Individet dog stedse kun er En blandt Mange, og da
Samfundets Livsinteresser udgøre et større Hele end Indbegrebet
af hans isolerede Livs Interesser, vil Hengivelsen kunne
træde i Modsætning til Selvhævdelsen. Hengivelsen til en større
Kreds af Interesser kan medføre en Udvidelse af Interessen,
der hindrer den individuelle Harmoni, som Selvhævdelsen stræber
efter. Selvhævdelsen har med Individets egen lille Verden
at gøre og betragter den som et afsluttet Hele, men Hengivelsen
fordrer denne lille Verden bragt i Sammenhæng med en
større Verden, og dette kan medføre en foreløbig Forstyrrelse
af den lille Verdens indre Harmoni. Det er jo ikke sagt, at
Ordningen af den lille Verden uden videre passer ind i den
store Verdensorden. Afslutning og Udvidelse kunne staa i Modsætning,
endog i Modstrid med hinanden. Derfor maa Selvhævdelse
og Hengivelse betragtes som to forskellige Tendenser i
Karakteren, to forskellige Dyder, og vi finde da ogsaa, at der
i Etikens Historie er lagt forskellig Vægt paa dem.

Dog er det ikke nødvendigt, at de komme i Strid med hinanden.
Den Trang til Enhed og Kontinuitet, som allerede ytrer
sig i Selvhævdelsen, forsaavidt den enkelte Personlighed betragtes
som et afsluttet Hele, kan føre ud over den individuelle Totalitet
og fremtræde som Trang til Sammenhæng med et mere
omfattende Hele. Den Enkelte kan maaske kun faa Sammenhæng
indenfor sin egen lille Verden ved at knytte sin Interesse
til noget Varigt, noget stedse Fremskridende, noget Uoverskueligt.
Han kan maaske kun være tro mod sig selv ved at være
tro mod Noget, der er større end han selv. Paa denne Maade
vilde Hengivelsen være Fortsættelsen af Selvhævdelsen.

Heraf følger da ogsaa, at det vilde være urigtigt at stille
Hengivelsen som den passive Dyd i Modsætning til Selvhævdelsen
som den aktive. Sympatien, den Følelse, der umiddelbart
fører til Hengivelse, er ikke en Tendens til passiv Hensynken,
der skulde danne Modsætning til den energiske Selvopholdelsestrang.
Man har i den nyeste Tid proklameret den
% --- p. 186 ---
\setcounter{footnote}{0}%
\opage{186}Stærkeres Ret, den egoistiske Hensynsløshed som det Højeste,
som „Overmenneskets“ Karaktermærke, og man har da betragtet
det som en Slaveopstand, at Menneskekærligheden proklameredes
som etisk Princip. I Modsætning til denne af \textsc{Nietzsche}
udtalte Paastand maa det hævdes, at Sympatien eller Menneskekærligheden,
hvor den er ægte og oprindelig, netop er et Udtryk
for Kraft, for aandelig Magt. Den forudsætter, at ikke al
Energi forbruges i de egne, rent individuelle Fornødenheders
Tjeneste, men at der er et Overskud, som gør det muligt at
føle Lyst eller Smerte ved Andres Skæbne, selvom denne ikke
griber ind i den egne, individuelt afsluttede Tilværelse. Sindet
raader da over en større Fylde end ved den isolerede Selvhævdelse.
Paa Grund af denne overskydende Fylde af Kraft og Interesse
har Individet i den ægte Sympati en sand Overlegenhed.
Dets Færd bliver i Hengivelsen bestemt indefra og er
uafhængig af Andres Had eller Kærlighed, af deres Foragt eller
Beundring. Den er --- for at bruge et Billede af \textsc{Marcus}
Aurelius --- som den rene og stærke Kilde, der lader sin
Strøm vælde frem, selvom man kaster Smuds og Sten i den;
den skyller det bort, og dens Renhed ændres ikke derved.

Allerede Urkristendommen beskrev Kærligheden som en Magt,
der er uafhængig af, om Andre elske os eller ikke. Den Kærlighed,
der kun er „reaktiv“, kun fremkaldes ved Andres Tjenester,
erklærer den for utilstrækkelig. Den sande Kærlighed
gør ingen Forskelle, omfatter Alle --- ligesom Himlens Sol og
Regn blive Alle til Del. Kærligheden bærer Alt, taaler Alt,
overvinder Alt! Kærlighedens Langmodighed lader den bestaa
trods den største Modstand. (Matthæusevangeliet V. --- 1. Korintherbrev
XIII. --- De imitatione Christi III, 5).

I den nyere Tid findes denne Opfattelse af Sympatien som
Udtryk for Kraft hos Spinoza og Rousseau. For Spinoza fremtræder
den aandelige Kraft, der udgør Dydens Væsen, dels som
Selvhævdelse eller Livsmod (animositas), dels som Højmod (generositas),
og den Højsindede søger at hjælpe Andre og forbinde
dem med sig i Venskab, ti „Sjælene besejres ikke ved Vaaben,
men ved Kærlighed og Højsind“. Rousseau forklarer Kærligheden
til Andre af den overstrømmende Kraft til Selvhævdelse,
% --- p. 187 ---
\setcounter{footnote}{0}%
\opage{187}der uvilkaarligt breder sig og omfatter Andre, naar de kun ligne
os, fordi den er for stor til eget Behov. Forskellen mellem os
og Andre agtes ikke; Livsfylden gaar ud over alle Skranker.
„La force d'une åme expansive midentifie avec mon semblable!“
Kærligheden er, som Rousseau udtrykker sig, en Følge af Kærligheden
til sig selv (amour de soi), hvilken han skelner fra
Egenkærligheden (amour propre). Egenkærligheden sætter Skranker,
idet Mennesket sammenligner sig med Andre og skiller sig
fra dem. Men derved blive vi netop afhængige af Andre, hvad
vi ikke ere i Kærligheden. Senere er en lignende Opfattelse
gjort gældende af \textsc{Jouffroy} og andre franske Filosofer, især
Guyau, til hvem \textsc{Krapotkin} slutter sig\footnote[1]{Kfr. min Psykologi. VI C, 3. 7. --- Den nyere Filosofis Historie'\textsuperscript{2}. I p. 326; 494 f. --- Jean Jacques Rousseau og hans Filosofi'\textsuperscript{1}, p. 101---105. Om Jouffroy se Bouillier: Du plaisir et de la douleur. Paris 1877. p. 169 f. --- Guyau: Esquisse d'une Morale sans obligation ni sanction. Paris 1885. --- Krapotkin: Anarkistisk Moral. (Dansk Overs. i ..Proletaren“ 279—284. 1896). --- Smlgn. ogsaa William J ames: Varieties of belief. p.}.

Der gives naturligvis en blødagtig og passiv Art af Følelse,
der ogsaa gaar under Sympatiens og Kærlighedens Navn. Og
der gives en Sentimentalitet, der nyder den egne, formentligt
store Sympati, ligesom en anden Art Sentimentalitet nyder de
egne, formentligt store Smerter. Men disse Former tør ikke
betragtes som typiske, og den sidstnævnte er endog snarest en
Art af egoistisk Selvnyden. ---

En karakteristisk Forskel mellem Selvhævdelse og Hengivelse
bestaar i, at i hin staar Individet som Formaal og virker
dog som Middel, idet Samfundet styrkes og udvikles ved
Enernes Selvhævdelse, --- i denne virker Individet som Middel
og bliver dog Formaal, idet dets personlige Liv bliver rigere
og fyldigere ved at faa et større Omraade at virke paa. Denne
Forskel godtgør tillige den indre Sammenhæng imellem dem.
Naar vi opfatte Retfærdighed paa den Maade, som tidligere
(III, 9) er angiven, kunne vi i den finde den harmoniske Enhed
af Selvhævdelse og Hengivelse, af Tendensen til at afslutte
sig i sig selv og til at aabne sig for en større Sammenhæng.
% --- p. 188 ---
\setcounter{footnote}{0}%
\opage{188}Hvad den psykologiske Oprindelse angaar, kan Retfærdigheden
have sin Rod i Selvhævdelsen, nemlig naar denne forbindes
med Anerkendelsen af Andres lige Berettigelse, en Anerkendelse,
der maaske først er aftvungen ved Magt og Autoritet,
men som ved Motivforskydning kan gaa over i Individets Kød
og Blod. Men den kan ogsaa have sin Rod i Hengivelsen,
nemlig naar denne forbindes med Forstaaelse af de enkelte
Personligheders (ogsaa den egne Personligheds) Ejendommelighed
og Betydning. Selvhævdelse og Hengivelse ere hver for
sig mere elementære Tendenser. Retfærdigheden er den mere
omfattende Tendens, der sætter Kronen paa den etiske Karakters
Udvikling. Retfærdigheden forudsætter, at Individet hævder
og udvikler sig selv, ikke blot uden at hæmme Andres
Udvikling, men saaledes, at denne derved fremmes, --- og at
Individet hengiver sig til omfattende Livsinteresser, saaledes at
dets egen Personlighed derved udvikles og hævdes\footnote[1]{Allerede i den græske Etik fremtræder Retfærdigheden i nøje Sammenhæng med Menneskekærligheden. Smlgn. Aristoteles: Ethica Nie. VIII, 1. 13. 14. Cicero: De finibus. V, 23, 65.,}.

Ikke blot Selvhævdelse og Hengivelse ere her i Harmoni,
men ogsaa Erkendelse og Følelse. Et stort Formaal staar for
Tanken, og denne søger ad Erfaringens og Undersøgelsens Vej
Midler til at virkeliggøre det. Det er et stort Fordelingsarbejde,
der skal gøres i Slægtens Tjeneste, og Fordelingen bestemmes
ved klar Forstaaelse af de enkelte Personligheders Trang og
Tarv. ---

Den individuelle Etik inddeles ifølge hele denne Betragtning
naturligt i Læren om Selvhævdelsen og Læren om Hengivelsen.
Naar disse to Tendenser skildres saaledes, at en harmonisk
Forbindelse mellem dem bliver mulig, vil dermed ogsaa
den etiske Hoveddyd, Retfærdigheden, være skildret.

2. Hos de forskellige Individer spille disse Tendenser en
forskellig Rolle. Der gives Naturer, som yde det Højeste, de
formaa, uden at den etiske Følelse rører sig hos dem som en
særegen Følelse ved Siden af de andre. Selvhævdelsen og Hengivelsen
udfolde sig hos dem uden særlige Motiver og uden
% --- p. 189 ---
\setcounter{footnote}{0}%
\opage{189}bevidst Anstrengelse. Muligheden af saadanne etiske Aladdinsnaturer
kan ikke benægtes. Der kan være dem, hos hvem en
saadan Aladdinsnatur kun er tilstede, hvad nogle, ikke hvad
alle Opgaver og Forhold angaar. Ligesaa let og ubevidst de
løse nogle af de etiske Opgaver, ligesaa megen Besindelse og
Anstrengelse maa de anvende paa andre. Men der gives ogsaa
Naturer, for hvem det er af den allerstørste Betydning, at den
etiske Følelse vækkes hos dem og kommer til at tale et afgørende
Ord med i de enkelte Tilfælde. Den umiddelbare
Trang til Selvhævdelse og til Hengivelse rører sig enten ikke
stærkt nok hos dem, eller ogsaa er der ikke det rette Forhold
mellem de to Tendenser. --- Det er af etisk-social Betydning,
at der findes saadanne Forskelligheder. Derved muliggøres et
større og mere righoldigt Arbejde, end hvis Alle mødte med
de samme Betingelser. Baade selve Mangfoldigheden og de
Modsætningsforhold, den gør mulige, virke udløsende og udviklende.

I den systematiske Fremstilling ville naturligvis alle disse
individuelle Forskelligheder ikke kunne udtømmes. Vi maa indskrænke
os til at give en ordnet og begrundet Fremstilling af
de vigtigste af de Karakteregenskaber, Velfærdsprincipet fordrer.

\chapterhead{X.}{DET ETISKE LIVS PERSONLIGE GRUNDLAG.}

1. Man har af forskellige Grunde benægtet, at Forberedelse
og Øvelse have nogensomhelst Værdi i etisk Henseende.

Især har Reaktionen mod den asketiske Tankegang været
medvirkende ved denne Paastand. Askesen --- hvorved her forstaas
frit valgt Anstrengelse og Lidelse, som blot tjener til Øvelse,
ikke til Opnaaelse af paaviselige Formaal --- kunde, mener man,
være konsekvent paa Nyplatonismens eller den middelalderlige
% --- p. 190 ---
\setcounter{footnote}{0}%
\opage{190}Kristendoms Standpunkt. For Nyplatonikerne gjaldt det om at
befri Sjælen for den Urenhed, den havde paadraget sig ved at
indtræde i den jordiske Verden. Fra det strenge kristelige Standpunkt
var Afholdenhed, Spægelse og Selvydmygelse Sonofre til
en vred Gud. Men hvor saadanne Anskuelser ikke længere
herske, kan man dér tillægge Askesen nogen Berettigelse og
Betydning?

Med Overgangen til en ny Livsanskuelse forbinder sig let
en Tilbøjelighed til, for en overfladisk „Konsekvens“'s Skyld,
at kaste Alt bort, hvad man paa det tidligere Standpunkt tillagde
Værdi. Hvad der historisk har udviklet sig under Indflydelse
af visse Motiver, kan meget godt beholde sin Værdi, fordi disse
Motiver falde bort og maa erstattes af andre. (Smlgn. I, 4). Ethvert
etisk Standpunkt maa stille Fordringer om at taale Anstrengelse
og Lidelse, og Øvelse vil derfor stedse være nødvendig.
Det er vel en Grundsætning af stor Betydning, at al Tilføjelse
af Smerte skal retfærdiggøres, men Frembringelsen af
Lyst har sin foreløbige Retfærdiggørelse i sig selv. Men hvor
der stræbes hen imod et fjernt og omfattende Formaal, vil megen
Forberedelse og mange Afsavn kunne blive nødvendige. Derfor
vil enhver alvorlig Livsanskuelse kunne lære af Askesens Heroer.
Ogsaa disse betragtede jo egentligt kun Lidelsen og Savnet
som et Gode, fordi de derved styrkedes og øvedes til at
opnaa det store og fjerne Maal, de higede imod. Deres Fejl
var den, at de satte Maalet i en anden Tilværelse, hvis Vilkaar
skulde være ganske andre end det nærværende Livs, saa at
Maalet ikke kunde naas gennem positivt Arbejde i den givne
Verden. Deres Askese hvilede paa en Mistillid til det naturlige
Liv og dets Kræfter. Men netop paa et Standpunkt, hvor man
ikke faar Reglerne for sin Færd gennem overnaturlig Aabenbaring
eller mystiske Indskydelser, men hvor der lægges Vægt
paa den menneskelige Tankes og Følelses selvstændige Myndighed,
netop dér stilles der store Fordringer til Individets Kræfter,
og det bliver af afgørende Betydning, at disse Kræfter
øves, for at Grundlaget for det etiske Liv kan blive fast og
sikkert. Der gives desuden Naturer, som i etisk Henseende
kun kende et Enten —Eller, enten streng Tugt eller fuldstændig
% --- p. 191 ---
\setcounter{footnote}{0}%
\opage{191}Slaphed. De kunne kun holde sig oppe ved at underkaste sig
en gennemført Disciplin. Afholdsbevægelsens store Udbredelse
viser, at der gives ikke faa saadanne Naturer. Ofte begaa disse
den Fejl at tro, at det, der er nødvendigt for dem, maa være
en Lov for Alle.

Man har indvendt, at det er et Tegn paa Ufuldkommenhed,
naar Handlinger foretages blot til Øvelse, da Middel og
Maal saa falde udenfor hinanden. Sligt kan vel, siger man,
være nødvendigt, naar man vil erhverve mekaniske Færdigheder;
men det etiske Liv skal være den hele Personligheds
Liv. Sligt kan ogsaa passe paa den mystiske og supranaturalistiske
Askeses Standpunkt, hvor man skyede det virkelige Liv
og i Ensomhed øvede Sjælens Kræfter, saaledes at man maaske
endog i Opgavernes Betydningsløshed fandt noget særligt Fortjenstligt,
som naar ægyptiske Munke vandede Stokke, der vare
plantede i Sand.

Denne Indvending gaar med Rette ud fra, at Livet selv er
den bedste Skole af alle, selvom det er den besværligste og
kostbareste. Det er utilstedeligt at unddrage sig dets Krav for
i Ensomhed at arbejde paa sig selv. Man kan ikke træde ud
af Nummer for at blive rigtigt Menneske, og saa begynde igen.
Enhver sand Opdragelse finder Sted gennem Vexelvirkning med
de virkelige Forhold. Hvad Selvopdragelsens Kunst her kan
gøre, er at vælge netop de Opgaver og de Forhold, som egne
sig til at øve og styrke Kræfterne, før der skrides til større Arbejder.
Man bør altid øve sine Kræfter paa Noget, som har
virkelig Betydning. Det bestemte Kald, som vælges, vil være
den naturlige Skole for Karakteren. Først er maaske Efterligning
og mekanisk Udøvelse af foreskrevne Funktioner det
Højeste, man kan drive det til; men der kan ved Gentagelse
og Øvelse uddanne sig Anlæg og Dispositioner, der gøre det
muligt at udføre Kaldsgerninger med Frihed og støttet til egen
Kraft. Da har der dannet sig en til Kaldet svarende Dyd. Men
det er ikke til alle Tider, at de virkelige Opgaver lægge Beslag
paa al vor Kraft. Tiden kan da anvendes til at forberede
sig overfor Mulighederne. Øvelsen bliver da som en Leg eller
en Kunst. I Kunsten have vi ikke med Livet selv, men med
% --- p. 192 ---
\setcounter{footnote}{0}%
\opage{192}en Afspejling af det at gøre, og dog have Kunsten og Livet
det samme Indhold; Brugen af Kræfterne i Kunsten kan derfor
være en Forberedelse til Brugen af dem i Livet. Desuden
er Mulighedernes Omraade større end Virkelighedens. I ydre
Handling kan ofte kun en ringe Del af det personlige Liv faa
Lejlighed til Udvikling og Dannelse. Paa det indre Livs Skueplads,
under Stemningers og Tankers Brydning og gennem spirende
Motivers Vexelvirkning foregaar ofte en væsentlig Del af
Arbejdet paa Karakterens Udvikling. \textsc{Schleiermacher} har i
sine „Monologer“ (Kap. 4) særligt udtalt dette. „Fra hvor
mange Sider“, spørger han, „vilde Mennesket ikke vedblive at
være ubestemt og udannet, naar hans indre Handlen kun gik
ud paa det, som virkeligt tilstøder ham udefra?\dots{} Jeg véd,
at mit ydre Liv aldrig vil kunne fremstille og fuldende mit
indre Væsen fra alle Sider\dots{} Jeg lever i stille Skjulthed ---
og dog paa Verdens store, daadrige Skueplads!“ Senere har S.
\textsc{Kierkegaard} udtalt en lignende Tanke: „Først den, der dannes
ved Muligheden, dannes efter sin Uendelighed.“ Han føjer
til, at de virkelige Situationer ikke ere saa klare og ubetingede,
som Mulighederne kunne tænkes, og at Prøven gennem Muligheder
derfor kan være strengere; men han har ogsaa Øje for,
at det er saa let at bedrage sig selv, naar man blot prøves ved
Muligheder. At Mulighederne ogsaa kunne suge Kraften af Villien,
er en Side af Sagen, han ikke dvæler ved\footnote[1]{Søren Kierkegaard som Filosof. p. 50 —51.}.

2. Alle aandelige Evners Udvikling kan tjene Befæstelsen
og Udviklingen af den etiske Følelse. Samvittighedens Kraft,
Omfang og Renhed bero ikke blot paa en Opdragelse af Følelseslivet,
men ogsaa paa Erkendelsens og Villiens Udvikling. Det
gælder jo paa den frie Samvittigheds Standpunkt ikke blot om
blind Lydighed, om at Buddet kan høres i Bevidstheden og
Handlingen mekanisk følge efter. Den frie Samvittighed søger
selv at stille sine Opgaver og indse deres Berettigelse. Det er
ikke blot Filosofen, der vender sig prøvende mod de overleverede
og instinktmæssigt billigede Regler; der kan for Enhver blive
stillet den Opgave at løsrive sig fra det hidtil Gældende og
% --- p. 193 ---
\setcounter{footnote}{0}%
\opage{193}prøve ubanede Veje. Erkendelsens Klarhed og Konsekvens
lader den Enkelte faa Øje paa de Tilfælde, i hvilke han af Svaghed
eller Egoisme unddrager sig sin egen anerkendte Lovs Gennemførelse,
og den lærer ham tillige at finde Huller og umotiverede
Skranker i den gængse Moral. Hvad han saaledes i de
rolige og lyse Tider har erkendt, gælder det derpaa at holde
fast i de mørke og bevægede Øjeblikke; kun gennem en Villiesanspændelse
vil Troskaben mod sig selv og mod den Sandhed,
man har erkendt, kunne bevares.

Kraft og Omfang faar Samvittigheden kun, hvor Tanke og
Fantasi ere levende og virksomme. Tænkningen maa følge Handlingerne
i deres Virkninger, som ofte forgrene sig ind i mange
Forhold, og Fantasien maa give anskuelige Billeder af disse Virkninger.
Sneverhed og Fordomme komme ofte ikke af en Mangel
i Følelseslivet, men af Mangel paa Fantasi eller paa logisk Konsekvens.
Den blotte Erkendelse er naturligvis ikke nok, den maa indøves
og gøres til en stadig Tanke, det vil sige til en Tanke, som i
Kraft af Forestillingsforbindelsens Love let og hurtigt gør sig gældende
paany, hvor der er Brug for den. Derved bliver den en Del
af vort reale Selv, af den Kreds af Tanker og Følelser, der
danner Tyngdepunktet i vor aandelige Natur, og til hvilken vi
stedse vende tilbage, saa meget end øjeblikkelige Indflydelser
og Tilskyndelser drage os bort derfra. Hverken Forestillingernes
eller Følelsernes Udvikling foregaar aldeles uden vor Selvvirksomhed,
og naar først Opmærksomheden er vakt, vil der
frembyde sig mange Lejligheder til at gribe regulerende ind.
--- En Ting skal blot særligt fremhæves. Det er ikke blot den
teologiske Etik, som har sine Opbyggelsesbøger. Der kan være
Bøger, som man stedse opsøger paany, for at styrke sine Tanker
og sin Følelse. Avislæsning og Underholdningslekture bør
ikke være al den aandelige Føde, man faar, hvis man ikke vil
svække sit indre Liv. Hvad Enhver her nærmest vil vælge,
beror paa hans Karakter og Aandsretning; men Ingen kan undvære
den Samling og Koncentration, som alvorlig Læsning medfører.
Vi føres derved ind i et Ideernes og Mulighedernes Rige,
fra hvilket vi kunne gaa med udvidet og styrket Følelse til de
praktiske Opgaver.

% --- p. 194 ---
\setcounter{footnote}{0}%
\opage{194}Hvad Samvittighedens Renhed angaar, er Selvprøvelse af
største Vigtighed, for at vi kunne blive klare over vore egne
Motiver og lære vore Kræfter at kende. Selverkendelsen har
stedse sine Grænser, idet Meget af det, der bevæger sig i vort
Indre, kun saare utydeligt træder frem for den klare Bevidsthed.
Der kan derhos ligge Muligheder og Anlæg i vor Individualitet,
som endnu ikke hverken indadtil eller udadtil have
ladet sig mærke. Vi kunne derfor aldrig være sikre paa, at
vor Bevidstheds Midtpunkt ogsaa er vor Individualitets Midtpunkt\footnotemark[1].
Men hvad de bevidste Motiver angaar, --- det, som vi
virkeligt have tænkt, følt og besluttet, --- maa vi kunne komme
paa det Rene. Der er vel hos de Fleste en Tilbøjelighed til
Selvskuffelse, til at ville tage sig bedre ud i egne Øjne, end
man i Virkeligheden er, til at besmykke sine Motivers Renhed
og sine Beslutningers Alvor for sig selv. Men det er af afgørende
Betydning, at man er ærlig overfor sig selv og lærer
at sønderrive de Illusioner, i hvilke man saa let spinder sig
selv ind. Denne indre Sandhed er Betingelsen for enhver kraftig
og sund Virksomhed. Paa en Løgn vil ingen holdbar Bygning
kunne opføres. --- Et Selvbedrag indtræder især let paa
saadanne Omraader, hvor ingen exakte Regler kunne opstilles,
og etiske Problemer blive meget hurtigt saa indviklede, at intet
almindeligt Princip kan løse Knuden for os. Selvom vi ikke
ere i Tvivl angaaende Kvaliteten af det Rette, saa kan der rejse
sig uløselige Tvivl om, hvilken Grad og hvilket Omfang vor
Aktivitet skal have. Havde endda de Ret, der mene, at man
har gjort sin Pligt, naar man har gjort, hvad „Samfundet“ kræver,
og at det ikke er Pligt at indlægge sig „Fortjeneste“! Men
der er dem, hvis indre Tilskuer ikke saa let lader sig tilfredsstille,
og hos hvem Frygten for Selvbedrag ved Fastsættelsen
af „Mandevilliens qvantum satis“ derfor stedse vil røre sig
paany.

Det moderne civiliserede Liv har en udad vendt Tendens,
fører let til en usund Selvforglemmelse. Optagen som man let
bliver af den skiftende Mangfoldighed af ydre Begivenheder og
% CHECK p. 194: note mark without a note

% --- p. 195 ---
\setcounter{footnote}{0}%
\opage{195}Forhold, kommer man til at betragte det, der foregaar i Ens
eget Indre, som ligegyldigt. Hint Ydre betragter man som det
egentlige Virkelige, de indre Oplevelser som et flygtigt Genskin
deraf. Men det egentlige Liv er dog stedse det indre Liv; det
er det, der giver det Ydre dets Betydning. Fra det kan der
desuden udgaa Kræfter, som kunne gribe ændrende ind i den
ydre Verden; men saadanne Kræfter fostres ikke, naar man
lader sig drive med Strømmen. Ethvert personligt Indre er et
af de aandelige Arnesteder i Verden, og det gælder at bevare
Ilden ren og kraftig. Der kan gives en Art aandeligt Selvmord,
som netop bestaar i, at man af Sløvhed lader Ilden paa det
indre Arnested gaa ud. Idealet slippes, intet Vurderingsmotiv
rører sig, og heller ikke Instinktet holder tilsidst oppe. Der
høres ved Selvmord af denne Art ikke engang et Plask som af
den, der drukner sig. \textsc{Sibbern} har med Rette beskrevet Samvittigheden
som en aandelig Selvopholdelsestrang, Udtryk for
en Art Selvomsorg. Selvomsorgen viser man her ved ikke at
slaa af paa de Fordringer, man stiller til sig selv. Selvets Niveau
maa ikke sænkes.

Selverkendelsen kan undertiden i høj Grad vanskeliggøres
ved det Sammensatte i den menneskelige Karakter og det Sporadiske
i den menneskelige Udvikling. Forskellige, endog indbyrdes
stridende Motiver og Tendenser kunne trænge sig frem;
snart raader den ene Tendens, snart den anden, --- og i hvilken
af dem skal da Mennesket genkende sit egentlige Jeg?
Hvad der rører sig hos os i de beg«jstrede, de „bedste“ Øjeblikke,
kan være højst forskelligt fra, hvad der kommer frem i
de onde Tider, og begge Dele høre dog med til vort Jeg. Eller
vi dømme og handle rigtigt og sikkert i enkelte Forhold, medens
vi i andre Forhold blindt lade os lede af selviske Lidenskaber.
Her ligger Beskyldningen for Hykleri saa nær, at den populære
Moral ikke bryder sig om, vel heller ikke forstaar at anstille
nogen indgaaende psykologisk Undersøgelse --- og endnu
mindre forstaar at holde sin Dom tilbage. I saadanne problematiske
Tilstande ytrer sig ofte en Gæring, af hvilken en klar og
harmonisk Karakter senere kan udvikle sig. Men der gives
ogsaa Naturer, som ere og blive problematiske, eller hos hvem
% --- p. 196 ---
\setcounter{footnote}{0}%
\opage{196}Striden mellem Sjælens Elementer tilsidst sprænger Karakterens
Enhed\footnote[1]{Smlgn. Jean Jacques Rousseau og hans Filosofi. p. 11. --- Joseph Butler har en Række interessante Bemærkninger om Selvbedrag (Sermon X).}-

3. Men ogsaa om Askese i snevrere Betydning, om Hærdelse
og Øvelse af Evnen til at udholde Anstrengelse og Lidelse
kan der blive Tale. Hos vilde og krigerske Folk bestaar en
stor Del af Opdragelsen, ogsaa Selvopdragelsen, i Udvikling af
Haardførhed, Evne til at døje Kulde, Sult og fysisk Smerte. I
det civiliserede Liv er en saadan fysisk Hærdelse vel mindre
nødvendig, men mange af det civiliserede Livs Misligheder stamme
dog sikkert fra, at man har tabt den for meget af Syne. Den
fysiske Hygiejne er for en stor Del tillige ogsaa en Skole for
Villien. Den Selvovervindelse og Taalmodighed, som maa anvendes
til Bedste for den fysiske Sundhed, kan komme den
aandelige til Gode.

Der gives dog ogsaa en direkte aandelig Hærdelse. Der
kan gives indre Kulde, Mørke og Afmagt ligesaavel som ydre.
Endog uden bestemt ydre Aarsag kan der komme tunge og
mørke Tider, hvor Intet vil lykkes, hvor neppe en Tankegang
kan holdes fast, og hvor Alt taber sin Glans og sin Friskhed
for os. Selv det Tilhold, som man ellers har i Samvittighedens
indre Sanktion, kan svigte. Man griber da let til ydre Midler
for at holde sig oppe, trøste eller stimulere sig 2), og dermed
er da ofte det første Skridt til Forfald gjort. Hvor det er rent
fysiske Depressionstilstande, der indtræde, kan det naturligvis
være fuldstændigt rigtigt at benytte ydre Midler. Men det kan
ofte være det eneste Rette at gøre Modstand, at sætte sin Villie
ind paa at holde sig oppe, for ikke at komme til at lide af
aandelig Blødagtighed. Det kan, som en gammel Mystiker siger,
være nødvendigt „at undvære al Trøst“. Villien øves derved i
at staa paa sine egne Ben.

Til de mørke og ængstende Tanker, som i saadanne Tider
% CHECK p. 196: note whose mark was not found
\footnote[2]{„Naar et Menneske begynder at blive lunken, da frygter han en ringe Anstrengelse og modtager gerne den ydre Trøst“. De imitatione Christi. 11, 4, 3. L}
% --- p. 197 ---
\setcounter{footnote}{0}%
% CHECK p. 197: notes paired with the marks in order (the layer misread a number)
\opage{197}kunne hjemsøge os, regner man i nyere Tid ogsaa gerne Tanken
om Døden. I Oldtiden var dette ikke saa almindeligt. Det
er først blevet almindeligt, efterat den kristelige Teologi særligt
har søgt at skærpe Dødsfrygten, „ligesom de Mødre, der styre
deres Børn ved at indbilde dem, at Mørket er fuldt af Spøgelser,
der gribe de Ulydige“\footnote[1]{Lecky: History of European Morals from Augustus to Charlemagne. I, p. 221 f.}. I Modsætning til den fremherskende
Plads og Betydning, som den teologiske Etik gav
Tanken om Døden, har \textsc{Spinoza} udtalt, at det frie Menneske, det
vil sige det Menneske, som ledes af den rette Erkendelse, ikke tænker
mindre paa Noget end paa Døden; dets Gransken er en
meditatio vitæ, har Livet til Genstand, ikke Døden\footnote[2]{Ethica IV, 67.}. Det er
ikke Tankeløshed, som her anbefales. Tankeløshed kan ofte
dække over en lønlig Frygt. Ingen alvorlig Livsanskuelse kan
undlade at kæmpe med Dødstanken. Men Forskellen mellem
den teologiske og den filosofiske Etik lægger sig her for Dagen
paa karakteristisk Maade. For den teologiske Etik er hele Livet
en Forberedelse til et andet Liv. Vi leve foran et Forhæng,
bag hvilket det sande Liv er skjult, og Alt, hvad vi gøre og
lade, skal tilsidst tjene til, at vi kunne være beredte, naar Forhænget
endeligt drages til Side. Tanken om Døden bliver derfor
her en Alt afgørende, en ledende Tanke. Den filosofiske
Etik bygger ikke paa Antagelser om, hvad der ligger udover
Erfaringen, men hævder, at Livet først og fremmest maa have
sit Formaal, sin Værdi i sig selv. Vi tillægge Livet Værdi,
ikke fordi det er Forberedelse til et andet Liv, som vi ikke
kende, men fordi det i sig selv indeholder Noget, som er skønt
og godt, Noget, som fortjener, at man lader sig opfylde af det
og slaar et Slag for det. Vi leve paa Virkeligheder, ikke paa
Muligheder. Døden er Livets Grænse, og den er bedst forberedt
til at nærme sig denne Grænse, som virkeligt har levet, det vil
sige, som har haft Del i det Bedste i det Liv, vi kende.

Om det Mulighedernes Rige, som begynder, hvor vor Erfaringsverden
hører op, ville Anskuelserne sandsynligvis stedse
% --- p. 198 ---
\setcounter{footnote}{0}%
\opage{198}være forskellige. En videnskabelig Afgørelse kan ikke tænkes.
Vi staa her overfor et af de aabne Spørgsmaal, og Enhver kan
her for egen Regning have sin Tro eller sit Haab. Hvad Filosofien
her gør os opmærksom paa, er blot den Kendsgerning,
at alle de Træk og Egenskaber, hvormed de Billeder, vi kunne
danne os af en anden Verden, ere udstyrede, ere hentede fra
Erfaringsverdenen, kun at de forestilles i langt højere Grader
og frigjorte fra alle Ufuldkommenheder. Hvis man ikke fra Erfaringen
kendte Noget, som var sandt, godt og skønt, vilde man
ikke kunne danne Forestillingen om en absolut Sandhed og Godhed,
en fuldstændig Salighed. Men saa er jo det Sande, Gode
og Skønne i Erfaringsverdenen en fælles Grund, hvor vi Alle
kunne mødes, hvad enten vi antage en anden Verden eller ikke.
Det er en Kendsgerning, at der kan leves et Liv, som er viet
til Arbejde i det Sandes og det Godes Tjeneste, uden at der
føles nogen Trang til at tro paa en anden Tilværelse. Og der
kan ikke nævnes nogen virkeligt etisk Egenskab, som kun skulde
blive mulig ved en saadan Tro. Selvom Troen paa en anden
Tilværelse er en personlig Fornødenhed for den Enkelte, saa
har han dog ikke Ret til at gøre den til en etisk Nødvendighed
for Alle.

Modstykket til Kampen mod de mørke Muligheder er Tanken
om de lysende Forbilleder. Saadanne Personligheder, hos
hvem det aandelige Liv rører sig med Kraft og Harmoni, idet
Selvhævdelsen og Hengivelsen virke sammen paa ethvert enkelt
Punkt, staa som Forbilleder for dem, der maa kæmpe møjsommeligt
for at hæve sig over, hvad der drager ned, eller for
at holde sammen paa, hvad der truer med at opløse sig i Strid. Et
tilsyneladende Paradox kan her fremkomme ved, at de Personligheder,
der staa som Forbilleder, ikke altid selv have vundet
Fuldkommenheden gennem Møje og Kamp. Energien og Harmonien
kunne have været Naturgaver, og dog kunne de tjene
som Tegn, der vise Retningen og styrke Haabet. De vise os,
hvilke Muligheder der ligge i den menneskelige Natur. Havde
hine Personligheder haft fuldt op at kæmpe med i deres Indre,
vilde deres Liv maaske ikke have frembudt os det Billede af
Kraft og Harmoni, som nu styrker og leder os. I den græske
% --- p. 199 ---
\setcounter{footnote}{0}%
\opage{199}Oldtid stod \textsc{Sokrates} med sin vidunderlige Harmoni og energiske
Aandsklarhed som det store Forbillede. Kræfter, der ellers
saa let træde i større eller mindre Modsætning til hinanden, ---
Tanken og Villien, --- vare hos ham, saavidt man kan se, forenede
i en oprindelig Enhed. Det var denne oprindelige Karakterstyrke
hos Sokrates, der gjorde den sokratiske Filosofi saa
ensidigt intellektualistisk og laa til Grund for den Lære, at Dyd
er Viden. Han kendte ikke af egen Erfaring Forskel mellem
Viden og Villen\footnote[1]{Karl Joël: Der echte und xenophontische Sokrates, Berlin 1893. I. p. 256. --- Det bekendte Sted i „Fajdros“,\textbackslash{} hvor Platon lader Sokrates sige, at han ikke véd, om han er et enkelt og mildt eller et mangfoldigt og flammende Væsen, maa skrives paa Platons Regning; først han, ikke Sokrates, skelnede mellem forskellige Dele i Sjælen. Smlgn. Aristoteles: Magna Moralia. c. 1. Joël anf. Skr. p. 277.}. Men hvad der gjorde hans Filosofi saa gaadefuld
ved sin Ensidighed, har netop gjort ham til et stort Forbillede.
Senere hen, da den græske Filosofi i \textsc{Nyplatonismen}
fik en mystisk-religiøs Karakter, og Dyden opfattedes som en
Stræben efter at blive Gud lig, indskærpedes med særligt Eftertryk
den Sætning, at hvad der bliver Dyd i Genbilledet, ikke
er Dyd i Forbilledet; ti som absolut Væsen kan Gud ingen Dyd
have. Og i den kristne Kirke er Kristus bleven opstillet som
Forbilledet, uagtet han efter Kirkens Lære var den efter sin
Natur Syndfrie. Der er den Sandhed heri, at Udviklingen paa
det etiske Omraade ligesom paa det organiske Omraade muliggøres
ved, at uvilkaarlige Variationer frembyde nye Muligheder
og føre ind paa nye Retninger. Allerede paa lave menneskelige
Udviklingstrin ville ikke blot de fysisk Overlegne, men
ogsaa de aandeligt Overlegne kunne komme til at staa som Forbilleder
til Efterlignelse. Der har hos dem i lykkelig og naturlig
Form aabenbaret sig et Karaktertræk, som føles i sin Værdi,
og som man tror vil berige Livet, om det kan forplantes videre.
Etiske Aladdiner have været Lysene paa den etiske Udviklings
ofte mørke Vej. Det Uvilkaarlige gaar stedse forud for det Vilkaarlige,
og de Egenskaber, der bleve ophøjede til Dyder, maatte
% CHECK p. 199: note whose mark was not found
\footnote[2]{Plotinos: Ennead. I, 2, 2. Smlgn. allerede Platon: Theaitetos p. 176 B.}
% --- p. 200 ---
\setcounter{footnote}{0}%
\opage{200}først uvilkaarligt aabenbare sig, før den etiske Udvælgelse kunde
finde Sted. Efterlignelsestrang, Beundring og Ærefrygt have øvet
denne Udvælgelse, sammen med hvilken den strenge Tugt fra herskende
Autoriteters og fra Samfundsmeningens Side stadigt har
virket; dog har en Vexelvirkning fundet Sted, da heller ikke Magthaverne
og Folkemeningen kunde unddrage sig Forbilledernes
Indflydelse. \textsc{Fichte}, \textsc{Goethe} og \textsc{Carlyle} have henvist til de
personlige Forbilleders store Betydning. „Det Gode“, siger
Goethe i en Samtale med Eckermann (1. April 1827), „er intet
Produkt af menneskelig Reflexion, men er medskabt og medfødt
skøn Natur. Mere eller mindre er det alle Mennesker medskabt,
men i høj Grad enkelte, ganske fortrinligt begavede Aander.
Disse have aabenbaret deres guddommelige Indre ved store
Gerninger eller Lærdomme, og det har saa ved Skønheden i
sin Fremtræden vundet Menneskenes Kærlighed og vældigt ledet
til Ærefrygt og Efterlignelse. Men ved Erfaring og Visdom
kunde det Sædeligt-Skønnes og Godes Værdi komme til Bevidsthed,
idet det Slette i sine Følger godtgjorde sig som Noget, der
forstyrrede den Enkeltes og det Heles Lykke, hvorimod det
Ædle og Rette viste sig som Noget, der hidførte og fæstnede
den særlige og den almene Lykke.“ Der ligger en stor biologisk
og etisk Sandhed i disse Goethes Ord\footnote[1]{Smlgn. mine Etiske Undersøgelser. p. 72. 82 f. --- Om Fichte og Carlyle se Den nyere Filosofis Historie'\textsuperscript{1}. Il. p. 151. 385; smlgn. ogsaa p. 358 (Comtes positivistiske Kalender).} —en Sandhed,
som ikke forringes ved, at vi maa lade Aarsagen til de spontane
Variationer (hvad Goethe kalder det Medskabte eller Medfødte
hos de fremragende Mennesker) staa hen som et Problem.

Hvad ingen Morallov og ingen etisk Deduktion kan yde,
det yde de store Forbilleder i levende og anskuelige Former.
Uden et Forhold til saadanne Forbilleder kan det personlige
Grundlag for det Etiske ikke udvikle sig til Inderlighed og
Fasthed. Hvilke Forbilleder der vælges, vil bero paa den vælgende
Individualitets Trang og Vilkaar. Maaske vil ikke det
samme Forbillede kunne vejlede os hele Livet igennem, ligesom
det ikke er de samme Stjerner, der vejlede den Søfarende paa
% --- p. 201 ---
\setcounter{footnote}{0}%
\opage{201}den nordlige og den sydlige Halvkugle. Der er et stadigt Vexelforhold
mellem de etiske Forbilleder og Erfaringen, saaledes
som Goethe saa klart har set det.

4. Historien viser, at den etiske Stræben hos de Enkelte
saavel som i Samfundet i enhver Periode og i ethvert Folk gaar
ud paa at fremkalde og udvikle visse Karakteregenskaber, der
da betragtes som Dyder. Et etisk Standpunkt betegnes ikke
blot ved det Følelsesgrundlag, der bestemmer Vurderingen, heller
ikke blot ved den objektive Maalestok, der anlægges, eller ved
de Handlingsmotiver, der fordres, men ogsaa ved de Karakteregenskaber,
der beundres og begunstiges. En Dyd kan, ligesom
en Pligt, begrundes paa forskellig Maade; men det faar en særegen,
baade praktisk og teoretisk Interesse at se, hvilke Dyder
der til forskellige Tider særligt drages frem.

Det var Grækerne, der først (i den europæiske Menneskehed)
opstillede en individuel Etik. Tidligt dannede der sig
en Række af Egenskaber, som det græske Folk beundrede og
fordrede. Først i Rækken kom de fire Hoveddyder Visdom,
Tapperhed, Selvbeherskelse og Retfærdighed til at staa. \textsc{Platon}
(Statens 4. Bog p. 427 E) tager det som givet, at Godhed bestaar
i disse fire Egenskaber, og forsøger saa dels ad psykologisk,
dels ad sociologisk Vej at vise, at al Dyd kan føres tilbage
til disse fire senere saakaldte Kardinaldyder (af cardo, et
Dørhængsel: altsaa de Egenskaber, om hvilke Alt drejer sig).
Det Mærkelige er nu, at Platons psykologiske Udledning er aldeles
selvstændig i Forhold til hans sociologiske Udledning. Det
var altsaa hans Overbevisning, at det enkelte Individ i og for
sig selv har en Opgave at løse. Denne Opgave bestod for Platon
i at gøre sig til et personligt Kunstværk, indenfor hvilket
de forskellige Evner og Tilskyndelser virke harmonisk sammen.
Den Vægt, der i hele den græske Etik, ikke blot hos Filosoferne,
lægges paa Harmoni, Begrænsning og Maadehold, bliver
kun ret forstaaelig, naar vi erindre, at græsk Kultur til sin
historiske Forudsætning havde en lidenskabelig og kaotisk Fortid.
Der var titaniske Kræfter, som skulde formes, før olympisk
Ro og Højhed kunde komme til at raade. De forskellige
Kardinaldyder udtrykke hver fra sin Side det saaledes vundne
% --- p. 202 ---
\setcounter{footnote}{0}%
\opage{202}Ideal. Visdommen lader Tanken være den ledende Magt, Tapperheden
holder den personlige Værdighed oppe overfor ydre
Modstand og indre Slappelse, Selvbeherskelsen bevirker, at de
sanselige Drifter begrænses uden at dødes, og Retfærdigheden
staar som Udtryk for alle sjælelige Elementers Sammenspil:
den betegner den sjælelige Tilstand, hvor hver enkelt Del af
Sjælen øver sin egen Funktion uden at hæmme de andres; den
er altsaa en sjælelig Harmoni. Skønt \textsc{Aristoteles} ikke optog
denne Lære om de fire Hoveddyder, opfattede han dog Dyden
som en individuel Harmoni, den Enkelte i og for sig, afset
fra hans Forhold til Samfundet, stræber efter at virkeliggøre.
Den klassiske græske Etik endte paa dette Punkt i det
store Problem, hvorvidt de Fordringer, Samfundet stiller til Individet,
og de, dette selv stiller til sin fulde Udvikling, kunne
bringes til at samstemme. Platon's og Aristoteles' individuelle
Etik passede desuden kun for dem, der uafhængigt af materielt
Arbejde kunde leve for deres personlige Udvikling. Og baade
for Platon og Aristoteles var der en Dyd, der tilsidst stod over
alle og truede med at sprænge den indre personlige Harmoni:
i Evnen til Tænken og Forsken, især til spekulativ Kontemplation,
saa de store græske Filosofer Menneskets egentlige Adelsmærke.
Den teoretiske Visdom blev for dem det Højeste.

De fire Kardinaldyder vedbleve dog at følge den etiske
Tænkning paa dens Vej. Stoikeren \textsc{Panaitios} optog dem (i det
2. Aarh. f. Kr.), uden at det klart kan ses, hvorledes han bragte
dem i Forbindelse med sin Lære om Menneskekærligheden, og
fra ham har igen \textsc{Cicero} faaet dem og gengivet dem i sit bekendte
Skrift „om Pligterne“ (De officiis). Foregik der sandsynligvis
allerede en vis Forskydning i den klassiske græske
Dydsliste, da den gik over fra Platonismen til Stoicismen, saa
foregik der ganske sikkert en Forskydning, da den fra græske
Tænkeres gik over i den romerske Retors Hænder. Borgerlig
Æresfølelse og Patriotisme kom til at spille en særligt stor
Rolle; honestum, det, en god Borger skal fyldestgøre, traadte i
Stedet for „det Skønne“, den sociale Pligt i Stedet for den personlige
Harmoni. Kun i politisk Færd, ikke i Privatlivet kunde
efter Romerens Opfattelse Sjælsstorheden godtgøres.

% --- p. 203 ---
\setcounter{footnote}{0}%
\opage{203}I Urkristendommen brydes, som vi allerede have set (II, 4),
Menneskekærligheden med den ekstatiske Forventning. I nøje
Sammenhæng med dette sidste Motiv staar Gudsforholdet i Tro
og Lydighed som Betingelse for Sjælens Salighed. Og ligesom
hos Platon og Aristoteles Visdommen tenderede til at hæve sig
over den personlige Harmoni, saaledes har indenfor den kristelige
Etik Troen en stadig Tendens til at fortrænge Kærligheden
fra Forrangen mellem Dyderne. Konsekvent maatte, jo mere
den kristelige Etik betonede Autoritetsprincipet, Troens Lydighed
blive den højeste Dyd, --- ligesom den første og højeste
Synd var Hovmod og Stolthed (VI, 1). Naar Kirkefædrene førte
al Dyd tilbage til Kærlighed, da tænkte de herved gerne paa
Kærlighed til Gud, som er Et med Tro. Og da vi nu kun ved
Kirkens Hjælp kunne lære Gud at kende, afløses den Sætning:
„Uden Tro ingen Dyd!“ snart af den Sætning: „Udenfor Kirken
ingen Dyd!“ Hedningenes Dyder bleve da kun at betragte
som glimrende Laster: ti det er Hovmod og Opblæsthed at ville
øve Dyden for dens egen Skyld!\footnote[1]{En interessant Sammenligning mellem Panaitios, Cicero og Ambrosius giver R. Thamin: St. Ambroise et la morale chrétienne au 4e siècle. Paris 1895.}

Alligevel følte man fra Kirkens Side paa Etikens Omraade,
ligesom paa Dogmatikens, Trang til at optage Elementer og Former
fra den græske Tænkning. I \textsc{Ambrosius} Skrift De officiis
ministrorum, der blev Middelalderens Haandbog i Etik,
blev Platon's og Panaitios' gamle Inddeling af Dyderne gennem
Cicero som Mellemled overført paa kirkelig Grund. Det kunde
naturligvis ikke ske uden forskellige Omtydninger og Forskydninger.
I Stedet for den individuelle Harmoni, som Platon kaldte
Retfærdighed, træder den levende Menneskekærlighed, der fører
til Hengivelse --- i Stedet for Visdommen træder Troen, Lydigheden
mod Autoriteten, --- Tapperheden viser sig især som
Taalmodighed, og i Stedet for Selvbeherskelsen sættes Renheden.
Senere har \textsc{Augustinus} forsøgt at vise, at de fire traditionelle

') Augustinus: De civitate dei XIX. 25. --- Smlgn. iøvrigt De
Wette: Lærebog i den christelige Sædelære og sammes Historie. Dansk
Overs. § 148. 156. 171.
% --- p. 204 ---
\setcounter{footnote}{0}%
% CHECK p. 204: notes paired with the marks in order (the layer misread a number)
\opage{204}Dyder hos den Kristne fremtræde som fire forskellige Former,
under hvilke Kærligheden til Gud ytrer sig\footnote[1]{De moribus ecclesiæ catholicæ. c. 15; 25. --- Augustinus henviser til «Visdommens Bog“ (VIII, 7), hvor de fire Dyder nævnes.}. Paa en mere
udvortes Maade forbandtes i Middelalderen, f. Ex. hos Thomas
Aqvinas, antik og kristelig Etik, idet de tre „teologiske“
Dyder, Tro, Haab og Kærlighed, anbragtes over de fire «filosofiske“
Dyder. Men trods al Omtydning og Omvurdering er
det dog et Tegn paa Sammenhængen i den etiske Udvikling, at
man saa længe har kunnet bruge, og endnu undertiden\footnote[2]{Victor Cathrein: Philos. mor. p. 98---111. --- S. Alexander: Moral Order and Progress. 2. ed. London 1891. p. 250 f. --- P. Natorp: Grundlinien einer Theorie der Willensbildung (Archiv für system. Philos. 1895).} den
Dag i Dag mener at kunne bruge den gamle græske Firdeling
af de etisk værdifulde Karakteregenskaber.

Med Renaissancen kommer den kraftige Selvhævdelse frem
som en individuel Tendens, der plejes og beundres. Livsglæde,
Magt og Nydelse af mere eller mindre ideel Art stilles i første
Linie. I den filosofiske Teori kommer hos en hel Række af
Tænkere fra det 16. til det 19. Aarhundrede Aandshøjhed (sublimitas)\footnote[3]{Dette Begreb havde sit Forbillede i et Begreb, der under forskellige Navne forekommer hos Platon, Aristoteles og Stoikerne, uden at det ret finder sin Plads i deres Inddeling af Dyderne. Smlgn Platon: Statens 6. Bog p. 486 A (διανοίας μεγαλοπρέπεια, Sindets Højhed). Aristoteles: Eth. Nic. IV. 7 (μεγαλοψυχία, det at have en stor Sjæl). Seneca: Dialogi I, 4---6; II, 11. 14 (magnanimitas). Baade Aristoteles og Seneca betragte denne Egenskab som den højeste.},
Livsmod (animositas) og Højmod (generositas) til
at staa som de største Dyder. De udledes af Selvopholdelsestrangen
og betragtes som dennes højeste og ædleste Former.
Især ere Telesio, Bruno, Campanella, Descartes og Spinoza her
at nævne. Kant og Fichte slutte sig i deres Lære om de etiske
Egenskaber til denne Række, idet for Kant den med Sjælskraft
hævdede personlige Værdighed som Modsætning til Kryberi, for
Fichte den ubetingede Selvvirksomhed som Modsætning til Træghed
stod som det Højeste paa Karakterens Omraade. I denne
Karaktervurdering virker Reaktionen mod den kristelige Etiks
% --- p. 205 ---
\setcounter{footnote}{0}%
\opage{205}Tendens til at gøre Lydighed og Ydmyghed til de højeste
Dyder.

Hos en anden Række af Tænkere virker Erindringen om
den stoiske Lære om Menneskekærligheden sammen med Kristendommens
saa ofte af Troen fortrængte Kærlighedsbud til at lægge
Hovedvægten paa Medfølelsen og Menneskelighedsfølelsen som
begrundede i selve den menneskelige Natur og som Karakteregenskaber
af højeste Værdi. Denne Retning repræsenteres af
Shaftesbury, Hutcheson, Hume, Adam Smith, Comte, Schopenhauer
og Herbert Spencer.

For denne sidste Retning er det aabenbart lettest at faa
en naturlig Sammenhæng mellem individuel og social Etik. I
Aandskraften og Aandshøjheden stod Individet som frit og uafhængigt,
baaret af Bevidstheden om at have sit Livs Betingelser
overvejende i sig selv. Et Hovedtræk, som særligt fremhæves
i disse Dyder, er, at de føre til at foragte den ydre Ære
i Bevidstheden om den indre Værdi. Ligesom i den antike Harmonilære
staar Individet her skilt fra Samfundet. I Højmodet,
som Descartes og Spinoza skildre det, er det dog et væsentligt
Element, at Individet har Bevidstheden om at høre til et større
Hele, hvis Interesse foretrækkes for de rent individuelle. Det
er den store, udover den individuelle Personlighed rækkende
Horisont, der betinger Højmodet. Men derfor danner dette ogsaa
Overgangen til de Karakteregenskaber, som den anden Række
af Filosofer lagde Hovedvægten paa. Og det forudsætter disse.
Til Højmodet kommer man kun ved at vandre ikke blot Selvhævdelsens,
men ogsaa Hengivelsens Veje. Den af Kant beskrevne
Hævden af personlig Værdighed er ikke blot betinget
ved, at Mennesket har Loven for sin Handlen i sit eget Indre,
men ogsaa derved, at denne hans indre Lov er universel, gælder
for alle Fornuftvæsener. Heller ikke denne Karaktertendens
er derfor psykologisk og historisk mulig, uden naar Individet
formaar at føle sig solidarisk med et større Hele.

Som vi tidligere have søgt at vise (se IX, 1), staa Selvhævdelse
og Hengivelse ikke nødvendigvis i Modstrid mod hinanden.
De udtrykke to Tendenser i den menneskelige Natur,
hvis indbyrdes Forhold kan være højst forskelligt hos de for%
% --- p. 206 ---
\setcounter{footnote}{0}%
\opage{206}skellige Individer, men som meget vel kunne staa i harmonisk
Forhold til hinanden, og som begge maa indgaa som Elementer
i den fuldkomne Karakter. Hengivelsen udvikler og udvider
Individets Væsen, og Selvhævdelsen dygtiggør det til at udfylde
sin Plads i det Hele, til hvilket det hører. Det platoniske Begreb
om Retfærdighed som personlig Harmoni kunne vi optage
til at udtrykke den indre Forbindelse, den højere Enhed af Selvhævdelse
og Hengivelse. Det antike Harmonibegreb maa udvides
og uddybes ved Hjælp af de etiske Erfaringer, der ligge
udtrykte i Stoicismens, Kristendommens, Renaissancens og den
moderne Humanitets Karaktervurderinger. Men en saadan Ændring
af Begrebets Indhold er meget vel mulig uden at sprænge
dets Ramme. Vor Etik er den græske Etik, som den senere
etiske Udvikling gennem sine mangehaande Svingninger har
kunnet ændre og berigtige, men som aldrig vil kunne fortrænges,
saalænge der skal bestaa Noget, der med Rette kan kaldes
menneskelig Etik.

Ved særlig Undersøgelse af Selvhævdelsen og Hengivelsen
hver for sig ville vi vise Muligheden og Nødvendigheden af at
forene dem til harmonisk Samvirken og derved at arbejde paa
Virkeliggørelsen af den højeste menneskelige Karakteregenskab.
Det er ganske vist kun en indirekte Beskrivelse af Retfærdigheden,
som naas ved at undersøge, hvorledes Selvhævdelsen og
Hengivelsen hver fra sin Side finde deres Fuldendelse i den.
Men Spørgsmaalet er, om det er muligt at naa videre, naar Idealet
ikke skal blive til en i Luften svævende Abstraktion. I hvert
enkelt Tilfælde udvikler Retfærdigheden sig enten overvejende
af Selvhævdelsen eller overvejende af Hengivelsen. Den faar
derved sin forskellige Klangfarve hos de forskellige Personligheder.
Overfor denne Rigdom vilde en ideal Skildring af den
højere Enhed af de to stadigt i Brydningsforhold staaende Retninger
i vor Natur faa en mat Tone. Det er da bedre, at vi
skærpe Problemet og holde os til de forskellige Veje, ad hvilke
Udviklingen kan nærme sig til den højeste Karakteregenskab.
Maaske vil den Tid komme, da det bliver muligt for en Etiker
at give en fuldkomnere og mere positiv Beskrivelse af den.

% --- p. 207 ---
\setcounter{footnote}{0}%

\opage{207}\chapterhead{XI.}{SELVHÆVDELSE.}

1. De Dyder og Pligter, som høre ind under Selvhævdelsen,
kunne henføres til tre Hovedformer: Selvopholdelse, Selvbeherskelse
og Selvstændighed. Ved dem alle vil det vise sig,
at Individets Stræben efter at hævde sig selv etisk set betinges
og begrænses ved Hensynet til hele Slægtens Bestaaen og Udvikling.
Som allerede bemærket er der i den nyere Tid, paa
Grund af de ivrige og betydningsfulde Emancipationsbestræbelser,
en vis Tendens til at overse dette, medens det er lettere
at paavise paa mere primitive Udviklingstrin. Men naar man
gaar tilbage til det Grundlag, fra hvilket vi her betragte Forholdet
mellem Individ og Samfund, vil det dog vise sig, at hin
Sætning er rigtig.

% CHECK p. 207: centred line: heading or verse
\begin{center}a. Selvopholdelse.\end{center}

2. Dersom det ikke i sig selv er et Gode at leve, er al
Etik urimelig. Direkte eller indirekte gaar al Etik ud paa at
bevare Livet, at frede det og udvikle det videre. Den fortsætter,
hvad Naturen selv indleder. Selv hvor en klar Bevidsthed
ikke er vaagnet, gør der sig hos alle levende Væsener et Selvopholdelsesinstinkt
gældende, som lader Individet søge, hvad der
er gavnligt for dets Bestaaen, og sky det Modsatte. \textsc{Kant} mente
endog, at Instinktet eller Driften her var saa stærk fra selve
Naturens Haand, at man ikke kunde tale om Pligt til Selvopholdelse:
ti hvad Enhver uundgaaeligt vil af sig selv, kan han
ikke forpligtes til\footnote[1]{Metaphysische Anfangsgründe der Tugendlehre. Einleitung § 4.}. Denne Anskuelse hænger for en Del sammen
med Kants urigtige Antagelse, at vi kun kunne forpligtes
til det, som vi ugerne ville. Men den hviler tillige paa den
ikke mindre urigtige Antagelse, at vi virkeligt stedse og uundgaaeligt
søge, hvad der tjener til vor Selvopholdelse. Det er
% --- p. 208 ---
\setcounter{footnote}{0}%
\opage{208}ikke en Gang saa hos Dyrene, hos hvem Instinkterne dog virke
i langt større Omfang og med langt større Magt end hos Menneskene.
Selv dér kan Instinktet lede vild, eller det kan svækkes
og ændres. Men for Mennesket bliver i hvert Tilfælde
Selvopholdelsen ikke blot Instinktets, men ogsaa den egentlige
Villies Sag.

3. Legemlig Sundhed og Kraft er det nødvendige Grundlag
for al videre og højere Udvikling af Livet. Vi raade til
enhver Tid kun over en begrænset Sum af Energi, og det bliver
derfor en vigtig Opgave at sørge for, at den bevares og forøges
og ikke gaar til Spilde. Vi kunne ikke frembringe Noget
af Intet; selv Helten er Intet uden Mad og Drikke. Askesen
med sine selvvoldte Smerter og Savn, al Slags Selvplageri og
ugrundet Mismod, Letsindighed og udsvævende Sanselighed kunne
hver paa sin Maade øde den Kapital, der kunde bruges til værdifulde
Formaal. Fordomsfuldhed, Kortsynethed og Egenraadighed
føre enten til at vrage de nødvendige Midler eller til at
forbruge dem uden Nytte. Naar det staar klart for den Enkelte,
at Følgerne af hans Adfærd i denne Henseende kunne ramme
ikke blot ham selv, men ogsaa Andre, dels dem, han selv maatte
sætte ind i Verden, og til hvem hans svækkede Natur kan gaa
i Arv, dels ogsaa dem, som maa undvære hans Hjælp og hans
Arbejde, saa vil Omsorgen for den fysiske Selvopholdelse staa
som en betydningsfuld Pligt.

Men heri ligger ogsaa Begrænsningen for den. Det vilde
være meningsløst at tilbringe Livet med at samle en Kapital,
der aldrig bruges, at gøre Livet til ét eneste stort Middel uden
Formaal. Saa er Asketen endda fornuftigere, da han trænerer
sig for at opnaa et fjernt Formaal. At leve er at bruge Kræfterne,
og de bruges kun ved at forbruges. En stor Fysiolog
har endog givet den paradoxe Definition: „Liv er Død!“ fordi
enhver Livsytring medfører Opløsning af det organiske Væv. En
snever Opfattelse af Selvopholdelsen vilde gøre den til Et med
Opsamlen af Stof og ikke se, at Forbruget er en nok saa væsentlig
Ting i Livet som Stofoptagelsen, selvom det er det, der tilsidst
gør Ende paa Livet. Det er endda ikke altid, at vi kunne
bevare det harmoniske Forhold mellem Optagen og Forbrug.
% --- p. 209 ---
\setcounter{footnote}{0}%
\opage{209}Der kan med etisk Nødvendighed komme et Misforhold frem.
En Opgave, som det er min Pligt at løse, kan nødvendiggøre
Tilsidesættelsen af, hvad Selvopholdelsespligten under sædvanlige
Forhold vilde fordre. Döbeln (vid Jutas) havde ikke Tid
til at være syg. Livet maa endog helt ofres, naar vi for at bevare
det maatte fornegte, hvad der giver det dets egentlige Værd.

At leve for sin Sundhed er en Uting. Det kan ske, at man
af Hensyn til sine svigtende Kræfter maa trække sig tilbage
til mere begrænsede Opgaver; men det at begrænse sine Formaal
er ikke det Samme som at gøre til Formaal, hvad der kun
skal være Middel. \textsc{Platon} spotter over Lægen Herodikos' Adfærd
under en uhelbredelig Sygdom: „Da han ingen anden
Virksomhed havde, levede han blot for at kurere paa sig selv
og var højst ulykkelig ved den mindste Afvigelse fra den tilvante
Levemaade; paa denne Maade opnaaede han ved sin Visdom
under en haard Dødskamp en høj Alder.“ Platon bemærker,
at for de Fattiges Vedkommende følger det af sig selv, at
de ikke kunne leve for at kurere paa sig; det er kun de Rige,
der kunne stille sig dette som Opgave; og han mener, at „den,
der ikke kan følge den naturlige Levevis, heller ikke bør helbredes,
da det hverken er til Gavn for ham eller for Staten“\footnote[1]{p. 122. ) D. Fr. Strauss in seinem Leben und seinen Schriften geschildert,}.
Vi vilde ikke gaa saa strengt frem. Vi have lært, at der under
og ved Lidelsen kan udfolde sig Egenskaber, som maaske hidtil
vare trængte tilbage, og som kunde give Livet en ny Værdi
for den Lidende selv og for dem, der paavirkes af hans Forbillede.
Hans Karakter faar maaske først nu sin Fuldendelse.
Vi tillægge Tankens og Følelsens Liv Værdi, afset fra de ydre
Virkninger. Vi se den enkelte Personlighed ikke blot som Medlem
af Staten, men ogsaa som Medlem af Familien, af Vennekredsen,
og se ham maaske selv i hans udadtil uvirksomme
Tilstand være til den største Glæde og Bestyrkelse for Andre.
Om den berømte Fritænker David \textsc{Strauss}' sidste lidelsesfulde
Dage, der vare saa opbyggelige for hans Venner, skriver hans
Biograf E. Zeller\footnotemark[1]: „Han fordrede ikke af Naturen, at den

') Staten. 3die Bog. Heises Overs. p. 178—182.
% CHECK p. 209: note mark without a note

% --- p. 210 ---
\setcounter{footnote}{0}%
\opage{210}skulde spare ham for Noget af det, dens Løb og dens Love
bragte med sig; men saa meget des mere fordrede han af sig
selv, at han skulde vide at finde sig i disse Love, at gøre ogsaa
Lidelsen til Stof for en sædelig og aandelig Virksomhed og
at finde det Velgørende midt i Smerten.“

4. Vi komme nu til det meget drøftede Spørgsmaal om
Selvmordets Berettigelse. Dette Spørgsmaal kan ifølge Sagens
Natur kun opkastes, forsaavidt Selvmordet er en virkelig Handling,
Frugt af Overvejelse og Beslutning. Dersom Selvmord
stedse udspringer af afgjort Sindssygdom, falder det ikke ind
under etisk Betragtning. Det Selvmord, som udspringer af en
fortvivlet Situation, i hvilken Individet ved sin egen Adfærd
har bragt sig, kan ogsaa være af den Beskaffenhed, at det ikke
er nogen overlagt Handling, men Virkning af en pludselig, al
Besindelse overvældende Stemning. Her falder dog den etiske
Vurdering ikke bort, men rettes (ligesom ved Handlinger, der
foretages i Drukkenskab) mod den Adfærd, ved hvilken Individet
har bragt sig i sin fortvivlede Tilstand. Saaledes ligger
hos den, der ved bedrageriske eller dog uforsvarlige Pengespekulationer
har styrtet sig i Ulykke og saa taget Livet af sig,
den egentlige Skyld før Selvmordet. Ligesaa hos \textsc{Mussets}
„Rolla“, der tager sig af Dage, efter at han har nedbrudt sin
fysiske og sin aandelige Personlighed ved Udsvævelser. Udsvævelser
kunne svække Modstandskraften, selvom de ikke
direkte frembringe fortvivlede Situationer. Alkoholisme og Selvmord
stige, som Statistiken viser, jevnsides.

I saadanne Tilfælde, hvor Selvmordet tydeligt udspringer
af, at man vil unddrage sig bestemte Forpligtelser, er der heller
ingen Grund til lang Debat. Det er en aaben Desertion, en
Handling af samme Art som Selvlemlæstelse for at slippe for
Krigstjeneste. Efter den romerske militære Straffelov medførte
Selvmordsforsøg Halshugning, idet det betragtedes som Forsøg
paa Desertion. Den, der myrder sig selv for at slippe fra en
pinlig eller farlig Situation, i hvilken han og hans Familie ere
komne, og som derved unddrager sine Nærmeste sin Hjælp,
udøver utvivlsomt en forkastelig Handling. Men det Uetiske
ligger da ikke i og for sig i, at han gør Ende paa sit Liv; det
% --- p. 211 ---
\setcounter{footnote}{0}%
\opage{211}ligger i, at han unddrager sig sine Forpligtelser, Noget, han
ogsaa kunde gøre paa anden Maade.

Spørgsmaalet er, om alle Tilfælde af Selvmord kunne henføres
under de her nævnte Arter. I saa Fald vilde Problemet
falde bort. Før jeg gaar nærmere ind paa at undersøge dette,
vil jeg berøre nogle Omstændigheder, der hænge sammen med
den, som det synes, stærke Tiltagen af Selvmord i den nyere Tid.

5. I sin Bog Om Selvmord i Kongeriget Danmark meddeler
\textsc{Kayser} en Række Breve, som ere skrevne af Selvmordere lige
før Døden. Den fælles Tone i disse Breve er Mangel paa Evne
til at udholde sin Skæbne, til at fortsætte Kampen for Tilværelsen
under de givne Vilkaar. Selvom man ikke deserterer, saa
kapitulerer man. Man føler sig staaende overfor en uovervindelig
Modstand. Det er hverken Mod eller Fejghed, der lægger
sig for Dagen i saadanne Selvmord, men en Svækkelse af Villieskraften,
en Kuelse af Livsdriften og Livslysten. Af en af
Statens statistiske Bureau udgiven Undersøgelse af Selvmordene
i Danmark i Tiaaret 1886— 1895 fremgaar det, at den vigtigste
Aarsag har været Tungsind, Livslede og Sindssyge, idet over
\textsuperscript{1}/\textsubscript{4} af alle Selvmordene, for Kvindernes Vedkommende endog
\footnotemark[2]\textsuperscript{2}/\textsubscript{5} af Selvmordene, ere begaaede af disse Motiver. Dette stemmer
med det Udtryk, man faar af de af Kayser udgivne Breve.
Med Hensyn til de forskellige Aldersklasser ses, at i Ungdommen
(fra det 15. til det 25. Aar) spiller ulykkelig Kærlighed
og Anger eller Frygt for Straf en betydelig Rolle som Selvmordsmotiver;
i Manddomsalderen træde Drikfældighed og økonomiske
Bekymringer frem i første Linie; i Alderdommen blive
legemlig Sygdom og Tungsind, Livslede eller Sindssygdom mere
og mere de overvejende Motiver, og for Kvindernes Vedkommende
tidligere end for Mændenes. Den Svækkelse af Modstandskraften
og den Blindhed overfor Livets givne Værdier
eller overfor lysere Muligheder, hvorom Selvmordenes Hyppighed
(--- i København et Selvmord næsten hver tredie Dag —)
vidner, have vel for en stor Del Grunde, der ligge langt tilbage,
i Temperamentet og i det uvilkaarlige Stemningsliv. Men
der kan sikkert ogsaa peges paa Omstændigheder, som følge
med hele den nyere Tids Kulturudvikling. Den moderne Kultur
% CHECK p. 211: note mark without a note
% --- p. 212 ---
\setcounter{footnote}{0}%
\opage{212}med sin ensidige Udvikling af de intellektuelle Evner paa de
andre Evners Bekostning, med sine vide Horizonter og sin
rastløse, men ogsaa ofte plan- og formaalsløse Foretagelsesaand
gør det vanskeligere at fastholde den Sluttethed og Koncentration,
den Begrænsning i Henseende til Formaal og Midler,
som er en nødvendig Forudsætning for et sundt og kraftigt
Villiesliv. For Tanke og Fantasi ere vide Horizonter tiltalende;
men Villien maa, naar den skal være Mere end blot Drift og
Ønske, rette sig mod noget aldeles Bestemt, og mod Noget,
som er indenfor Rækkevidde. Den moderne Tid er en Tvivlens
Tid, ikke særligt gunstig for Harmoni mellem de sjælelige
Kræfter. Det er den aandelige Hygiejnes Opgave at raade
Bod paa den Disharmoni og Sygelighed, som under disse Forhold
saa let udvikler sig. Pædagogen og Lægen maa virke sammen.
Alt, hvad der tjener til kraftigere Udvikling af de aktive
Evner --- hvortil ikke blot den fysiske Arbejdskraft, men ogsaa
Evnen til selvstændig Iagttagelse og Tænkning og den kunstneriske
Fantasi høre --- vil virke lægende overfor Villiessygdommen.
Ikke ved at afbryde Kulturarbejdet og kue Tvivlen
--- hvad der da ogsaa er en Umulighed --- men ved at styrke
Selvvirksomheden og vække Tilliden til de menneskelige Kræfter,
hvor disse virke paa deres rette Omraade, bidrager man
til at faa Bugt med den holdningsløse, rugende eller blaserede
Tvivl, der følger den nyere Tids storartede Udvikling som Skyggen
følger Legemet.

Dertil kommer endnu en anden Ting. Den moderne Kultur
gør paa én Gang Individerne mere afhængige --- ved at
vække flere Fornødenheder hos dem --- og mere isolerede ---
ved at „emancipere“ dem og stille dem paa deres egne Ben.
Og Isolationen naar sit Højdepunkt, naar Individet med sine
mange utilfredsstillede Fornødenheder staar ensom i den store
Travlhed rundt om. De enkelte Individer maa fægte mere hvert
for sig end under den gamle Tingenes Orden. De nye Associationer
modvirke vel denne Isolation; men de ere endnu i deres
Barndom og have neppe endnu haft tilstrækkelig Indflydelse
paa Sindelaget hos de Enkelte. Den Ulykkelige og Lidende føler
sig overladt til sig selv. En Dyreflok mishandler eller dræber
% --- p. 213 ---
\setcounter{footnote}{0}%
\opage{213}de Individer, der ikke kunne følge Flokkens Fart; i den menneskelige
Verden maa de ofte skøtte sig selv, om de end ikke
mishandles og nedtrampes. Et levende Sympatiforhold, hvor den
Enkelte føler Deltagelsen fra Andres Side og selv med Deltagelse
følger deres Færd, vil gøre, at hans Livsmod ikke saa let knækkes,
og hans Horizont ikke ganske formørkes. Den, der ikke
slipper Andre, vil heller ikke saa let føle sig forladt af Alle.
Blaseretheden og Trætheden ville ikke saa let træde i Friskhedens
og Interessens Sted. --- Det synes, som om Selvmord
aftage under store politiske Kriser, maaske fordi Sindene da
ere optagne af de store Begivenheder og de fælles Interesser.
Hvis det er saa, maa stadig Opgaaen i fælles Opgaver være en
Magt, der kan hindre Mismod og Selvopgivelse. Den Resignation,
som skal kunne bære ud over en stor Skuffelse eller en
stor Nød, bliver kun psykologisk mulig, hvor der er en stor
Livsinteresse, man kan gaa op i og leve for, naar det, hvori
man hidtil har fundet sit egentlige Jeg, maa opgives. Selvmordet
udspringer af, at man betragter det, hvori man hidtil har
fundet sit Jeg, som det Eneste. Dette tidligere Jeg maa ofte
helt forsages, for at Selvmordet skal hindres. Drift til Selvmord
er Udtryk for en Selvhævdelsestrang, der kun viger for den
højere Art Selvhævdelse, der fremtræder som Forsagelse af sig
selv. Det er i denne Forstand, \textsc{Mme} de \textsc{Stael} har sagt: Le
renoncement å soi méme est en tout I'opposé du suicide.

Den, der har boet som ensom Fremmed i en stor By, véd,
hvilken trykkende Forladthedsfølelse der kan betage Sindet, og
hvilken Opmuntring det kan være blot at faa et venligt Ord af
en Forbigaaende, som man spørger om Vej eller viser en Tjeneste.
Saaledes vil ofte en sympatisk Ytring, en venlig Tiltale
kunne standse Selvmorderen paa hans Vej. Han vil føle, at
Baandet, som knytter ham til Slægten, ikke er brudt, og vil
kunne føle Trang til ogsaa fra sin Side at knytte det fastere
igen. I gamle Dage vare Klostrene Tilflugtssteder for ulykkelige
og modløse Naturer, som dér kunde finde Støtte og Opgaver
i et fælles Liv. Hjemløsheden for den Art Naturer er
bleven større, efter at Klostrene ere blevne lukkede. Det er nu
lagt i deres egen Haand at finde Opgaver og Hjemsted. --- Det
% --- p. 214 ---
\setcounter{footnote}{0}%
\opage{214}gælder da om at sætte Kræfter i Bevægelse, som kunne ophæve
Isolationen af de enkelte Evner i Menneskets Indre og
mellem Menneskene i Samfundet. Herved hænger dette specielle
Spørgsmaal i sin etiske Grund sammen med de store, almindelige
etiske og sociale Spørgsmaal.

6. Skyldes Selvmord nu altid Sygdom i Villien, ligefrem
Sindssygdom eller ogsaa et bevidst Ønske om at unddrage sig
bestemte Forpligtelser? Eller kan der være en Ret, maaske
endogsaa en Forpligtelse til at begaa Selvmord?

Hvor man har hævdet Individets fulde Frihed til at raade
over sit Liv, er man stedse gaaet ud fra et rent individualistisk
Standpunkt. --- Saaledes lærte Stoikerne, at den Vise, der er
sig selv nok og sin Skæbnes absolute Herre, ikke blot vil lade
sit Liv, naar Fædrelandets Vel kræver det, eller naar han kun
kan bevare det ved at begaa en uetisk Handling, men at han
vil forlade Livet, naar han paa Grund af ydre Tilskikkelser ikke
kan føre det, som han vil. De senere Stoikere udstrakte denne
Ret til „at lede sig ud“ af Livet meget vidt. „Ikke blot i den
yderste Nød“, siger \textsc{Seneca} (Epist. 70), „men saasnart Skæbnen
begynder at blive noget betænkelig“, vil den Vise tage sig
af Dage. \textsc{Marcus} \textsc{Aurelius} fordrer blot, at man skal forlade
Livet uden Vrede og Sorg, saaledes som man gaar ud af et
Værelse, hvor der er Røg (Comment. V, 29; X, 8). At ende
sit Liv af Fejghed og Blødagtighed ansaa Stoikerne for usædeligt,
--- skønt de paa den anden Side mente, at fejge og blødagtige
Mennesker vare uværdige Gæster ved Livets Fest, og at
de derfor maaske gjorde bedst i at fortrække. --- I den nyere
Tid forsvaredes Selvmord fra et extremt individualistisk Standpunkt
af en Række Forfattere i det 18. Aarhundrede. \textsc{Montesquieu}
f. Ex. begrunder Selvmordets Berettigelse ved, at da
Samfundet er grundet paa gensidig Fordel, maa man have Lov
til at trække sig tilbage derfra, naar man kun har Besvær af
at leve i det; man overtræder altsaa ingen Pligt ved at gøre
Ende paa sit Liv, naar man finder det uudholdeligt (Lettres
persanes. Nr. 76).

Selvom man nu vilde antage, at det enkelte Individ havde
fri Ret til at forlade Livet, naar dette var uudholdeligt, vilde
% --- p. 215 ---
\setcounter{footnote}{0}%
\opage{215}det være meget vanskeligt at godtgøre, at Uudholdeligheden
virkeligt var tilstede i det enkelte Tilfælde. Ti hvorledes kan
Individet vide, at det har naaet Grænsen af sin Kraft og har
brugt alle Midler? --- Man har hertil svaret, at selve det, at
det naturlige Selvopholdelsesinstinkt er kuet og den naturlige
Angst for Døden overvunden, vidner om, at Mennesket maa
være bragt til det Yderste. Saaledes siger \textsc{Holberg} i sin 135.
Epistel: „Erfarenhed lærer, at Mennesker have Livet kær; hvorudover
man kan slutte, at de, som bedrive saadan Gerning
[Selvmord], maa være overvældede af Angst, Lidelse og Genvordighed,
ja udi saadan Stand, som heller opvækker Medlidenhed
end Vrede“. Men fordi jeg overvældes af Angst og Smerte,
saa min Modstandskraft knækkes, derfor kan jeg godt befinde
mig i en Illusion. Instinktet leder mig ikke altid sikkert, og
Instinktets Ophør kan derfor heller ikke altid være noget sikkert
Kendemærke. En øjeblikkelig Følelse af Angst og Nød er
intet paalideligt Udtryk for de virkelige Forholds Beskaffenhed.
Maaske befinder Individet sig netop paa det laveste Punkt under
Skæbnens rytmiske Svingning, saa at det vil gaa opad igen for
ham, naar han blot kan holde lidt ud endnu. Et er, at Individet
fortvivler, et Andet, om Sagernes Stilling virkeligt er fortvivlet.
En fuldstændigt begrundet Forvisning vil aldrig kunne
haves herom. Naar der efter Nogles Mening kan spores en Aftagen
af Selvmordernes Antal, naar Udvandringen tiltager\footnote[1]{Legoyt: Le suicide ancien et moderne. Paris 1881. p. 257 f.}, saa
er der dermed givet et Exempel paa, hvorledes nye Muligheder
og Udveje kunne holde Modet oppe. Slappelse af Energi og
Mod kan gøre, at man ikke opdager og benytter Udvejene; men
et objektivt Bevis for, at de ikke ere der, vil være vanskeligt,
om ikke umuligt at føre.

Dersom Individet er absolut suverænt (smlgn. III, 5---7),
vil det her heller ikke komme saa meget an paa, om en saadan
objektiv Bevisførelse er mulig. Men naar den højeste Maalestok
for den etiske Vurdering søges i det enkelte Individs Sammenhæng
med Slægten, dets Betydning for Slægten, kan der
ikke slaas af paa Fordringen om strengt Bevis. Det var der%
% --- p. 216 ---
\setcounter{footnote}{0}%
\opage{216}for en rigtig Tanke, naar i nogle græske Stater (Massilia og
Keos) Dadlen af Selvmord faldt bort, dersom Selvmorderen først
overbeviste Øvrigheden om, at der var tilstrækkelig Grund til
at foretage et saadant Skridt. I Thomas More's Utopia opfordre
Præster og Øvrighedspersoner de uhelbredeligt Syge, som
ikke mere kunne arbejde og ere en Byrde for dem selv og
Andre, til frivilligt at ende deres Liv, hvorimod et Selvmord,
der mangler saadan Autorisation, medfører Vanære. Der ligger
her den Forudsætning til Grund, at Individet ikke lever for sig
selv alene, men er stillet overfor Opgaver, som ere af Betydning
for Andre end det selv. Det kan være fuldstændigt rigtigt,
at Livet i det enkelte Tilfælde er uudholdeligt, naar Individet
ikke kender nogen anden Maalestok end sin egen egoistiske
Nydelse. En Tilnærmelse til en saadan Indskrænkning af Interessen
vil man finde hos alle Selvmordere, der handle med
Overlæg. De have sat Alt paa ét Kort, have gjort deres Liv
til et Hasardspil. Deres Lidenskab er koncentreret paa et eneste
Punkt, med hvilket hele Livet for dem staar og falder. Dette
gælder saaledes om Oldtidens beundrede Selvmord. De „sidste
Romere“ dræbte sig, fordi deres Horizont ikke strakte sig
videre end til Roms aristokratisk-republikanske Liv. Vi kende
baade mere omfattende og mere inderlige menneskelige Forhold,
som kunne stille Krav til os. Og selv Cato anerkender dette,
naar han (ifølge Plutark) raadede sin Søn, som han ikke vilde
have med sig i Døden, til at holde sig fjernt fra Politiken ide
Tider, som nu maatte komme. Naar Sønnens Ære ikke led ved
at følge dette Raad, vilde Faderen selv kunne have fulgt det
uden at fornægte sig selv. Den Kraft og Trods, som udgør det
Heroiske i det bevidst valgte Selvmord, har ofte Noget af det
Teatralske ved sig. Den vilde anvendes bedre, hvis den rettedes
mod det Maal: at gøre Livet taaleligere, om ikke for sig
selv, saa for Andre. --- Ogsaa hvor Selvmord udspringer af den
ovenfor (5) skildrede Villiessvækkelse, bukker Individet under
for Modstand paa et enkelt, objektivt set ofte ubetydeligt Punkt.
Men dér skyldes denne Koncentration paa et enkelt Punkt en
forudgaaende Svækkelse af Energien. I de her skildrede Tilfælde
er derimod i Forvejen Interessen med Lidenskabens hele
% --- p. 217 ---
\setcounter{footnote}{0}%
\opage{217}Energi samlet paa ét Punkt, og Alt falder derfor samtidigt med
det. To helt modsatte psykologiske Processer kunne saaledes
føre til samme Udgang.

Det Haab, den Enkelte har tabt for sit eget Vedkommende,
vil han kunne fastholde for Slægtens Vedkommende, især naar
han fatter den store Lære om de smaa Virkningers Betydning.
Den nyere Videnskab har aabnet vort Blik for, hvad uafbrudt
Ophoben af smaa Virkninger kan udrette. De mægtigste Koralrev
ere opstaaede ved Aflejring af utallige Smaavæseners Skaller.
Hvad den største Revolution ikke vilde kunne udrette, tilvejebringes
ved den stille Virken af dagligdags, uanselige Kræfter.
Den beskedneste Virken i den snevreste Kreds yder en
Skærv til Slægtens store Fællesliv. Intet behøver at gaa til
Spilde. Enhver Ytring af inderlig Følelse, af ideal Livsretning,
af Karakterfasthed kan faa sin Betydning og gribe ind paa en
gavnlig Maade, selvom vi ikke faa den teatralske Følelse af at
udøve noget Stort. \textsc{Rousseau} siger med Rette (i La nouvelle
Heloise), at hvis den, der føler sig fristet til Selvmord, vilde
undersøge, om der ikke endnu var nogle gode Handlinger, han
kunde udføre, en Fattig, han kunde hjælpe, en Ulykkelig, han
kunde trøste, en Undertrykt, han kunde forsvare, saa vilde han
sikkert holdes tilbage fra at udføre sit Forsæt.

Ligesaa lidt, som det kan godtgøres, at alle Muligheder
ere udtømte, ligesaa lidt vil det altsaa kunne godtgøres, at alle
Forpligtelser ere opfyldte. Jo alvorligere og dybere Livet opfattes,
des flere Muligheder og Forpligtelser fremstiller der sig.
Dette er den Hovedbetragtning, Etiken maa gøre gældende i
denne Sag. Iøvrigt maa de etiske Bestræbelser overfor Selvmord
snarere gaa ud paa at forebygge det ved at styrke Modstandskraften,
Villien til Livet og Sympatien med de Levende,
end paa at formulere strenge Domme over de udførte Selvmord.
Det gælder at bringe det dertil, at ethvert Individ føler sig
baade som Formaal --- ved at være et selvstændigt Led af
Slægten, --- og som Middel --- ved at være et Led, hvis Virken
kan blive af Betydning for andre Led. Fæstnen af Livets
Grundlag er bedre end Moraliseren, og da især, naar denne
kommer bagefter, og overfor en Handling, der i Regelen er
% --- p. 218 ---
\setcounter{footnote}{0}%
\opage{218}Virkning af voldsom Forblindelse. I alle Handlinger, som ere
etisk forkastelige, virker Forblindelse med; men ingen af disse
Handlinger motiveres saa hyppigt af virkelig Ulykke og Nød,
som den, hvorved Individet gør Ende paa sit Liv. ---

Der maa endnu spørges, om det under visse Forhold ikke
kan blive ligefrem Pligt at berøve sig selv Livet. Selvmord og
Selvopofrelse kunne grænse saa nær op til hinanden, at det
kan være vanskeligt, om ikke umuligt at drage Grænsen mellem
dem. Individet kan mene, at han staar i Vejen for Andres
Lykke, eller at han ved at forlade Livet befrier Andre for en
Fare. Saaledes naar En frygter for, at han under et skarpt
Forhør skal komme til at røbe Hemmeligheder af stor Betydning,
og derfor dræber sig for ikke at volde en ubodelig Skade.
Eller naar En, som er bidt af en gal Hund og mærker Sygdommens
Udbrud, dræber sig for ikke at komme til at bide
Andre. Eller naar en fattig Familiefader tager Livet af sig,
fordi han véd, at der, hvis han dør (og kun da), vil blive sørget
for hans Familie, som nu er udsat for Hunger og Nød.
Dersom saadanne Tilfælde gives, kan der jo ikke være Tale
om, at Individet unddrager sig Forpligtelser ved at tage sig af
Dage. Dets Handling udtrykker netop en Anerkendelse af, at
det ikke føler sig og sin Existens som det eneste eller højeste
Formaal, men betragter sig som Led af et større Hele, for
hvilket det maa ofre sig. Derimod staar stadigt den Vanskelighed
tilbage, hvorvidt virkelig alle Muligheder ere udtømte. Hvor
det gælder et Menneskeliv, hvor det gælder Tilintetgørelsen af
et af Slægtens levende Elementer, maa der kræves en bestemt
Paavisning af Handlingens Nødvendighed, for at den etisk kan
anerkendes.

I den gamle Kirke, som ellers var saa streng overfor Selvmord,
gjorde man dog en Undtagelse med de Kvinder, der
dræbte sig selv for ikke at blive voldtagne af Forfølgerne; de
æredes som Helgener\footnote[1]{Smlgn. Barbeyrac: Traité de la morale des péres de Véglise. Amsterdam 1728. p. 242 ff. --- Kun Augustinus var noget betænkelig ved Sagen; men han hjalp sig med en Antydning af, at hine Kvinder statistiken i „Tilskueren“ 1884. p. 466 ff.}. Aldeles afset fra de overspændte
% --- p. 219 ---
\setcounter{footnote}{0}%
\opage{219}asketiske Forudsætninger, hvorfra Kirken gik ud, kunne vi her
finde et Exempel, hvor Selvmordet er en etisk Handling. Selvmordet
slaar her fast, at der gives en Mishandling af Mennesker,
som i Brutalitet staar lige med Mord, og det bidrager
derved mægtigt til at skærpe og højne Forestillingerne om den
kvindelige Renheds Betydning.

Hvor hyppige Selvmord af ædle og højhjertede Bevæggrunde
ere, er vanskeligt at afgøre, som det idethele er forbundet med
store Vanskeligheder at faa Oplysninger om de Motiver, som i
de enkelte Tilfælde føre til Selvmord\footnote[1]{Der Selbstmord. Leipzig 1881. p. 270 ff. Smlgn. ogsaa Oettingen: Moralstatistik. 3. Aufl. p. 780 ff.}. Morselli\footnotemark[1] bestrider
den af nogle Forfattere opstillede Paastand, at det hos Mennesker,
der ikke lide af Sindssygdom, for det Meste er ædle Motiver,
som bevirke Selvmord. I vore Dage i hvert Tilfælde er,
mener han, Selvmord væsentligt Udslag af Egoisme. „Dog mangler
der“, tilføjer han, „ikke heller Udslag af den bedre Del af
vor Natur, og det især hos Kvindekønnet\dots{} Hos Mændene
fremtræder den egne Interesse som aldeles herskende Motiv,
og da kun en Fjerdedel eller en Femtedel af Selvmordene begaas
af Kvinder, er ædle Motivers Sjældenhed ved Selvmord
allerede sikret ved dette Talforhold“. Det maa overlades Statistikerne
at bedømme denne Antagelse.

7. Selvmordet kan kun betragtes fra et medicinsk, et psykologisk
og et etisk Standpunkt. Det kan ikke med Rette drages
ind under juridisk Betragtning. Af Staten er Individet kun
Medlem, saalænge det lever, og Statens Straf kan ikke ramme
Selvmorderen selv. Straffen kommer her i Virkeligheden til at
gaa ud over Andre end Individet selv, hvad enten den bestaar
i Inddragelse af Arven (som i den romerske Kejsertid, naar
Motivet var, at man vilde unddrage sig Straf for Majestætsforbrydelse)
eller i foragtelig Behandling af Liget (som det fra
Middelalderen af har været almindeligt i alle Lande). --- Endnu
kunde have handlet i Følge en umiddelbar guddommelig Aabenbaring
(De civitate dei I, 26).

 Smlgn. Marcus Rubins kritiske Bemærkninger om 
% CHECK p. 219: note mark without a note
Selvmords%
% --- p. 220 ---
\setcounter{footnote}{0}%
\opage{220}Biskop Martensen\footnote[1]{lnd.ividu.el Etik p. 461.} fandt i, hvad man kalder „den Slaphed,
Staten lægger for Dagen ved Selvmorderes Begravelse“, et Tegn
paa, i hvilken Grad „Kristendommens Aand er tilbagetrængt“.

% CHECK p. 220: centred line: heading or verse
\begin{center}b. Selvbeherskelse.\end{center}

8. Medens Selvopholdelsen især beror paa Overvindelse af
ydre Modstand, bestaar Selvbeherskelsen i Overvindelse af Modstand,
som stammer fra Menneskets egen indre Natur. Selvbeherskelse
forudsætter, at der er givet to eller flere forskellige
Tendenser i Mennesket, og at den ene af dem maa underordnes
eller fortrænges af den anden. Hvis der i det menneskelige
Selv ikke kunde gives flere indbyrdes modstaaende Drifter
eller Lidenskaber, vilde Selvbeherskelse være umulig. Det er
i Selvbeherskelsen Følelse, der viger for Følelse, Drift, der viger
for Drift, Lidenskab, der viger for Lidenskab. Man bliver ikke
Herre over sig selv ved sin „Fornuft“ alene, men ved at samle
Bevidsthedens Energi i en stærk Følelse eller Lidenskab, som
er rettet mod, hvad der staar for os som det Højeste. Selvbeherskelse
er i og for sig en rent formel Dyd, hvis Værdi beror
paa, i hvilken Sags Tjeneste den virker. Den, der elsker
Penge over Alt, formaar ved denne Lidenskab at holde andre
Drifter og Tilskyndelser nede. Hos Andre er Erkendelsesglæden
og Forskerlidenskaben, hos atter Andre Elskovsfølelsen eller
Fædrelandskærligheden den raadende Kraft. Ad disse og mange
flere Veje udvikles Evnen til Selvbeherskelse. Historisk have
især de ved Forholdet til Autoriteter (i Familie, Stat og Kirke)
bestemte Følelser haft stor Indflydelse paa Udviklingen af Evnen
til Selvbeherskelse. Forholdet til en Myndighed, som Mennesket
føler langt overgaar hans egen Magt og Evne, vækker en Følelse
af Bundethed, retter alle Kræfter mod ét Maal, underordner
alle Hensyn under et eneste, og virker saaledes som koncentrerende
Magt i Bevidstheden. Under Autoriteternes Herredømme,
med Lydigheden som Hoveddyd og med Haab og Frygt
eller --- paa højere Trin --- begejstret Ærefrygt som herskende
Lidenskaber har Menneskeheden gennemgaaet en Selvbeherskel%
% --- p. 221 ---
\setcounter{footnote}{0}%
\opage{221}sens Skole og vundet en Øvelse, som ikke behøver at gaa tabt
fordi de oprindelige Motiver afløses af andre.

Kun gennem Selvbeherskelsen vindes den indre Frihed og
Enhed, der gør fuld og samlet Virksomhed mulig. Det skal
have været Kynikeren Antisthenes, der først indskærpede Selvbeherskelsen
som en egen Dyd, idet han lærte, at den, der
ikke behersker sig selv, er Slave, hvad enten han udvortes er
fri eller ikke. Det var et etisk Synspunkt, ud fra hvilket Forskellen
mellem den frie Mand og Slaven tabte sin Betydning\footnote[1]{Smlgn. Karl Joël: Der echte und der xenophontische Sokrates. II, 1. p. 39—47. --- Med Hensyn til Selvbeherskelsens Teori maa især mærkes flere ældre Tænkeres stærke Betoning af, at en Drift kun kan hæmmes ved en anden Drift, --- en Analogi til Inertiprincipet eller til Energisætningen paa det fysiske Omraade. Smlgn. Platon: Fajdon. Heises Oversættelse. p. 23---25. Baco: De augmentis scientiarum. VII, 3. Spinoza: Ethica IV, 7 (kfr. V, 20. Schol.).}.
--- Den indre Frihed er Betingelsen for den ydre Frihed og
Uafhængighed. De forskellige Drifter og Tilskyndelser i vor
Natur ere ligesaa mange Veje, ad hvilke den ydre Verden kan
blive Herre over os. Dersom den uden videre kan vække hvilkesomhelst
Tilskyndelser hos os i ethvert Øjeblik, spiller den paa
os som paa et Instrument, hvilken Melodi den vil, og vi staa
afmægtige overfor den. Det gælder da om, at vi have et Centrum
i os selv, en Kreds af Tanker og Følelser, der stedse bestemme
Hovedretningen i vort Liv. Selvbeherskelsen skal ikke
gøre os følelsesløse, men vore Følelser skulle ikke svaje hid og
did for enhver Vind. --- Det behøver ingen nærmere Paavisning,
hvor stor en Betydning Selvbeherskelsen har for Samvittighedens
Udvikling og for Selvopholdelsen. I det Foregaaende
have vi derfor paa mange Punkter allerede forudsat den.

9. Det aandelige Liv udvikler sig sporadisk, begynder ud
fra forskellige Anlæg og Udgangspunkter, og først efterhaanden
kan der opnaas Harmoni mellem de forskellige Tendenser, som
røre sig. Forskellige Kræfter og Drifter bevæge Bevidstheden i
forskellige, ofte modsatte Retninger. Først gennem en Gærings- og
Kampperiode lykkes det ofte at bringe Orden og Enhed tilveje.
Det er en Betingelse for Livets højere Udvikling, at der
% --- p. 222 ---
\setcounter{footnote}{0}%
\opage{222}er givet et rigt Stof at bearbejde. Men det rige Stof kan være
genstridigt og vanskeligt at forene. Askesen overhuggede Knuden
ved at kue og tilbagetrænge de naturlige Drifter. Den stillede
sig mistænksom ikke blot overfor de sanselige Instinkter,
men ogsaa overfor Erkendelsestrangen og de æstetiske Fornødenheder.
Den store Kunst er derimod den at give de umiddelbare
og uvilkaarlige Tilskyndelser det frie Løb, som er en Betingelse
for et sundt og frisk Liv, og dog være Herre over
Løbets Retning, saa at denne ikke strider mod, hvad man betragter
som sit egentlige Maal. Død og mekanisk Orden og
selvplagerisk Ængstelighed ere Tegn paa Ufuldkommenhed. Det
er desuden ikke altid sundt at drage de halvt bevidste Tilskyndelser
og Fantasier helt frem i Bevidsthedens Dagslys; man
bliver lettere af med dem igen, naar de blive ude i Bevidsthedens
Udkant. Megen Moraliseren henleder netop Opmærksomheden
paa det, som skal fortrænges.

Selvbeherskelsen antager en forskellig Karakter, ikke blot
efter Beskaffenheden af det, som den tjener, men ogsaa efter
det, den fortrænger. Det kan være Nydelse, Fordel, Ære, Hævn,
som er Genstand for den opblussende Affekt, det gælder om at
modarbejde. Enhver Affekt eller Sindsbevægelse hindrer klar
og rolig Overvejelse og tillader ikke Forestillingerne at forbinde
sig paa den Maade, der ellers vilde være naturlig for dem. Særlige
Dispositioner hos det enkelte Individ kunne gøre Kampen
haardest for det i den ene eller den anden Retning. Der kan
være medfødte Anlæg til Mismod, Vredagtighed, Forfængelighed
og Sanselighed. Det Arbejde, den Enkelte her maa gøre, det,
han maa lide for at kunne opnaa indre Frihed og Harmoni, vil
derfor være højst forskelligt og kan ikke foregaa efter nogen
bestemt Recept. Fordringen om Selvbeherskelse lyder ens til
Alle, men betyder for Enhver noget højst Forskelligt baade i
Grad og i Art. Ligesom en og samme Byrde tynger forskelligt
paa forskellige Skuldre, saaledes udkræves der meget forskelligt
Arbejde af de forskellige Individer, for at de kunne holde sig
paa samme Udviklingstrin. Mangt et Menneske maa kæmpe
haardt blot for ikke at synke ned under det gængse Jevnmaal
af etisk Udvikling. For den ydre Betragter bliver en Kamp af
% --- p. 223 ---
\setcounter{footnote}{0}%
\opage{223}denne Art ikke let synlig eller beundringsværdig. Men Leonidas
var alligevel ikke mindre Helt end Alexander den Store,
uagtet denne kunde bruge sin Tapperhed til at erobre Verden,
medens hin maatte kæmpe til Døden under forgæves Forsøg
paa at holde Fjenden borte fra sit Land.

Allerede Aristoteles har i sin berømte Definition af Dyden
som den rette Midte, hvormed han mener et harmonisk
Forhold mellem de menneskelige Drifter, klart anerkendt disse
individuelle Forskelligheder i Selvbeherskelsens Art og Grad.
Ti den rette Midte ligger ifølge Aristoteles ikke paa samme
Punkt hos alle Individer\footnote[1]{Eth. Nic. II, 5---9}. Da deres Tilbøjelighed til de forskellige
Affekter er forskellig, kan det, der hos et Individ er
Tegn paa stor Selvbeherskelse, hos et andet Individ være Noget,
der sker aldeles uvilkaarligt. Hvad der hos En er Tegn paa
Gerrighed, kan være Sparsommelighed hos en Anden og Ødselhed
hos en Tredie. Dog taler Aristoteles mere som Psykolog
end som Etiker; han ser i hvert Tilfælde ikke sin Læres etiske
Konsekvenser. Han har fuldstændig Ret i, at Betingelserne for
at naa den indre Harmoni ere højst forskellige hos de forskellige
Individer som Følge af deres forskellige Natur. Men han
overser, at denne indre Harmoni ikke altid udtømmer det, som
fordres. Selvbeherskelsens Betydning beror jo paa, at vi skulle
virke med samlet Kraft i Arbejdet for vore Opgaver. Men hvis
nu min Opgave fordrer en større Grad af Selvbeherskelse, end
der stemmer med „den rette Midte“, den indre Harmoni? Der
kan være Ligevægt i min Natur, uden at min Natur derfor
fyldestgør de Krav, Forholdene stille. Jeg kan maaske bekæmpe
min Vredagtighed i en saadan Grad, at jeg viser større Selvbeherskelse
end den, der af Naturen er blid og fredsommelig;
og dog er det ikke sagt, at jeg er sikret imod at rives med af
den opblussende Affekt paa en forkastelig Maade. Der gør sig
her hos Aristoteles den tidligere (III, 5) omtalte individualistiske
Tendens gældende, som staar i en mærkelig Modsætning til
den store Vægt, den græske Etik ellers lægger paa Stat og
Samfund.

% --- p. 224 ---
\setcounter{footnote}{0}%
\opage{224}Baade Selvbeherskelsen og Selvopholdelsen tage sig anderledes
ud fra et rent individualistisk Standpunkt end fra den
humane Etiks Standpunkt (III, 10). Individualisten maa ganske
vist stræbe efter Orden og Harmoni mellem de forskellige Sider
i hans Væsen (III, 6), men fra den sociale Etiks Standpunkt
stilles der Fordringer til Selvbeherskelsen, som Individualisten
ikke vilde have nogen Grund til at stille. Dersom han blot bevarer
Ro og Harmoni i sit Indre, hvorfor skulde han saa ikke
følge sin Erhvervsdrift, sin Hævndrift? Han vender fra Forholdet
til Andre stedse tilbage til sig selv og har ingen Anledning
til at forfølge sine Handlingers Virkninger, naar disse ikke
gaa ud over ham selv.

Imidlertid ligger der en blivende Sandhed i den aristoteliske
Lære. Det er uimodsigeligt, at en og samme etiske Fordring
kan nødvendiggøre et højst forskelligt indre Arbejde hos
forskellige Individer. Det Problem, som Aristoteles gik let hen
over, fordi han ikke bragte den individuelle Etik i Forbindelse
med den sociale, bevarer sin Betydning ogsaa fra en social Etiks
Standpunkt, --- saafremt dette Standpunkt da virkeligt skal være
etisk og ikke falde sammen med det juridiske eller med den
offentlige Menings Standpunkt. Naar der er et indre Forhold
mellem Individ og Samfund, saa kan det ikke være ligegyldigt,
hvilke Forudsætninger det enkelte Individ møder med i etisk
Henseende. Kun en Dogmatisme, der vilde paastaa, at alle Mennesker
ere ens udstyrede med Evne til Selvbeherskelse, vilde
kunne hævde, at den rette etiske Lov var funden, naar der
fordredes samme ydre Kvantum Selvbeherskelse af Alle. (Smlgn.
IV, 2 og VIII, 6). Hellere end at indrømme dette, der vilde
føre til Barbari og Farisæisme, maa Etiken anerkende sin Grænse
i videnskabelig Henseende. Overalt hvor de individuelle Forskelligheder
træde i Kraft paa afgørende Maade, er der en
Grænse for Videnskaben. Hvad der kan fordres af Enhver, af
hvem der overhovedet (af ham selv eller Andre) kan fordres
Selvbeherskelse, er, at hans Handling ligger i en Række, der
ved Fortsættelse vilde føre til fuldkomment Herredømme over
sig selv; men paa hvilket Punkt i Rækken det enkelte Individs
Handling i det enkelte Tilfælde skal ligge, det er Noget, den
% --- p. 225 ---
\setcounter{footnote}{0}%
\opage{225}etiske Tænkning med de Midler, der staa til dens Raadighed,
ikke formaar at bestemme saaledes, at baade individuelle og
sociale Hensyn komme til deres Ret. Og det af den simple
Grund, at de forskellige Individer ikke begynde ud fra samme
Punkt.

10. Forskellen mellem det Standpunkt, hvorpaa vi her have
stillet os, og det individualistiske ses ogsaa af den Maade, hvorpaa
de forskellige Arter af Selvbeherskelse vurderes. Vi betragte
ikke det enkelte Individ som en lille afsluttet Verden for sig,
men stedse i dets Forhold til Andre. Spørges der da: er det
vigtigst at beherske sin Hævnlyst eller sin Havesyge eller sin
sanselige Nydelsesdrift? --- kan der naturligvis gives et Svar
herpaa efter den Maade, hvorpaa disse Drifter gribe ind i Individets
eget Indre; men dette Svar vil ikke være tilstrækkeligt.
Sanselighed og Nydelsessyge gøre Mennesket afhængigt af det
Ydre og hæmme derved hans frie Udvikling. De hindre ogsaa
de ideelle Følelsers Udvikling og svække Villiens Energi, komme
maaske endog i Strid med, hvad fysisk Selvopholdelse kræver.
Af størst Betydning i etisk Henseende er det, at Hengivelse til
sanselige Nydelser indsnevrer Sindet. Sanselig Nydelse kan ofte
ikke deles med Andre, ja kan maaske endog kun naas paa Andres
Bekostning. Det er maaske først en af Affekterne fremkaldt
Tankeløshed, som lader det nydende Individ se bort fra
de Følger, dets Nydelser kunne have for Andre; men denne
Tankeløshed slaar ved Vane over til Hjerteløshed. Udsvævende
Mennesker ere ofte mere haardhjertede og kolde end Mennesker,
der forse sig paa Andres Gods eller lade sig rive hen af opblussende
Vrede. Etisk og juridisk Betragtning fører her til
forskelligt Resultat. Den, der stjæler for sine sultende Børns
Skyld, staar etisk set langt over den, der i Udsvævelser tilintetgør
den materielle eller aandelige Kapital, som kunde være
kommen Mange til Gode. Den positive Moralitet og især den
offentlige Mening nærmer sig her, som saa ofte, for meget den
blot juridiske Betragtningsmaade. Tyven er udelukket fra det
«gode Selskab“, men den udsvævende Egoist indtager maaske
Hæderspladsen.

Hvad særligt Kønsforholdet angaar, kan det først indenfor
% --- p. 226 ---
\setcounter{footnote}{0}%
\opage{226}den sociale Etik finde sin fulde Bedømmelse. Indenfor den individuelle
Etik er den vigtigste Regel den, at for den Rene er
Alt rent. Det Urene fremkommer, naar Bevidstheden er hildet
i sanselige Forestillinger, ikke kan frigøre sig for dem. Skyhed
og Snerperi kunne være betænkelige Tegn, fordi de vidne om,
at slige Forestillinger saare let indfinde sig hos Individet. Ikke
ved at pleje en mystisk Sky for Kønslivets Fænomener, men
ved at betragte dem som alle andre Naturfænomener med uhildet
Blik og klar Forstaaelse, arbejder man paa en sund og naturlig
Renhed i Sindet. Der gives kun to sunde Maader at staa
overfor Kønsforholdet paa, dels paa den antydede Maade som
fysiologisk og psykologisk Iagttager, dels som den, der erfarer
dets Magt gennem inderlig Hengivelse, altsaa gennem mere end
blot sanselig Nydelse. Forskellig fra begge Dele er Leflen, Raahed
og uterlig Fantaseren.

Forholdet til Kønslivet stiller et af de allervanskeligste etiske
Problemer. Natur og Kultur støde her sammen paa den skarpeste
Maade. Som Naturvæsen er Mennesket moden Mand eller Kvinde,
længe førend dets Existens som Kulturvæsen endnu er befæstet
saaledes, at det kan sørge for sit Afkom. Der rejses her stadigt
stærke Fordringer fra Naturens Side, som ikke kunne tilfredsstilles
i det Omfang, i hvilket de trænge frem, uden at Kultur- og
Samfundslivet sprænges. Her, som saa ofte, rører Naturtrangen
sig med Ubændighed og Ødselhed, ret som om Naturen
for enhver Pris vilde sikre Livets Bestaaen paa Jorden, hvorledes
det saa gaar med det enkelte Individ. Det er Naturens
sorgløse Udstrøen af Livsspirer, som her viser sig fra den
grelleste Side og stiller de Væsener, hos hvem Bevidsthedslivet
er vaagnet, overfor de største Opgaver. Naturen bruger den
individuelle Organisme som Middel; men naar et personligt Liv
har udviklet sig, vil det ikke være blot Middel. Der staar en
Kamp mellem de Instinkter, vi have fælles med Dyrene, og de
Interesser, som føre til en højere Udvikling af det menneskelige
Liv. \textsc{Kant} drog endog heraf den Slutning, at Mennesket af
Naturen ikke var anlagt til det civiliserede Liv, men at Naturens
Hensigt kun var Bevarelsen af Menneskeslægten som Dyreart.
Han fandt her en dyb Disharmoni og en Kilde til Elendighed,
% --- p. 227 ---
\setcounter{footnote}{0}%
\opage{227}som kun kan overvindes, naar Menneskene naa et Kulturtrin,
fra hvilket de endnu ere langt fjernede\footnote[1]{Kant: Muthmasslicher Anfang der Menschengeschichte. (1786) (Vermischte Schriften. Halle 1799. III, p. 48 f.). — Smlgn. hermed Elie Metchnikoff: Études sur la nature humaine. Paris 1903 chap. 5 (Désharmonies dans l'organisation et le functionnement de l'appareil de la reproduction).}. Den store Tænker
har her paa én Gang peget paa den store Vanskelighed og paa
den eneste Maade, paa hvilken man etisk set kan stille sig til
den. Det er let nok at overhugge Knuden og uden videre
fordre ubunden Tilfredsstillelse af de naturlige Instinkter. En
saadan er netop kun mulig, hvis vi igen ville gaa tilbage til
det dyriske Trin. Og selv paa det dyriske Trin hersker der ikke
fuldstændig Ubundethed, hvad Parringsvalget og Parringskampen
noksom vise. Selvom Askesen i sit Naturhad og den officielle
Moral i sit Snerperi ender i Usandhed og Unatur, gaar Reaktionen
mod dem paa sin Vis over Stregen, naar den vil hævde,
at Kønsforholdet er neutralt eller i hvert Tilfælde af underordnet
Betydning i etisk Henseende. Det er i Kraft af et rigtigt Instinkt,
at man til alle Tider har betragtet Kønsforholdet og dets
Ordning som et af de allervigtigste etiske Spørgsmaal. Her have
vi endnu kun at gøre med det fra det enkelte (skønt ikke suveræne)
Individs Standpunkt og kunne derfor ikke give nogen
fuldstændig Behandling af dette Spørgsmaal; en saadan kan først
Læren om Familien give. Men saa Meget er klart, at da Kønsinstinktet
kun kan tilfredsstilles ved Hjælp af et andet Individ,
kan der her ikke være Tale om nogen ubegrænset individuel
Frihed. Dersom Instinktets Tilfredsstillelse medfører Ulykke og
Nedværdigelse for et andet Individ, eller endog forudsætter Tilværelsen
af en hel Klasse af Mennesker, der leve af at være
Midler for Andres sanselige Lyst, saa behøves der ingen nærmere
Paavisning af, at vi her ingenlunde staa overfor noget
Ligegyldigt, en ren Privatsag. Kønsinstinktet skjuler i sig Spiren
baade til det Højeste og Ædleste og til det Laveste og
Nederdrægtigste i den menneskelige Natur. Det kan føre Individet
ud over sig selv i Begejstring og Glæde over et andet
% --- p. 228 ---
\setcounter{footnote}{0}%
\opage{228}Menneske og derved vise sig beslægtet med al Hengivelse og
Opofrelse; men det kan ogsaa føre til, at Individet gør Andre
til blotte Midler for sin Nyden, hildet som det er i sine Tilskyndelser.
Mellem disse Poler svinger Instinktet, og der hører
en mærkelig Blindhed til for at mene, at Selvbeherskelse her
ikke skulde have nogen Betydning. Selvbeherskelse viser sig
netop her at være Menneskets Adelsmærke, Betingelsen for, at
det kan fortsætte den Udviklingsgang, ved hvilken det har hævet
sig op af en dyrisk Existens.

Selvbeherskelse er ogsaa mulig i langt højere Grad, end
man ofte mener. Man taler her ofte om fysiologiske Nødvendigheder,
som ikke ere tilstede, men som man ganske vist kan
fremelske ved ideligt at tale om dem. Der er Individer, som
uden Skade for deres aandelige og fysiske Livs Sundhed og
Kraft bevare deres Renhed. Hvor langt det enkelte Individs
Evne til Selvbeherskelse gaar, kan kun det selv prøve; men
ligesaa lidt som Nogen har Ret til at kaste den første Sten,
ligesaa lidt har Nogen Ret til at kalde sin egen Svaghed for
en almenmenneskelig Nødvendighed\footnote[1]{Der føreligger ingen pædagogisk Grund til at sænke Niveauet. Esquirol bemærker: „Si la continence, dans quelques cas trés-rares, a causé l'aliénation mentale, le libertinage est une cause plus fréquente''. Des Maladies mentales. Bruxelles 1838. I. p. 24. --- I meget stærke Udtryk angriber den engelske Læge Lionel Beale (Our Morality and the Moral Question, chiefly from the medical side. London 1887) dem, der paastaa hin fysiologiske Nødvendighed. Han anfører (p. 99) en Udtalelse i samme Retning af Sir James Paget. --- Smlgn. ogsaa S. Ribbing: Om den sexuela hygienen (Stockholm 1888). (Forelæsninger for svenske Studenter). --- Albert Heim: Das Geschlechtsleben des Menschen vom Standpunkte der natürlichen Entwickelungsgeschichte. Zürich 1900 (Et Foredrag for Schweizerstudenter). --- Erik Pontoppidan: Kønssygdommenes Betydning og Forebyggelse. København 1903. (Forelæsninger for danske Studenter). I dette sidste Værk fremhæves særlig (p. 65 f.; 98), hvorledes etiske og sanitære Hensyn her falde sammen: „I Spidsen for alle gode Raad til Ungdommen angaaende disse Sygdomme og deres Forebyggelse sættes en Opfordring om Personlighedens Opdragelse og Indordnen under faste sædelige Principer, for at styrke Karakteren og sætte en mandig Villie mod de kønslige Fristelser.“}. Den, som maa kæmpe
haardt for at holde sig oven Vande, kan gøre det med den Be%
% --- p. 229 ---
\setcounter{footnote}{0}%
\opage{229}vidsthed, at han tager sin Del af den Lidelse, Slægten maa gaa
igennem for at naa et højere Trin, og at han yder sin Skærv
til Fremme for Udviklingen derhen imod. Vi ere Alle Led af
den store Udviklingsrække, der fører fra Dyret til Mennesket.
Og i ingen Henseende er det maaske saa vigtigt som her, at
Trangen hindres i at synke tilbage til den rent dyriske Form.
Den vældige Trang, som Talen her er om, ytrer sig i Regelen
hos Dyret blindt, isoleret og ligegyldigt overfor sin Genstand,
skønt der i Dyreverdenen findes Træk, der vidne om, at en
højere Form allerede her er mulig. Humaniseret bliver Trangen,
jo mindre den skilles fra andre Elementer i Personligheden,
jo mere den forbindes med og bestemmes ved sin Genstands
særegne, individuelle Egenskaber, og jo mere den bliver bærende
Grundlag for sympatiske Følelser og ideelle Interesser\footnotemark[1].

Ligesaa lidt her som overfor de Motiver, der føre til Selvmord,
kan Kampen altid føres direkte. Selvbeherskelsen kan
ikke altid øves i det enkelte, yderste Øjeblik. Ved Legemsøvelser,
ved sund og kraftig Næring for Fantasien og ved
Entusiasme for ideelle Formaal hindres Instinktet fra at brede
sig og fremkalde de fixe Ideer og den Svimmelhed, som fører
til Fald.

11. Medens det i den positive Moralitet ikke mangler paa
farisæisk Holdning overfor de Forvildelser, der fremkomme ved
yppige Instinkters og Tilskyndelsers Udfoldelse, har den ikke
saa meget Blik for, at Selvbeherskelse kan blive af meget stor

:) I en interessant Diskussion om „écheation sexuelle“ i det franske
filosofiske Selskab gjorde baade Modsætningen mellem den biologiske og
den sociologiske Betragtning af Kønsforholdet og Modsætningen mellem
den teologiske og den filosofiske Betragtning af dette Forhold sig gældende.
Overfor den rent biologiske Betragtning gjordes fra sociologisk
Side gældende, at Kønsakten maa ses i dens Sammenhæng med Familielivet,
og at den kun staar som ret menneskelig, naar den er Udtryk for
inderligt og stadigt Fællesskab. Overfor den teologiske Betragtning gjordes
fra filosofisk Side gældende, at religiøse Motiver ikke ere nødvendige
for en etisk Holdning overfor Kønsforholdet, men at Selvbeherskelse her
motiveres dels ved menneskelig Værdighed, dels ved Retfærdighedsfølelse.
(Bulletin de la Société Française de Philosophie. 1911). (1913).
% CHECK p. 229: note mark without a note

% --- p. 230 ---
\setcounter{footnote}{0}%
\opage{230}Betydning ogsaa under helt andre Forhold. Er der i den nys
omtalte Periode Fare for at løbe vild, kan der i en senere Periode
være Fare for at løbe fast. Og Mange, som løbe vild i
Ungdommen, løbe sig senere saa meget desmere fast. Mange
unge Himmelstormere ende som Filistre; Manges Liv deler sig
i en Udsvævelsernes og en Blaserethedens Periode. Den, der
forbliver sig selv tro gennem de Kampe, de urolige Tilskyndelser
føre ind i, vil ogsaa lettere bevares for at stagnere.

Vanens og regelmæssige Livsforholds Indflydelse stiller ofte
ikke ringe Krav til Selvbeherskelsen, for at den aandelige Frihed
og Friskhed skal kunne bevares og en ny Væxt stedse
blive mulig. Det Nye maa finde os nye; men Forudsætningen
derfor er en stadig Bestræbelse for at holde Blikket aabent og
klart. Vor Væxt kan ganske vist ikke gaa i det Uendelige. Der
er sørget for, at Træerne ikke voxe op i Himlen. Men det gælder
om at faa alle værdifulde Væxtmuligheder udviklede. Der
er intet smukkere Syn end en kraftig og livlig Udvikling i en
Alder, hvor al Udvikling ellers plejer at være afsluttet. Stadigt
aandeligt Arbejde og varm Sympati for alt Stort og Skønt er
det bedste Middel til at undgaa Stivnen\footnote[1]{Smlgn. min Karakteristik af F. C. Sibberns Personlighed og Udviklingsgang i Mindre Arbejder.}. Kejser \textsc{Marcus}
Aurelius stiller i sine Optegnelser (VI, 30) følgende Opfordring til
sig selv: „Se til, at du ikke forkejsres (ὅρα μὴ ἀποκαισαρωθῇς)!
Sørg for, at du ikke faar Farve efter din Stilling! Det sker jo
saa let. Bevar dig derfor ligefrem, god, ren, alvorlig, uden Prunk,
retfærdig, from, fuld af Velvillie og Kærlighed, standhaftig i at
opfylde dine Pligter. Kæmp for, at du kan vedblive at være et
saadant Menneske, som din Stræben efter Visdom vilde gøre
dig til!“ --- Den Fare for at blive en Filister, som den filosofiske
Kejser frygtede, truer enhver, hvad enten han er lavt
eller højt paa Straa, og Midlet imod den er overalt det, som
han anbefalede: at være tro imod sig selv og imod sin bedste
Stræben.

% --- p. 231 ---
\setcounter{footnote}{0}%
% CHECK p. 231: centred line: heading or verse
\opage{231}\begin{center}c. Selvstændighed.\end{center}

12. Baade Selvopholdelsen og Selvbeherskelsen ere Betingelser
for den personlige Selvstændighed. Det enkelte Individ
skal hævde sig udad til og indad til for at komme til at staa
som et ejendommeligt og selvstændigt Led af Slægten. Det har
Noget i sig, som ganske paa samme Vis ikke kommer igen hos
noget andet Individ. Trangen til selvstændig Udvikling, til at
gøre sine Anlæg og Evner gældende er en af de vigtigste Kræfter,
som føre til det menneskelige Livs Fremgang. Uden Tro
paa sig selv og Tillid til, at der er Noget i Ens Selv, som fortjener
at leve og udfolde sig, kan Intet udrettes. Mismod og
Mistillid til sig selv lammer al Villen og alt Arbejde, ja er ofte
Tegn paa en Villiessvækkelse, der kan nærme sig til Idioti.
Allerede Aristoteles har indskærpet, at den, der har for liden
Selvfølelse, ikke kommer til at udrette alt det Skønne og Gode,
han ellers kunde, fordi han ikke anser sig værdig dertil. En
saadan „Mikropsyki“ (som Aristoteles kalder den) hindrer Udfoldelsen
af, hvad der ligger i Individets Natur. Den sande Selvfølelse
(Makropsykien) er ganske vist --- som Aristoteles ogsaa
viser --- vanskelig og sjelden; den forudsætter, at man ikke
blot agter sig selv værdig til noget Stort, men ogsaa virkeligt
er værdig dertil. Den, som besidder denne Følelse, har en fast
Maalestok i sig selv. Han agter ydre Ære og ydre Goder kun
som noget Underordnet. Han kender sine Grænser, søger at
udrette Alt, hvad han formaar indenfor dem, taaler intet Overgreb
fra Andres Side, men begaar heller ikke selv noget Overgreb.
Han gaar tryg og frejdig gennem Livet\footnote[1]{Aristoteles's Lære om Makropsyki og Mikropsyki findes i Eth. Nic. IV, 5—9. --- Smlgn. fra en nyere Tid: Hume: Treatise III, 3, 2 (Of greatness of mind». og Adam Smith: Theory of moral sentiments. VI, 3 (Slutning).}.

Ydmyghed og Beskedenhed ere ikke saa store Dyder, som
de ofte ere udgivne for at være. Den store Ros, som er bleven
dem til Del, kan kun forklares dels som Reaktion imod den
overdrevne og hensynsløse Selvfølelse, som er saa hyppig, dels
% --- p. 232 ---
\setcounter{footnote}{0}%
% CHECK p. 232: notes paired with the marks in order (the layer misread a number)
\opage{232}som Eftervirkning af den asketiske Tankegang, for hvem Lydighed
er den højeste Dyd. Hvad det gælder om, er at kende sig
selv og saa stræbe at faa Alt frem, hvad man kan gøre. Naar
Idealet blot sættes tilstrækkeligt højt, behøves der ingen særlig
Lovprisning af Ydmyghed. Vi ville da nok faa vore Skranker
at mærke og behøve ikke at svække vort Mod ved at dvæle
for meget ved vor Afmagt\footnote[1]{Spinoza definerer Ydmyghed (humilitas) som „den Sorg, der opstaar ved, at Mennesket betænker sin Afmagt eller Svaghed“, en Følelse, der let lammer og fører til Foragt for sig selv (abjectio). Eth. III. Aff. Def. 26—29. --- Aristoteles vil heller ikke betragte Undseelighed (αἴδως) som en egentlig Dyd. Eth. Nic. IV, 15.} i Stedet for at samle og anvende
den Magt, vi virkeligt have. Den, der kender sine Grænser,
har paa én Gang Tro paa sin Ejendommelighed og Blik for sin
Lidenhed\footnote[2]{Smlgn. Goethes Selvvurdering (i en Samtale med Eckermann. 30. Marts 1824). Man havde villet sidestille Tieck med ham som Digter. Dertil siger Goethe, at det vilde være det Samme, som hvis han selv vilde stille sig ved Siden af Shakespeare.}. Han føler sig paa én Gang stor og lille. Den sande
Selvfølelse faar derved et Præg af Humor.

13. Man hævder ikke sin Selvstændighed ved ængsteligt
at holde al fremmed Indflydelse borte. Jo mere Stof et Menneske
kan behandle og tilegne sig, desmere vil han kunne
lægge sin Ejendommelighed for Dagen. Den sande Selvstændighed
viser sig netop ved at kunne fordybe sig i et rigt Indhold.
De mest originale Mennesker ere gerne mest villige til at anerkende,
hvad de have lært af Andre. \textsc{Marcus} \textsc{Aurelius} begynder
sit Skrift med at opregne alle dem, som have havt Indflydelse
paa hans Udvikling. \textsc{Goethe}, der som Faa havde Ret til at
føle sin Selvstændighed, siger (i en Samtale med Eckermann):
„Man taler om Originalitet; men hvad vil det sige! Saasnart
vi blive fødte, begynder Verden at virke paa os, og det gaar
saaledes lige til Enden. Og hvad kunne vi overhovedet kalde
vort Eget uden Energien, Kraften, Villien! Hvis jeg kunde sige
Alt, hvad jeg skylder store Forgængere og Medlevende, saa
vilde der ikke blive Meget tilbage“.

Ofte mener Individet kun at kunne hævde sin Selvstændig%
% --- p. 233 ---
\setcounter{footnote}{0}%
\opage{233}hed ved at trække sig tilbage til Følelsen og beraabe sig paa
den. Følelsen har noget mere Individuelt ved sig end Tanke
og Handling. Men naar Individualitet ikke skal være det Samme
som Isolation, og naar den Ensidighed og Begrænsning, som
hvert enkelt Individ altid lider af, skal kunne suppleres, saa
maa det ikke blive ved den mystiske Beraabelse paa Følelsen,
men Ejendommeligheden maa lægge sig for Dagen i Tankens
og Villiens Arbejde. --- Overspændt Følelse af den egne Personlighed
uden Forstandskraft som Modvægt er et hyppigt Karaktertræk
hos Sindssyge\footnote[1]{Maudsley: Pathologie de Vesprit. Trad. de l'anglais. p. 228 ff.}.

14. Ydre Selvstændighed er nødvendig for, at den indre
ret kan være tilstede. Afhængighed af Andre medfører Skranker
for den frie Bevægelse og Udfoldelse af Kræfterne. En
moden etisk Personlighed vil derfor stræbe efter ydre Frihed.
Vi se ogsaa, at Verdenshistorien egentligt er en stor Befrielseshistorie.
Hvad der kæmpes om i Historien, er ikke blot om
Midlerne for at leve, men ogsaa om Midlerne for at kunne leve,
som man vil. \textsc{Hegel} har træffende sagt: i Orienten var Én
fri, i Grækenland vare Nogle frie, i nyere Tid ere Alle frie.
Vi kaldes nu alle Herrer; hvad der først betegnede den Enes
Overhøjhed over den Anden, betegner nu den Enkeltes personlige
Selvstændighed. Det var en asketisk Tendens, der gjorde
sig gældende, naar baade Buddhismen, Stoicismen og den ældste
Kristendom hævdede Friheden som noget rent Indre og erklærede
det for ligegyldigt, om man i det Ydre var fri eller
Træl. Vel kan Slaven have sin indre Helligdom, hvor Ingen
kan trænge ind; men den fuldstændige Modsætning mellem
Indre og Ydre er ikke blot pinlig, men kan skade selve det
etiske Liv ved at hindre de indre Kræfter i at blive anvendte
udad til. Naar Historien fra én Side set er en stor Frigørelsesproces,
saa er den etiske Betydning deraf, at der skabes selvstændige,
personlige Udgangspunkter for Slægtens Liv. Slægten
lever kun i de enkelte Individer; jo kraftigere og friere disse
røre sig, des fyldigere og rigere vil ogsaa Slægtens Liv blive.
Den individuelle og sociale Etik mødes altsaa her i samme
% --- p. 234 ---
\setcounter{footnote}{0}%
\opage{234}Fordring. (Smlgn. VIII, 6). Grænsen for den Enkeltes personlige
Frihed maa bestemmes ved, hvad Andres lige Frihed
kræver, og det vil være Opgaven saavidt muligt at faa disse
Grænser til at falde sammen med Grænserne for Individets
Evne og Drift.

Ære og Ejendom kunne betragtes som ydre Forlængelser
af den personlige Frihed.

For at kunne naa vort Maal er det ikke nok, at vi have
Tro paa os selv; vi behøve ogsaa en vis Anerkendelse fra
Andres Side. Ogsaa Andre maa betragte os som et Væsen
med en berettiget Ejendommelighed. Æren kommer under sunde
Forhold af sig selv, som en Forudsætning, der ikke behøver at
tildrage sig særlig Opmærksomhed, men som dog giver vor
Optræden overfor Andre Tryghed og Fasthed. Det er usundt,
naar Opmærksomheden for meget henledes paa det Billede,
der aftegner sig af os i Andres Bevidsthed. Hvad der skal
være Betingelse for vor Virksomhed, bliver da let til Formaal
for den, og Skinnet kommer da til at raade.

Den personlige Udvikling behøver dernæst ogsaa materielle
Midler. Saalænge Individet lever fra Haanden i Munden, kan
ingen højere Udvikling komme i Stand; og saasnart det har
Mere end fra Haanden i Munden, har det ogsaa Ejendom. Her
i den individuelle Etik betragte vi Ejendommen kun som Middel
for personlig Selvstændighed, som en Forlængelse af Personligheden.
Ved enhver Ordning af Samfundet, der ikke gør
den Enkelte til en blot Maskine, et villieløst Redskab, maa han
have Raadighed over de Midler, han under visse Betingelser
kan skaffe til Veje. Hvorledes disse Betingelser nu end skulle
formuleres, og hvor stor en Kontrol man end vil have ført med
hans Raaden over Midlerne, maa der stedse blive et Omraade
tilbage, hvor han selv tager Bestemmelserne. Og dette Omraade
bør være saa stort, som det er foreneligt med Hensynet til
Andres Tarv. Enhver Begrænsning af Raadigheden er i sig
selv et Onde. Hvorfor skal jeg ikke bruge mine Kræfter til
at bearbejde det Stof, jeg forefinder, eller min Magt til at lægge
Beslag paa saa Meget, jeg formaar, naar jeg ikke skader Nogensomhelst
derved? I og for sig er det et Gode at kunne brede
% --- p. 235 ---
\setcounter{footnote}{0}%
\opage{235}sig i Tilværelsen; Ondet kommer først, naar det kommer til
Sammenstød mellem Flere, der alle have denne Trang til at
brede sig. Bevisbyrden paahviler altsaa den, der vil indskrænke
min Magtsfære, ligesom Bevisbyrden overhovedet (i Kraft af
Velfærdsprincipet) paahviler den, der vil have, at vi skulle føle
Smerte istedetfor Lyst. Denne Trang til Ejendom og Magtudfoldelse
er ikke nødvendigvis Egoisme. Det er jo ikke sagt,
at man gør Besiddelse og Magt til sine højeste Formaal. De
kunne bruges, ligesom Sundhed, fysisk og aandelig Kraft, i
Menneskehedens Tjeneste.

Paa alle Trin af menneskelig Existens finde vi denne Trang
til Ejendom, i hvert Tilfælde rørlig Ejendom. Det Ejendomsfællesskab,
især af Jord, som er saa almindeligt paa mere primitive
Udviklingstrin, forsvinder overalt ved den kraftigere
Udvikling af de enkelte Personligheder, som Civilisationen fører
med sig. Det viser sig uforeneligt med en selvstændigere og
friere Udvikling af det personlige Liv. Under Ejendomsfællesskabet,
som f. Ex. i de sydslaviske Zadruga'er, er Enhvers
Skæbne fastslaaet og kan ikke blive væsentligt forskellig fra de
Andres. Men der opstaar naturligt en Trang til at „leve efter
eget Tykke, arbejde for sig selv, drikke af sit eget Glas“, ---
til at „nyde det uafhængige Livs Fortryllelser og trodse dets
Farer“\footnote[1]{mærker: „Fællesskabet havde som alle saadanne gamle Institutioner haft sin naturlige Berettigelse\dots{}. Men i Tidens Løb mistede det sin Be-}. Det er meget vel muligt, at Ejendomsretten, dens
Begrundelse og dens Udstrækning vil og maa undergaa mangfoldige
Forandringer. Men hvilken Ordning der saa anvendes,
vil en væsentlig Del af den Værdi, som kan tilskrives den,
bero paa det Omfang, i hvilket den tilfredsstiller Selvstændighedstrangen
hos saa Mange som muligt. I Selvstændigheden
ligger en af de vigtigste Betingelser for, at Anvendelsen kan
blive god\footnotemark[1]. ---

*) Laveleye: Om Ejendomsretten og dens primitive Former. Dansk
Overs. p. 120. --- Om Ejendommens Historie smlgn. Spencer: Political
Institutions. Chap. 15. --- Cl. Wilkens: Sociologi p. 284—295.

\footnotemark[1] V. Fal be Hansen (Stavnsbaandsløsningen og Landboreformerne.
Set fra Nationaløkonomiens Standpunkt. I. Kbhvn 1888. p. 69 f., 73) 
% CHECK p. 235: note mark without a note
be%
% --- p. 236 ---
\setcounter{footnote}{0}%
% CHECK p. 236: indent after a broken word?
\opage{236}Paa lignende Maade som Ære og Ejendom maa de borgerlige
Rettigheder betragtes. Historien (især Roms og Englands)
viser os, af hvor stor Betydning det har været for et Folks, ja
for hele Menneskehedens Udvikling, at den Enkelte staar fast
paa sin gode Ret. „Fuldstændig Opgivelse af Retten“, siger
Jhering\footnote[1]{Kampen for Retten. Dansk Overs. p. 20. --- Smlgn. ogsaa Der Zweck im Recht. I, 2. Aufl. p. 74 f. 259. --- Steinthal (Allg. Ethik p. 154) bemærker: „I Tyskland, hvor Retsudviklingen i Aarhundreder har været saa aldeles hæmmet, ere Karaktererne sparsomt udsaaede og Retsfølelsen svag.“}, „er et moralsk Selvmord“.

\chapterhead{XII.}{HENGIVELSE.}

% CHECK p. 236: centred line: heading or verse
\begin{center}a. Kærlighed til andre Væsener.\end{center}

1. Kærlighed arbejder saa direkte som muligt for det, der
er al Etiks sidste Maal. Var der kun denne ene Følelse i Mennesket,
og var den ikke saa ofte blind, vilde hverken praktisk
eller teoretisk Etik være nødvendig eller mulig. Men med et
saa simpelt og fuldkomment Kompas ere Menneskene nu engang
ikke udstyrede.

Kærligheden forudsætter en Evne til at genkende et Følelsesliv
hos andre Væsener. Den strækker sig saa vidt, som det
rettigelse og blev i høj Grad skadeligt, navnlig fordi det hindrede Forandring
og Fremskridt. Den Enkelte maatte ved sin Jords Dyrkning følge
de Andre; han maatte ikke alene vedblive med samme Driftssystem, men
han kunde heller ikke i Enkelthederne uden Vanskelighed foretage større
Forandringer\dots{}. Udskiftningen brød gammel Skik og Slendrian; den
tvang Bonden, der tidligere tankeløs og villieløs havde fulgt med i Landsbyens
nedarvede Fællesdrift, til selvstændig, individuel Tænken og Virken.
Den gjorde det muligt for den Enkelte at gaa frem uden at hindres af
Andre.“
% --- p. 237 ---
\setcounter{footnote}{0}%
\opage{237}er os muligt at sætte os i Andres Sted, at føle og lide med
dem. Den almindelige Menneskekærlighed har langsomt udviklet
sig gennem successiv Udvidelse af Kredsen af de Væsener,
med hvilke man kunde sympatisere. Følelsens Omfang er
nu dog ikke uden Indflydelse paa dens Art og Styrke. I de
snevreste Familie- og Venskabsforhold optræder den anderledes
end i de løsere og fjernere Forhold. Ved Spørgsmaalet om
Kærlighedens rette Art og Styrke i de forskellige Forhold bliver
det klart, hvad tidligere er vist, at Kærligheden gaar over til
Retfærdighed. Retfærdigheden er den efter sit eget Princip ordnede
Kærlighed. Den almindelige Menneskekærlighed (den
uinteresserede og universelle Sympati) fører, naar den bliver
klar over sig selv, til den Fordring, at der, netop for hele
Slægtens Velfærds Skyld, skal vises den største Kærlighed i
de snevreste Kredse, hvor Menneskene kunne være og gøre
Mest for hverandre. Ved et umiddelbart og stadigt Livsfællesskab,
som i Familien eller i et nøje Venskabsforhold, aabnes
gensidig Adgang til saa dybt Indblik i Personlighederne, som
idethele er muligt. Enhedsfølelsen kan her blive saa stor, som
den idethele kan blive. Kærlighedens Styrke, hvad enten den
har Heftighedens eller Inderlighedens Karakter, kan ikke blive
saa stor i de vide Kredse som i de snevre. Der skelnes her
ikke mellem Mit og Dit, mellem Given og Modtagen, og dog
lægge netop Personlighedernes Ejendommeligheder sig for Dagen,
da de intetsteds kunne udfolde sig saa frit og frejdigt som i et
saadant Forhold. I saadanne snevre Forhold fødtes Sympatien
fra først af i Menneskeslægten, og fødes den stedse paany.
Selvom derfor den etiske Vurdering for at faa alle Hensyn
med stiller sig paa den almindelige Velfærds Standpunkt, saa
fordrer den dog netop fra dette Standpunkt de smaa Kredses
Bestaaen. I dem naas i mange Henseender det Højeste, som
kan naas. De blive Forbilleder for de større Kredse. Og ikke
blot det, men fra dem udstrømmer ogsaa den Kraft, som kan
besjæle de større Kredse.

Der kan ofte være et Misforhold mellem Kærlighedens
Styrke og Omfang. Den kan begrænse sig, saaledes at der opstaar
en Gruppeegoisme, naar Familiens, Standens, Nationens
% --- p. 238 ---
\setcounter{footnote}{0}%
\opage{238}eller Racens Interesser dyrkes paa Bekostning af store almenmenneskelige
Interesser. Ligesom den Enkeltes Selvopholdelsestrang
maa underordnes Hensynet til Alles Tarv, saaledes maa
ogsaa den begrænsede Sympati undertiden ofres for den mere
omfattende. Vore Pligter som Familiemedlemmer eller som Venner
staa ikke over vore Pligter som Patrioter og som Mennesker.
Men Misforholdet kan ogsaa være af omvendt Art. Der
kan gives en Opofrelse og Begejstring for store nationale og
almenmenneskelige Interesser, der fører til hensynsløs og brutal
Optræden i de mindre Kredse. „Man mangler“, siger Maudsleyl),
„ikke Exempler paa, at Martyrerne for Menneskehedens
Sag selv kunne frembringe Martyrer blandt dem, til hvilke de
staa i inderligt og dagligt Forhold. Det daglige Livs ydmyge
og kedelige Pligter og Selvfornægtelser fordre en rolig og bramfri
Selvbeherskelse\dots{} men men gøre ikke Krav paa offentlig Opmærksomhed
og Deltagelse.“ Der er i denne Ytring af Maudsley
rigtigt antydet, at den Sympati, der fører til Virksomhed i
større Kredse, saare let faar en Tilsætning af Ærgerrighed og
Forfængelighed. Ofte kan det ogsaa være Rastløshed, der fører
ud over de snevrere, men inderligere Forhold.

Naar Hengivelse og Kærlighed idethele opstilles i Etiken
som Dyd, der skal øves, eller Pligt, der skal opfyldes, forbinder
sig dermed ejendommelige psykologiske Vanskeligheder, der
ofte ere blevne fremhævede fra filosofisk Side --- i Oldtiden
især af Stoikerne, i nyere Tid især af Kant og hans Disciple.

Ideen om almindelig Menneskekærlighed fremtræder, som
vi have set (III, 8; IX, 1), hos græske Filosofer ide sidste
Aarhundreder før Kristendommens Fremtræden. En særligt inderlig
Form faar denne Følelse hos Kejsertidens Stoikere. \textsc{Marcus}
\textsc{Aurelius} fordrer saaledes, at vi ikke blot skulle føle os som
Del (meros), men som Led (melos) af Menneskenes Rige. Man
skal altsaa forholde sig til Andre, ikke som en Sten ien Dynge
forholder sig til andre Sten i Dyngen, men som et Organ forholder
sig til andre Organer i samme Organisme. Og hermed
hænger den inderlige Betagethed sammen, hvorved Kærligheden
% CHECK p. 238: note whose mark was not found
\footnote[1]{Pathologie de Vesprit. p. 260.}
% --- p. 239 ---
\setcounter{footnote}{0}%
\opage{239}bliver en Berigelse af det egne indre Liv. „Hvis du blot er
Del, ikke Led“, siger Marcus Aurelius, „saa elsker du endnu
ikke Menneskene af Hjertet: det at gøre Andre vel glæder dig
endnu ikke saaledes, at det helt betager dig, men du gør det
som blot og bar Pligt, endnu ikke med den Følelse, at du derved
gør dig selv vel.“ --- Hvad der alligevel hindrede, at Kærligheden
vandt den Anerkendelse af de stoiske Filosofer, som
den vandt i Kristendommen, det var en anden Tendens, som
var ejendommelig for Stoicismen —og endnu mere ejendommelig
for den end Tendensen til Hengivelse, nemlig Bestræbelsen
for aandelig Selvstændighed og Uafhængighed af alt Ydre. Den,
der hengiver sig, gør sig jo afhængig; han raader ikke mere
over sit Følelsesliv; hans Glæde og Sorg betinges ved, hvad
der skerudenfor ham, og hans Følelser kunne sættes i stærkere
Svingninger, end det er foreneligt med Bevarelsen af den indre
Harmoni, der for den græske Individualisme var det Højeste.
Stoicismen er betænkelig ved al stærk Sindsbevægelse, men
især, naar denne fremkaldes udefra. Hos \textsc{Epiktet} fremtræder
tydeligst denne Side af Sagen. Vel skulle vi have Medfølelse
med Andres Sorg og ytre denne Medfølelse, sukke med dem:
„men“, tilføjer Epiktet, „vogt dig vel for, at du ikke kommer
til at sukke i dit Inderste!“\footnote[1]{Garbe: Die Sankhyaphilosophie. p. 143). --- Se min Sammenligning mellem Buddha og Jesus: Religionsfilosofi p. 270 og 272.} --- I Modsætning til denne Frygt
for at blive afhængig gennem Kærligheden staar Kristendommens
Tro paa, at den aandelige Frihed kan bestaa med Følelseslivets
stærkeste Svingninger, ja endog først ret bliver til hos

') Ifølge Stoikerne hørte Medlidenheden til Affekterne (τὰ πάθη) og
var en Last. Den Vise føler ingen Medlidenhed. Diogenes Laërt. VII,
111. 123. --- Medens --- siger Seneca --- kraftige Naturer ere tilbøjelige
til Vrede, ere milde Naturer tilbøjelige til ringere Fejl, som Medlidenhed.
Ligesom det er uværdigt for den Vise at blive afhængig af Andres Slethed
ved at blive vred, saaledes er det uværdigt for ham at blive afhængig
af Andres Ulykke ved at føle Medlidenhed. De ira II, 7. 15 (kfr. I, 7---8).
--- Ogsaa for den indiske Visdom indeholdt Medlidenhed en Fare. Den
kongelige Vismand Bharata var lige ved at opnaa den forløsende Erkendelse,
men gik glip af den, da han plejede en lidende ung Gazelle. (R.
% --- p. 240 ---
\setcounter{footnote}{0}%
% CHECK p. 240: notes paired with the marks in order (the layer misread a number)
\opage{240}den, der helt og fuldt gennemlever de store Modsætninger af
Sorg og Glæde, Frygt og Haab. Kristendommen staar her med
et realistisk Præg i Forhold til den græske Filosofi. Den frygter
ikke for, at Hengivelsen til Livets Virkelighed, Individets fulde
Indleven i Slægtens Skæbne skal berøve Sjælelivet dets Frihed
og Kraft. Den sætter Livet til for at vinde det. Her er et
Element i den urkristelige Etik, som den filosofiske Etik maa
tilegne sig. Det er et etisk Experiment, som har en blivende
Betydning, uafhængigt af de specielle Forhold, under hvilke
det foretoges\footnote[1]{Se Sammenligningen mellem Stoicismen og Kristendommen i min Afhandling Hedenske Sandhedssøgere (Mindre Arbejder. I. p. 174—179).}.

\textsc{Immanuel} \textsc{Kant} og hans Skole saa ikke Sagen paa denne
Maade. Naar Kant mente at være i Overensstemmelse med den
kristelige Etik, saa var det, fordi han opfattede dennes Kærlighedsbud
som Opstilling af et Fuldkommenhedsideal, til hvilket
det var muligt at nærme sig, men som intet endeligt Væsen
fuldstændigt kunde virkeliggøre. Kærlighed er en Følelse,
--- og en Følelse frembringes ikke paa Befaling, siger Kant; derfor
maa det, der befales ved hint Bud, være praktisk Udøvelse
af Pligten i den ideale Form, hvor Udøvelsen udspringer af inderste
Hjerte og Sindelag. Men hos Mennesket, fortsætter Kant, ytrer
der sig stedse andre Motiver end det, der bestaar i Agtelse
for Loven, og Selvtvang er derfor stedse nødvendig. Idealet er
helligt, men Mennesket kan kun drive det til Dyd\footnote[2]{Kritik der praktischen Vernunft. Kehrbachs Udg. p. 100—102. Smlgn. Tugendlehre § 25---26 (Menschenliebe, ikke som „Liebe des Wohlgefallens“, men som „Maxime des Wohlwollens“).}. En af
Kants danske Disciple, M. G. \textsc{Birckner}, nærede særligt den
Betænkelighed ved det kristelige Kærlighedsbud, som det udtaltes
af Kristus, at det fører til at begrunde Moralen paa Egennytte:
„De Egenskaber, der opvække Følelsen af Kærlighed
hos os, ere netop kun saadanne, som enten umiddelbart bidrage
eller dog kunde bidrage til at glæde og gavne os selv....
Kærlighed er altsaa i Grunden, efter sin Oprindelse, en egennyttig
Følelse\dots{}. Som Kærlighed har sit Udspring af Egen%
% --- p. 241 ---
\setcounter{footnote}{0}%
\opage{241}nytte, saa har derimod Agtelse for Personer sit Udspring fra
den rene, uegennyttige Agtelse, vi i dem føle for Moralloven
selv.“ Birckner løser Vanskeligheden paa lignende Maade som
Kant. Hvad der i Kristendommen menes med Ordet Kærlighed,
er, mener han, „Lyst til det moralske Gode for dets egen
Skyld“\footnote[1]{M. G. Birckner: Efterladte Skrifter. København 1800. p. 175---184.}.

Følelsens Psykologi giver Kantianerne baade Ret og Uret
i deres Betænkeligheder. Det er i en vis Forstand rigtigt, at
Kærlighed ikke kan kommanderes frem. Det nytter ikke at
fordre Kærlighed, hvor der ingen Kærlighed er. Men hvis man
gaar ud fra, at Etik er Mere end Jura, der er „slaaet ind“, og
at det Etiske kan bevare sin Gyldighed, selvom Kommandotonen
aflægges, vil man indse, at Følelser meget vel kunne
fremkaldes, og at der derfor meget vel kan fordres en Stræben
efter at tilvejebringe Betingelserne for Kærlighedens Opstaaen.
Ofte kan en Følelse kun fremkaldes indirekte og ad Omveje;
det gør Opgaven indviklet, men gør ikke dens Løsning umulig.'
Det er et vigtigt Hensyn ved Ordningen af Samlivsforhold, ved
Opdragelse og Selvopdragelse, at der sørges for, at Betingelserne
for kærlighedsfuld Hengivelse komme tilstede. Den uvilkaarlige
Samleven og Samvirken vil, naar den faar Lov at udfolde
sig i Ro, stille og langsomt virke i denne Retning. Og Henblikket
paa store Forbilleders inderlige og kærlighedsfulde Liv
vil vække til begejstret Efterligning af den Aandskraft, som udtrykker
sig deri. Sikkert faa Kantianerne ogsaa Ret i, at der
her kun kan være Tale om Tilnærmelse. Den sande Kærlighed
er en saadan Kraft, at Enhver i Sammenligning med den vil
have stadig Lejlighed til at føle sin Afmagt og sin Afstand fra
Idealet. Men dette gælder jo dog om Forholdet til alle Idealer,
til Selvhævdelsens Idealer ikke mindre end til Hengivelsens.
Hvormange etiske Karakteregenskaber gælder det overhovedet
om, at de kunne bringes til Veje umiddelbart, efter kategorisk
Befaling?

Det er grundet paa falsk Psykologi, naar Kantianerne mene,
at al Følelse med Undtagelse af Agtelsen for den moralske Lov
% --- p. 242 ---
\setcounter{footnote}{0}%
% CHECK p. 242: notes paired with the marks in order (the layer misread a number)
\opage{242}er af egennyttig Art. Til Egennytte eller Egoisme udfordres
en bevidst Henføren af Alt til sig selv som Formaal. Det er
ikke nok, at en Følelse er forbunden med Selvtilfredsstillelse,
for at den skal kunne kaldes egoistisk (naar man ikke giver
dette Ord en meget vid Betydning). Deraf, at Hengivelsen medfører
eller betinges ved Velbehag ved den Genstand, til hvilken
vi hengive os, kan ikke sluttes, at vi gøre os selv til Formaal
og kun behandle Genstanden som Midlet til vor Nydelse. Hengivelsen
vilde være psykologisk umulig, dersom Genstanden
ikke kunde vække vort Velbehag\footnote[1]{Smlgn. min Psykologi VI C, 7.}. Og den Agtelsesfølelse, Kant
beskriver, er en psykologisk Umulighed, naar det, der agtes,
ikke kan vække nogetsomhelst Velbehag hos os. Er Kærlighed
Egennytte, saa er Agtelse det ogsaa.

2. I hvert enkelt Tilfælde kan Kærligheden antage forskellige
Former, som hver for sig kunne have deres etiske Interesse\footnote[2]{Aristoteles har i 8de og 9de Bog af den nikomakeiske Etik underkastet disse forskellige Former en Undersøgelse, som endnu har klassisk Betydning. (Dansk Oversættelse af disse Bøger ved E. Bojesen i Program fra Sorø 1859).}.
Den kan ikke altid have Karakteren af en umiddelbar Enhedsfølelse,
selvom de Individer, den gælder, ikke staa os fjernt.
Den ene Part kan staa saaledes over den anden, at Forholdet
fra denne sidstes Side faar et Præg af Beundring og Ærefrygt,
som dog ikke behøver at ændre noget i den gensidige Kærligheds
Ligelighed. Der kan føles Fællesskab og Sammenhøren
trods Forskellen i Karakter og Evner. Hos den overlegne Part
ytrer Kærligheden sig som Højsindethed. Det Høje i Højsindetheden
beror ikke paa, at Individet føler sig staaende paa et
højere Stade end andre; men netop paa, at han ikke føler sig
staaende paa et højere Stade, uagtet han virkeligt gør det. Den
inderlige Hengivelse lader alle saadanne Forskelle forsvinde.
Højsindetheden kan fremtræde som Humor, naar der vel er
Blik for Andres Ringhed og Begrænsning, deres Selvmodsigelser
og Ufuldkommenheder, men dette ikke hindrer Sympatien med
dem. Det Høje i Højsindetheden beror fremdeles paa, at Med%
% --- p. 243 ---
\setcounter{footnote}{0}%
\opage{243}følelsen bestaar uafhængigt af de Andres Færd, trods deres
Utaknemlighed og Ligegyldighed. Og endeligt beror det Høje i
Højsindetheden paa, at Blikket er fæstet mod store, omfattende
Livsinteresser, saa at Sympatien ikke er begrænset til saadanne
snevrere Kredse, hvor den Enkelte umiddelbart træder overfor
Andre. Den Højsindedes Kærlighed til Menneskeheden bestemmes
ved en omskuende Erkendelse af Menneskeslægtens Livsbetingelser;
det er hans Stræben at hævde disse, selvom han i det
Enkelte maa krænke, hvad en til snevrere Kredse bundet Sympati
maatte fordre. Med Urette mener derfor \textsc{Adam} Smith\footnote[1]{Theory of Moral Sentiments. IV, 2.},
at Højsindethed (generosity) og Menneskekærlighed (humanity)
ere to helt forskellige Følelser. Han slutter dette af, at den,
der sætter sit Liv til for at redde en Anden, maaske ikke vilde
føle egentlig Sorg, om denne Anden var død af en Aarsag,
hvis Virkning det ikke var muligt at forebygge. Men denne
Slutning er ikke rigtig. Det er jo dog ganske naturligt, at Menneskekærligheden
rører sig, naar der er Mulighed for Handlen,
selv hvor der ikke forud existerer noget nærmere Forhold til
det andet Individ. Det er kun i de snevrere Kredse, at den
hyppigt har Lejlighed til at optræde som aktuel; overfor fjernere
staaende Væsener holder den sig, afset fra specielle Anledninger,
i potentiel Form. \textsc{Schopenhauer} (Neue Paralipomena \$ 244)
har træffende bemærket, at der gives en Art Mod, der udspringer
af samme Rod som Hjertensgodheden, nemlig deraf, at
Mennesket føler sin Existens næsten som identisk med Andres.
„Denne Bevidsthed frembringer Mod derved, at Mennesket hænger
mindre fast ved sin individuelle Tilværelse, da det næsten lever
ligesaa meget i alle Væseners Existens og derfor bekymrer sig
mindre for sit eget Liv, og hvad dette angaar. Dette er ikke
altid Modets Kilde; ti Modet kan udspringe af forskellige Aarsager.
Men hint Mod er den ædleste Art af Mod, hvilket viser
sig i, at det her er forbundet med stor Blidhed og Taalmodighed.“
--- Der er i den højsindede Kærlighed en saadan Kraft
og Fylde, at den i sig selv hører til de Ting, der give Livet
Værdi. Den er værdifuld, ikke blot gennem sine Virkninger,
% --- p. 244 ---
\setcounter{footnote}{0}%
\opage{244}men ved selve sin Existens, og dens Virkninger ere kun dens
uvilkaarlige Expansion. ---

Sand Kærlighed kan ikke vises ovenfra nedad. Det faar
især Betydning for Velgørenheden. Ofte rækkes Gaven ned fra
Skyerne, og Medlidenheden er hyppigt forklædt Hovmod. Den
højsindede Kærlighed udsletter Forskellen mellem Giver og Modtager,
idet den anerkender den Andens Ret til at tage, saavel
som sin egen Pligt at give. Den arbejder hen til at gøre den
Lidende og Trængende til et selvstændigt Led af Slægten. Den
behandler heller ikke Modtageren som blot Middel til at faa
Udladning for blinde sympatiske Instinkter eller for et sentimentalt
Hang. Det blinde Instinkt tør vi her ofte ligesaa lidt følge,
som Vrede, Ærgerrighed og andre Tilskyndelser, selvom vi i
Tvivls Tilfælde gør bedst i at lade Hjertet tale. Vi give jo
ikke blot, fordi vi føle Trang dertil, men fordi vi ville bringe
virkelig Hjælp. \textsc{Malthus} er den Første, der har paavist, hvor
vanskelig en Kunst Velgørenhed er. Her viser sig atter Kærlighedens
nøje. Sammenhæng med Retfærdigheden. Under det
sympatiske Instinkts Virken maa Hensynet til større Kredses
Tarv holdes fast.

Hos den, der fæster Blikket paa den store Kløft mellem
de Trængende og Lidende paa den ene Side og dem, der have
forholdsvis let Adgang til Livets materielle Goder, paa den
anden Side, opstaar let en Svimmelhed, som kan svække baade
Sympatien og Energien. Man skal naturligvis ikke lukke Øjnene
for denne Kløft. Men det store Problem om Fattigdom og Lidelse
er langt mere indviklet og forgrenet, end at det kan løses
ved direkte Indgriben. Det nyttede ikke, om vi Alle gav alt
vort Gods til de Fattige. Vi maa søge vor Retfærdiggørelse i,
at vi ved den Gerning, vi have valgt, yde det bedste Bidrag,
vi formaa, til Slægtens Velfærd. Men fuldstændigt trygge kunne
vi aldrig føle os. Naar have vi gjort nok? Er vor Iver tilstrækkelig?
Skyldes vor gunstigere Stilling ikke for en Del
Uretfærdigheder i Godernes Fordeling i Samfundet? Arbejde vi
virkeligt i Retning af en bedre Fordeling? --- Spørgsmaal som
disse ville bevare deres Brod. De fysiske og moralske Ulykker
% --- p. 245 ---
\setcounter{footnote}{0}%
\opage{245}i Verden føles undertiden stærkest netop af dem, som ikke selv
direkte lide under dem. ---

Der gives ikke blot en fysisk, men ogsaa en aandelig Velgørenhed,
og denne sidste er nok saa vigtig som hin. Ogsaa
her gælder det om at hjælpe Andre og hævde sig selv. Den
aandelige Hjælp er endnu vanskeligere at bringe end den fysiske.
\textsc{Kant} negtede endog, at man kunde arbejde for andre Menneskers
Fuldkommenhed, hvorimod han nok vilde indrømme, at
man kunde arbejde paa deres Lyksalighed. Hans Tanke var,
at et Menneskes væsentligste Egenskaber maa være hans eget
Værk, udspringe af hans egen Villen. Heri ligger den Sandhed,
at al aandelig Hjælp kun kan bestaa i Vækkelse af Egenvirksomhed.
Men dette gælder --- skønt Kant ikke vilde indrømme
det --- „Lykken“ ligesaa vel som „Fuldkommenheden“. Lykken
har kun fast Grund, naar den er selverhvervet. Dog gælder
det ganske særligt ved indre, aandelige Goder. Den Givende
maa her, endnu mere end ved den fysiske Hjælp, gøre sig til
Middel. Uden Utaalmodighed og uden Egoisme maa han sætte
sig ind i Modtagerens Tilstand for at kunne give ham eller
være for ham, hvad der trænges til. Ovenfra nedad kan heller
ikke her Noget udrettes. Der skal, som \textsc{Sokrates} udtrykte det,
ydes en aandelig Fødselshjælp; de Spirer, Individet skjuler i
sig, skulle lokkes frem. Denne sokratiske Metode er den eneste
frugtbare. Den er Kærlighedens egen Metode. --- Aandeligt
forkomne Mennesker hjælper man allerede ved at vise dem, at
man tror dem i Stand til at udrette Noget. Den Tillid, man
viser dem, kan blive Spiren til deres egen Selvtillid. ---

Den højsindede Kærlighed er ikke betinget ved, hvad Individet
selv har modtaget eller lidt. Det er det fulde og rige
indre Liv, der gør, at den Højsindede ikke ændrer sin Færd
paa Grund af Andres Raahed, Utak eller Fjendskab. (Smlgn. IX,
1). Den er Følelsesudtrykket for Slægtens Enhed; men denne
Enhed kan ikke ophæves ved, hvad der end kan ske. Det var
et vigtigt Gennembrud i Slægtens etiske Udvikling, det vigtigste
Skridt paa Vejen til den almindelige Menneskekærlighed, da
man begyndte at hæve sig over den Sætning, som udtrykker
det barbariske Udviklingstrins Standpunkt: du skal elske din
% --- p. 246 ---
\setcounter{footnote}{0}%
\opage{246}Ven og hade din Fjende! Det var en Dualisme af Kærlighed
og Had. Indenfor den europæiske Verden finde vi i \textsc{Platons}
Skrifter (Kriton, Gorgias, Statens første Bog) for første Gang\footnote[1]{Der er paa dette Punkt en Modstrid mellem Xenofons og Platons Fremstillinger af Sokrates. Ifølge hin hyldede Sokrates det gamle Standpunkt, at det er Tegn paa en dygtig Mand at elske sine Venner og skade sine Fjender. (Memor. II, 3, 14; IV, 2). --- Løsningen kunde maaske ligge i, at Sokrates selv ikke har udtrykt sig bestemt om dette Punkt, og at Striden om det først er opstaaet i hans Skole, hvor saa Xenofon (i Tilslutning til Antisthenes) har hævdet den gængse Opfattelse, medens Platon, i Modsætning til denne og ud fra sin ideale Etik, har fortolket Mesterens mere ubestemte Udtalelse saaledes, at Modsætningen mellem Ven og Fjende ikke gælder paa det højeste Standpunkt. Denne Hypotese er opstillet af Karl Joël: Der echte und der xenophontische Sokrates. I. p. 396 f. (Kfr. II, 2. p. 1008 f.). --- Allerede Demokritos havde iøvrigt udtalt (fragm. 48 Natorp), at den, der øver Uret, er ulykkeligere end den, der lider Uret. Denne Udtalelse hænger nøje sammen med, at Demokritos lægger Vægten paa det indre Sindelag, Sindets indre Tilstand, i Modsætning baade til ydre Skæbne og til ydre Handlen. Heri mødes den demokritiske og den platoniske Etik. Smlgn. Natorp: Die Ethica des Demokritos. p. 102 f.}
en helt anden Grundsætning opstillet: „man skal aldrig paaføre
Ondt uden som opdragende Straf, og det er bedre at lide Uret
end at gøre Uret“. Sætningen indføres med stort Eftertryk og
med fuld Bevidsthed om dens Modsætning til den gængse Tænkemaade.
„Det véd jeg godt“, siger \textsc{Sokrates} i Platons Kriton
(p. 49 D), „at det baade nu kun er nogle ganske faa, der mene
dette, og at det altid vil være saaledes. Hvem der da nu har
denne Anskuelse, og hvem der ikke har den, disse kunne ikke
forhandle sammen; nej, de maa uundgaaeligt foragte hinanden,
naar de gensidigt se hinandens Tænkemaade.“ Det Nye er, at
selv Forholdet til Fjenden nu falder ind under Kærlighedsforholdet,
idet Hævninstinktet underordnes Hensynet til en højere
Retfærdighed. Deraf følger ikke en asketisk Passivitet, saaledes
som der vilde følge af Bjergprædikenens bekendte Bud, hvis
man tog det efter Ordlyden. Man gør sig ikke til blot Objekt
for Andre, fordi man ikke lader den blinde Hævnlyst raade,
som ingen Grænser kender uden sin egen Mættelse, og som
ikke har den Andens Vel til sidste Maal. --- I Læren om Statens
% --- p. 247 ---
\setcounter{footnote}{0}%
\opage{247}Straffemyndighed ville vi igen møde denne Modsætning mellem
Hævnens (Gengældelsens) Grundsætning: „Ondt for Ondt!“ og
den opdragende Kærligheds Grundsætning: „intet Ondt uden
for det Godes Skyld!“

3. Naar vi definere Sympati som Lyst ved Andres Lyst,
Ulyst ved deres Ulyst, saa ligger deri ikke forudsat Mere om
disse Andre, end at de have Evne til at føle Lyst og Ulyst.
Den uinteresserede og universelle Sympati maa derfor strække
sig ud over den menneskelige Verden, og Velfærdsprincipet vil
kunne finde Anvendelse ogsaa ved Forholdet til andre end menneskelige
Væsener; det fordrer jo i sin Almindelighed blot det,
at Lystfølelsen saavidt muligt skal tiltage, Ulystfølelsen aftage
i Verden. Der kan derfor meget godt tales om Pligter mod
Dyrene og om Rettigheder for Dyrene. Det er jo ikke høj
Fornuft og stor Fatteevne, men blot Evne til at føle og lide,
som er Betingelsen for at være Objekt for etisk Pligt og forsaavidt
Subjekt for etisk Ret. Saasnart Medlidenheden faar en
afgørende Betydning som Element i Vurderingsmotivet, vil det
følge af sig selv, at der anerkendes Pligter mod Dyrene, og at
disse ikke kunne behandles blot 'som Midler, saa at hvilken
som helst Behandling af dem skulde være forsvarlig. Efterat
først \textsc{Clarke} havde talt Dyrenes Sag udfra sit ejendommelige
Moralprincip, der fordrede, at ethvert Væsen skulde behandles
efter sin Natur, og efterat hos os Fr. \textsc{Eilschow} udfra Ideen
om et stort Samfund af alle levende Væsener havde hævdet,
at der gaves Pligter mod Dyrene, var det \textsc{Hutcheson} og Rousseau,
som udledte Pligterne mod Dyrene af Medlidenhedsfølelsen:
da Evnen til at lide er fælles for Dyr og Mennesker, maa unødvendig
Tilføjelse af Lidelse ligesaa vel være forkastelig overfor
hine som overfor disse. Og senere erklærede Bentham: „Spørgsmaalet
er ikke: kunne Dyrene tænke? eller kunne de tale?
men: kunne de lide?“\footnote[1]{Fr. Eilschow: Philosophiske Breve. København 1748. p. 90 f. (Eilschow polemiserer særligt mod Pufendorf, der i sit Skrift Et Menneskes og en Borgers Pligter efter Naturens Loi'. Dansk Overs. 1742. p. 257 havde udtalt, at „efterdi Bæsterne ikke havde Fornuft, uden hvilken man ikke kan gøre sig nogen Concept om nogen egentlig saa kaldet Ret og Forpligtelse, er der følgeligen ingen fælles Ret for Menniskene og Bæsterne“). --- Hutcheson: System of Moral Philosophy. London 1755. I. p. 314. --- Rousseau: Discours sur Vinégalité. Amsterdam 1755. p. LXII f. --- Bentham: Principles of Morals and Legislatlon. XVII, 1, 4. Note. --- Om Clarke's Lære se Ueberweg-Heinze: Grundriss der Geschichte der Philosophie. 8 Aufl. 111, 1. p. 159.} --- Naar vi alligevel ikke regne Dyrene
% --- p. 248 ---
\setcounter{footnote}{0}%
% CHECK p. 248: notes paired with the marks in order (the layer misread a number)
\opage{248}med til den etiske Verden, der kun omfatter Menneskeslægten,
kommer det af forskellige Grunde.

For det Første ere Dyrene vel Objekter for Pligt, men
ikke selv Subjekter for Pligt i samme Forstand, som Menneskene
ere det. De kunne altsaa ikke som selvstændigt Led
udfylde en Plads i den etiske Verden, og deres etiske Ret kan
derfor heller ikke være den samme som Menneskets. For det
Andet har Dyrets Skæbne paa Grund af dets ringere og mere
begrænsede Natur ikke den store Betydning, som et Menneskes
Skæbne har. Og for det Tredie kan den Smerte, Dyr føle,
neppe naa en saadan Grad af Styrke, som menneskelig Smerte
kan\footnote[1]{Se Panum: Indledning til Fysiologien. 2. Udg. p. 7. --- Smlgn. min Psykologi I, 6. --- Dette indrømmer endog den ivrige Dyreven Schopenhauer (Die beiden Grundprobleme der Ethik. 2. Aufl. p. 245).}. Naar Mennesket overfor Dyret gør Brug af den Stærkeres
Ret og benytter det som Middel for sine Formaal, saa
maa dette retfærdiggøres ved de menneskelige Formaals Betydning
i Sammenligning med den Lyst eller Smerte, Dyret kan
føle. Men at Dyret paa intet Punkt kan betragtes som blot
Middel, kommer netop af, at det kan lide. (Smlgn. VIII, 6).
Bevisbyrden maa paahvile den, der vil have Dyrene behandlede
anderledes end Menneskene.

Naar man ikke direkte, udfra Velfærdsprincipet, vil begrunde
Pligterne mod Dyrene, forsøger man det ad en Omvej.
Man siger da f. Ex. med Kant\footnote[2]{Tugendlehre. § 16—17.}, at grusom Behandling af
Dyrene strider mod Menneskets Pligt mod sig selv, „fordi Medfølelsen
med deres Skæbne derved sløves, og altsaa et for den
moralske Adfærd overfor andre Mennesker meget gavnligt naturligt
Anlæg svækkes og efterhaanden udryddes“. Denne Betragtning
har sin Betydning som underordnet Motiv. Men Med%
% --- p. 249 ---
\setcounter{footnote}{0}%
\opage{249}følelse med Dyrene har ikke blot den Betydning at være Forberedelse
og Skole for Medfølelse med Menneskene; den har
ogsaa en umiddelbar Værdi ved at lindre en Del af den Smerte,
der føles i Verden. Den, der hjælper en Fugl af Klemme, gør
en god Gerning, ikke blot fordi han styrker sin Evne til at
have Medfølelse med andre Mennesker. Fra retsfilosofisk Side
begrundes Forbud mod Mishandling af Dyr gerne paa Hensynet
til den menneskelige Følelse, der fordærves eller krænkes derved.
Det er her snarere Andres Følelse end den Handlendes egen,
der tages Hensyn til\footnote[1]{638. Enkelte tyske Stater have i deres Straffelov optaget dette Synspunkt}. Men ogsaa dette er en Omvej eller et
Bihensyn, som ikke maa træde i Stedet for den direkte Begrundelse.

% CHECK p. 249: centred line: heading or verse
\begin{center}b. Sandhedskærlighed.\end{center}

4. Kærligheden kan ikke blot være rettet mod enkelte bestemte
Væsener, men ogsaa mod Ideer, der betage os og gøre Krav
paa vor Opofrelse. Der kan opstaa en Modstrid mellem de to
Arter af Kærlighed. Der kan gives en lidenskabelig Hengivelse
til Ideer (videnskabelige, kunstneriske, politiske eller religiøse
Ideer), som ikke tager i Betænkning at træde de virkelige levende
og følende Væseners Vel under Fødder eller maaske slet ikke
ænser det. Paa den anden Side kan Trangen til at hjælpe i de
enkelte, personlige Forhold ogsaa gøre blind for de almindelige
Interessers og Bestræbelsers Betydning, paa hvilke dog al højere
og friere Udvikling af det menneskelige Liv beror. Man har
kaldet Optagetheden af almindelige Ideer Alsind og stillet den

') Jhering: Der Zweck im Recht. Il. p. 138 ff. Goos: Retslære.
I. p. 169. --- Den tyske Rigsstraffelov forbyder kun offentlig og forargelig
Mishandling af Dyr (§ 360); den danske Straffelov forbyder raa Mishandling
eller anden grusom og oprørende Mishandling af Dyr (§ 297). ---
Derimod fremhæver Lon ing som Formaal med Forbud mod Dyrplageri
ikke blot Beskyttelse af Befolkningens sædelige Følelse, men ogsaa „Beskyttelse
af Dyrene selv mod unyttige og derfor usædelige Mishandlinger“.
Sittlichkeitspolizei (Schönbergs Handb. der polit. Oekon. 1. Ausg. Il). p.
% --- p. 250 ---
\setcounter{footnote}{0}%
% CHECK p. 250: notes paired with the marks in order (the layer misread a number)
\opage{250}i Modsætning til Humanitet i Betydning af Omsorg for hver
enkelt menneskelig Existens\footnote[1]{Gabriel Sibbern: Om Humanitet og Alsind. København 1855. Smlgn. F. C. Sibbern: Moralfilosofi § 22.}.

De enkelte Menneskers Trang til Lykke, Fred og Trøst er
det Første og tilsidst det Vigtigste. Alt Fremskridt kan jo tilsidst
kun maales ved, om saadan Trang kan tilfredsstilles bedre
og mere end før. Alsindets, eller som vi her ville kalde det,
Sandhedskærlighedens Betydning forsvinder derfor ikke. Sandheden
betegner jo ikke Andet end Tilværelsens Sammenhæng,
saavidt vi kunne fatte den. I denne Sammenhæng kunne og
skulle vi leve os ind; gøre vi det ikke, bygge vi paa Sand.
Hvad de ovenfor nævnte Ideer vise os, er jo kun, at vi leve
under mere omfattende og anderledes beskafne Forhold, end vi
havde tænkt os.

Den Fordring: søg Sandheden! udspringer af Selvopholdelsestrangen.
Menneskets første Møde med Sandheden som Udtryk
for Tingenes faste Sammenhæng finder Sted, naar det første
Gang indser, at det uundgaaeligt behøver visse Midler for at
naa sine Formaal\footnote[2]{Psykologi VI F, 3 (smlgn. V D, 4).}. Her er der altsaa ikke endnu Modsætning
mellem Sandhed og Følelse, mellem Teori og Praxis, selvom
Ønsket om strax at naa Maalet rører sig med Utaalmodighed.
Men ved videre Udvikling af Erkendelsen kan et Fjendtlighedsforhold
indtræde. Det er ikke sagt, at Erkendelsens Resultater
umiddelbart tilfredsstille Følelsen. Nye Ideer kunne derfor ofte
kun føres igennem med en vis Hensynsløshed. For at bringe
nye Synspunkter og de større Horizonter frem maa man ofte
foreløbigt vække Forstemthed og Uvillie, Lidelse og Strid. Det
gaar her, som naar en Statsmand styrter sit Folk i en Krig:
han véd, at han vil volde Tusinder af Mennesker Smerte, Død
og Sorg; saa sikker kan han ikke være paa, at han ved Krigen
vil sikre og forhøje Folkets Fremtidslykke; og dog kan det være
nødvendigt at vove den, dersom større Ulykker ellers ville styrte
ind. Selv den, der stiftede Menneskekærlighedens Religion, sagde:
„Mene I, at jeg er kommen for at give Fred paa Jorden? Nej,
% --- p. 251 ---
\setcounter{footnote}{0}%
\opage{251}siger jeg Eder, men Tvedragt. Ti fra nu af skulle fem i et
Hus være splidagtige, tre imod to og to imod tre“\footnote[1]{Lukas' Evang. XII, 51—52. -) Se min Bog S. Kierkegaard som Filosof. p. 138—140.\textsuperscript{1} Lecky: History of European Morality. I. p. 52.\textsuperscript{1} Psykologi V D, 1. 5. (Smlgn. V B, 4).}. --- Intet
Under, at den, der i Sandhedens Navn vil rejse en stor og
bitter Strid, en Stund kan holdes tilbage ved Tanken om den
Lykke og Livstryghed, han vil forstyrre ved de skærpede Fordringer,
han føler sig forpligtet til at hævde! Sin egen Lykke
kunde han ofre, og maaske har han ikke en Gang nogen Lykke
at ofre; men de Andres Lidelse bliver en Vanskelighed, der
for en Stund kan holde ham tilbage. I Søren Kierkegaards Dagbøger
før hans voldsomme Kamp mod den sædvanlige Opfattelse
af Kristendommen findes denne Vanskelighed skildret paa en
inderlig og ophøjet Maade\footnotemark[1].

Man har ment, at den ubetingede Søgen efter Sandheden
ikke skulde kunne begrundes ved Velfærdsprincipet. „Vi skylde“,
har man sagt\footnotemark[1], „vore Illusioner Mere end vor Erkendelse. Vor
stedse arbejdende Fantasi bidrager sandsynligvis Mere til vor
Lykke end Fornuften, som paa det spekulative Omraade hovedsageligt
er kritisk og nedrivende“. --- Hertil er at svare, at
Fantasien meget godt kan virke, uden at vi hilde os i Illusioner.
Illusionerne opstaa først, naar vi forvexle Fantasibilleder og
Virkelighed. Al Sandhed er tilsidst praktisk, er et Lys, der skal
lede Villien. Hvis Illusionen ligesaa godt som Sandheden kan
lede Villien, saa mangle vi ethvert Skelnemærke mellem Illusion
og Sandhed. Ti det eneste Kendemærke, vi have paa Sandheden,
er det, at vi kunne drage alle Konsekvenser af den,
kunne bygge paa den baade i vor Handlen og i vor Tænken\footnotemark[1].
Grænsen mellem Illusion og Sandhed kan ikke drages én Gang
for alle. Det er en uendelig Opgave, der er sat vor Videnskab:
at danne et tro og sammenhængende Verdensbillede som Udtryk
for vore Iagttagelser. Arbejdet paa denne Opgaves Løsning er
--- selv afset fra det praktiske Udbytte --- forbundet med en
Følelse af Tilfredsstillelse, som langt kan overgaa den Glæde,
man føler ved at vugge sig i Illusioner.
% CHECK p. 251: note mark without a note

% --- p. 252 ---
\setcounter{footnote}{0}%
\opage{252}Pligten at søge Sandhed udspringer af det psykologiske Forhold
mellem Erkendelse og Villie. Allerede de mest primitive
Villiesytringer, Instinktet f. Ex., som man saa ofte har prist for
dets Sikkerhed, kunne tage fejl. Instinktet fremlokkes stedse
ved Fornemmelser; men disse Fornemmelser kunne optræde,
uden at derfor Betingelserne for Instinktets heldige Virken ere
tilstede. Dyret, der lokkes i Fælden ved Lugten af Lokkemaden,
kommer til at lide, fordi Instinktet hindrer det i at undersøge
alle Omstændighederne. Vi ere idethele tilbøjelige til at tillægge
vore Forestillinger større Gyldighed, end de have, og
det er først de praktiske Konsekvenser, der føre til den fornødne
kritiske Begrænsning. Den, som ikke vaager over sine
Tankers Sandhed og Klarhed, vaager heller ikke over den Retning,
hans Villie tager; og da hans Villen ikke blot medfører
Konsekvenser for ham selv, men ogsaa for Andre, kan Manglen
paa Kærlighed til Sandhed blive til Haardhjertethed og Ligegyldighed
for Andres Vel.

Nærmere beset er der tre Forhold, ved hvilke Sandhedskærligheden
kan vise sig: Forholdet mellem Tanken og Ordet,
Forholdet mellem Tanken og Personligheden, Forholdet mellem
Tankens Indhold og dens objektive Gyldighed.

Ordet skal svare til Tanken, naar man idethele taget vil
give sin Tanke Ord. Derved udelukkes, at man mener Et og
siger et Andet. Det vanskelige Punkt ligger her, som paa saa
mange Punkter i Etiken, i det pædagogiske Hensyn, det kan
være nødvendigt at tage. Muligvis vilde et Ord, der fuldt ud
udtrykte Tanken, ikke lede til saa fuld og selvstændig Forstaaelse,
som en halvkvædet Vise eller en ironisk Modsigelse.
En saadan Pædagogik udspringer netop af Sandhedskærligheden.
Mod denne strider det kun, at man ved Akkommodationer og
Omtydninger vedbliver at bruge Ord, der kun have Betydning,
naar der bag dem ligge helt andre Tanker end de, man faktisk
har. Paa det religiøse, det politiske og det sociale Omraade er
der rigelig Anledning til at vise denne Art af Sandhedskærlighed
--- Ærlighedens Sandhedskærlighed --- i Modsætning til
uhæderlige Kompromiser.

Men dybere end dette Forhold ligger Forholdet mellem
% --- p. 253 ---
\setcounter{footnote}{0}%
\opage{253}Tanken og Personligheden. Det kommer her an paa, om Tanken
virkeligt er Personlighedens Ejendom. Sandhed betyder her personlig
Sandhed, Overensstemmelse mellem Tanken og Personligheden,
en Overensstemmelse, der fremkommer ved, at Tanken
mødes med levende Tilslutning fra det Centrale i Personligheden,
og at dens Anerkendelse er Frugt af Selvvirksomhed,
ikke af blind, passiv Modtagen. Poul MøLLER gik endog (i sine
Fragmenter om Affektation) saa vidt, at han sagde: „Ingen Livsytring
har Sandhed, uden deri ligger skabende Selvvirksomhed“.
Affektationen bestaar for ham i, at man optager Tanker i sig
uden at gøre dem til sin virkelige Ejendom, eller rettere til sit
eget Værk. Men Affektationen udelukkes jo dog, hvor Tilslutningen
er selvvirksom, uagtet Selvvirksomhed ikke altid kan
siges at være „skabende“. Der maa jo desuden ofte være Overgangstider,
hvor de nye Tanker (hvad enten de ere frembragte
af Personligheden selv eller af Andre) endnu ikke ere blevne
personlig Ejendom, og det vilde være uretfærdigt at stemple
Tilstanden i saadanne gærende og søgende Tider som Affektation.
\textsc{Søren} \textsc{Kierkegaard} har fortsat Poul Møllers fragmentariske
Tænkning ved sin lidenskabelige Betoning af det subjektive Forhold
til de Tanker, der bestemme Livsanskuelsen. Han fordrede
paa disse Punkter en Tænken „med Hjertekulen“, --- en „subjektiv
Tænken“, en \textasciitilde{}Tænken i Existens“, en Tænken, der gaar
i Et med en Villen. Derved har han indført en stor Værdimaaler:
Livsanskuelserne ere at prøve ikke blot efter deres
objektive Indhold, men særligt efter den Rod og det Tag, de
have i de Personligheder, som hylde dem; --- og omvendt bestaar
det at hylde en Livsanskuelse nu ikke mere i at vugge
sig i visse Tanker eller Trosartikler, men i at leve og dø paa
dem som det inderligste Tilhold, man kan finde. Og i sit eget
Liv og sin egen Tænken gav Kierkegaard selv et stort Forbillede
paa en saadan existentiel Tænken\footnote[1]{Søren Kierkegaard som Filosof. p. 24—26; 44—53; 127—129- 146.}. Derved har han
bevist Menneskeheden den største Tjeneste, han efter sin Natur
og sine Livsvilkaar kunde yde den: at foretage det store Forsøg
paa at prøve en overleveret --- ofte paa meget udvortes Vis
% --- p. 254 ---
\setcounter{footnote}{0}%
\opage{254}overleveret --- Livsanskuelses Forhold til det personlige Liv i
vore Dage, at undersøge dens Bærekraft og Muligheden af at
fyldestgøre dens virkelige Fordringer. Et i sin Art enestaaende
Experiment, der maaske giver vigtigere Bidrag til en Livsanskuelses
Vurdering end de fleste lærde Undersøgelser om
dens objektive Indhold!

Men ogsaa Gyldigheden af Anskuelsernes objektive Indhold
maa prøves uforbeholdent. De subjektive Tænkere, der lægge
saa stor Vægt paa Personlighedens Forhold til Ideerne, ere i
Regelen!) tilbøjelige til at undervurdere den objektive Kritiks
og Undersøgelses Betydning. I deres Iver for at gøre Sandheden
personlig glemme de den Pligt at undersøge, om det nu
ogsaa er Sandheden, de forholde sig til. Imod deres Villie komme
de da derved let til at betrygge Andre i deres vanemæssige
Forhold til det Overleverede. De overse, at Selvvirksomheden
netop skal vise sig i at prøve og finde Sandheden. Som den
tredie Art af Sandhedskærlighed maa derfor nævnes den intellektuelle
Redelighed, der hindrer, at man uden tilstrækkelig Prøvelse
antager og udtaler Noget som objektiv Sandhed. Det er
gennem denne Dyd, at Sandhedskærligheden fra den Enkeltes
Stræben efter at bevare Sammenhængen mellem Ord og Tanke
og mellem Tanke og Personlighed kan føre over til videnskabelig
Undersøgelse. Denne Dyd har man med Rette kaldet „den
yngste og vanskeligste af alle Dyder“. Den kunde ikke indskærpes,
før der havde dannet sig bestemte videnskabelige Metoder,
og før adskillige „Sandheder“ havde forvandlet sig til
ligesaa mange Problemer. I sin skarpeste Form fremstilles den
af W. K. \textsc{Clifford} i hans Afhandling The Ethics of Belief
(Contemporary Review 1877). Han hævder, at et Menneskes
Meninger eller Trosanskuelser ingenlunde ere en Privatsag, saa
at han skulde være berettiget til kun at tage Hensyn til, hvad
der tiltaler eller trøster ham. Sandheden er en hellig Arv fra
Slægt til Slægt, og Enhver har et Bidrag at yde til den. At
gøre sig selv lettroende er derfor en Forbrydelse mod 
% CHECK p. 254: note whose mark was not found
\footnote[1]{Smlgn. foruden Kierkegaard: Pascal, Rousseau, Jacobi, Carlyle. (Se de vedkommende Afsnit i „Den nyere Filosofis Historie“).}Menneske%
% --- p. 255 ---
\setcounter{footnote}{0}%
% CHECK p. 255: notes paired with the marks in order (the layer misread a number)
\opage{255}slægten. Kan man ikke vinde selve Sandheden, skal man prøve
Sandsynligheden, og kan man ikke selv prøve Sandheden, skal
man prøve de Autoriteter, man holder sig til, --- prøve dem
ikke blot med Hensyn til deres Ærlighed, men ogsaa med
Hensyn til Muligheden af, at de skulde have fundet Sandheden.
Traditionen mister ikke derved sin Betydning: den skal give os
Spørgsmaal at stille og Midler til at prøve paa at besvare dem,
men ikke færdige Resultater; den skal muliggøre, at vi kunne
stille nye Spørgsmaal. Og gaa vi i vore Antagelser ud over,
hvad Erfaringen lærer os, da tør det kun ske i Kraft af den
Forudsætning, at hvad vi ikke kende, har Lighed med, hvad vi
kende. --- Man har fundet disse Fordringer for strenge, naar
de stilles til alle Mennesker, og man har med Rette hævdet,
at uvilkaarlig Tro og Tillid nu engang gaa forud for Tvivl, og
at der kan høre lang Erfaring til, før en selvstændig Tvivlen
kan opstaa\footnote[1]{Disse Betænkeligheder ere fremførte af en Mand, der selv praktisk er et af de smukkeste Exempler paa Sandhedskærlighed (baade som Ærlighed, som Stræben efter personlig Sandhed og som intellektuel Redelighed) i vor Tid: Christoph Schrempf. Se hans Kritik af Clifford i hans Tidsskrift Die Wahrheit. Stuttgart 1894. I. p. 163 ff.}. Ogsaa fra en helt anden Side har man kritiseret
Clifford's Paastand, idet man har fremhævet, at Modet til at
vove er en Betingelse for at finde Sandheden. Hellere narres
mange Gange end miste Muligheden for at gætte rigtigt! Bærer
Clifford sig ikke ad som en General, der siger til sine Soldater,
at det er bedre at holde sig fra Kampen end at blive saaret?
Og man har draget en Analogi mellem de uvilkaarlige Indskydelser,
der ofte ligge til Grund ved Opdagelser, og de spontane
Variationer, der spille saa stor en Rolle under Kampen for Tilværelsen\footnote[2]{William James: The Will to believe. New York 1897. p. 18 f.; 93; 249. --- Smlgn. hermed Leslie Stephens Bemærkninger om Mill's Bog „Om Friheden“, i hvilken en med Clifford's beslægtet Anskuelse udvikles. The Utilitarians. III. p. 255—259. Stephen fremhæver især, at Forskningen er en stor, fælles Opgave, ved hvis Behandling den Ene maa have Hjælp fra den Anden. Mill isolerer de enkelte Individer for meget og lægger for stor Vægt paa, at selv vel begrundede Antagelser finde Modsigelse.}.
--- Men ikke blot overfor den kritikløse Tro, Kirken
% --- p. 256 ---
\setcounter{footnote}{0}%
\opage{256}har indskærpet som den højeste Dyd, ogsaa overfor den store
Popularisering af videnskabelige Resultater i fortællende og beskrivende
Form, er det af stor Betydning, at den intellektuelle
Redelighed kommer frem i Dydernes første Række. Det er Betingelsen
for, at det aandelige Arbejde ikke skal falde bort eller
blive undervurderet. Og særligt af dem, der i Kirke og Skole
skulle være offentlige Lærere, bør denne Dyd kræves i en ganske
anden Udstrækning, end det hidtil har været Tilfældet, --- og
kræves i Sandhedskærlighedens Navn.

Det intellektuelle Fremskridt hænger ikke blot meget nøje
sammen med Fremskridt paa alle andre Omraader, men det er
selv den Form for Fremskridt, som er lettest at paavise. Selv
de, der ellers ikke antage, at der finder noget Fremskridt Sted,
indrømme det dog paa Erkendelsens Omraade. Følelse og Villie
ændres og udvikle sig langsommere, og deres Fremgang er for
en meget stor Del bunden til Erkendelsens. Den en Tid gængse
Sætning, at Ideerne regere Verden, er ikke psykologisk rigtig;
Historiens Spil drives af Menneskenes Følelser og Lidenskaber;
--- men Tankernes Klarhed og Sandhed er en af de allervigtigste
Betingelser for Følelsers og Lidenskabers sunde Udvikling. ---

Sandhedsinteressen har det fælles med den etiske Følelse,
at den fører os til at se bort fra vor tilfældige Individualitet
og vore tilfældige Forhold for at spørge om det Almengyldige.
Den, der ikke kan taale Sandheden, kan ikke se ud over sin
egen isolerede Individualitet; derfor fører Sky for Sandheden
til Egoisme, hvis den ikke fra først af udspringer deraf. Den
religiøse Etik har ofte virket fordærveligt ved at sætte Tro som
det eneste Fortjenstlige og stemple Tvivl som Synd. Den har
ikke set, at en ærlig og fornuftig Tvivl baade er sund og god.
Den kommer let til at sætte en Præmie paa Dumhed; ti det
er ikke al Enfoldighed, der er „hellig Enfold“. Sættes der blot
én Skranke for Sandhedsbestræbelsen, saa vil den snart føre
flere efter sig. \textsc{Coleridge}, som selv var en troende Kristen,
har sagt: „Den, som begynder at elske Kristendommen mere
end Sandheden, vil siden elske sin egen Sekt eller Kirke mere
end Kristendommen og vil ende med at elske sig selv mere
end Alt“.

% --- p. 257 ---
\setcounter{footnote}{0}%
\opage{257}Det er ikke altid den glødende Tro, der fører til at betragte
Tvivl som noget Ondt. I vore Dage er det meget ofte Frygten
for det Ubekendte, for at mangle Fodfæste, naar man slipper de
vante Forestillinger, for den Uro og Møje. det koster at vinde
en ny Anskuelse. Overfor den, der ikke kan opgive en Mening,
fordi han i den finder sin eneste Trøst, sit eneste Tilhold, vil
Ingen rejse Tvivlen og Diskussionen. Man skal ikke berøve
den Fattige hans eneste Lam. Men derfra er der et stort Spring
til den Antagelse, at det beror paa vor Trang og vore Ønsker,
hvad der er Sandhed, --- et Spring, som dog foretages af Mange
med en forbavsende Lethed.

5. Deraf, at Sandheden altid skal søges, følger, som allerede
antydet, ikke, at den altid skal siges. Sandhedspligten gaar
ud paa at bringe Sandheden til at herske; men dette Formaal
kan ofte hindres ved at sige Sandheden. Ved at udtale Sandheden
for tidligt bærer man sig ad som Barnet, der trækker
Planten op af Jorden for at se, hvor dybt dens Rod er trængt
ned. Man konstaterer, hvad der er naat --- men kan derved
hindre, at der naas Mere! Ukaldede Sandhedsapostle synde ofte
paa denne Maade. Pligten at tale Sandhed er pædagogisk begrænset.
Sandheden er saa stor og omfattende, at Mennesket kun
Skridt for Skridt kan blive indviet i den. Den Del af Sandheden,
som det endnu ikke kan fatte, miskender og foragter det let.
Det gælder da at sige Sandheden saaledes, at denne Rest i hvert
enkelt Tilfælde bliver saa lille som muligt. Og det gælder at
finde den Form, i hvilken Sandheden bedst kan finde Indgang,
og den gradvise Afsløringsmaade, der kan sikre dens Forstaaelse.
Især gælder dette, naar Talen er om etiske „Sandheder“. Der
gives megen ukaldet Moraliseren i Verden, og det overses stadigt,
at Udtalelsen af etiske Domme selv er etisk betinget, og at al
Vurdering staar i et Formaals Tjeneste, er et Opdragelsesmiddel
(se III, 11). At udtale etiske Domme, hvor de ikke behøves,
er Barbari eller Farisæisme. Indirekte Opdragelse er ofte den
rette. Dog kan det undertiden netop være pædagogisk rigtigt
at rejse Forargelsen og Forbitrelsen og opstille et Modsigelsens
Tegn. Det er ofte kun gennem en lidenskabelig Kamp for og
imod, at Sandheden ret træder frem for os.

% --- p. 258 ---
\setcounter{footnote}{0}%
\opage{258}Hermed er Spørgsmaalet om Nødløgns Berettigelse væsentligt
afgjort. Ligesaa lidt som man er forpligtet til af sig selv
at sige Alt, hvad man tænker, ligesaa lidt er man forpligtet til
at besvare alle Spørgsmaal. Morderen, der leder efter sit Bytte,
eller den, der vil aflokke mig en Hemmelighed, som det er min
Pligt at bevare, er jeg baade berettiget og forpligtet til at lede
vild. Det Samme kan gælde i mange mindre Forhold. Den,
som afslaar en Indbydelse, kan ikke altid angive den virkelige
Grund, fordi Indbyderen ikke vilde forstaa den, eller fordi den
vilde saare ham. Men vi tage os det dog sikkert i det daglige
Liv altfor let i denne Henseende. Dels af utidig Høflighed,
dels af Magelighed, dels af Mangel paa Mod udstrække vi Brugen
af Nødløgn langt videre end nødvendigt og berettiget.

6. Naar vi fastholde, at Sandheden er i sin Vorden, blive
Tolerance og Pietet let forenelige med Sandhedskærligheden. Vi
befinde os Alle, hver for sig, paa forskellige Trin af Tilnærmelse
til Sandheden. En skarp Grænselinie mellem Nogle, som have
Sandheden, og Andre, som ikke have den, kan ikke drages.
Sandheden er saa mægtig og saa nøje forbunden med Virkeligheden,
at Ingen, selv med sin bedste Villie, fuldstændigt kan
isolere sig fra den. Det er jo Virkeligheden, Berøringen og
Vexelvirkningen med de virkelige Forhold, der opdrager til
Sandhed, og Ingen kan helt slippe for at faa Del i denne Opdragelse.
Men Opdragelsen kan være lang og besværlig, og
hvor langt den Enkelte naar, beror paa mangehaande Betingelser.

Tolerancen (for at beholde dette stygge Ord for en smuk
Ting) slaar ikke af paa Sandheden, men udspringer af Indsigt i
Erkendelsens haarde Tilblivelseskaar. Den, der slutter saaledes:
„der kan kun være én Sandhed; hvis jeg finder mig i, at Andre
mene og lære, hvad der strider mod min Overbevisning, indrømmer
jeg, at det, jeg antager, maaske ikke er Sandhed!“ ---
han ender som Forfølger. Det var i Kraft af denne Tankegang,
at den gamle Kirke tidlig kom ind paa Kætterforfølgelsens Veje\footnote[1]{Se allerede 2det Johannes' Brev 10—11: „Dersom Nogen kommer til Eder og fører ikke denne Lærdom, ham annammer ikke til Huse og byder ham ikke velkommen. Thi hvo som byder ham velkommen, bliver delagtig i hans onde Gerninger“. --- I det 2. Aarhundrede maatte ret-}.
% --- p. 259 ---
\setcounter{footnote}{0}%
\opage{259}Man kan meget godt være fuldt overbevist om en Sandhed og
dog indse, hvorfor Andre ikke antage den. Og indser man ikke
Aarsagen til, at de ikke antage den, vil man ikke kunne finde
den rette Maade at overbevise dem paa. Ved at forske efter
denne Aarsag, vil man tilsidst komme ud over Tænkningens
Omraade. Man vil da finde, ikke blot, at der er Følelser, som
kunne hindre den klare Erkendelses Udvikling, men ogsaa, at
der er Følelser, som kunne bestaa til en vis Grad uafhængigt af
Erkendelsens Udvikling. Netop de vigtigste Følelser --- Menneskekærlighed,
Samvittighed o. s. v. --- ere forenelige med de mest
forskellige teoretiske og religiøse Antagelser. Hvad Fritænkeren
betragter som et meningsløst Dogme, er for den Troende en
uundværlig Tanke, Udtryk for en Følelsestrang, den alvorlige
Fritænker paa sin Vis maaske ogsaa kender, skønt han giver
den Luft paa anden Maade. Ved at opsøge det Følelsesliv, der
ligger bagved Erkendelsen, vil man finde det Fælles og bane Vejen
for en Sympati, der paa Afstand kunde synes umulig. Tolerancen
hviler, hvor den er Mere end ligegyldig Passivitet, paa en Tro
paa den menneskelige Kærne i Alle, en Tro, der er Et med
Kærligheden selv.

Ligesom Tolerancen er Følelsen overfor dem, der samtidigt
med os kæmpe for Sandheden i en anden Lejr, saaledes knytter
Pieteten et Baand mellem os og Fortidens Repræsentanter.
troende Kristne ikke samtale med Banlyste og Kættere. (Lecky: History
of European Morality I, p. 451). --- Gregor fra Nazianz (i' 4. Aarh“)
erklærede, at man ikke kunde tillade Kætterne at holde Forsamlinger
uden at indrømme deres Læres Sandhed. (Barbevrac: Traité de la
Morale des Péres. p. 170). --- I det 17. Aarhundrede skyede mange ivrige
Lutheranere de Reformerte i den Grad, at de ikke vilde sidde til Bords
med dem eller bo i Hus med dem. (Gass: Die Lehre vom Gewissen.
p. 164). --- Dr. Johnson blev meget forarget over at høre, at en af Hume's
teologiske Modstandere havde spist til Middag med Hume, trykket hans
Haand og udvexlet Besøg med ham. \textasciitilde{}„En Vantro“, sagde han bør ikke
behandles pænt af en Kristen, blot fordi han bestræber sig for at
begaa sit Rov med en vis Oprigtighed“. (Denne Udtalelse, saavel som en hel
Række Udtalelser af Johnson i samme Tone om Rousseau, Voltaire, Hume
og A. Smith, findes i de af Griesebach udgivne Edita und Inedita
Schopenhaueriana. Leipzig 1888. p. 127---132). Exempler fra vore Dage
% CHECK p. 259: note whose mark was not found
\footnote[1]{ville ogsaa kunne anføres.}
% --- p. 260 ---
\setcounter{footnote}{0}%
\opage{260}Vi staa i Pietetsforhold til det, der har været Autoritet for os,
og hvis Skole vi uden Bitterhed have forladt. Vi have en Forstaaelse
af dets Betydning, som den ikke kan have, der kommer
som Fremmed og ser udefra derpaa. Vi forholde os til det som
til vor egen Fortid; selvom vi have brudt med den, give vi den
dog ikke ganske til Pris, saalænge vi fastholde Kontinuiteten i
vor Udvikling. Paa lignende Maade er Pieteten den Følelse, i
hvilken Slægtens Kontinuitet lever; den danner Baandet mellem
de skiftende Generationer. Ofte har netop den nærmeste Generation
ondt ved at bevare dette Baand. Det 18. Aarhundrede
havde ikke stor Pietet overfor Middelalderen; den kom først
(og i Overmaal) med det 19. Aarhundredes Romantik, som til
Gengæld ikke havde stor Pietet for det 18. Aarhundrede; og i
vor Tid, der er Romantiken fjendtlig, føler man sig saa igen
mere pietetsfuld overfor det 18. Aarhundrede. --- Pieteten kan
føles baade overfor enkelte Personer og overfor en hel Retning
eller en hel Slægt. Den, der i sit Forhold til de Autoriteter,
hvis Herredømme han har unddraget sig, endnu ikke er naat
saa vidt, at han kan erkende og sympatisere med deres historiske
Betydning, han er endnu ikke fri, men er afhængig af sin Modsætning
til Fortiden. Pieteten er den sande historiske Følelse.
Alt, hvad der har haft Værdi, lever i den et fortsat Liv. Den
er en af de vigtigste Former, under hvilke den Enkelte lever
sig ind i Slægten.
% --- p. 261 ---
\setcounter{footnote}{0}%

\opage{261}\partitle{SOCIAL ETIK}

% --- p. 263 ---
\setcounter{footnote}{0}%

\opage{263}\chapterhead{XIII.}{INDLEDNING OG INDDELING.}

1. Den sociale Etik adskiller sig fra Sociologien (Kulturhistorien)
paa lignende Maade, som den individuelle Etik adskiller
sig fra Psykologien. Medens Psykologien undersøger de
menneskelige Forestillinger, Følelser og Villiesytringer som rent
naturhistoriske Fænomener og blot vil forstaa dem i deres Oprindelse
og Udvikling, er det den individuelle Etiks Sag at
vurdere dem efter deres Forhold til Idealet for Individets Liv.
Og medens Sociologien studerer det menneskelige Samfundsliv
i dets forskellige Former, til de forskellige Tider og hos de forskellige
Folkeslag og stiller sig overfor disse Former som Naturforskningen
overfor Naturfænomener, vil den sociale Etik vurdere
dem efter deres Forhold til Idealet for det menneskelige Samfundsliv.
Etiken benytter overalt det Givne som Udgangspunkt
for en ny Udvikling, og naar den anstiller en Vurdering af det,
der rører sig i Samfundet, er det, fordi denne Vurdering øver
en bestemmende Indflydelse paa Udviklingens fremtidige Gang.
Hvor langt denne Indflydelse kan strække sig, kan kun Erfaring
lære. Saasnart vi vurdere, forudsætte vi Muligheden af at gribe
ændrende ind. Det gælder ikke blot for Etiken, men for enhver
praktisk Videnskab. Naar den almindelige Retslære ikke blot
undersøger, hvilke Forhold der ere ordnede retsligt, men ogsaa,
hvilke der egne sig dertil, saa ligger deri et Ideal, som kan
behøve megen Kamp for at virkeliggøres. Nationaløkonomien
% --- p. 264 ---
\setcounter{footnote}{0}%
\opage{264}undersøger ikke blot den faktiske Produktion og Fordeling, men
kritiserer denne, forsaavidt den ikke stemmer med Betingelserne
for et sundt og lykkeligt Samfundsliv. Baade Retslære og Nationaløkonomi
føre tilbage til Etiken som deres almindelige
Grundlag. (Smlgn. III, 1. 14). —

Ligesom enhver Karakteregenskab og enhver Villiesakt, som
den individuelle Etik fordrer, maa være psykologisk mulig, saaledes
maa ogsaa enhver Samfundsordning, som den sociale Etik
fordrer, være sociologisk (historisk) mulig. Dog kan et etisk
Ideal, selvom det ikke i lange Tider eller maaske nogensinde
kan virkeliggøres, have stor Betydning ved den Retning, det
giver det menneskelige Sind og de menneskelige Kræfter, naar
blot fantastisk Overspændthed og Miskendelse af de virkelige
Forhold holdes borte. De nye Udviklingsmuligheder opdages
ofte kun af den, hvis Blik ikke er fængslet til det øjeblikkeligt
Givne, men som formaar at hæve sig over det. Det er et
vigtigt Vendepunkt, hvor bevidst Sætten af Formaal og Idealer
afløser den blinde, instinktmæssige Handlen. Dette Vendepunkt
naas gennem historisk Udvikling, og den bevidste Vurdering,
som derved bliver mulig, virker igen tilbage paa den fortsatte
Udvikling.

2. Fra en dobbelt Side set faar al Etik, særligt al social
Etik, en historisk Karakter. Det er allerede tidligere (III, 13)
fremhævet, at den vurderende Subjektivitet stedse gaar ud fra
et psykologisk Grundlag, som ikke er det samme til de forskellige
Tider. Men her se vi nu, at tillige det Stof og de Genstande,
den etiske Vurdering og det etiske Arbejde angaar og
behandler, frembyde forskellige Egenskaber til de forskellige
Tider. Udfra det samme Grundlag vil man altsaa til forskellige
Tider naa et forskelligt Resultat. Hvad der er muligt paa ét
Udviklingstrin, er ikke muligt paa et andet. Man kan ikke
udfra det almindelige Velfærdsprincip konstruere en Ordning af
Familie eller Stat, som uden videre skulde kunne indføres overalt.
En hvilkensomhelst Samfundsordning forudsætter visse indre og
ydre Betingelser givne hos dem, der skulle udgøre Samfundets
Medlemmer. Paa bar Grund kan ingen social Bygning opføres.
Det er, hvad Historien noksom viser, lettere at omstøde gamle
% --- p. 265 ---
\setcounter{footnote}{0}%
\opage{265}Institutioner og „indføre“ nye end at ændre de Anlæg og Drifter
i den menneskelige Natur, som gav de gamle Institutioner
Livskraft. Vi have paa det aandelige og sociale Omraade altid
gamle Læderflasker til den nye Vin, og denne faar altid en vis
Smag efter dem. Ofte maa den nye Vin sprænge de gamle
Læderflasker, før de kunne blive erstattede med nye.

Den historiske Relativitet, som saaledes nødvendigvis præger
ethvert af Etikens Resultater, har man, efter at man én Gang
havde faaet Øjnene op for den, ofte overdrevet og misforstaaet.
Den betyder ikke, at Alt kan være lige godt eller lige slet. Den
betyder, at hvad der skal være godt under visse bestemte Forhold,
maa svare til dem. De etiske Idealer individualiseres paa
forskellige Maader til forskellige Tider og paa forskellige Steder;
og dette er saa langt fra at stride mod deres Væsen, at et Ideal,
som ikke kunde individualiseres, vilde være tomt Fantasteri.
Men de etiske Idealer have mere Karakteren af almindelige
Tendenser og Retninger end af Formler, som uden videre kunne
anvendes under hvilke Forhold som helst. Hvad vi kunne opnaa,
er, at vore Handlinger og de Former, i hvilke Samfundslivet
ordnes, komme til at ligge i en Række, der forløber i Retning
af Idealet, og saa langt som muligt fremme i denne Række. Vi
komme ikke udover den stadige Stræben. Relativiteten betyder
netop, at kun en stadig Tilnærmelse er mulig. (Smlgn. III, 12;
IV, 3—4 og VII, 3). Herved maa det ikke overses, at Idealerne
selv under den fremskridende historiske Udvikling undergaa Ændringer,
idet der baade opdages nye Muligheder, nye Fordringer
--- og nye Skranker\footnote[1]{Smlgn. nærmere om den historiske Relativitet: Etiske Undersa «eiser p. 47---51.}.

3. Den etiske Udviklingsproces er dels en Humaniseringsproces,
dels en Emancipationsproces. --- Paa de lavere Udviklingstrin
herske de dyriske Instinkter, til hvilke Individets og
Slægtens Opholdelse er knyttet. Og de høre ikke op at virke,
fordi ideelle og sympatiske Følelser udvikle sig. De have endog
i deres Natur en Voldsomhed, mod hvilken de højere og mildere
Følelser ofte ere afmægtige. Naar højere Interesser eller
% --- p. 266 ---
\setcounter{footnote}{0}%
\opage{266}Værdier --- det vil sige: saadanne, der pege ud over Øjeblikket,
ja maaske over Individets og Slægtleddets hele Liv --- gøre sig
gældende, sker dette ikke ved, at der skabes ny Energi. Etiken
har ingen Grund til at fornegte Energisætningen. Men Energien
kan omsættes fra sine mere elementære Ytringsformer til
Former, i hvilke den kan tjene højere Værdier, og tillige vil der
kunne udløses Energi, som der ikke var Brug for i de mere
begrænsede Interessers Tjeneste. Det gælder om at overføre
Energien, som fra først af tages i Beslag af de dyriske Instinkter,
til de ideelle, mere menneskelige Følelser, og derved vil der
ogsaa blive Anledning til, at hidtil slumrende Kræfter vækkes.
Den strenge Askese er her et revolutionært Princip, idet den
uden videre forlanger hine Instinkter dødede. Det er først
gennem en langsom Udviklingsgang, at Instinkterne under Livsforholdenes
Indflydelse, der ingenlunde gør bevidst Villiesanspændelse
overflødig, kunne formildes og forædles. Foruden
denne Humaniseringsproces, ved hvilken vi fra Dyr efterhaanden
ere blevne og blive Mennesker, foregaar der ogsaa en stadig
Emancipationsproces, ved hvilken vi efterhaanden blive frie Mennesker.
Legemlig og aandelig Tvang er karakteristisk for de
lavere Trin af etisk og social Udvikling; Autoritetsforholdet
hersker i dets strengere eller mildere Former. Heller ikke det
kan --- som det 18. Aarhundrede i sin revolutionære Iver mente
--- pludseligt afskaffes. Gennem Autoritetens opdragende Indflydelse
og gennem det frie Valg af Autoriteter som Mellemled
foregaar den selvstændige Personligheds Udvikling.

Ethvert Individ, enhver Generation og ethvert Folk befinder
sig paa et vist Stadium af denne Udviklingsproces (i dens
dobbelte Form), og alle Institutioner, Sæder og Livsformer ville
faa Præget deraf, hvorledes man saa bærer sig ad. Hvad det
kommer an paa, er, hvilke Institutioner og Livsformer der bedst
kunne tage de forhaandenværende Kræfter og Drifter i deres
Tjeneste, saaledes at de tillige føre Udviklingen opad mod et
højere Stadium.

4. Paa et absolut individualistisk Standpunkt kan Samfundet
kun opfattes som opstaaet ved Sammenslutning af Individer.
I Individernes mere eller mindre bevidste Tanker og Følelser
% --- p. 267 ---
\setcounter{footnote}{0}%
\opage{267}finder man da Forklaringen af Samfundslivets Art og Form.
Man forudsætter selvstændige Individer givne forud for Samfundets
Opstaaen. Men ethvert Individ forefinder faktisk et
Samfund, indenfor hvilket det udvikler sig, og først efterat
Individets Udvikling har faaet sit bestemte Præg ved dette
givne Samfunds Forhold og Traditioner, kan Individet selv igen
virke bestemmende tilbage paa Samfundet. Før det bliver sig
klart bevidst, lever og virker det, tænker og føler det kun som
Led af Samfundet. Det Uvilkaarlige gaar stedse forud for det
Vilkaarlige. Der foregaar en ved Samfundsforholdenes ubevidste
Indflydelse betinget Opdragelse af Villien, der er nok saa vigtig
som den Opdragelse, der gaar for sig gennem bevidst Indgriben
og Paavirkning.

Ogsaa dette hører til de psykologiske og etiske Sandheder,
som \textsc{Aristoteles} allerede har haft et klart Blik for. Hans
Tankegang er omtrent følgende:\footnote[1]{Eth. Nic. li, 2.} Egenskaber (f. Ex. etiske
Dyder) frembringes ved, at man udøver de tilsvarende Handlinger.
Ved Bygningsarbejde blive vi dygtige Bygmestre, og
ved at øve Retfærdighed blive vi retfærdige. Hvad vi bevidst
skulle udføre med Færdighed, maa vi først indøves i. Derfor
er det af stor Betydning, om Samfundsordningen straks fra først
af fører den Unge i det rette Spor, leder ham til Virksomhed
i den rigtige Retning. Samfundsordningens etiske Betydning
beror paa, hvilke Vaner den bibringer de Enkelte. --- Vi ville
i det Følgende meget ofte faa Brug for denne Sætning: at uvilkaarlig
Indøvelse gaar forud for fuldt bevidst Udøvelse, og ville
for Kortheds Skyld kalde den det aristoteliske Princip. Dette
Princip finder Anvendelse forud for det tilsyneladende modsatte
Princip, at hvad der først øves med Bevidsthed, senere kan
øves ubevidst, som Vane, eller som Instinkt. Ti fra først af
véd Individet ikke, hvorfor det indøver visse Forestillinger eller
Handlinger; det sker uvilkaarligt og tildels ubevidst. Det Uvilkaarlige
gaar stedse forud for det Vilkaarlige, og selve Overgangen
fra det Uvilkaarlige til det Vilkaarlige sker uvilkaarligt.
„Villiens Førstetanke“ (se VI, 1) opdukker uden Forsæt. Vi
% --- p. 268 ---
\setcounter{footnote}{0}%
\opage{268}raade ikke for vore første Forestillingers Opstaaen, hvorpaa igen
vor bevidste Villen, vor vilkaarlige Handlen beror.

Ligesom der saaledes er en stadig Vexelvirkning mellem
det Bevidste og det Ubevidste, og mellem det Vilkaarlige og det
Uvilkaarlige, saaledes foregaar der ogsaa en Vexelvirkning mellem
Villien paa den ene Side, dens Midler og dens Resultater
paa den anden Side. Hvad Villien først har grebet som blot
Middel, kan senere blive til Formaal for den, og det Resultat,
i hvilket den mente at havne, kan blive Udgangspunkt for ny
Stræben. Der kan foregaa Motivforskydninger, og der kan
foregaa Værdiforskydninger.

Ifølge Motivforskydningens Lov kan det, der først udføres
af et Motiv, senere udføres af et helt andet, idet det oprindelige
Middel er blevet til Formaal. Billedligt kan dette udtrykkes ved,
at Individets Interesse har flyttet sig. (Smlgn. I, 4). Da det nu
ikke er sagt, at det oprindelige Motiv ganske forsvinder, kan
der med Forskydningen forbinde sig en Sammensmeltning, en
Art kemisk Sammensætning af forskellige Motiver, ved hvilken
der opstaar et nyt Motiv, en ny Følelse, hvis Egenskaber ikke
umiddelbart kunne udledes af de tidligere Følelser, af hvilke
den er opstaaet. Ved Hjælp af disse psykologiske Love\footnote[1]{Smlgn. om disse psykologiske Love min Psykologi VI B, 2 d; C, 2 og 5; E, 4—5 (kfr. II, 6 d; VB, 2). --- Motivforskydningens Lov er især paavist af Spinoza, Hartley og James Mill, Sammensmeltningsfænomenet af Hartley og F. C. Sibbern.} forstaas
det, at de sociale Livsinteresser umiddelbart kunne staa
for Individet som de højeste og blive det centrale Synspunkt,
til hvilket han i sin Vurdering og i sin Livsplan tilsidst vender
tilbage. Det er et saadant Standpunkt, som vi her have valgt
som Grundlag for vor systematiske Vurdering. (Smlgn. III, 9—13).
Kun paa et saadant Standpunkt følger det af sig selv, at den
sociale Etik omfatter den individuelle. Den Proces, hvorved
dette Standpunkt bliver til, er nu af stor Interesse for selve den
sociale Etik, da det naturligvis maa være en Hovedopgave for
denne at forstaa det Grundlag, af hvilket dens egen principielle
Stilling afhænger.

% --- p. 269 ---
\setcounter{footnote}{0}%
\opage{269}Erfaringen viser fremdeles, at der af de Enkeltes uvilkaarlige
eller vilkaarlige, egenkærlige eller opofrende Handlinger
udspringe langt flere og ofte helt andre Virkninger, end der
blev forudset eller kunde forudses, endsige tilsigtes af de Handlende.
Af de individuelle Interessers og Lidenskabers Vexelvirkning
kunne udspringe Resultater af blivende og omfattende
Betydning. Ofte staar Noget som et sidste afsluttende Formaal,
der begrænser Villiens Horizont, og det kan saa dog bagefter
vise sig at kunne bruges som Middel til endnu fjernere Formaal
eller at besidde Værdier af anden Art og større Rækkevidde,
end man forud havde tænkt sig. Der gives jo baade Kræfter,
som ville skade og dog komme til at gavne, og Kræfter, som
ville gavne og dog komme til at skade. Men Hovedsagen er
den Udvidelse af Værdiernes Horizont, som kan ske, og som
stadigt sker i Historien. Et Exempel frembyder Humanitetsfølelsens
Udvikling efter Alexanders Erobringer. For Alexander
var Stiftelsen af et Verdensrige det sidste Maal; men det viste
sig, at da dette Maal var naat, var der derved tilvejebragt Betingelser
for en saadan Vexelvirkning, et saadant Samkvem
mellem de forskellige, hidtil adskilte Folkeslag, at en almindelig
Menneskelighedsfølelse kunde udvikle sig. De Værdier (potentielle
eller aktuelle), der frembringes i Løbet af den menneskelige
Udvikling, bero for en meget stor Del paa denne objektive
Værdiforskydning (som man kunde kalde den til Forskel fra
Motivforskydningen som en subjektiv Værdiforskydning). I den
geniale Indledning til sin Philosophie der Geschichte har \textsc{Hegel}
paavist dette Fænomen, „at der i Verdenshistorien gennem
Menneskenes Handlinger bliver noget Andet til foruden det, de
tilsigte og opnaa, ja overhovedet vide og ville“. Han kalder
det Fornuftens List, at den [Verdensfornuften] lader Lidenskaberne
virke for sig: „Ideen [Verdensfornuften] betaler Tributen til
Livet og Forgængeligheden ikke af sin egen Lomme, men ud
af Individernes Lidenskaber“.\footnote[1]{Denne ubevidste Produktion af mere omfattende Værdier kalder Wundt med et fra Zoologien laant Udtryk „die Heterogonie der Zweeke“ Ethik.\textsuperscript{3} I p. 275.} De Resultater, der saaledes
% --- p. 270 ---
\setcounter{footnote}{0}%
\opage{270}realiseres midt igennem Menneskenes bevidste Arbejde, vedkomme
dog kun Etiken, forsaavidt de virkeligt kunne kaldes
Værdier og som saadanne træde i Forhold til Villie og Følelse
(III, 1 og 14). Hvis man strækker Betragtningen saa vidt, at
de sidste Resultater, Endemaalene for al menneskelig Udvikling
skulle ligge ud over Alt, hvad der muligvis kunde blive Formaal
for menneskelig Villen og staa som værdifuldt for menneskelig
Følelse, saa har man overskredet Etikens Grænse og er gaaet
over til Historiens Filosofi og Religionsfilosofien (eller, om man
vil, Metafysiken). Det er da et transscendent, al Erfaring overskridende
Resultat, der staar som Endemaalet for den menneskelige
Udvikling. I denne Retning gaar \textsc{Hegels} Opfattelse og
i den nyeste Tid \textsc{Wundts}; dog indrømmer denne sidste, at
hint transscendente Endemaal ikke er eller kan blive os bekendt.
Men selvom man ikke vil forlade Etikens Omraade og gaa over
i den spekulative Tænknings Regioner, indeholder Værdiforskydningens
Lov en Tanke af stor Betydning, især overfor en kortsynet
og stakaandet Opfattelse af Velfærdsprincipet.

Ved Motivforskydningen er Sammenhængen mellem det
Gamle og det Nye nøjere end ved den objektive Værdiforskydning.
Hist er der en Handlemaade eller en Livsordning, der i
det Væsentlige vedbliver at være som den var, men fastholdes
og anerkendes af nye Grunde; her er det en helt ny Handlemaade
eller Livsordning, der opstaar, og de Grunde, som gøre
den værdifuld, kunne saa være saadanne, der forhen slet ikke
kendtes. Det sociale Resultat, som Alexanders Erobringer hidførte
ved Nedbrydelsen af de nationale Skranker omkring det
østlige Middelhav, kunde ikke blot ikke forudses, men heller
ikke anerkendes fra det klassisk græske Standpunkt. Fælles for
de to Arter af Forskydning er det, at Anerkendelsen af de nye
Værdier ikke nødvendigvis medfører Anerkendelse af den Maade,
paa hvilken de ere blevne til. Anerkendelse af Privatejendommens
Betydning medfører ikke Anerkendelse af den Maade, paa
hvilken Privatejendom historisk har udviklet sig, og Humanitetsfølelsen
behøver ikke at føre til ubetinget Beundring for Alexanders
Erobringslyst. Der kan her opstaa en skarp Modsætning
mellem den etiske og den historiske Betragtningsmaade. Løs%
% --- p. 271 ---
\setcounter{footnote}{0}%
\opage{271}ningen af denne Strid maa søges deri, at den etiske Vurdering
efter sin Natur er rettet mod Fremtiden, mod det, der skal ske.
Den tager Livet op paa det Grundlag, som nu en Gang faktisk er
bragt tilveje. Ligesaalidt som Hensigten altid helliger Midlet, ligesaalidt
behøver der fra de nye, opnaaede Værdier at straale
nogen Glans tilbage paa de Veje, ad hvilke de bleve naaede.
Naar Mefistofeles vil det Onde, men stedse kommer til at frembringe
det Gode, saa takke vi ham ikke derfor; vi kunne i det
Højeste finde os i, at der gives den Slags Fyre. Men Etiken
bryder sig heller ikke om at holde Dom over Fortiden; den
betragter det Frembragte som den Jordbund, hvorpaa der nu
skal arbejdes. Efter Jordrevolutioner, der have ført til Sænkninger
og Hævninger af Jordlagene, breder Plantevæxt og Dyreliv
sig over det nye Terræn, hvor ujevnt og omkalfatret det
end maatte være. Det staar ikke i Livets Magt at gøre Revolutionen
om igen (selvom maaske meget Herligt er gaaet til Grunde);
men Livet formaar ved sin stille Magt at frembringe en ny
Verden paa den gamles Ruiner, --- naar kun Livsspirerne have
holdt sig trods Katastrofen. I Analogi hermed maa Forholdet
mellem det etiske Liv og de historiske Verdensprocesser opfattes.

Motivforskydningen taler snarest til Grund for den konservative
Opfattelse, den objektive Værdiforskydning til Gunst for den
revolutionære Opfattelse, idet hin viser Muligheden af nye Grunde
til Anerkendelse af det Gamle, denne Muligheden af, at det
pludseligt opstaaede Nye ogsaa vil finde nye Vurderingsmotiver.

Der kan være en indre Sammenhæng mellem de to Arter
af Forskydning, idet Betingelsen for, at et Middel bliver Formaal,
kan bestaa i, at Resultatet foregribes under Arbejdet og
derved giver dette Del i den Værdi, det selv har. En Vej, vi
gaa til et glædeligt Maal, kan selv blive glædelig. Men det
gælder ikke for alle Tilfælde. Ofte er det selve Arbejdet, der
bliver tiltrækkende ved den Brug af Kræfterne, det medfører.
Det er da ikke Resultatet, der giver Arbejdet Del i sin Værdi,
men det er Arbejdet, der selv vinder en ny Værdi *). Men om

') Ehrenfells (System der Werttheorie. I. p. 137 ff.) betegner det,
% CHECK p. 271: note whose mark was not found
\footnote[1]{jeg kalder Motivforskydning, som „Wertbewegung nach abwarts“, o: til- bagegaaende Værdibevægelse, i Modsætning til den objektive Værdiforskydning som „Wertbewegung nach aufwarts“. Men det er ikke al Motivforskydning, der sker „abwårts“. --- Smlgn. iøvrigt min Anmeldelse af Ehrenfells' Bog i \textasciitilde{}G6ttinger gelehrte Anzeigen“. 1900. Nr. 9.}
% --- p. 272 ---
\setcounter{footnote}{0}%
\opage{272}begge disse Underarter af Motivforskydning gælder det, at Kontinuiteten
er større ved dem end ved den objektive Værdiforskydning.

Fælles for det aristoteliske Princip og de to Forskydningsformer
er det, at den bevidste Anerkendelse kun kommer bagefter.
Der sættes derved snevre Grænser for enhver deduktiv
Etik; Erfaringens Nødvendighed indskærpes med Eftertryk. Og
tillige medføre alle tre Processer en afgørende Vanskelighed for
Forhandling mellem Standpunkter, af hvilke det ene ligger forud
for, det andet efter en eller flere af de beskrevne Metamorfoser.
(Smlgn. Kap. III).

5. Man har ofte ment at finde Vejledning til Forstaaelse
af Forholdet mellem Individ og Samfund i Analogien mellem
Samfundet og en Organisme. Under forskellige Former gennemføres
og anvendes denne Analogi hos Platon, Hobbes og Spencer.
Samfundet betragtes da som et stort Menneske eller som
en Organisme, hvis forskellige Organer eller Celler de enkelte
Individer ere. Der vil ogsaa, som især \textsc{Spencer}'s Sociologi
viser, kunne anstilles mange interessante Betragtninger ved Hjælp
af denne Analogi. Men den har sin bestemte Grænse, som vi
her ville dvæle ved\footnote[1]{Den fremdrages iøvrigt allerede af Spencer: Principles of Sociology. I, p. 478 ff.}, da Forholdet mellem individuel og social
Etik derved nærmere belyses. (Smlgn. III, 17 og VIII).

I Organismen ere de enkelte Elementer ubevidste, og kun
til de nervøse Centralorganer er der knyttet Bevidsthed. Denne
Centralisation er desmere fremtrædende, jo højere Organismen
staar. De enkelte Cellers eller Organers Tilstand og Virksomhed
vurderes tilsidst efter deres Indflydelse paa disse Centraldeles
Tilstand, der i Bevidstheden fremtræder som Lyst eller Smerte.
Idet Centraldelene kunne siges at repræsentere hele Organismen,
blive altsaa alle Celler og Organer underordnede Midler
for den hele Organismes Velfærd. --- I Samfundet er det
% --- p. 273 ---
\setcounter{footnote}{0}%
\opage{273}derimod netop alle Elementerne, de enkelte Medlemmer, som
have Evne til at føle Lyst og Smerte; denne Evne er ikke indskrænket
til et Centrum. Kun ved fantastisk Mystik kan man
tillægge Samfundet som Helhed Bevidsthed, afset fra de enkelte
Individer. Samfundets Velfærd er de enkelte Individers Velfærd.
Samfundet er en Forbindelse af personlige Væsener, men ikke
selv et personligt Væsen.

Kun tilsyneladende indeholder denne Betragtning en Modsigelse
mod Underordningen af den individuelle Etik under den
sociale (VIII). Skønt Samfundet og Slægten bestaar af Individer
og paa hvert givet Sted kun træder frem for os gennem visse
bestemte Individer, have dog Begreberne Samfund og Slægt den
store Betydning, at de overfor det enkelte Individ repræsentere
Indbegrebet af alle Individer (det selv indbefattet som et af de
mange). Begreberne Samfund og Slægt have, som tidligere
(VIII, 4) bemærket, den etiske Betydning at pege dels paa det
Uoverskuelige i de Livsinteresser, der skal tages Hensyn til, dels
paa Nødvendigheden af at tilvejebringe Muligheder, Betingelser,
potentielle Værdier, og ikke blive staaende ved det, der kun har
aktuel Værdi og maaske kun har denne en stakket Tid. De
muliggøre altsaa det mest omfattende Synspunkt for Vurderingen
af det enkelte Individs Villen og Handlen. De fordre, at
Individets Villen og Handlen ikke blot betragtes fra dets eget
Standpunkt, men ogsaa fra et højere, ligesom Kopernikus fordrede,
at man ikke blot skulde betragte Jorden fra dens eget
Standpunkt, hvor Alt drejer sig om den, men fra Solens Standpunkt,
hvor den selv ses at være med i Bevægelsen.

6. Den sociale Etiks højeste Ide er Ideen om Humanitetens
Rige, et Samfund af harmonisk og rigt udviklede Personligheder.
Et saadant Samfund er des fuldkomnere, jo mere ejendommelig
og selvstændig hver enkelt Personlighed er, og jo
inderligere og fastere Forbindelsen mellem de enkelte Personligheder
er, --- jo mere altsaa paa én Gang Mangfoldigheden og
Enheden komme til deres Ret. (Smlgn. III, 10). Denne Ide udspringer
af Velfærdsprincipet. Ti den største Velfærd maa være
til Stede, hvor hvert enkelt Individ udvikler sig paa selvstændig
Vis og derved tillige baade bevidst og ubevidst hjælper Andre
% --- p. 274 ---
\setcounter{footnote}{0}%
\opage{274}til en lignende Udvikling ud fra deres Stade. Den individuelle
Etiks Hoveddyder: Retfærdighed, Selvhævdelse og Hengivelse,
udspringe umiddelbart heraf. Og ikke blot muliggøres herved de
Enkeltes Karakterudvikling og deres indbyrdes Harmoni. Der
muliggøres ogsaa Frembringelsen af Resultater, som de Enkelte
hver for sig ikke vilde kunne naa, og som dog ere af Betydning
for Slægtens Udvikling idethele. Samfundets Ordning og
Kulturens Værker ere saadanne Resultater. De blive kun forstaaelige
derved, at der --- i Kraft af det aristoteliske Princip,
af Motivforskydningen og af Værdiforskydningen (4) --- kan
komme Andet og Mere ud af Individernes Arbejde, end disse i
deres Filosofi tænke paa.

Ideen om Humanitetens Rige dannes ved simpel Kombination
af de to Begreber Samfund og Velfærd. Men de speciellere
Former for det etiske Samfundsliv kunne ikke udledes af den.
Vi kunne bruge den som Maalestok ved Vurderingen; men da
Udviklingsmulighederne ere forskellige paa de forskellige Trin,
ville ogsaa Fordringerne i det Enkelte kunne blive meget forskellige.

De forskellige Trin af menneskelig Udvikling adskille sig,
hvad Samfundslivet angaar, ikke mindst derved, at der paa lavere
Trin kun findes liden Forskel mellem forskellige Arter af Samfund;
det er her som paa andre Omraader Differentiationen,
der især er karakteristisk for de højere Trin.

7. De forskellige Arter af Samfund, som man paa højere
Udviklingstrin kan skelne imellem, adskille sig fra hverandre
dels ved BeskafFenheden af de Kræfter, som forbinde Individerne
i Samfundet, dels ved de Formaal, som tilstræbes i Samfundet,
dels ved den større eller mindre Kreds af Individer,
som Samfundet omfatter.

I Familien er det den instinktmæssige Sympati, som er
det Forbindende. I Forholdet mellem Moder og Barn, Familiens
Kærne, viser den sig tydeligst og stærkest. Det er Blodets Baand,
som her knytter Menneskene sammen. Men ogsaa andre Motiver
end de elementære sympatiske Instinkter virke i Familien
paa dens forskellige Udviklingstrin. Paa primitive Udviklingstrin
betragtes Kvinder og Børn overvejende som Mændenes Ejen%
% --- p. 275 ---
\setcounter{footnote}{0}%
\opage{275}dom, og som underlagte deres Herredømme; og forsaavidt minder
den primitive Familie om Staten. Den minder om Kirken, forsaavidt
som den har sin Kultus og sine Ofringer til Forfædrenes
Aander. Og den er et Samfund til Kulturens Fremme, forsaavidt
som den ikke blot sørger for fysisk Pleje og Opholdelse,
men desuden for Udviklingen af Evner og Færdigheder. Ogsaa
paa højere Udviklingstrin er Familien ikke blot den stadige Kilde
til Næring for de sympatiske Følelser, men ogsaa den første
Magt, der fører Mennesket ind i Kulturen og vækker Bevidstheden
om en fast social Orden.

En anden Art Samfund opstaar, naar fælles eller dog forbundne
Interesser og Formaal forene Menneskene. Den Ene opnaar
ofte kun sine Formaal, naar han hjælper den Anden at
opnaa sine. Saaledes er Forholdet mellem Sælger og Køber, ved
al Ombytning. Til den Enes Tilbud svarer den Andens Efterspørgsel;
den Ene har tilovers, hvad den Anden mangler og
ønsker. Her supplere Individerne altsaa hverandre. Deres Interesser
ere forbundne, skønt der ellers ikke behøver at være nogen
nærmere eller yderligere Forstaaelse og Sympati imellem dem.
--- Men ofte ere Formaalene fælles og kunne kun naas ved forenede
Kræfter. Sikkerhed f. Ex. er et fælles Gode, som naas
ved Sammenslutning. --- Samvirken af det ene eller det andet
af disse Motiver frembringer (i Kraft af det i 4 Bemærkede)
efterhaanden en Følelse af Fællesskab, en Samfølelse, der kan
blive til uinteresseret Sympati, selvom det fra først af er Egoisme,
der har drevet Individet til at søge Samkvem med Andre.
Det paa Interessernes indbyrdes Suppleren eller paa Interessernes
Fællesskab beroende Samkvem kan saaledes virke opdragende.
Det er af stor Vigtighed, at allerede de ydre Livsvilkaar
virke forbindende. --- Det er dog ikke blot egoistiske Formaal,
som paa denne Maade kunne grundlægge Samkvem og
Samfund. Det er ogsaa saadanne Formaal, som ikke umiddelbart
spille en Rolle i den egoistiske Kamp for Tilværelsen. Videnskab
og Kunst, Religion og Menneskekærlighed sætte Menneskene
i fælles Bevægelse og føre dem sammen om fælles og
uselviske Formaal. --- I begge Tilfælde, baade hvor det fra først
af er de Enkeltes Kamp for Selvopholdelse, der fører dem sam%
% --- p. 276 ---
\setcounter{footnote}{0}%
\opage{276}men, og hvor det er ideelle Interesser, der er det Forbindende,
kan der opstaa en Broderskabsfølelse, dels bestemt ved den personlige
Sammenslutning, dels blot ved Bevidstheden om at staa
i de samme store Formaals Tjeneste. --- Denne Art Samfund,
hvor den forbindende Kraft ligger i Formaalene, og hvor hverken
elementære sympatiske Instinkter eller ydre Tvang raade,
kunne vi kalde det frie Kultursamfund. Dets Grænser ere ikke
saa snevre som Familiens; ja, de strække sig langt videre, lige
saa langt som Kulturen plejes; og det vil igen sige, at de tilsidst
falde sammen med Menneskeslægtens Grænser.

Endeligt gives der et Samfund, som adskiller sig fra de to
foregaaende ved, at det ikke blot bygger paa naturlig Sympati
eller paa Interessernes forbindende Kraft, men i sidste Instans
støtter sig til Magtanvendelse og Tvang. Hvad Formaal angaar,
har Staten, eller kan den have, megen Lighed med Familien
og det frie Kultursamfund. Den vil ikke blot opholde og sikre
Livet for sine Medlemmer, men ogsaa virke for deres Fremgang
i Kultur, baade materiel og ideel Kultur. Men overalt, hvor den
virker, staar Magten i Baggrunden. I Retten, som gælder i Staten,
haves Indbegrebet af de for Magtanvendelsen fastsatte Regler.
Ogsaa indenfor Familien og det frie Kultursamfund kunne
vi tale om Ret, naar vi derved forstaa Sædvane og Skik, saaledes
som de under Udviklingens Gang umærkeligt have dannet
sig og uvilkaarligt føre til at behandle nye Forhold og Tilfælde
paa lignende Maade som de tidligere af lignende Art. Historisk
finde vi vel neppe Familie eller Kultursamfund uden ydre Tvang
af en eller anden Art, og Staten betragter det paa den anden
Side som en af de vigtigste Opgaver at værne om Kulturen og
Familien. De tre Samfund kunne ikke holdes aldeles ude fra hverandre,
men betegne forskellige Synsmaader for det menneskelige
Samfund. Staten har ganske vist Magten som sidste Argument,
men ogsaa kun som det sidste. Den søger at drage Fordel
af de samme Kræfter, som raade i Familien og Kultursamfundet;
netop derfor er det den ikke ligegyldigt, hvorledes disse udvikle
sig, og dens egen Eksistens bliver saameget des fastere grundlagt,
jo mindre det er den blotte Magt, som raader. I Omfang
skiller Staten sig baade fra Familien og det frie Kultursamfund:
% --- p. 277 ---
\setcounter{footnote}{0}%
\opage{277}den er større end hin, men mindre end dette. En stor Familie
kan kaldes en lille Stat. Man har sagt, at en Familie er en Stat,
naar den ikke kan undertvinges uden Krig\footnote[1]{Thomas Hobbes: Leviathan. Cap. 20.}. Deri ligger, at
det er Magten, der er ejendommelig for Staten, og at den især
finder Anvendelse mod ydre Fjender. Og deri ligger tillige Forskellen
mellem Staten og Kultursamfundet. Staten forudsætter
et Folk, en Menneskegruppe, der føler sig som en Enhed i
Modsætning til andre Grupper. Men Kultursamfundet kan
brede sig ud over alle saadanne Grupper og forbinde dem ved
de harmonerende eller fælles Formaal, selvom ingen Magt omspænder
dem alle. ---

Disse tre Former for Samfundsliv ville vi nu nærmere betragte
med særligt Hensyn til deres etiske Betydning og til den
Aand og Retning, i hvilken deres videre Udvikling bør foregaa,
naar vi anvende den i det Foregaaende opstillede Maalestok.

% --- p. 278 ---
\setcounter{footnote}{0}%

\opage{278}\grouphead{A. FAMILIEN.}

\chapterhead{XIV.}{FAMILIENS ETISKE BETYDNING.}

1. Intet giver et bedre Exempel paa, hvorledes Naturen
kan bane Vej og danne Grundlag for, hvad Etiken fordrer eller
billiger, end den Omstændighed, at det inderligste og fuldkomneste
af alle menneskelige Samfund stiftes ved nogle af de stærkeste
Instinkter i den menneskelige Natur. Det Humanitetsrige, som
er Etikens højeste Ideal, har ikke blot sin første Spire og sin
stadige Kilde i Familieforholdet, men er i Familielivet, naar
dette har naat sin højeste Form, virkeliggjort paa en Maade,
som andre Samfundsformer ikke kunne vise Sidestykke til.
Højdepunkterne af alle andre Samfundsformer maales ved den
Grad, i hvilken de minde om Familieforholdets Inderlighed og
Styrke. Humanitetsriget vilde have naat sin Fuldendelse, naar
et almindeligt Broderskab forenede Alle; og til at betegne et
inderligt Forhold mellem Herre og Tjener, Mester og Lærling,
Øvrighed og Folk har man det mest træffende Udtryk i Forældres
Forhold til deres Børn.

Fra flere forskellige Sider vil Familieforholdets store etiske
Betydning kunne paavises.

2. I andre Samfund deltager Mennesket kun med en Del
af sit Væsen, men i Familien kan det finde Næring for alle Sider
% --- p. 279 ---
\setcounter{footnote}{0}%
\opage{279}af sin Natur. Kun der lever det egentligt som helt Menneske.
De mest primitive Instinkter og de mest ideelle Følelser finde
der deres Tilfredsstillelse. Livsfællesskabet strækker sig (eller kan
strække sig) fra de rene Naturinstinkter, ved hvis Virken Individet
ofte kan synes at være blot Middel for Slægtens Villie til
Livet, gennem alle materielle og aandelige Omraader. Alle Interesser
kunne finde deres første Pleje i Familien. Familien er
en lille Verden, som sætter alle Kræfter i Virksomhed. Hermed
hænger det igen sammen, at Familien mere end noget andet
Samfund forener (eller kan forene) Personlighedernes Selvstændighed
med deres inderlige Forbindelse. Netop fordi Familien
ikke er noget specielt Samfund, men et almindeligt Livsfællesskab,
kan den største Ejendommelighed og Selvstændighed hos
de enkelte Medlemmer frit røre sig og dog møde Forstaaelse og
Sympati. I andre Samfund er dybere gaaende Selvstændighed
og Ejendommelighed mere eller mindre nødt til at skjule sig
eller ogsaa til at holde Andre paa Afstand. I Familien kunne
paa Grund af det fuldstændige Livsfællesskab selv de tilsyneladende
sære og paradoxe Egenskaber møde Sympati, fordi de
opfattes og forstaas i Sammenhæng med Individets hele Natur.
Fremmedheden er ophævet. Derfor er der ikke Hvile og Styrke
i noget Forhold som i dette.

I Familieforholdet leves der ikke blot med fuld Bevidsthed,
men en Mængde ubevidste eller halvt bevidste Indflydelser gøre
sig uafbrudt gældende. I de andre Samfundsforhold spiller bevidst
Iagttagelse og Reflexion og bevidst Beslutning og Handling
en langt større Rolle. Der stilles Krav til anspændt Opmærksomhed
og til selvbevidst Optræden. Men i Familielivet
raader det Uvilkaarlige og Ubevidste i langt højere Grad. Gennem
utallige Paavirkninger, Erindringer og Stemninger voxer
Følelsen ved Hjemmet til en saadan Højde, at den, naar den
angribes, kan fremtræde som en Affekt. Om de Følelser, der
ere knyttede til Hjem og Familie, gælder det i højere Grad end
om andre, at de næres og øges ved smaa Tilvæxter, som først
ved at summere sig sammen komme til at gøre sig gældende
i Bevidstheden. Følelseslivet faar derved en Fasthed, som det
% --- p. 280 ---
\setcounter{footnote}{0}%
\opage{280}mangler, hvor det fremtræder som en Række af heftigt opblussende,
men ligesaa hurtigt forsvindende Affekter\footnote[1]{Smlgn. Psykologi III, 7; VI E, 4—5}.

Endeligt forbinder Familien forskellige Generationer ved Naturens
Baand. Den slaar Broen mellem Slægtens Fortid og dens
Fremtid. Den gør en Forstaaelse mulig, hvor den mellem fjernere
Staaende vilde være umulig. Striden mellem det Gamle
og det Nye dæmpes indenfor Familien ved den dybere liggende
Sympati, som her fører det første Ord.

3. Man kunde indvende imod Familien, at den er for snevert
et Samfund, at den koncentrerer Følelse og Interesse om en lille
Kreds, saa at hvad der staar udenfor denne, bliver ligegyldigt.
Der kan uddanne sig en Familieegoisme, ganske vist af videre
Omfang end den rent individuelle Egoisme, men dog ikke mindre
end denne en Hindring for Udviklingen af den almindelige Menneskekærlighed.
Hertil er at svare, at Sympatien først maa udvikle
sig i de snevre Kredse, før den kan brede sig i vide Kredse.
I de Følelser, der udvikle sig og finde stadig Næring indenfor
Familien, have vi de første og stærkeste Midler til at dæmpe
og opdrage den individuelle Egoisme. Den almindelige Menneskekærlighed
er kun en Udvidelse af en Følelse, der opstaar indenfor
Familien, en Udvidelse, som ikke altid gaar for sig uden
Modstand, men som dog stedse forudsætter, at Spiren er lagt i
det Smaa. Der er derfor ingen nødvendig Modsigelse mellem
Familiekærligheden og den almindelige Menneskekærlighed. Og
dertil kommer (smlgn. XII, 1), at der er et omvendt Forhold
mellem Sympatiens Styrke og Omfang. Udvides Omfanget, gaar
det let ud over Styrken. De færreste Mennesker ere (i hvert
Tilfælde endnu) i Stand til at omfatte fjerne og vide Forhold
med en saadan Inderlighed og Kraft, som hvad der ligger dem
nær. Dersom det da ikke skal gaa ud over Følelsernes Styrke,
maa der være snevre Kredse, i hvilke de voxe sig stærke. ---
Familien er her paa én Gang Formaal og Middel. Den yder den
højeste Tilfredsstillelse for de Enkeltes Trang, og den er Hjemstedet
for Kræfter, der blive af Betydning for hele Slægten. Jo
% --- p. 281 ---
\setcounter{footnote}{0}%
\opage{281}flere Hjem der findes, des flere Arnesteder for den lid, der
holder Slægtens Liv vedlige. ---

En anden Indvending mod Familien er den, at den ved sin
Begrænsning og ved sit paa Overlevering og Gentagelse byggede
Liv let fører til Stillestaaen, Ensformighed og Sløvhed. Derved
hindres ofte det, der netop skulde være et af Familielivets Fortrin,
at den Enkelte her kunde udfolde sig efter alle Sider i sit
Væsen. Den mangesidige Udfoldelse hæmmes ved det snevre
Indhold og ved en ejendommelig Sky for at optræde uforbeholdent
i Ord og Handling overfor de Nærmeste. Derved forstaas,
at mange Individer i deres Hjem vurderes ringere end udenfor
det --- skønt det Omvendte vilde være naturligst. --- Disse Betænkeligheder
hidrøre dog kun fra Hjemmets eller Familiens
Isolation fra Kultursamfund og Stat. Fra disse større Verdener
komme de friske Strømninger, der kunne forjage den forgemte
Luft, der ved Hjemmets Isolation kan opsamles indenfor dets
Vægge\footnote[1]{Smlgn. min Indledningsafhandling til Værket „Vort Hjem“.}.

Ganske overvindes Vanskelighederne kun, hvor Familien har
naat sin højeste Form. Virkeligheden staar her, som overalt,
langt tilbage for Idealet. Men den frembyder Tilnærmelser,
som det gælder at udvikle videre. --- Til nærmere Belysning
af de Spørgsmaal, som her frembyde sig, ville vi først betragte
Ægteskabet, dernæst Kvindens Stilling og Vilkaar, og endeligt
Forholdet mellem Forældre og Børn.

\begin{quote}

\grouphead{1. ÆGTESKABET.}

\chapterhead{XV.}{SOCIOLOGISKE DATA.}

\end{quote}

1. Paa de laveste Trin, vi kende, har Forbindelsen mellem
de to Køn Karakteren af et Magt- og Tvangsforhold, hvor den
% --- p. 282 ---
\setcounter{footnote}{0}%
\opage{282}svagere Part blot er Middel for Tilfredsstillelsen af den andens
Lyster eller maa trælle for hans Behov. Nogle nyere Forfattere
(\textsc{Bachofen}, \textsc{Mc} \textsc{Lennan}, \textsc{Lubbock}) have ment, at der oprindeligt
herskede fuldstændig Parringsfrihed mellem de to Køn (Promiskuitet,
almindelig Hetærisme), saa at enhver Kvinde tilhørte
enhver Mand, som et Øjeblik kunde bemægtige sig hende. Dette
er dog sikkert en Overdrivelse af den store Frihed, som rundt
om paa Jorden hos vilde Folk finder Sted i Kønsforholdet. Der
har neppe været nogen Tid, hvor fuldstændig Hetærisme har
hersket. Selv hvor ingen Skik eller Lov hindrede den stadige
Skiften af kønslige Forbindelser, vilde dog individuel Forkærlighed
saavel som Lyst til Ejendom og Magt føre til, at Forbindelserne
bleve mere end rent øjeblikkelige\footnote[1]{At Beviserne for en oprindelig Hetærismes eller Promiskuitets Bestaaen ere uholdbare, er fra forskellige Sider blevet paavist, navnligt af C. N. Starcke (Die primitive Familie. Leipzig 1888) og Edward Westermarck (The History of Human Marriage. I. Helsingfors 1889. --- Il. London 1891). (Tysk Oversættelse 1902). --- Smlgn. min Anmeldelse af Starckes Bog i „Tilskueren“ 1888.}. Kvinden betragtedes
desuden tidligt ikke blot som Middel til Kønsinstinktets Tilfredsstillelse;
hun blev tillige den første Slave og fik derved en Værdi,
som maatte bevirke, at fuldstændig Frihed ikke kunde tilstedes.
Man kan ikke bevise, at alle andre Former for Ægteskab have
udviklet sig af en oprindelig Hetærisme, selvom man finder store
Tilnærmelser til en saadan rundt om paa Jorden\footnotemark[1]. Da Monogamiet
ogsaa findes hos Dyr, er der i og for sig Intet til Hinder
for, at det ogsaa paa rent primitive Trin kan findes hos
Mennesker, og dette er ogsaa Tilfældet. Saa tiltalende det end
kunde synes, saa kan man dog ikke paavise en sammenhængende
Udviklingsgang, der fra Hetærismen som laveste Trin fører
gennem Polygami --- dels Polygyni, dels Polyandri, dels Gruppeægteskaber
(hvor flere Mænd have flere Kvinder fælles) --- til
Monogami. Udviklingen har sikkert hos forskellige Stammer haft
en meget forskellig Karakter. Ægteskabets Ordning hænger overalt
sammen med saa mange andre sociale Forhold, at der neppe
kan paavises nogen enkelt, tydelig og simpel Udviklingsgang.

') Smlgn. Spencer: Principles of Sociology. I. p. 662 ff.
% CHECK p. 282: note mark without a note

% --- p. 283 ---
\setcounter{footnote}{0}%
\opage{283}Men én Ting kan dog paavises at have øvet en væsentlig og
nødvendig Indflydelse, nemlig den Grad, i hvilken den individuelle
Personligheds Ret er bleven følt og anerkendt. For denne
Ret har baade den sanselige Drift og Herskerdriften maattet
vige. Det er den, der har gjort, at Monogamiet under Udviklingens
Gang mere og mere er blevet anerkendt som den højeste
og den sande Form for Forbindelse mellem Mand og Kvinde.

Selvom der skulde kunne paavises Tider eller Steder, hvor
Flokken eller Stammen havde Kvinder og Børn fælles, ligesom
Tilfældet var med anden Ejendom, saa maatte dette Fællesskab
dog falde bort, efterhaanden som den Enkelte begyndte at føle
sin Selvstændighed. Trangen til at have Noget for sig selv
maatte ogsaa ytre sig paa dette Omraade. Det er fra først af
en rent selvisk Trang, der kun ytrer sig eller i hvert Tilfælde
kun vinder Anerkendelse for de Stærkes, og det vil sige for
Mændenes Vedkommende. Polygyni forekommer derfor hyppigere
og holder sig længere end Polyandri og Gruppeægteskab, der
især forekomme, hvor Kaarene ere trange. Monogamiet er ofte
ogsaa kun et Udslag af Trangen til Særeje, altsaa et Udslag af
Magt, kun at denne af en eller anden Grund, som ikke altid
er let at paavise, indskrænker sit Omraade. Forholdet mellem
Mand og Kvinde i Ægteskabet har her endnu overvejende Karakteren
af et Tvangsforhold. Det stiftes maaske endog ved Vold,
Rov eller Køb. Købet betegner et Fremskridt i Forhold til Vold
og Rov; ti derved træder det tydeligt frem, at Kvinden har en
økonomisk Værdi, og hun skaanes som enhver kostbar Vare.
Men selv hvor Køb ikke længere finder Sted, gives Ægtefællerne
(især Kvinden) ofte bort uden frit Valg; Ægteskabet betragtes
som et Anliggende, ikke imellem de to Individer, men
mellem de to Slægter, som forbindes ved dem, som forplantes
gennem deres Børn, og hvis Ejendom gaar i Arv til disse.
Ægteskabet er her ikke en personlig Pagt, men et „Ætteforbund“\footnote[1]{ter. Il. Kristiania 1867. p. 306. --- Michel de Montaigne (Essais III, 5) udtaler sig (i Slutn. af d. 16. Aarh.) saaledes om Ægteskabet: „On ne se marie pas pour soy, quoy qu'on die; on se marie autant, ou plus,}.
I det gamle Ægteskab var ikke det personlige For-

!) Smlgn. for Nordens Vedkommende: R. Keyser: Efterladte Skrif%
% --- p. 284 ---
\setcounter{footnote}{0}%
% CHECK p. 284: notes paired with the marks in order (the layer misread a number)
\opage{284}hold mellem Mand og Kvinde Hovedsagen, men Grundlæggelsen
af en ny Husstand\footnote[1]{Leist: Alt-arisches Jus civile. I. Jena 1892. p. 154. 166.}. At Ægteskabsbrud straffes saa strengt
i tidligere Tider, hænger sammen med, at Kvinden betragtes
som Slægtens eller Mandens Eje\footnote[2]{A. H. Post: Die Grundlage des Rechts und die Grundziige seiner Entwickelungsgeschichte. Oldenburg 1884. p. 374 ff. --- Fu stel de Coulanges: La cité antique. 4. éd. p. 109.}. --- Først hvor de to Individer
vælge hinanden efter fri Tilskyndelse og med fri Beslutning,
have vi den højeste Form for Ægteskab: det frie Monogami.

2. Hos Mange opstaar der en vis Svimmelhed, naar de
første Gang blive opmærksomme paa de store Forskelligheder i
Ordningen af Ægteskabsforholdene rundt om paa Jorden og til
de forskellige Tider. Tilfældighederne synes at herske, og det
Ene synes at kunne være ligesaa godt som det Andet; i hvert
Tilfælde synes det, der én Gang er blevet praktiseret, allerede
derved at faa en vis Autoritet for sig. Dersom f. Ex. Hetærisme
har været almindelig, eller dog Tilnærmelser til Hetærisme,
skulde det saa ikke kunne tænkes, at vi vendte tilbage til den
igen? Hertil er at svare, at Prøvekortet, som Sociologien viser
os over de forskellige Ordninger af Ægteskabet, vel er broget;
men at der dog, som allerede antydet, viser sig en tydelig Sammenhæng
mellem Udviklingen af det frie Monogami og Anerkendelsen
af den individuelle Personligheds Betydning. Denne
Sammenhæng skal i det følgende Kapitel nærmere paavises. Den
bliver oprindeligt til ved Motiv- og Værdiforskydninger. I det
Væsen, Manden knytter til sig for sin Lysts Skyld, opdager han
en økonomisk Værdi, og senere opdager han, at hans kvindelige
Slave har en Sjæl som han selv, og at hun maa staa for ham
ikke blot som Middel, men som Formaal. --- Sociologien kan
give Etiken værdifulde Fingerpeg, men den er ikke absolut bepour
sa posterité, pour sa famille; I'usage et I'interest du mariage touche
nostre race, bien loing pardelå nous: pourtant me playt cette facon
qu' on le conduise plustost par main tierce que par les propres“. ---
Ogsaa i andre end adelige Kredse findes denne Opfattelse lige op til
vore Dage.
% --- p. 285 ---
\setcounter{footnote}{0}%
\opage{285}stemmende for Etiken. Selvom man kunde paavise en Udviklingsgang,
der begyndte med Hetærismen og endte med det frie
Monogami, vilde der endnu staa den Opgave tilbage for Etiken
at paavise dettes Værdi i Kraft af Velfærdsprincipet. Den historiske
Udvikling kunde jo meget godt have ført i en fordærvelig
Retning, saa at Etiken maatte fordre en Forandring, saa
vanskelig end en saadan maatte være at gennemføre. Hvad der
er virkeligt, er derfor ikke fornuftigt.

Det virkelige Forhold er nu dog det, at det frie Monogami
vel i de mest civiliserede Lande er officielt anerkendt som den
eneste rette Form for Ægteskab, men at de andre Former dermed
ikke ere forsvundne fra Livet. Hetærismen breder sig
rundt om os, selvom den nødes til at holde sig i Mørket. Og
hos en stor Del af Befolkningen holder den sig ikke engang i
Mørket, men udfolder sig frit og uvilkaarligt. De laveste Former
for kønslig Forbindelse fremtræde stedse paany, ligesom de
laveste Dyreformer kunne vedblive at bestaa, selv efter at højere
Typer forlængst have udviklet sig. Hvad der betegnes som
„sædelig Slappelse“\footnote[1]{Hafstrøm: Om Sædelighedsforholdene i det danske Folk. København 1888. p. 67.}, er meget ofte gamle Livsformer, der ikke
ere forsvundne. Dog mødes her paa en mærkelig Maade Naturen
og Overcivilisationen. Den stadigt florerende Hetærisme
skyldes ikke blot den friske, ubændige Naturs Drifter, som endnu
ikke ere mildnede og harmoniserede. Den skyldes ogsaa for en
Del en ensidig Uddannelse og Forfinelse af nogle Aandsevner
paa den hele Karakters Bekostning. Intellektuel og æstetisk Kultur
kan være forenet med stor Løshed i kønslig Henseende.
Man har sine, maaske alvorlige Interesser paa andre Omraader
og forholder sig væsentligt nydende og legende paa det kønslige
Omraade, uden at dettes etiske Betydning gøres til Genstand
for Besindelse. Med ensidig intellektuel og æstetisk Kultur
følger desuden ikke sjældent en blaseret og raffineret Tænkemaade
og en Sans for det Pikante, der særligt finder Næring
paa det kønslige Omraade. Ogsaa to andre modsatte Motiver
kunne mødes: den bitre Nød og den overstrømmende Livskraft.
% --- p. 286 ---
\setcounter{footnote}{0}%
\opage{286}Nøden fører til at give sig hen, og Nydelsestrangen fører til at
benytte dem, der nødes til at sælge sig, som Midler til et Øjebliks
Lyst. --- Af saa forskellige Aarsager kan stadigt den Isolation
af Kønsdriften fra de andre Elementer i den menneskelige
Personlighed fremgaa, som er Brodden i det sexuelle Problem.
At denne Drift kan fremtræde isoleret og dog med saa stor
Vælde, hænger sammen med det Overskud af Energi, Naturen
overalt sætter ind paa at hævde Slægternes Bestaaen. Isolationen
bliver betænkeligere, jo mere Kulturlivet afløser Naturlivet.
Det sexuelle Problem er derfor vævet nøje sammen med hele
Kulturproblemet og med hele det sociale Problem og kan kun
forstaas ret, naar denne Sammenhæng ikke glemmes.

\chapterhead{XVI.}{DET FRIE MONOGAMI.}

1. Et Samfund mellem menneskelige Personligheder kan
kun da være fuldkomment, naar ingen af Deltagerne staar som
blot Middel for de Andre, og naar ingen Del af den enkelte Deltagers
Væsen bliver ensidigt fremdragen eller skudt tilside. Alle
polygame Ægteskabsformer stride mod dette Princip. I Polygamiet
deler en Mand sig mellem flere Kvinder eller en Kvinde
mellem flere Mænd. Den fulde Hengivelse kan da ikke være
tilstede overfor nogen Enkelt: Forholdet er stedse delt, og den
fulde Forbindelse af Personlighederne indtræder ikke. Paa den
ene Side stykker man sig ud, paa den anden Side nøjes man
med et Brudstykke. „I Polygamiet“, siger \textsc{Kant} (Rechtslehre
§ 26), „vinder den Person, der giver sig bort, kun en Del af
den Person, til hvilken den ganske hengiver sig, og den gør
sig derved til en blot Ting“. De Mange gøres til Midler for den
Ene, eller den Ene for de Mange. Og dette medfører igen, at
der indenfor den enkelte Personlighed indtræder en Sondring
% --- p. 287 ---
\setcounter{footnote}{0}%
\opage{287}af det, der skulde virke udelt sammen. Hvor Elskovsfølelsen er
mere end dyrisk Brynde, virke et fysisk og et ideelt Element
uadskilleligt sammen, og Hengivelsen omfatter den hele Personlighed.
(Smlgn. XI, 10). Men i Polygamiet maa nødvendigvis
det rent fysiske Element blive særligt fremdraget. Det Brudstykke
af sig selv, man kan give til Mange, er her kun det rent
fysiske Element. Som rent fysisk Forhold, som rent dyrisk Akt
har Kønsforholdet ogsaa kun fysiske Grænser. Den fysiske Side
af Kønsforholdet er den mindst individualiserede; jo mere den
drages frem, des mere forsvinder Forskellen mellem Objekterne;
og omvendt, jo flere forskellige Objekter Kønsinstinktet retter
sig imod, des mere maa dets fysiske Side blive fremherskende.
Fluernes Kønsliv er et let tilgængeligt Exempel. Den rent fysiske
Forbindelse er momentan, og den Standsning af den indbyrdes
Kamp for Tilværelsen, den medfører, varer ofte kun i selve Parringsøjeblikket.
Hos nogle Edderkoppearter vaagner Rovdyrnaturen
igen strax efter Parringen, saa at Han og Hun nu kun
se et muligt Bytte i hinanden. Det viser, hvor isoleret og momentant
Kønsinstinktet kan optræde. I den egentlige Elskovsfølelse
ytrer sig derimod foruden det elementære Instinkt Billedet
af det andet Individ i dets Ejendommelighed og den inderlige
Glæde over denne. Instinktet virker da ikke som en rent
blind Magt, men aabner Øjet for det andet Individs Natur og
baner derved --- gennem en Forskydning af Motiver og Værdier
--- Vejen for en Hengivelse, der er mere end blot fysisk
Forbindelse. Den etiske Betydning af Kønsinstinktet vil paa
højere Trin være den, at der opdages og vurderes personlige
Ejendommeligheder, der ellers vilde forblive ubemærkede. Der
vækkes en Opmærksomhed, som virker med en Naturkrafts Magt.
Forbindelsens Inderlighed beror da ikke blot paa, at Slægtens
Selvopholdelsestrang knytter et Individ til et Individ af det andet
Køn, men ogsaa derpaa, at Egenskaber og Livsytringer, der kun
kunne fremtræde ved et varigt Samliv, øve deres Indflydelse.
Ikke blot Kærligheden, ogsaa Troskaben er et Organ til at gøre
aandelige Opdagelser.

Kun Monogamiet kan derfor være Formen for kønslig Forbindelse
mellem selvstændige Personligheder, der indtræde i For%
% --- p. 288 ---
\setcounter{footnote}{0}%
\opage{288}holdet hver efter sin hele Natur, ikke for at tilfredsstille en
enkelt, isoleret Fornødenhed.

2. Det hører med til den personlige Hengivelse, at den
ikke blot omfatter Personligheden efter dennes hele Omfang,
men at den ogsaa omfatter hele Livet, saalænge dette varer for
begge Parter. En Følelse, der ikke uvilkaarligt tror paa sin egen
Vedvaren, er ikke ægte. Det vilde være unaturligt, om der,
naar Følelsen overfor et andet Individ er i den stærkeste Bevægelse,
dog gjordes et Forbehold og var en udtrykkelig Bevidsthed
om, at det Hele kun skulde gælde for en vis Tid. Er
Følelsen ægte, vil en saadan Mulighed ikke kunne træde frem
for Tanken. Fordi Følelsen tror og maa tro paa sin egen Vedvaren,
behøver den ganske vist ikke at vare ved. Det Superlativiske
ligger i Følelsens Væsen; den udelukker, naar den er
stærk, Begrænsning og Sammenligning. Men selvom det, der
mentes at skulle have evig Bestaaen, viste sig at være forgængeligt,
saa vilde Følelsen dog ikke kunne udfylde den Tid, den
varede, paa den rette Maade, naar den ikke var forbunden med
Troen paa sin Vedvaren. Illusioner ere hyppige; men et Forhold,
som stiftes med udtrykkelig Bevidsthed om eller endog Beregning
af, at man altid vil kunne drage sig ud af det, et saadant
Forhold er ikke blot i Regelen et Bedrageri overfor den
anden Part, men kan vanskeligt blive noget helt og fuldt Forhold.

Et personligt Væsen kan ikke udstykkes i løsrevne og skiftende
Øjeblikke. Personlighed er kun tilstede, hvor der gennem
de enkelte Øjeblikke drager sig en indre Enhed og Sammenhæng,
og hvor der er visse bestemte Følelser, som udgøre den
faste Kærne. Men et personligt Væsen kan heller ikke udtømmes
i et enkelt Øjeblik. Det indeslutter i sig en saadan Rigdom,
at der, naar virkelig Sympati er tilstede, kan være nok at opdage
et helt Liv igennem. Det er en stor Illusion, naar man
mener, at en Række skiftende kønslige Forbindelser skulde afgive
et rigt Stof for Menneskekundskab og menneskelig Udvikling\footnote[1]{I saa Fald maatte en af Hieronymus omtalt romersk Kvinde have drevet det til en høj Grad af Menneskekundskab: hun var gift med sin treogtyvende Mand og var hans enogtyvende Kone.}.
% --- p. 289 ---
\setcounter{footnote}{0}%
\opage{289}Øjebliksforbindelser føre ikke ind i den inderste Helligdom;
denne oplader sig kun for den stadige, trofaste Sympati.

Hvad man med et ofte misbrugt Ord kalder Læren om
„den frie Kærlighed“, er et Forsøg paa at sætte Ubestandigheden
i System og proklamere den som et væsentligt Moment
i Elskovsfølelsen. \textsc{Enfantin}, St. Simons bekendte Discipel, mente,
at det nu bestaaende Ægteskab kun passede for de stadige og
skinsyge Naturer, ikke for de ustadige og lunefulde, --- for
Othello'erne, ikke for Don Juan'erne. Men da alle Slags Naturer
have Krav paa Tilfredsstillelse --- ogsaa „les personnes vives,
coquettes, séduisantes, attrayantes, changeantes“ --- saa fordrer
han Skilsmissen som Institution, som Anerkendelse af Retten
til Ustadighed, ikke blot som en Udvej i Konflikttilfælde\footnote[1]{London 1875. p. 371.}. En
lignende Tankegang findes hos den anonyme Forfatter til „Elements
of Social Science“. Naar denne sidstnævnte Forfatter dog
ikke vil afskaffe Begrebet Selvbeherskelse, og naar han hævder,
at Kønsinstinktet maa tilfredsstilles saaledes, at det ikke bliver
til Skade for Andre, saa fordrer han en Underordning af de
øjeblikkelige Tilskyndelser under mere omfattende Hensyn, som
strider mod den Tankegang, der fører til Læren om den frie
Kærlighed. Ti Lidenskaben er absolutistisk, og dens Eneherredømme
ophæves, naar andre Følelser, f. Ex. Medfølelsen, gøre
sig gældende ved Siden af den. Der er neppe Spørgsmaal om,
at naar en levende Medfølelse med andre Væsener er vakt, før
den sexuelle Trang kan melde sig i sin fulde Kraft, vil denne
Trang undergaa en betydelig Metamorfose. Det er paa Isolationen,
en væsentlig Del af dens Styrke beror. --- Iøvrigt maa
man, for ikke at gøre denne sidstnævnte Forfatter Uret, mærke
sig, at han ved Ægteskab mener det uopløselige Ægteskab. Hvor
der er nogenlunde let Adgang til Skilsmisse, som i Tyskland,
dér --- mener han --- er Ægteskabet faktisk afskaffet\footnotemark[1]. Man
ser, hvor let det Hele bliver en Strid om Ord. Mange af de
Angreb, der, især af Æstetikere, rettes imod Ægteskabet i al
p Smlgn. Paul Janet: St. Simon et le St. Simonisme. p. 138---142.

') Elements of Social Science. By a Doctor of Medicine. 13th ed.
% CHECK p. 289: note mark without a note

% --- p. 290 ---
\setcounter{footnote}{0}%
\opage{290}Almindelighed, skyldes sikkert kun Reminiscenser fra franske
og engelske Forfattere, som polemisere mod det uopløselige
eller kun af ganske enkelte Grunde opløselige Ægteskab.

Den Ubestandighed, som undertiden regnes for et væsentligt
Moment i Forholdet mellem Mand og Kvinde, kan ikke
undgaa at føre til Smerte og Ulykke, saalænge den ikke ligger
i Alles Natur, og --- maa der føjes til --- saalænge Trangen til
Afvexling ikke hos begge Parter indtræder (ved en mærkelig
harmonia præstabilita) paa samme Tid. En dansk Forfatter\footnote[1]{Den anonyme Forfatter til En Livsanskuelse, grundet paa Elskov (Kbhvn. 1881) og Forholdet mellem Mand og Kvinde belyst gennem Udviklingshypotesen (Kbhvn. 1884). Smlgn. ogsaa hans Af handling i „Tilskueren“ (1885): Om en Reaktion mod den moderne Stræben efter større sexuel Sædelighed. --- (Tv},
som med stor Dygtighed og Alvor har drøftet Forholdet mellem
Mand og Kvinde fra et rent psykologisk og etisk Standpunkt,
har paavist, at saalænge der foruden de polygame Naturer, for
hvem Ubestandighed og Skiften hører til Elskovsfølelsens Væsen,
ogsaa gives monogame Naturer, som ikke kunne tænke sig sand
Elskov uden Bestandighed og Troskab, saalænge vil den Trang
til Afvexling, som den frie Kærligheds Forkæmpere lægge saa
megen Vægt paa, ikke kunne tilfredsstilles uden at medføre
Smerte og Sorg for Andre. Det vil ofte gaa, som med et Par
i \textsc{Daudets} „Sappho“.\textbackslash{} Forholdet stiftedes med den Bagtanke, at
det skulde være paa Tid; men for den unge Kvinde er Forholdet
blevet til Alvor, og hun tager sin Død, da Adskillelsen
skal indtræde: On meurt donc quelquefois de ces ruptures! ---
Er det letsindigt at lege med Ild, saa er det tusinde Gange
letsindigere at lege med et Menneskes Velfærd.

3. Den individuelle Personlighed er i Reglen ikke fuldt
færdig i den Periode af Livet, i hvilken Elskovsfølelsen spiller
den største Rolle. Naar denne Følelse nu har ført til ægteskabelig
Forbindelse, saa kommer det an paa, om de to Personligheders
fortsatte Udvikling kan foregaa i indbyrdes Harmoni.
Dette er af saa meget des større Betydning, som det ikke blot
er de to Individers Karakter idethele, men ganske særligt den
Følelse, der forbinder dem, der i Tidens Løb kan og vil under%
% --- p. 291 ---
\setcounter{footnote}{0}%
\opage{291}gaa Forandringer. Hvor Udviklingen foregaar lykkeligt, vil Elskovsfølelsen
fra en opblussende Affekt blive til en inderlig
Følelse, der --- hvor det gælder --- kan have den samme
Styrke, om ikke den samme øjeblikkeligt betagende Voldsomhed
som i sin første Frembryden.

Den Udvikling, som Individerne og særligt den Følelse,
der forbinder dem, kunne undergaa under det fortsatte Samliv,
er ikke uafhængig af deres Villie. Det er en af de gængse
psykologiske Misforstaaelser, at Villien er Noget, som er helt
forskelligt fra Tanke og Følelse, en Magt, der først særligt skal
appelleres til, for at den som en deus ex machina skal træde
frem og løse Vanskelighederne. Under sunde og naturlige Forhold
udvikler og fæstner Villien, hvad Tanken har grebet og
Følelsen sluttet sig til. Villien maa være med fra først til sidst;
kun saa kan den ogsaa hjælpe i det enkelte, skæbnesvangre Øjeblik.
Det er ikke en blind Skæbne, de to Individer ere undergivne.
Sagen er for en stor Del lagt i deres egen Haand; det
kommer an paa den Alvor, hvormed de føre deres Liv. Ægteskabet
fordrer som ethvert Samliv Selvbeherskelse og Anstrengelse
for at kunne bestaa. Hvis enhver Misstemning, enhver
Karaktermodsætning faar Lov at blive bestemmende for Forholdets
Fremtid, vil denne kun blive kort. Hvis der ikke kendes nogen
anden Tilfredsstillelse end den, som den opblussende Affekt og
Indflydelsen af det Nye og Overraskende giver, vil den egentlige
ægteskabelige Kærlighed ikke kunne udvikle sig. Der gives en
Form for æstetisk Livsanskuelse, for hvem det Højeste er Følelsernes
skiftende Spil, og som derfor stedse søger Forandring og
Nyhed. (Smlgn. III, 4). En saadan Livsanskuelse kan ikke gaa
ind paa Ægteskabet; men den kan konsekvent heller ikke gaa
ind paa nogetsomhelst fast Forhold. Fastheden kommer kun,
naar Villien sættes ind, og naar der kendes anden Form for
Følelsesliv end Øjeblikkets Affekter. Hvad det gælder om, er,
at den erotiske Følelse optages som Led af det hele personlige
Liv og ikke vedbliver at staa som et tilfældigt og løsrevet Element.
Og for en saadan inderlig Forbindelse med Livets andre
Elementer har den maaske større Mulighed end andre Følelser.
Der var noget Træffende i, naar en af de senere Stoikere (hos
% --- p. 292 ---
\setcounter{footnote}{0}%
\opage{292}hvilke man, som det er blevet bemærket\footnote[1]{Bonhöffer: Die Ethik des Stoikers Epiktet. Stuttgart 1894. p. 88. -) Hafstrøm: Om Sædelighedsforholdene i det 'I danske Folk. p. 82.}, finder den mest ideale
Opfattelse af Ægteskabet, Oldtiden kender) beskriver den ægteskabelige
Forbindelse som en kemisk Forbindelse, en „fuldstændig
Blanding“, i Modsætning til den mere mekaniske Forbindelse
ved andre Slægtskabs- og Venskabsforhold. Hvo der vil bevare
det sexuelle Element i dets Isolation for derigennem at sikre sig
en egen Kilde til Nydelse, vil frygte den Metamorfose, hint Element
undergaar i det rette Ægteskab. Og den, der vel ønsker
en saadan Metamorfose, men frygter for, at den ikke vil lykkes,
han vil ogsaa nære Betænkelighed ved Ægteskabet. Men det vil
kunne medføre en helt anden Opfattelse af Forholdet mellem de
to Køn, om man kun kender det uden saadan Metamorfose eller
efter en saadan. En erfaren Iagttager af Straffefanger og deres
Forhold-) bemærker: „Skønt Straffefanger jo i Reglen ere udgaaede
fra de daarligste Samfundslag, vise deres Brevvexlinger
jevnligt, hvor trofaste og udholdende de fleste Hustruer ere,
medens deres Mænd afsone Straf for Forbrydelser\dots{} Ogsaa
de fleste Mænd, hvor dybt sunkne de end ellers kunne være,
hænge dog med Kærlighed ved deres Hustruer og Børn, og helt
underligt er det at se, hvor de gifte Fanger i Regelen have et
helt andet Syn paa Kvinden end de ugifte. I Ægteskabet lære
de hurtigt, at Kvinden har en ganske anden Betydning og Værdi
end blot at være Redskab for Mandens Sanselighed“.

Det vilde være upsykologisk at paastaa, at Villien kan gøre
Alt. I det Livsforhold, som her er Tale om, virke saare mange
Elementer sammen, og det er ikke dem alle, som selv den alvorligste
Villie kan beherske. Uden egentlig Skyld paa nogen
af Siderne kan Udviklingen tage en uheldig Retning. Der kan
opstaa Disharmoni i Karaktererne og i Livsanskuelserne; Forholdet
mellem det fysiske og det ideelle Element i selve Elskovsfølelsen
kan udvikle sig paa forskellig, ja modstridende
Maade hos de to Individer; overfor de nye Forhold og Spørgsmaal,
som Livet fører med sig, kunne de føle sig i helt modsat
Stemning. Det maa endog, naar man ret betænker, hvor mange
% --- p. 293 ---
\setcounter{footnote}{0}%
\opage{293}ydre og indre Betingelser der her virke sammen, betragtes som
et stort Held, naar den oprindelige Følelse staar sin Prøve og
undergaar sin Metamorfose under Individernes fortsatte Udvikling
uden at stanse paa noget af de mange Skær; Stoikeren
Antipater kaldte forsaavidt med Rette Ægteskabs Indgaaelse for
en Heltegerning (ἡρωϊκόν). Derfor indeholder Ægteskabet Muligheden
af saa mange Dramaer og fører til saa mange Tragedier.
Det Uvilkaarlige og det Vilkaarlige, Skæbne og Skyld kunne
her, mindre end paa noget andet Punkt, holdes skarpt ude fra
hinanden. Det vilde være doktrinært at paastaa, at naar Følelsen
blot var alvorlig fra først af, saa vilde den ogsaa vare ved.
Den kan have været alvorlig fra først af, men den senere Udvikling
kan have unddraget den en Næring, uden hvilken den
ikke kan bestaa.

4. Af væsentlig Støtte ved den Metamorfose, som Elskovsfølelsen
undergaar under det fortsatte Samliv, ere de fælles Opgaver,
der stilles de to Individer, og den fælles Virksomhed, de
udøve. Der opstaa nye Værdier, som forud ikke anedes. --- De
ere fælles om Arbejdet for deres materielle Udkomme. De have
at kæmpe med de samme ydre Tilskikkelser. De kunne i Fællesskab
arbejde paa deres videre aandelige Udvikling. Naar de
følges ad i deres Søgen, ville de lettere følges ad i deres Finden.
--- Men størst Betydning har fælles Omsorg for Børnene.
Fælles Bekymring og fælles Opofrelse lade dem komme hinanden
endnu nærmere end i den nye Følelses lette og glade Tid.
Sympatien imellem dem inderliggøres og forhøjes. Ansvaret for
at bevare Forholdet føles saa meget des stærkere, naar ikke
blot et enkelt Individs Skæbne, men flere hjælpeløse Væseners
Fremtid beror paa, hvilken Alvor og Selvbeherskelse der vises\footnote[1]{Charles Darwin fortæller om sin Fader, der var Læge: „Paa Grund af hans Evne til at vinde Tillid modtog han mange forunderlige Bekendelser om Ulykke og Skyld. Han talte ofte om, hvor mange ulykkelige Hustruer han havde kendt. I flere Tilfælde vare Mand og Hustru komne meget godt ud af det sammen i nogle og tyve Aar, men kom saa til at hade hinanden bittert; dette forklarede han ved, at de havde mistet et fælles Baand derved, at deres Børn vare blevne voxne.“ Life and Letters. I. p. 13). --- Andre fælles Baand end dette ere altsaa nødvendige.}.
% --- p. 294 ---
\setcounter{footnote}{0}%
\opage{294}Her er tillige et Punkt, hvor Monogamiets Fortrin for Polygamiet
tydeligt viser sig. Kun ved det monogame Ægteskab kan
Barnet faa den hele og fulde Forældrekærlighed. Kun dér kan
den fulde Harmoni og Enhed raade i Hjemmet. Og selv hvor
Elskovsfølelse ikke har været det første Motiv til Ægteskabets
Indgaaelse, men hvor „Fornufthensyn“, især økonomiske Hensyn
eller Ønsket om at efterlade sig Afkom have været bestemmende,
vil der ved en Motivforskydning kunne opstaa den
inderligste Forening. De fælles Opgaver og den gensidige Forstaaelse
af Personlighederne, som Samlivet kan medføre, kunne
muliggøre en Sympati, der ofte er nok saa inderlig som den,
der opstaar ved Elskovsfølelsen og dens Metamorfoser. --- Hvor
Ønsket om Afkom ikke har været det Motiv, der førte til Ægteskabet,
staar Afkommet som en ny Værdi, det indgaaede Forhold
afføder. --- Ægteskabets Udviklingsgang i de enkelte Tilfælde
kan saaledes frembyde Exempler paa alle de tre Arter af
Metamorfose, som ovenfor omtaltes (XIII, 4).

5. Ved den i det Foregaaende givne Karakteristik af Ægteskabet
udelukkes ikke blot Polygamiet, men ogsaa det ufrie
Monogami, hvor Kvinden ikke fuldt ud staar som Mandens Lige
i Ægteskabet, dels fordi hun ikke vælger med Frihed, men
gives bort af sin Slægt, dels fordi hun indenfor Ægteskabet
ikke har samme Ret som Manden, og dels fordi hun ikke i
aandelig Henseende er ham jevnbyrdig. Det sidstnævnte Punkt
indeholder Grunden til de to første. Naar Kvinden væsentligt
fører et dunkelt Følelses- og Instinktliv, og naar hendes Kald
tænkes udtømt med Husholdning og Barnepleje, kan hun kun
være et Vedhæng til Manden, udfylde en ganske enkelt Side
af hans Liv. Til at være Mere mangler hun den Selvstændighed
i Tanke- og Villieslivet, som er Betingelsen for at kunne
vælge og gøre sig gældende efter Valget. Ægteskabet bliver da
paa de væsentligste Punkter et Forhold mellem en aktiv og
passiv Part og faar ikke den Fuldkommenhed, som hvor det
er et Forhold mellem to Parter, der hver paa sin Vis ere aktive
og kunne følge hinandens Virksomhed med forstaaende Sympati.
Arbejdsdelingen imellem dem vil da kunne rette sig efter deres
Individualitet, og Overlegenheden vil ikke være saa godt som
% --- p. 295 ---
\setcounter{footnote}{0}%
\opage{295}selvskreven paa den ene Side. Det maa ikke skyldes et legaliseret
Magtsprog, hvem der skal være den Første. Men dette vil
kun kunne opnaas ved en mere alsidig Udvikling af den kvindelige
Begavelse, end man indtil den nyeste Tid har anset for
mulig og forsvarlig. Da først vil Monogamiet naa sin Fuldendelse.
Den Gerning, som Naturen én Gang har anvist Kvinden,
vil hun naturligvis ikke skyde fra sig; en virkelig Naturbestemmelse
lader sig ikke fornægte. --- En nærmere Betragtning
af dette Spørgsmaal vil i et senere Afsnit blive anstillet.

Den forskellige Stilling, Mand og Kvinde efter den hidtil
gængse Betragtningsmaade indtage i Ægteskabet, lægger sig
særligt for Dagen i den forskellige Dom, som rammer deres
Utroskab. At Hustruens Utroskab dømmes strengest, kunde
synes at finde sin Begrundelse ved, at den medfører langt mere
indgribende Følger for Familien end Mandens, hvis Eventyr
slet ikke behøve at medføre Følger, der mærkes i Familien.
Men dette taber sin Betydning overfor den Fordring om fuldkommen
lige Ret, som det frie Monogami stiller. Hvor denne
Fordring ikke tilfredsstilles, er der i Virkeligheden et Polygami,
og den forskellige Bedømmelse af mandlig og kvindelig Utroskab
stammer ogsaa fra et Udviklingstrin, hvor Kvinden betragtedes
som Mandens Ejendom. (Se XV, 1). Den forskellige Bedømmelse
har holdt sig, efter at den officielle Anerkendelse af
det frie Monogami er sket.

En italiensk Forfatter, Mantegazza\footnote[1]{Die Physiologie der Liebe. Oversat fra Italiensk). Jena 1877. p. 351 ff.}, kæmper drabeligt for
den ulige Bedømmelse. Hans Tankegang bygges især paa to
Grunde. For det Første paa, at Kvindens hele etiske Opgave
ligger i Familieforholdet. Medens der — siger han --- af Manden,
der maa arbejde og kæmpe i saa mangfoldige Forhold,
kræves hundrede Dyder, er Troskab den eneste [!] Dyd, der
forlanges af Kvinden: „er det maaske for Meget?“ —Men dernæst:
Manden er af Naturen tilbøjelig til Polygami, han er
mere troløs, brutal, lunefuld og sanselig end Kvinden, der ikke
saa let overfaldes af en Sanserus. Denne Forskel udmaler
Mantegazza med saa stærke Farver, at det bliver ubegribeligt,
% --- p. 296 ---
\setcounter{footnote}{0}%
\opage{296}hvorledes et saa rent Væsen nogensinde kan falde. I hvert
Tilfælde vilde det dog være rimeligt at vende den strengeste
Dom ikke mod den faldne Engel, men mod den Djævel, der
volder hendes Fald, især hvis de have Ret, der mene, at Kvindens
Fald langt oftere end Mandens skyldes et højere Motiv
end blot Sanselighed. --- At Forskellen er tilstede til en vis
Grad, kan vel ikke nægtes. Men den er her skruet op til en
unaturlig Højde, samtidigt med, at Kvindens hele Opgave som
Menneske indskrænkes til Familieforholdet. Og selv indenfor dette
er der dog Lejlighed nok for Kvinden til at lægge flere Dyder
for Dagen end hin ene, ganske vist overmaade betydningsfulde.

Baade den Indvirkning, som Løshed i ægteskabelige Forhold
har paa Familielivet idethele, og den Maade, paa hvilken
Utroskab griber ind i den Enkeltes personlige Liv, ville naturligvis
være forskellige efter de psykologiske og sociale Forhold.
Saaledes bemærker \textsc{Burckhardt} (Kultur der Renaissance. 4. Aufl.
II. p. 183 ff.): „Hvad der staar som ejendommeligt for Renaissancens
Italien, er, at man i højere Grad og i hvert Tilfælde
med større Bevidsthed end andre Steder træder Ægteskabet og
dets Ret under Fødder, og at man ligefrem udtaler den Grundsætning,
at Ægteskaber kun skulle sluttes paa Tid, nemlig saalænge
Hustruen behager Manden. De højere Stænders unge Piger
komme ikke i Betragtning; al Lidenskab har gifte Kvinder til
Genstand. Men mærkeligt er det, at Ægteskaberne dog ikke
ses at være aftagne, og at Familielivet langtfra led den Forstyrrelse,
som det under lignende Forhold vilde lide i Norden. . . .
\dots{} Den højere dannede, individuelt udviklede Kvinde raader
over sig selv med en ganske anden Suverænitet end i Norden,
og Utroskaben frembringer ikke noget frygteligt Brud i hendes
Liv, naar hun kan sikre sig mod de ydre Følger. Ægtefællens
Ret til hendes Troskab har ikke den faste Grund, som den faar
hos Nordboerne ved den Poesi og den Lidenskab, Bejlingen og
Trolovelsen fører med sig: efter det flygtigste Bekendtskab og
umiddelbart fra Forældrenes eller Klosterets Værge træder den
unge Kvinde ud i Verden, og nu først udvikler sig hendes
Individualitet, og det med vidunderlig Hurtighed. Især af denne
Grund er Ægtefællens Ret meget betinget.“ --- I de sydlige
% --- p. 297 ---
\setcounter{footnote}{0}%
\opage{297}Lande have i denne Henseende Forholdene ikke ganske forandret
sig siden Renaissancens Dage.

6. Saaledes som det frie Monogami i det Foregaaende er
skildret, er det vel et Ideal, men ikke et, der svæver i Luften.
Ethvert lykkeligt og alvorligt Ægteskab viser os en Tilnærmelse
dertil. Det fyldestgør paa bedste Maade Familiens Formaal,
ikke blot at yde fuld Tilfredsstillelse for Medlemmerne, men
ogsaa at være Sympatiens Arnested og den bedste Skole for
den nye Slægt. Dets Hævdelse og Udvikling bliver derfor en
betydningsfuld Opgave.

Men uagtet det frie Monogami, forsaavidt det virkeligt er
til, staar som en af den historiske Udviklings vigtigste Frugter,
finde vi dog ved Siden af det overalt de lavere, ja de allerlaveste
Former for Kønsforbindelse i fuld Flor. Hetærismen,
Øjebliksforbindelserne og de løse Forbindelser ere ikke forsvundne,
fordi det frie Monogami har vundet Anerkendelse.

Det vilde være en uhistorisk Opfattelse at antage det frie
Monogami ikke blot for den ideale Form for Ægteskab, men
ogsaa for den oprindelige Form, og betragte Hetærismen og de
løse Forbindelser som opstaaede ved Forfald fra og Fornægtelse
af den. Man vil derved ikke blot komme i Strid med Sociologien,
men ogsaa komme i Forlegenhed med at forklare Muligheden
af et saadant Fald; og man vilde staa haabløs overfor
de virkelige Forhold. Har det frie Monogami derimod udviklet
sig historisk, saa have vi blot at gaa videre i samme Retning,
at hævde og udvikle, hvad der er indledet. Hetærismen og de
løse Forbindelser ere da Tegn paa, at Menneskenaturen endnu
ikke er tillempet efter højere sociale Livsvilkaar, --- at den
endnu har for Meget af Dyret i sig til at kunne virkeliggøre
den egentligt menneskelige Form for Kønsliv. Det vil her stedse
vise sig, at de to Bestræbelser: Humaniseringen og Emancipationen
(XIII, 3) høre nøje sammen. Jo brutalere og lavere en
Karakter Kønslivet har, des mere staar det ogsaa tilbage i
Henseende til Personlighedernes Frigørelse. Kun aandeligt umyndige
Individer lade sig bruge som blotte Midler for Andres Lyst.

Det mangler ikke paa Lidelser og Ulykker formedelst denne
ufuldendte Udvikling. Der føres ofte her en fortvivlet Kamp
% --- p. 298 ---
\setcounter{footnote}{0}%
% CHECK p. 298: notes paired with the marks in order (the layer misread a number)
\opage{298}mellem de forskellige Tendenser i den menneskelige Natur. Og
Kampen vil blive lang. Udviklingen gaar paa det etiske Omraade
langsomt fremad, ikke mindst, hvor det gælder om at indordne
en af de mægtigste Naturtilskyndelser under etisk-sociale
Love, der svare til et højere Livstrin end det, paa hvilket Mennesket
skilte sig fra Dyret. Men den værste Løsning af alle
vilde være at slaa af paa Idealet. Det vilde kun være Begyndelsen
til en Tilbagegang. --- Lovbud, Tvang og farisæisk
Moraliseren ville kun udrette lidet. En sund og harmonisk Opdragelse,
en levende Sans for det personlige Livs Ret hos Alle,
en varm Menneskekærlighed ville give den bedste Støtte i
Kampen'. Det gælder især at hindre, at Kønslivet isoleres fra
det øvrige sjælelige Liv. Det skal paa en Gang øve sin dragende
og begejstrende Magt og have sin rette Begrænsning som en
enkelt af Livets Former. Enhver, der i sin Livsførelse stræber
i denne Retning, arbejder med paa et væsentligt Stykke af
Menneskeslægtens Velfærd. ---

Hetærismens stadige Bestaaen er dog, som allerede bemærket,
ikke blot at forklare af, at Instinktet endnu har bevaret
sin ubændige Karakter, men ogsaa af, at materiel Nød og
legemlig og aandelig Forkommenhed hindre et rent menneskeligt
Liv. Dersom det er sandt, at i London hver syvende og i
Hamborg hver niende (yngre) Kvinde er en Skøge\footnote[1]{Oettingen: Moralstatistik. 3. Aufl. p. 197. --- I København er lidt over 1 pCt. af Kvinder mellem 20 og 30 Aar Skøger. M. Rubin i Nationaløkonomisk Tidsskrift 1887. p. 40. --- I Paris er der ca. 5000 indskrevne Skøger, men Tallet paa den hemmelige Prostitution angives til 50,000. E. Pontoppidan: Kønssygdommene. p. 71.}, saa er det
ikke Naturinstinktet, som maa bære hele Skylden, men det er
materielle og sociale Forholds Virkninger, vi her have for os. Vi
staa her overfor en Forgrening af det store sociale Spørgsmaal,
denne Afgrund, i hvilken vi fra saa mange Sider se ned. Det er især
den bitre Nød, som driver saa Mange til at kaste sig seiv bort\footnote[2]{Prosper Despine: Psychologie naturelle. Paris 1868. III. p. 215. (Ofte er det ikke egen Nød, men Forældres eller Børns Nød, som faar Kvinder til at søge Erhverv som Skøger). Kfr. Ed v. Ehlers: Bidrag til Diskussionen om Prostitutionsspørgsmaalet. Kbhvn.' 1896. p. 9: „Den,}.
% --- p. 299 ---
\setcounter{footnote}{0}%
\opage{299}Bedre og selvstændigere Udvikling af de kvindelige Evner vil
baade give Kvinden større Modstandskraft ved at styrke hendes
Selvfølelse, og skaffe hende bedre Udsigter til at fægte sig frem
i Verden. --- Ogsaa for Mandens Vedkommende gribe de sociale
og økonomiske Forhold ind her, idet den Alder, i hvilken Udkommet
tillader Indgaaen af Ægteskab, for Middelstandens Vedkommende
bliver senere og senere\footnote[1]{sociale Lagdeling- Kbhvn. 1890. p. 47 f., 51—53.}.

\chapterhead{XVII.}{OM ÆGTESKABETS INDGAAELSE OG OPLØSNING.}

1. Alt, hvad der deler Menneskene i forskellige Lejre,
kan blive en Hindring for Ægteskab mellem Individer, der ellers
vilde føle Tilbøjelighed for hinanden. Men den egentlige Grund
ligger dog ogsaa i saadanne Tilfælde i de Karakterforskelligheder,
der betinge eller betinges ved de modstridende Retninger. Ti
Erfaringen viser, at hverken religiøse eller nationale eller politiske
Modsætninger ere absolute Hindringer for lykkelig ægteskabelig
Forbindelse. Det kommer ikke saa meget an paa Overensstemmelse
i Anskuelser, som paa Overensstemmelse i Karakterer,
saa at disse enten ligne hinanden eller supplere hinanden. Fra
et strengt teologisk Synspunkt ere „blandede“ Ægteskaber naturligvis
forkastelige. Gamle Kirkelærere betragtede dem som Hor;
Katolicismen forbyder dem; moderne protestantiske Teologer
lade sig nøje med at betegne dem som „monstrøse“. Det kan
heller ikke nægtes, at Muligheden af ægteskabelig Forbindelse
mellem Individer af forskellig Tro tyder paa en Svækkelse af
der tror, at Prostitutionen kan udryddes, maa ogsaa kunne tro paa Fattigdommens
Udryddelse“.

 Rubin og Westergaard: Ægteskabsstatistik paa Grundlag af den
% --- p. 300 ---
\setcounter{footnote}{0}%
% CHECK p. 300: notes paired with the marks in order (the layer misread a number)
\opage{300}den konfessionelle Lidenskab. Men den humane Etik ser ingen
Ulykke i en saadan Svækkelse, naar der kun ikke følger aandelig
Lunkenhed og Blaserthed med. Den betragter saadanne Forbindelser,
hvor de lykkes, som Udtryk for Naturens Sejr over
unaturlige Skranker. Den antager, at der meget vel bag forskellig
Trosbekendelse kan ligge et dybt aandeligt Slægtskab,
og den finder i Elskovsfølelsens Magt til at lade dette dybere
liggende aandelige Slægtskab komme til sin Ret en af de betydningsfuldeste
Sider ved denne mægtige Kraft. To stræbende
Naturer ville forståa hinanden\footnote[1]{Schleiermachers Hustru, der stod i saa inderligt et Forhold til ham, trængte dog til en mere ortodox Kristendom end den, hendes Mand forkyndte. Det gav Anledning til en smuk Tankeudvexling mellem dem. (Aus Schleiermachers Leben. In Briefen. Il. p. 434 —436). Derimod berettes det om den ultramontane Romantiker Gorres' Hustru, at hun aldrig gik i Kirke, og dog var ogsaa her Ægteskabet overordentligt lykkeligt. (Sepp: Görres und seine Zeitgenossen. p. 68). Smlgn. hermed Bjørnsons vidunderligt gribende Skildring af Forholdet mellem Præsten Sang og hans Hustru (Over Evne. I).}. Men paa den anden Side kan
som sagt forskellig Tro hænge sammen med Karakterforskelligheder,
eller Trosforskellighederne kunne være saa dogmatisk
tilspidsede, at et blandet Ægteskab bliver et betænkeligt Vovestykke\footnote[2]{Efter Erfaringer fra Schweiz ere blandede Ægteskaber mere udsatte for Opløsning end Ægteskaber, hvor begge Parter ere Katoliker eller Protestanter. (H. Westergaard i Nationaløkonomisk Tidsskrift. 1887. p. 30). --- En af de Ting, der bragte Ulykke ind i Forholdet mellem Carlyle og hans Hustru, skal have været, at den højt begavede Kvinde vel var enig med sin Mand i hans Kritik af de positive Religioner; men ikke i den Tro, han selv holdt sig til. (Froude: My relations with Carlyle. p. 7 f.). <J. 331 f. --- Rubin og Westergaard (Befolkningsstatistik etc. p. 100 ff.) have paavist, at jo flere Born, der fodtes i et Ægteskab, eller jo hurtigere Børnenes Fødsel følger paa hinanden, des større er forholdsvis Dødeligheden.}.
Sin inderligste Form faar Ægteskabet kun, hvor
Tilnærmelse og Sammenslutning i den hele Maade at opfatte
og føre Livet paa er mulig.

2. Indgaaelse af Ægteskab er naturligvis kun etisk berettiget,
hvor der haves Udsigt til at kunne ernære en Familie.
Men denne etiske Betingelse vil ikke kunne hævdes ad retslig
Tvangs Vej, saaledes som i tidligere Tid forsøgt i nogle tyske
Stater. For det Første er det betænkeligt og utaaleligt at lade
det bero paa en offentlig Autoritets Skøn, om to Individer, der
% --- p. 301 ---
\setcounter{footnote}{0}%
\opage{301}ville gifte sig, have Udsigt til at komme ud af det i økonomisk
Henseende. Den Enkelte kan bedst selv afgøre det, og det er
gavnligt, at han selv faar hele Ansvaret. Der aabner sig desuden
ogsaa Mulighed for alskens Myndighedsmisbrug. Og det
har vist sig, at man ved saadanne Foranstaltninger kun opnaaede
at gøre Ondt værre, idet de løse Forbindelsers og de uægte
Børns Antal toge til *),

Ikke blot for Indgaaelsen af Ægteskab, men ogsaa for de
Børn, som i dette sættes i Verden, have Individerne et etisk
Ansvar. Indre og ydre Forhold kunne gøre det utilstedeligt at
efterlade sig Afkom. Saaledes naar man selv lider af arvelige
Sygdomme (som Sindssygdom, Spedalskhed o. s. v.). Eller naar
de ydre Forhold ere saa udsigtsløse, at Børnene fordømmes til
Lidelse og Undergang. Den store Dødelighed blandt nyfødte
Børn, som man har kaldet den moderne Moloksdyrkelse -), kan
kun hindres, naar der ikke sættes fl ere Børn i Verden, end der
kan fin de Pleje og Næring under de givne Forhold. Børn blive
jo dog ikke til af sig selv uden Forældrenes Vidende og Villie.
Hvorledes den Enkelte kan fyldestgøre det Ansvar, som her
paahviler ham, er et Spørgsmaal, der snarere hører under Medicinens
end under Etikens Afgørelse.

3. En dybt liggende Følelse, hvis første Kilde og Motiv
det ikke er let at paavise, opfylder nu de fleste Mennesker
med Afsky og Rædsel ved Tanken om en kønslig Forbindelse
mellem Forældre og Børn eller mellem Søskende indbyrdes.
Hos vilde Folk ere derimod saadanne Forbindelser hyppige\footnote[1]{Spencer: Principles of Sociology. I. p. 636 f. 680.};

*) Smlgn. E. Löning: Armenpflege und Armenpolizei. (Schonbergs
Handbuch der politischen Oekonomie. 1. Udg. II), p. 581. 595. --- Stuart
Mill sympatiserer mærkværdigvis med saadanne Forbud og finder i dem
Intet, der gaar Individets Frihed for nær, skønt han antyder, at de maaske
paa Grund af stedlige Forhold ikke altid vise sig praktiske. (Om Friheden.
Dansk Overs. p. 196).

-) Rümelin: Reden und Aufsätze. Freiburg u. Tiibingen. 1875.
% --- p. 302 ---
\setcounter{footnote}{0}%
\opage{302}dog fremtræder tidligt forskellig Skik hos forskellige Stammer\footnote[1]{(Die primitive Familie. p. 238---247). De gamle Forbud udledes af, at man vilde undgaa den Forvirring i de paa Slægtskabet hvilende juridiske Forhold, som vilde opstaa ved Ægteskab mellem Beslægtede. Derved forklares dog ikke den stadige Eftervirken og Bestaaen af den Følelse, der oprindeligt skulde skyldes rent økonomiske Forhold.}.
Strenge Straffe for ægteskabelig Forbindelse med Moder, Datter
og Svigerdatter fastsættes af den babyloniske Lovgiver Hammurabi
(2200 Aar f. Kr.); ligeledes i Mose Lov, hvor ogsaa Forholdet
til Søsteren forbydes\footnotemark[1].

Sociologisk kan hin Følelse maaske sættes i Forbindelse
med den hos saa mange Folk herskende Skik, at Hustruen
skal hentes fra en anden Stamme end Mandens egen; men
selve denne Skiks (Exogamiets) Opstaaen er endnu ikke ret
forklaret. I den indiske og i den middelalderlige Lovgivning
findes som bekendt strenge Forbud mod Ægteskab ogsaa med
fjernere Slægtninge indtil sjette eller syvende Grad. Efter den
i protestantiske Lande raadende Skik er kun Ægteskab mellem
de allernærmeste Slægtninge utilstedeligt. Hvilken saa den
sociologiske Forklaring heraf maatte være\footnotemark[1], vil den etiske Begrundelse
findes i følgende Betragtning.

Forholdet mellem Søskende og mellem Forældre og Børn
vilde miste sin frie og trygge Karakter, dersom Muligheden af
et kønsligt Forhold og de dermed forbundne Lidenskaber var
tilstede. Den fuldstændige Tillid, som hint Forhold maa hvile
paa, hvis det skal have sin fulde Værdi, vilde mangle, naar
Elskovdriften trængte sig ind her. Under det nære og stadige

:) Snorre Sturleson: Ynglingesaga. Kap. 4: \textasciitilde{}„Da Njord var hos
Vanerne, havde han sin Søster til Ægte; ti saa tillod Loven der\dots{}. Men
hos Aserne var det forbudt at have saa nære Frænder til Ægte.“

) Smlgn. Fr. Buhl: Kong Hammurabis Lovsamling. Nordisk
Tidsskrift 1903. p. 468.

\footnotemark[1] Det er neppe rimeligt, at det er Erfaringer om Afkommets Svækkelse
ved Ægteskab mellem Nærbeslægtede, som have virket bestemmende.
En saadan Svækkelses Nødvendighed er det ikke hidtil lykkedes at paavise.
--- Findes der først i en Familie sygelige Anlæg, ville disse naturligvis
forplantes og forøges ved gentagne Forbindelser mellem Familiemedlemmer.
--- Et Forsøg paa sociologisk Forklaring giver Starcke
% CHECK p. 302: note mark without a note

% --- p. 303 ---
\setcounter{footnote}{0}%
\opage{303}Samliv vilde der hyppigt blive Anledning til lidenskabelige Forhold,
der kunne virke oprivende og opløsende\footnote[1]{Smlgn. allerede Maimonides, en jødisk Filosof i det 12. Aarh. (Ståudlin: Geschichte der Vorstellungen und Lehren von der Ehe. Göttingen 1826. p. 457). Smlgn. af senere Forfattere Grotius: De jure belli et pacis. 11, 12—13. Hume: Enquiry concerning the principles of Morals. Sect. IV.}. En sjælelig
Umulighed vil her være det bedste Værn. Og først ved en saadan
Umulighed bliver det ejendommelige, rige og inderlige Forhold
mellem Broder og Søster til; en ny Værdi vindes for Livet.
Ogsaa Forholdet mellem Fader og Datter, Moder og Søn bliver
først udformet i sin Ejendommelighed, naar det træder frem i
afgjort Forskellighed fra Kønsforholdet. --- Og paa den anden
Side er der noget Naturligt i, at Elskovsaffekten ikke opstaar
under den hjemlige Tilvanthed og Fortrolighed, den Følelse af
oprindelig Sammenhøren, som raader i Forholdet mellem Forældre
og Børn og mellem Søskende indbyrdes. I Elskovsaffekten
ytrer sig en Trang til Supplering, som bedst tilfredsstilles
ved Forbindelse med et Individ af en anden Familie. En Affekt
opstaar ved det hidtil Ukendte, ved det, der bringer noget hidtil
Uerfaret. Derfor udelukkes Elskovsaffekten under sædvanlige
Omstændigheder mellem Broder og Søster, der ikke føle sig som
saa forskellige, og hvis Forhold ikke først skal indledes og stiftes
ved Vækkelsen af en særlig Drift. Karakteristisk for Elskovsfølelsen
er det, at en fremmed Personlighed opdages; derfor
stiftes Elskovsforholdet ved Erklæring og Valg. Søskendeforholdet
og Forholdet mellem Forældre og Børn skyldes derimod
ikke Opdagelse og Valg, men er givet som et Naturforhold.

4. Ægteskabets Indgaaelse medfører naturligt nye Pligter
og Rettigheder for Parterne i deres Forhold til hinanden og
til Andre. Det Fællesskab, som stiftes, omfatter alle Sider af
Livet, og det er derfor af Vigtighed, at det bliver konstateret.
Dannelsen af en ny social Celle interesserer Flere end de to
Individer selv. Ikke blot deres Slægt, men ogsaa Staten, hvis
Opgave det kan blive at hævde Rettigheder, som stiftes ved
Ægteskabet, er interesseret i, at Forbindelsen officielt erklæres.
Dette krænker heller ikke Ægteskabets Frihed. Det giver det
% --- p. 304 ---
\setcounter{footnote}{0}%
\opage{304}kun en ydre Ramme. Hvor den personlige Følelse har talt og
knyttet Forbindelsen, dér vil det kun være at tage alle Konsekvenserne,
naar den ydre sociale og retslige Side bringes i Orden,
og ingen nok saa romantisk Frihedsfølelse kan krænkes derved.
Oppositionen mod en offentlig Erklæring om Indgaaelsen af et
ægteskabeligt Forhold stammer især fra, at det synes at ligge
i de sædvanlige Former --- baade ved borgerlig og kirkelig
„Vielse“ --- at Forbindelsen skulde have sin Sanktion udefra
og ikke være helliget ved selve den inderlige Følelse, der har
grundlagt den.

Ægteskabets etiske Gyldighed beror ikke paa Samfundets
og Statens „Sanktion“ af det. En saadan Opfattelse kan kun
hævdes fra Autoritetsprincipets Standpunkt. Ægteskabet faar sin
etiske Gyldighed ved den frie Erklæring mellem de to Personer
selv, og Statens Opgave kan kun være den at fastsætte og hævde
de Rettigheder og Pligter, de derved overtage. Det er en Forvexling
af Skal og Kærne, naar man betragter den officielle
Konstatering af Ægteskabet som dets etiske Stiftelse. Dette gælder,
hvad enten den officielle Konstatering sker paa kirkelig eller
paa borgerlig Vis. Der kan gives et Ægteskab i Aand og Sandhed
uden disse ydre Former, og det vilde være Barbari, om
man --- som det er blevet foreslaaet --- ved Politiets Magt vilde
paatvinge Nogen saadanne Former. Trangen til at konstatere Forholdets
Indgaaen hang oprindeligt sammen med Ægteskabets
sociale Betydning som Grundlæggelse af en ny Husstand. Med
Civilrettens Udvikling blev den offentlige Erklæring i mundtlig
eller skriftlig Form Hovedsagen, og Brudens Overførelse til Brudgommens
Hus, der i ældre Tider var Tegnet paa den nye Husstands
Grundlæggelse og Hovedpunktet ved Bryllupet, blev efterhaanden
blot til en festlig Skik\footnote[1]{Leist: Alt-Arisches Jus Civile. II. Jena 1896. p. 120 ff. risk Tidsskrift. Femte Række. Første Bind). ---}. Den civilretlige Akt, der
traadte i Stedet for den uvilkaarligt udviklede Form for Konstateringen
af Ægteskabets Indgaaen, fik ikke strax en rent borgerlig
Form, men omgaves med religiøse Symboler. Den kristne
Kirke har dog længe fastholdt Mindet om den mere primitive
Form for Ægteskabs Indgaaelse, idet den i lange Tider ikke
% --- p. 305 ---
\setcounter{footnote}{0}%
\opage{305}betragtede den kirkelige Akt som nødvendig Betingelse for
gyldigt Ægteskab. Adams og Evas Ægteskab, Forbilledet for
alle andre Ægteskaber, var jo indgaaet uden Vidner! Efter den
katolske Lære var Ægteskabet et Sakrament, der forrettes
af Mand og Kvinde, selv naar de indgaa deres Forbindelse
uden Præst og uden Vidner; og det var derfor en Inkonsekvens,
naar Tridentinerkoncilet fastslog den kirkelige Vielses
Nødvendighed ').

De Bestemmelser, som Statens Love indeholde angaaende
Ægteskabet, angaa ikke blot Betingelserne for, at retslige Virkninger
kunne flyde af det, men sigte ogsaa til at fastsætte Forholdet
mellem Ægtefællerne indbyrdes. De bære endnu i høj
Grad Præget af den gamle Forestilling om, at Manden er Kvindens
Herre, og at hun er hans Tugt undergiven.

I det frie Monogamis Natur ligger det, at Kvindens retslige
Stilling i Ægteskabet maa være sideordnet med Mandens.
Der bør ikke uden hendes Samtykke kunne raades over hendes
Arv eller erhvervede Ejendom. Og hun bør have Medbestemmelsesret
ved alle vigtige Afgørelser, der angaa fælles Interesser.
Faktisk øver hun vel i Regelen en stor Indflydelse og har et etisk
Ansvar for de Beslutninger, som tages; men dette Ansvar vil
skærpes, naar hun har juridisk Medbestemmelsesret. Hvor hun
ikke kan bifalde Mandens Dispositioner, bør hun idetmindste
kunne sikre sine Børn og sig selv mod truende Fordærv\footnote[1]{Hvor Meget Kvinden kan udrette selv uden juridisk Ret, ses af}. Jo

') Denne Inkonsekvens blotter Paola Sarpi paa en meget klar
Maade i sin Tridentinerkoncilets Historie. (Fransk Overs. ved Amelot
de la Houssaie. Amsterdam 1686. p. 765). Kirkelig Indvielse af Ægteskabet
blev tidligt Skik i den gamle Kirke; fastsat ved Lov blev den
først af de østromerske Kejsere for den græske Kirkes Vedkommende
og af Tridentinerkoncilet for Romerkirkens Vedkommende. --- At hemmelige
Ægteskaber ogsaa i den første protestantiske Kirke betragtedes
som gyldige, ses af Peter Pl ad es Visitatsbog (udg. af S. Grundtvig,
p. 82 f.), hvor den offentlige Konstatering i Kirken blot anbefales ved
Vigtigheden i retslig Henseende af, at der er Vidner om Ægteskabets
Indgaaelse. --- Smlgn. om dette og dermed beslægtede Spørgsmaal:
J. Nellemann: Retshistoriske Bemærkninger om kirkelig Vielse. (Histo%
% --- p. 306 ---
\setcounter{footnote}{0}%
\opage{306}mere aandeligt udviklet hun bliver ved en forbedret Opdragelsesmaade,
jo større Selvstændighed og praktisk Erfaring hun vinder
ved rigere Lejlighed til at lære Livet omkring sig at kende,
--- des bedre vil en saadan Ordning kunne bestaa uden at medføre
Hindring for fornuftige og berettigede Foretagender. --- Naar
Hustruens Rettigheder fastsættes ved Ægteskabets Indgaaelse som
noget Selvfølgeligt, vil der ikke ligge Noget deri, som kan krænke
det inderlige Forhold mellem Ægtefællerne, og der vil dog være
tilvejebragt et Værn, dersom Uenighed og Misstemning skulde
indtræde. Man afviser ofte de Tanker, her ere udtalte, dermed,
at Ægteskabet skal være et inderligt Tillidsforhold, ikke et Retsforhold;
Enheden skal være det Afgørende, ikke Forskellen. Men
et Retsforhold er det dog ogsaa, naar al Myndighed tillægges
den ene Part. At slaa denne Myndighed fast for Manden, er
jo dog at betone Forskellen mellem Ægteskabets to Parter langt
stærkere, end det vilde ske, naar der i økonomisk og juridisk
Henseende tildeltes dem lige Myndighed. Det kan være, at en
saadan Ordning har sine tekniske Vanskeligheder; men det vil
sikkert lykkes Fremtiden at komme ud over disse\footnote[1]{Smlgn. N. Tsakni: La Russie sectaire. Paris 1888. p. 153 f.}. Hvad den
aandelige Myndighed angaar, kan den naturligvis hverken fastsættes
ved Tradition eller ved Lov, og her staar Myndighedsfordelingen
da ogsaa ofte (og vil mere og mere komme til at
den Maade, hvorpaa Arbejdernes Hustruer i Rochdale benyttede den
berømte Brugsforening der. „Mange gifte Kvinder blive Medlemmer,
fordi deres Mænd ikke ville gøre sig den Ulejlighed; andre træde ind
i den til Selvforsvar, for at hindre Mændene i at øde deres Penge til
Drik. Manden kan ikke tage de Sparepenge, som staa paa Konens
Navn, ud, naar hun ikke undertegner Anvisningen. Vel kan efter den
gældende Lov Manden ved en lovlig Proces komme i Besiddelse af
Pengene; men en Proces tager Tid, og før Loven kan sættes i Gang,
bliver Manden ædru og kommer paa bedre Tanker.“ (Stuart Mill:
Principles of Political Economy. IV, 7, 6).

') Hos en russisk Sekt, Molokanerne, er der Ejendomsfællesskab
mellem Mand og Hustru, og Manden kan ikke uden Hustruens Samtykke
pantsætte eller sælge Ejendelene. Den molokanske Kvinde vænnes fra
den spæde Barndom af til Uafhængighed, og hendes aandelige Liv udvikles
(indenfor Sektens puritanske Ramme) paa selvstændig Maade.
% --- p. 307 ---
\setcounter{footnote}{0}%
\opage{307}staa) i et komisk Misforhold til den traditionelle og juridiske Fordeling
af Myndigheden udad til. --- løvrigt er der i den nyeste
Tid i Tyskland --- det Land, hvor der gøres størst Modstand mod
Kvindens Selvstændighed --- gjort nogle ikke ubetydelige Fremskridt
med Hensyn til Kvindens Retsstilling indenfor Ægteskabet
ved Udarbejdelsen af den nye borgerlige Lovbog for det tyske
Rige. Vel forkastedes et Forslag om, at ved Meningsforskelligheder
i økonomiske Sager skulde den af Ægtefællerne have Afgørelsen,
af hvis Formue Familiens Omkostninger for Størstedelen
bestredes, og samme Skæbne havde et Forslag om, at Hustruen
skulde have Ret til at beholde sit Familienavn. Men gennem
Meningsfæller i Lovkommission og Rigsdag tilkæmpede de tyske
Kvinder sig dog Retten til at kunne blive beskikkede til Formyndere
og at være Medlemmer af Familieraadet. Tillige blev
det fastsat, at Manden kun med Formynderrettens Billigelse kan
forbyde Hustruen at udøve et Kald eller et Haandværk. Og om
selve den Indflydelse, de tyske Kvindeforeninger gennem deres
Forsamlinger og Petitioner have øvet paa den Form, i hvilken
Rigets nye borgerlige Lovbog blev vedtagen, skriver en tysk
Kvinde, der selv er Doktor i Lovkyndighed: „Aldrig før have
Kvinderne taget saa livlig Del i Lovgivningen som i Anledning
af Forhandlingerne om Udkastet til den borgerlige Lovbog. Aldrig
før have de udøvet saa stor Indflydelse paa Formuleringen
af de Retsbestemmelser, som angaa dem\dots{}. Det er vist et
enestaaende Tilfælde i den moderne Kulturhistorie, at hver enkelt
af Kvindernes og deres Talsmænds Fordringer blev forhandlet
og prøvet i Kommissionerne og fandt Udtryk i Loven, forsaavidt
de kunde forenes med Medlemmernes Anskuelser“\footnote[1]{Dr. jur. Emilie Kempin: Die deutschen Frauen und das biirgerliche Gesetzbuch. (Schweizerische Blätter für Wirtschafts- und Socialpolitik. 1896. p. 679).}. Det
er dog kun et lille Skridt, der her er gjort fremad i Retning
af det, som Retfærdigheden og som den voxende Udvikling af
Kvindens Evne til praktisk og teoretisk Arbejde kræver. Tyske
Kvindeforeninger have energisk protesteret mod den Begræns%
% --- p. 308 ---
\setcounter{footnote}{0}%
\opage{308}ning af Kvindens Rettigheder og Pligter, som den nye tyske
Lovbog fastholdt.

5. At Monogamiet hersker i Europa, skyldes især den græske
Filosofi, Romerretten og Kristendommen. Denne sidstes Indflydelse
har dog været mere indirekte; thi intetsteds i det nye Testamente
er det gammeltestamentlige Polygami ophævet. Luther og
Melanchton vare ogsaa klare over, at de som Teologer ikke kunde
modsætte sig Dobbeltægteskab; det var i deres Øjne den verdslige
Øvrigheds, ikke Kirkens Sag at ordne Ægteskabsforholdene\footnote[1]{--- Hos Jøderne skal Polygamiet have holdt sig til omtrent Aar 1000. Det blev da opgivet --- „som en Indrømmelse til Occidenten.“}

Vort Ægteskab er det romerske, med Afskaffelse af den vil
kaarlige Ret til Skilsmisse, som fra Slutningen af Republikens
Tid spillede saa stor en Rolle\footnotemark[1]. Efter den strenge kristelige
Lære er Skilsmisse utilstedelig, hvor der ikke foreligger fysisk
Utroskab. Men det gaar her en overspændt idealistisk Retning,
som saa ofte, at den. netop kommer til at lægge overdreven og
unaturlig Vægt paa den rent fysiske Side af Sagen. Der kan
være en Utroskab i Sind og Villie, der er fordærveligere for det
hele Forhold end den fysiske Utroskab. Ægteskabet kan efter
dets indre Aand og Kærne være opløst ved, at den Følelse, der
skulde bære og besjæle det, er borte. Naar Livet er borte, kan
den ydre Form vel holdes oppe ved Tvang, men den har da
ingen Værdi. Ærlighed og Oprigtighed ere Livsbetingelser i det
ægteskabelige Forhold; hvor de mangle, kan den frie, ubetingede
Hengivelse ikke finde Sted. Resignation og Selvbeherskelse kunne
udrette store Ting; men der er en Grænse for deres Evne, og
naar selve den besjælende Følelse er helt udtørret, kan den ikke
erstattes ved noget Andet. Der behøver, som allerede tidligere

:) Smlgn. Hausrath: Kleine Schriften religionsgeschichtlichen
Inhalts. p. 239 f.

\footnotemark[1] Smlgn. Henry Maine: Early History of Institutions. p. 60.
Lecky: History of European Morals. II, p. 316. Renan: Marc
Auréle. p. 547. Leist: Alt-ar. Jus civ. I. p. 438. Robert Bartsch
(Die Rechtsstellung der Frau als Gattin und Mutter. Leipzig 1903.
p. 39) fremhæver den græske Indflydelse: „Es ist sicher, dass das auf
hellenistischem Milieu entstandene Christenthum wesentliche Stücke seines
familienrechtlichen Inhalts der griechischen Philosophie verdankt.“
% CHECK p. 308: note mark without a note

% --- p. 309 ---
\setcounter{footnote}{0}%
\opage{309}(XVI, 3) vist, ikke at være Skyld paa nogen af Siderne, og
Forholdet kan dog blive grundulykkeligt. De ulykkelige Følger
strække sig over hele Hjemmet. Fjernhed og Kulde eller endog
Bitterhed og Fjendskab mellem Forældrene vil for Børnene
være den fordærveligste Atmosfære at opvoxe i.

Dersom den retslige Ordning skal være en Støtte for det
etiske Familieliv, maa den derfor ikke sætte for store Hindringer
for Skilsmisse. Derved vilde den nedværdige Ægteskabet
til en Tvangsinstitution. Af Hensyn til de forskellige Individers
Interesser maa der være foreskrevet bestemte Former for Ægteskabets
Opløsning saavel som for dets Indgaaelse. Ved disse Former
beskyttes Ægtefællerne ikke blot mod hinandens Overgreb,
men ogsaa hver for sig mod sin egen Overilelse. Der bør gaa
en vis Tid, inden den fuldstændige Skilsmisse indtræder; og
denne Tid bør vare længere, naar det kun er den ene Part, der
ønsker Skilsmissen. Troskaben har sin Ret og bør have Tid til
at prøve sin Magt. Det var en for vidt dreven Individualisme,
naar \textsc{Wilhelm} v. \textsc{Humboldt} i sit Ungdomsskrift om Statens
Grænser ansaa den ene Parts erklærede Ønske for tilstrækkeligt
til at ophæve Ægteskabet.

Der vil i den Lidelse, et ægteskabeligt Liv, som ender med
Skilsmisse, medfører, ligge Advarsel nok for dem, der paa Grund
af Letheden ved at opnaa Skilsmisse vilde vise Letsindighed i
at indgaa Ægteskab, saa man behøver ikke at frygte for, at en
liberal Skilsmisselovgivning skulde svække den almindelige Agtelse
for Ægteskabet. I Staten Indiana, som i mange Aar har
haft liberale Skilsmisselove, er Familielivet ligesaa smukt og
varigt, som i New York og Syd Karolina, hvor det kun kan opløses
af én Aarsag. Hos nogle russiske Sekter (Doukhoborerne,
Molokanerne), hvor Skilsmisser ere hyppige og meget lette at
iværksætte, hersker der høje og ædle Forestillinger om Ægteskabet,
og idethele skal hos dem det moralske Niveau være højere
end hos deres ortodoxe Naboer\footnote[1]{History of Women Suffrage. Boston 1878. I. p. 742. --- Lecky: History of European Morals. Il. p. 325 f. --- Tsakni: La Russie sectaire. p. 143 f. 152. ---}. Strenge Ægteskabslove (særligt:
% --- p. 310 ---
\setcounter{footnote}{0}%
% CHECK p. 310: notes paired with the marks in order (the layer misread a number)
\opage{310}Forbud eller Hindringer mod Fraskiltes Giftermaal) have aldrig
vist sig istand til at hævde og fremme Sædeligheden. De kue
og binde de rene og ædle Naturer; de letfærdige lade sig ikke
binde, men finde altid deres Udveje. I katolske Lande som f. Ex.
Italien, hvor Skilsmisse ikke er tilladt, ere de ægteskabelige Forhold
i mange Kredse i højeste Grad nedværdigende, saa at man
forlanger Ret til Skilsmisse for at værne om Ægteskabets Hellighed
og Værdighed\footnote[1]{Mantegazza: Die Physiologie der Liebe. Kap. 20 og 21.}. Skilsmissens Hyppighed afhænger, efter
en fransk Statistikers Mening, ikke af Lovene, men af Sæder,
Religion, Race, Stand. Skilsmisser synes tillige at være hyppigst
i de Lande, hvor ogsaa Selvmord er hyppigst, og hvor der altsaa
findes flest Mennesker, der have Vanskelighed ved at bevare
Ligevægten i Livets forskellige Forhold (individus mal equilibrés).
Det er hyppigst Kvinden, som ønsker Skilsmisse; og denne er
forholdsvis sjælden, hvor der er Børn\footnote[2]{J. Bertillon i „Dictionnaire des sciences anthropologiques“. p. 384—385. --- Smlgn. hermed H. Westergaards Afhandling i „Nationaløkonomisk Tidsskrift“, 1887. --- W. F. Willcox: The Divorce Problem. New York. 1891.}.

\grouphead{2. KVINDENS STILLING OG VILKAAR.}

\chapterhead{XVIII.}{SOCIOLOGISKE DATA.}

1. fcn Undersøgelse af Kvindens etiske Stilling danner et
naturligt Supplement til den foregaaende Lære om det frie Monogami.
Det viste sig, at dette først ret virkeliggøres, naar ikke
blot Manden, men ogsaa Kvinden har Evne og Ret til at føre
% --- p. 311 ---
\setcounter{footnote}{0}%
\opage{311}en selvstændig menneskelig Tilværelse. Hun kan kun fyldestgøre
sin Plads i Ægteskabet, naar hun har Muligheden for ogsaa
at kunne udfylde en Plads udenfor det. Kun da kan hendes
Valg blive frit, naar begge Muligheder frembyde sig for hende.
Hvad hun saa vælger, vælger hun det ikke blot som Middel for
at sikre Existensen, men som en Livsopgave. --- Den Tid vil
sikkert komme, da det vil være unødvendigt at gaa ind paa en
nærmere Begrundelse af Kvindens Ret til selvstændig Udvikling
og til frit Valg af Livsopgave. Ligesaa lidt som vi i Etiken
nu have et særligt Afsnit om Mandens Stilling og Vilkaar, vil
man i Fremtidens Etik behøve et Afsnit om Kvindens Stilling
og Vilkaar. Men endnu er den Forestilling ikke forsvunden, at
Kvinden ved sin Natur kun er anlagt til at være Hustru og
Moder, og at enhver Bestræbelse for at muliggøre en Virksomhed
af anden Art for hende beror paa en betænkelig Vildfarelse.

2. Med Hensyn til Kvindens Natur staa to Anskuelser
skarpt overfor hinanden, den ene (som i Oldtiden hyldedes af
Platon, i Nutiden af Stuart Mill), at Forskellen mellem Mandens
og Kvindens Natur og Evner kun er en Gradsforskel, om den
idethele er tilstede, --- den anden (som bl. a. hyldes af Spencer),
at Forskellen fra Naturens Haand er saa dyb og fast, at
den stedse maa betinge en Artsforskel mellem de to Køns
Stilling og Virksomhed.

Man maa stedse være forsigtig med at tale om Naturejendommeligheder
og Naturforskelligheder, som om de vare evige
og uforanderlige\footnote[1]{De nyeste Undersøgelser om Mandens og Kvindens Natur findes benyttede i A. Clod-Hansens interessante Bog Mand og Kvinde. Kbhvn. 1895. Og det fremhæves netop her paa flere Punkter, hvor vanskeligt det er at afgøre, hvad der skyldes Natur, og hvad der skyldes Opdragelse og Øvelse.}. Naturen er ien stadig Udvikling, især paa
de levende Væseners Omraade. Den Natur, vi staa overfor, er
stedse selv bleven til, og en ny Natur vil igen udvikle sig af
den. Men paa den anden Side gaar denne Udvikling langsomt
for sig, og vi maa vel agte paa, hvor langt vi til den givne
Tid ere komne.

% --- p. 312 ---
\setcounter{footnote}{0}%
\opage{312}Derhos staa Kvindens Natur og hendes Vilkaar i Vexelvirkning
med hinanden. Hendes Vilkaar bestemmes ved hendes
Natur, men virke ogsaa tilbage paa denne. Derfor beror hendes
Natur for en stor Del paa, hvad der fordres af hende, og
hvad man giver hende Ret til. Hvad det her, som overalt,
gælder om, er at afpasse Fordringer og Rettigheder saaledes,
at de lade Naturen ske Fyldest og tillige udvikle den i den
heldigste Retning.

3. Vi finde, naar vi betragte Kvindens Natur og Stilling
til forskellige Tider og paa forskellige Steder, forskellige Former
for Arbejdsdelingen mellem Mand og Kvinde.

Paa det mest elementære Trin af Udvikling, som kendes,
er der ingen Arbejdsdeling. Mand og Kvinde sørge hver for
sig for Livets Ophold, og kun Kønsforholdet bringer dem ien
kortere eller længere Forbindelse\footnote[1]{sørgeligste Emne i Historien, og at den uskrevne Historie her --- hvis vi kendte den --- vilde vise sig at være endnu sørgeligere. Kannibalisme og Mishandling af Fanger have kun lejlighedsvis fundet Sted; men den brutale Mishandling af Kvinder er universell}. Naar der behøves mere forskelligartet
Arbejde for at leve, læsser den Stærke naturligt det
ubehageligste Arbejde over paa den Svage, og Kvinden bliver
da den første Slave. Hun maa besørge alt det Arbejde, som
Manden ikke vil paatage sig. Hun maa bære og slæbe, sørge
for Bolig, Klæder og Jordens Dyrkning. Hun vurderes først
og fremmest efter sin Arbejdskraft\footnotemark[1]. Allerede som Barn sættes
hun til Side for Brødrene. Skønt hun er det svage Køn,
maa hun dog bruge sine Kræfter ret forsvarligt. Men hun raader
ogsaa over mere Kraft end i det civiliserede Liv. For blot at
anføre et Exempel: en Indianerkvinde forløser hyppigt sig selv
og tager strax efter sin Nedkomst igen fat paa sit strenge Arbejde.
Efter nogle Forskeres Mening var hos de første Mennesker,
hvis Spor vi finde, Forskellen i Bygning og Styrke ikke

') Smlgn.Westermarck: Geschichte der menschlichen £7zé>\footnotemark[2]. p.218f.

\footnotemark[1] Se f. Ex. Th. Waitz: Die Indianer Nordamerika's. p. 99. ---
Fr. Müller: Allgemeine Ethnographie. p. 161. 288 ff. --- Spencer
(Principles of Ethics § 428) paastaar, at Kvindernes Behandling er det
% CHECK p. 312: note mark without a note
% CHECK p. 312: note mark without a note
% --- p. 313 ---
\setcounter{footnote}{0}%
\opage{313}saa stor, som den senere --- gennem forandret Arbejdsdeling
--- er bleven.

En ny Arbejdsdeling træder i Kraft, naar det haardeste
Arbejde kan læsses over paa den overvundne, men skaanede
Fjende. Kvinden faar da kun Arbejdet i Husets Indre. Hun
faar da sit eget Omraade, som hun indtil den nyeste Tid har
beholdt under forskellige Forhold, alt eftersom Polygami eller
Monogami var den herskende Form for Ægteskab.

Endeligt er der i vore Dage i adskillige Lande aabnet
Adgang for Kvinder til en hel Række af Virksomheder, som
man tidligere kun kunde tænke sig besørgede af Mænd. De
arbejde som Læger, Advokater, Præster, videnskabelige Forskere,
Ingeniører, Kontorister, Stationsforvaltere o. s. v., og de have
faaet Valgret til kommunale og politiske Forsamlinger. Der er
hermed banet Vej for en helt anden Udvikling og helt andre
Vilkaar for Kvinden, end ældre Tider ansaa for de eneste
naturlige. Dog rejses der endnu af Mange Tvivl om denne
sidste Arbejdsordnings Nødvendighed og Gavnlighed.

Det vil saameget desmere være af Interesse at undersøge
dette Spørgsmaal, som Kvindens Stilling kan siges at afgive en
Maalestok for det etiske Fremskridt i Menneskeslægten. Grækeren
roste sig overfor Barbaren af sin Behandling af Kvinden;
Romeren roste sig af det Samme overfor Grækeren, og den
Kristne igen overfor Romeren.

Dog viser Historien, at Kvindens Stilling ikke blot bestemtes
til de forskellige Tider ved den økonomiske og sociale Arbejdsdeling,
men ogsaa ved den almindelige Dannelses og Aandslivets
hele Karakter.

Indenfor den græske Kultur have især de stoiske Filosofer
hævdet Kvindernes etiske Ligeberettigelse med Mændene. Selvom
nogle Opgaver laa nærmere for Manden, andre for Kvinden, saa
gaves der dog --- mente de --- ingen berettiget Stræben, fra
hvilken det ene eller det andet Køn skulde være udelukket.
De senere Stoikere vilde have de unge Kvinder, ikke blot de
unge Mænd underviste i Filosofi. Det var disse samme Filosofer,
der havde den ædleste Opfattelse af Ægteskabet som et inderligt
Livsfællesskab; en videre gaaende aandelig Udvikling var i
% --- p. 314 ---
\setcounter{footnote}{0}%
\opage{314}deres Øjne en Betingelse for, at Kvinden paa rette Maade kunde
udfylde sin Stilling i sin huslige Kreds. Heri laa et betydeligt
Fremskridt, naar man erindrer, at de ateniensiske Kvinder vare
udelukkede fra aandelig Uddannelse, saa at Mændene kun i
Hetærernes Kredse kunde tilfredsstille Trangen til Omgang med
dannede Kvinder.

Ogsaa Kristendommen bidrog til en højere Stilling for Kvin
den gennem et dybere gaaende aandeligt Fællesskab. Tro og
Haab vare jo fælles, og Kvindens Liv saavel som Mandens blev
baaret af Forventningen om de store Ting, der skulde ske, og
snart ske. Forventningen om „denne Verdens“ snare Afslut
ning gav alle Enkeltes Liv et højt, af deres særegne Kaar ganske
uafhængigt Maal. Kvinden kunde da, selvom hun ikke var Hustru
og Moder, naa sin personlige Bestemmelse. Af asketiske Grunde
havde Urkristendommen endog Forkærlighed for ugift Stand,
især for Jomfruelighed. Ganske afset fra de her til Grund lig
gende Motiver var det af stor Betydning, at Kvindens Liv fik
en stor, af hendes Stilling til Familien uafhængig Opgave. Overfor
denne principielle Opfattelse var det af ringere Vægt, at
Kvindens Underordning paa orientalsk Vis blev strengt fastholdt,
og at Kvinden skulde tie i Menigheden og lade sig belære
hjemme af sin Mand. Hovedsagen var det store Princip, Urkristendommens
uforgængelige Hæder, at det Højeste stod aabent
for Alle, uden Forskel mellem Køn, Stand og Race.

I Renaissancen blev den nye, store Strøm af Dannelse Kvinderne
til Del, saavel som Mændene, og den stærke Udvikling
af Individualiteterne foregik i den kvindelige Verden saavel som
i den mandlige. \textsc{Burckhardt} (Kultur der Renaissance in Italien.
4. Aufl. I. p. 124) siger herom: „Om en særlig, bevidst Emancipation
er der ikke Tale, da Sagen forstod sig af sig selv.
Den Kvinde, som var af Stand, maatte den Gang ganske ligesom
Manden stræbe efter en afsluttet, i enhver Henseende fuldendt
Personlighed. Den samme Proces i Aand og Hjerte, der
gør Manden fuldkommen, skulde ogsaa gøre Kvinden fuldkommen.“
Ejendommeligt for denne Tid var det, at denne Udvikling
af Personligheden først blev mulig for den gifte Kvinde.
De unge Piger opdroges for en stor Del i Klostrene, og deres
% --- p. 315 ---
\setcounter{footnote}{0}%
\opage{315}aktive Dannelse begyndte først efter Ægteskabet. Der herskede
i denne Henseende --- og hersker endnu i de romanske Lande
--- en mærkelig Dualisme i Kvindens Stilling.

--- Saa Meget se vi af disse historiske Notitser, at naar der
rører sig store Ideer, som gribe dybt ind i Menneskelivet, rører
der sig ogsaa en Tendens til at betragte Mænd og Kvinder som
ligestillede. Overfor en stor Udvidelse af den aandelige Horizont
forsvinder den store Forskel, man ellers havde gjort mellem
de to Køns Stilling.

\chapterhead{XIX.}{KVINDENS ETISKE STILLING.}

1. Det er af Naturen betrot Kvinden at bære Spiren til
den nye Slægt i sig og nære den, indtil den kan føre sit eget
Liv. At dette Naturhverv er bestemmende for hendes hele
Organisation, kan der neppe være Tvivl om. Der kan da ikke
blive al den Energi tilovers til andre Funktioner, som Manden
kan have til sin Raadighed. Herved bliver Kvinden det svage
Køn. Og herved bliver hendes Natur ogsaa renere og bedre
end Mandens. Selvom Dyriskhed og Raahed herske i alle andre
Forhold, findes den første Spire til Humanitet dog i Forholdet
mellem Moder og Barn. I Dyrkelsen af Moderen med Barnet
 (Kurotrofos; Madonna) bevaredes en Erindring om denne Spire
til alt Godt i Menneskeverdenen.

Anerkendelsen af Kvindens Ret maatte begynde med, at
man fik Øje for den store Betydning af den Funktion, hun som
Naturvæsen skal udføre for Slægten. Det var et stort Fremskridt,
da hun (ved den anden Arbejdsdeling) blev erkendt som
det svagere Køn, og de tungeste Byrder borttoges fra hendes
Skuldre. Det er endda ikke fuldstændigt sket endnu. I de fattige
Klasser maa hun ofte arbejde saaledes, at det gaar ud over
% --- p. 316 ---
\setcounter{footnote}{0}%
\opage{316}hendes Pligter som Moder, og det er en vigtig Side ved det
store sociale Spørgsmaal at bringe det dertil, at hun kan ofre
sig for sit Hjem og sine Børn\footnote[1]{(Methods of Social Reform. London 1883).}. --- Ved Anerkendelsen af Kvindens
store Betydning i Hjemmet er der gjort et Skridt, som
ikke kan eller skal gøres tilbage. Den Plads, Kvinden udfylder
i Familien, kan ikke erstattes ved nogetsomhelst Fremskridt i
anden Retning. Familielivets store Betydning for Slægtens Udvikling
beror først og fremmest paa Kvindens Stilling i det som
Moder, Hustru eller Søster. Der er her en Kraft, som ikke kan
undværes. Men netop fordi det Hverv, Kvinden her udøver, har
en dyb og fast Grund i Naturen, behøve vi ikke ængsteligt at
vogte paa, at det bliver taget op. Naturen trænger ikke til vor
Beskyttelse; den melder sig nok selv. Og de Forskelligheder,
som virkeligt ere tilstede mellem Mand og Kvinde, ville stedse
gøre sig gældende, naar de ikke skyldes forbigaaende, mere eller
mindre kunstige Forhold. \textsc{Stuart} \textsc{Mill's} Paastand (i hans Skrift
om „Kvindernes Underkuelse“) var, at de Forestillinger, man
havde dannet sig om Kvindens Natur og Evner, kun beroede
paa Vane og Tradition, ikke paa virkelig Erfaring; hvad vi kalde
„kvindelig Natur“, var efter hans Overbevisning kun et Kunstprodukt.
Netop derfor maatte efter hans Overbevisning de Skranker,
man hidtil havde sat for de kvindelige Evners Udvikling
og Brug, ophæves, forat man kunde lære den virkelige kvindelige
Natur at kende. Det gjaldt nu om --- for første Gang i
Verdenshistorien --- at lade Erfaringen afgive sit Vidnesbyrd i
denne Sag! --- Mill ventede, at den kvindelige Begavelse under
friere Forhold vilde vise sig i en overordentlig Glans. Det er
tvivlsomt, om Erfaringen berettiger til en saadan Forventning,
og denne behøves heller ikke for at begrunde den Fordring,
Mill vilde opstille.

Det er, som allerede bemærket, netop Ønsket om, at Kvinden
paa endnu fyldigere og selvstændigere Maade end hidtil
maa kunne udfylde sin Plads som Hustru og Moder, der først
og fremmest fører til Fordringen om en alsidigere Udvikling af

*) Smlgn. Jevons' Afhandling: Married Women in Factories.
% --- p. 317 ---
\setcounter{footnote}{0}%
\opage{317}hendes Evner. Selv hvor hun ikke deltager i det egentligt produktive
Arbejde, vil dog hendes Stilling i Huset som „Konsumtionens
Leder“ udkræve en Uddannelse af forskelligartede Evner.
Og for at hun skal kunne følge Mandens Virksomhed med Forstaaelse
og Sympati og være hans Raadgiver, saavel som for at
hun skal kunne lede sine Børns Opdragelse, maa hun være
saa godt orienteret paa det aandelige og det sociale Omraade
som muligt. Den store Indflydelse, som Kvinden har i Familien,
og som hun derved øver paa hele den etiske, sociale,
religiøse og politiske Udvikling, nødvendiggør en saa rig Uddannelse
af hendes Evner som muligt.

Dertil kommer endnu, at huslig Virksomhed og produktivt
Arbejde ikke altid udelukke hinanden. Der vil ogsaa kunne
blive Brug for kvindelig Arbejdskraft i Produktionen, og det
vil ikke være noget Onde, dersom det ikke gaar ud over Familielivet\footnote[1]{Smlgn. Marcus Rubin: Om Kvindens Adgang til Erhverv. Kbhvn. 1886. (Udg. af Dansk Kvindesamfund). p. 10 ff. --- Mary Gilliland: Women in the Community and in the Family. (International Journal of Ethics. 1895).}.
--- Selv hvor Kvinden slaar ind paa Veje, som hidtil
alene Manden har vandret ad, vil hun --- hvis der virkeligt
er den store Forskel i de to Køns Evner og Natur, som ofte
hævdes --- ikke fornægte sin Ejendommelighed under de nye
Forhold. Hvorledes Ejendommeligheden vil ytre sig, kan først
Erfaring vise.

2. Forhold, hvorover Kvinderne ikke raade, bevirke, at de
ikke altid kunne naa den eneste Bestemmelse, Modstanderne af
Kvindeemancipationen anvise dem. Dels Dødelighedsforholdene,
dels Udvandringen gøre, at der i Europa og i de østlige amerikanske
Stater findes flere Kvinder end Mænd. I Tyskland var
der for nogle Aar siden 800,000 flere Kvinder end Mænd; i
Staten Massachusetts var der i Aaret 1880 over 66,000 flere
Kvinder end Mænd. Paa Landet er Antallet af ugifte Kvinder
forholdsvis noget talrigere end af ugifte Mænd, fordi „Døden
bortriver forholdsvis flere Mænd, saa at der bliver ringere
Chancer for Kvinderne for at blive gifte“. Det store Overskud
af Kvinder i Byerne finder for Størstedelen sin Forklaring i,
% --- p. 318 ---
\setcounter{footnote}{0}%
% CHECK p. 318: notes paired with the marks in order (the layer misread a number)
\opage{318}at flere Kvinder end Mænd indvandre til Byerne fra Landet.
Faktum er, at der ikke blot idethele findes flere Kvinder end
Mænd, men at dette ogsaa specielt er Tilfældet for de Alders
klasser, i hvilke det især er naturligt at indgaa Ægteskab (Tiden
mellem 20 og 40 Aar). Der levede i Danmark i Aaret 190
over 60,000 flere Kvinder end Mænd; i Alderen fra 20 til
40 Aar var der 30,000 flere Kvinder end Mænd\footnote[1]{Rümelin: Bevölkerungslehre. (Schönbergs Handbuch der politischen Oekonomie 1). p. 1207—1209. --- Rubin: Om Kvindens Adgang til Erhverv. p. 21. --- Tidsskriftet „Kvinden og Samfundet“. 1886. p. 126. --- Rubin og Westergaard: Ægteskabsstatistik. p. 67. --- Statistisk Aarbog. 8. Aargang. p. 10—11.}. Det er da
en Selvfølge, at ugifte voxne Kvinders Antal er meget stort.
Man vil maaske hertil sige med en tysk Statistiker, at et
Overskud af kvindelig Befolkning er „et socialt Onde“; men
dermed skaffer man det ikke af Vejen, og de Individer, som
udgøre dette Overskud, har man ikke Ret til at betragte som
overflødige Mennesker. Det maa da dreje sig om at forvandle
det sociale Onde til et socialt Gode ved at tage de tilsyneladende
overflødige Kræfter i Beslag. Overflødige ere de kun,
naar man bliver staaende ved den anden Arbejdsdeling og ikke
vil anerkende den tredie. \textsc{Malthus} var den Første, som krævede
større Respekt for de ugifte Kvinder, og det var hans
Undersøgelse af Befolkningsspørgsmaalet, som førte ham dertil.
Men for ikke længe siden har dog en tysk Filosof hævdet, at
det er egoistisk og etisk forkasteligt, naar en ung Pige, for at
bevare sin Frihed til at vælge sin Ægtefælle efter sin Tilbøjelighed,
uddanner sig saaledes, at hendes Liv kan blive værdifuldt,
selvom hun ikke gifter sig!\footnote[2]{A. Dorner: Das menschliche Handeln. Berlin 1895. p. 422 f.}

3. Man har paastaaet, at der i Kvindens fysiske og sjælelige
Natur skulde mangle Betingelserne for en saadan Uddannelse
og Udvikling, som kunde gøre hende ligestillet med
Manden.

At hun i Reglen er fysisk svagere end Manden, er allerede
omtalt. Men hendes Fortid viser, at hendes Kræfter forslaa
til Mere, end man er vant til at antage. Det civiliserede
% --- p. 319 ---
\setcounter{footnote}{0}%
\opage{319}Liv har i sin Ensidighed ogsaa udviklet hende paa en ensidig
Maade. En sundere og naturligere Opdragelse og Undervisning
vil engang stille Sagen i et andet Lys. Det maa særligt
mærkes, at den Retning, i hvilken alle sunde Reformplaner for
Mænds Opdragelse og Undervisning gaa, netop vil gøre det
lettere at give begge Køn den samme Uddannelse. Større Vægt
paa Anskuelsesevnens Udvikling, paa Tænkeevnens Selvvirksomhed,
paa Muskelsystemets Øvelse er netop, hvad Kvinden
paa sin Vis trænger til ligesaa vel som Manden. Den Forskel
i fysiske Kræfter, som vel stedse bliver tilbage, behøver ikke
at være større end den, der kan udvikles mellem Mænd indbyrdes\footnote[1]{Dog skal det have vist sig, at generel Parese (en med delvis Lammelse forbunden Nervesvækkelse) i de senere Aar er bleven almindeligere end før hos Kvinden i de civiliserede Stater, hvor de deltage i Kulturarbeidet. (Clod-Hansen p. 181). Maaske er dette dog kun et Overgangsfenomen.}.

Naar man mener, at Følelsen er Kvindens psykologiske
Hovedevne, og at hun derfor ikke er anlagt til højere og selvstændigere
intellektuel Udvikling, saa gaar man ud fra en
upsykologisk Modsætning mellem Følelse og Erkendelse. Følelsens
Udvikling hindrer ikke nødvendigvis en Udvikling af Erkendelsen.
Det er kun Følelsen som stærk Sindsbevægelse, som
Affekt, der staar i Modsætning til Erkendelsen. Men en Følelse,
der mere har Inderlighedens end Voldsomhedens Karakter, behøver
ikke at hindre Udviklingen af Tænkeevnen, men kan endog
begunstige den. Levende Sympati fører saaledes naturligt til Fordybelse
i Genstandene, til at optage dem i sig efter deres hele
Ejendommelighed, og en saadan Sympati er nøje beslægtet med
Forskerstemningen. For at vinde Klarhed over sig selv driver
Følelsen idethele til at forske efter de Aarsager, der frembringe
den, og sætter saaledes Erkendelsen i Bevægelse. Selvom hin
Karakteristik af Kvinden som Følelsesmenneske gjaldt for alle
Kvinder uden Undtagelse, vilde den derfor ikke nødvendigvis
udelukke dem fra Retten til at forsøge en højere Udvikling af
deres intellektuelle Evne.

Mange Ejendommeligheder, som man med større eller
% --- p. 320 ---
\setcounter{footnote}{0}%
\opage{320}mindre Ret anfører som karakteristiske for den kvindelige Natur,
skyldes sikkert kun de Forhold, under hvilke Kvinden i lange
Tider har levet, og ville kunne ændres, naar hendes Forhold
ændres. Dette gælder om den nys omtalte psykologiske Ejendommelighed.
Ti den tidligere sædvanlige kvindelige Opdragelse
har ikke været skikket til at udvikle Kvindens Intelligens
og Villie, men har næret det ubestemte Følelsesliv paa de andre
Evners Bekostning. Dermed hænger det sammen, naar Kvinder
ere mere paavirkede ved religiøs Kultus end Mænd. „Vi læse,
at hos Grækerne vare Kvinder lettere at bringe i religiøs Begejstring
end Mænd. Sir Rutherford Alcock fortæller os om
Japaneserne, at det i Templerne er meget sjældent at se en
Forsamling, der bestaar af Andre end Kvinder og Børn; Mændene
ere i hvert Tilfælde faa og høre til de lavere Klasser.
Om Pilegrimmene til Juggernauts Tempel berettes, at idetmindste
fem Sjettedele, og ofte ni Tiendedele, ere Kvinder. Og om
Sikherne høre vi, at Kvinderne tro paa flere Guder, end Mændene
gøre“ *). At det ikke blot er i Asien, men ogsaa i Evropa, at
Kirkerne have deres bedste Støtter og ivrigste Besøgere i Kvinderne,
er en bekendt Sag. Selv „Fornuftens Kultus“ under den
franske Revolution havde sine ivrigste Tilhængere blandt Kvinderne,
og disse vedbleve at deltage i den, efter at Mændene
vare begyndt at afholde sig fra den. Et Øjenvidne beretter:
„Forhen saa man i Kirkerne flere Kvinder end Mænd; det
Samme er Tilfældet i Fornuftens Templer“\footnote[1]{Spencer: The Study of Sociology. Chap. 15. --- Dette Træk i den kvindelige Natur fremhæves meget stærkt af Bachofen som et vigtigt Element til Forklaring af det Kvindeherredømme, han mener har fundet Sted paa et vist Trin af Kulturudviklingen. (Das Mutterrecht. Stuttgart 1861. p. XIII—XVIII).}. I det nuværende
Frankrig har Taine\footnotemark[1] ment at kunne beregne, at medens i Paris
kun 1 Mand af 50 og i Provinserne kun 1 Mand af 12 opfylde
deres Paaskepligt ved at skrifte og gaa til Alters, opfyldes denne
Pligt i Paris af 1 Kvinde af 12 og i Provinserne af 1 Kvinde
% CHECK p. 320: note whose mark was not found
\footnote[2]{A. Schmidt: Pariser Zustände wåhrend der Revolutionszeit. III. Jena 1876. p. 239.}
% CHECK p. 320: note whose mark was not found
\footnote[3]{La Regime Moderne. II. Paris 1894. p. 148.}
% CHECK p. 320: note mark without a note

% --- p. 321 ---
\setcounter{footnote}{0}%
% CHECK p. 321: notes paired with the marks in order (the layer misread a number)
\opage{321}af 4. I København gik i Aaret 1897 (efter en Angivelse i en
Avis) 20,018 Mænd --- men 49,505 Kvinder til Alters. I London
(hvor Daily News et halvt Aar igennem lod Kirkegangen undersøge
ved Hjælp af en Stab af Undersøgere) var der ligeledes
blandt Kirkegængerne (som udgjorde 16 pCt. af Befolkningen)
flere Kvinder end Mænd. En amerikansk Religionspsykolog har
ment at finde en kvalitativ Forskel mellem unge Mænds og
Kvinders Religiøsitet indenfor en og samme Konfession: hos
hine spillede Villie og Tanke, hos disse Frygt og Følelse af
Ufuldkommenhed en overvejende Rolle\footnote[1]{Starbuck: Psychology of religion. p. 225.}. Men selvom dette
hænger sammen med den kvindelige Naturs Ejendommelighed,
behøver denne dog ikke nødvendigvis altid at ytre sig paa
samme Maade; man kan ikke deraf slutte, at Kvinden stedse
maa være autoritetstroende. Eduard v. Hartmann\footnote[2]{Ed. v. Hartmann: Die Phånomene des sittlichen Bewusstseins. p. 521 ff. --- Rousseau har allerede sagt, at Kvinden først har sin Moders, siden sin Mands Religion, og han fandt det i sin Orden. (J. J. Rousseau og hans Filosofi. p. 141).} mener, at
en Kvinde enten maa være autoritetstroende eller have sig en
Mand, der er Fritænker. Det er altsaa to Slags Autoritetsforhold,
hun faar Valget imellem. Men hvis hun vælger det
sidste (eller idethele vælger en Mand med en anden Tro end
den, hun er opdragen i), vil hendes intellektuelle Aktivitet nødvendigvis
være sat i Bevægelse; hun maa da have Evne til at
frigøre sig fra det, som fra først af lagde Beslag paa hendes
Udvikling; og hvorfor skulde hun da ikke kunne naa frem til
en selvstændig Overbevisning i religiøs Henseende, hvad enten
saa denne gaar i den ene eller i den anden Retning? --- Den
samme Forfatter mener, at hun paa Grund af sin levende Sans
for individuelle Tilfælde og individuelle Personligheder ikke formaar
at betragte og behandle en Sag fra en streng og universel
Regels Standpunkt, og at hun derfor nok kunde egne sig til at
være Advokat, men ikke til at være Dommer. Selvom det
skulde være saa, er det dog altid et Fremskridt, naar hun
anses for dygtig til Advokat. Der er mange Mænd, som heller
ikke egne sig til Dommere, og der har været Tider, hvor
% --- p. 322 ---
\setcounter{footnote}{0}%
\opage{322}Kvinden ikke en Gang ansaas for skikket til at være Vidne,
eller hvor hendes Vidnesbyrd kun gjaldt halvt imod Mandens!\footnote[1]{Post: Die Grundlagen des Rechts. p. 451.*) Levned. Dansk Overs. p. 256.}
Kvindens Liv er saa længe hengaaet i Familiens Indre, hvor
den rent personlige Side af alle Spørgsmaal og Forhold i Følge
Sagens Natur staar i første Række, medens de fjernere, mere
upersonlige Hensyn træde tilbage. Hos Mænd, hvis Udvikling
foregaar under lignende Vilkaar, vil man finde den samme
Ejendommelighed. Og den Sans for det Individuelle, hvis
Skyggesider man bruger som Indvending mod Kvinden, har
paa den anden Side store Fortrin, der kunne blive af megen
Betydning i det praktiske Liv. \textsc{Stuart} \textsc{Mill} siger om sin
Hustru: „I intet var hun af større Værd for min aandelige
Udvikling end ved sin korrekte Opfattelse af forskellige Betragtningers
relative Betydning“\footnotemark[1]. En saadan korrekt Opfattelse er
der i Livets forskellige Forhold stadig Brug for, og den er
ikke hyppig. ---

4. Erfaringen har jo allerede aflagt sit Vidnesbyrd, idet
Kvinder faktisk virke paa en hel Mængde Omraader, fra hvilke
de tidligere holdtes ude. Ingen nægter, at de udfylde deres Plads
godt. Om der, som nogle begejstrede Tilhængere af Kvindeemancipationen
mene, hermed vil begynde en ny Æra, der skal
aabenbare uanede og usete Vidundere, vil Tiden vise. Men
Forventningen behøver ikke at spændes saa højt. Naar man
ofte har sagt, at ingen Kvinde hidtil paa noget Omraade har
ydet Noget af allerhøjeste Rang, og deri har set en Indvending
mod Kvindens Adgang til selvstændig Virksomhed, saa er dertil
at svare, at den Maalestok anlægger man dog ikke, naar
Talen er om en ung Mands Ret til at udvikle de Evner, han
føler hos sig. De fleste Mænd vilde være ilde farne, naar en
saa ideal Maalestok skulde anlægges paa dem, og de allerfleste
Mænd maatte være glade, hvis deres aandelige Udvikling,
baade hvad Intelligens og Karakter angaar, havde naat en
Højde som hos Sophie Germain, en George Sand eller en
George Eliot.
% CHECK p. 322: note mark without a note

% --- p. 323 ---
\setcounter{footnote}{0}%
\opage{323}Der er heller ingen Grund til med en tysk Filosof\footnote[1]{Lotze: Grundzuge der praktischen Philosophie. § 49.} at
slaa fast, at hvad der kan gøres uden Kvinderne, ikke bør
gøres af dem. Den Betragtning anlægger man dog heller ikke
paa Mænd. Det vilde være en lang Prøvelse, man havde at
anstille, dersom man kun var berettiget til at slaa ind paa en
saadan Virksomhed, der ikke kunde bestrides af nogen Anden.
Den sunde Tankegang er dog den, at ethvert menneskeligt Individ,
Kvinde som Mand, søger at blive klar over sin Evne og sin
Drift og vælger sin Vej derefter. Erfaringen maa da vise, om
Valget var det rette. Ethvert saadant Valg er et Vovestykke,
men hvo Intet vover, Intet vinder.

Naturforskelligheder, som hænge sammen med varige og
uforanderlige Livsvilkaar, ville ikke kunne udviskes. Men det
er ikke rimeligt, at alle de Forskelle, man mener at kunne
paavise mellem Mands og Kvindes Natur og Begavelse, ere af
denne Art. Hvilke de dybest liggende Forskelligheder ere, vil
først lægge sig for Dagen, naar begge faa Lov at bruge deres
Kræfter. Det vil da maaske vise sig, at Lighederne ere flere
og Forskellighederne andre og finere, end vi nu tænke os.

Et stort Fremskridt er sket i Sammenligning med tidligere
Tider. Man ansaa det da for et betænkeligt Tegn, naar en
Kvinde af Middelstanden kunde læse og skrive\footnotemark[1], og selv \textsc{Goethe}
Zweite Epistel) ønskede de unge Piger opdragne i Køkken,
\textbackslash{}Kælder og Systue, helst uden anden Lekture end Kogebogen:

Wünscht sie dann endlich zu lesen, so wählt sie gewisslich
ein Kochbuch.

Bevægelsen er gaaet frem Skridt for Skridt, og Kvindens Adgang
Til de enkelte Omraader har ikke mødt saa megen Modstand

om det almindelige Princip om hendes „Emancipation“. Hvad
der sker langsomt og efterhaanden, vækker ikke saa bestemt
 tydelig en Bevidsthed og derfor heller ikke saa megen 
% CHECK p. 323: note whose mark was not found
\footnote[2]{'“') „Bei den virginibus“, skrev en gammel Skolelærer i Aaret 1772, ist das Schreiben nur ein vehiculum zur Luderlichkeit“. Seiv Justus}
% CHECK p. 323: note whose mark was not found
\footnote[3]{mente, at som Mand af Folket vilde han ikke ægte en Pige, der kunde læse og skrive. (G. Schmoller: U'eber einige Fragen des Rechts und der Volkswirthschaft. Berlin 1874. p. 120).}
% CHECK p. 323: note mark without a note
Mod%
% --- p. 324 ---
\setcounter{footnote}{0}%
\opage{324}stand som det Nye og Pludselige. Dog er der Lande, hvor ---
især paa Grund af den herskende Militarisme og Bureaukratisme
--- kun meget smaa Skridt have kunnet gøres.

5. Hidtil er der kun talt om Muligheden og Berettigelsen
af Kvindens selvstændige Udvikling og Virksomhed. Men her
kan ikke blot være Tale om Ret, ogsaa om Pligt. Ethvert enkelt
Individ har den Opgave at faa det mest Mulige ud af sin
Natur for derved paa bedste Maade at udfylde sin Plads som
Led af Slægten. Hvad Kvinden fordrer, naar hun vil „emanciperes“,
er egentligt Ret til at komme til at gøre sin fulde Pligt
i Menneskehedens Tjeneste, til at arbejde med paa de fælles
Opgaver. I den nordamerikanske Kvindebevægelse træder denne
Side af Sagen frem paa en smuk og interessant Maade. Den
amerikanske Kvinde fordrede først sin Ret, da det var nødvendigt
for hende, for at hun kunde gøre sin Pligt. Virksomheden
for Kvindens friere Stilling har nemlig historisk udviklet
sig af Virksomheden for Negeremancipationen. Lige fra Begyndelsen
toge de amerikanske Kvinder ivrig Del i denne Virksomhed.
Men deres Ret til offentlig Optræden, selv i en saa
stor og smuk Sag som denne, mødte en forbitret Modstand, der
naaede sit Højdepunkt paa det store Antislaverimøde i London
(1840), hvor man ikke vilde modtage de kvindelige Udsendinge
fra Amerika, fordi saadan Optræden af Kvinder stred „mod
Landets Sæder og den guddommelige Lov“\footnote[1]{History of Women Suffrage. I, p. 55.\textsuperscript{1} Jeannette Schwerin zum Gedåchtniss. Berlin 1899. p. 25.}. Da først begyndte
den Agitation for Kvindens Frigørelse, i hvilken saa mange
fremragende Kvinder toge Del. Et andet Exempel frembyder
en tysk Dame, den i 1899 afdøde Fru Jeannette Schwerin, der
virkede ivrigt for Kvindernes retslige Selvstændighed. Ved sit
praktiske Arbejde blandt de Fattige stod hun ofte hjælpeløs, idet
hun var afskaaren fra at hjælpe Kvinder, der behandledes brutalt
af deres Mænd, som havde Loven til Støtte for deres Magt.
Det var saadanne Erfaringer, der førte hende ind paa Arbejdet
for Kvindernes selvstændige Stilling\footnotemark[1].

6. Naar Kvinden baade har den Mulighed, den Ret og
% CHECK p. 324: note mark without a note
% --- p. 325 ---
\setcounter{footnote}{0}%
\opage{325}den Pligt at arbejde med paa de almenmenneskelige Opgaver
og uddanne sin Personlighed paa selvstændig Maade, vil hun
heller ikke kunne holdes ude fra kommunal og politisk Stemmeret.
De indre og ydre Betingelser, som denne forudsætter, kan hun
ligesaavel være i Besiddelse af som de fleste Mænd, og hun
har ikke mindre end Mændene den største Interesse af, at de
offentlige Sager ledes vel. Den Øvelse og Erfaring, hun mangler,
kan hun i Følge Sagens Natur kun faa gennem praktisk Deltagelse
i det offentlige Liv. Allerede nu øver hun stor Indflydelse
i politisk Henseende, men da hun holdes borte fra
praktisk Erfaring, bliver denne Indflydelse ensidig og bestemmes
ved snevre Synsmaader. Hun har tillige ikke den Ansvarsfølelse,
som Stemmeret og Stemmepligt give. Naar hun faar denne
Ret, ville Mændene blive nødte til at gøre alvorligere Rede for
deres Stemmegivning og ville ikke saa let kunne gaa paa Akkord
med deres Overbevisning\footnote[1]{Stuart Mill: Om den repræsentative Regering. Dansk Overs. p. 241 f. --- Herbert Spencer frygter, at Kvindens Stemmeret paa Grund af hendes Ærefrygt for al Autoritet og hendes Tilbøjelighed til at sætte Filantropi over Retfærdighed vil faa skadelige Følger. Dog hænger dette sammen med vor Overgangstids Vilkaar, mener han; der vil komme den Tid, da kvindelig Stemmeret vil være gavnlig. \{Princ. of Eth. IV. § 108).}. Selvom Mand og Hustru skulde
stemme paa forskellige Sider, saa vilde deri ikke ligge en meget
større Ulykke, end der allerede nu kan være til Stede, naar
Hustruens politiske Overbevisning er en anden end Mandens.
Den større Opfordring, som politiske Rettigheder ville give hende
til at tænke selvstændigt, vil ogsaa frigøre hende fra Præsters
og Skriftefædres Myndighed, saa at det at give Kvinderne
Stemmeret ikke i Længden vil være det Samme som at skænke
Præsterne flere Stemmer, end de allerede have. --- Hvor Kvinden
allerede har faaet denne Ret, synes den at have udøvet en god
Indflydelse paa det offentlige Liv og ikke at have medført de
Misligheder, man ofte har frygtet for.

% --- p. 326 ---
\setcounter{footnote}{0}%
\opage{326}\begin{quote}

\grouphead{3. FORÆLDRE OG BØRN.}

\chapterhead{XX.}{SOCIOLOGISKE DATA.}

\end{quote}

1. Selvom Moderkærligheden rører sig paa de allerlaveste
Trin af menneskelig Existens, selvom vi finde Forældres Kær-;
lighed til deres Børn som et hyppigt Træk hos de Vilde, ---
saa kan det dog siges, at Barnets Behandling og Vilkaar ligesaavel
som Kvindens afgive en Maalestok for den etiske Udvikling.
Barnet er paa lavere Kulturtrin aldeles i Forældrenes
Vold. Familiefaderen har Hals og Haand over det, kan sælge
eller dræbe det og staar Ingen til Regnskab derfor. Hvilken
Behandling Barnet faar, vil meget bero paa, under hvilke ydre
Vilkaar Familien eller Stammen lever. Naar vi finde Børneudsættelse,
især af vanføre Børn, af Trillinger og af --- Pigebørn\footnote[1]{Man har altsaa tidligt næret den Mening, at for mange Kvinder ere et „socialt Onde“.},
næsten hos alle vilde og barbariske, ja endog hos nogle
i andre Henseender civiliserede Folk, saa er det fra først af
den haarde Nød, som har ført til denne Skik. En omstrejfende
Flok Vilde, der stedse er i Fare for Hungersnød eller for Overfald
af Fjender, skiller sig af Selvopholdelsestrang ved Børn,
Syge og Gamle, der hæmme Farten og tære paa Forraadet.
Barnemordene komme hos de Vilde og Barbarerne især paa
Fædrenes og Mændenes Regning; men de tage dog Bestemmelsen
om den Nyfødtes Skæbne paa Familiens eller Flokkens
Vegne. Det er i hvert Tilfælde --- hvad enten det er Moderen
eller Faderen, der dræber --- Hensynet til Stammens Velfærd,
som er det overvejende Motiv til Barnemord, saavelsom til, at
man især dræber Pigebørn. „Der hvor Barnemord er Skik og
Brug, vil Tilværelseskampen jo blive mindre streng, og alle
Stammens Medlemmer ville have næsten lige god Udsigt til at
opføde deres faa overlevende Børn. I de fleste Tilfælde bliver
% --- p. 327 ---
\setcounter{footnote}{0}%
% CHECK p. 327: notes paired with the marks in order (the layer misread a number)
\opage{327}der dræbt et større Antal Pigebørn end Drengebørn; ti det er
indlysende, at de Sidstnævnte ere mere værd for Stammen, da
de, naar de ere blevne voxne, ville hjælpe til at forsvare den
og kunne sørge for sig selv. Men den Vanskelighed, som
Kvinder af Erfaring vide, at der er ved at opføde Børnene,
det at denne Vanskelighed indvirker skadeligt paa Kvindernes
Skønhed, og at de leve bedre, naar der er faa af dem, det angive
Kvinderne selv, og forskellige Iagttagere med dem, som
yderligere Motiver til at dræbe Børnene“\footnote[1]{Darwin: Menneskets Afstamning. II. p. 360.}. Under vore „civiliserede“ Forhold er det oftest den forladte, ensomme Moder,
der begaar Barnemord, og hendes Motiv er Skamfølelse, nok
saa meget som Nød. Moderkærligheden viser sig ofte stærkere
end Nøden, medens den bukker under for Skamfølelse. Man
har ment at finde, at Barnemord ere sjeldne i Lande, hvor
Dommen over Kvinder, der føde udenfor Ægteskab, er meget
mild\footnote[2]{Lecky: History of European Morals. Il. p. 37.}. Dog vise sørgelige Erfaringer, at Nød og Overlæsselse
med Arbejde kan svække og udrydde Moderkærlighed\footnote[3]{Jevons: Married W omen in factories, (Methods of social reform). p. 157 f.; 166.}. --- Hos
mange Kvinder ytrer Moderkærligheden sig først, naar de have
set Barnet og begyndt at vise Omsorg for det. Despine anfører
(i sin Étude psychologique sur les personnes qui commettent
l'infanticide, i 3. Bind af hans „Psychologie naturelle“) en Række
Exempler derpaa. Det bliver derved forstaaeligt, hvorledes den
forladte Moder kan tage Livet af sit Barn strax efter Nedkomsten:
Kontramotivet kommer for sent. --- For Skøger er det en Ære
at faa Børn; de begaa derfor meget sjældent Barnemord. --- I
Oldtidens Stater holdt Forestillingen om Faderens absolute Myndighed
sig, og dermed forbandt sig for Grækernes Vedkommende
den Antagelse, at en Stat stod sig bedst ved en meget begrænset
Tilvæxt i Befolkning. Derved forklares det, at de gamle
Filosofer og Lovgivere hævde Retten eller endog Pligten til at
berøve det undfangne eller allerede fødte Barn Livet. De indførte
derved ikke noget Nyt, men stadfæstede en meget gammel
% --- p. 328 ---
\setcounter{footnote}{0}%
\opage{328}Skik, som først Kristendommen med sin Omsorg for hvert enkelt
Menneskes Sjæl førte til at omstyrte. Først nu blev det
ufødte og det nyfødte Barns Ret hævdet\footnote[1]{vidt, at naar ved Barnsnød Moderen eller Barnet maatte ofres, blev Moderen ofret, for at det udøbte Barn ikke skulde gaa til Helvede. Lecky ib. p. 25.}.

2. Ligesom for Kvindens Vedkommende var det ogsaa for
Barnets tildels økonomiske Grunde, der førte til bedre Vilkaar
for det. Førend Tanken om det enkelte Individs Værdi og Betydning
kunde gøre sig gældende, maatte den Hjælp og Nytte,
Børnene kunde yde, bevirke, at de bleve skaansomt behandlede.
Børnene forrettede Arbejde ligesom Trællene. Men dertil
kom endnu, idetmindste for Sønnernes Vedkommende, ogsaa
andre Hensyn. I sin Søn saa Faderen den, der skulde hævne
ham, hvis han kom voldeligt af Dage, og den, der skulde forrette
Ofrene ved Familiealtret efter hans Død. Kun naar han
efterlod sig en Søn, kunde han betale sin Gæld til sine Forfædre
ved at sikre Familiekultusens Vedvaren. Men Forholdet
vedblev dog til den nyeste Tid at være et Myndigheds- og
Underdanighedsforhold, idet Forældrenes Autoritet anerkendtes
paa en Maade, som endnu kunde minde om den primitive Familie.
Efterat Staten har sat snevrere Grænser for Forældremyndigheden,
og efterat i det øvrige Samfund Friheds- og
Lighedsideer ere blevne herskende, er Forholdet mellem Forældre
og Børn langt mere blevet et rent Sympatiforhold end
et strengt Autoritets- og Lydighedsforhold. Familien har derved
selv udviklet sig til i langt højere Grad end paa det primitive
Trin at være Forbilledet for et inderligt og fuldkomment Samfund,
hvor intet Medlem blot er Middel, men Alle ere selvstændigt
medvirkende Led.

') Smlgn. for Nordens Vedkommende: R. Keyser: Efterladte Skrifter. II, p. 315. --- Den strenge kristelige Betragtningsmaade gik endog saa

% --- p. 329 ---
\setcounter{footnote}{0}%

\opage{329}\chapterhead{XXI.}{ETIK OG PÆDAGOGIK.}

1. Allerede Fysiologien giver os et tydeligt Vink om Forældrenes
etiske Stilling overfor Børnene\footnote[1]{Smlgn. med Hensyn til det Følgende: Psykologi VI C 2—3.}. Forældrene „sætte“
ganske vist Barnet ind i Verden, men de ere dog selv kun
Mellemled i den Naturorden, der længe før deres Bevidsthed
og deres Drifter udviklede sig, havde grundlagt den første Spire
til det nye Individ. Paa deres fysiske Forhold til Barnet kan
der derfor ikke grundes nogen Ejendoms- og Herskerret. Deres
Myndighed overfor Barnet grunder sig derimod paa det Hverv,
der paahviler dem som de Nærmeste: at pleje Spiren og lede
Væxten. Deres Autoritet grunder sig paa Opdragelsens Nødvendighed.
Forældreautoriteten kan, som al anden Autoritet,
uvilkaarligt eller vilkaarligt gøre sig selv til Formaal istedetfor
Middel. Det sker, naar den Magt, der er lagt i Forældrenes
Haand, nydes som Magt, saa at Børnene blive Midler for Magttrangens
Tilfredsstillelse. Det sker ogsaa, naar Forældrene give
deres ved Børnenes Adfærd vakte Følelser frit Udtryk, uden at
betænke, at baade Ros og Dadel ere vanskelige Midler at anvende,
og at Hovedsagen ikke er, at Forældrene faa Luft for
deres Følelse, men at de Domme, de udtale, kunne faa den
rette motiverende Indflydelse paa Udviklingen af Børnenes Karakter.
Retten til at udøve Magt og udtale Domme beror paa
Pligten til at være opdragende Forsyn. Denne Pligt krænkes
paa endnu mere uvilkaarlig Maade ved enhver Mangel paa Selvbeherskelse,
der enten direkte, ved at kue og plage Børnene,
eller indirekte, ved det Forbillede, som gives dem, kan stride
mod Opdragelsens Formaal.

Hin Pligt bæres under sunde Forhold frem af en af de
stærkeste naturlige Følelser, som næres og styrkes, jo stadigere
Lejlighed der er til dens Tilfredsstillelse. Den instinktmæssige
Virksomhed gaar forud for den egentlige, inderlige Følelse, og
% --- p. 330 ---
\setcounter{footnote}{0}%
\opage{330}denne tiltager med den fortsatte Virksomhed. Samlivet frembringer
Samfølelsen. (Smlgn. det aristoteliske Princip). Den almindelige
Erfaring, at Velgørere elske dem, de gøre vel imod,
højere end de selv elskes af disse, bekræfter sig her. Den
Følelse hos Børnene, som svarer til Forældrenes Kærlighed,
kan vanskeligt naa den samme Grad af Styrke. Homer bruger
ofte det Udtryk om unge Helte, som falde, at de ikke fik betalt
deres Forældre „Fostringsløn“, Gengæld for, hvad de havde
modtaget af dem. Men saa stærk Børns Taknemlighed og Pietet
end maatte være, ville de dog aldrig kunne naa at betale den
fulde Fostringsløn. Det har sin Grund i, at Barnets Blik naturligt
er rettet fremad. Selvom det suger nok saa megen Næring
af Fortiden, er det dets Drift og dets Kald at se frem og ikke
tilbage, og naar Forældrenes Sympati er uinteresseret, ville de
ogsaa ønske det saaledes. .De ville vide, at dem bør det at
dale, Børnene at stige. Dog vil, hvor det rette Forhold er tilstede,
Forældrenes Sympati med Børnenes Fremskriden mødes
med Børnenes Pietet mod Forældrene, og derved den Harmoni
mellem den ældre og yngre Slægt opstaa, der er en af de betydningsfuldeste
Sider ved Familien. (Smlgn. XIV, 2 og XII, 6).

Foruden det Baand, der forener Forældre og Børn indbyrdes,
har ogsaa det, der forener Søskende indbyrdes, en stor
etisk Betydning. Fællesskab i Udspring, i Livsvilkaar, i Erindringer,
i Opdragelse og Virksomhed danne --- foruden det naturlige
Slægtskab i Temperament og Anlæg --- et fast Grundlag
for Fællesskab og Forstaaelse, som kan holde sig, selvom Livets
senere Erfaring og Virksomhed bliver højst forskellig. Man har
ikke kunnet finde noget bedre Udtryk for det Samfund, der
bæres af almindelig Menneskekærlighed, end almindeligt Broderskab.
---

2. Selvom det nu kan siges at staa fast, at Barnet ikke
blot er Forældrenes Ejendom og Middel for deres Formaal, saa
kunde det dog synes, som om Barnet i en anden Henseende
maatte behandles som Middel og ikke som Formaal. Barndommen
er jo Livets Begyndelse, en Overgang og Forberedelse
til det fulde og egentlige Menneskeliv. Barnet maa da --- kunde
det synes --- betragtes som Middel, da det gælder om at ud%
% --- p. 331 ---
\setcounter{footnote}{0}%
\opage{331}danne det til et saa kraftigt og fuldkomment Menneske som
muligt: altsaa det maa være Middel for sit eget fremtidige Menneske!

Men ligesaa lidt som hele Livet maa betragtes blot som
Middel, ligesaa lidt maa nogen enkelt Del af Livet betragtes
blot som Middel for en anden Del, naar der ikke er en tvingende
Nødvendighed tilstede. Barndommen staar først og fremmest
som en selvstændig og ejendommelig Livsperiode og Livsform,
der har sine egne Vilkaar og skal bedømmes efter sine egne Love.
Den er ikke mere en blot Overgangstid end enhver Del af Livet
er det. Barnet maa først og fremmest betragtes som Barn, ikke
blot som vordende Voxen. Af Larven udvikler sig vel i Tidens
Fylde et Insekt; men derfor gaar det ikke an at behandle Larven
som et Insekt. Larven har sine ejendommelige Evner og
Fornødenheder. Bortset fra asketiske Retninger, der betragte al
sund Livsglæde og naturlig Virksomhed med Mistanke, har Barndommen
ikke værre Fjender end en udvortes Nyttelære, der blot
spørger om, hvad Barnet vilde have Nytte af, og ikke hvad det
har Nytte af.

Først og fremmest gælder det, at Barnet faar Lov til og
Hjælp til (ti det kan behøve Hjælp) at være et rigtigt Barn,
at gennemleve Barndomstiden med fuldstændig Hengivelse, uden
at forstyrres af Bekymring, Savn og Møje. Derved samler det
legemlig og aandelig Kraft til den fremtidige Udvikling. --- Og
dernæst gælder det, at det Arbejde, som fordres af Barnet, saavidt
muligt bliver dets eget Arbejde, saa det kan føle Glæden
ved at bruge sine egne Kræfter. --- Naar disse Betingelser fyldestgøres,
løses derved ogsaa det Problem: at behandle Barndommen
som en selvstændig og ejendommelig Livsperiode, og dog at gøre
den til en Forberedelse til det følgende Liv.

Det er \textsc{Jean} \textsc{Jacques} Rousseaus store Fortjeneste at
have fremsat disse Sandheder saaledes, at de ikke ville kunne
glemmes igen\footnotemark[1]. Derved blev han paa en Gang Barndommens
Evangelist og Pædagogikens Reformator.

3. For Barnet er ligesaa lidt som for den Voxne Lydighed

Se min Bog J.J. Rousseau og hans Filosofi.\textsuperscript{2} p. 89—91; 133—148.
% CHECK p. 331: note mark without a note

% --- p. 332 ---
\setcounter{footnote}{0}%
\opage{332}den højeste Dyd. Lydighed har ingen Værdi i sig selv, men er
kun et Middel. Det er nødvendigt, at Barnet lærer Lydighed,
da det ellers ikke kan have sin fulde Frihed. Man kan kun give
Barnet frit Spillerum, hvis der er et elastisk Baand, der holder
tilbage, naar det gaar for vidt. Naar Barnet nærmer sig en Afgrund
uden at ane Faren, kan kun Lydighed mod den advarende
Stemme frelse det; der bliver ikke Tid til at give Grunde og
nærmere Forklaringer, men Tilliden til den Røst, som høres,
maa være tilstrækkeligt Motiv. Forældrene repræsentere overfor
Barnet den modne Erfaring og den selvstændige Dom. Indtil
Barnet opnaar Selvstændighed, ledes det af Trangen til at finde
Billigelse af sin Færd hos dem, det elsker højest; af denne Trang
udvikler sig saa efterhaanden Trangen til Selvagtelse og dermed
den etiske Selvstændighed\footnote[1]{Dette er meget smukt paavist af J. G. Fichte: Reden an die deutsche Nation. Berlin 1808. p. 317—329. --- Preyer (Die Seele des Kindes. 3. Au fl. p. 227) bemærker om sit Barns Iver for at efterligne, hvad de Voxne gjorde: „Sein Nachahmungstrieb erscheint hier fast wie Ehrgeiz.“}. Barnet har Samvittigheden udenfor
sig, før det kan faa den i sig selv. Men Lydigheden og Tilliden
til Andre føre kun til selvstændig Overbevisning derved,
at de føre til Handling. Vi mindes atter her det aristoteliske
Princip, at man kun bliver god ved at handle godt. Den uvilkaarlige,
instinktmæssige Maade, hvorpaa Barnet følger Ældres
Færd og efterligner dem, den Iver, hvormed det slaar ind paa
de Veje, man aabner for dets stærke Trang til at bruge sine
friske Kræfter, har nok saa Meget at sige som den bevidste Lydighed.
Der skabes, uden at Barnet selv mærker Noget til det,
Præcedenser, som bestemme Fremtiden og afgive et bedre Grundlag
end Formaninger og Trusler, Belønninger og Straffe. Erfaringer,
som det senere faar Brug for, kan det paa denne Maade
gøre i det Smaa saaledes, at det føler Glæde ved sin egen Virksomhed.
Naar den Tid kommer, da det forstaar Buddene og Formaningerne,
har det allerede udøvet dem; og det er egentligt
den eneste Maade, hvorpaa det ret kan forstaa dem. Ligesom
man efter nyere Metoder i Sprogundervisningen ikke begynder
med grammatiske Regler, men med praktisk Øvelse i Sproget,
% --- p. 333 ---
\setcounter{footnote}{0}%
\opage{333}saaledes bør man i den praktiske Etik heller ikke begynde med
moralske Regler, men med Øvelse af Kræfterne.

Den ubevidste Opdragelse er nok saa vigtig som den bevidste.
Der er Tanker og Følelser, som først blive mulige, naar
Virksomhed gaar i Forvejen. --- Men ved ubevidst Opdragelse
kan man ogsaa forstaa den Opdragelse, som Forældrene øve,
uden at de selv vide af det. Deres ubevogtede Øjeblikke virke
nok saa meget paa Barnet som de Øjeblikke, i hvilke de optræde
overfor det med klar pædagogisk Bevidsthed; i sine ubevogtede
Øjeblikke taler og handler man jo ofte mest energisk.
Det nytter derfor ikke at lægge det fortræffeligste pædagogiske
System til Grund, naar talrige ubevogtede Øjeblikke rive ned,
hvad Systemet har bygget op. Overhovedet bestaar en stor
Del af Barnets Opdragelse i uvilkaarlig Efterligning, der fører
til Indøvelse i Virksomheder og Arbejder, det senere med Bevidsthed
skal udøve. Gennem aktiv Efterligning lærer Barnet
Livets forskellige Sider at kende. I sin Leg gentager det de
Voxnes Færd, optræder i forskellige Værdigheder og Situationer, i
hvilke det ser de Voxne fremtræde, er snart Herre, snart Tjener,
snart givende, snart modtagende. Det lærer Livets Modsætninger
at kende, og tillige erfarer det, at selv de Voxne, der
ellers synes at være frigjorte for alle Skranker, have Regler og
Hensyn, efter hvilke de rette sig. Skridt for Skridt føres det
ind i Livets store Sammenhæng, og den moralske Udvikling
kan her gaa Haand i Haand med den intellektuelle. Evnen til
Medfølelse, til i levende Sympati at glædes og lide med Andre
næres gennem dette af uvilkaarlig Efterligningstrang fremkaldte
Erfaringskursus, som først i den nyeste Tid er bleven nøjagtigere
opfattet og beskrevet\footnote[1]{Josiah Royce: On Certain Psychological Aspects of Moral Training. (International Journal of Ethics. 1893). --- Mark Baldwin: Mental Development in the Child and the Race. Methods and Processes. New York 1895. --- Samme: Social and ethical interpretations in mental development. A study in social psychology. New York 1897. (En af det danske Videnskabernes Selskab prisbelønnet Afhandling.)}, og som er et vigtigt Exempel paa de
aristoteliske Princip (XIII, 4).

4. Paa den intellektuelle Opdragelses Omraade synder
% --- p. 334 ---
\setcounter{footnote}{0}%
\opage{334}man mod Sætningen om Barndomslivets selvstændige Betydning,
naar man vil udvikle og øve Barnets Evner ved Midler, der
ikke have nogen Interesse eller Værdi i sig selv. Saaledes
naar man begrunder Nytten af grammatiske Ramser med, at
de styrke Hukommelsen. Hvad der bruges til Udvikling af de
formelle Evner, bør være Noget, som ogsaa ved sit Indhold
kan have Interesse, saaledes at der kan føles Tilfredsstillelse
ved at have tilegnet sig det.

Og ved den intellektuelle Opdragelse er det af den største
Betydning, at Selvvirksomheden vækkes. Saa forskellige end
Barnealderen og den voxne Alder kunne være, saa stemme de
dog overens i, at det Højeste, som kan naas, er at virke med
sin egen Kraft i en betydningsfuld Sags Tjeneste. Det er muligt
for Barnet paa dets Vis, saavel som for den Voxne paa hans
Vis. Det Stof, man fordrer, at Barnet skal tilegne sig, maa
derfor lægges saaledes til Rette, at Barnet kan blive selvvirksomt,
hvad enten det er Fantasien eller Forstanden, det særligt
kommer til at arbejde med. Ved Selvvirksomhed kommer
Barnet til at føle Glæde ved Arbejdet, uden at ane, at det derved
tillige lægger den bedste Grund til sin egen Fremtid. Det
Arbejde at lære ligger, som \textsc{Duhring} træffende har bemærket\footnote[1]{Der Werth des Lebens. 2. Ausg. p. 95.},
imellem Legen, der overvinder selvskabte Hindringer, og det
egentlige Arbejde, der overvinder virkelige, objektive Hindringer;
de Hindringer, Lærearbejdet skal overvinde, ere de, der ligge
i Trægheden og Uøvetheden.

Kun et idealt pædagogisk System vilde fuldstændigt kunne
forene de to Hensyn: at give Barndommen dens selvstændige
Betydning og dog lade den være en Forberedelsestid. Men alle
Reformtanker i Pædagogiken pege i den Retning. At der vil
være store Vanskeligheder for disses fulde Gennemførelse, følger
af sig selv. Det ligger allerede i, at det er Voxne, som opdrage;
Insektet kan aldrig ganske forstaa Larven. Foreløbigt er
det et stort Gode, at Opmærksomheden er saa stærkt vakt for
Opdragelsens Betydning. Er det for de Voxne saa overordentligt
vanskeligt at faa Evner og Opgaver til at passe sammen,
% --- p. 335 ---
\setcounter{footnote}{0}%
\opage{335}saa maa det dog kunne lykkes at afpasse dem efter hverandre
for Barnet, da man selv er Herre over, hvilke Opgaver man
vil stille det; og lykkes det engang at løse Opdragelsens
Problem, ville derved ogsaa de Voxne hjælpes i Kampen med
deres Problemer.

5. Barnets Selvvirksomhed har naturligvis sine Grænser.
Det kan ikke selv paany gøre alle de Erfaringer, Slægten paa
sin lange Vej har gjort. Erfaringen opdrager os for en stor Del
gennem Skuffelser, gennem Forventninger, der ikke opfyldes.
Derved lære vi at fastsætte Grænserne mellem det Virkelige og
det blot Mulige\footnote[1]{Die Seele des Kindes. 3. Aufl. p. 216.}. Men det er Opdragerens Pligt at sørge for,
at Skuffelserne ikke blive for store for Barnet, saa de kue dets
Mod. En for brat Overgang fra Forventning eller Illusion til
Realitet kan let kvæle den Selvtillid og Tryghedsfølelse, som
Barnet allermindst kan undvære. En lille Iagttagelse, \textsc{Preyer}
har gjort, oplyser dette. „I den 62. Uge kunde Barnet endnu
ikke staa længere end et Øjeblik, uden at det understøttedes eller
berørtes. Denne Mangel paa Evne beror ikke mere paa Vanskeligheden
ved at bevare Ligevægten, men paa Mangel af Selvtillid;
ti kun da kan det ikke staa alene, naar det véd, at man
ikke holder ved det. Men naar det ikke véd, at Faderens støttende,
stedse mindre trykkende Haand er fjernet fra dets Ryg,
saa kan det staa i flere Sekunder uden Understøttelse“\footnotemark[1]. Hvad
der gælder den legemlige Gang, gælder ogsaa den aandelige Udvikling.
Barnets Forestillinger og Følelser udvikle sig i tryg Tillid
til den foreløbige Ramme, det har dannet sig for sin Verdensopfattelse,
dels uvilkaarligt, dels ved Optagelse fra sine Omgivelsers
Ytringer; og sprænges denne Ramme for tidligt, kan det
hæmme Udviklingen og bringe en Usikkerhed ind, som kan skade
Livets Sundhed. Barnet begynder med naivt at tillægge alle sine
Forestillinger Gyldighed. Det begynder med en Sangvinitet, der
foregriber Erfaringen. En brutal Sønderrivelse af Illusionerne,
en altfor pludselig Indskærpelse af Forskellen mellem Fantasiens
og Virkelighedens Verden kan Barnet ikke taale. Denne For-

') Smlgn. Psykologi V B, 4.
% CHECK p. 335: note mark without a note

% --- p. 336 ---
\setcounter{footnote}{0}%
\opage{336}skeis Indøvelse strækker sig jo ogsaa gennem hele Livet og
kan ikke tilendebringes paa én Gang.

Opdragelsens Kunst bestaar i saavidt muligt at lade Barnets
Sangvinitet og Fantasi have sit frie Spil, og dog at øve en Kritik,
der beskærer det Fordærvelige og betoner den sunde Kærne.

Intetsteds finder dette en vigtigere Anvendelse end paa de
religiøse Forestillinger. Barnets Maade at opfatte og forklare sig
Tingene paa har Meget fælles med Naturmenneskets. De have
begge en Tendens til at finde Forklaringen af, hvad der foregaar
omkring dem, i personlige Væseners Indgriben. De personificere
Naturen og forstaa egentligt kun personlige Aarsager. Det er
ikke altid sagt, at der kommer en poetisk Naturopfattelse ud deraf.
Netop Tendensen til at personificere kan give Forklaringen en
udvortes og mekanisk Karakter. Naar Barnet forklarer Stjernernes
Tilsynekomst om Aftenen ved, at Vorherre tænder dem, ligesom
„Manden“ tænder Gaslygterne, eller naar det mener, at der om
Foraaret kommer en Mand og sætter Bladene paa Træerne, da
er der heri ikke saa megen Poesi som i den naturvidenskabelige
Forklaring. Men Barnet lægger sig nu én Gang Tingene
tilrette paa sin Maade og appellerer til saadanne Aarsager, som
ere det bedst bekendte. Derfor optager det ogsaa med Lethed
de bibelske Fortællinger, som passe saa godt til dets hele Maade
at forestille sig Verden paa. Denne hele mytologiske Periode
maa Barnet have Ro til at gennemgaa, uden hvert Øjeblik at
forstyrres af en Kritik, der gaar ud fra helt andre Forudsætninger.
Baade dogmatiske og antidogmatiske Opdragere forse sig
ofte herimod. Ligesom Missionærer ofte ubarmhjertigt erklære
Hedningernes Tro for Djævelens Værk, saaledes sønderrives ofte
Barnets vildt voxende mytologiske Forestillinger af ængstelige
Forældre, fordi de ikke stemme med Katekismen. Og Fritænkere
mene sig ofte af Sandhedskærlighed nødte til at stanse alle mytologiske
Forestillinger hos Barnet i Væxten. I begge Tilfælde
hæmmes Barnets frie naturlige Udvikling. Man maa give det
Lov til at gennemløbe en lignende mytologisk Periode som den,
Slægten har gennemgaaet, ligesom det i sin tidlige Fosterudvikling
har frembudt Former, der minde om lavere Dyrearter. Hvad
man kan gøre, er kun at forkorte og lette Vejen; men det maa
% --- p. 337 ---
\setcounter{footnote}{0}%
\opage{337}selv arbejde sig ud af denne Periode, hvis Udviklingen skal
være sund. Man kan hjælpe det til at drage alle Konsekvenserne
af dets Forestillinger og derved prøve disse i det Enkelte.
Derved udvikler man dets Kritik og dets Evne til at slutte.
Man kan give det Lejlighed til at gøre saa mange frugtbare
Erfaringer som muligt. Man kan bibringe det en Forestilling
om, hvor vanskeligt det er at naa en nøjagtig Viden, og hvor
Meget det endnu har at lære for ret at kunne behandle de
interessante Spørgsmaal, det opkaster. Og man kan vise det, at
det Vigtigste er at have et godt Hjerte og en ren Samvittighed,
og at de kunne findes, selv hvor der er meget forskellige religiøse
Forestillinger.

For tidlig direkte Kritik vil let give Forstandslivet en Overvægt
over de andre Sider af Sjælelivet. Erkendelsen har idethele
en Tendens til at udvikle sig hurtigere end Følelsen; og
denne Tendens medfører let en sygelig Ensidighed, naar man
forstærker den ved at paatvinge Barnet, hvad man anser for
Sandhed, istedetfor at lade Barnet selv finde det ad sin egen
Vej. Der vil da opstaa en sjælelig Disharmoni, en Delthed og
Lammethed i Sindet. Og der vil da senere let komme en Reaktionsperiode
af en usund Beskaffenhed. Ved en Art Kontrastvirkning
higer da Individet senere efter at opleve en Kreds
af Følelser, som det ikke lærte at kende i den naturlige Tid.

Kunne Selvtilliden og Selvvirksomheden kues ved for tidlig
Kritik, saa kues de naturligvis endnu mere, naar man systematisk
indgyder Barnet den Forestilling, at det ikke tør stole paa
sine egne Sanser og sin egen Forstand, men skal lære at tage
dem til Fange, da det, som bør tros, ikke kan forstaas. Men en
saadan Adfærd ligger for langt fjernet fra det hele Standpunkt,
hvorfra vi her gaa ud, til at den behøver nogen nærmere Kritik.
En Opdragelse, der ledes i etisk Aand, kan ikke lægge
absolut Vægt paa nogensomhelst Tro eller Ikke-Tro. Dens Maal
er Karakterens Dygtighed og Selvstændighed, Følelsernes Inderlighed
og Sundhed, Forstandens Kraft og Klarhed. Alle konfessionelle
Modsætninger ere af forsvindende Betydning i Forhold
til dette Maal. Det er et Menneske, der skal komme ud
af Opdragelsen, ikke en Troende eller en Ikke-Troende.

% --- p. 338 ---
\setcounter{footnote}{0}%

\opage{338}\chapterhead{XXII.}{STATEN OG BØRNENE.}

1. Familien er paa det moderne Kulturtrin ingen aldeles
afsluttet Verden, saaledes som den patriarkalske Familie var, der
kun gennem sit Overhoved stod i Forbindelse med det større
Samfund. De enkelte Individers, ogsaa Kvindens og Barnets,
Ret til selvstændigt Liv er blevet anerkendt, og det er Statens
Sag at hævde denne Anerkendelse. De Krav, Staten stiller overfor
Familien, bero dels umiddelbart paa dens Kald til at beskytte
de Svage, dels derpaa, at Familiens Medlemmer ogsaa ere
Statens Medlemmer og maa fyldestgøre de Betingelser, som derved
udkræves.

Saaledes hævder Staten det ufødte og nyfødte Barns Ret
til at existere og tillader Ingen at erklære noget andet Individ
for overflødigt eller for uberettiget eller uskikket til at være til.
Selv hvor Udsigterne for det Individ, der begynder sin Tilværelse,
synes at være nok saa mørke, kan en saadan Myndighed til at
overskære den én Gang anlagte Livstraad ikke tillægges Nogen.

Staten tager sig fremdeles af de forladte Børn og holder
Forældrene til at gøre deres Pligt, hvor deres egen Samvittighed
og Forældrefølelse ikke driver dem dertil. Den beskytter
Børnene mod Forældrene. Det viser sig nemlig, at den moderlige
og faderlige Følelse kan sløves og udryddes af Nød, Skam
og andre Aarsager, og hvor Familieforholdet saaledes ikke yder
tilstrækkelige Motiver til at fyldestgøre, hvad Velfærdsprincipet
fordrer, træder det større Samfund til, skønt det ikke kan løse
Opgaven saa godt, som det snevrere under sunde Forhold kan.
Det er i Kraft af dette Princip, at Staten fører Tilsyn med og
indskrænker Anvendelsen af Kvinder og Børn i Fabriker, skønt
den næppe endnu har gjort nok for „at sikre Barnets Ret til
dets Moders Bryst“. I Kraft af dette Princip burde Staten ogsaa
tvinge Forældrene til at vedkende sig deres Børn, hvad den
imidlertid for de uægte Børns Vedkommende fordetmeste ikke
gør. Overfor uægte Børn følger Staten en barbarisk Regel: mater
semper certa, Moderen er sikker, hende holde vi os til. Byrden
% --- p. 339 ---
\setcounter{footnote}{0}%
\opage{339}lægges overvejende over paa Moderen. Barnets Ret til Faderens
Navn anerkendes ikke; og selv hvor Faderen paavises, bestemmes
hans Bidrag til Barnets Underhold efter Moderens, ikke
efter hans egen sociale Stilling. Endnu i det tyske Riges nye
borgerlige Lovbog hedder det (\$ 1567): „Et uægte Barn og
dets Fader gælde ikke som beslægtede“; --- i det oprindelige
Udkast stod endog: „Mellem et uægte Barn og dets Fader bestaar
intet Slægtskab“. Et slaaende Exempel paa Forskellen
mellem juridisk og moralsk Tankegang! --- Naar en saadan Ordning
forsvares ved, at ellers vilde løse Forbindelser begunstiges
og Ægteskabets Anseelse svækkes '), overser man, at der her
foreligger en umiddelbar og naturlig Adkomst for Barnet, som
ikke kan skydes til Side for fjernere liggende og ubestemte Hensyn.
Man lader her faktisk Barnet lide for Forældrenes Synder,
--- atter et barbarisk Princip, som den mere udviklede etiske
Følelse ellers har forkastet. Barnets Ret til Faderens Navn vil
desuden gøre, at der føles større Betænkeligheder ved at sætte
Børn i Verden uden for Ægteskab\footnote[1]{Goos: Almindelig Retslære. I, p. 544—55]}. En lettere Adgang til at
indgaa lovformeligt Ægteskab vil indskrænke Antallet af de Forbindelser,
af hvilke uægte Børn udspringe; Indførelsen af borgerligt
Ægteskab skal saaledes i flere Lande have medført Aftagelse
af uægte Barnefødsler. Det maa derhos ikke glemmes, at saadanne
Forbindelser langtfra er den værste, mest nedværdigende
Form, under hvilken Menneskenes Mangel paa Evne til at virkeliggøre
Kønsforbindelsens Ideal, det frie Monogami, viser sig.

2. Ogsaa Barnets aandelige Udvikling maa Staten værne
om, baade fordi den skal beskytte de Svage, og fordi Barnet
ellers ikke kan komme til at fyldestgøre, hvad Staten fordrer
af sine Borgere. Hvor Forældrene ikke kunne eller ikke ville
sørge for deres Børns Undervisning, griber Staten derfor ind
med Magt. Den maa ogsaa her ofte beskytte Børnene mod 
% CHECK p. 339: note whose mark was not found
\footnote[2]{Dog synes de statistiske Erfaringer her at modsige hverandre. Smlgn. Oettingen: Moralstatistik. 3. Aufl. p. 323 med Rümelin: Reden und Aufsätze. Neue Folge. p. 620. Rümelin hævder, at den franske Regel: la recherche de Ia paternité est interdite, styrker Kvindens Modstandskraft overfor Forføreren.}For%
% --- p. 340 ---
\setcounter{footnote}{0}%
\opage{340}ældrene, nåar disse for at drage Fordel af Børnenes Arbejde
holde dem borte fra Skolen I>. Hvor Forældrene af egen Drift
sørge for Børnenes Undervisning, har Staten kun at forvisse
sig om, at det Minimum af Kundskab, den fordrer, er tilstede.

Naar man har ment, at Staten herved begaar Indgreb i
Forældrenes Ret, gaar man ud fra, at Børnene ikke have nogen
selvstændig Ret, som Forældrene kunne krænke dels af Ligegyldighed,
dels af Uvidenhed, dels af Egennytte. Erfaringen viser,
at en Kontrol med Forældrene i denne Henseende ikke kan
undværes, og den ligger utvivlsomt (hvad dog først Læren om
Staten nærmere kan begrunde) indenfor Kredsen af Statens Opgaver.
Staten maa have Myndighed til at fordre, at ethvert normalt
Barn skal have de Kundskaber, som ere nødvendige for at
være Borger i den.

At det er Barnets Ret, det her gælder om, er ogsaa udtalt
i den danske Grundlov, hvis § 85 lyder: „De Børn, hvis Forældre
ikke have Evne til at sørge for deres Oplærelse, ville erholde
fri Undervisning i Almueskolen“. I det første Udkast til
denne Paragraf var der derimod Tale om Forældrenes Ret til
at faa fri Undervisning for deres Børn, naar de ikke selv kunde
bestride denne. Omarbejdelsen motiveredes paa den grundlovgivende
Rigsdag udtrykkeligt ved, at „det dog ikke er Forældrene,
men Børnene selv, hvem den her omhandlede Ret tilkommer“.
--- I den tilsvarende Bestemmelse i tyske Forfatninger
(stammende fra den tyske Rigsforfatning af 1849) er der
ogsaa Tale om en Ret for Børnene. Forældremagten indskrænkes,
for at en almindelig menneskelig Grundrettighed kan beskyttes\footnote[1]{u. Aufsatze. Neue Folge). p. 475.}.

Man behøver ikke at frygte for, at man derved skal formindske
Forældrenes Ansvar. \textsc{Spencer} har, ledet af denne Betænkelighed,
gjort Indsigelse mod al Statsundervisning. Han for-

i) De engelske Arbejdere have indset dette Punkts Betydning baade
for Børnene og for hele Arbejderstanden. De engelske Fagforeninger
have ivrigt virket for Indførelsen af Skoletvang. (L. Br en tan o: Die Arbeitergilden
der Gegenwart. Il, p. 100).

 Smlgn. Rümelin: Ueber das Object des Schulz w anges. (Reden
% --- p. 341 ---
\setcounter{footnote}{0}%
\opage{341}drer Familie og Stat aldeles holdte ude fra hinanden og er overbevist
om, at begge kun kunne trives ved at staa hinanden saa
fjernt som muligt\footnote[1]{Principles of Sociology. I. p. 738 ff.}. Hertil er at sige, at man jo dog ikke tager
det etiske Ansvar fra et Menneske, fordi man sørger for, at
Noget bliver gjort, som han vel er den Nærmeste til at gøre,
men som han alligevel ikke kan eller vil gøre. Staten udøver
her en hjælpende, supplerende Virksomhed. Familie og Stat ere
ganske vist meget forskellige i deres Væsen, men de have for
en Del fælles Formaal, da Familiens Medlemmer ogsaa ere Statens.
Det større Samfund træder til, hvor det mindre af en
eller anden Grund ikke udfører sit Hverv. Det, som fordres,
er et Minimum, og det staar stedse Forældrene frit for at gaa
saa langt ud over dette, som de kunne og ville. ---

En videre gaaende Uddannelse kan Staten tilbyde, og kan
stille den som Betingelse for at opnaa offentlige Hverv, men
maa ikke paatvinge den. Det samme gælder om den religiøse
Livsanskuelse. Staten kan benytte sin Indflydelse til at virke
for, at de forskellige aandelige Retninger i Folket frit kunne udfolde
sig, men maa lade det være de enkelte Hjems Sag, i hvilken
Retning Barnets religiøse Liv skal ledes. Dens offentlige
Læreanstalter bør nærme sig saa meget til Konfessionsløshed,
som det idethele er muligt.
% --- p. 342 ---
\setcounter{footnote}{0}%

\opage{342}\grouphead{B. DET FRIE KULTURSAMFUND.}

\chapterhead{XXIII.}{FRIHED OG KULTUR.}

1. Allerede ved Familielivet maatte Maalestokken for Udviklingen
være den: i hvilken Grad det enkelte Individ stod ikke
blot som Middel, men som Formaal. Denne Maalestok maa finde
endnu større Anvendelse ved det frie Kultursamfund. Til Familien
hører Individet, ogsaa førend det har naaet sin Modenhed;
men det frie Kultursamfund, der bliver til ved Individernes Sammenslutning
om fælles eller forbundne Formaal, forudsætter, at
Individerne møde med modne og udviklede Kræfter til at arbejde
i deres Formaals Tjeneste.

Indenfor Familien arbejdes der ogsaa for Kulturformaal,
men der er Forbindelsen mellem Individerne det Første, det til
Grund Liggende, og de fælles Bestræbelser ere en Følge af denne
Forbindelse. I Kultursamfundet er det omvendt de fælles Formaal,
som vække fælles Bestræbelser og derved foraarsage Forbindelser
mellem Individerne.

Friheden og Kulturen ere de to Begreber, som i Forening
tydeliggøre det frie Kultursamfunds Væsen. De kunne synes at
staa i en vis Modsætning til hinanden: Friheden isolerer, Kulturen
forener; ved Friheden betones Selvstændigheden, ved Kulturen
Hengivelsen. Men de ere meget nøje knyttede til hinanden.
% --- p. 343 ---
\setcounter{footnote}{0}%
\opage{343}Saasnart Livets Gang fører det med sig, at Mennesker ere
fælles Betingelser underlagte, have fælles Skæbne, fælles Fare,
fælles Fjender, fælles Arbejde eller fælles Leg, vil der udvikle
sig en Solidaritetsfølelse mellem dem, som gør, at Frihed og
Hengivelse gaa i Et. Den Enkeltes Selvstændighed hindres da
ikke, men muliggøres netop ved, at han virker sammen med
de Andre eller for de Andre. Kulturens Udvikling er ikke
mulig uden en saadan Solidaritet, da den stiller Opgaver, der
kun kunne løses ved et Arbejde, der overstiger den Enkeltes
Kræfter. Men tillige ere Kulturopgaver saadanne, som ikke
blot komme den Enkelte til Gode. Kulturen frembringer Goder,
der direkte eller indirekte kunne faa Betydning for Alle, og
gennem Kulturarbejdet virker den Enkelte saaledes for Noget,
der er mere omfattende end hans egen lille Verden. Ved at
blive delagtigt i Kulturarbejdet udvikles Mennesket derfor baade
gennem den Solidaritet, Arbejdsfællesskabet frembringer, og
gennem den personlige Indleven i universelle Formaal, som er
Kulturens bedste Frugt.

Ofte kan der være en lang Skole at gennemgaa, inden
Fællesskabet i Arbejde og i Formaal ret bliver til. Der foregaar
da en Opdragelsesproces, som kan være streng og haard. Trægheden
og Raaheden, Egoismen og Blindheden maa overvindes,
men kunne gøre stærk Modstand. Livet i Kultursamfundet er
fra denne Side set en Fortsættelse af den Opdragelse, der finder
Sted indenfor Familien. Og det er stedse Kulturlivets vigtigste
Betydning at virke opdragende og udviklende paa det personlige
Liv. Hele Menneskehedens Historie kan betragtes som en
Opdragelsesproces, der aldrig bliver afsluttet, men særligt træder
frem, naar det er nye Formaal, der skulle arbejdes frem, og
naar Fællesskabet i Virken for disse Formaal derfor først skal
frembringes. Naar der allerede har dannet sig et Samfund om
et vist Kulturformaal som Midtpunkt, er Opgaven for det enkelte
Individ at leve sig ind i den færdige Kultur. Det bestemmes
da i sin Udvikling ved Traditionen; Efterligning og Suggestion
spille en stor Rolle, uden at det mærkes; færdige Forestillingsforbindelser
og Forskydningsresultater komme det til Gode og
bestemme det Niveau, fra hvilket det begynder sin Udvikling,
% --- p. 344 ---
\setcounter{footnote}{0}%
\opage{344}og udfra hvilket der saa --- hvis Selvstændighed og Kritik vaagne
--- kan indledes nye Opdragelses- og Forskydningsprocesser.

2. Man opstillede i det 18. Aarhundrede den personlige
Frihed som en selvindlysende Menneskerettighed\footnote[1]{I Teorien har Kant bedst formuleret Retten til personlig Frihed: „Frihed (Uafhængighed af enhver Andens tvingende Vilkaarlighed), saavidt den kan bestaa med enhver Andens Frihed efter en almindelig Lov, er den eneste oprindelige, ethvert Menneske som Menneske tilkommende Ret“. (Rechtslehre, 2. Ausg. Königsberg 1798. p. XLV). --- I Praxis proklameredes denne Menneskerettighed i de nordamerikanske Fristaters Uafhængighedserklæring, hvor det hedder: „Vi anse følgende Sandheder for selvindlysende: at alle Mennesker ere skabte lige; at de af deres Skaber ere begavede med visse uafhændelige Rettigheder, og at blandt disse findes Liv, Frihed og Stræben efter Lykke.“ --- (Det behøver vel ikke at bemærkes, at Frihed i den Betydning, hvori Ordet her bruges, ikke har noget at gøre med den „Villiens Frihed“, som blev drøftet i Kap. V).}. Men en
selvindlysende og absolut Ret kan Friheden ikke siges at være.
Den er et afledet Princip. Den behøver en Begrundelse, og
denne sker ved Slutninger ud fra Velfærdsprincipet, af hvilket,
som vi allerede have set (VIII, 6), den frie Personligheds
Princip udspringer gennem en psykologisk Betragtning. De
personlige Væsener ere de Centralsteder i Verden, hvor Livets
Værdi føles; en Krænkelse af et saadant Væsens uhindrede og
harmoniske Livsudfoldelse strider derfor mod Velfærdsprincipet.
Ethvert Indgreb i Personlighedens frie Udvikling maa begrundes
ved Hensynet til selve denne Udviklings fortsatte Forløb; det
kunde jo være, at en Hæmning i det enkelte Øjeblik eller paa
det enkelte Punkt var nødvendig, for at den hele Udvikling
kunde gaa saa frit og harmonisk for sig som muligt. Paa Grund
af den nøje Sammenhæng mellem Frihed og Velfærd, er Friheden
selv et Formaal. Men den er tillige et Middel, der anvises af
Velfærdsprincipet. Ti uden et selvstændigt og harmonisk Liv
hos de Enkelte kan Slægtens Liv idethele ikke gaa frem.
Initiativet, den nye Begyndelse til det Store og Gode, bliver
kun til, naar Livet faar Lov at røre sig frit hos de enkelte
Individer. Tvang og Autoritet ville her være utilstrækkelige;
de kunne forslaa til en Tid at holde det en Gang Erhvervede
% --- p. 345 ---
\setcounter{footnote}{0}%
\opage{345}fast, --- og dog vil der til selvstændig Værdsættelse af det
Overleverede og til Begejstring for at hævde det kræves en
aandelig Frihed og en Uafhængighed af materielle Skranker,
som ikke blive til, hvor de Enkelte ere Herredømme og Formynderskab
underlagte.

Men skønt vi saaledes nok kunne skelne mellem Friheden
som Maal og Friheden som Middel, hænge de nøje sammen og
betinge hinanden. For at kunne have Friheden som Middel,
maa man tildels have opnaaet den som Maal; for at kunne virke
som Kraftcentrum maa man have opsamlet Kraft i sig. Og
omvendt: kun gennem fri Brug af sine Kræfter bliver man en
virkelig Personlighed. Vi genfinde atter her det aristoteliske
Princip: vi blive kun frie ved at handle med Frihed. Heri
ligger netop det Vanskelige ved Frihedsprincipets Anvendelse.
Ikke blot er Friheden et afledet Princip, men det kan ofte ---
baade som Maal og Middel --- kun indirekte anvendes, da de
Individer, der skulle gøres frie, kunne behøve Støtte og Omsorg
for at naa saa vidt, at de kunne virke paa egen Haand.
Anvendelsen af aandelig og fysisk Magt kan være berettiget i
Frihedens Navn; men Anvendelsen af Magt eller Tvang, Autoritetsprincipet
i enhver Form, er dog stedse afledet i Forhold til
Frihedsprincipet, ligesom selve Frihedsprincipet er det i Forhold
til Velfærdsprincipet. Velfærd er Maal; Frihed er baade
Maal og Middel; men Autoritet er egentligt kun Middel, et
Middel, som skal forsvinde, naar Maalet er naat.

Kun naar Autoriteten indtager den herved anviste Plads,
er den af etisk Natur. Mellem Autoritetsprincipet og Frihedsprincipet
staar den store Kamp i Historien. Det er de to Kræfter,
hvis Vexelvirkning bestemmer den sociale Udvikling. De
have hver sin Tid og sit Omraade, som bestemmes ved Velfærdsprincipet.
--- Det gør en meget stor Forskel, om Autoriteten anvendes
som Opdragelsesmiddel eller som ubetinget Princip. Den
hele Aand og Maade bliver en anden, hvor den betragtes som
Middel, end hvor den betragtes som Maal. Naturligvis forudsættes
det her, at Anerkendelsen af, at Autoritetsanvendelsen kun
er Middel, er oprigtig. Ti det er meget vel muligt at have en
abstrakt Begejstring for Friheden som et fjernt Maal og dog i
% --- p. 346 ---
\setcounter{footnote}{0}%
\opage{346}sin virkelige Færd anvende Magt og Myndighed paa den mest
brutale Maade. Det er en af de mange Former, i hvilke den altid
betænkelige og ikke altid ærlige Distinktion mellem Teori og
Praxis fremtræder. Maalestokken for Kulturens og Kultursamfundets
Udvikling bestaar i den Grad, i hvilken Friheden ikke
staar blot som en teoretisk Ide, men --- baade som Maal og som
Middel --- ligger til Grund ved de enkelte Livsforholds Ordning.

Ved at opstille Frihedsprincipet som det underordnede Princip,
der først og fremmest udledes af det almindelige Velfærdsprincip,
skille vi os fra \textsc{Bentham}, for hvem Sikkerhed og Lighed
stod som de første afledede Principer. Men det er klart,
at Sikkerheden kun faar sin Værdi som Betingelse for fri Besidden,
Nyden og Virken. Usikkerhed hæmmer, kuer og deler.
Det er ikke Friheden, der er Betingelse for Sikkerheden, men
omvendt. Og hvis man vil Sikkerhed for enhver Pris, kommer
det let til at gaa ud over Friheden. Der maa voves, for at Udviklingen
kan gaa fremad, og det er ved ethvert saadant Vovestykke
Tilliden til de frie Kræfters Virkninger, som holder oppe.
Ved at fæste sit Blik paa Sikkerheden i første Linie kommer
man let til at tilstoppe Udviklingens Kilde. Af Frygt for Faren
ved at fortsætte Livet stanser man Livet. Ogsaa hvad Ligheden
angaar, er det klart, at den er en underordnet Livsbetingelse.
Hvad det gælder om, er at faa det personlige Liv frem i de
Enkelte. Men da disse, jo mere deres Natur kendes, vise sig
at frembyde kvantitative og kvalitative Forskelligheder, ville de
kun i rent udvortes Forhold kunne behandles aldeles ens. Den
fordelende Retfærdighed kræver, som vi allerede have set (III,
9; IV, 2; XI, 9), en dybtgaaende Individualisering. Ogsaa Ligheden
er da en Betingelse, der er Friheden underordnet. Fordelingen
af materielle og aandelige Goder og Fastsættelsen af
de Fordringer, der med Hensyn til materiel og aandelig Kultur
ere at stille til de Enkelte, maa ske med Henblik til deres
Evne til at fyldestgøre Frihedsprincipet, altsaa til deres Evne
til at udvikle sig frit og harmonisk og til at virke med Frihed
i Kulturformaalenes Tjeneste.

3. Det var det 18. Aarhundredes store Ære, at det opdagede
Frihedsprincipet og dermed fastslog en af de vigtigste
% --- p. 347 ---
\setcounter{footnote}{0}%
\opage{347}Betingelser for en højere menneskelig Samfundsorden. Men i
sin Begejstring over det nye Princip oversaa man ikke blot
dets afledede Natur og de Begrænsninger, det underligger; man
opstillede ogsaa som selvindlysende og oprindelig Rettighed,
hvad der først efterhaanden har arbejdet sig frem i Historien
og stedse kun tilnærmelsesvis kan virkeliggøres gennem historisk
Udvikling.

Ufriheden opstaar ved Arbejdets Deling. En af de simpleste
Former for Arbejdsdeling er den, at den Stærkere læsser
alt ubehageligt Arbejde paa den Svageres Skuldre. Kvinder og
Børn var de første Slaver. Den overvundne Fjende gøres til
Slave, naar de kannibalske Instinkter ikke føre til, at han
strax fortæres. Den økonomiske Interesse virker, hvor etiske
Motiver endnu ikke kunne virke. Medens et omvankende Jægerfolk
ofte vil være nødt til at skille sig ved de Overvundne, vil
et Nomadefolk eller en agerdyrkende Stamme finde sin Regning
ved at benytte sig af deres Arbejdskraft. Hvad der først
grundlægges af egoistiske Motiver og betegner et økonomisk
Fremskridt, bliver til et etisk Fremskridt, idet Samlivet og den
Omsorg for Andre, som egoistisk og økonomisk Interesse motiverer,
udvikle Sympatien, i hvert Tilfælde til en vis Grad. Det
er i Herrens egen Interesse, at Slavens Sundhed og Kraft ikke
lide. Dette Forhold træder især frem, hvor der findes Livegenskab
eller Vornedskab, ikke personligt Slaveri. Ogsaa dette skyldes
Egoisme og økonomiske Motiver. Livegenskabet kan opstaa dels
ved Erobring, dels ved Overgang fra personlig Slavestilling til
Bundethed ved Jorden, dels ved at frie Mænd i farefulde Tider
give sig under Mægtigeres Beskyttelse, og dette Beskyttelsesforhold
saa --- maaske ved vilkaarlig Tvang fra Beskytterens
Side --- gaar over til Stavnsbundethed. Dets Udvikling forudsætter,
at der henligger store Omraader af udyrket Jord, som
kun paa denne Maade kunne gøres frugtbringende '). Den 
% CHECK p. 347: note whose mark was not found
\footnote[1]{Motivet til Stavnsbaandets Indførelse i Danmark var Ønsket om at hindre Bønderne i at flytte fra de ufrugtbare til de frugtbare Egne, hvorved mange Godser vilde komme til at mangle Arbejdskraft. V. Falbe Hansen: Stavnsbaandsløsningen og Landboreformerne. I, p. 6.}Liv%
% --- p. 348 ---
\setcounter{footnote}{0}%
\opage{348}egne har et større personligt Spillerum end Slaven, flere Motiver
til Arbejde og Selvudvikling. Og Forholdet mellem Jorddrot
og Livegne behøver ikke blot at være et Magtforhold. Jorddrotten
har den Pligt at værne om sine Undergivne, og han
staar som Mellemled mellem dem og den øverste Statsmagt.
Samfundet frembyder en Række af Trin, hvor det højere Trin
værner om det laveres materielle og aandelige Vel, medens det
lavere med Troskab og Tillid følger det højere i Kamp og Arbejde
for materielle og aandelige Formaal. Arbejdsdelingen bliver
nu den, at Nogle skulle øve Sværdets og Aandens Værk, beskytte
Samfundet og tænke for det, medens Andre til Gengæld
skulle besørge det materielle Arbejde. En saadan Arbejdsdeling
fandtes allerede --- med personligt Slaveri som Grundlag --- hos
Grækerne, og opstod ogsaa --- med Livegenskab som Grundlag
--- i den kristelige Middelalder, kun at her Sværdets og Aandens
Arbejde øvedes af forskellige Stænder. Da Ridderen hørte op
at kæmpe, og da Andre end Præsten begyndte at tænke, var
den etiske Berettigelse af denne sociale Ordning forbi. Frihedsprincipet
havde allerede den stoiske Filosofi\footnote[1]{Den første principielle Bekæmper af Slaveriet synes dog Antisthenes, Stifteren af den kyniske Skole, at have været. Han kom til sit Resultat især ad to Veje. Mennesket maa begrænse sine Fornødenheder for at bevare sin aandelige Frihed: men saa maa det ogsaa kunne undvære Slaver! Og selve den aandelige Frihed kan findes hos Slaven; det sande Slaveri er Afhængighed af sanselige Lyster; den ydre Forskel mellem Slave og Herre har ingen væsentlig Betydning. Det er muligvis især Kynikerne, Aristoteles polemiserer imod i sin Statslære. (Joël: Der echte und der xenophontische Sokrates. II, 2. p. 569—574. Kfr. I, p. 406 f.). Den kyniske Tankegang blev fortsat af Stoikerne, især efterat Karneades havde vist, at Slaveriet stred mod den stoiske Etik. (Smlgn. Schmekel: Die Philosophie der mittleren Stoa. p. 228. 378 f.). Paa en med Kynisme og Stoicisme analog Maade virkede Vedantalæren og Buddhismen i Indien. (Deussen: Geschichte der Philosophie. I, 1. p. 163).} og Kristendommen
udtalt uden at drage de praktiske, sociale Konsekvenser af det
(hvilket igen hang sammen med, at de kun betragtede Friheden
som Maal, ikke ogsaa som Middel). Nu kom den Tid, da det
kunde proklameres som socialt Princip. Men at dette kunde
ske paa praktisk Maade, skyldes den Omstændighed, at Friheden
% --- p. 349 ---
\setcounter{footnote}{0}%
\opage{349}allerede længe havde rørt sig. Paa Industriens og Handelens
Omraade havde der virket Mennesker, som hverken vare Livegne
eller Jorddrotter. Det var deres Erfaringer og den ved
disse Erfaringer vundne Tillid til de frie Kræfter, som bevirkede
den store Tilslutning til Frihedsprincipet. Ved den store
Række Emancipationer, der begynder i sidste Halvdel af det
18. Aarhundrede, virke økonomiske og etiske Motiver sammen,
ligesom paa tidligere Udviklingstrin.

Kun paa raa Trin af Menneskevirksomhed kunne Slaver
og Livegne gøre Fyldest. De have ikke den store Interesse for
Arbejdet som de frie Mænd. Fra en af de Egne i Danmark,
hvor Bønderne i det 18. Aarhundrede stod højest, og hvor
Hoveriforholdene vare bedst ordnede, blev det udtalt, at ti Hoveriarbejdere
ikke udrette saa Meget som to lejede Arbejdere.
Slaver og Hoveriarbejdere staa sig ikke ved at gøre Mere, end
der lige netop forlanges af dem; det vilde jo kun føre til, at
Fordringerne for Fremtiden forhøjedes. De have ingen Grund
til at vise særegen Omsorg for Redskaber og Materiale, som
ikke tilhører dem. Hvad der spildes, er kun til Tab for Herren.
Emancipationerne have derfor nok saa meget været til Fordel
for Arbejdsgiverne som for Arbejderne, og Bondeemancipationen
har medført en større Produktivitet af Jorden\footnote[1]{Smlgn. Roscher: Die Grundlagen der Nationalökonomie § 71. --- v. d. Goltz: Landwirthschaft. (Schönbergs Handbuch. 1. Udg. I). p. 580. --- I et (hos Henry George: Social Problems. Kap. 15 aftrykt) Brev fra en tidligere Slaveejers Søn hedder det: „Planterne ere fornøjede med Forandringen. De sige: \textasciitilde{}„Hvor dumt var det ikke af os at føre Krig for Slaveriets Skyld! Vi faa Arbejdet billigere nu, end da vi ejede Slaverne.“ Hvorfor faa de det billigere? Fordi de i Form af Renter tage Mere af Negerens Arbejde, end de kunde under Slaveriet; ti den Gang vare de nødte til at yde ham tilstrækkelig Føde, Klæder og Lægebehandling for at holde ham rask, og de nødtes af deres Samvittighed og af den offentlige Mening saavel som af Loven til at underholde ham, naar han ikke længere kunde arbejde. Nu ophører deres Interesse og deres Ansvar, naar de have faaet alt det Udbytte ud af ham, de formaa.“ --- Falbe Hansen: Stavnsbaandsløsningen og Landboreformerne. I, p. 145: „Hoveriets Fastsættelse og Afløsning, Udskiftningen, Afløsningen af Naturafgifterne m. m. vare Reformer, der ligesaa meget vare i Godsejernes Interesse som i}.

% --- p. 350 ---
\setcounter{footnote}{0}%
\opage{350}Paa den anden Side, --- hvor Arbejdet var af den Natur,
at det øvede Arbejdernes Kraft og Intelligens, dér fremtraadte
baade Trangen til og Nødvendigheden af personlig Frihed. Haandværkerne
have derfor i Regelen opnaat denne før Landarbejderne.
Den aandelige Selvstændighed vækkes ved Følelsen af
Selvvirksomhed og Tilliden til sine egne Kræfter. Den frie
Stræben overføres fra det ene Omraade til det andet. Naar det
fysiske Arbejde gaar bedst med Frihed, saa maa det aandelige
Arbejde ogsaa gaa bedst med Frihed.

Men en vis aandelig Udvikling er overalt nødvendig, for at
Trangen til Frihed skal føles. Ufrihed føles kun som et Onde
af den, der stræber ud over de Skranker, der ere satte ham.
Den Slave, der har en mild Herre, kan ofte langt bedre og
sikrere faa sine materielle Fornødenheder tilfredsstillede end
den frie Mand. Slaven føler sig ofte som et Barn og ønsker
ligesaa lidt Frihed, som Barnet ønsker at forlade sit Hjem. Han
holdes borte fra enhver Lejlighed til at anstille Sammenligninger\footnote[1]{fredse.'' Paa lignende Maade finder man undertiden i Lande, hvor der hersker Polygami, at en Kvinde er stolt af sin Mands mange Hustruer: det er Tegn paa, at det er en mægtig og rig Mand, hun er gift med.},
og nye Ønsker og Drifter opstaa derfor ikke saa let hos ham.
(Smlgn. VII, 3). Da Trangen til Frihed mangler, føles Ufriheden
ikke som et Onde. Trællens Stilling nærmer sig paa tidligere
Kulturtrin til Barnets og behøver ikke at være forbunden med
Bondens. De bevirkede, at Jordens Udbytte steg, og at dens Værdi forøgedes;
men det maatte jo nærmest komme dens Ejer, Godsejeren, til
Gode.“ (Se ogsaa p. 68. 101. 149). Det er gaaet paa samme Maade paa
det industrielle Omraade; især er Arbejdsdagens Forkortelse ogsaa kommen
Fabrikejerne til Gode (se nedenfor XXVI, 14).

') Det var i de nordamerikanske Slavestater under streng Straf forbudt
at lære en Slave at læse og skrive, endsige at indføre Bøger, der
handledeom Slaveemancipationen. Lyell:  Reisen in Nordamerika. Deutsche
Uebers. p. 118. 120 --- Lyell anfører ogsaa Exempler paa, hvor glade og
stolte Negerslaver kunde være ved deres Stilling. „Jeg tilstaar,“ siger den
berømte Naturforsker, „at man kan tænke og filosofere over denne ejendommelige
og interessante Art af Forfængelighed, indtil man i den finder
Beviset for den yderste sociale Nedværdigelse; men det første Indtryk,
den gjorde paa mig, var meget trøstende, saa at jeg umuligt kunde føle
en smertelig Medlidenhed med Folk, der følte sig saa overordentligt til%
% --- p. 351 ---
\setcounter{footnote}{0}%
\opage{351}den Nedværdigelse, som man paa et senere og højere Kulturtrin
finder i denne Stilling. Tværtimod kan der, især naar den
hele Ordning betragtes som en Naturorden, som Noget, der
ikke kan være anderledes, og naar fælles Erfaringer i Lykke
og Ulykke komme til at virke, udvikle sig et hjerteligt Forhold
mellem Trællen og hans Herre, som mellem Kammerater og
Medstridende. Overfor et saadant Forhold vilde Forkyndelsen
af Frihedsprincipet virke opløsende. Intetsteds viser det sig som
her, hvorledes de almindelige Principers Værdi betinges i hvert
enkelt Tilfælde ved de specielle Forhold, de skulle anvendes
paa. Hvad enten Ufrihedstilstanden har hin idylliske og hjertelige
Karakter, som antydedes, eller den bliver en Undertrykkelsestilstand,
hvor Magten og Herredømmet ikke bruges til
Værn, men til at tilfredsstille Lunefuldhed, Magtsyge og Nydelsessyge,
saa vil en lang Ufrihedstilstand --- for den Enkelte
og for Slægten --- ikke paa én Gang kunne afløses af en Tilstand,
hvor Friheden kan ytre sig baade som Middel og Maal.

Da Udviklingen hæmmes, saalænge Ufriheden bestaar, idet
baade Drift og Midler mangle til den, kan man ikke vente, at
den blotte Proklamation af Friheden strax skal hidføre en helt
ny Tilstand. Frihedens første Virkninger ere sjeldent aldeles
heldige. Der foregaar jo en Overgang til helt nye Livsvilkaar,
som Individets Natur endnu ikke har kunnet lempe sig efter.
De Frigivne havde et ilde Ry paa sig i Rom. Hverken Negerslaverne
eller de russiske Bønder kunde strax bruge deres Frihed
paa rette Maade. Om Virkningen af Landboreformerne paa de
danske Bønder udtales der paa Grundlag af en historisk Undersøgelse\footnote[1]{Falbe Hansen. p. 87.}:
„Reformernes væsentligste Betydning for Nationens
Økonomi var den, der fremkom ved, at Bondestandens moralske
og intellektuelle Niveau blev højnet. Landboreformerne virkede
ganske vist ogsaa umiddelbart til at ophjælpe Landbruget . . .
men endnu større var vistnok den middelbare, den indirekte
Indflydelse, der fremkom gennem Reformernes opdragende Indvirkning
paa Bondestanden og paa Befolkningen som Helhed.
Det er imidlertid indlysende, at denne indirekte Virkning først
% --- p. 352 ---
\setcounter{footnote}{0}%
\opage{352}kunde vise sig sent. En Slægt af forkuede, trælbundne Bønder
bliver ikke paa én Gang forandret, fordi den blotte Frihedserklæring
udstedes; der maa en helt ny Generation til for at se
Virkningen heraf give sig Udslag i Landbedriften. Og hvad t. Ex.
Hoveriet angaar, da kunde en af dettes værste Ulemper, at
Bonden vænnede sig til at arbejde trægt, og at Husbonden
ødslede med Arbejdskraften, først sent forsvinde; man har endog
ment at spore Hoveriets Eftervirkninger i Landmændenes Arbejdsmaade
mere end et halvt Aarhundrede efter dets Afskaffelse.“
--- Med ét Slag forandrer Menneskenaturen sig nu engang ikke.
Først efterhaanden, gennem Overgangstilstande, der maaske vare
i flere Slægtled, vil Modenheden kunne naas. Og naas fuldstændigt
kan den kun gennem selve Brugen af de frie Kræfter.
Det gælder om at finde foreløbige Former, i hvilke Selvvirksomheden
kan bevæge sig, indtil den formaar at paatage sig
de større Opgaver.

4. Man satte ved Proklamationen af Menneskerettighederne
Formen over Indholdet, ja mente endog, at Formen kunde gøre
det Hele. Man oversaa, at den Grad, i hvilken Friheden til en
given Tid kan virkeliggøres og gennemføres, beror paa en hel
Række forskellige sociale Forhold. Indenfor den gamle Samfundsorden
vare Grænserne snevert dragne, og der var tunge Byrder
at bære. Men af dem blandt de Privilegerede, der opfattede
deres Magt som et overdraget Hverv, udvistes der en Omsorg
og en faderlig Følelse, som gjorde, at de Afhængige aldrig
kunde føle sig aldeles forladte. Nu derimod, da Pligten faldt
bort med Magten, stødtes det emanciperede Individ ud i Verden
med sit Frihedsbrev i Haanden uden maaske at vide, hvad det
egentligt skulde bruge sin Menneskerettighed og sin Frihed til.
Et Beskyttelses- og Pietetsforhold blev pludseligt forvandlet til
et Retsforhold. Dette traadte ikke mindst frem ved Ophævelsen
af de gamle Lav og de dertil knyttede Institutioner. Man ophævede
derved enhver Organisation af Arbejdet og overlod til
de Enkelte at fægte sig frem med egne Kræfter. --- Karakteristisk
nok var det den samme Tid, i hvilken man gjorde en
skarp Adskillelse mellem Etik, Nationaløkonomi og Retslære og
mente her at have tre. aldeles sondrede Omraader for sig. I
% --- p. 353 ---
\setcounter{footnote}{0}%
\opage{353}Nationaløkonomi skulde der kun tages Hensyn til Erhvervsdriften;
lod man blot den frit Spil, vilde Menneskenes økonomiske
Interesser af sig selv komme i Harmoni. I Retslæren
skulde der kun søges Betingelserne for en saadan ydre, mekanisk
Ordning, at Sikkerhed og Frihed for Alle blev mulig; den
skulde holde sig til den ydre Handlen (Legaliteten) og ganske
se bort fra Sindelaget (Moraliteten).

Man gik tillige ud fra en Opfattelse, der betragtede ydre
Forhold, Opdragelse og sociale Vilkaar som eneste Aarsag til
alle Forskelligheder mellem Mennesker. Alle Mennesker ---
mente man --- vare lige i Natur og Begavelse; det var kun
de ydre Forhold, under hvilke de udviklede sig, der medførte
Forskellighederne. \textsc{Adam} \textsc{Smith}, hvis Anskuelser fik saa stor
Indflydelse paa den almindelige Opfattelse af sociale Spørgsmaal,
gik saaledes ud fra, at Forskelle i Karakter og Begavelse ikke
saa meget skrive sig fra Naturen som fra Vane, Levevis og
Opdragelse, ikke saa meget ere Aarsager til Arbejdets Deling i
Samfundet, som Virkninger af denne\footnote[1]{Om Nationalvelstands Natur og Aarsag. Dansk Overs. I, p. 25. — Smith havde denne Lære fra Helvetius. \{De l'esprit. III, chap. 26—27).  --- Endnu James og Stuart Mill hyldede den.}. Naar man da aabnede
alle Bomme og holdt al Tvang borte, syntes der at være Udsigt
til, at Alle kunde naa frem. Retfærdighed maatte fordre,
at Betingelserne blev lige for Alle, og heller ikke Naturen gav,
mente han, den Ene noget Forspring for den Anden.

Man lededes herved endnu et Skridt videre. Man forbød
endog alle frie Associationer. Eller rettere, man kunde, efter
hvad man havde erfaret ved de gamle Lavs Udartning, ikke
tænke sig nogen Association uden Tvang og Tyranni. Turgot
fik ved et Edikt af 1776 Lavene ophævede. Han saa (som det
udtales i Indledningen til dette Edikt) Kilden til alle Onder paa
det industrielle Omraade i den Ret, Haandværkerne af samme
Fag havde haft til at samles og forene sig i Foreninger. Enhver
fik nu Ret til at udøve, hvilket Haandværk eller hvilken
Handel han vilde; men det blev forbudt alle Mestre, Svende
% --- p. 354 ---
\setcounter{footnote}{0}%
\opage{354}og Arbejdere under nogetsomhelst Paaskud at danne Samfund
og Foreninger. Da Turgot blev styrtet, blev denne Reform ophævet,
ligesom alle hans andre Reformer; men Revolutionen
optog den igen. En Lov af 14. Juli 1791, altsaa fra Revolutionens
første Tid, erklærede Tilintetgørelsen af alle Korporationer
af dem, der havde samme Stand og Gerning, for en
af den franske Forfatnings Grundvolde og forbød Genoprettelse
af saadanne Samfund under nogensomhelst Form. Man havde
ikke Lov at samle sig og under bestemte Former raadslaa om
sine „foregivne fælles Interesser“ (interéts prétendus communs)!\footnote[1]{Loven af 14. Juni 1791 som et Angreb af Bourgeoisiet paa Arbejdernes Frihed.}.
Dette Skridt er karakteristisk for den abstrakte og overspændte
Opfattelse af Frihedsprincipet. Man frygter Associationerne, fordi
de let kunne frembringe nye Forskelligheder, nye ejendommelige
sociale Grupper, og i sin Iver for at værne Friheden mod
denne Fare gaar man saa vidt, at man indfører en af de allerværste
Indskrænkninger i Friheden ved Forbuddet mod fri
Sammenslutning med Andre. Ja, man tillægger endog Statsmagten
Myndighed til at afgøre, hvorvidt man idethele har
Interesser sammen med Andre! Den nævnte franske Lov er et
betegnende Exempel paa Tilbøjeligheden til kun at have den
Fare for Øje, som ligger nærmest, og at vise saa stor Iver
for at afværge den, at man lægger store Hindringer i Vejen for
den fortsatte Udvikling. Man mente i det 18. Aarhundrede, at
Alt var gjort, naar Skranker og Baand afskaffedes. Man tænkte
ikke paa, at fordi de gamle Lav havde virket skadeligt, var det
ikke sagt, at alle Associationer vare skadelige. Man hildede
sig i den store Modsigelse, at man i Frihedens Navn forbød
en af de vigtigste Anvendelser, der kan gøres af Friheden:
Indgaaen af frit Samfundsforhold med andre Individer.

Men denne Overspændthed hindrer dog ikke, at det var
et overordentligt betydningsfuldt Princip, som proklameredes.

•) Léon Say: Turgot. Paris 1887. p. 150---159. --- Taine: La Révolution.
I. p. 221 f. --- Det skyldes sikkert en agitatorisk Brug af Historien,
naar Karl Marx (Das Kapital. 2. Udg. I. p. 772) fremstiller

% --- p. 355 ---
\setcounter{footnote}{0}%
\opage{355}Hvad det gælder om, er, at Friheden saavidt muligt bliver
bestemmende ved Arbejdsdelingen, saa at denne ikke, som paa
de primitive Trin, sker ved Tvang og saaledes, at det ubehagelige
Arbejde læsses over paa de Svage. Den Enkelte maa kunne
følge sin Drift og Evne, saavidt han er sig dem bevidst, og saavidt
han derved ikke krænker større Interesser. Friheden bør
være Mulighed for at binde sig til noget Bestemt, for at vælge et
Kald. I den gamle Samfundsorden herskede sociale Forskelligheder,
der næsten havde Kasteforskelles Karakter, og hver Enkelts
Plads var ham anvist fra først af. I det frit valgte Kald bestemmer
derimod Individet sin egen Plads, det Punkt, udfra hvilket
han vil og kan virke for Slægtens Opgaver.

Derved bliver en fri Samfundsdannelse mulig. Alle Tanker
og Formaal maa først opstaa i enkelte Menneskers Bevidsthed,
før de gøre deres Virkning i den øvrige Verden. De virke samfundsdannende
ved at samle dem om sig, der forstaa deres Betydning;
og naar Tanken saaledes er prøvet og anvendt i mindre
Kredse, kommer den Tid, da Statens fastere, i sidste Instans
paa Tvangsmidler beroende Organisation kan træde til. I det frie
Arbejde paa Kulturen udfoldes de produktive Kræfter; det er
da Statens Sag med sin Retsorden og sine Magtmidler at beskytte,
hvor Beskyttelse er mulig og gavnlig. Statens Forhold
til de frie Kultursamfund er altsaa analogt med dens Forhold
til Familien. Intet af Stederne frembringer den de egentligt
virkende Kræfter; begge Steder staar den overfor Noget, som
ikke er af den Verden, i hvilken den selv hører hjemme; ---
men begge Steder kan dens Virksomhed være nødvendig til at
organisere og hjælpe. Og det er ingenlunde saa, at Organisationen
først er Statens Sag. Levekraftige Organisationer maa,
som Historien viser, voxe frem gennem Sammenslutning af frie
Kræfter, der under Livsforholdenes og Livsinteressernes Indflydelse
finde hverandre og mere eller mindre bevidst supplere
hverandre eller arbejde mod samme Maal. De Former, disse
Organisationer antage, kunne kun blive fuldkomne ved at prøves
overfor de virkelige Forhold. Nogle af de værdifuldeste Erfaringer,
Menneskeslægten har gjort, hidrøre fra denne uvilkaarlige
eller vilkaarlige Tillempelse af Samlivsformerne efter de faktiske
% --- p. 356 ---
\setcounter{footnote}{0}%
\opage{356}Forhold. Det er Erfaringer, hvis Virkninger ikke blot angaa de
ydre Former, men forgrene sig ind i Sind og Tanke. Der er en
stadig Vexelvirkning mellem det Indre og det Ydre. Et af de
vigtigste Problemer i den sociale Etik er det, hvorvidt den sociale
Udvikling kan overlades til den frie Tillempelse, af hvilken den
frie Association er en af de betydningsfuldeste Former, og hvor
det Punkt ligger, hvor Staten maa gribe ind med sin tvingende
Myndighed. Idet Revolutionen proklamerede ikke blot Friheden
og Ligheden, men ogsaa Broderskabet, udtalte den, at den blotte
Emancipation ikke er nok, men at der behøves positiv Sammenslutning.
Men Forholdet mellem disse to Sider har vist sig at
være langt mere indviklet, end man i den første Begejstring saa,
og de sociale Problemer komme netop frem, naar man overtænker
Forholdet mellem de forskellige Led i den revolutionære
Formel. De stille sig dog naturligvis i noget forskellig Form
paa Kulturens forskellige Omraader.

5. Der gives lige saa mange forskellige Kulturformaal, som
der gives forskellige Omraader, hvor der kæmpes for Livet.
Livet hviler ikke blot paa den fysiske Selvopholdelse, Tilvejebringelsen
af de materielle Betingelser for at bestaa; paa højere
Trin er Livets Værdi knyttet til ideelle Formaal, Tilfredsstillelsen
af Tankens, Fantasiens og Følelsens Trang; og der er tillige en
Trang til at lade de materielle og aandelige Goder falde i saa
Manges Lod som muligt. Til disse tre Arter af Kultur svare
tre Arter af frie Kultursamfund.

Den materielle Kultur udspringer umiddelbart af Selvopholdelsen.
--- Dersom der ikke var det stadige Tryk paa Menneskene,
som „Villien til Livet“ volder, vilde de neppe underkaste sig
den rastløse Arbejden og det møjsommelige Slid i den materielle
Kulturs Tjeneste. Dette Tryk har drevet Menneskene frem fra
lavere til højere Trin, endnu før ideelle Motiver kunde gøre sig
gældende. --- Paa det laveste Trin af menneskelig Existens bestaar
Næringen af vilde Dyr og Planter, som man opsøger. Man
kender maaske ikke engang Ilden. Kun ved Sprogets Evne og
ved Brugen af simple Redskaber adskiller Mennesket sig her fra
Dyret. Dette Trin har man kaldet det egentligt vilde Stadium.
--- Et noget højere Trin er det barbariske, hvor Ilden kendes,
% --- p. 357 ---
\setcounter{footnote}{0}%
% CHECK p. 357: notes paired with the marks in order (the layer misread a number)
\opage{357}og hvor man kan danne sig Metalredskaber og øve Agerdyrkning
og Kvægavl. --- Den egentlige Civilisation begynder, naar
Skriftsproget er opfundet, og dermed en sikker og udbredt Erindring
og Overlevering i Slægten bliver mulig\footnote[1]{Smlgn. Tylor:  Anthropology. p. 24. 179 f.}. Skriftsproget
har vel først tjent rent praktiske Formaal, som Regnskaber, Fortegnelser,
Meldinger. Men ved en Værdiforskydning fik dette
Middel til Bevarelse og Meddelelse af Erfaringer og Forestillinger
en Betydning, der fuldendte den Overgang fra materiel til
ideel Kultur, som allerede var begyndt, hvor der var Interesse
for mundtlig Overlevering. --- Det er af stor etisk Interesse, at
jo mere den materielle Kultur udvikler sig, des mere gør den
det muligt at tilvejebringe og styrke Forbindelse og Samkvem
mellem Menneskene, og at den tillige nærmer sig den ideelle
Kultur og fører over til den. I disse Hensyn maa Maalestokken
for den materielle Kulturs Udvikling og særligt for Fuldkommenheden
af de Former, hvorunder den udvikler sig, søges.
Jo mere den materielle Kultur virker forbindende, og jo mere
den fører over til ideel Kultur, des højere staar den.

Den ideelle Kultur bliver til, saasnart der opstaar Formaal,
der gaa ud paa Andet og Mere end det at opholde Livet.
Selvom det er de samme Kræfter, der virke her og i den
materielle Kultur, saa anvendes de her for deres egen Skyld,
formedelst den umiddelbare Tilfredsstillelse, der er forbunden
med deres Brug. Denne Brug af Kræfterne for selve Virksomhedens
Skyld er et Tegn paa, at et sandt menneskeligt Livstrin
er naat. At det naas, skyldes en Motivforskydning (se XIII, 4).
Ikke blot det „Nødvendige“, men ogsaa det „Skønne“ (for at
bruge Grækernes Udtryk)\footnote[2]{Ved det „Nødvendige'' forstod Grækerne det, der er Middel for noget Andet, ved det \textasciitilde{}Skønne“ det, der i sig selv er Formaal. Smlgn. Aristoteles: Rethor. 1, 9.} bliver nu Genstand for Interesse
og Bestræbelse. Mennesket stræber nu ikke blot for at leve,
men lever for at stræbe. Nu opstaar der Videnskab og Kunst,
et æstetisk og et religiøst Følelsesliv.

Den filantropiske Kultur har Tilfredsstillelsen af Menneske%
% --- p. 358 ---
\setcounter{footnote}{0}%
\opage{358}kærligheden til Genstand. Den gaar særligt ud paa at støtte
dem, der kæmpe paa de laveste Trin af materiel og ideel Kultur.
Den vil raade Bod paa legemlig og aandelig Forkommenhed,
hvor den findes. Særligt vil den søge at hindre, at der udvikler
sig en dualistisk Forskel mellem materiel og ideel Kultur. Men
den er rettet mod al Lidelse og Nød, legemlig og aandelig,
hvor den findes. Ofte kan der trænges til en Haandsrækning,
netop hvor der er fuldt op af materielle og aandelige Goder.
--- Den filantropiske Kultur behøver ikke at træde frem som
en særlig Bestræbelse ved Siden af de to andre Arter af Kultur,
men kan være forbunden med dem og betegne den Aand og
Retning, i hvilken de dyrkes. Men skjult eller aabenbart maa
den være tilstede, dersom Udviklingen skal være sund og kraftig.

\grouphead{1. DEN MATERIELLE KULTUR.}

\chapterhead{XXIV.}{SOCIALE MODSÆTNINGER.}

1. Om materiel Kultur bliver der først Tale, naar Mennesket
maa arbejde for at opholde Livet. Hvor det finder Midlerne
til Selvopholdelse i Naturen, f. Ex. Frugter, der trives uden Dyrkning,
bliver det staaende paa et halvt dyrisk Trin. Men naar
den materielle Kultur har gennemløbet sine første Stadier, kommer
der et Punkt, hvor man med Ringeagt ser ned paa de Virksomheder,
der gaa ud paa at skaffe de første og uundværligste
Livsbetingelser tilstede. Det hænger sammen med den Arbejdsdeling,
der tidligt indtræder. De Stærke vælge selv det Arbejde,
de ønske at udføre, og overlade Resten til de Svage. Paa det
barbariske Trin er det Jagt og Krig, som ene anses en Mand
% --- p. 359 ---
\setcounter{footnote}{0}%
\opage{359}værdige. Ved mere fremskridende Kultur sættes aandelig Virksomhed
ved Siden af disse Idrætter. Der gør sig nu en Dualisme
gældende mellem en Kreds af Virksomheder, der vælges
for deres egen Skyld, og som er istand til at fylde og udvikle
det personlige Liv, og en anden Kreds af Virksomheder, som
paatvinges ved Magt eller ved Nød, og som gøre den enkelte
Personlighed til blot Middel for Andre. Det anses da som en
Nedværdigelse at arbejde for det Nødvendige. Gennem den
klassiske Oldtid og Middelalderen kan denne Tankegang paavises.
Det var dels det mekaniske Arbejdes ringe Sammenhæng med
Personligheden, dels den Afhængighed, som fremkommer, hvor
der arbejdes for Løn, der gjorde, at Haandværket ansaas for at
være en fri Mand uværdigt: „Intet Ædelt“, siger \textsc{Cicero}, „kan
trives i Værkstedet“ (nec enim qvicqvam ingenuum potest habere
officina)\footnote[1]{Cicero: De officiis. I, 150. --- Cicero følger her som sædvanligt Panaitios, der i dette Punkt var mere Aristokrat end tidligere og senere Stoikere. Smlgn. Bonhöffer: Die Ethik des Stoikers Epictet. Stuttgart 1894. p. 73 f. 233—243.}. Indenfor den stoiske Skole hævdedes Arbejdets Værdighed,
idet den Vises Sindelag adler enhver menneskelig Livsvirksomhed.
Og gennem Kristendommen udbredtes denne Tankegang
i store Kredse. Men det paa Aristokrati og Hierarki byggede
middelalderlige Samfund kom igen bort fra denne Opfattelse.
Endnu i Aaret 1781 udsatte Akademiet i Madrid den
Prisopgave: „at vise, at de nyttige Haandteringer ikke havde
noget Ærerørigt ved sig.“

Det er dels æstetiske og religiøse, dels etiske Betragtninger,
til hvilke denne længe herskende Opfattelse har støttet sig. ---
For Grækernes æstetiske og for Middelalderens religiøse Idealisme
var Arbejdet med det materielle Stof nedværdigende og
nedad dragende. Det hindrer Mennesket i at gøre sin Personlighed
til et harmonisk Kunstværk eller i at afdø fra den sanselige
Verden. --- Og idet den materielle Produktion først og
fremmest tjener til Individets egen Selvopholdelse, syntes den
at hvile paa Egoisme. Den moderne Nationaløkonomi har ofte
begunstiget denne Opfattelse ved stærkt at betone de egoistiske
Motiver til Erhvervsvirksomheden.

% --- p. 360 ---
\setcounter{footnote}{0}%
\opage{360}Naar der i den nyere Tid udvikler sig en større Anerkendelse
af det materielle Arbejde, saa hænger dette sammen med
Tanker, som ere ejendommelige for den moderne Livsanskuelse.
--- Denne er overbevist om den nøje Sammenhæng mellem det
aandelige og det materielle Liv. Hvorledes man end tænker sig
det nærmere Forhold mellem Aand og Materie, saa staar det
fast, at Aandens Liv er knyttet til den højeste og fineste Form
for materielt Liv. Et højt begavet Menneskes Hjerne er det
fineste materielle Produkt, vi kende. Og der behøves materielt
Arbejde til at frembringe saa gode Hjerner som muligt, og Arbejdet
for materiel Kultur kan derved ogsaa blive Arbejde for
aandelig Kultur. --- Det materielle Arbejde kommer dernæst
mere og mere til at staa som Anvendelse af Love, der ere opdagede
ved Tankens Arbejde. Vi bearbejde og beherske Naturen
i Kraft af den Overlegenhed over den, som Naturvidenskaben
giver os. Der er herved lagt en Bro mellem Haandens og Aandens
Verden, som Oldtid og Middelalder ikke kendte. Cicero
skelnede mellem Arbejde (opera) og Kunst (ars); hint var uværdigt,
denne værdig for en fri Mand. Men for mange Arter af
Arbejde falder denne skarpe Modsætning efterhaanden bort. Naar
Jagt, Krig og Tankearbejde stod som ædle Virksomheder, fordi
de personlige Kræfter frit kunde røre sig i dem, saa er nu den
Mulighed ikke udelukket, at ogsaa Arbejdet med de materielle
Stoffer kan udvikle Personligheden, jo mere det lægger Beslag
paa Forstand og Villie og ikke blot er en udefra paalagt mekanisk
Virksomhed. Endeligt have vi lært at se, at den Enkeltes
Arbejde, saa uanseligt det end kunde synes at være, dog spiller
sin Rolle i den store sociale Husholdning. Han skaffer sig derved
ikke blot Midlerne til sit eget Livs Opholdelse, men yder
ogsaa sin Skærv til Slægtens Liv, forøger de Goder, der staa til
Samfundets Raadighed. Hans Energi og hans Sparsommelighed
kunne --- formedelst Solidariteten i Individernes Kamp for Tilværelsen
--- komme Flere end ham selv til Gode. Han kan
derfor ogsaa udføre sin Gerning med et Sind, der er rettet mod
en større Horizont end den, som bestemmes ved hans rent individuelle
Fornødenheder. Videnskaben har lært os, at Naturen
baade paa det aandelige og det materielle Omraade ofte naar
% --- p. 361 ---
\setcounter{footnote}{0}%
\opage{361}store Resultater ved Opsummeren af smaa Virksomheder. Ved
at være tro i det Lidet kan den Enkelte arbejde for det
store Hele.

2. Men der er ved denne forandrede Betragtning af Arbejdet
langt snarere opstillet et Problem end naat et virkeligt
Resultat. Den højere Vurdering af det materielle Arbejde forudsætter
et ideelt Forhold mellem Arbejder og Arbejde, som
endnu langtfra er tilstede, men som er Indholdet af den Fordring,
der fra Etikens Side maa stilles til den sociale Udvikling.
Denne Fordring afgiver Maalestokken for Vurderingen af den
materielle Kulturs Udvikling til enhver given Tid. Før vi gaa
nærmere ind herpaa, ville vi dog først fremdrage de sociale Modsætninger,
der vidne om, at den primitive Arbejdsdeling og dens
Misligheder endnu ikke ere overvundne.

a. Modsætningen mellem besiddende og blot arbejdende Individer
er en historisk Fortsættelse af den primitive Arbejdsdeling
mellem Herre og Træl. Den Besiddende nyder ofte kun Frugten
af Fortidens Arbejde eller af Samtidiges Arbejde; i saa Fald
er der en fuldstændig Modsætning mellem ham og Arbejderen.
Men selvom han arbejder, er Modsætningen skarp nok. Han
raader nemlig baade over Arbejdsmidlerne og over det færdige
Produkt. Den besiddelsesløse Arbejder derimod, som kun har
sine to Hænder, maa stille dem til den Besiddendes Raadighed
for at faa Stof at bearbejde og Andel i det Udbytte, som vindes.
Uden Stof kan intet Arbejde udføres; men selvom Arbejdskraft
ikke kan faas, kan Stoffet dog maaske nydes og forbruges af
Besidderen. Saa uundværligt et Mellemled Arbejdet end er ved
Overgangen fra det raa Stof til Produktet, saa betragtes det dog
fra Produktionens og Ombytningens Synspunkt blot som et Middel.
Arbejdslønnen betragtes væsentligt som en Udgift, det gælder
at reducere til det mindst Mulige. Fra sit Standpunkt maa
Arbejderen naturligvis betragte Arbejdslønnen som et Formaal.
Dens Størrelse betegner hans sociale Stilling, og paa en Maade
ogsaa hans Stilling som Menneske. Men denne Betragtning savner
han Magt til at gøre gældende saalænge han er overladt til sig
selv. Ligesom hans Arbejde væsentligt betragtes som blot Middel
for Produktionsprocessen, saaledes bliver hans Person ogsaa
% --- p. 362 ---
\setcounter{footnote}{0}%
\opage{362}let betragtet som blot Middel. En talrig Arbejderbefolkning er
nødvendig, for at Arbejdet skal kunne faas til saa billig Pris som
muligt. For at Nogle skulle lade sig nøje med den mindst mulige
Løn, maa der være Andre, som slet ingen Løn faa. Den konsekvente
Gennemførelse af den Betragtning, at Arbejdslønnen blot
er Udgift, fører altsaa til, at en Mængde menneskelige Væsener
trykkes ned til det laveste Trin af Existens.

Denne Afhængighed, i hvilken Arbejderen, trods den personlige
Frihed, han officielt besidder, kommer til at staa overfor
den Besiddende, strækker sig videre end til Vilkaarene for
hans Arbejde og Underhold. Ikke blot Lønnen, men ogsaa Arbejdstiden
og de Forhold, under hvilke Arbejdet udføres, bestemmes
overvejende uden hans Villie. Det lyder som en Haan,
naar det siges ham, at han er en fri Mand, og at det beror paa
ham selv, om han vil overtage et Arbejde eller ikke\footnote[1]{Da der i en engelsk Regeringskommission blev stillet Fordring om nøjagtigere Tilsyn med Kulgruberne paa Grund af de hyppige Ulykkestilfælde, sagde Arbejdsgivernes Repræsentant: „Beror det da ikke paa Minearbejderne selv, om de ville stige ned i Gruberne ?“ Et Vidne svarede: „Ganske vist. Men det beror ogsaa paa dem selv at suite ihjel, naar de ikke stige ned.“ L. Brentano: Die Arbeitergilden der Gegenwart. Il, p. 17. --- Smlgn. p. 165 om en Arbejdsgivers Glæde over de mange ubeskæftigede Arbejdere, der gøre Lønnen billig.}. Den
Vare, han har at udbyde --- hans Arbejdskraft --- er uadskilleligt
knyttet til hans Person, saa at han ved at sælge den paa
Vilkaar, han ikke raader over, ogsaa sælger sin Person. Afhængigheden
af Arbejdsgiverne strækker sig let fra det ydre,
materielle Omraade til det indre, saa at han heller ikke overfor
sociale, politiske og religiøse Spørgsmaal kan optræde som en
fri Mand.

b. De færreste Arbejdere have endnu kunnet nyde godt
af, hvad Videnskaben har udrettet i Industriens Tjeneste. Tankens
og Musklens Arbejde, det intelligente og det fysiske Arbejde
ere endnu ikke forenede, men ere fordetmeste forskellige
Individers Sag. Anvendelsen af Videnskabens Opdagelser, Anlægget
af Arbejdsplanen og Tilvejebringelsen af Midlerne til at
% --- p. 363 ---
\setcounter{footnote}{0}%
\opage{363}udføre den kræve Evner og Betingelser, som de fysiske Arbejdere
ikke raade over. Det fysiske Arbejde udøves saa godt
som i Blinde, uden Indsigt i de Naturkræfter, det virker med,
eller i Betydningen af de Resultater, der opnaas. --- Ogsaa
indenfor de fysiske Arbejderes egen Kreds gennemføres mere og
mere en Arbejdsdeling, som tilsidst indskrænker den enkelte
Arbejder til mekanisk Gentagelse af en og samme ganske simple
Virksomhed. Der foregaar ved Arbejdets Udvikling dels en Opløsning
og Udstykning af et sammensat Arbejde til flere forskellige
Arbejdere, dels en Kombination af flere forskellige Arbejdsgrene
til Frembringelse af en og samme Genstand. Det
er de samme to Udviklingsformer --- Isolation og Kombination
--- som fremtræde paa alle Omraader for menneskelig Virksomhed.
Fabrikationen af en Knappenaal var allerede paa Adam
Smiths Tid fordelt paa 18 Hænder. I Aaret 1860 var der kun
fire Arbejdere til at lave en Stol; men Aaret 1895 behøvedes
der 23 Arbejdere til at lave samme Slags Stol ved Hjælp af
Maskiner. Ved saadant Samarbejde, hvad enten det opstaar ved
Udstykning eller Kombination, mister den enkelte Arbejder Overblikket.
Han indøves i en enkelt simpel Virksomhed, og kun
en ganske enkelt Side af hans Personlighed og hans Evner udvikles.
Det menneskelige Liv udvikles ikke hos ham i dets fulde
Omfang; han kommer ogsaa fra denne Side set til at staa som
Middel, ikke som Formaal. Jo mere mekanisk han kan øve sin
Gerning, des bedre griber den maaske ind i den hele Virksomhed.
Hans Ensidighed gør ham til et saa meget des bedre Led
i det hele Maskineri. Og Ensidigheden gør ham ogsaa mere afhængig,
da han ikke er istand til at paatage sig nyt, mere
sammensat Arbejde. Ganske vist gør Udstykningen det let at
gaa over fra et simpelt Arbejde til et andet; men denne Fordel
opvejes ved, at saa godt som Alle kunne paatage sig saadanne
simple Virksomheder, saa at Konkurrencen bliver stor
og Lønnen ringe. --- Det bliver da kun til liden Gavn for Arbejderen,
om han uddanner sin Aand og sine Evner, da der til
den Virksomhed, der falder i hans Lod, ikke behøves nogen
videre Uddannelse. Uden Glæde ved sin Virksomhed, uden Haab
% --- p. 364 ---
\setcounter{footnote}{0}%
\opage{364}om Fremgang og uden Drift til Udvikling synker han let ned
til en halvt dyrisk Existens. ---

De Modsætninger, som her ere skildrede i deres stærkeste
Former\footnote[1]{Rousseau har først med Energi fremdraget de uheldige Sider ved Arbejdets Deling og derved stillet det sociale Spørgsmaal. Smlgn. min Bog J. J. Rousseau og hans Filosofi. p. 99—106; 125 f.; 130 f. --- Kort efter Rousseau belyste Adam Ferguson (Essay on the History of Civil Society. Edinbourgh 1767. IV, 2: Of the subordination consequent to the separation of art and professions) det samme Spørgsmaal fra et kulturhistorisk Synspunkt. --- I sit paa grundige sociologiske og psykologiske Studier byggede Skrift Gemeinschaft und Gesellschaft (Leipzig 1887, 2. Udg. 1912) har Ferdinand Tönnies givet en lærerig Fremstilling af den sociale Patologi. Smlgn. min Anmeldelse af dette Skrift i Afhandlingen Social Pessimisme i „Mindre Arbejder“ I.}, betinge det sociale Spørgsmaal.

\chapterhead{XXV.}{DET SOCIALE SPØRGSMAAL.}

1. Naar man taler om det sociale Spørgsmaal, gælder det
at holde den Misforstaaelse borte, at det er et aldeles enkelt og
simpelt Spørgsmaal, og at det skulde kunne løses paa én Gang
og én Gang for alle. Det opstaar ved de stærke Modsætninger,
der true med at splitte Samfundet; men disse Modsætninger betinges
igen ved mange forskellige Forholds Vexelvirkning. Det
er ikke blot det materielle Arbejde, hvis Ordning det her gælder;
ogsaa den ideelle Kultur, det aandelige Livs Udvikling
bliver af Betydning; og endeligt vil ogsaa Statens Forfatning
faa væsentlig Indflydelse baade paa Spørgsmaalets Skikkelse og
paa dets Løsning. Der er en ligesaa nøje Forbindelse mellem
disse forskellige Sider af det sociale Spørgsmaal, som der er mellem
Hjerne, Hjerte og Tarm. Sygdom i et enkelt af disse Or%
% --- p. 365 ---
\setcounter{footnote}{0}%
\opage{365}ganer kan volde det hele Legemes Død, og en Fejl i det ene
Organ vil faa Følger for de andres Virksomhed, der saa igen
virke tilbage paa hint. --- En Løsning af Spørgsmaalet én Gang
for alle er usandsynlig, da det hænger sammen med saa mange
Forhold, og da det, selvom det træder tydeligst og skarpest frem
i vor Tid, dog under en eller anden Form har bestaaet til alle Tider.

Grunden til, at det først er i vore Dage, man taler om det
sociale Spørgsmaal, maa søges i flere Omstændigheder. Det er
neppe, fordi der er større Egoisme og mere Misundelse i Verden,
end der før har været. Det er heller ikke, fordi Nøden
og Lidelsen i og for sig skulde være større nutildags end
forhen; snarere turde det Modsatte paastaas. Man kan maaske
endog sige, at dersom Forholdene ikke havde forbedret sig, vilde
man ikke tale om et socialt Spørgsmaal. Ti den største Nød
slaar til Jorden, kuer Tanken og vækker ikke Driften til at
komme videre. Det er forsaavidt netop en Fremgang i Arbejdernes
Kaar, der gør det muligt, at det sociale Spørgsmaal kan
stilles af dem selv. Det er et Tegn paa, at deres Natur alligevel
ikke er fuldstændigt kuet. Kun forsaavidt maa Forholdene siges
at have forværret sig, som den Ensidighed og Afhængighed, den
stedse videre gennemførte Arbejdsdeling bevirker, maa føles saa
meget des stærkere, jo mere Arbejderen i Kraft af det 18. Aarhundredes
Principer nu staar som fri Mand. Saalænge han kun
var Slave eller Tjener, følte han ikke saa meget disse Misligheder,
især da den gamle Samfundsorden ved Siden af den Afhængighed,
hvori den holdt ham, paa mange Maader ydede ham
Beskyttelse og Hjælp. Dertil kommer endnu den forøgede Oplysning.
Tænkeevnen er vakt, og Sammenligninger anstilles. Man
bøjer sig ikke længere overfor den givne Samfundsorden og
Livskaarenes Forskellighed som overfor Noget, der skulde være
berettiget, blot fordi det er. Man spørger om Berettigelsen, og
man opstiller Idealer. Og jo større Modsætningen mellem Ideal
og Virkelighed er, des stærkere sættes Fantasien i Bevægelse.

Men det er ikke blot den vaagnende Sammenligning og
Eftertanke, der stille det sociale Spørgsmaal saa skarpt. Det
stilles jo ikke blot af dem, der umiddelbart lide under Forholdene;
ogsaa Andre føle dets Brod. Medfølelsen og Retsfølelsen
% --- p. 366 ---
\setcounter{footnote}{0}%
\opage{366}have udviklet sig til større Finhed og større Omfang. Den Uro,
man klager over i Nutiden, kommer for en stor Del af Sympati.
Kunde man være mere egoistisk og tankeløs, saa vilde
man for sin egen Person leve lykkeligere. Den Pessimisme, som
er saa almindelig i vor Tid, udspringer vel for en Del af Blaserethed;
men dens ædleste Kilde er den levende Medfølelse
med den megen Lidelse og Disharmoni, som findes i Verden,
og som tidligere Tider gik mere let hen over.

2. Samtidigt med, at Bevidstheden om det sociale Spørgsmaals
Betydning er bleven skærpet, og Driften til at arbejde paa
dets Løsning er bleven vakt, have vi tillige lært, hvor dyb en
Grund det har i den menneskelige Natur og dens Vilkaar. Den
barnlige Historiens Filosofi, der udleder alle store Bevægelser
af enkelte Individers Optræden, ere vi ved at komme ud over;
den strander allerede paa, at det staar som ubegribeligt, hvorledes
de enkelte Individer selv ere faldne paa at tænke og
handle, som de gjorde. Baade hos Konservative og Radikale kan
dog denne barnlige Opfattelse endnu paavises. Efter den konservative
Optimisme vilde Alt være godt, naar der kun ingen
radikale Agitatorer fandtes; efter den radikale Optimisme vilde
Alt være godt, naar der kun ingen Konger og Præster fandtes.
Begge overse dels, at vi i det sociale Spørgsmaal have en historisk
Arv, en Fortsættelse af tidligere Tiders sociale Vanskeligheder,
dels at der i den menneskelige Natur og de menneskelige
Livsvilkaar ligger Aarsager, der maaske stedse ville føre til,
at det sociale Spørgsmaal stilles, selvom det bliver i nye Former.

Det blev i det Foregaaende omtalt, at Arbejderbefolkningens
store Talrighed var en Aarsag til, at Arbejdet indtog en
forholdsvis saa underordnet Plads blandt Produktionens Faktorer.
Hvis der ikke var saa mange Arbejdere, vilde Arbejdernes
Afhængighed af dem, der raade over Stof og Produkt, ikke
blive saa stor. Det er Befolkningens stærke Tiltagen, som volder
Kampen for Tilværelsen paa det sociale Omraade ligesom i
hele Naturen. Paa alle det organiske Livs Omraader opholdes
Arterne ved, at der kastes langt flere Spirer ud, end der under
de givne Forhold kan udvikle sig. For at nogle faa Spirer
kunne slaa Rod og udvikle sig, maa mangfoldige forgaa. Det
% --- p. 367 ---
\setcounter{footnote}{0}%
\opage{367}synes, som om Livet kun kan bestaa, naar en vældig Kraft sættes
ind paa at frembringe et Overskud af Individer. Ogsaa Menneskeslægten
har en Tendens til at forplante sig hurtigere, end
Næringsmidlerne altid kunne følge med. Denne Sætning, som
med Overdrivelse og Ensidighed blev fremsat af Malthus (Essay
on population. 1798) mod den herskende Optimisme, der
udledte alle sociale Misligheder af de menneskelige Indretninger,
har henledet Opmærksomheden paa en væsentlig Side ved det
sociale Spørgsmaal. Samfundslivet viste sig nu at hænge sammen
med Naturlivet idethele. \textsc{Darwin} genfandt i hele Plante- og
Dyrelivet de Forhold, Malthus havde antaget for Menneskelivets
Vedkommende. Vi forstaa ikke det sociale Spørgsmaal,
naar vi ikke se det i dets Sammenhæng med hele Naturlivet.
Naar det er saa dybt liggende Kræfter og Tendenser, der virke
med, vil man ikke saa let mene, at Spørgsmaalet skulde kunne
løses ad én eneste Vej, ved én enkelt Formel og for alle Tider.

Malthus' Ensidighed bestod for det Første --- som allerede
Auguste Comte bemærkede\footnote[1]{Cours de philosophie positive. IV. 2. éd. p. 455.} —i, at han ikke fremhævede,
at en talrig Befolkning er nødvendig, for at det sociale Liv skal
kunne udfolde sig i sine forskellige Former. Arbejdsdeling, livlig
Ombytning, social Organisation ere kun mulige ved talrig
Befolkning. Den talrigere Befolkning er tillige en tættere Befolkning,
og derved fremmes Samlivet og Alt, hvortil det igen
fører. Ogsaa overfor ydre Farer kan Talrigheden være af Vigtighed.
--- For det Andet bestod Malthus' Ensidighed i, at han
ikke saa, at Misforholdet mellem Befolkningens Størrelse og de
til Raadighed staaende Næringsmidler ogsaa kan skyldes en
ufuldkommen Udvikling af den menneskelige Virksomhed. Større
Energi eller Kløgt vilde maaske kunne finde Anvendelse for
den tilsyneladende overflødige Arbejdskraft, enten ved Hjælp
af nye Opfindelser eller ved en bedre Brug af Kapitalen. Eller
nye og lettere tilgængelige Næringsmidler kunne bringes tilveje.
Misforholdet kan altsaa i de enkelte Tilfælde have forskellige
Aarsager, snart i en Mangel ved Produktionen, snart i
en Mangel ved Fordelingen, snart i en Mangel ved Forbruget.
% --- p. 368 ---
\setcounter{footnote}{0}%
\opage{368}--- Men hvad Malthus har Ret i, er, at hver Gang der er frembragt
en vis Harmoni mellem Befolkning og Næringsmidler,
vil Befolkningens stærke Formeringstendens stedse paany ophæve
Ligevægten og gøre nye Anstrengelser nødvendige for at
bevare det Livstrin, som er naat. Det er dette, der stedse paany
driver Mennesker ind paa Arbejdets og Kulturens Vej. Malthus
siger, og man behøver ikke at være Pessimist for at give ham
Ret deri, at hvis Næringsmidlerne af sig selv forøgedes i samme
Grad som Menneskeslægten, indser han ikke, hvilken Drivfjeder
der kunde være stærk nok til at overvinde Menneskets Indolens
og føre det videre paa Kulturens Vej\footnote[1]{Versuch iiber die Bedingungen und die Folgen der Volksvermehrung. Aus dem Engl. von Hegewisch. Altona 1807. Il. p. 152.}. Trægheden, i hvilken
vi have fundet Spiren til alt Ondt (VI, 1), behøver en stærk
Modvægt for at overvindes. Saalænge ikke højere Motiver have
større Magt i den menneskelige Natur, end hidtil er Tilfældet,
vil det af Malthus fremhævede elementære Motiv stedse paany
være nødvendigt. Der behøves et stærkt Pres til at vække
Driften og Kraften. Først naar det har gjort sin Virkning, naas
der et Trin, hvor mere ideelle Motiver kunne virke.

Der er ingen Grund til at føle Beundring for den Naturorden,
vi her staa overfor. Men den maa tages som noget Givet.
Og det maa paa den anden Side ikke glemmes, at den vældige
Kraft, som fører til at stille det sociale Spørgsmaal, er den samme,
der ved at lede til Familiens Stiftelse grundlægger det inderligste
menneskelige Samfund, Kilden til al Sympati i Menneskeslægten,
og den samme, der ved sin Indflydelse paa Fantasien
og ved at vække Entusiasmen sætter Evnen til Hengivelse og
Opofrelse i Virksomhed.

3. Det sociale Spørgsmaal er et etisk Spørgsmaal. Dette
ligger allerede i, hvad der ovenfor er sagt, at kun formedelst
den vakte Sympati fremtræder det i nyere Tid saa skarpt og
saa tilspidset. Tænk Etikens Maalestok, den etiske Tanke om et
idealt Samfund (III, 10; XIII, 6) borte --- og „Spørgsmaalet“
vil da med det samme falde bort, eller rettere, det vil være
reduceret til et rent Magtspørgsmaal; der vil blot stilles den
% --- p. 369 ---
\setcounter{footnote}{0}%
\opage{369}Enkelte den Opgave at hytte sit Skind eller mele sin Kage saa
godt som muligt under den almindelige Kamp af Alle mod Alle.
I Tanken om det ideale Samfund, om Humanitetsriget ligger
derimod Fordringen om, at ethvert menneskeligt Væsen skal
være Mere end Middel, skal indtage sin ejendommelige og selvstændige
Plads i det store Menneskerige\footnote[1]{Tanke og dens Betydning for det sociale Spørgsmaal.}. Det strider mod
Idealet af et menneskeligt Samfund, naar et større eller mindre
Antal menneskelige Væsener blot staar som passiv Masse, som
underordnede Midler og Redskaber, hvis Nydelse og Lidelse
ikke tages med i Beregning, naar Regnskabet for socialt Fremskridt
eller Tilbagegang skal gøres op. Det maa herved vel
erindres, at Nydelsen eller Lidelsen ikke blot betinges ved de
Goder eller Onder, der udefra falde i de Enkeltes Lod, men
ogsaa ved de Opgaver, de kunne stille sig og arbejde paa, og
de Pligter, de formaa at overtage. Det kommer ikke blot an
paa Udbyttets, men ogsaa paa selve Arbejdets Fordeling. Proletarens
Stræben efter at opnaa et højere Livsniveau finder meget
ofte sin Berettigelse ved, at han kun, naar et saadant Niveau
er naat, kan komme til at arbejde paa de Opgaver og fyldestgøre
de Pligter, der svare til hans Evner. Saaledes fremtraadte
ogsaa Kvindens Ønske om „Emancipation“ ofte som begrundet
ved Trangen til aktiv Deltagelse i Menneskehedens store Værk.
(Smlgn. XIX, 5).

Saalænge Begrebet Masse endnu kan anvendes paa menneskelige
Væsener, saalænge er Maalet ikke naat og mere eller
mindre betænkelige Disharmonier ere tilstede. Masse er indenfor
et Samfund det, der ikke er optaget med i Organisationen,
ikke er nogen levende og virkende Bestanddel af Samfundet. I
Massen udviskes og forsvinde de enkelte Personligheder og staa
ikke som ejendommelige Kraftpunkter. Hvor der er Masse, kommer
Personligheden ikke til sin Ret. Hverken den Enhed eller
den Mangfoldighed, som et fuldkomment Samfund forudsætter,
kan derfor være tilstede: Enheden ikke, fordi Samfundet spal-

) I min Afhandling Forestillingen om Retfærdighed i dens historiske
Hovedformer (Nationaløkonomisk Tidsskrift 1873) har jeg udviklet denne
% --- p. 370 ---
\setcounter{footnote}{0}%
\opage{370}tes i Dele, der ikke hænge organisk sammen, men maaske
endog staa fjendtlige overfor hverandre; --- Mangfoldigheden ikke,
fordi de enkelte Personligheders Selvstændighed forsvinder. Den
frie Personligheds Princip, den vigtigste Sætning, der kan udledes
af Velfærdsprincipet (VIII, 6), krænkes, idet personlige
Væsener nedsættes til at være Midler uden tillige at betragtes
som Formaal.

4. Men kan det nu være anderledes? --- Der gives dem,
som besvare dette Spørgsmaal med Nej. Deres Tankegang er
omtrent følgende. --- Der hersker i Samfundslivet som i Naturlivet
bestemte Love, som ikke kunne trodses. En saadan Lov
er Loven om Tilbud og Efterspørgsel. Det nytter ikke at opbyde
den største Anstrengelse, naar det Resultat, man opnaar,
ikke er efterspurgt; vi faa for vort Arbejde kun, hvad der kan
gives paa Markedet. Under denne Lov maa særligt de komme
til at lide, der møde paa Markedet med den blotte fysiske Arbejdskraft.
De ville ikke kunne naa en betryggende materiel
Existens, og de ville ikke have Tid og Kraft tilovers til at blive
delagtige i den aandelige Kultur. Dette kan ikke være anderledes,
og det har aldrig været anderledes. Paa den anden Side
er det nødvendigt, at det materielle Arbejde udøves. Samfundets
materielle Fornødenheder maa skaffes tilveje af den store
Masse, for at en lille Skare kan arbejde paa den aandelige Kultur.
Den højere Kultur har til alle Tider kun været tilgængelig
for en lille Kreds, der har været begunstiget af Lykken og
har kunnet udvikle sig harmonisk og alsidigt uden at behøve at
trælle for det materielle Behov. For at Lyset kan falde klart
over denne Menneskeslægtens Elite, maa den store Masse ligge
i Mørket. Det vil ikke engang være heldigt, om der spredes
for meget Lys fra den snevre Kreds ud i den store. Det er en
naturlig Arbejdsdeling, at Nogle tænke, Andre arbejde, og Tænkning
vil let gøre Arbejdet mindre godt eller vil vække Utilfredshed
og Uro og derved forstyrre den sunde Orden. Alligevel give
de Udvalgte Mængden Mere, end de modtage fra den, saa man
kan ikke tale om Hjerteløshed og Umenneskelighed hos hine. Og
forsaavidt de, der høre til den store Masse, ikke finde det Liv,
de ere henviste til at føre, værdt at leve, maa de henvises til
% --- p. 371 ---
\setcounter{footnote}{0}%
\opage{371}Kirketroen og trøste sig ved Haabet om et andet og bedre
Liv, i hvilket de kunne finde Oprejsning for, hvad de her
have lidt.

Som Modstykke bliver ofte følgende Betragtningsmaade fremført.
--- De saakaldte sociale Naturlove, man beraaber sig paa,
ere kun Udtryk for en Samfundsorden, som tildels endog er bragt
istand ved Vold og Overgreb. At nogle Individer have saa
stort Forspring i den store Konkurrence, skyldes en Uretfærdighed,
det maa være muligt at korrigere. Med hvilken Ret
taler man egentligt her om en Samfundsorden? I Konkurrencen
mellem de menneskelige Individer hersker egentligt kun
den dyriske Kamp for Tilværelsen, blot i en noget maskeret
Form. Det nuværende Samfund er aldeles ikke bygget paa
etiske og humane Principer, men paa Magt og Egoisme. Det
maa aldeles opløses. Der maa skaffes bar Grund, saa kunne
vi altid tale om, hvad der skal bygges i Stedet. Medens det
materielle Arbejde nu staar som en underordnet, tjenende Kraft,
maa den Indsigt trænge igennem, at Arbejdet er Kilde til al
Rigdom og Kultur. «Overfor Arbejderklassen,“ hedder det i et
af tyske Socialdemokrater (i Gotha Maj 1875) udstedt Program
„ere alle andre Klasser kun en reaktionær Masse“.

Som man raaber i Skoven, faar man Svar. Fra begge
Sider raaber man „Masse!“ til Modstanderen. Fra begge Sider
mødes man i den Paastand, at der bestaar en skarp Modsætning,
en Dualisme indenfor Samfundet, kun at man fra den ene Side
anser den for gavnlig eller i hvert Tilfælde nødvendig, fra den
anden Side for umenneskelig og for let at ophæve, naar blot
Villien er god.

5. Efter den almindelige Opfattelse, jeg tidligere har gjort
gældende (III, 19---20; XIII, 1---4), skal den etiske Udvikling
føre det videre, som den naturlige, mere eller mindre ubevidste
Udvikling har indledet. Den menneskelige Villie kan ikke skabe
af Intet, men kan kun bygge videre paa den Grundvold, Naturen
har lagt. Det gælder altsaa at finde Udgangspunkter og
Fremtidsspirer i det Givne. Hvor saadanne aldeles mangle, er
Tilstanden haabløs. Men efter begge de nys fremdragne Opfattelser
ende vi i Haabløshed. For den først omtalte Opfattelses
% --- p. 372 ---
\setcounter{footnote}{0}%
\opage{372}Vedkommende er dette konsekvent; ti da den ansér Dualisme
i Menneskeslægten for en normal Tilstand, er der for den ingen
Grund til at hige ud over den. Overfor den anden Opfattelse
derimod maa der stilles det Spørgsmaal: dersom der i det nuværende
Samfund aldeles ikke findes Noget af Værdi, hvorfra
faas da Kræfterne til at arbejde paa den nye Samfundsbygning?
De Kræfter, vi bruge, har Fortiden stedse bragt tilveje. Det
er derfor ikke let at se, hvorledes en absolut Fordømmelse af
det historisk Overleverede kan forenes med et stort Haab til
Fremtiden.

Der er desuden en Ting, som begge Parter overse. Dualismen
er i Virkeligheden ikke fuldstændig. Naar det sociale
Spørgsmaal i vor Tid er stillet skarpere end før, er det, efter
hvad vi have set, for en stor Del, fordi Bevidstheden om dets
Betydning er vaagnet. Denne Bevidsthed er ikke indskrænket
til nogen enkelt Stand og Kreds. Det er, som \textsc{Stuart} \textsc{Mill}
har sagt, det Ejendommelige for vor Tid, at alle Klasser tage
Del i Debatten om de vigtige Spørgsmaal, og at mere end nogensinde
før de lidende Klasser have en vigtig Stemme med\footnote[1]{History of Trade Unionism. London 1894. p. 197, 228 f., 246 f.), haves ogsaa et Vidnesbyrd om, at den arbejdende Klasse ikke fører sin Kamp uden Hjælp og Sympati fra andre Klasser, Noget som Marxisternes stadige Tale om Klassekampen lægger Skjul paa. Selve de socialistiske Ideer ere desuden først fremsatte af Mænd fra de „højere“ Klasser. --- At Positivisterne naturligt arbejde sammen med Proletarerne, ses af det i Den nyere Filosofis Historie \textbackslash{} II. p. 350 og 359 Anførte.}. De
to modstaaende Dele af Samfundet ere altsaa dog ikke aldeles
isolerede, men maa have en fælles Grund, siden de kunne drøfte
Problemet med hinanden. Herved er Samfundets aandelige Enhed
hævdet, og tillige et væsentligt Middel til at komme videre
bragt tilveje. Ti Historien viser klart, at ingen Reform og
ingen Forbedring fører til Maalet, naar den ikke paa aktiv Maade

») Principles of political Economy. III, 1, 2. --- Lassalle hævdede,
at hans Agitation indledede et Forsoningsværk, idet den gav den besiddende
og den arbejdende Klasse Mulighed for at mødes i Forhandlingen om
det sociale Spørgsmaal. (Arbeiterlesebuch. p. 54 f.). --- I den Rolle,
Positivister og „kristelige Socialister“ have spillet under de engelske
Arbejderes Kamp for deres Stands Ret (se Sidney and Beatrice Webb:
% --- p. 373 ---
\setcounter{footnote}{0}%
\opage{373}støttes af dem, der skulle hjælpes. De véd bedst, hvor Skoen
trykker, og uden deres Medvirken kan Trykket ikke afhjælpes.
Hverken teoretisk eller praktisk kan det sociale Spørgsmaal behandles
paa frugtbar Maade, uden at Arbejderne selv tage Del
i Drøftelsen og Afgørelsen. Og hertil er nu Begyndelsen gjort. ---

Men der er en lang Vej fra Drøftelse til Afgørelse, og den
Trøst, som blot hentes fra, at Alle kunne være med at diskutere,
kunde med Rette synes noget betænkelig. Vi maa da undersøge,
hvilke praktiske Udviklingsmuligheder der frembyde sig.

\chapterhead{XXVI.}{UDVIKLINGSMULIGHEDER.}

1. Hvad der gør en Menneskemængde til en blot Masse,
er, at de enkelte Personligheder ikke faa selvstændig Betydning,
ikke hver for sig udvikle sig paa ejendommelig Maade. Men
med denne Udvisken af de personlige Ejendommeligheder følger
ogsaa en Mangel paa indre Forbindelse, paa organiseret Samfundsliv.
Massens enkelte Dele staa ligegyldige overfor hverandre.
Istedetfor et virkeligt Samfund, i hvilket Individerne ere
forbundne ved fælles Formaal, faa vi da et blot Selskab\footnote[1]{Denne Modsætning mellem Begreberne Samfund (Gemeinschaft) og Selskab (Gesellschaft) er udviklet af Tönnies i hans ovenfor (XXIV, 2) nævnte Skrift.}, hvor
det ene Individ betragter det andet som blot Middel. Ydre Interesser
forbinde Menneskene, uden at de derfor staa i nogen
væsentlig indre Forbindelse. Den Ene vil med det mindst mulige
Offer fra sin Side opnaa det mest Mulige af det for ham Attraaværdige,
som den Anden besidder. Naar den ene yder den
Anden Midler til hans Formaal, er det kun for at faa Midler til
% --- p. 374 ---
\setcounter{footnote}{0}%
\opage{374}sine Formaal i Stedet. Forholdet mellem Individerne antager
stedse mere en upersonlig Karakter; ti det Upersonlige benytte
og betragte vi blot som Middel for vore egne Formaal. Hver
Enkelt følger sine Instinkter og sine Formaal. Men idet hver
Enkelt følger sin Interesse, uden at der bestaar noget virkeligt
Fællesskab med Andres Interesser, hvor de ikke kunne gøres
til Midler for ham, kommer det nødvendigvis til Sammenstød.
Hver Enkelt vil leve uden at indse Nødvendigheden af, at Andre
ogsaa leve. Det er især paa de materielle Goders Omraade, at
Sammenstødene maa komme. Et materielt Gode kan kun tilhøre
en Eneste. Hvad der mætter mig, mætter ikke min Nabo.
Jeg maa skyde hans Selvopholdelsesinstinkt til Side for at kunne
tilfredsstille mit eget, og jo flere materielle Goder jeg tilegner
mig, des færre maa der blive til ham. Et og samme Atom kan
ikke samtidigt være Bestanddel af min Organisme og af hans.
Derfor opstaar her let en Kamp af Alle mod Alle, og den Opgave
stiller sig at fordele Goderne saaledes, som Alles Velfærd
kræver. Den fordelende Retfærdigheds (III, 9) Fordring er en
Fordring om en social Organisation, ved hvilken Individerne
kunne supplere hverandre istedetfor at bekæmpe hverandre.
Det store Spørgsmaal er, ad hvilken Vej en saadan skal tilvejebringes,
og hvilke Muligheder til en saadan Erfaringen anviser
os. Hvilke Kræfter raader Menneskeslægten over til Løsning
af dette Problem?

Dette Spørgsmaal drøfte vi her, i den sociale Etik, fra en
noget anden Side end i Nationaløkonomien. Vi have at anstille
en etisk Vurdering af de forskellige Maader, paa hvilke en social
Fordeling og Organisation kan frembringes. Nationaløkonomien
holder sig derimod mere til en Undersøgelse af, hvorvidt Mulighederne
virkeligt ere tilstede; den drøfter mere den tekniske
Side ved Sagen. Dog kan der, som allerede tidligere bemærket
(III, 14), ikke drages nogen skarp Grænse mellem Etik og Nationaløkonomi,
idet denne sidste ikke blot er en Lære om Frembringelse
og Ombytning af materielle Goder, men ogsaa en
Lære om deres Fordeling.

Der vil nu, som Erfaringen viser, være to Grupper af
Kræfter, der kunne virke i Retning af en med Velfærdsprin%
% --- p. 375 ---
\setcounter{footnote}{0}%
\opage{375}cipet stemmende Fordeling, og derved tillige i Retning af social
Organisation paa den materielle Kulturs Omraade. Dels kunne
nemlig de individuelle Kræfter virke sammen i frie Associationer,
fremkaldte ved Interessefællesskab og Sympati, dels kan
Staten med sin tvingende Centralmagt gribe ind og fastslaa en
Ordning af det materielle Kulturarbejde. Vi behandle disse to
Grupper hver for sig. De frembyde begge igen forskellige Former.

% CHECK p. 375: centred line: heading or verse
\begin{center}a. Arbejdets Ordning gennem fri Association.\end{center}

2. Frihedsprincipets første Virkning var opløsende og isolerende.
Dets Opstilling virkede som en Mine, der sprængte
den gamle Samfundsbygning, og Verden er endnu ikke kommen
til Ro efter denne Katastrofe. Som vi have set, betragtede man
endog i revolutionær Iver Isolation som en nødvendig Betingelse
for Frihedens Bevarelse. Men Friheden lægger sig ligesaa vel
for Dagen ved Valget af Forbindelse med Andre som ved Valget
af Sondring fra Andre. Den betyder kun det, at der ingen
vilkaarlige Hindringer stilles for Individets Virksomhed. Den
frie Association maa netop siges at være den naturlige Fortsættelse
af Emancipationen. Hvis Friheden har medført Ulykker
og Misligheder, saa er den --- for en Del idetmindste --- istand
til selv at læge de Saar, den har slaaet. Frihedens Proklamation
sker i Kraft af Velfærdsprincipet, fordi den frie Brug af
Kræfterne baade har Værdi som Formaal og som Middel, og
det er ligeledes i Kraft af Velfærdsprincipet, at Emancipationen
fortsættes af Associationen. Dette vil træde klart frem, naar
vi betragte de Hovedformer for Arbejderassociationer, der historisk
ere fremtraadte.

3. Allerførst gælder det at ophæve Isolationen og Striden
indenfor Arbejdernes egen Kreds. Arbejderens Afhængighed overfor
den, der køber hans Arbejde, hænger jo især sammen med
den stærke Konkurrence mellem Arbejderne selv. Arbejdslønnen
synker formedelst det stærke Tilbud om Arbejde. Ved at
optræde i Fællesskab overfor Arbejdskøberne blive Arbejderne
istand til at stille deres Betingelser, baade hvad Arbejdslønnen,
Arbejdstiden og Arbejdsvilkaarene angaar. Først derved blive
% --- p. 376 ---
\setcounter{footnote}{0}%
\opage{376}de virkeligt frie Mænd, idet de have Mulighed for at faa deres
Fordringer anerkendte. Saalænge Friheden ikke er forbunden
med den mindste Magt, er den kun et Ord; og Magten naas
kun ved Sammenslutning og Organisation. Som saa ofte i Historien
maa Magten udfoldes, for at Retten kan blive anerkendt.
Længe efter at Trældom og Livegenskab vare ophørte, og efter at
Lavene vare ophævede, betragtede Arbejdsgiverne sig som absolute
Herrer og Autoriteter overfor deres Arbejdere og ansaa
den ubetingede Myndighed over dem for en Grundbetingelse
for industriel Organisation. De gjorde Fordring paa det gamle
Regimentes Rettigheder, efter at de tilsvarende Pligter vare faldne
bort. Og ikke blot dette: de beraabte sig samtidigt paa det
Frihedsprincip, Revolutionen havde proklameret, og fordrede
kun at træde i et frit Kontraktforhold til den enkelte Arbejder;
de benyttede Friheden til at isolere. Saa meget mere gjaldt det
da for Arbejderne om at forene sig; det blev en Grundbetingelse
for dem i Kampen for Tilværelsen.

Derved fremkaldtes Fagforeningerne, der især i England
have udviklet sig paa en ejendommelig og betydningsfuld Maade\footnote[1]{Jeg holder mig i det Følgende særligt til de engelske Fagforeninger (Trades Unions), hvis Historie er interessantest og bedst kendt, og støtter mig herved til Lu jo Bre n tan o: Die Arbeitergilden der Gegenwart (Leipzig 1871---1872) og til Sidney and Beatrice Webb: History of Trade Unionism. (London 1894). For Danmarks Vedkommende henvises til J.Jensen og C. M. Olsen: Oversigt over Fagforeningsbevægelsen i Danmark 1871—1900. (København 1901).}.
Selve deres Udviklingshistorie frembyder en stor etisk Interesse.
De havde en haard Kamp at bestaa, dels med Statsmagten, der
ikke vilde anerkende dem, dels med Arbejdsgiverne, dels med
deres egne Tilhængeres Raahed. Man nægtede Arbejderne Ret
til at slutte sig sammen om deres fælles Interesser, og jo mere
man stillede dem udenfor Loven, des mere optraadte de med
Voldsgerninger, Ødelæggelse af Maskiner og Mord. Ved den
Klogskab og det Maadehold, de mere og mere lagde for Dagen,
have de dog efterhaanden vidst at skaffe sig Anerkendelse; og
jo mere denne Anerkendelse af deres Formaal og deres Virksomhed
trængte igennem, des mere ophørte Raaheden hos deres
% --- p. 377 ---
\setcounter{footnote}{0}%
\opage{377}Tilhængere. Deres Stræben gaar navnlig ud paa at opnaa en
jevn og sikker Stigning af Arbejdslønnen og i hvert Tilfælde at
forebygge en Tilbagegang i Arbejdernes Vilkaar. De indse Vigtigheden
af, at Arbejderne vænnes til at stille visse Fordringer
til Livet, vænnes til at kæmpe for en vis Maade at leve paa
(standard of life), for at de ikke skulle lade sig trykke ned til
den blotte dyriske Selvopholdelse. Et saadant Niveau frembringes
ikke ved pludselige og stærke Svingninger, men ved regelmæssige
og varige Forhold. De lægge derfor ogsaa større Vægt paa Forkortelse
af Arbejdstiden, Tilvejebringelse af sundere ydre Forhold,
Sikkerhedsforanstaltninger o. s. v. end paa Forhøjelse af
Lønnen. Herhen hører ikke mindst Bestræbelsen for at lade
det usunde Hjemmearbejde afløses af Arbejde i Værksteder.
Denne Bestræbelse er med Rette bleven anført som Vidnesbyrd
om Arbejdernes intellektuelle og moralske Udvikling. „En Arbejderklasse,
som er nedsunken i Uvidenhed og socialpolitisk
Afhængighed, vil ikke ofre Noget paa at befri sit Hjem fra
Værkstedets Uhygge, men udelukkende rette sine Bestræbelser
paa at faa lidt mere Betaling for sit Arbejde og lade sig nøje
hermed“\footnote[1]{Jensen og Olsen: Fagforeningsbevægelsen i Danmark p. 69.}. --- Ved det Overblik over Markedet og over Arbejdsforholdene
paa de forskellige Steder, de ved deres udbredte
Organisation forstaa at skaffe sig, blive de istand til at afgøre,
hvilke Fordringer de have Udsigt til at sætte igennem. Kampen
for Arbejdernes Fremgang føres nu ikke længere i Blinde, men
støttet til klar Indsigt ide virkelige Forhold. Denne Organisation
fra Arbejdernes Side har igen gjort Oprettelsen af Forligs- og
Forhandlingskamre mulig, saa at Repræsentanter for Arbejdere
og Arbejdsgivere kunne mødes til Drøftelse af fælles Anliggender.
Der er her indenfor det enkelte Folk sket noget Lignende, som
naar Voldgift mellem stridende Folkeslag træder i Stedet for
Krig. Og først gennem Fagforeningernes Kamp blev Arbejdernes
borgerlige Selvstændighed anerkendt. Det er betegnende, at
den engelske Lov, der regulerede Forholdet mellem Arbejdsgivere
og Arbejdere, ligetil 1875 kaldtes „Lov om Husbonder og Tjenere“
(Masters and Servants Act); først i det nævnte Aar afløstes denne
% --- p. 378 ---
\setcounter{footnote}{0}%
\opage{378}Benævnelse af „Lov om Arbejdsgivere og Arbejdere“ (Employers
and Workmen Act\footnote[1]{S. and B. Webb: History of Trade Unionism. p. 274 f.}. Et Bevis paa, at Friheden først bliver
gennemført ved Associationens Hjælp, skønt den selv er Betingelse
for Associationens Opstaaen! I deres Kamp for Foreningsfriheden
staa Arbejderne derfor ogsaa sammen med den radikale
Individualismes Tilhængere. Men medens disse mente, at Alt
var opnaaet, naar Friheden --- Foreningsfriheden indbefattet ---
var gennemført, var Associationen for Arbejderne Grundlag
for en videre gaaende Udvikling.

Associationen medfører Virkninger af stor etisk Betydning.
--- Den Enkeltes Selvstændighedsfølelse holdes oppe, naar han
ser, at han ved sin Deltagelse i Fagforeningen kan lægge et
Lod i Vægtskaalen, medens han ved at staa isoleret Intet formaar
at udrette. Der stilles tillige Fordringer til hans Dygtighed,
ti Fagforeningerne fordre, at enhver Deltager skal have
lært Arbejdet og virkeligt kunne fortjene en sædvanlig Dagløn.
Udygtige Arbejdere vil man ikke vide af, da de bevirke, at
Lønnen falder. --- Den kammeratlige Følelse, Sympatien med
Andre voxer. Den Enkelte ser, at hans Adfærd under og udenfor
Arbejdet, hans Dygtighed, Flid, Selvbeherskelse og Sparsommelighed
blive af Betydning for langt Flere end ham selv. Der dannes
saa at sige en etisk Sfære omkring ham, en stor Familie, af
hvilken han føler sig som Led. Han lærer Broderskabet at kende,
som før næsten glemtes for „Frihed og Lighed“. Han lærer at
underordne sine egne Interesser under de fælles. Han føler sig
solidarisk med sine Fagfæller og --- gennem de forskellige Fagforeningers
Forbund --- med andre Arbejdere, ja med andre
Landes Arbejdere. Hans Horizont udvides; han faar Evne til
at stille sig større Formaal, og han voxer gennem Forholdet til
disse større Formaal. Og han faar gennem de fælles Erfaringer
og det Overblik over Handels- og Fabrikforholdene, der bestemmer
Fagforeningernes Politik, en klarere Opfattelse af Arbejdernes
Stilling til andre Samfundsklasser, lærer baade sine
Rettigheder og Pligter som Led af Slægten bedre at kende. ---
Det er kort sagt en hel Opdragelse fra Egoisme til Sympati, fra
% --- p. 379 ---
\setcounter{footnote}{0}%
\opage{379}blind Raahed til klart seende Kraft, fra Kamp til fredelig Forhandling,
som her foregaar. Og alt dette sker ad Frihedens Vej.
Der gives intet bedre Svar til dem, der betragte vor Tid som
en blot Opløsningens Tid, end at henvise dem til den Samfundsdannelse
og den etiske Udvikling, som her gaar for sig. Og den
foregaar efter de forhen (XIII, 4) fremdragne psykologiske Love.
Samfølelsen udvikles gennem Samliv og Samvirken, gennem
fælles Skæbne og fælles Arbejde. Det aristoteliske Princip
lægger sig her tydeligt for Dagen. En ejendommelig Motivforskydning
fremtræder derved, at skønt det ofte er den rent
individuelle Kamp for Selvopholdelsen, som fører Arbejderen
til at søge Association, kan han dog leve sig saaledes ind i
de fælles Livsinteresser, at disse umiddelbart blive Formaal for
ham. Og endeligt frembringes der gennem disse nye Samfundsformer
Virkninger af Betydning langt ud over Arbejdernes Kreds.
De andre Klasser i Samfundet opdrages ved dette nye Element,
der aktivt griber ind i Udviklingen; de lære deres Grænser bedre
at kende, og de faa tillige nye Formaal og Opgaver, idet det
nye Element medfører en successiv Omdannelse af Samfundslivet
idethele.

Ikke alle Fagforeninger have naat det samme Udviklingstrin.
Baade i Organisation og i Fremgangsmaade gør der sig
store Forskelligheder gældende. Vi have her holdt os til det
Billede, de mest fremskredne Foreninger frembyde. --- For den
enkelte Arbejder kan der under en af Fagforeningen besluttet
Arbejdsnedlæggelse komme en alvorlig etisk Konflikt, idet han
kan blive stillet mellem sin sultende Familie og det, han maa
anse for sin Stands Ære og Velfærd. Fagforeningernes Færd
overfor de saåkaldte „Skruebrækkere“ har ofte været streng;
men man maa erindre, at der her foreligger en etisk Konflikt.
Hvis Arbejdsnedlæggelsen virkeligt er i hele Standens Interesse,
er det utvivlsomt Pligt for den Enkelte --- en Pligt, som Solidaritetsfølelsen
indskærper overalt, hvor den rører sig --- at holde
ud saa længe som muligt. Selvom han ikke hører til Fagforeningen,
vil han høste Gavn af dens Sejr, og han vil derfor ikke
kunne skille sin Sag fra dens i Kampens Tid. Det er et stort
Ansvar, som de paatage sig, der erklære Krigen; men er Krigen
% --- p. 380 ---
\setcounter{footnote}{0}%
\opage{380}erklæret, maa den Enkelte finde sig i den uundgaaelige Lidelse.
Og i disse Kampe er der utvivlsomt i de snevre og skjulte Forhold
lagt Egenskaber for Dagen, som paa en større Skueplads
vilde have ført historisk Berømmelse med sig. En nationaløkonomisk
Forfatter, der langtfra at være en ubetinget Beundrer af
Fagforeningerne kritiserer dem meget skarpt, \textsc{Stanley} \textsc{Jevons},
siger om dem: „Jeg tvivler ikke om, at dersom Strikers og
Arbejderstridigheders Historie blev fuldstændigt skreven, vilde
den frembyde ligesaa mange Exempler paa Troskab og Heltemod
og paa frygtløs Lidelse af Elendighed, ja endog af Døden,
som mangen en Krig, der beskrives i Historien“\footnote[1]{Smlgn. Brentano: Arbeitergilden. 11, p. 233—244.}.

Man har indvendt mod Fagforeningerne, at de ved at stræbe
efter højere Løn og gunstigere Arbejdsvilkaar sætte Prisen op
paa Varer, som Arbejdere i andre Fag have Brug for. Man fordrer
altsaa her af Arbejderne, at de skulle tage et Hensyn til
Forbrugerne, som Arbejdsgiverne ikke tage, naar de sætte Prisen
paa deres Varer. Man foreholder dem et Ideal, som man
ikke falder paa at foreholde Handlende og Fabrikanter. Men
ogsaa afset herfra er Indvendingen übegrundet\footnotemark[1]. Naar Arbejdslønnen
stiger for et Fags Vedkommende, bliver det ogsaa til
Fordel for Arbejdere af andre Fag, idet de bedrestillede Arbejdere
have større Købeevne. Følgen er altsaa kun den, at en
større Del af den nationale Indtægt kommer den arbejdende
Klasse til Gode. Den fordelende Retfærdighed har gjort et Skridt
fremad.

Man vil med større Ret kunne indvende mod dem, at de
grundlægge et Aristokrati indenfor Arbejderklassen ved kun at
ville optage dygtige og udlærte Arbejdere. En vis aristokratisk
Tendens kan ogsaa spores mellem de forskellige Klasser af Arbejdere
paa samme Fabrik, f. Ex. naar de faglærte Arbejdere
forlange at have en Spisesal for dem selv og ikke tillade ulærte
Arbejdere at tage Plads imellem dem. Men al social Udvikling

*) Trades Societies, their objects and their policy. (I Skriftet: 
% CHECK p. 380: note without a star: runs over from the page before?
\footnote[1]{thods of Social Reform). p. 115. Smlgn. hermed en lignende Udtalelse af Kropotkin: Mutual Aid. London 1902. p. 270.}
% CHECK p. 380: note mark without a note
Me%
% --- p. 381 ---
\setcounter{footnote}{0}%
% CHECK p. 381: notes paired with the marks in order (the layer misread a number)
\opage{381}foregaar nu engang lagvis. Hele Massen kan ikke organiseres
paa én Gang. Det er et stort Fremskridt, naar foreløbigt de
øverste Lag komme med i Udviklingen. Og Fagforeningernes
Historie viser, at selvom de skulle have Tilbøjelighed til at afslutte
sig og danne et Aristokrati, virker det Princip, Associationsbevægelsen
skyldes, i bestandigt nye Lag. Ogsaa de „ulærte“
Arbejdere begynde at organisere sig; Londoner Dokarbejdernes
Strike (1889) danner i denne Henseende Epoke i Associationens
Historie. Kort efter organiseredes Arbejdsmændene i København,
og deres Fagforening anerkendtes efter nogen Kamp fra Arbejdsgivernes
Side\footnote[1]{Jensen og Olsen: Fagforeningsbevægelsen i Danmark, p. 51. --- Smlgn. med Hensyn til den nyeste Udvikling i Udlandet: Harald Scavenius: Fagforeningsbevægelsen i Frankrig (1911) og A f den nyere engelske Arbejderbevægelses Historie (1912).}.

4. Med alle deres Fortrin repræsentere Fagforeningerne dog
kun en enkelt Stands Interesser. Og naar de ved deres Organisation
af Arbejderne gøre Mægling mellem disse og Arbejdsgiverne
mulig, saa forudsætter Mægling altid forudgaaende Strid og Modsætning.
Da nu paa den anden Side Mægling kun kan lykkes,
naar der kan paavises fælles Interesser, saa ligger det nær at
spørge, om der ikke kunde gives frie Associationer mellem Arbejdere
og Arbejdsgivere, ligesaavel som mellem Arbejdere indbyrdes.
Og saadanne finde vi ogsaa i den nyeste Tid. Kloge
og sympatisk stemte Arbejdsgivere, som ikke blot ærgre sig
over Tabene ved Striker og Lock-out'er, men sørge over den
Bitterhed og Ufred, den nuværende Arbejdsordning saa ofte fører
med sig, have indført at give Arbejderne Del i Udbyttet for
derved at knytte dem nøjere til deres Gerning\footnote[2]{Stanley Jevons: On Industrial Partnerships. (Methods of Social Reform. p. 122—155). --- Stuart Mill: Principles of political Economy. IV, 7, 5. --- I sin Kritik af Fællesforeningerne viser Beatrice Webb, at de kunne blive til Skade for Fagforeningernes Stræben. Die britische Genossenschaftsbewegung. Leipzig 1893. p. 139—145.}. En vis bestemt
Del af Udbyttet er forbeholdt Arbejdsgiveren; af Resten
fradrages derpaa en Sum til Reparation af Maskiner og Udvidelse
af Bedriften; og hvad derefter bliver tilbage, deles i to Dele,
% --- p. 382 ---
\setcounter{footnote}{0}%
\opage{382}af hvilke den ene tilfalder Arbejdsgiveren, den anden fordeles
mellem Arbejderne i Forhold til den Løn, de oppebære. Desuden
gives der Arbejderne Adgang til at købe sig Andele i
Forretningen. --- Saaledes er idetmindste Ordningen ved nogle
af disse Associationer.

Saadanne Fællesforeninger skyldes Arbejdsgivernes Initiativ;
men deres Opstaaen er dog ikke en Naadessag. I Fag, som
egne sig for denne Ordning, --- og det er især dem, hvor Arbejdernes
Paalidelighed ved Arbejdets Udførelse og ved Behandlingen
af Maskiner og Redskaber er af stor Betydning, --- i
saadanne Fag ville ogsaa Arbejdsgiverne staa sig ved at stifte
Samfund med deres Arbejdere. Ligesom Livegenskabets Ophævelse
blev til Fordel for Besidderne, saaledes vil Fremskridtet
ogsaa her lønne sig. Arbejdernes forøgede Flid og Omhu,
den Maade, hvorpaa de indbyrdes kontrollere hverandre, den
Fred og Tillid, som hersker hos Alle, Undgaaelsen af Arbejdsstandsninger,
--- alt dette bringer materiel og aandelig Vinding
med sig. --- Det er ikke blot Arbejderne, det er ogsaa Arbejdsgiverne,
der trænge til at opdrages. Det sociale Spørgsmaal kan
kun bringes sin Løsning nærmere, naar Arbejdsgiverne betragte
deres Stilling som en social Opgave, der medfører sociale Pligter.
Til disse Pligter hører ikke blot det at levere visse Produkter
saa godt og billigt som muligt, men ogsaa det at lede
det lille Samfund, i hvis Spidse de staa, frem til større materiel
og aandelig Velfærd. Arbejdsgivere kunne anvende en Art social
Rovdrift af fordærvelig Art. Saaledes naar de anlægge Fabriker,
trække en Mængde Arbejdere til, høste det størst mulige
Udbytte af Driften --- og derpaa stanse det Hele, enten fordi
de have tjent nok, eller fordi det ikke mere betaler sig. Arbejderne
afskediges og staa nu, maaske med store Familier, brødløse og
arbejdsløse. Hele den Aand, i hvilken Arbejdsgivernes Virksomhed
udøves, er derfor af største Betydning, og de sociale Forholds
Fremtid beror for en stor Del paa deres, ikke blot, som
det saa ofte farisæisk hedder, paa Arbejdernes etiske Egenskaber.
Det er ikke Velgørenhed og store Opofrelser, der kræves. Det
Vigtigste vil her være et skarpt, af Sympati ledet Blik for de
fælles Interesser. Det er muligt, at det har lange Udsigter, inden
% --- p. 383 ---
\setcounter{footnote}{0}%
\opage{383}saadanne Egenskaber udvikle sig. Det er nu én Gang saaledes,
at de Utilfredse snarere forbedre sig end de Tilfredse. Og det
vilde i hvert Tilfælde ikke være heldigt, om Arbejderne vilde
indskrænke sig til at vente taalmodigt, indtil det sker. Derved
vilde deres egen Udvikling og deres Selvstændighed lide. Maaske
har ogsaa \textsc{Stuart} \textsc{Mill} Ret, naar han mener, at „længe førend
de højere Klasser kunne blive tilstrækkeligt forbedrede til
at kunne udøve et beskyttende Regimente, ville de lavere Klasser
være saa meget forbedrede, at de ikke kunne regeres paa
denne Maade“.

5. Der gives en Art fri Association, som endnu mere end
de hidtil nævnte gaar ud paa at ophæve Forskellen mellem Arbejdsgiver
og Arbejder, og mellem intelligent og fysisk Arbejde.
Det er Produktionsforeningerne, Forbindelser af Arbejdere, der
med sammensparede eller overladte Penge selv købe Produktionsmidlerne
og siden fordele Udbyttet mellem sig\footnote[1]{Stuart Mill: Principles of political Economy. IV, 7, 6. --- L. Brentano: Die christlich-soziale Bewegung in England. Leipzig 1883.}. Disse Foreninger
vidne om den Energi, Intelligens og Opofrelsesevne, der
kan findes hos Arbejderne, og indeholde allerede derved et godt
Varsel for Fremtiden. For at naa deres Maal underkaste de Arbejdere,
der stifte disse Foreninger, sig Savn og Møje og desuden
en Tvang og en Disciplin, som Mænd, der arbejde i Andres
Tjeneste, ikke vilde finde sig i. De have ofte formaat at
arbejde sig frem trods al den Modstand, Autoriteterne lagde
dem i Vejen. Der er hos Deltagerne i disse Foreninger vakt
en Interesse for Oplysning, saavel som en Moralitet, der neppe
vilde kunne frembringes ved Prædikener eller gennem Afholdsforeninger.
--- Ofte ere Produktionsforeninger grundlagte af Fag- eller
Brugsforeninger.

De forudsætte dog Egenskaber, som foreløbigt kun ville være
tilstede hos ganske Faa. Det har vist sig, at de kun trives, hvor
der begyndes med selvsparede Midler, medens Understøttelse fra
Staten eller fra Private ikke har virket heldigt. Det vil --- idetmindste
foreløbigt --- kun blive en lille Elite, som paa denne
Maade vil arbejde sig frem. Og ofte ende Produktionsforenin%
% --- p. 384 ---
\setcounter{footnote}{0}%
\opage{384}gernes heldige Stiftere, naar de ikke besjæles af stor og varig
Entusiasme, som Kapitalister og Arbejdsgivere, der tage andre
Arbejdere i deres Tjeneste paa de gængse Vilkaar.

6. Gennem Fag-, Fælles- og Produktionsforeningerne søge
Arbejderne at skaffe deres Stand Del i Udbyttet af Produktionen.
Men der gives en anden Vej, ad hvilken Sammenslutning
kan føre til at hæve Standens Kaar. Det gælder jo ikke blot
om, hvor mange Penge man tjener ved sit Arbejde, men ogsaa
--- og nok saa meget --- om, hvor Meget man kan faa for sine
Penge. Og Arbejderbefolkningen i sin Helhed er jo Samfundets
største Køber; den har da et Magtmiddel i selve sit Forbrug,
der ingenlunde --- som Karl Marx mente --- er forsvindende i
Forhold til Kapitalisternes. Det gælder kun om at drage Fordel
af selve dette Forbrug. Ved Dannelse af Konsum- eller
Forbrugsforeninger bliver det muligt at faa Del i den Gevinst,
der ellers tilfalder Mellemhandlerne. Forbrugerne danne en Association,
hvis første Formaal er Opnaaelsen af billigere Næringsmidler
derved, at Vareomsætningen saa direkte som muligt ledes
fra Producenter til Konsumenter. Men som det gaar ved de
andre Associationer, gaar det ogsaa ved disse: der udvikler sig
nye værdifulde Egenskaber som Følge af Samvirksomheden, og
højere, mere omfattende Formaal stilles foruden de oprindelige
og elementære. Hvor Forbrugsforeningerne ret have udviklet
sig (især i England og Schweiz) \textbackslash{} have de været mægtige Aarsager
til Udvikling af Solidaritetsfølelsen. De have desuden krævet
Redelighed og Uegennyttighed, Klogskab og Energi, demokratisk
Sindelag og Evne til Selvstyre. De ere ved Siden af Fagforeningerne
blevne den vigtigste Skole for Arbejderstanden, det
vigtigste Middel til social Organisation af denne Stand, der efter
% CHECK p. 384: note whose mark was not found
\footnote[1]{Beatrice Webb: Die britische Genossenschaftsbewegung. Leipzig 1893. --- Hans Müller: Die schweizerischen Konsumgenossenschaften. Basel 1896. --- Først efter Socialistlovenes Ophævelse fik de tyske Arbejdere ret Blik for den Betydning, Brugsforeningerne kunde have for Arbejderstandens Udvikling. Først fra Begyndelsen af 90'erne have de ret taget Opsving, tildels under voldsom Bekæmpelse fra Detailhandlernes Side. Smlgn. hermed R. Riehn: Das Konsumvereinswesen in Deutschland. Stuttgart und Berlin 1902.}
% --- p. 385 ---
\setcounter{footnote}{0}%
\opage{385}Stavnsbaandets og Lavenes Ophævelse stod som en kaotisk Masse.
Af særlig Betydning er det blevet, at Arbejdet i saadanne Foreningers
Tjeneste har frembragt en Stab af Mennesker med administrativ
Evne og med Forstaaelse af omfattende Interesser.
Det er en Udvikling, der er foregaaet fra neden, gennem Øvelse
af Kræfterne i de mindre Kredse og ofte paa tilsyneladende ubetydelige
og materielle Opgaver. Det er dog ikke blot Forskaffelsen
af billige Næringsmidler, der staar som Formaal. Man har
ogsaa ideelle Formaal, særligt Medlemmernes intellektuelle og
æstetiske Dannelse, for Øje, og mange Forbrugsforeninger have
i deres Bygninger Foredragssale og Biblioteker. Ved Hjælp af
fælles Kapital gaar man desuden ofte over til fælles Produktion
af Nødvendighedsgenstande, og det har vist sig, at Produktionsforeninger
af denne Art føre en sundere Tilværelse end de, der
dannes paa bar Grund, og at de ikke saa let udarte til blotte
Aktieselskaber. Denne hele Bevægelse vil maaske blive den vigtigste
Modvægt mod den extreme Arbejdsdeling, hvortil Udvikling
af Produktionsvirksomheden har ført. Den forener ligesaa
meget, som denne har adskilt. Og den vil mere og mere kunne
omfatte alle Omraader, de ideelle saavel som de materielle. Med
Rette har man henvist til den frikirkelige Bevægelse som Forbilledet
for Brugsforeningerne og til det fælles Maaltid som det
gamle Billede paa inderligt Samfund mellem Mennesker\footnote[1]{M Ch. Gide: La solidarité économique. (I Skriftet „Philosophie de la solidarité“). Paris 1902. p. 229—231.}.

Det er ikke nødvendigt nærmere at udføre, hvorledes her,
ligesom ved de tidligere omtalte Associationer, de tre Hovedlove
for etisk Udvikling (XIII, 4) lægge sig for Dagen. Hvilken Fremtid
der maatte være disse og lignende Associationer beskikket,
saa indeholder deres Historie sikkert de mest betydningsfulde
historiske Erfaringer, Menneskeslægten i det sidste Aarhundrede
har gjort. De Samfundsformer, der udvikle sig paa Frihedens
Grund, have den store Fordel, at de skyldes Krav, der af sig
selv gøre sig gældende og ikke kunstigt fremlokkes. De prøves
først i mindre Kredse, før de brede sig i større. De vise os,
% --- p. 386 ---
\setcounter{footnote}{0}%
\opage{386}hvorledes Selvbeherskelse, Sympati og Sans for ideelle Goder
vaagne, saasnart Menneskene i Fællesskab kæmpe for Livet og
ikke staa hvert for sig med deres blinde Selvopholdelsestrang.

b. Arbejdets Ordning ved Statens og Kommunens

% CHECK p. 386: centred line: heading or verse
\begin{center}Indgriben.\end{center}

7. De i det Foregaaende skildrede sociale Fænomener høre
afgjort ind under det frie Kultursamfunds Begreb. De bestod i
fri Sammenslutning, betinget ved Stræben efter fælles Formaal.
Deres etiske Betydning bestod i de psykologiske Ændringer og
Forskydninger, som Livet i disse Associationer fremkalder, og i
de nye Pligter, Opgaver og Konflikter, som opstaa ved Forholdet
til dem. Men det Spørgsmaal opstaar naturligt, hvorledes
disse frie Kulturforbund forholde sig til de Samfundsformer, hvis
Ordning ikke overlades til frie Kræfters uvilkaarlige eller vilkaarlige
Forbindelser, men om fornødent gennemføres med Magt
og Tvang. Hvis man ved Socialisme forstaar den Opfattelse, at
Staten eller Kommunen bør raade over alle Produktionsmidler
og bestemme deres Anvendelse saavel som Fordelingen af Produktionens
Udbytte, saa er det det frie Kultursamfunds Forhold
til Socialismen, vi her staa overfor. Men det maa vel mærkes,
at Udtrykket Socialisme bruges om meget forskellige sociale Opfattelser,
Noget, som især de, der skræmmes ved Ordets blotte
Lyd, ikke maa glemme. Under alle sine Former betegner Socialismen
Hævdelsen af og Beundringen for et socialt Fremtidsideal,
men de enkelte Træk i dette Idealbillede, og endnu mere
hele Fremtidsidealets Forhold til den nærværende Tilstand saavel
som den Maade, paa hvilken Overgangen skal ske fra den
nærværende Tilstand til Fremtidstilstanden, kunne fremtræde paa
højst forskellig Vis i de forskellige Opfattelser. Vi maa derfor
i Korthed karakterisere de forskellige Hovedformer af Socialisme.
Der viser sig, naar vi stille dem i historisk Rækkefølge, en rytmisk
fremskridende Udviklingsgang i de socialistiske Ideer, fremkaldt
ved Vexelvirkningen med de historiske Forhold.

a. \textsc{Platon} har i Staten skildret et idealt Samfund, i hvilket
den private Ejendomsret var ophævet for de regerende Klassers
Vedkommende, for at de kunde leve for offentlige Interesser og
% --- p. 387 ---
\setcounter{footnote}{0}%
\opage{387}intellektuelle Bestræbelser. Paavirket af ham, men især under
Indflydelse af de sociale Forhold i England ved Middelalderens
Slutning, skrev \textsc{Thomas} \textsc{More} sin Utopia (1516). Og et Aarhundrede
efter skrev CAMPANELLA i det Fængsel, i hvilket
hans Deltagelse i Syditaliens sociale og politiske Bevægelser
havde bragt ham, sin Solstat (Civitas solis. 1623)\footnote[1]{i Se om dette Skrift Den nyere Filosofis Historie'\textsuperscript{2}. I. p. 152 f.*) Smlgn. mine Mindre Arbejder. p. 108—112. 25'}. Baade i
den platoniske Stat, i Utopia og i Solstaten er den sociale Ordning
fastslaaet ved tvingende Love. Der skelnes endnu ikke
mellem Samfund og Stat, og man kender ikke de frie Kræfters
Magt. Derved adskille de sig fra de senere Utopier, som f. Ex.
Fr. \textsc{Chr}. \textsc{Sibberns} (i Meddelelser af Indholdet af et Skrift fra
Aaret 2135), der først lader det fuldstændige Fællesskab indtræde,

efterat alt Statsliv og al tvungen Myndighed er ophævet\footnotemark[1].
--- Interessen ved disse ideale Konstruktioner beror
især paa den Kritik af de faktiske sociale Forhold, de forudsætte.
Det er Uretfærdigheden og Nøden, der ved Kontrastvirkning
fremkalde de sociale Ideer. Til Grund ligger herved
den samme Maalestok, som vi i det Foregaaende have opstillet,
uagtet denne Maalestoks Anvendelse ikke sker med fuld Konsekvens.
Mest karakteristisk for den i de nævnte Skrifter fremtrædende
Form for Socialisme er det, at Fremtidsbilledet uden
videre stilles i Modsætning til den faktiske Tilstand. Der gøres
intet Forsøg paa at vise, hvorledes Overgangen fra Nutiden til
den ideale Fremtid kan gøres. Til alle Tider har den kritiske
Holdning overfor den faktiske Tilstand og den begejstrede Tro
paa Idealet været Socialismens Kendemærke, saaledes at snart
det kritiske, snart det ideale Element havde Overvægten. Og
sin store Indflydelse skylder den dels Kritiken og Indignationen
overfor det Givne, dels Fremtidsbilledets dragende Magt. Menneskene
behøve ikke blot Kritik, men store Billeder, der kunne
fylde deres Sind og give deres Trang til Haab et bestemt Indhold,
og dette Behov fyldestgøre hine Utopier i Forhold til
deres historiske Betingelser. Andre Former for Socialisme søge
paa forskellig Maade at finde Mellemled mellem Tvivlen og
% CHECK p. 387: note mark without a note

% --- p. 388 ---
\setcounter{footnote}{0}%
\opage{388}Troen, mellem det mørke Nutidsbillede og det lyse Fremtidsbillede.

β. En Række Mænd fra de første Aartier af det 19. Aarhundrede
ventede sig store Ting af frie Kræfters Samarbejde.
Saint-Simon\footnote[1]{Se om denne Mand, hvis Ideer bleve af stor Betydning for Filosofiens Udvikling: Den nyere Filosofis Historie\textsuperscript{2}. 11. p. 313—317.}, \textsc{Charles} \textsc{Fourier} og \textsc{Robert} \textsc{Owen} mødtes i
den Grundtanke, at den Nød og Uretfærdighed, der opstod ved
Menneskenes gensidige Udbytten af hverandre, vilde afhjælpes,
naar de sluttede sig sammen for i Fællesskab at udbytte Naturen.
Her peges altsaa hen paa en bestemt Vej til det ideale
Maals Opnaaelse. Men Begrænsningen hos disse Mænd ligger
deri, at de ikke igen kende noget Middel til at komme ind
paa Vejen, --- til Dannelse af de frie Sammenslutninger og
Broderskaber. De tænke nærmest paa Foreninger af Mennesker,
der i filantropisk Begejstring direkte ville realisere det sociale
Ideal. Denne filantropiske Socialisme træder altsaa nok Virkeligheden
nærmere end hin ældre Retning, som man kunde kalde
den utopiske Socialisme, men mangler endnu adskillige Mellemled.
Alligevel --- og dette er ikke det mindst Interessante ved den
--- var dens Begejstring ikke uden Nytte, men blev en vigtig
virkende Aarsag indenfor den sociale Udviklingsgang. Det var
for en stor Del de Disciple, som Saint-Simon, Fourier og Owen
vandt indenfor Arbejderklassen, hvis Begejstring og praktiske
Sans førte til Grundlæggelsen af Fagforeninger, Forbrugsforeninger
og Produktionsforeninger. Produktionsforeningerne ere især Saint-Simonisternes
Værk; fra dem optog de engelske „kristelige
Socialister“ dem. I de engelske Fagforeningers og Forbrugsforeningers
Historie spille Owens Ideer og Owens Disciple en
væsentlig Rolle. Den schweiziske Forbrugsforeningsbevægelse er
beaandet af Fourierske Ideer, senere paavirket af den engelske
Bevægelse. Det godtgjordes saaledes, at det ikke er nok at
begynde med Ideerne, men at Livets uvilkaarlige Trang og Udfoldelse
maa mødes med Idebevægelsen. Konstruktion og Erfaring
maa virke sammen. Dette blev muligt ved, at der --- som Følge
af det 18. Aarhundredes Reformer og Opdagelser --- kunde
% --- p. 389 ---
\setcounter{footnote}{0}%
\opage{389}skelnes mellem Samfund og Stat. Der kunde nu finde en fri
social Udvikling Sted, ved hvilken nye Ideer kunde prøves
nedefra, i Livets mindre Kredse, saa at Vejen oppefra, gennem
Lov og Tvang, ikke var den eneste mulige.

γ. Medens de filantropiske Socialister nærmest tænkte paa
fri Organisation af de ved den store Revolution frigjorte Kræfter,
betegner den tredie Hovedform af Socialisme en Tilbagevenden
til den utopiske Socialisme, forsaavidt den lærer Statens Overtagen
af alle Produktionsmidler. Men den mener at være
kommen ud over Utopismen derved, at hvad der i denne kun
stod som Trosgenstand og Fantasibillede, nu bliver udledet og
bevist ad videnskabelig Vej. Denne Retning stiller sig overfor
Utopismen, omtrent som den spekulative Filosofi i Romantikens
Periode stillede sig overfor den positive Religion: enig i Indholdet,
men forskellig ved den videnskabelige Form\footnote[1]{Smlgn. Den nyere Filosofis Historie'\textsuperscript{2}. II. p. 183—187.}. Det er
neppe noget Tilfælde, at denne Retnings Grundlæggere ere udgaaede
fra den Hegelske Filosofi. Den kalder gerne sig selv
den videnskabelige Socialisme, men bør hellere kaldes den
spekulative Socialisme. Den videnskabelige Betydning, som
\textsc{Friedrich} \textsc{Engels} og \textsc{Karl} Marx, Grundlæggerne af denne
Retning, utvivlsomt have, bestaar især i, at de have hævdet
de nationaløkonomiske Begrebers og Loves historiske Karakter,
medens man ellers har været tilbøjelig til at se evige Ideer og
Sandheder i dem. Dette hænger sammen med, at Engels og
Marx bringe det sociale Spørgsmaal i Forbindelse med hele
den historiske Udvikling. Men de opfatte denne Forbindelse
saaledes, at det er de økonomiske Spørgsmaal, hvorom hele
Historien drejer sig. Al Samfundsordning, al Moral og Ret, alle
politiske og religiøse Ideer afhænge ifølge Marx af de økonomiske
Forhold. Hele den ideelle Kultur bestaar kun i den
Maade, paa hvilken Menneskene blive sig deres økonomiske
Existensforhold bevidst. Særligt er den sociale Videnskab selv
et Produkt af den historiske Bevægelse i Retning af nye økonomiske
Forhold, ikke --- som hos Utopister og Filantroper ---
et mere eller mindre vel udtænkt System. Den sande sociale
% --- p. 390 ---
\setcounter{footnote}{0}%
% CHECK p. 390: notes paired with the marks in order (the layer misread a number)
\opage{390}Videnskab opstaar, naar man bliver sig bevidst, hvad der er
sket eller er i Færd med at ske. Og det, der er sket, er, at
der som Følge af den historiske Magtfordeling er indtraadt en
Deling af Arbejdet, ved hvilken Arbejderen skiltes fra Arbejdsmidlerne
(Jorden og Redskaberne), idet Kapitalen koncentreredes
paa forholdsvis faa Hænder. Ved Maskinernes Indførelse fuldendtes
Arbejderens Afhængighed af Arbejdsgiveren. Men efter at
Kapitalisterne som Klasse have exproprieret Arbejderne, exproprieredes
de smaa Kapitalister af de store. Vi staa nu ved et
Vendepunkt, hvor der er et afgjort Misforhold mellem den produktive
Kraft, som ene er at søge i Arbejdet, og Produktionens
Ordning, der lader den allerstørste Del af Udbyttet tilfalde dem,
der selv ikke arbejde, nemlig Kapitalisterne. I dette Misforhold,
der stedse bliver større og større, bestaar det sociale Spørgsmaal.
Løsningen kan kun komme, naar den nuværende Ordnings
Muligheder ere udtømte. Dette vil ske, naar Koncentrationen
af Kapitalen er saa vidt fremskreden, at Staten med Lethed
kan expropriere de faa Kapitalister. Omslaget sker da ved en
Negationens Negation: de hidtil Exproprierende blive nu selv
exproprierede\footnote[1]{„Forvandlingen af den paa Individernes eget Arbejde beroende, spredte Privatejendom til kapitalistisk Ejendom er naturligvis en ulige mere langvarig, haard og vanskelig Proces end Forvandlingen af den faktisk allerede paa social Drift beroende kapitalistiske Privatejendom til social Ejendom. Hist drejede det sig om Folkemassens Expropriation ved nogle faa Usurpatorer; her drejer det sig om Expropriationen af nogle faa Usurpatorer ved Folkemassen.“ Karl Marx: Das Kapital, I. 2. Aufl. Hamburg 1872. p. 793.}. Men Betingelsen herfor er, at Arbejderne have
samlet sig i store Associationer, beredte til at overtage Magten,
og tillige, at store tekniske Fremskridt have gjort Fællesdrift
baade mulig og --- paa Grund af Arbejdsmidlernes Kostbarhed
--- nødvendig\footnote[2]{I et af sine ældre Skrifter udtaler Marx sig klart og interessant om sin Opfattelse at' Socialvidenskabens Opgave. „Socialisterne og Kommunisterne ere Proletariatets Teoretikere. Saalænge Proletariatet ikke er tilstrækkeligt udviklet til at konstituere sig som Klasse, og Proletariatets Kamp med Bourgeoisiet derfor ikke har nogen politisk Karakter, og saalænge Produktionskræfterne hos selve Bourgeoisiet endnu ikke ere til-}. Det gælder først og fremmest om, at Arbejderne
% --- p. 391 ---
\setcounter{footnote}{0}%
\opage{391}vaagne til Bevidsthed om deres Mission og om deres Magt, og
derfor var det Marx's Opraab: „I Proletarer i alle Lande, forener
Eder!“ Kun gennem Klassekampen kan Magten fravristes
dem, der sidde inde med den. Og her er da det Punkt, hvor
Marx's spekulative Tankegang munder ud i Agitationen.

Det mest Karakteristiske for denne hele Opfattelse er den
rent deduktive Maade, paa hvilken den begrundes. Den udledes
af en almindelig historiefilosofisk Teori, ifølge hvilken de økonomiske
Forhold ligge til Grund for al Kultur. Den ideelle
Kultur er kun en Virkning eller Afspejling af den materielle.
Derfor betragtes Arbejdet, hvorved menes det materielle Arbejde,
som Kilde til al Værdi, og derfor nægtes det, at den
historiske Udvikling paa nogen Maade bestemmes ved ideelle
Faktorer. Det er en Kamp om Føde, det ene og alene gælder,
hvormange „ideologiske“ Forklædninger saa dette oprindelige
Motiv optræder i.

Saa simpelt er nu Forholdet mellem Økonomi og Etik og
mellem materiel og ideel Kultur ikke. Det er rigtigt, at materiel
Kultur er Forudsætning for ideel Kultur, forsaavidt som man
maa leve for at kunne leve ideelt, --- og forsaavidt som man
kan leve, uden at leve ideelt. Men deraf følger ikke, at den
ideelle Kultur blot skulde være Virkning og Afspejling af den
materielle. Naar Livets materielle Krav ere nogenlunde tilfredsstillede,
kan den Energi, som ikke benyttes dertil, bruges til
ideelle Sysler, til Tanke og Stemning, til Kunst, Religion og
Videnskab. Men denne Brug er ikke selvfølgelig, og der forestrækkeligt
udviklede til at aabenbare de materielle Betingelser, der ere
nødvendige til Proletariatets Befrielse og Dannelsen af et nyt Samfund,
saa længe ere disse Teoretikere kun Utopister, der udtænke Systemer og
søge en regenererende Videnskab til Afhjælpen af de undertrykte Klassers
Trang. Men efterhaanden som Historien skrider frem, og derved Proletariatets
Kamp aftegner sig tydeligere, behøve Teoretikerne ikke mere at
søge Videnskaben i deres eget Hoved; de behøve kun at gøre sig selv
Regnskab for, hvad der sker for deres Øjne, og at gøre sig selv til Organer
derfor\dots{}. Fra dette Øjeblik af bliver Videnskaben bevidst Produkt
af den historiske Bevægelse; den er ikke længere doktrinær, men er
bleven revolutionær.“ Misere de la philosophie (1847). Tysk Overs. 
% CHECK p. 391: note whose mark was not found
\footnote[1]{gart 1885. p. 121. (Et mod Proudhon rettet Stridsskrift).}Stutt%
% --- p. 392 ---
\setcounter{footnote}{0}%
\opage{392}ligger her Problemer, som Marx ikke har ænset. Desuden
virker den ideelle Kultur, naar den har udviklet sig, tilbage
paa den materielle. Dette sker allerede under saa primitive
Forhold, at vi ikke kende noget Folk, hvis Kamp for Tilværelsen
ikke er bestemt ved Religion, Tradition og Moral. „Ideologien“
bliver meget tidligt en Faktor i Udviklingen, som ingen historiefilosofisk
Teori tør overse. Selvom de Tanker, Tilværelseskampen
fremkalder, først blot staa som Midler og Omveje for at naa det
materielle Maal, saa blive de dog meget snart selv Formaal,
idet der tillægges dem Værdi for deres egen Skyld. Det kan
f. Ex. være, at den enkelte Arbejder først slutter sig til sine
Kammerater for at opnaa en bestemt Fordel derved; men Standens
Ære og Fremgang kan senere blive et Formaal, han forfølger
uden egoistisk Bagtanke. Saadanne Motivforskydninger gøre Forholdet
mere indviklet, end Marx's rent deduktive Fremstilling
kan anerkende.

Marx's Teori vil skildre den sociale Udvikling ganske uafhængigt
af alle ideelle Motivers Indflydelse. Men det sociale
Problems og selve Socialismens Historie godtgør, at saadanne
Motiver faktisk virke med. Og i det Foregaaende (IV, 6) saa
vi, at selve Samvittigheden dog ogsaa er en Kraft, et Led i
den Aarsagsrække, der bestemmer Udviklingen. Efter Marx skal
al Teori, al „Ideologi“ kun være Bevidsthed om, hvad der sker;
men hvilken Værdi har det, der saaledes rører sig i Hjernen
hos flere eller færre Individer, naar det ikke faar nogen praktisk
Virkning? Marx mener da ogsaa, at det har mere end teoretisk
Interesse at finde et Samfunds økonomiske Bevægelseslov.
„Selvom et Samfund“, siger han i Fortalen til „Das Kapital!“,
„er kommet paa Spor efter sin Bevægelses Naturlov, kan det
hverken overspringe eller bortdekretere naturlige Udviklingsfaser.
Men det kan forkorte og formilde Fødselsveerne“. Allerede en
saadan Forkorten og Formilden er Mere, end Marx konsekvent
kunde indrømme. Men Marx har idethele flere Forudsætninger,
end han vil indrømme. Hans Teori er i Grunden en etisk
Teori; det Resultat, han kommer til, hviler paa et etisk Postulat,
som han idetmindste et enkelt Sted kommer til at røbe, det
Postulat nemlig, at et Menneske ikke blot skal behandles som
% --- p. 393 ---
\setcounter{footnote}{0}%
\opage{393}Middel, men stedse tillige som Formaal. Dette Postulat, der
indeholder Begrundelsen af den Indignation, som spores gennem
hans Fremstilling, saa lærd og tung denne end er, udtales i
følgende Sætning: „Ved den kapitalistiske Produktionsmaade er
Arbejderen til, for at foreliggende Værdier kunne finde Anvendelse,
istedetfor at omvendt den objektive Rigdom er til for at
tjene Arbejdernes Trang til Udvikling“. (Das Kapital. 2. Aufl. I.
p. 646). Hvorfra mon Marx véd, hvad Rigdommen er til for?
Ad rent historisk Vej, med tilhørende Deduktion, kan han ikke
vide Noget derom, kun Noget om, hvad Rigdom efter en nødvendig
Udviklingslov vil blive brugt til engang. Deraf, at Arbejderen
kun har ringe Del i den Kultur, han hjælper med til
at frembringe, følger i og for sig selv ingenlunde, at han skal
have en større Andel, selvom det er muligt. Dette følger kun,
hvis man gaar ud fra det etiske Princip, at enhver Personlighed
er Formaal, ikke blot Middel. Det er da intet Under, at dette
Princip uvilkaarligt trænger sig frem hos Marx, skønt han ingen
særlige etiske Forudsætninger vil vedkende sig, men i „videnskabelig“
Fornemhed overlader slige Forudsætninger til Utopister
og Filantroper. Ved en Udtalelse som den anførte fordrer han
udtrykkeligt en Indgriben i Tingenes Gang af etiske Motiver.
Hin Sætning indeholder en hel Etik i Spiren. Naar den praktiske
Vurdering, som saaledes ligger til Grund, tages med, forstaar
man bedre, hvorledes Erkendelsen af Udviklingsloven kan
virke „forkortende og formildende“. Det er Indignationen over
Afstanden mellem Ideal og Virkelighed, som virker æggende og
styrkende i Kampen mod Modstanden. Marx's egen „Ideologi“
er et mægtigt Vaaben i hans og hans Tilhængeres Haand, og
den er maaske fra først af bleven til, fordi den kunde tjene
som Vaaben, ikke blot af rent teoretisk Interesse. Forgæves
søger Marx at fornægte eller dølge den Idealisme, som ligger
bag hans egen Optræden, og som tillige er en nødvendig Betingelse
for, at der i Historien skal kunne udrettes saa store
Ting som Frembringelsen af en ny social Ordning. Han staar
paa et etisk Standpunkt uden at ville være ved det.

I sin Paavisning af Vejen til at naa Maalet lægger Marx
Vægten paa Klassekampen. Den er for ham det vigtigste Fænomen
% --- p. 394 ---
\setcounter{footnote}{0}%
\opage{394}i den moderne sociale Udvikling\footnote[1]{I den Fremstilling af Marxismen, Werner Sombart gav i sine Ziiricherforedrag, blev denne Side af Sagen ganske særligt betonet. Foredragene ere udkomne under Titel: Socialismus und sociale Bewegung im 19. Jahrhundert. Bern 1897. Ved den Diskussion, der fulgte efter disse Foredrag, havde jeg Lejlighed til at udtale mig om Marxismen, saaledes som Sombart havde fremstillet den. (Se det nylig anførte Skrift p. 56—59). Det har interesseret mig at se, at Sombart senere har bebrejdet Marx, at Fagforeningerne for ham kun vare Midler til Revolution, ikke havde selvstændig Betydning for Arbejdernes Udvikling. Sombart har fremdeles senere søgt at vise, „at Kapitalisme og Socialisme ikke ere Modsætninger, der udelukke hinanden, men at deres Idealer til en vis Grad meget vel kunne være virkeliggjorte i et og samme Samfund.“ De geniale Banebrydere, de kongelige Købmænd kunne ikke undværes, hvis det økonomiske Fremskridt skal være muligt. Men Arbejderforholdet skal gaa over fra Privatrettens til den offentlige Rets Omraade, og der skal sikres Arbejderstanden Medbestemmelse ved økonomiske og politiske Afgørelser. (Dennoch! Aus Theorie und Geschichte der gewerbschaftlichen Arbeiterbewegung. Jena 1900. p. 92—94). I sit store Værk Der Kapitalismus (1901) skelner han --- under aabenbar Indflydelse af en mere kritisk Filosofi end den, der laa til Grund hos Marx --- mellem tre forskellige Synspunkter, fra hvilke økonomiske Fænomener maa betragtes: 1) det kausale Synspunkt, som opsøger de faktisk virkende Aarsager og Motiver; 2) det teleologiske [bedre: etiske] Synspunkt, der gaar ud fra bestemte socialpolitiske Formaal; 3) det kritiske Synspunkt, som fører til Drøftelse af de Grundbegreber, det kausale og det teleologiske Synspunkt anvende, med Hensyn til deres Begrundelse og deres Begrænsning. (Der Kapitalismus. I. p. XXXII). --- Med denne Opfattelse er egentligt Overgangen gjort fra den spekulative til den empiriske Socialisme.}. Dog er det neppe rigtigt at
holde sig til den negative Side af Sagen. Hvor et nyt socialt
Lag skal arbejde sig frem, har det naturligvis en haard Kamp
at bestaa med de Lag, der hidtil sad inde med hele den sociale
Magt, en Kamp, som ikke altid er ublodig. Men Vægten maa
ikke udelukkende lægges paa Modsætningen til de andre Klasser.
Der erhverves nye Egenskaber ved den Sammenslutning, det
Broderskab mellem Arbejderne indbyrdes, som Kampen nødvendiggør.
Foreningen gør det muligt, at den Enkelte ved den
Dannelse og Underordning, som fordres i de fælles Formaals
Tjeneste, bedre lærer sine Pligter og Rettigheder som Led i
Menneskehedens Tjeneste at kende. Ved denne stigende Selv%
% --- p. 395 ---
\setcounter{footnote}{0}%
\opage{395}følelse saavel som ved den materielle Magt, Associationerne
raade over, og ved deres politiske Indflydelse komme Arbejderne
mere og mere til at staa som Formaal, ikke blot som Midler i
den sociale Proces. At Marx ikke bestemtere fremhæver denne
Side af Sagen, kommer af, at ifølge hans Teori skal med Kapitalens
stigende Koncentration „Elendigheden, Trykket, Tyranniseringen,
Nedværdigelsen og Udplyndringen“ stige, og derved
den Katastrofe indtræde, hvorved „Expropriatorerne exproprieres“.
Men den opdragende Indflydelse, Associationerne øve paa Arbejderne,
og den stadige Forbedring af de materielle Livsvilkaar,
som disse samme Associationer muliggøre, stemme ikke godt
med Katastrofeteorien. Hvis Marx havde tillagt Associationerne
positiv Betydning, vilde det have været ham umuligt at deducere
Katastrofen.

Trods den videnskabelige Karakter, Marx mener at have
givet Socialismen, har han dog ikke ganske brudt med Utopien.
Dette viser sig ikke blot ved den Sikkerhed, med hvilken han
forudsiger Katastrofen, men ogsaa i de, ganske vist faa og korte
Antydninger, han giver af den sociale Tilstand, der vil indtræde
efter Katastrofen. Efter de faa Kapitalisters Expropriation ved
hele Folkemassen skal der ikke mere existere nogen Klasseforskel.
„Vil der“, spørger Marx i et af sine ældre Skrifter\footnote[1]{Misère de la Philosophie. Tysk Overs. 1885. p. 181 f.},
«efter det gamle Samfunds Omstyrtning gives et nyt Klasseherredømme,
som tilspidses i en ny politisk Magt? Nej. Betingelsen
for den arbejdende Klasses Befrielse er Afskaffelsen
af enhver Klasse\dots{}. Der vil ikke mere gives nogen egentlig
politisk Magt, da netop den politiske Magt er det officielle Udtryk
for Klassemodsætningen indenfor det borgerlige Samfund“.
Et Samfund uden „nogen egentlig politisk Magt“ er aabenbart
en Utopi, og det endog en Utopi, der overgaar Platons, Mores
og Campanellas; ti i deres Idealstater er der en organiseret
politisk Magt. I sine senere Skrifter udtaler Marx sig kun ubestemt
og negativt om den fremtidige sociale Tilstand. Han og
hans Tilhængere mene, at det altid er tidsnok at tale om den,
naar den nuværende Tingenes Orden er forbi. Ogsaa dette
% --- p. 396 ---
\setcounter{footnote}{0}%
\opage{396}hænger sammen med Katastrofeteorien, ifølge hvilken „Omslaget“
eller „Revolutionen“ skal føre til Noget, der er helt modsat det
Bestaaende, men det strider mod den Erfaring, at det Fremtidige
ikke blot negativt, men ogsaa positivt forberedes i det
Forudgaaende. Marx har desuden ikke den naive Tillid til Fantasiens
Konstruktioner, som de gamle Utopister havde, og de
vage Antydninger, han tillader sig, fremkomme snarere til Brug
ved Agitationen end som nødvendige Led i Teorien.

δ. I Modsætning baade til den utopiske, den filantropiske
og den spekulative Socialisme gør der sig i den nyeste Tid,
især i England, en Bestræbelse gældende, som med Rette kan
føre Navn af den empiriske Socialisme. Den gaar ikke ud paa
at konstruere et Fremtidsbillede, og den indser, at filantropiske
Ønsker ikke ere nok, men at der ædrueligt maa spørges og
prøves, hvad der under de givne Betingelser kan opnaas; heller
ikke stoler den paa historiefilosofiske Deduktioner, men vil gaa
induktivt frem, prøve Mulighederne paa Grundlag af bestemte
Erfaringer. Og denne Begrænsning paalægger den sig ikke
blot af kritisk Betænksomhed og af Uvillie mod demagogiske
Deklamationer\footnote[1]{Selskab er opkaldt efter den romerske Feltherre Fabius, der udrettede Mere ved at tøve til det rette Øjeblik og ved at forberede det, end adskillige Brushoveder, der mente, at det gjaldt om at gribe til strax.}, men ogsaa og især, fordi den mere end nogen
af de andre Former for Socialisme anerkender Friheden baade
som Middel og Maal. En social Organisation har kun virkelig
Værdi, naar den bestaar i frie Kræfters Sammenslutning. Derfor
lægger man stor Vægt paa de frie Arbejderorganisationer, som

!) Smlgn. Fabian Tracts Nr. 41. (Bernard Shaw: The Fabian
Society. What it has done, and how it has done it). p. 5. „Vor Forkærlighed
for praktiske Vink og praktisk Kritik og vor Utaalmodighed overfor
alle almindelige Udtalelser af Sympati med Arbejderklassens Bestræbelser,
og derhos vor Lyst til at spøge med vore Modstandere istedetfor at erklære
dem for Fjender af Menneskeheden, stødte flere varmhjertede og
veltalende Socialister fra os\dots{}. Der var altfor megen Lighed og Fortrolighed
mellem Fabierne, til at noget Medlem skulde kunne staa op og
prædike for de Andre paa den Maade, paa hvilken de arbejdende Klasser
endnu finde sig i, at deres Ledere prædike for dem.“ --- Det fabiske
% --- p. 397 ---
\setcounter{footnote}{0}%
\opage{397}i det Foregaaende ere omtalte. De lavere Grader og Former
af Frihed skulle benyttes til at frembringe de højere. Man bebrejder
det kapitalistiske System, at det hindrer Personlighedens
Udvikling hos saa Mange ved at trykke Existensbetingelserne
ned og gøre dem usikre. Det gælder om at hævde et vist
økonomisk Niveau (standard of life, „Levefod“) for Arbejderne,
for at de kunne erhverve og hævde et højere aandeligt Niveau.
Ved Socialisme (eller som man i den nyeste Tid ofte siger:
Kollektivisme) forstaar man her en Lære, efter hvilken det er
Samfundets Sag at hævde saadanne Betingelser for Arbejdet,
at Arbejdernes fysiske og aandelige Udvikling ikke hæmmes.
Der skal tillige tages Hensyn til Individernes forskellige Kræfter;
den Svage skal arbejde i Forhold til sin ringe Kraft, ligesom
den Stærke i Forhold til sin store Kraft. Der skal sørges for
Stadighed i de økonomiske Forhold, for at den lammende Usik kerhedsfølelse
kan falde bort. De ydre, fysiske og sociale Forhold,
under hvilke Mennesket lever, bestemme for en stor Del dets
Karakter. Derfor bør der sørges for, at det sociale Maskineri
i saa høj Grad som muligt formes saaledes, at det kan medføre
heldige Virkninger paa Karakteren. Her kan det ikke nytte at
stole paa Arbejdsgivernes filantropiske Sindelag. Og Arbejderne
selv formaa først ved Erfaring at vurdere sunde, rene og ædle
Livsforholds Betydning. Derfor er det Samfundets, Statens og
Kommunens Sag at tilvejebringe saa gunstige hygiejniske og
moralske Forhold for den store Befolkning som muligt. Arbejderklassen
og dens Talsmænd maa benytte de borgerlige Rettigheder,
som de skylde den radikale Individualisme, til at faa
stigende Indflydelse paa det offentlige Liv gennem Skoleraad,
Kommunalraad og Parlament. Derved kan der høstes Erfaringer
og gøres Forsøg, som kunne vejlede ved den fortsatte Udvikling.
Statssocialismen begynder fra oven og vil virke gennem et
bureaukratisk Regimente og paa et doktrinært Grundlag. Den
empiriske Socialisme lægger Vægt paa Selvstyrelsen i Fag- og
Forbrugsforeninger, i Kommune og Stat; den bevæger sig gennem
de smaa Kredse til de store og søger at udvikle Folkets Kræfter
gennem Arbejdet paa smaa Opgaver til at magte de store Op%
% --- p. 398 ---
\setcounter{footnote}{0}%
% CHECK p. 398: notes paired with the marks in order (the layer misread a number)
\opage{398}gåver. Det er karakteristisk, at den empiriske Socialisme\footnote[1]{Den bedste Fremstilling af den, jeg kender, er Sidney Ball's Afhandling: The Moral Aspect ofSocialism (International --- Journal of Ethics VI. Kfr. den ved denne Afhandling fremkaldte Diskussion i samme Tidsskrift 6. og 7. Bind). Udtrykket „empirisk Socialisme“ har jeg fra Sidney and Beatrice Webb: History of Trade Unionism. p. 403. Disse Forfatteres historiske Værker vise, hvorledes den empiriske Socialisme har udviklet sig. Smlgn. ogsaa Hans Müllers Værk om de schweiziske Konsumforeninger. --- Fra 1883 virker Fabian Society i London gennem Skrifter og Foredrag i samme Retning. (Fabian Tracts, en Række smaa Piecer, klart og koncist skrevne, faas ved Henvendelse til The Fabian Office. 276 Strand. London, W. C.). --- Den empiriske Socialismes Standpunkt er ikke meget forskelligt fra det, paa hvilket John Stuart Mill stillede sig overfor det sociale Spørgsmaal. (Den nyere Filosofis Historie'\textsuperscript{1}. II. p. 427—430). Heller ikke Eugen Diihrings Standpunkt synes at være meget forskelligt (anf. Skr. p. 564 f.).} har
sit Hjemsted i England, medens den spekulative Socialisme har
sit Hjemsted i Tyskland\footnote[2]{Marxismen er en Udløber af den tyske spekulative Filosofi. Endnu i 1891 skrev Fr. Engels under sit Billede: „Vi tyske Socialister ere stolte over at nedstamme --- ikke blot fra Saint-Simon, Fourier og Owen --- men ogsaa fra Kant, Fichte og Hegel. Den tyske Arbejderbevægelse er den tyske klassiske Filosofis Arvetager.“ (Werner Sombart: Friedrich Engels. Ein Blatt zur Entwickelungsgeschichte des Sozialismus. Berlin 1895. p. 13). --- Smlgn. om Marxismens Forhold til den hegelske Filosofi Den nyere Filosofis Historie'\textsuperscript{2}. Il. p. 266. 282 f.}. Medens den spekulative Socialisme
nærmest minder om Utopisterne, hænger den empiriske Socialisme,
baade hvad Karakter og hvad historisk Oprindelse angaar,
nærmest sammen med de filantropiske Socialister.

8. Fra det Standpunkt, hvorfra vi her betragte den sociale
Udvikling, maa vi paa Forhaand sympatisere med Socialismen
i to Punkter: i dens vigtigste Grundtanke og i Hovedtrækkene
af dens Kritik af den nuværende Samfundsorden. Og disse to
Punkter ere fælles for alle Former af Socialisme.

Ofte optræder Socialismen ganske vist som en Drøm, i
hvilken Sindet har fundet Hvile, naar det var pint og skrækket
ved Tidens Ulykker. Men selv da indeholder den en virkelig
etisk Grundtanke: Ideen om en fordelende Retfærdighed, om et
fuldkomment Samfund, hvor Enhvers Evne og Trang kommer
% --- p. 399 ---
\setcounter{footnote}{0}%
\opage{399}til sin Ret. Den optræder særligt i vort Aarhundrede som en
nyttig Modvægt mod den ensidige Individualisme, der spalter
Samfundet i isolerede Individer. Den hævder i Virkeligheden
selve den sociale Etiks Grundtanke: at Individets Stilling i Samfundet
bør være bestemt ved hele Samfundets Tarv (Individets
eget indbefattet). Og det er ud fra denne Grundtanke, at
Socialismen underkaster den nuværende Samfundsorden en skarp
Kritik. Hvad Selvprøvelse og Selverkendelse er for det enkelte
Menneske, er en saadan Kritik for Samfundet. Den blotter
Mislighederne og Lidelserne, og det er den første Betingelse
for, at der kan raades Bod paa dem.

Men man kan anerkende baade Grundtanken og Kritiken
uden at anerkende de Midler og Veje, ved hvilke Tanken skal
virkeliggøres og Mislighederne afhjælpes. Et er at paavise
Sygdommen, et Andet er at anvise Lægemidlet. Mod saadanne
Former for Socialisme, der gaa prøvende og experimenterende
frem og arbejde ved Hjælp af frie Associationer eller af de bestaaende
kommunale og politiske Organisationer, til hvilke Arbejderne
mere og mere faa Adgang, kan der Intet indvendes.
Det skulde da være, at man --- hvad der iøvrigt ikke er usædvanligt
--- anser den nuværende Ejendoms- og Arbejdsfordeling
for Indbegrebet af al Visdom. I hvor høj Grad denne Ordning
kan afløses af en anden, maa Historien vise\footnote[1]{Smlgn. hvad en kyndig og konservativ Iagttager som A lex is de Tocqueville udtaler som Resultat af sin Forsken og sin Erfaring: Ce qu'on appelle les institutions nécessaires, ne sont souvent que les institutions auxquelles on est accoutumé, et en matiére de constitution sociale, le champ du possible est bien plus vaste que les hommes qui vivent dans chaque société ne se I'imaginent. Souvenirs. Paris 1893. p. 111 f.}. Mulighederne
kunne vi ikke gennemskue forud --- saa meget mindre, som de
stadige Forskydninger (XIII, 4) kunne gøre, at der ofte kommer
noget helt Andet ud af de bevidste Bestræbelser, end man paa
Forhaand kunde ane. Motiv- og Værdiforskydningerne gøre Indgrebene
i den sociale Udvikling til Spring i Mørket. Men Intet
at gøre er her ogsaa at springe. Ethvert Indgreb eller Ikke-Indgreb
i Tingenes Gang maa motiveres ved den bedst begrundede
Overbevisning, som kan vindes.

% --- p. 400 ---
\setcounter{footnote}{0}%
% CHECK p. 400: notes paired with the marks in order (the layer misread a number)
\opage{400}Den empiriske Socialisme stiller sig det Maal at ændre
den bestaaende Arbejds- og Ejendomsordning i social Aand; den
lader det henstaa uafgjort, hvor langt Udviklingen vil kunne
gaa. En energisk Medarbejder i den schweiziske Forbrugsforeningsbevægelse
har defineret Socialismen som „Læren om
en saadan Indretning af det menneskelige Samfund, at der raades
Bod paa Misforholdet mellem overvættes Rigdom og Massearmoden“\footnote[1]{xj Se Citatet hos Hans Miiller: Die schweizerischen Konsumgenossenschaften, p. 453.}.
Der stilles her en Opgave, som der Skridt for
Skridt kan arbejdes paa. I Modsætning hertil staar Marxismens
Definition af Socialisme: „Ophævelse af den kapitalistiske Arbejdsordning
ved Socialisering af Produktionsmidlerne, og Klassekamp
som Vejen til dette Maal“\footnote[2]{Werner Sombart: Friedrich Engels. p. 25. Der tilføjes: „Diese programmatische Kernpunkte werden mehr und mehr Gemeingut des kampfenden Proletariats“.}. Det er klart, at de to Opfattelser
stille det frie Kultursamfund i et forskelligt Forhold til
Staten; efter den første skal det frie Kultursamfunds Aand efterhaanden
indpodes i Statsmekanismen; efter den anden gælder
det saa hurtigt som muligt at faa Raadighed over Statsmekanismen
for at kunne bestemme det frie Kultursamfunds Ordning.

I den nærmere kritiske Prøvelse af Socialismen have vi i
det Følgende især denne sidste, spekulative eller marxistiske
Opfattelse for Øje. Det vil vise sig, at Kritiken af den fører
til at lægge saa meget des større Vægt paa den empiriske Socialisme.

9. Statsmagten øver under alle Forhold en stor Indflydelse
paa Fordelingen af Arbejds- og Forbrugsmidler, og man vil forsaavidt
kunne sige, at enhver Forfatning er en relativ Socialisme.
Det vil idethele vise sig umuligt at afstikke absolute Grænser
for Statens Virksomhed. Statens Væsen er ikke uforanderligt;
det udvikler sig med Menneskenaturen og med de historiske
Forhold, og Ingen kan sige, hvorledes Staten engang i Fremtiden
vil blive beskaffen. Men ved at overdrage al Fordeling
til Statsmagten forudsætter man Egenskaber hos de Mennesker,
der udøve Statsmagten (og det er jo altid bestemte Mennesker,
% --- p. 401 ---
\setcounter{footnote}{0}%
\opage{401}der udøve den), som disse hidtil ikke kunne siges at have vist
sig i Besiddelse af. Hvorledes kunne Mennesker blive saa fuldkomne,
at de ikke misbruge en saa uhyre Magt, naar de allerede,
hvad Historien noksom viser, misbruge den mindre Magt,
der hidtil har været de Regerende betrot? Det er karakteristisk,
at i England, det Land, hvor der kan øves den skarpeste Kontrol
med Statsmagtens Anvendelse, nærer man den største Utilbøjelighed
til at udvide dens Omraade. Statssocialismen forudsætter
en Overtro paa Staten, en Tillid til, hvad der kan gøres
ovenfra nedad, som overser, at Statens Ror dog bestandigt holdes
af Mennesker, ikke af Guder. Og dette vil ikke forandres, selv naar
Folkets Majoritet bestemmer Regeringens Sammensætning. Naar
Menneskene blive saa gode og dygtige, som Socialismen forudsætter,
vil det sociale Spørgsmaal allerede være ude af Verden.

Ikke blot moralsk Fuldkommenhed, ogsaa Alvidenhed maatte
Statsmagtens Bærere i den socialistiske Stat besidde. For at kunne
fordele Arbejdet og Udbyttet maatte de kende de forskellige Individers
Evner og Fornødenheder. Men baade Evner og Fornødenheder
ere i stadig Udvikling, og Individet er den, der bedst
kan opdage dem, naar de blot faa Lov at udfolde sig saa frit
som muligt, for at de kunne blive prøvede. Staten har vel ogsaa
nu den Opgave at vælge de til visse Virksomheder bedst
skikkede Individer og at imødekomme forskellige Fornødenheder.
Men der er vel Faa, der mene, at den øver dette Hverv
til en saadan Grad af Fuldkommenhed, at man uden tvingende
Grunde skulde ønske det udstrakt til Alt. Og Staten har, som
den nu er stillet, den frie Udviklingstrang, det frie Initiativ hos
de enkelte Individer at støtte sig til: den kan vælge mellem
dem, der af egen Trang have udviklet sig i en vis Retning.
Den har desuden paa mange Omraader private Foretagender at
konkurrere med. Hvor udvidet man end vil tænke sig Statens
Virksomhed, vil den dog ikke kunne undvære det frie Initiativs
og den private Virksomheds Konkurrence, saafremt den ikke
skal stivne i Doktrinarisme og Slendrian. Det gælder baade
paa den materielle og den ideelle Kulturs Omraade.

Hvad særligt Udbyttets Fordeling angaar, saa opstaar her
det store Spørgsmaal, hvad retfærdig Fordeling egentligt vil sige.
% --- p. 402 ---
\setcounter{footnote}{0}%
\opage{402}De socialistiske Forfattere dele sig her i to Grupper, idet Nogle
ville have Enhvers Andel bestemt ved det Arbejde, han har ydet,
medens Andre ville have Fordelingen bestemt ved Enhvers Trang.
Det første Standpunkt er udtrykt i St. \textsc{Simons} Sætning: „Chacun
doit étre classé selon sa capacité et retribué selon ses
oeuvres!“ --- det andet i det socialdemokratiske Gothaerprograms
Ord: „Jedem nach seinen vernunftgemässen Bedürfnissen
--- i første Tilfælde støder man ikke blot paa den Vanskelighed
at dele lige mellem de forskellige Arter af materielt og
aandeligt Arbejde, som Samfundet behøver, men ogsaa paa det
store Problem, hvorledes man kan sikre sig, at Udbyttets Værdi
svarer til det anvendte Arbejde. Det kommer jo dog an paa,
om Arbejdsproduktet har tilstrækkelig Brugsværdi, det vil sige,
om det tilfredsstiller virkelige Fornødenheder i Samfundet. Værdien
af et Produkt bestemmes ikke blot ved det Arbejde, dets
Frembringelse har kostet, eller den Tid, dette Arbejde har taget.
Det kommer an paa, om der er Brug for det, som er frembragt,
altsaa om der er en Trang, en Fornødenhed tilstede, som kan
tilfredsstilles ved det. Derfor kan man kun regulere Arbejdet,
naar man kan regulere Fornødenhederne. Den første Teori maa
altsaa ligesom den anden tillægge Staten Myndighed til at bestemme
Individernes Fornødenheder. Og man maa tillige sørge
for, at der ikke produceres Andet og Mere, end hvad der kan
forbruges i Landet selv. Ti over de udenlandske Fornødenheder
er den enkelte Stat ikke Herre. Derfor maa Verdenshandelen,
der fører til, at der produceres Mere, end Landet selv behøver,
og derved gør Landet afhængigt af Udlandet, begrænses og kontrolleres
af Regeringen. Denne Konsekvens saa J. G. \textsc{Fichte}
og var ikke bange for at drage den (i sit Skrift Geschlossener
Handelsstaat. 1800). Derved vilde ganske vist tillige Kilden
til de moderne sociale Disharmonier være stoppet; ti disse gaa
historisk tilbage til det 13. og 14. Aarhundrede, da der begyndte
at danne sig et Verdensmarked, saa at Produktionen blev bestemt
ved andre end de hjemlige, let overskuelige Fornødenheder *). ---
% CHECK p. 402: note whose mark was not found
\footnote[1]{Lujo Brentano: Ueber die Ursachen der heutigen socialen Noth. Leipzig 1889. Smlgn. samme Forfatters Die Arbeitergilden der Gegenwart. I. p. 58 ff., II. p. 320.}
% --- p. 403 ---
\setcounter{footnote}{0}%
\opage{403}Naar den anden Teori vil have Fordelingen bestemt ved Enhvers
fornuftige Fornødenheder, er det klart, at det ikke kan
være det enkelte Individs egen Fornuft, der skal afgøre, om hans
Fornødenheder ere „fornuftmæssige“. Det maa være Statsmagtsudøvernes
Fornuft, og de enkelte Individer gøres altsaa umyndige.
Selvom en Regulering af Fornødenhederne er nødvendig,
idet Staten jo maa bestemme sine Embedsmænds Levefod (standard
of life), vil det dog ligesaa lidt her som ved Produktionen
være heldigt, om denne Reguleren ikke har en fri Udfoldelse
og Tillempelse ved Siden, for at Sammenligning og Valg kunne
ske i saa stort Omfang som muligt.

10. En Samling Individer, som hverken have Ret til at
afgøre, hvilke Evner eller hvilke Fornødenheder de besidde, men
som maa finde sig i at blive tilskaarne efter et af Autoriteter
fastslaaet Skema, --- en saadan Samling Individer er en blot
Masse, ikke noget organiseret Samfund. Og det er ligegyldigt,
om man tænker sig Samfundets fremtidige Ledere som Genier'
eller om man tænker sig dem som Idioter, --- om man tænker
sig dem som enevældige Monarker og Diktatorer eller som udgaaede
af hele Folkets Valg. Hvad der giver Livet dets Værdi,
den frie Udfoldelse af Evner og Drifter, er i alle Tilfælde bortfaldet.

Trangen til selv at afgøre, hvilke Evner og Fornødenheder
man besidder, og hvilke af dem der fortjene at udvikles og tilfredsstilles,
er ikke blot en egoistisk Trang. Den er, som vi have
søgt at vise, Betingelsen for, at der i Samfundet kan findes produktive
Kræfter, som frembringe noget Nyt og ikke vedblive at
gaa i de gamle Spor. Selve den sædvanemæssige Virksomhed
øves kun ret, naar Individet selv bestemmer sig til den. Men
naar det gælder om at slaa ind paa nye Veje, er Tilfredsstillelsen
ved at følge sit eget Initiativ til det, man anser for godt
og gavnligt, ofte den eneste Løn, der gives. Store Opfindere
bryde sig ofte slet ikke om den Fordel, de selv kunne have af
deres Opfindelser, og de bedrages da ogsaa ofte nok for den.
Hvad der forbitrer dem Livet og gør deres Skæbne haard, er
de Hindringer, der stille sig i Vejen for den frie Brug af deres
Kræfter i den Retning, de ønske. Hvorledes vil det da ikke
% --- p. 404 ---
\setcounter{footnote}{0}%
\opage{404}gaa, naar alle Arbejdsmidler tages i Beslag af Statsmagten! Hvorfra
skulle da Midlerne komme til de private Experimenter, som
Kulturen skylder saa Meget, hvis den ikke skylder dem Alt?

--- \textsc{Schäffle} har vel hævdet, at man gør Socialismen Uret ved
at mene, at den nødvendigvis ophæver al fri Bevægelse og al
fri Disposition over materielle Goder. Han skelner skarpt mellem
Nydelsesmidler og Produktionsmidler og søger at paavise,
at Socialismen kun ophæver Ejendom som Produktionsmiddel,
ikke som Nydelsesmiddel. Over de Nydelsesmidler, som i det
socialistiske Samfund tilfalde os, kunne vi frit disponere. Vi
kunne opsamle Indtægter til Brug i vore egne private Øjemed
eller for at give Gaver og Hjælp til Andre\footnote[1]{Socialismens Kvintessens. Dansk Overs. p. 84—97. --- Smlgn. ogsaa Gustav Bang: Den socialistiske Fremtidsstat. i København 1903.}. Men det er dog
ikke noget stort Raaderum, den frie Bevægelse her faar. Forbruget
af det Opsparede foregaar med stor Hurtighed, naar jeg
ikke har Lov at gøre det frugtbringende ved at overlade det til
Andre, som savne det. Dersom disse til Gengæld for at faa Lov
at bruge, hvad jeg havde opsparet, gav mig et vist Vederlag for
den Tid, i hvilken jeg overlod dem Brugen, saa kunde jeg anvende
denne Tid til Sysler, som ikke vare umiddelbart produktive,
skønt de kræve alvorligt og langvarigt Arbejde. Og denne
Tid vil være saameget desbedre anvendt, som den maaske kan
bruges til Arbejder, hvis Værdi ingen Anden, særligt ikke de
i Samfundet herskende Myndigheder anerkende, enten fordi det
ikke indses, at de svare til virkelige Fornødenheder, eller fordi
de svare til Fornødenheder, der først skulle vækkes, og hvis
Vækkelse vil gøre Livet rigere. I den socialistiske Stat, der forbyder
al Rente, vil der kun være Plads for saadanne Sysler,
som Statsmagten vil begunstige. Ikke blot vil den individuelle
Produktion falde bort, men den individuelle Nydelse vil blive
holdt indenfor visse snevre Grænser. Der vil ikke kunne danne
sig friere og ejendommelige Retninger og Bestræbelser, naar der
leves af Statsmagtens Forgodtbefindende. Renten er ganske vist
en Indtægt, som Besidderen faar uden selv øjeblikkeligt at pro%
% --- p. 405 ---
\setcounter{footnote}{0}%
\opage{405}ducere den; men den har den sociale Betydning at gøre andre
Virksomheder end den øjeblikkeligt nyttebringende mulige og at
begunstige Sparsommelighed. Den Sparsommelige véd nu, at han
kan opnaå ikke blot at betrygge sin egen Existens, men ogsaa
paa varig Maade at fremhjælpe Interesser og Bestræbelser, han
sætter Pris paa. Renten kan naturligvis misbruges til Lediggang;
men al Dispositionsret kan misbruges: Statens Dispositionsret
kan jo ogsaa misbruges. --- Ligesom Socialismen forudsætter en
Menneskehed, der kan stille fuldkomne Regenter, saaledes forudsætter
den ogsaa en Menneskehed, hvis Virkelyst og Opfindelsesevne
ikke svækkes, fordi dens Initiativ hæmmes, og fordi
dens Fornødenheder bestemmes af Andre.

Gennem det private Initiativ kommer der frisk Blod ind i
det sociale Liv. Den Enkelte maa derfor ikke være absolut udelukket
fra at producere. Det vilde være absurd, om han vel
skulde have Lov til at fortære sine Sparemidler ved Udsvævelser,
men ikke til at bruge dem som Produktionsmidler. Denne
Skranke for Friheden til at forsøge og til at vove, vil aldrig
kunne udholdes. Og det vil ikke blot være en Skranke for den
Enkeltes Frihed, men ogsaa for Samfundets egen Udvikling.
Bureaukratisk og parlamentarisk Selvgodhed vil hæmme Fremskridtet
paa utaalelig Maade. Ikke blot den materielle, men ogsaa
den ideelle Kultur vil lide derved. Gennem Arvelovgivningen
og gennem Statens Overtagelse af private Foretagender af almennyttig
Karakter, naar de et vist Antal Aar ere komne Grundlæggerne
til Gode, og naar de egne sig til at ledes af det Offentlige,
vil det kunne hindres, at der af den for hele Samfundet
saa vigtige private Foretagsomhed og Nydannelse udvikler
sig en Arveadel. Forholdet kunde ordnes i Analogi med de nyere
Ordninger af Forfatter-, Kunstner- og Opfinderrettigheder.

Den spekulative Socialisme vil tilstoppe Kilden til Goderne,
for at Fordelingen kan blive desbedre. Men den hilder sig derved
i en Modsigelse, da der tilsidst ingen Goder vil blive at
fordele, naar Fremskridtets Kilde tilstoppes. Den utopiske Socialisme
var konsekventere: den forudsatte ikke blot en begrænset
eller lukket Stat; men den fordrede ogsaa Fornødenhederne
direkte regulerede. Saaledes f. Ex. Platon og Campanella. Den
% --- p. 406 ---
\setcounter{footnote}{0}%
\opage{406}nyere Socialisme gør den individuelle Frihed den Indrømmelse
at tilstede Privatejendom af Forbrugs- og Nydelsesmidler og
vil kun have Forbud mod at bruge disse som Produktionsmidler.
Jo mere den sociale Erfaring skrider frem, desmere ville
sikkert ogsaa den strenge Socialismes Tilhængere indse, at den
sociale Etik for sin egen Skyld maa fordre og fremme den individuelle
Frihed, ikke blot til Forbrug og Nydelse, men ogsaa
til Arbejde og Produktion. Idetmindste forudsætter Ophævelsen
af Produktionsfriheden som Løsning af det sociale Problem psykologiske
og sociale Forhold, der ere i den Grad forskellige fra
de nu foreliggende, at det vilde være dogmatisk Forvovenhed
at udtale sig afgørende om, hvad de ville gøre muligt eller ikke
muligt.

11. Det er en socialistisk Hovedsætning, at Arbejdet er
Kilden til al Værdi og al Kultur. Gothaerprogrammet begynder
med denne Sætning. Men den indeholder en vis Tvetydighed,
som ogsaa ligger i Ordet Arbejde. Der gives jo ikke blot et
materielt, men ogsaa et aandeligt Arbejde. Og selvom vi ved
Betragtningen af den materielle Kultur især tænke paa det materielle
Arbejde, saa har det dog vist sig, at den forandrede og
forbedrede Stilling, Arbejdet i den materielle Kulturs Tjeneste
indtager i den nyere Tid, hænger sammen med, at den moderne
Industri i saa høj Grad er en Anvendelse af den moderne Videnskabs
Resultater. Det materielle Arbejde forudsætter altsaa
her aandeligt Arbejde, Musklens Arbejde forudsætter Hjernens.
De Forsøg, Tanker og Planer, som aandeligt Arbejde har ført
til, have givet utallige materielle Arbejdere nok at gøre. Man
vil derfor ikke kunne gennemføre en Samfundsordning, naar man
begynder med, som de socialdemokratiske Programmer fra Gotha
(1875) og Gent (1877), at stille Arbejderklassen (hvorved menes
de materielle Arbejdere: Genterprogrammet kalder den Proletariatet)
i Modsætning til alle andre Klasser. Det kan være berettiget
i Kampens Hede; men dersom Følelsen af Modsætning
og Sondring overfor det øvrige Samfund næres for stærkt, spærres
den eneste Vej til bedre Tilstande.

Ikke fysisk Arbejde alene er Kilde til Rigdom og Kultur.
Det største Kulturarbejde er øvet paa det aandelige Omraade.
% --- p. 407 ---
\setcounter{footnote}{0}%
\opage{407}Naturligvis maa de materielle Fornødenheder tilfredsstilles, for
at man kan arbejde. Og meget aandeligt Arbejde behøver materiel
Hjælp. Organisten vilde være hjælpeløs uden Bælgetræder.
Men det er dog ikke Bælgetræderen, der er „Kilde“ til
Musiken. --- Socialismen sætter ogsaa fri Tænken og Forsken
og offentlig Undervisning paa sit Program. Den maa da anerkende,
at hele den aandelige Atmosfære, i hvilken Arbejderen
lever, er af den største Betydning for ham. Den moderne Filosofi
og Socialvidenskab førte først til Ophævelse af de gamle
Lav og dernæst af de Bestemmelser, der hindrede Arbejderne
i at danne lovlige og lovbeskyttede Associationer. Den ensomste
Forsker kan sende Tanker ud i Verden, som ved deres Indflydelse
paa den almindelige Livsopfattelse og den offentlige
Mening kan bestemme Kulturens Gang i langt højere Grad end
mange Tusinders materielle Arbejde. Dette er nu én Gang saa,
og ingen Programmer kunne hindre det. Hvis Fremtidsudviklingen
skal gaa i en sund Retning, maa man stræbe efter at
forminske Afstanden mellem materielt og aandeligt Arbejde.
Men en saadan Bestræbelse kan ikke lykkes, naar Arbejderklassens
Modsætning til alle andre Klasser saa skarpt betones,
som det i Regelen sker. --- Det skal indrømmes, at de „andre“
Klasser, de Klasser, der hidtil næsten udelukkende have været
stillede saaledes, at de have kunnet paatage sig aandeligt Arbejde,
ikke altid have vist den Adfærd overfor Arbejderklassen, som
deres Pligt paabød. Fordomme af forskellig Art, Aandshovmod
og Mangel paa Sympati have hindret dem i at anerkende dens
Ret. Her ligger den vigtigste Skyld for Misforholdet. Men vi
spørge her ikke først og fremmest efter Skylden; vi undersøge,
hvor stor Modsætningen mellem Klasserne i Samfundet faktisk
er og etisk-socialt set bør være.

Hine socialdemokratiske Programmer stride egentligt mod
Marxs Lære, efter hvilken al Klasseforskel skal bringes til at
høre op, ligesom Standsforskellen mellem Adel og Borgere faldt
bort ved Revolutionen. Arbejderklassen staar dog ikke endnu
overfor de andre Klasser, som den tredie Stand ved Revolutionens
Udbrud stod overfor den Tids „højere“ Stænder. Den
Klasse, der nu med Haan kaldes Bourgeoisiet, havde skabt den
% --- p. 408 ---
\setcounter{footnote}{0}%
\opage{408}moderne Industri og Handel, Videnskab og Kunst, og af den
var i de germanske Lande den protestantiske Frihedsbevægelse
fremgaaet. Trods al den Beundring, der maa føles overfor Arbejderklassens
Udvikling i det sidste Aarhundrede, kan det ikke
siges, at den nu har naat et lignende Punkt, som den tredie Stand
havde naat for hundrede Aar siden. Dens Udvikling er endnu
ikke afsluttet. Det er ikke dens egen Skyld; ti Betingelserne
ere først nyligt komne til Stede. Men saa meget des mindre er
det da berettiget at betone Klassemodsætningen saa stærkt, som
Marxismen gør det. I hvert Tilfælde have de modstaaende Klasser
gensidigt Meget at lære af hinanden. Og den Arbejderklasse,
der forhaabentligt engang vil komme til at omfatte Alle, vil
ikke være den nuværende, ufuldkomment udviklede Arbejderklasse,
der ikke blot paa Grund af sine Interesser, men ogsaa
paa Grund af sin ufuldkomnere Udvikling staar i Modsætning
til de andre Samfundsklasser. Ufuldkomment udviklet vil den
vedblive at være, saalænge den endnu --- paa Grund af Klassekampen
--- maa optræde som Masse, uden at der kan foregaa
en fri Individualisering indenfor den. Saalænge man f. Ex. betragter
Forholdstalsvalgmaaden som en rent aristokratisk Indretning,
er det klart, at Forstaaelsen af den personlige Selvstændigheds
Betydning mangler. De dogmatiske Stridigheder indenfor
det tyske Socialdemokrati, der have ført til Udstøden af Personligheder,
hvis intellektuelle Selvstændighed ikke kunde bøje
sig for Programmets ene saliggørende Sandheder, pege i samme
Retning. Hvis det er rigtigt, at den socialdemokratiske Partistyrelse
i Tyskland literært har større Magt og Indflydelse over
sine Partifæller end Statsmyndighederne over deres Undersaatter,\footnote[1]{B. Friedlånder: Die vier Hauptrichtungen der modernen socialen Bewegung. Berlin 1901. I, p. 302. --- Smlgn. ogsaa Bruno Wille: Die Philosophie der Befreiung. Kap. 13 (Parteiherrschaft).}
vidner dette om stor aandelig Umyndighed indenfor Klassen.
De teoretiske Anarkister gøre her et gavnligt Arbejde ved
deres Kritik; til Gengæld betragtes de som Frafaldne og regnes
til „Bourgeoisiet“.

12. Problemet vil dog stedse blive stillet paany, paa Grund
af en Omstændighed, som Socialisterne i Regelen skyde til Side,
% --- p. 409 ---
\setcounter{footnote}{0}%
\opage{409}nemlig det Naturinstinkt, som bevirker en Tendens til en stærkere
Befolkningsforøgelse, end der altid svarer til Næringsmidlernes
Forøgelse. Selvom Malthus skulde have overdrevet denne Tendens,
saa virker der dog her aabenbart en Kraft, som stedse
paany vil føre udover enhver Ligevægtstilstand, man kunde tænke
sig tilvejebragt. (Smlgn. XXV, 2). Hvis det lykkedes at indrette
en Ordning, ved hvilken Enhver trygt kunde se Fremtiden
imøde, vilde denne Tryghed blandt Andet lægge sig for Dagen
ved, at der stiftedes talrige Familier, saa at der snart igen kom
overflødig Arbejdskraft, flere Hænder end der maaske var Arbejde
til, og flere Munde, som kræve Føde. Dette er en simpel
Følge af den samme Aarsag, som gør, at ogsaa nu Ægteskabernes
Antal stiger, ikke blot naar Kornpriserne falde, men allerede,
naar der blot er Udsigt dertil, eller naar en forhaabningsfuld
Stemning raader]).

Det sociale Forhold bestemmes til enhver Tid ved Forholdet
mellem Befolkningens Tilvæxt og den Grad, i hvilken forøget
Kløgt og Energi kunne forøge Jordens Produktivitet. Hvor
forskellige Kræfter virke sammen, vil der indtræde en rytmisk
Bevægelse, idet snart den ene, snart den anden Tendens faar
Overhaand. En Tilværelse, der ikke er underlagt en saadan
rytmisk Skiften, kende vi ikke, og forstaa vi ikke. Det er de
forskellige Kræfter og den Maade, hvorpaa de brydes mod hverandre,
der gør Livet til en Kamp og volde Smerte, især hver
Gang Svingningernes Udslag bliver større. Det er ikke sagt, at
Menneskenaturen bestandigt vil medføre den samme stærke Forøgelsestendens;
men en afgørende Forandring i denne Hen
seende ligger for fjernt til, at den for os kan faa nogen etisk
Betydning. Igennem den haarde Kamp med de givne Vilkaar
% CHECK p. 409: note whose mark was not found
\footnote[1]{I „Efter en god Høst pleje Vielsernes og Fødslernes Antal at tiltage i betydelig Grad, og ligeledes omvendt at aftage efter en slet Høst. I det første Tilfælde er det endnu mere Haabet end den virkelige Besiddelse, der tilskynder til ny Familiestiftelse, hvorfor man ikke finder den stærkeste Tilvæxt ved de absolut laveste Kornpriser, men ved dem, der staa i den mest paafaldende Modsætning til et foregaaende Uaar“. Roscher: Die Grundlagen der Nationalokonomie. § 240.}
% --- p. 410 ---
\setcounter{footnote}{0}%
\opage{410}har Menneskenaturen udviklet sig til sit nuværende Trin; det
er den, der gør de strenge Fordringer om Selvbeherskelse, Klogskab
og Sympati nødvendige; og Vilkaarene ville, saavidt vi
kunne se, ikke foreløbigt forandres. (Smlgn. med det her udviklede
XI, 10; XVII, 2 og XXV, 2).

Vi se derfor ogsaa, at der, naar et socialt Lag har arbejdet
sig frem til gunstigere Forhold, til en højere Levefod, neden
under det danner sig et nyt Lag, hvis Levefod trænger til at
hæves. Efter Landboreformerne i Danmark i Slutningen af det
18. Aarhundrede voxede paa Grund af Udparcelleringerne og
det forøgede Behov af Arbejdskraft Husmændenes og Landarbejdernes
Antal --- men voxede i en Grad, der stillede et
nyt socialt Spørgsmaal i Stedet for det, der nyligt var løst.
„Selv for de Husmænd, der havde Jord, var Stillingen ingenlunde
altid god. Og de Husmænd, der vare jordløse eller vare
Indsiddere, og det var omtrent Halvdelen af det samlede Antal,
de sad næsten altid i meget ringe Kaar; de vare snarere gaaet
tilbage end fremad“\footnote[1]{Falbe Hansen: Stavnsbaandsløsningen og Landboreformerne. I. p. 139. 373---393.\textsuperscript{1} Sidney and Beatrice Webb: History of i Trade Unionism. p.}. Paa lignende Maade omfattede Fag- og
Forbrugsforeningerne fra først af kun de mest intelligente Arbejdere;
de Fag, der ikke kræve nogen videre Specialuddannelse,
bleve ikke organiserede. Arbejderbevægelsen truede en Tid med
at ende i Dannelsen af et Arbejderaristokrati, der --- som det
er blevet sagt --- blev betragtet af de „ulærte“ Arbejdere med
lignende Følelser som Overhuset i den parlamentariske Verden.
Der maatte da en Bevægelse i Gang for at organisere de „ulærte“
Arbejdere\footnotemark[1]. Desuden ville uheldigt stillede Arbejdere naturligt
søge hen til Egne og Lande, i hvilke der allerede er opnaaet
gunstigere Stilling for Arbejderstanden. Der foregaar saaledes
i Europa en stadig Vandring af Arbejderskarer fra Øst til Vest
og fra Syd til Nord; --- og helt ude i den østlige Horizont
true den gule Races flittige og nøjsomme Arbejdere. --- Problemet
vil i lange Tider følge Slægten paa dens Vej; en Løsning
én Gang for alle er usandsynlig.
% CHECK p. 410: note mark without a note

% --- p. 411 ---
\setcounter{footnote}{0}%
\opage{411}13. Den socialistiske Teori er en afgjort idealistisk Teori,
forsaavidt den hviler paa Overbevisningen om den menneskelige
Villies Evne til at rydde alle Vanskeligheder for Dannelsen af
et harmonisk Menneskesamfund af Vejen\footnotemark[1]. Den betragter jo Arbejdet
som en Kilde til al Rigdom og Kultur og ser bort fra
alle Naturaarsager, der kunne fremme eller hæmme Kildens Løb.
Ved denne Idealisme, der ikke lader sit Haab og sin Begejstring
svækkes ved ængstelige Hensyn til de naturligt og historisk
givne Betingelser, er Socialismen i Slægt med det Bedste i det
18. Aarhundredes Tankegang, skønt den paa den anden Side
ved sin Anti-Individualisme betegner en Reaktion mod det 18.
Aarhundrede. Saa mange teoretiske Fejltagelser og Illusioner,
man end maatte kunne paavise i Socialismen, saa fremtræder
den dog i Praxis som en af de betydningsfuldeste etisk-sociale
Bevægelser i Nutiden. Den har forstaaet at vække og begejstre
Arbejderne; den har rettet deres Tanker mod Idealer og Opgaver,
der omfatte langt mere end den snevre Kreds, i hvilken
den isolerede individuelle Selvopholdelsesdrift bevæger sig. Den
har stillet Millioner af Mennesker et stort fælles Ideal, der løfter
den Enkelte op over den snevre Interessekreds, rettet Blikket
mod et Fremtidsbillede, der viser Guldalderen foran os, ikke
bagved os; --- og den har samtidigt formaaet at faa et fast Tag
ide rent praktiske og materielle Interesser, at bruge disse som

) I en særligt idealistisk Form fremtræder Fremtidsidealet i Syndikalismens
Teori. Ikke blot Kapitalismen, ogsaa Staten skal falde bort.
Et økonomisk Forbund af frie Værksteder skal afløse Staten. Arbejdsglæde,
Pligtopfyldelse og Sympati skulle træde istedetfor Autoritet og Tvang.
Gennem kunstnerisk og moralsk Udvikling skulle Arbejderne modnes til
at kunne virkeliggøre dette Ideal. Syndikalismen stiller sig baade i Modsætning
til de sædvanlige Fagforeninger og til Socialismen, der vil gøre
Staten til eneste Arbejdsgiver. Smlgn. (foruden de ovenfor (4) nævnte
Skrifter af H. Scavenius) Challaye: Le Syndicalisme revolutionnaire.
(Revue de Métaphysique et de Morale. 1907). Paul Louis: Le Syndicalisme
contre l'état. 1910. --- Ved Forkastelsen af Staten minder Syndikalismen
om Utopier af Sibberns Art (i Modsætning til Platons, More's
og Campanellas); ved sin Betoning af Sympatien minder den 'om den
filantropiske Socialisme; ved den Betydning, den tillægger Striken (tilsidst
Generalstriken) som Virkemiddel, er den trods al Modsætning beslægtet
med Marxismen. (1913). b
% CHECK p. 411: note mark without a note

% --- p. 412 ---
\setcounter{footnote}{0}%
\opage{412}Motiver, der kun føre den Enkelte ud over ham selv. Uden
et stort Fremtidsbillede kan ingen social Bevægelse gaa for sig.
Undertiden --- saaledes for faa Aar siden i England --- var
Fagforeningsbevægelsen ved at gaa i Staa, og det var da den
begejstrede Tro paa det socialistiske Ideal, der vakte Driften ti
Fremgang igen. Fra oven kan Udviklingen ikke ledes; selvom
der kan komme frugtbare Tankespirer og Impulser fra de andre
sociale Lag, maa Selvvirksomheden fremfor Alt vækkes, og dertil
behøves Idealer, som staa i naturligt Forhold til den Trang,
som føles, og til den intellektuelle og moralske Horizont, der
staar aaben. Hvis de andre sociale Lag ikke synes om disse
Idealer, og hvis de finde Trangen for lav og Horizonten for
snever, saa er det deres Opgave at ændre Livsvilkaar og Samfundsforhold
saaledes, at Trangen kan højnes, Horizonten udvides,
og dermed ogsaa Idealet faa en anden Karakter. Den
socialistiske Bevægelse er en Fortsættelse af den store Frihedsbevægelse,
som mange Liberale ventede afsluttet med Menneskerettighedernes
Erklæring, med Parlamentarisme og Frihandel.
Den er den positive Side af den store Emancipationsproces,
der begynder i det 18. Aarhundrede. Den Interessernes
Harmoni, som ældre Nationaløkonomer og Utilitarianere troede
paa, søger den at gøre til en Kendsgerning. Og samtidigt vil
den republikanisere og demokratisere den Menneskelighedsfølelse,
som det gamle Samfund væsentligt kun kendte i Form
af Herrens patriarkalske Forhold til de Afhængige. Den vil
udvikle Menneskekærligheden til Retfærdighed. Skyggesiderne
og Illusionerne i Fremtidsbilledet ville kunne rettes under den
videre Udvikling. Socialismen maa lære at skelne mellem Ideal
og Virkelighed. Idealet behøver ikke at miste sin drivende Kraft,
fordi man ikke tilslører og overser de virkelige Forhold. Og
den maa navnligt lære, at hvis den har en Fremtid, kommer
denne væsentligt gennem det frie Arbejde og den frie Association,
selvom Staten meget vel i langt højere Grad, end vi nu
ere vante til at tænke os, kan træde beskyttende, hjælpende,
udlignende og opdragende til.

14. Som allerede berørt (9), griber Statsmagten uafbrudt
ind i den sociale Udvikling, selv hvor man ikke 'I bliver sig det
% --- p. 413 ---
\setcounter{footnote}{0}%
\opage{413}tydeligt bevidst. Der gives ikke en eneste Side af den politiske
Ordning (Forfatning, Administration, Finans-, Rets- og Militærvæsen,
Kirke- og Undervisningsvæsen), som ikke paa en eller
anden Maade bliver af bestemmende Indflydelse paa Arbejdsforholdenes
Ordning. Et stort Fremskridt vilde det allerede være,
om man blev sig dette tydeligere bevidst, saa at denne saare
væsentlige Side ved de politiske Spørgsmaal blev dragen mere
frem. Det er dog de politiske Spørgsmaals sociale Betydning,
som først giver dem virkelig Interesse. Hvor denne Betydning
helt' er tabt af Syne, er den politiske Strid enten en rent personlig
Strid om Magten eller en Strid om blotte Formaliteter.
--- Vi skulle her nævne nogle Punkter, hvor Staten uden at
krænke Frihedsprincipet kan udrette Meget i Henseende til Arbejdets
Ordning og Udbyttets Fordeling, og hvor den ogsaa
allerede --- ved en ubevidst eller uvilkaarlig Socialisme --- har
grebet ind.

a) Allerførst maa fremdrages Fordringen om, at der vises
simpel Retfærdighed. De store Tyve ere endnu i mange Henseender
heldigere stillede end de smaa. Man ser med mistænksomme
Blikke paa Bevægelser i Arbejdernes Kredse, som man
i andre Kredse af Samfundet ikke vilde lægge Hindringer i
Vejen. Den Modstand, som er bleven gjort mod Arbejdernes
Ret til personlig Frihed og deres Forsamlings- og Foreningsret,
har ikke kunnet stemme dem venligt overfor de herskende Stænder.
Adam Smith klagede i sin Tid over, at Mestrene maatte
indgaa Forbund om at nedsætte Arbejdslønnen, medens Arbejderne
ikke maatte indgaa Forbund om ikke at arbejde under
en vis Betaling\footnote[1]{Om Nationaivelstand. Dansk Overs. I, p. 202 f.}. Den tyske Lovgivning anerkendte paa lignende
Maade endnu i den nyeste Tid Arbejdsgivernes Lavsforeninger,
men ikke Arbejdernes tilsvarende Foreninger, og den
megen Tale om at genoprette de gamle Lav gaar i Virkeligheden
ud paa at gøre Arbejderne afhængige af Arbejdsgiverne -).

b) Staten har erkendt det som sin Pligt at beskytte Arbejdernes
Frihed, Sundhed, Sikkerhed og Moralitet overfor 
% CHECK p. 413: note whose mark was not found
\footnote[2]{L. Brentano: Die gewerbliche Arbeiterfrage. (Schonbergs Handbuch. 1. Udg. I). p. 931 f. 970.}Ar%
% --- p. 414 ---
\setcounter{footnote}{0}%
\opage{414}bejdsgivernes Vilkaarlighed. Efter en haard Kamp gik --- først
i England --- den saakaldte Fabriklovgivning igennem, der regulerer
Sikkerheds- og Sundhedsforanstaltninger i Fabriker og
Bjergværker, bestemmer Arbejdstidens Længde, fastsætter Grænser
og nærmere Bestemmelser for Kvinders og Børns Arbejde,
forbyder Arbejdslønnens Udbetaling paa Værtshuse og fordrer
den udbetalt i rede Penge\footnote[1]{The Utilitarians. III. p. 175 —177. il}. --- Ligesom det har vist sig, at
det frie Arbejde er mere produktivt end Slavearbejdet, saaledes
har det ogsaa vist sig, at de gunstigere Betingelser, som denne
Lovgivning, der efter engelsk Mønster ogsaa blev gennemført i
andre Lande, har skaffet Arbejderne, langt fra at skade Produktiviteten
netop har forøget denne. Den har ikke blot, som
selv \textsc{Karl} \textsc{Marx} indrømmer, bevirket „Fabrikarbejdernes fysiske
og moralske Genfødelse“; men den større Friskhed og
Kraft, hvormed der nu kunde arbejdes, forøgede endog ofte Arbejdets
Produktivitet. Arbejdsgiverne selv fandt deres Regning
derved og opgave deres tidligere Modstand mod, hvad de ---
karakteristisk nok --- havde kaldt Indgreb i deres personlige
Frihed! --- Særligt har Indskrænkningen af Arbejdstiden været
et stort Gode, og det baade for Arbejdsgivere og for Arbejdere.
Som de engelske Fabrikinspektører sige i en Beretning: forhen
havde Herren ikke Tid til at tænke paa Andet end Penge, og
Arbejderne ikke Tid til at tænke paa Andet end Arbejde. Den
lange Arbejdstid gjorde Arbejderne til rent fysiske Væsener,
idet al deres Fritid gik med til Søvn og Hvile for at faa Kraft
til at tage fat paany. Man nærmer sig ved den successive Begrænsning
af Arbejdstiden til \textsc{More's} Utopia, hvor „Forfatningens
vigtigste Formaal er at regulere Arbejdet efter Folkets
Fornødenheder, saaledes at der kan blive Tid tilovers til Udvikling
af Aanden, hvori Utopianerne mene, at Livets Lykke
bestaar“. \textsc{Fichte} har med Rette sagt, at et Folks sande Rig-

') Gneist: Das Self-Government in England. 3. Aufl. p. 314 f. ---
L. Brentano: Die gewerbl. Arbeiterfrage. (Schonbergs Handbuch I). p.
973. --- K. Marx: Das Kapital. 2. Aufl. I. p. 224—314. --- Karakteristisk
er den Modstand, Utilitarianerne gjorde mod Fabriklovgivningen,
undtagen forsaavidt den gjaldt Beskyttelse, af Børnene. Leslie Stephen:
% --- p. 415 ---
\setcounter{footnote}{0}%
\opage{415}dom bestaar i den Fritid, det raader over\footnote[1]{p. 4. --- Fordelingen 8 Timers Arbejde, 8 Timers Søvn og 8 Timers Fritid er allerede foreslaaet af Comenius, og senere af Hu fe land.}. Ti den Tid, der
kan afses fra Arbejdet for de materielle Fornødenheder, kan
anvendes til en højere og friere Udvikling, til ædel Livsnydelse,
til at virke i den ideelle Kulturs Tjeneste og føle sig som Mere
end et Hjul i et stort Maskineri. --- Det kommer naturligvis
an paa, hvorledes Fritiden benyttes, og man kan ikke undre
sig over, at den ikke lige strax eller altid benyttes saa godt
som muligt. Der maa vækkes Drift og skaffes Midler til højere
Udvikling. Her bliver Undervisningsvæsenet, baade det af Staten
ledede og det af frie Kræfter besørgede, af største Betydning,
og ikke mindst Alt, hvad der kan gøres for at udbrede
og tilfredsstille Sansen for Kunst- og Naturskønhed. Gennem
Deltagelse i det politiske Liv udvikles Følelsen af borgerlige
Pligter og Rettigheder og af at arbejde med i det hele Samfunds
Tjeneste. Man kan ikke vente med at begrænse Arbejdstiden,
til der er opnaaet Sans for at bruge Fritiden godt; ti denne
Sans opstaar først, naar Fritiden er der. Og naar man klager
over, at Arbejderne ikke forstaa at bruge deres Fritid, saa er
der langt mere Grund til at klage over den Maade, paa hvilken
de velhavende Klasser bruge deres Fritid. Naar man ikke er
vant til at have Fritid, er det intet Under, at man ikke har
lært at bruge den; men det er sørgeligt at se, hvor tomt og
lavt ofte det er, hvormed de saakaldte dannede og højere Klasser
udfylde deres rigelige og ofte ufortjente Fritid. Arbejderne
maa netop have større Fritid, for at den Maade, paa hvilken de
nu ofte misbruge Fritiden, kan falde bort. Den Otte-Timers Arbejdsdag
forlanges, for at den „blaa Mandag“ kan falde bort\footnotemark[1],
c. Staten kan fremdeles øve en udlignende Indflydelse
paa de sociale Modsætninger gennem Ordningen af Finans- og

') System der Rechtslehre (1812). Nachgelassene Werke. II, p. 543:
„Das Eigenthum des Ganzen „ist] die Musse, die allen nach vollbrachter
Arbeit bleibt“. --- Smlgn. L. Brentano: Die Arbeitergilden der Gegenwart.
II. p. 356: „Die Frage nach der Lange des Arbeitstages ist eine
Frage nach dem Stand der Civilisation“.

 Robert Seidel: Der achtstundige Arbeitstag. Zurich 1896.
% CHECK p. 415: note mark without a note

% --- p. 416 ---
\setcounter{footnote}{0}%
% CHECK p. 416: notes paired with the marks in order (the layer misread a number)
\opage{416}Skattevæsenet. Man har endog\footnote[1]{Adolf Wagner: Direkte Steuern. (Schonbergs Handbuch. 1. Udg. II). p. 169. 259. Smlgn. ogsaa Jhering: Der Zweck im Recht. 2. Aufl. I, p. 533.} dateret en ny Periode i Skattevæsenets
Historie fra den Tid af, da man begyndte at fremhæve
denne udlignende og fordelende Virkning af Skatteordningen
som et Hovedsynspunkt. Det er den „socialpolitiske“ Skatteteori,
som anlægger dette Synspunkt ved Siden af det blot
finansielle. \textsc{Henry} \textsc{George} og hans Tilhængere vente endog
det sociale Spørgsmaals fuldstændige Løsning gennem Beskatning
af Jordrenten. Uden Staten vilde Jordrenten ikke existere, og
derfor er den det rette Skatteobjekt\footnote[2]{Smlgn. dog Westergaard: Nationaløkonomien i Hovedtræk. (1908). p. 84---89.}. --- Herhen hører ogsaa
en stærkere Beskatning af Arv og Afskaffelse af Arv gennem
Sidelinier.

d) Et Skridt videre fører Spørgsmaalet om, hvorvidt Staten
skal oprette Forsikringsanstalter for Arbejderne. Dersom
det skal være en Tvangssag for Arbejderen at benytte sig af
saadanne Anstalter, berøves han let sin personlige Uafhængighed.
Ti naar Staten ikke kan garantere ham stadigt Arbejde
for en vis Løn, kan en Svingning i Forholdene gøre ham arbejdsløs
og derved ude af Stand til at betale Indskud, hvorved han igen
mister hvad han tidligere har indskudt. Det vil da staa i Arbejdsgiverens
Magt at paalægge ham Betingelser, han ikke vilde
finde sig i, naar han ikke frygtede for at miste sine tidligere
Indskud\footnote[3]{L. Brentano: Die gewerbl. Arbeiterfrage. (Schönbergs Handbuch. 1. Udg. I). p. 985 f.}. Staten vil sikkert her, som paa saa mange andre
Punkter, virke bedst ved at virke indirekte, saaledes at den
støtter og kontrollerer de Organisationer, der have udviklet sig
ved Arbejdernes egen frie Sammenslutning.

e) Der er mange Virksomheder, baade paa Produktionens
og paa Omsætningens og Samfærdselens Omraade, som Stat og
Kommune allerede have overtaget og ville kunne overtage efterhaanden.
Hvor Grænsen ligger i denne Henseende, kan kun
fremskridende Erfaring afgøre. Her har netop den empiriske
% --- p. 417 ---
\setcounter{footnote}{0}%
\opage{417}Soc.alisme sin store Betydning. Men for at kunne benytte Erfaringen
maa man have sikre Kendsgerninger. Og her er et
Omraade, hvor Staten kan gøre uberegnelig Nytte: ved at tilvejebringe
en nøjagtig Statistik. De nordamerikanske Fristater
ere i denne Henseende gaaede i Spidsen. Baade i de enkelte
Stater og for hele Unionens Vedkommende er der oprettet
Anstalter for Arbejderstatistik. I den Lov, ved hvilken i Aaret
1888 the United States Department of Labour oprettedes, angives
dets Opgave at være den, at samle og i de forenede
Staters Befolkning udbrede Oplysninger „angaaende Arbejderspørgsmaal
i dette Ords mest omfattende Betydning, særligt om
Arbejdets Forhold til Kapitalen, om Arbejdstiden, om mandlige
og kvindelige Arbejderes Fortjeneste og om Midlerne til at
fremme Arbejdernes materielle, sociale, intellektuelle og moralske
Den sociale Lovgivning og den offentlige Debat i
de forenede Stater har allerede høstet Gavn af denne Foranstaltning\footnotemark[1].

Den nærmere Udvikling og Begrundelse af de her nævnte
Punkter (og hvilke flere der kunde føjes til) tilkommer Nationaløkonomien.
Her ere de blot anførte, fordi de alle forudsætte
en etisk-social Opfattelse af Staten og dens Virksomhed og
en Erkendelse af, at Statshjælp og Selvhjælp ikke udelukke
hinanden. Det kommer kun an paa at anbringe det rette Forhold
mellem dem. De frie Kræfters Omraade er det egentligt
produktive, det, hvor Initiativerne trives; Staten kan kun yde
Beskyttelse, Form og materiel Støtte for det, der selv har arbejdet
sig frem. Forholdet mellem Statshjælp og Selvhjælp bør
netop være det omvendte af det, som den spekulative Socialisme,
i mærkelig Overensstemmelse med Bureaukrati og Absolutisme,
fastsætter. Og da Erfaringerne gøres bedst i mindre
Kredse, vil kommunal Ordning af Produktionen*) --- som ogsaa
den empiriske Socialisme hævder --- --- have mange Fortrin fremfor
den egentlige „Statssocialisme“. Man er ogsaa fra social-Tidsskri\textasciicircum{}isi)\footnotemark[8]\footnotemark[6]“
\footnotemark[1]\footnotemark[3]\footnotemark[1]\footnotemark[1]“ "“*“****** "* Stastistiken --- (Nordisk

Tidsskrift. 1893). Smlgn. F. Pio: Kommunesocialisme. (Tilskueren. 1903).
% CHECK p. 417: note mark without a note

% CHECK p. 417: note mark without a note

% CHECK p. 417: note mark without a note

% CHECK p. 417: note mark without a note

% --- p. 418 ---
\setcounter{footnote}{0}%
\opage{418}istisk Side begyndt at komme bort fra at tænke sig Staten ---
efter den Betydning, i hvilken vi nu tage dette Ord --- som
den Samfundsform, der skal overtage hele Produktionen. Man
tænker sig Organisationer af højst forskellig Art og Omfang,
alt efter Opgavernes Art, og man tænker sig en Samvirken
mellem disse forskellige Associationer\footnote[1]{Gustav Bang: Den socialistiske Fremtidsstat.  p. 64.}. Sandsynligvis maa
denne Betragtning, der --- i Forbigaaende sagt --- fører bort
fra den egentlige Marxisme, udvides saaledes, at det frie,
individuelle Initiativ faar sin Plads ved Siden af det Selvstyre,
man indrømmer de smaa og store Organisationer. Der maa
være en stadig Vexelvirkning mellem de Enkelte, de frie Kultursamfund
og Staten, for at Udviklingen skal blive sund og
kraftig.

15. Den spekulative Socialisme ophæver ikke absolut Privatejendom,
men indskrænker den til Brugs- og Nydelsesmidler:
Der vilde kun være fuldstændig Modstrid mellem Privatejendom
og Socialisme, dersom man antog Ejendomsretten for ubetinget.
Men en ubetinget Ejendomsret kan ikke paavises noget Steds i
Historien. Ideen om den ubetingede Ejendomsret er en Fantasi,
som Individualismen opstillede mod en vilkaarlig Statsmagts
Luner. Til enhver Tid har Statsmagten grebet ind i den Enkeltes
Besiddelse, naar fælles Interesser syntes at kræve det.
Grænserne for disse Indgreb ere skiftende, ændre sig med de
historiske Forhold.

Selve Begrebet Privatejendom i den Udstrækning, i hvilken
vi nu tage det, er forholdsvis nyt. Paa de primitive Kulturtrin
hersker Fællesskabet. Jagtgrundene ere fælles for hele Stammen,
og hos agerdyrkende Stammer faar hver Familie aarligt sin Lod
anvist, som Familiens Medlemmer dyrke i Fællesskab. Det vilde
naturligvis være urigtigt at tænke sig dette Fællesskab som en
Kommunisme i egentlig Betydning; det betyder kun, at Trangen
til gennemført Fordeling endnu ikke er opstaaet, og at den Enkelte
umiddelbart virker sammen med Familien eller Horden,
ikke endnu har nogen Anledning til at forsøge paa at skille sin
Existens fra den. Paa den anden Side vilde man ved at karak%
% --- p. 419 ---
\setcounter{footnote}{0}%
\opage{419}terisere den «primitive Kommunisme“ som en Assurance, det
primitive Fællesskab som „en rumlig Samværen af individuelle
isolerede Interesser“\footnote[1]{B. W. Leist: Alt-Arisches Jus Civile. I. Jena 1892. p. 512 f,} overføre moderne Begreber, der skylde
Individualismens Mekanik deres Oprindelse, paa de gamle sociale
Forhold. Man kan ikke sammenligne det gamle Fællesskab
med et moderne Aktieselskab, der netop er Udtrykket for, at
hidtil isolerede Interesser mødes for rent udvortes og upersonligt
at virke sammen. Hvad særligt Jordejendom angaar, er det
af Betydning, at den sociale Organisation i Klan og Familie i
Tiden gaar forud for Bosættelse og stadig Jordbesiddelse\footnotemark[1]. Den
første egentlige Privatbesiddelse er det Vildt, den Enkelte nedlægger,
den Frugt, han høster, de Redskaber og Klæder, han
selv skaffer sig. Den rørlige Ejendom gaar altsaa forud for den
urørlige. Først senere gaar Jord over til Privatejendom. Nødvendigheden
af regelmæssig, efter en omfattende Plan gennemført
Dyrkning af Jorden, Arbejdets Deling, den vaagnende Selvfølelse
og Selvstændighedstrang hos de enkelte Individer (smlgn.
XI, 14 og XXIII, 3), alt dette var Aarsager, der maatte medføre
en Individualisering af Ejendommen. Den primitive Kultur
kender hyppigst kun en fælles Brugsret indenfor Stammen eller
Familien, og en Ejendomsret kun dér, hvor Stammer eller Familier
støde sammen. Først med Statens Udvikling opstaar den
individuelle Ejendomsret som den af Statsmagten anerkendte og
beskyttede Besiddelse af materielle Goder. Men selve denne
Anerkendelse og Beskyttelse, som er et uundværligt Moment i
Ejendommens Begreb, viser, at Ejendomsretten ikke er ubetinget.
Statsmagten knytter stedse bestemte Betingelser til sin Anerkendelse
og sin Beskyttelse. Den sætter Skranker for Magten
til at skalte og valte med Besiddelsen. Man har ikke Lov at
bebygge sin Grund, dyrke sin Jord eller borttestamentere sine
Penge, som man selv vil, og man yder af, hvad man besidder,
større eller mindre Bidrag til fælles Formaal. Den Magt, som
ved sin Anerkendelse og Beskyttelse gør den blotte Besiddelse
til Ejendom, er den samme, som paalægger Skranker og Pligter;

*) C. N. Starcke: Samvittighedslivet. I. p. 235.
% CHECK p. 419: note mark without a note

% --- p. 420 ---
\setcounter{footnote}{0}%
\opage{420}disse fremtræde altsaa ingenlunde som Indgreb i en fra først af
absolut Ejendomsret.

Det er tilsidst etiske Hensyn, der begrunde baade Anerkendelsen
og Indskrænkningen. --- En Bestemmelse af Grænserne
for de enkelte Individers Raadighed over materielle Goder er
for det Første nødvendig, for at Staten skal kunne løse sin mest
elementære Opgave, nemlig at hævde Freden og Sikkerheden.
Det er hertil ikke nok at fastslaa Grænsen mellem Stammers
eller Familiers Raaderum. Eftersom Udviklingen skrider frem,
faar Staten at gøre med de enkelte Individer og deres indbyrdes
Forhold og maa have Midler til at konstatere, hvor der har
fundet Overgreb Sted imellem dem. Dersom der f. Ex. hos en
hidtil nomadiserende eller af Plyndringstog levende Stamme begynder
at udvikle sig Agerdyrkning, vil denne fra først af være
Enkeltes Sag. Men jo Flere der slaa ind paa denne Vej, des
lettere vil der indtræde Sammenstød, idet Flere vilde dyrke
samme Stykke Jord. De Mægtigste vilde søge at lægge Beslag
paa de største Stykker Jord; men herimod vilde saa den store
Mængdes Interesse rejse sig, da man var vant til at benytte
den fælles Jord til Græsning for hele Stammens Kvæg. Kun
en bestemt Fordeling og Afgrænsning fra Statsmagtens Side vilde
her kunde skaffe Fred og Sikkerhed til Veje\footnote[1]{D. M. Wallace: Rusland. Dansk Bearbejdelse. II. p. 57—61.}. --- --- Dernæst er
anerkendt og beskyttet Besiddelse en Betingelse for stadig, selvstændig
og opofrende Virksomhed for Kulturens Formaal. At
besidde Ejendom er at have et socialt Tillidshverv. Vel tilfredsstilles
ogsaa den Enkeltes Selvstændigheds- og Magtfølelse derved;
men den sociale Etik kan ikke betragte den Begrundelse
som' tilstrækkelig, der blot betragter Ejendommen som Personlighedens
Forlængelse ud i Omverdenen. En saadan Forlængelse
eller Breden er naturligvis berettiget, hvor den ikke fører
til Strid og Sammenstød. Men den positive, sociale Betydning
af den træder først frem, naar den Enkeltes Selv- og Magtfølelse
tages i Samfundets Tjeneste. Ikke af den individuelle
Personligheds Begreb i og for sig, men af den sociale Nødvendighed
af at have saa mange frie, aktive Udgangspunkter for
% --- p. 421 ---
\setcounter{footnote}{0}%
\opage{421}Kulturvirksomheden som muligt følger Privatejendommens etiske
Karakter. Den Enkelte har end ikke selvindlysende Ret til det,
han frembringer og erhverver ved eget Arbejde; ti dette Arbejde
er ikke en Skabelse af Intet; det er ikke blot betinget
ved hans egen Evne og Villie; --- at det kunde øves, og at det
kunde blive frugtbringende for ham, skyldes Samfundets Beskyttelse
og Hjælp. Som han altsaa med Taknemmelighed og Pietet
maa føle, hvad han skylder Samfundet, maa han ogsaa, naar
han ser Tingene fra det etiske Synspunkt, opfatte sin Besiddelse
som Middel til at øve en Virksomhed i Samfundets Tjeneste.
Han er ligesaavel Embedsmand som de, Staten særligt beskikker,
og har sit etiske Ansvar for Anvendelsen af, hvad han ejer.
Etisk set er der derfor ogsaa snævrere Grænser for Benyttelsen
af Ejendommen, end der er fra den ydre Retsordens Synspunkt.
De materielle Goder, som ere i den Enkeltes Besiddelse, skal
han gøre saa frugtbringende --- ikke blot for sig selv, men for
Slægten --- som muligt. De tilhøre ham nu, men de kunne
senere blive Andre til Del. Han har derfor f. Ex. --- selvom
Lovgivningen ikke altid kan hindre det --- ikke Ret til at udpine
Jorden ved Rovdrift og derved forringe dens produktive
Evne for Fremtiden. Det er en egoistisk Hensynsløshed overfor
de følgende Slægter. For en øjeblikkelig personlig Fordels Skyld
skyder man Slægtens blivende Gavn til Side. Ofte er det dog
ligefrem Uvidenhed, der ligger til Grund, og vi have da et interessant
Exempel paa, hvorledes forøget Indsigt kan forøge de
etiske Pligter. Det var først Liebigs Lære om Planternes Trang
til Mineralstoffer, som godtgjorde Rovdriftens Skadelighed og
paaviste Slægtens Solidaritet i Henseende til Jordens Brug.

Meget, som nu er Privatejendom, kan tænkes en Gang i
Tiden at blive Fællesejendom. Det kunde f. Ex. tænkes, at det
viste sig hensigtsmæssigt og muligt at inddrage al Jordejendom
under Staten, for at Jordrenten kunde komme hele Samfundet til
Gode. Hvis det da skete, vilde det ikke være noget Overgreb;
det vilde være i Konsekvens af den Tankegang, paa hvilken
Privatejendommens etiske Begrundelse beror. Hvad der taler
mod denne Tanke, er, at der sandsynligvis bringes større Udbytte
ud af Jorden ved privat Dyrkning end ved Statsdrift.

% --- p. 422 ---
\setcounter{footnote}{0}%
\opage{422}16. Privatejendom af materielle Goder gør Handel, ombyttende
Virksomhed nødvendig. I det socialistiske Samfund vilde
der ingen Handel være: fra Centralmagtens Sæde (i hele Staten
eller i den enkelte Kommune) vilde Produkterne og Nydelsesmidlerne
fordeles omkring til de enkelte Individer, men der vilde
ikke være nogen mellem de utallige enkelte Individer foregaaende
Byttevirksomhed. En Centralisation af Byttevirksomheden
vil nu ogsaa være i Samfundets Interesse, idet overflødige
Mellemhænder derved kunne undgaas. Den stærke Blomstring
af Forbrugsforeningerne (XXVI, 6) viser, at der er langt
flere overflødige Mellemhænder, end man paa Forhaand skulde
tænke. Da Samfundet ikke existerer for Handelens Skyld, men
Handelen for Samfundets Skyld, er en Bestræbelse som Forbrugsforeningernes
fuldstændigt berettiget og er et væsentligt Led i
den empiriske Socialisme. Men efter de overflødige Mellemhænder
bør man ikke dømme om Handelens Betydning. Handel be
ror paa Virksomhed af de frie, individuelle Kræfter, der her vise
sig nok saa værdifulde som ved Produktionen. Det gælder for
den, der sidder inde med Goder, han ikke selv bruger, at finde
et Sted, hvor der er Trang til dem. Handelsaanden bestaar
langtfra blot i den Kunst at sælge for Mere, hvad man har købt
for Mindre; den væsentligste Side ved den er Evnen til at finde,
hvor der er Fornødenheder at tilfredsstille. Der udføres et produktivt
Arbejde, naar der udfindes et Sted, hvor et Produkt kan
blive til Nytte, og naar Produktet flyttes derhen fra et Sted,
hvor det ikke er til Nytte. Socialistiske Teoretikere have ikke
villet anerkende dette. Og dog er det lige saa klart, som at det
er produktivt Arbejde, Bonden foretager, naar han kører sit
Korn til Købmanden i Stedet for at lade det ligge paa Laden
eller staa paa Marken. Det er naturligvis Egeninteresserne, som
fra først af sætte hin Evne i Virksomhed, og dermed hænge de
uheldige Sider ved Handelen og de Handlende sammen. Naar
man, især i tidligere Tider, saaledes de græske Filosofer, de
kristelige Kirkefædre og Skolastikere, saa ned paa Handelen som
en umoralsk Erhvervsvirksomhed, da var det, fordi den mentes
at frembringe Smaalighed, Pengegriskhed, Troløshed og Hjerte%
% --- p. 423 ---
\setcounter{footnote}{0}%
\opage{423}løshed\footnote[1]{Platon: Lovene. Fjerde Bog. p. 705 f.}. Man oversaa den store etisk-sociale Betydning, Handelen
har ved at knytte Menneskene sammen gennem deres Interesser.
Handelen er et Udtryk for, at det enkelte Menneske
og det enkelte Folk ikke er sig selv nok, hvorfor de af fri Tilskyndelse
søge Supplering. Den socialistiske Stat, der skal udelukke
Handel, maa enten (som \textsc{Fichte} tænkte sig den) være en
lukket Stat, der i økonomisk Henseende er sig selv nok, eller
(som \textsc{Rodbertus} tænkte sig den) en Verdensstat, der omfatter
alle Folkeslag. Indtil en saadan Verdensstat har udviklet sig,
maa Handel være en Nødvendighed, og Handelen vil endog
være et Hovedmiddel for dens Udvikling. Og hvad den lukkede
Stat angaar, saa skulde denne i Følge Fichte kun staa i intellektuelt,
ikke i økonomisk Samkvem med andre Stater. Videnskaben,
ikke Handelen skulde betinge Sammenhængen mellem
Menneskene. Men den ideelle Sammenhæng afhænger her netop
af den materielle; denne baner Vejen for hin. Naar Baandet
en Gang er knyttet af økonomiske Grunde, kan der derved være
forberedt en videregaaende Forbindelse. Allerede en varig Handelsforbindelse
forudsætter et Tillidsforhold, som bygger ikke blot
paa de egoistiske Interesser, men ogsaa paa Karaktererne. Livlig
Vexelvirkning og faste Baand mellem Individer af samme
Folk og mellem forskellige Folkeslag knyttes paa denne Maade.
Handelen har derfor spillet en stor Rolle i Kulturens Historie.
Den har indledet Forbindelse og bragt gensidigt Kendskab til
Veje mellem Mennesker, der ellers ikke vilde være komne i
Berøring med hverandre. Den har ofte udvidet den snevre Horizont,
indenfor hvilken Blikket bevæger sig, før Trangen til og
Muligheden af Vexelvirkning med fjernere Dele af Menneskeheden
er fremkommen. Især hvor der finder Søhandel Sted med
fjerne Lande, fremmes et større Blik paa Livet og dets Forhold;
mangehaande nye Tanker og Planer røre sig; nye Sæder og
Indretninger blive bekendte, og der indledes en Frigørelse fra
det Tilvante og det Overleverede. Med Skibene og Karavanerne
følge de usynlige Passagerer: gennem den materielle Forbin%
% --- p. 424 ---
\setcounter{footnote}{0}%
\opage{424}deise bringes ogsaa Menneskenes Følelser og Tanker i Forbindelse.
Den materielle Omsætning faar derved ogsaa Betydning
for det aandelige Liv. (Smlgn. XXIV, 1).

\grouphead{2. DEN IDEELLE KULTUR.}

\chapterhead{XXVII.}{MATERIEL OG IDEEL KULTUR.}

1. Den materielle Kultur gaar nærmest ud paa Tilvejebringelse
af et stort System af Midler. Men til Midler hører
Formaal. Den materielle Kultur peger derfor stedse ud over sig
selv. Det sociale Spørgsmaal kommer jo netop frem derved, at
Midlerne synes at brede sig paa Formaalenes Bekostning, idet
en stor Del af Menneskene synes at være henviste til at holde
Livet oppe, men undvære det, der først gør Livet værdt at leve.
Vi synes at stønne under Vægten af et Apparat, som dog er
bragt til Veje for at gøre det lettere at leve. Den etiske Undersøgelse
af den materielle Kultur gik ud paa Muligheden af at
komme ud over denne Mislighed og at hævde den Grundtanke,
at en menneskelig Personlighed aldrig maa betragtes blot som
Middel. Allerede denne Grundtankes Anvendelse fører ud over
den materielle Kultur, idet den viser os, at al Kultur skal tjene
Udviklingen af det personlige Liv. I den ideelle Kultur, som
bestaar i den frie Udvikling af Tænkning, Fantasi og Følelse,
er Personligheden Mere end Middel; det er Personlighedens egne
Kræfter, som her sættes i Virksomhed, og det ikke af anden
Grund, end fordi der er en umiddelbar Trang til at bruge dem.
I Kraft af en højere Selvopholdelsestrang udfolde Tanke, Fantasi
og Følelse sig og skabe ejendommelige Former, i hvilke de faa
% --- p. 425 ---
\setcounter{footnote}{0}%
\opage{425}Luft og Udtryk. Her er intet fjernere Maal, som skal naas: den
klare Forstaaelse, det levende Fantasibillede, den inderlige og
dybe Følelse have deres Værdi i sig selv.

Forsaavidt kunde Forholdet mellem materiel og ideel Kultur
synes at være ganske simpelt: hin synes at forholde sig til
denne som Middel til Maal. Men hvor skal Grænsen sættes?
Gives der nogensomhelst Udvikling af Tanke, Fantasi og Følelse,
der ikke igen kan virke tilbage paa den materielle Kultur og
lette Frembringelsen af materielle Goder direkte eller indirekte?
Der finder et stadigt Kredsløb Sted mellem ideel Kultur og
materiel Kultur. En højere Intelligens, en mere levende Fantasi,
en inderligere Følelse forandrer Arbejdets Vilkaar og bestemmer
dets Retning. Frie Bevægelser paa det aandelige Livs Omraade
vække Haab, Dristighed og Energi til ogsaa at gaa fremad
paa den materielle Kulturs Omraade. Og paa den anden Side
er det materielle Arbejde ikke altid blot Middel til at frembringe
et Produkt; det kan blive en Skole for Villien, en Form for
Øvelse af Kræfterne. Hvis der er en Fremgang i det materielle
Arbejdes Vilkaar, maa den  bestaa i, at Arbejdet, selvom det
ikke er en Leg, dog forbindes med en umiddelbar Tilfredsstillelse
ved Kræfternes Brug. Mennesket er ikke blot et Nervevæsen,
men ogsaa et Muskelvæsen; det har ligesaa vel en Trang
til at bruge sine legemlige Kræfter som til at bruge sine aandelige
Kræfter, og Tilfredsstillelsen af denne Trang er Mere end
blot Middel, ligesaa vel som Tilfredsstillelse af Erkendelsestrangen
er det. Dersom det materielle Arbejde stedse kunde forbindes
med en saadan umiddelbar Tilfredsstillelse, vilde Modsætningen
mellem materiel og ideel Kultur falde bort i et væsentligt
Punkt. Historien viser, at der paa ethvert Trin af menneskelig
Udvikling er baade ideel og materiel Kultur til Stede, skønt
de kunne være udviklede i forskellig Grad og til forskellige
Tider staa i forskelligt Forhold til hinanden.

Tilsidst gives der kun en eneste Vurdering. Baade materiel
og ideel Kultur faa deres Værdi ved at være Betingelser eller
Former for Fremgang til den største Velfærd for saa Mange
som muligt. Alt hvad der direkte eller indirekte virker i denne
Retning, er etisk set produktivt. Kulturudviklingens Sundhed
% --- p. 426 ---
\setcounter{footnote}{0}%
\opage{426}beror paa det rette Forhold mellem materiel og ideel Kultur,
mellem Arbejdet paa at sikre Livets materielle Grundlag og Arbejdet
paa Tankens, Følelsens og Fantasiens Udvikling. Den
etiske Maalestok for Rigdom og Kultur er Fritiden, d. v. s. den
Tid, som kan anvendes, og som virkeligt anvendes, til ideelt
Kulturarbejde. (Smlgn. XXVI, 14).

2. Vi ere dog endnu ikke komne videre, end at det allermeste
Arbejde gaar ud paa at skaffe Midlerne til at leve. Det
store Arbejde, som anvendes herpaa, er kun i meget ringe Grad
umiddelbart tilfredsstillende og udviklende for Arbejderens Personlighed.
Og den Fritid, vi have, og som altfor Faa af os have,
forstaa vi endnu ikke at bruge paa rette Maade.

Paa laveste Trin af Existens gaar næsten al Virksomhed
ud paa Individets og Slægtens Opholdelse. Den Tid, som bliver
tilovers derfra, er væsentligt Hviletid, som anvendes til at samle
Kræfter til nyt Opboldelsesarbejde. Det betegner et betydningsfuldt
Vendepunkt, naar Hviletiden kan anvendes til Virksomhed,
som ikke er nødvendig. Da kan Livet faa en naturlig Rytme,
ikke blot af fysisk Arbejde og af fysisk Hvile, men af fysisk
Arbejde og aandeligt Arbejde. Herpaa beror den kulturhistoriske
Betydning af Grækernes Fester og Jødernes Sabbat. Især er
denne sidste af stor Betydning, da den synes at have fastslaaet
Længden af den for de Fleste naturlige Rytme\footnote[1]{uj Under den franske Revolution prøvede man paa at lade den syv Dages Rytme afløses af en ti Dages (Ugen af Dekaden), som man fandt mere rationel. Men den strandede af forskellige Grunde, ikke blot religiøse. Kvinderne i Forstaden St. Marcel erklærede, at en Arbejder ikke kunde arbejde ni Dage i Træk uden at hvile. --- Ad. Schmidt: Pariser Zustånde während der Revolutionszeit von 1789—1800. III, p. 247.}- Men det er
ikke altid, at Hviletiden anvendes paa bedste Maade. Ligesom
det er en Maalestok for Nationalrigdommen, hvor stor Fritiden
er, saaledes er det en Maalestok for Nationaldannelsen, hvorledes
denne Fritid anvendes. Aristoteles siger om Spartanerne,
at deres Stat gik tilbage, fordi de ikke forstode at bruge
deres Fritid paa rette Maade; de havde skaffet sig Magt og Rigdom,
men af Mangel paa Aandsdannelse havde de Intet at bruge
% --- p. 427 ---
\setcounter{footnote}{0}%
% CHECK p. 427: notes paired with the marks in order (the layer misread a number)
\opage{427}disse Midler til\footnote[1]{Polit. 11, 9.}. Om saare Mange ide saakaldte højere Klasser
gælder det Samme endnu. Netop i disse Klasser hengiver
man sig ofte til sanselige Nydelser af laveste Art, fordi der
mangler Drift til ædlere Nydelser. Man kan da ikke undre sig
over, at det Samme er Tilfældet i saa høj Grad i de fysiske
Arbejderes Klasse, hvor den stærke Anspændelse i Arbejdstiden
ofte ikke lader nogen Energi tilovers til Hviletiden. Idethele
synes dog en Sammenligning mellem Fornøjelsernes Art i ældre
og nyere Tid at føre til det Resultat, at Raaheden er bleven
mindre end i tidligere Tider\footnote[2]{L. Felix: Der Einfluss der Sitten und Gebråuche auf die Entwicklung des Eigenthums. p. 197.}.

3. Af to indbyrdes nøje sammenhængende Grunde faar
den ideelle Kultur en højere Værdi end den materielle. --- Den
staar i inderligere Forbindelse med Menneskets Personlighed. Der
er ikke den Forskel mellem Arbejde, Arbejdsmidler og Arbejdsprodukter,
som ved den materielle Kultur. Vi arbejde her med
vor egen Aand, og hvad vi frembringe, kan ikke skilles fra
vor egen Aand, men tilhører den for bestandigt. Derimod opstod
jo det sociale Spørgsmaal paa Grund af, at Arbejdskraften
ikke altid var forbunden med Arbejdsmidlerne eller raadede
for Produktet\footnote[3]{Smlgn. den nærmere Udvikling heraf i Afhandlingen Aandelig Kultur (Mindre Arbejder. Tredie Række).}. —Og uagtet Arbejdet her er mere forbundet
med selve Personligheden, arbejder Individet dog ikke for sig
selv alene, men for hele Slægten. Hvis det lykkes ham at frembringe
nye Tanker, nye Billeder eller nye Former for Følelseslivet,
har han dermed forøget Slægtens aandelige Kapital, og
han bliver ikke selv fattigere, fordi Alle dele den med ham. En
materiel Partikel kan paa én og samme Tid kun være Del af
én eneste Organisme; derfor staar Kampen saa heftigt om de
materielle Goder. Men en Tanke kan frembringes i saa mange
Bevidstheder, det skal være. Her gælder ingen isoleret Ejendomsret.
Og ikke blot er Fællesskab lettere muligt paa det
aandelige Omraade end paa det materielle; men Trangen til det
er ogsaa større. Uden Vexelvirkning og fælles Stræben udvikler
% --- p. 428 ---
\setcounter{footnote}{0}%
\opage{428}det højere aandelige Liv sig ikke. Det er ikke blot de færdige
Tanker, Fantasier og Følelser, som føre til Vexelvirkning mellem
Individerne; det gælder nok saa meget om de vordende. De
nye Tanker og Følelser kunne kun udvikles, naar forskellige
Bevidstheder paavirke hverandre. Det gælder dog ikke blot om
Frembringen af nye Tanker og Aandsformer, men ogsaa om
Tilegnelse, selvstændigTilegnelse. Ogsaa derigennem gaarAandslivet
frem, og der er ingen skarp Sondring mellem Frembringelse
og Tilegnelse, Aktivitet og Passivitet. Enhver, der har
Mulighed, det vil sige: Evne, Kraft og Tid, til at udvikle sit
Tanke- og Stemningsliv, har ikke blot en individuel, men en
social Forpligtelse dertil, hvad enten saa Tilegnelsen eller Frembringelsen
hos ham er det Overvejende. Der er en indre Kunst
i Sindet, som er det Væsentlige hos de store Mestre, og som
ogsaa kan øves af dem, der lære af dem. „At udfinde Noget
selv, er skønt“, siger \textsc{Goethe}, og han føjer til: „men naar du
har fattet og vurderet, hvad Andre have fundet paa lykkelig
Maade, er det saa ikke lige saavel at kalde dit eget?“ Aaben
Sans og ideel Begejstring kunne i enhver Kreds bane Vej for
aandelige Goder. Ingen er udelukket. ---

Dog har den ideelle Kultur ogsaa sine Skyggesider. Den
udvikler sig ofte paa en ensidig og usund Maade; egoistiske
Lidenskaber tumle sig ofte paa dens Enemærker; og Arbejdet i
dens Tjeneste er ofte forbundet med store Lidelser. --- Vi betragte
hvert af disse tre Punkter for sig.

Hvor Tanke, Fantasi og Følelse ikke udvikle sig Haand i
Haand med Villien og Muskelkraften, faar den ideelle Dannelse
let Karakteren af magelig Nydelse, Leg med Tanker og Billeder,
Svælgen i Fantasier og Følelser, Sentimentalitet og Reflexionssyge.
Der gives en Art aandelig Gourmandise, som vurderer
Tanker og Følelser blot efter deres øjeblikkelige Smag, ikke
efter deres virkelige Næringsværdi. Og der gives en død Lærdom
og en Opgaaen i Specialiteter, der ikke mindre kunne hæmme
den egne personlige Udvikling. Saadanne Misligheder ere at
sammenligne med den forkrøblede og forkomne Skikkelse, det
legemlige Liv antager under uforholdsmæssigt Tryk af materielt
% --- p. 429 ---
\setcounter{footnote}{0}%
\opage{429}Arbejde. Ogsaa den ideelle Kultur har sine Trælle, hvem den
paatrykker Ensidighedens Stempel.

Den ideelle Kultur udvikler sig ligesom den materielle
under Kamp, og Kampen sætter ogsaa her Lidenskaberne i Bevægelse.
Ærgerrighed, Herskesyge og Skinsyge have her en
frugtbar Jordbund. Sekt- og Partivæsen blomstrer, og Modsætningerne
faa ofte her en Dybde, som de ikke have under Kampen
om materielle Goder. Det gælder jo her om Noget, som langt
inderligere, end materielle Goder kunne være, er knyttet til
selve Personligheden; Dommen om et Menneskes Tanker, Fantasier
og Følelser er i ganske anden Forstand en Dom om ham
selv end Dommen om hans materielle Arbejde. Den ideelle
Kultur udvikler desuden Forskellighederne mellem de menneskelige
Karakterer i langt højere Grad end den materielle Kultur.
Frembringelsen af materielle Goder lægger kun Beslag paa
elementære Kræfter, der i det Hele og Store ere ensartede hos
Alle; men ved den frie Udfoldelse af det aandelige Liv fremtræde
Forskelle og Nuancer, som forhen vare umærkelige.
Individualitetsforskellighederne stige med stigende ideel Kultur.
Men dermed stige ogsaa Mulighederne for Strid og Disharmoni.
Der gives kun én Sandhed; Striden maa da blive heftig mellem
dem, der trods deres store Forskellighed hver for sig mene at
besidde den. Og naar det endda blot gjaldt selve Sandheden.
Men Æren for at finde Sandheden spiller ofte en større Rolle
end selve Sandheden, og den kan kun En faa. --- Er der Meget
i den ideelle Kultur, som gør den til en mægtig forbindende
Kraft, saa er der ogsaa Mulighed for, at netop her store Modsætninger
kunne opstaa.

Endeligt følger der af det inderlige Forhold mellem aandeligt
Arbejde og Arbejderens Personlighed Muligheden af en Art
Lidelse, som kun den aandelige Arbejder kender. Det er ikke
til alle Tider, at det indre Arbejde lykkes. Der kan være indre
Modstand at overvinde. Der kan være en Disharmoni mellem
Trang og Evne, der føles des pinligere, jo mere det aandelige
Arbejde føles som en personlig Sag, og jo mere ideal en Maalestok
man anlægger paa sin Stræben. Det gælder ikke blot om
at sætte Noget ind i Verden, der kan fyldestgøre Markedets
% --- p. 430 ---
\setcounter{footnote}{0}%
\opage{430}Fordringer, men om at tilfredsstille sin Genius. Ofte kan den
aandelige Arbejder misunde den materielle Arbejder, der trygt
forlader sit Værksted, naar han i den bestemte Tid har ydet et
vist Kvantum mekanisk Arbejde. Den aandelige Arbejder har
ingen saadan Tryghed, især naar hin Disharmoni gør sig gældende.
Mange mørke Timer og Tider fremkaldes derved; der
kan føles en Angst ved Stansningen af det aandelige Liv i os,
som minder om den organiske Angstfølelse, naar Aandedrættet
eller Blodomløbet hæmmes. Tanken vil ikke klare sig, og Følelsen
synes at være tørret hen. Lykken i de gode Tider købes
dyrt med Angsten og Tvivlen i de mørke Tider. Naar Sindet
har erfaret, hvad det er at fyldes med en stor ideal Interesse,
føler det Slappelsen og Uformuenheden i de onde Tider langt
mere, end hvis det slet ikke havde kendt saadanne Interesser.
Dertil kommer den ydre Modstand, Misforstaaelse, Spot og Kulde,
som det aandelige Værk saa ofte møder, ikke mindst naar det
er ejendommeligt og originalt. --- Denne Lidelse falder ikke blot
i de store Geniers Lod, men kan ogsaa blive dem til Del, der
paa ejendommelig og selvstændig Maade tilegne sig, hvad hine
have frembragt. Ogsaa de kunne erfare baade den indre og
ydre Modstand.

\grouphead{A. INTELLEKTUEL KULTUR.}

\chapterhead{XXVIII.}{VIDENSKABELIG ERKENDELSES ETISKE BETYDNING.}

1. Den Antagelse, at Alt i Verden er underlagt en rytmisk
Bevægelse, og at hvis der er et Fremskridt, gaar det i hvert
Tilfælde ikke i lige Linie, bekræftes maaske intetsteds tydeligere
end ved Dommene om den intellektuelle Udviklings Betydning.
% --- p. 431 ---
\setcounter{footnote}{0}%
\opage{431}--- For Grækerne stod Tanken eller Fornuften som Menneskets
Adelsmærke. Kun den var et sandt Menneske, der kendte det
Skønne og Gode. For de græske Filosofer stod Tænkning og
Erkendelse som den højeste af alle Virksomheder. I den kristelige
Middelalder gjaldt det derimod at tage Fornuften fangen.
Den naturlige Forstand var en Hedning, der skulde bøje sig for
Troens Myndighed. Da man senere havde brudt med Autoritetsprincipet,
ventede man sig af Videnskabens og Oplysningens
Udvikling en fuldstændig Forandring, en Fuldkommengørelse af
det menneskelige Liv. Allerede hos Bacon og Descartes kan
denne store Forventning mærkes; men hos mange af det 18.
Aarhundredes Forfattere blev den til Lidenskab, til fanatisk Tro.
Den revolutionære Grusomhed finder for en Del sin Forklaring
i, at man ikke kunde tænke sig andet end en ond Villie som
Aarsag til, at ikke Alle bøjede sig for det nye Evangelium:
Fuldkommenheden laa jo saa nær, naar man blot vilde aabne
Øjnene! Og dog ytrer der sig allerede under selve Revolutionsperioden
en helt modsat Strømning, dels Eftervirkningen af
Rousseaus Kamp for Følelsens Ret overfor Forstanden, dels
Frygt for, at der skulde danne sig et nyt Aristokrati, hvis Videnskaben
begunstigedes. Reaktionen i den første Del af det 19.
Aarhundrede vilde vel ikke, som de revolutionære Fanatikere,
afskaffe Videnskaben; men den forlangte, hvad der var endnu
betænkeligere, at Videnskaben skulde „vende om“, d. v. s. fornægte
sine egne Principer, og den ansaa desuden Oplysning
for en farlig Sag. --- Hos den nyeste Tids Tænkere finder man
Forsøg paa en mere uhildet og alsidig Undersøgelse af Forholdet
mellem Erkendelsen og Sjælelivets andre Sider. I \textsc{Stuart}
\textsc{Mill}'s Selvbiografi kan man se, hvor megen Anstrengelse det
har kostet en ædel og dygtig Aand at frigøre sig for den ensidigt
intellektualistiske Anskuelse, i hvilken han var bleven
opdragen.

Disse Svingninger i Vurderingen af den intellektuelle Kulturs
etiske Betydning ville forstaas, naar man ser hen dels til det
psykologiske Forhold mellem Erkendelsen og de andre Sider af
Sjælelivet, dels til de historiske Erfaringer, som foreligge om
den intellektuelle Udviklings Virkninger.

% --- p. 432 ---
\setcounter{footnote}{0}%
\opage{432}2. Psykologien lærer os, at Erkendelsen udvikler sig hurtigere
end Følelse og Villie. Vore Iagttagelser og Forestillinger
kunne opstaa og udfolde sig i Bevidstheden, uden at tilsvarende
Følelser og Drifter strax røre sig i deres hele Styrke. Vi kunne
omspænde langt Mere i Tanke og Fantasi, end vi nogensinde
kunne omfatte med levende Følelse eller virkeliggøre gennem
Villiens Arbejde. Dette kommer dels af, at vi til enhver Tid
kun raade over en vis bestemt Sum af Energi, og naar intellektuel
Virksomhed lægger Beslag paa en stor Del af denne
Sum, bliver der saa meget des Mindre tilbage til de andre sjælelige
Virksomheder; --- dels af, at Følelsen og Villien ere mere
konservative i deres Natur end Erkendelsen, holde mere fast
ved den Retning, de én Gang ere komne ind paa. Det tager
Tid, inden Erkendelsesresultaterne nedfældes i Kød og Blod,
saa at et harmonisk Forhold bringes til Veje mellem Sjælelivets
forskellige Sider.

Der indtræder derfor ofte som foreløbig Følge af intellektuel
Udvikling en Tilstand af Disharmoni og Splid i Bevidstheden.
Den udvidede Ramme, Tanken har bragt i Stand, kunne Følelse
og Villie ikke udfylde. Og den instinktmæssige Sikkerhed,
hvormed Livet førtes, saalænge Horizonten endnu var begrænset,
giver Plads for Tvivl og Uro. --- Hvor denne Disharmoni ikke
er til Stede, ser man til Gengæld ofte en Svækkelse og Afstumpning
af Følelseslivet indtræde. En vis Mathed breder sig;
Livet mister i Varme, hvad det vinder i Klarhed.

Den revolutionære og den reaktionære Fanatisme kunne
forsaavidt beraabe sig paa Psykologien. Kundskabens Træ er
ikke uden videre Livets Træ. Men derfor er det jo ikke sagt,
at man strax skal hugge det om eller beskære det.

3. Disse psykologiske Resultater bekræftes forsaavidt af
Historien, som der ikke kan paavises nogen sikker og almindelig
positiv etisk Virkning af den stigende Oplysning. Denne gør
sig mere mærkelig ved den Uro, den vækker, end ved sin opdragende
Indflydelse. Sindssygdom og Selvmord ere særligt
hyppige i Lande, hvor Undervisningsvæsenet staar forholdsvis
højt, og man kan i hvert Tilfælde ikke paavise, at Forbrydelserne
tage af i saadanne Lande. Den Forhaabning, som man
% --- p. 433 ---
\setcounter{footnote}{0}%
\opage{433}har haft, at man skulde kunne lukke et Fængsel for hver
Skole, man aabnede, er ikke gaaet i Opfyldelse. Selvom man
nu\footnote[1]{Forbrydelsernes Tiltagen. --- Rubin har i „Tilskueren“ (1884) p. 466 ff. vist, paa hvor usikkert Grundlag Oettingen har konstrueret dette Begreb s} vilde indskrænke sig til at sige, at der ikke statistisk kan
paavises nogen Aarsagssammenhæng mellem intellektuel og moralsk
Udvikling, saa maa dette dog allerede vække Forundring
og Skuffelse, naar man ser hen til det store Arbejde, der anvendes
paa Undervisning og Oplysning. Og hvis man vil sige,
at det kun er den hidtil saa ufuldkomne Oplysning og Undervisning,
som er Aarsag til Mislighederne og til de ringe etiske
Virkninger, saa rejser der sig en Betænkelighed af en anden
Art. Dersom kun den Kundskab, der vindes og tilegnes paa
selvstændig Maade, kan have gode Virkninger i etisk Henseende,
saa synes det at være farligt at give Undervisning til Andre
end dem, der have indre og ydre Betingelser for at kunne naa
Maalet fuldstændigt. Men hvor Mange vilde da ikke blive udelukkede!
Og hvilken Modsætning mellem Vidende og Uvidende
vilde da ikke opstaa! Samfundet vilde spaltes paa en ikke mindre
betænkelig Maade end det enkelte Individs Væsen spaltes ved
en Erkendelse, der ikke har slaaet Rod i hans Natur. Det vil
kun være en daarlig Trøst, at Historien stedse viser os en
saadan Modsætning mellem Vidende og Uvidende under forskellige
Former. De Vilde havde deres „Medicinmænd“ og Regnmagere;
de gamle Ariere havde deres Sangere, der kunde lokke
Guderne ned til Ofrene før Kampen; Inderne og Ægypterne
havde deres Præstekaste, Kineserne have deres Mandariner.
Overalt er der en Kundskab, som agtes højt, og som kræver
tnl,nltSmtr ltmeUnl Zusam\textasciicircum{}kang der sittlichen und
intellektuellen Bildung. (Reden und Aufsatze. Neue Folge). p. 4. --- Mondière
i „Dictionnaire des sciences anthropologiques“. Art. Instruction.

--- R. Garofalo: La Criminologie. Paris 1888. p. 135 f. 162. --- Buckle
gaar i sin Englands Civilisationshistorie et Skridt videre, idet han negter,
at der overhovedet finder nogen moralsk Udvikling Sted (men vel en
intellektuel). Og Oettingen gaar endnu et Skridt videre, idet han (Woralp
ff,\footnotemark[3] Aufl, P ---\footnotemark[6] f“\footnotemark[7]\footnotemark[6]\footnotemark[2] ff° giver „Halvdannelsen“ Skylden for
% CHECK p. 433: note mark without a note
% CHECK p. 433: note mark without a note
% CHECK p. 433: note mark without a note
% CHECK p. 433: note mark without a note
% --- p. 434 ---
\setcounter{footnote}{0}%
\opage{434}særegen Indvielse og Øvelse, saa at den kun kan findes hos
nogle Faa, der øve et aandeligt Herredømme over hele Resten.
Vil det ikke stedse gaa saaledes med vor videnskabelige Viden,
jo mere speciel og dybtgaaende denne er? Videnskaben, saa
synes det, er og maa være aristokratisk.

Der kunde i denne Henseende synes at være sket et Tilbageskridt
siden Middelalderen. Den Gang var der vel ogsaa
en lærd Stand; men Troslærens Grundsætninger vare fælles for
den Enfoldigste og for den lærdeste Skolastiker. Enhver Landsby
kunde kaldes et Athen: de højeste Sandheder kunde forkyndes
og forstaas der. Er noget Tilsvarende nu muligt? --- Videnskaben
har specialiseret sig i den Grad, at den enkelte Forsker
kun kan beherske et meget lille Omraade paa selvstændig Maade.
Hvorledes kan der da være Tale om en fælles, almindelig udbredt
intellektuel Dannelse? ---

De her berørte Spørgsmaal ere de vigtigste, den sociale Etik
har at behandle med Hensyn til den intellektuelle Udvikling.

4. Vi ere sikkert i den nyere Tid komne ind paa en
Overvurdering af de intellektuelle Evner. Vi udvikle dem uden
at forvisse os om, hvorvidt Udviklingen sker i Harmoni med
Sjælelivets andre Sider. Det er Erkendelsens naturlige Rolle at
bestemme Følelsens og Villiens Indhold og Retning. Men kun
den Erkendelse kan lede os overfor Virkeligheden, som selv er
udsprungen af Virkeligheden. Klagerne over den blotte Forstandsoplysning
og dens Unyttighed eller Skadelighed burde derfor
særligt gaa ud paa, at den baade i Udspring og i Anvendelse
staar det virkelige Liv for fjernt. Det var ogsaa i mange Maader
Tilfældet med det 18. Aarhundredes Oplysning. Den var begrænset
til en lille Kreds; den store Mængde skulde laane sit
Lys fra denne lille Kreds, og en selvstændig Erhvervelse kunde
der derfor ikke være Tale om. Og selve den lille Kreds af
virkeligt oplyste og selvtænkende Mænd stode det virkelige Liv
for fjernt. Enevælden holdt dem ude fra al praktisk Deltagelse
i Folkets Anliggender. Derved maatte den hele Literatur faa
et ensidigt, tildels fantastisk Præg. Folket blev uselvstændigt
paa Grund af den herskende Centralisation og manglede baade
Trangen til Oplysning og Evnen til at bruge denne. Folke%
% --- p. 435 ---
\setcounter{footnote}{0}%
\opage{435}oplysningen skal endog have været større i Frankrig i det 16.
end i det 17. og 18. Aarhundrede\footnotemark[1].

Videnskaben voxer stedse paany frem af Livet. Kunstige
Midler formaa hverken at hindre den, naar der rører sig en
Trang til den, eller at holde den oppe, naar en saadan Trang
mangler. Naturvidenskaben voxer ud af det praktiske Livs Krav,
af Trangen til at beherske den ydre Natur; Filosofien bliver
(som i nyere Tid især Filosofiens Historie i England viser) til
ved Spørgsmaal, som Livet selv fører til at stille; Historien
bliver til ud af Trangen til i Erindringen at bevare Folkets og
Slægtens Liv. Den Kundskab, som forbliver virkningsløs, er
den færdige, den udefra meddelte Kundskab. Dér hvor man
kan faa et Frøkorn til at spire, kan man ikke altid faa et helt
Træ til at fæste sin Rod.

Det er den pædagogiske Kunsts Ideal at vække Trangen
til Kundskab, før Kundskaben meddeles, og kun meddele Kundskaber
i det Omfang, i hvilket der føles Trang til dem. Jo mere
det lykkes at virkeliggøre dette Ideal, des mere vil den ovenfor
skildrede Disharmoni forsvinde. Der er ingen Grund til at
vende om; det gælder kun om at fortsætte den intellektuelle
Udviklingsproces paa saa gode Betingelser som muligt. Denne
Udviklingsproces er intet Kunstprodukt, men er fremgaaet af
Livets egen Trang. ---

Den Oplysningsstræben, som blot bedømmer et Menneske
efter dets Forstandsudvikling, og hvis hele Psykologi indskrænker
sig til at skelne mellem Dunkelhed og Klarhed, har Rousseau
givet Dødsstødet*). Saa stor Betydning vi end tillægge Vi-

) Smlgn. Tocqueville: L'ancien Régime et la Révolution. 7 éd.
p. 207 f. Adolf Schmidt: Pariser Zustände. III, p. 335—337. --- Alle
rede Rousseau har fremhævet den Skade, Erkendelseslivet lider ved at
staa det praktiske Liv for fjernt: „Tant que Ia puissance sera seule d'un
côté, les lumiéres et la sagesse seules d'un autre, les savans penseront
rarement de grandes choses, les princes en feront rarement de belles, et
les peuples continueront d'étre viles, corrompus et malheureux.“ (Discours
sur les sciences et les arts).

-) Smlgn. min Bog: Rousseau og hans Filosofi\footnotemark[2], især p. 49—78
(Bruddet med Encyklopædisterne og med Voltaire).
% CHECK p. 435: note mark without a note

% CHECK p. 435: note mark without a note

% --- p. 436 ---
\setcounter{footnote}{0}%
\opage{436}denskab og Kundskab, saa er der dog andre Sider af Menneskets
Væsen, til hvilke vi først og fremmest se hen ved Dommen
om et Individ. Overfor det Liv, der udfolder sig i Følelsen
og Villien, er Erkendelsen kun en Forgaard, i hvilken vi kunne
færdes, uden at komme ind i den egentlige Helligdom. Her
er det, Slægtskabet mellem det enfoldigste Menneske og den
største Tænker stadigt kan være tilstede trods al Forskel i intellektuel
Dannelse. Det er ingen Ringere end \textsc{Kant}, som har
sagt: „Jeg er selv Forsker, og er det af egen Drift. Jeg føler
den hele Tørst efter Erkendelse og den begærlige Uro efter at
komme videre deri, saavel som Glæden over ethvert Fremskridt.
Der var en Tid, da jeg troede, at alt dette kunde udgøre Menneskehedens
Ære, og jeg foragtede Pøbelen, som Intet véd.
Rousseau har ført mig paa den rette Vej. Dette blændende
Fortrin forsvinder. Jeg lærer at ære Menneskene og vilde anse
mig selv for langt unyttigere end de simple Arbejdere, naar jeg
ikke troede, at selve denne Betragtning kunde virke til at
hævde Menneskehedens Rettigheder“\footnote[1]{Fragmente. Werke udg. af Rosenkrantz. XI, p. 240.}. Kant havde i sine første
Skrifter trøstet sig over Nøden og Kampen i Verden ved den
Tanke, at Oplysningen dog gik fremad hos en lille udvalgt
Skare. Studiet af Rousseaus Skrifter frembragte en hel Revolution
i hans Anskuelser og lod ham søge det egentligt Menneskelige
dybere end i Forstanden. ---

Naar Videnskaben voxer ud af Livet og stedse bevarer
Sammenhængen med Livet, og naar desuden det aandelige Livs
Centrum ikke ligger paa det intellektuelle Omraade, men i Følelsen
og Villien, saa formindskes Faren for den Disharmoni
hos den Enkelte og den Spaltning i Slægten, som den intellektuelle
Udvikling saa let syntes at føre til. Og ikke mindre
vil derved den Fare undgaas, som ligger i Udviklingen af et
intellektuelt Proletariat. Hvor den intellektuelle Udvikling staar
i nøje Forhold til et bestemt videnskabeligt Arbejde, eller til
en udpræget Trang til Klarhed i Livsanskuelsen, eller til Virksomhed
for bestemte praktiske Formaal, dér er den sund og
nødvendig. Men den kan blive en Ulykke, hvor Livsforholdene,.
% --- p. 437 ---
\setcounter{footnote}{0}%
\opage{437}de indre og de ydre, ikke tillade denne Forbindelse at bestaa.
Det er en Opgave for pædagogisk Kunst, saavidt muligt at
bevare Harmonien i denne Henseende hos den Enkelte og i
Folket.

5. Det kan ikke være Andet, end at der bliver en Modsætning
mellem forskellige Trin af intellektuel Udvikling. Selvom
det ikke er mellem Vidende og Uvidende, saa er det mellem
mere og mindre Vidende. Det følger af Arbejdets Deling. Skal
Videnskaben dyrkes paa tilbørlig Maade, maa der være dem,
der vie den hele deres Liv. Der vil danne sig en lærd Kreds
eller Stand i Folket. Det gælder da her om to Ting, hvis Udviklingen
skal være sund. For det Første maa der ikke danne
sig en lærd Kaste med Tilgang fra nogle Klasser af Samfundet,
men ikke fra alle. Den lavere og den højere Undervisning
maa ordnes saaledes, at der bliver en jevn Overgang fra de
laveste til de højeste Trin\footnote[1]{En saadan jevn Overgang var til Stede, da Teologi var Hovedfaget ved Universitetet ligesom Katekismen i Folkeskolen. Nu er Teologien reduceret til et Minimum i den lærde Skole og til en Specialitet ved Universitetet, men Katekismen giver endnu Folkeskolen dens Præg. Martensen hævder, at saaledes skal det være: „Det er Religionen, som gør Folkeskolen til Folkeskole“ (Social Etik, p. 355), en Ytring, der er paafaldende endog for hans eget Standpunkt. Her slaas en Dualisme fast, som vil faa sine meget betænkelige Følger, hvis der ikke i Tide raades Bod derpaa. Og det bliver da neppe Videnskaben, der kommer til at „vende om“.}, og at Vanskelighederne ved at
gennemløbe alle Trin ikke blive uovervindelige for den, hos
hvem alvorlig Drift og Evne rører sig. Og dernæst maa den
lærde Stand betragte Forskningen ikke blot som Tilfredsstillelse
af sin egen Drift, men som en social Gerning, en Gerning, der
øves paa hele Samfundets Vegne. Den enkelte Forsker er sat
paa sin Udkigspost for at iagttage Alt, hvad han dér kan faa Øje
paa. Naar dette Sindelag besjæler ham, bevarer han Enheden
med Slægten, selvom han føler sig isoleret paa sin ensomme
Post, og selvom han længe bliver miskendt og misforstaaet,
fordi de Andre ikke kunne faa Øje paa, hvad han har set. Hans
Tro paa Sandheden er tillige en Tro paa Sandhedens Betydning
for Slægten.

% --- p. 438 ---
\setcounter{footnote}{0}%
\opage{438}Videnskabens fremskridende Specialisering gør det sikkert
vanskeligere og vanskeligere selv for Videnskabens egne Dyrkere
at skaffe sig et selvstændigt Overblik over dens Resultater.
Men intellektuel Kultur forudsætter heller ikke, at Alle
vide Alt. Der kan være en selvstændig intellektuel Dannelse
i store Kredse, naar blot selve Tænkeevnen er udviklet. Trods
Videnskabernes Forskellighed er det dog en og samme Tankekraft,
med hvilken de alle arbejde, og det er en og samme
Verden, de alle, hver fra sin Side, søge at trænge ind i. Derfor
vil den, der paa fornuftig Maade er ført ind i et enkelt
Omraade, let kunne bringes til at forstaa de Opgaver og Vanskeligheder,
som frembyde sig paa andre Omraader. Der er
desuden visse Hovedtanker, som komme igen overalt, hvor menneskelig
Erkendelse arbejder, og som det er Filosofiens Opgave
at fremdrage og udvikle. En saadan Tanke er Ideen om Tilværelsen
som en lovbunden Sammenhæng. De forskellige Videnskaber
søge hver at gennemtrænge sit Stykke eller sin Side
af én eneste stor Aarsagssammenhæng. Til Ideen om Aarsagssammenhængen
slutter sig Ideen om Udviklingen, der i sine
Hovedtræk genfindes paa de forskellige Erfaringsomraader. I
en lille Del af Verden kunne vi derfor studere den hele Verden
og vinde en Forstaaelse af den. Der danner sig for vor
Bevidsthed et almindeligt Verdensbillede, hvis Huller de enkelte
Videnskaber uafbrudt arbejde paa at udfylde. Selv den, der
ikke kan ofre sig for videnskabelig Forsken, vil dog kunne
danne sig en Forestilling om Hovedlovene i den Tilværelse, af
hvilken han er et Led, og faa en Vejledning til at anlægge en
mere universel Betragtningsmaade paa sig selv og sit Liv, end
det er muligt, saalænge han tror, at Alt drejer sig om ham og
hans Interessers snevre Horizont. Ved at erkende sin Plads i
Tilværelsen vil han ogsaa vinde i Selverkendelse. Derhos vil
han styrkes i sin Tro paa Menneskeslægtens Enhed, naar han
faar et Indblik i det store Fællesarbejde paa at danne Menneskeslægtens
Verdensbillede, et Arbejde, der Slægt efter Slægt
lægger Beslag paa de ædleste Kræfter, den mest udholdende
og opofrende Stræben. Den Enkelte arbejder her ikke for sig
% --- p. 439 ---
\setcounter{footnote}{0}%
\opage{439}selv alene; hans egen Forsken vilde ikke føre ham langt. Men
han kan lægge sin Sten til den store Bygning, som er Menneskeslægtens
fælles Verdensanskuelse, der langsomt udvikler
sig gennem Tiderne. Og han kan maaske fra sin Arbejdsplads
overse saa Meget af Bygningen, at han kan danne sig en
Anelse om, hvorledes den vilde se ud, hvis den blev helt
færdig\footnote[1]{samt en Række Foredrag af franske Historikere og Filosofer, udgivne under Titel: L'éducation et la démocratie. Paris 1903.}.

6. Videnskaben virker samfundsstiftende ved Nødvendigheden
af Samarbejde paa Slægtens Verdensbillede. Derved forbindes
ikke blot de samtidigt Forskende, men Tanken om den
Gæld, i hvilken vi staa til Fortiden, bliver levende. Intet Steds
kan en fremskridende Udvikling saa tydeligt paavises som paa
den videnskabelige Erkendelses Omraade. Intet Steds er det
lettere at vise, hvorledes den ene Slægt staar paa den andens
Skuldre og derved faar et videre Blik end den. Men da Videnskaben
er i stadig Udvikling, kan Samarbejdet ikke altid være
ganske fredeligt. Skoler og Partier staa overfor hverandre og
bekæmpe hverandre. Her dannes da ikke ét stort Samfund,
men flere mindre i Strid med hverandre. Men saadan Strid
kan være Middel til Fremgang, naar den er Mere end personligt
Kævleri. Den opstaar da dels af, at selve Sagen frembyder
flere forskellige Sider, dels af, at Forskerne møde med forskellige
Forudsætninger. I begge Henseender vil Partidannelsen
kunne virke frugtbringende. Naar man slutter sig sammen med
Andre for at udtømme Alt, hvad der kan udfindes om Sagen
fra den Side, man ser den, eller ud fra de Forudsætninger,
man gaar ud fra, vil det ske med større Grundighed og Iver.
Selve den Affekt, som Partimodsætningen medfører, kan skærpe
Blikket. Selv rent personlige Stridigheder kunne paa denne
Maade blive frugtbringende. Det fælles Værk fremmes ikke
blot gennem Individernes uinteresserede Hengivelse, men ogsaa
gennem Brydningen mellem deres Egoismer. De arbejde rast-

:) Smlgn. mit Foredrag Videnskab og Folkeliv (Tilskueren 1903),
% --- p. 440 ---
\setcounter{footnote}{0}%
\opage{440}løst for at faa Ret, for at hævde deres Ære, og kunne derved
ofte fremme et Formaal, som er langt mere omfattende end
det, de i Stridens Hede have for Øje.

\chapterhead{XXIX.}{DEN INTELLEKTUELLE KULTURS FRIHED OG SELVSTÆNDIGHED.}

1. Frihed er den vigtigste Betingelse for intellektuel Kultur.
Videnskaben kender ingen udefra afstukne Grænser, Intet,
der er for højt eller for lavt til at falde ind under dens Behandling.
Den afstikker sig selv sine Grænser, idet den udfra
selve Erkendelsens Natur og Virkemaade afgør, hvilke Spørgsmaal
den formaar at drøfte, og hvilke den ikke formaar at afgøre
Noget om. Og selv med Hensyn til disse sidste, for Videnskaben
uløselige, Spørgsmaal tiltager Videnskaben sig Ret
til at undersøge de Svar, man har givet paa dem, --- efterforske
Oprindelsen til disse Svar og Grunden til, at de have fundet
Tilslutning. Hvor Erkendelsesteorien ender, tager den historiske
Kritik og Psykologien fat.

Friheden er her, som overalt (se VIII, 6; XXIII, 2), baade
Formaal og Middel. Der er en umiddelbar Tilfredsstillelse ved
den uhindrede Brug af Tankekraften, saa at et unødvendigt
Indgreb her er en Forsyndelse mod selve Velfærdsprincipet.
Men Friheden er ogsaa Middel. Kun den frie Forsken kan se
Sagen fra alle Sider, opdage de forskellige Muligheder og drage
alle Konsekvenser. Dog møde vi her det Særsyn, at selv de,
der afsky enhver Begrundelse af etiske Regler ved Hensynet
til Nytten, ofte appellere til den paa dette Punkt. „Den frie
Forsken“, siger man da, „rokker ved Alt, ogsaa ved Meninger,
Menneskene ikke kunne undvære. Sandheden er ikke skadelig;
men Tvivl og Undersøgelse uden fast Grund under Fødderne
% --- p. 441 ---
\setcounter{footnote}{0}%
\opage{441}virke skadeligt. Den Enkelte kan have Lov at tvivle og forske,
hvis han kan forsvare det for sin Samvittighed; men lad ham
beholde sin Tvivl hos sig selv og ikke drage Andre ind paa
sin vaklende Grund!“ --- Den Berettigelse, en saadan Tankegang
under visse Forhold kan have, have vi allerede i en anden
Sammenhæng anerkendt (XII, 4—6, smlgn. ogsaa VII, 4
og XXI, 5). Men her, hvor Talen er om den intellektuelle
Kulturs Betydning for Slægtens Liv idethele, maa Hensynet til
Slægtens varige Velfærdsbetingelser stilles over Hensynet til
de Enkeltes Ømtaalelighed. Det vil være fordærveligt at tage
Hensyn til den øjeblikkelige Smerte, som voldes\footnote[1]{Stuart Mill har med Rette fremhævet det som karakteristisk for den nyeste Tid, at man taler mere om Meningernes Nytte end om deres Sandhed. „I Nutiden“, siger han (Om Friheden. Dansk Overs. p. 39), „ere Folk ikke saa meget overbeviste om deres Meningers Sandhed, som om, at de ikke vide, hvad de skulle gøre uden dem''. Mill, den ivrige Utilitarianer, føler Betænkelighed herved. Sagen er, at Sandheden for ham er en Livs- og Velfærdsbetingelse af første Rang, og hvad han bekæmper, er den kortsynede og begrænsede Opfattelse af, hvad der er nyttigt.}, og derover
glemme de Følger, Forholdelse af Sandheden --- og ogsaa Tvivlen
er en Sandhed, naar den er en begrundet Tvivl --- kan have
for store Kredse og for lange Tider. De pædagogiske Hensyn er
det især den Enkeltes Sag at tage overfor andre Enkelte; i
Samfundet maa de ofte træde helt tilbage. --- De, som fra et
socialt Standpunkt ere betænkelige ved den frie Forsken, mene
desuden i Regelen, at de selv besidde Sandheden, og istedetfor
at forhandle med dem om Berettigelsen af at udtale Tvivlen
vilde det derfor være mere paa sin Plads at forhandle med dem
om selve Tvivlens Berettigelse. De dække sig bag Følgerne af
Tvivlens Udtalelse, men det er dog egentligt selve Tvivlen, de
ville til Livs. I hvert Tilfælde maa de indrømme, at Nytten
af at udtale Tvivlen maa kunne drøftes frit og aabent. De
kunne dog ikke mene at sidde inde med en ufejlbar Erkendelse
af, hvad der er nyttigt eller ikke. Og kan Tvivlens Nytte
eller Skade drøftes offentligt, uden at derved selve Tvivlen
bliver offentligt udtalt og bekendt? --- --- Hvorledes man vender
% --- p. 442 ---
\setcounter{footnote}{0}%
\opage{442}og drejer sig, vil man ikke kunne undgaa denne Fare, hvis
det er en Fare.

Fri Forsken fører sikkert Farer med sig. Men kun gennem
Vildfarelser gaar Vejen til Sandhed. „En levende Vildfarelse“,
har man sagt, „er bedre end en død Sandhed“. Det
Levende i Vildfarelsen er en Tanke, der kan arbejde sig bort
fra den gale Vej; men i en forlængst afgjort Sandhed, der ikke
længere kan sætte Sindets Kræfter i Bevægelse, er der ingen
Livsspirer. Naar vi mindes dem, der fandt de Sandheder, paa
hvilke vi bygge, saa maa vi ogsaa mindes dem, der udviklede
de „levende“ Vildfarelser og derved byggede Mellemstationer
paa Vejen til Sandheden, selvom denne Vej er en Omvej. ---
En levende Sandhed er naturligvis det Allerbedste. ---

2. Den intellektuelle Kultur virker ikke blot samfundsstiftende
ved at føre de forskende Individer sammen til fælles
Undersøgelse og Tankemeddelelse, men ogsaa ved at føre til
Dannelsen af Læreanstalter. Historien viser os her nogle af
de smukkeste Exempler paa fri Samfundsdannelse. I den græske
Oldtid dannedes de filosofiske Skoler ved, at unge Mænd sluttede
sig sammen om en Tænker, der havde vakt deres Interesse.
I Middelalderen opstod der Skoler ved Klostrene, hvor
unge Mænd studerede i Fællesskab. Universiteterne opstod i
Middelalderen ved fri Sammenslutning af lærebegærlige Mænd;
Ordet Universitet betyder som bekendt netop Samfund eller
Korporation af Lærere og Studerende. Først senere fik Kirke
og Stat Indflydelse paa disse videnskabelige Anstalter, der paa
den Tid havde en saa meget des større Betydning, som de
vare de eneste Midler til højere Dannelse; Bøger og andre videnskabelige
Midler fattedes saa godt som aldeles. I vore
Dage have vi set et lignende Fænomen herhjemme i Højskolebevægelsen.

Hvor en saadan fri Sammenslutning sker, er det Statens
Pligt ikke at lægge den Hindringer i Vejen, men efter Evne
at støtte den. Staten har, som tidligere (XXII) omtalt, den Ret
og den Pligt at sikre sig et Lavmaal af Kundskab hos enhver
af sine fremtidige Borgere. Men det er ogsaa en af Statens
Opgaver at bruge sine Midler og sin Organisation til Bedste
% --- p. 443 ---
\setcounter{footnote}{0}%
\opage{443}for den højere intellektuelle Kulturs Fremgang, afset fra, at den
af dem, der ville være dens Embedsmænd, maa fordre en vis
videnskabelig Dannelse. Den frie Sammenslutning faar let en
vis aristokratisk Karakter. Der er Mange, som ere stavnsbundne
af Trang og ikke uden Hjælp kunne arbejde med paa Tilegnelsen
og Fortsættelsen af den intellektuelle Kultur. Og det
vil ikke være heldigt, om kun de rige og velhavende Klasser
bære den intellektuelle Kultur. Staten har her, ligesom paa
den materielle Kulturs Omraade (XXVI, 14), en fordelende Virksomhed
at udøve. Den skal ikke blot søge at begunstige Dannelsen
af en Kapital af Kundskab som et fælles Gode for hele
Samfundet, men ogsaa at fremme den bedst mulige Fordeling
af denne Kundskab. Og det gaar ikke her, som det saa let
gaar ved Fordelingen af materielle Goder, at der ikke kan
gives til den Ene, uden at der tages fra den Anden.

Den Fare vil dog stedse ligge nær, at Staten gør Videnskaben
til sin Tjener, istedetfor selv at være Videnskabens Tjener.
Staten er jo til enhver given Tid repræsenteret ved Statsmagtens
Ihændehavere, hvis Partiinteresser eller personlige Interesser
kunne føre dem til at begunstige eller hindre visse
Retninger eller Personer uden Hensyn til den intellektuelle
Kulturs virkelige Tarv. Man har endog villet opfatte Sagen
saaledes, at Bestræbelserne for den intellektuelle Kultur stedse
skulle underordnes andre Hensyn. Man har sagt: de egentlige
menneskelige Samfund ere Hjem, Kirke og Stat; Skolen (Samfundet
til Videnskabens Fremme og Tilegnelse) skal tjene dem,
ikke omvendt; Hjemmet, Kirken og Staten have i Forening at
bestemme, hvorledes Skolen skal være\footnote[1]{toritåt und Kirche. Mainz 1862. p. 209. Den katolske Kirke fordrer i Kraft af sit guddommelige Udspring Ret til at have Tilsyn med alle andre Fag ere kun at betragte som Tillæg. Konfessionsløse Skoler (scholæ neutræ) ere en Uting: Jesus har jo seiv sagt, at den, som ikke er med ham, er imod ham! Smlgn. Cathrein: Philosophia moralis, p. 396. Gury: Theologia moralis. I. p. 314. --- Ikke blot hos katolske Forfattere træffer man en saadan Tankegang.}.

') „Skolen er ingen selvstændig Anstalt ved Siden af Hjemmet, Staten
og Kirken, men en af dem afhængig Hjælperske. Det er Skolens
væsentlige Stilling, som Natur og Religion har anvist den, og deri viser
den grundløse Usandhed sig af de moderne Bestræbelser, som ville stille
Skolen uafhængigt af Hjem og Kirke. Hjem, Stat og Kirke maa have
Skoler, der svare til deres Aand og deres Fordringer. De kunne ikke
uden stor Uretfærdighed hindres deri“. Biskop Ketteler: Freiheit, Au%
% --- p. 444 ---
\setcounter{footnote}{0}%
% CHECK p. 444: indent after a broken word?
\opage{444}Man nægter her Skolen, det paa Videnskabsdyrkelse grundede
Samfund, ethvert Initiativ. Men hvilken Selvmodsigelse
indvikler man sig ikke derved i! Ti hvad er det, Staten, Kirken
og Hjemmet ville opnaa gennem Skolen? Det maa vel
være Sandheden. Men de foreskrive selv i Forvejen, hvorledes
Sandheden skal være beskaffen. De maa da mene i Forvejen
at være i Besiddelse af Sandhedsreglen, saa at Skolen blot skal
anvende den i det Enkelte. Men deraf komme kun Autoritetsmeninger,
ikke videnskabelige Sandheder. Og derhos overses,
at Videnskaben jo blandt Andet ogsaa studerer Hjemmets, Kirkens
og Statens Oprindelse, Væsen og Betydning. Hvilken Betydning
har Videnskaben, naar den ikke tilsidst, maaske ad
mange Omveje, men dog uundgaaeligt fører til Forandringer i
det Liv, der føres i Hjem, Kirke og Stat? En intellektuel Kultur,
som er berøvet Muligheden af Tilbagevirken paa Livet, er,
som vi have set (XXVIII, 4), usund og unyttig. Man bliver
derfor nødt til at indrømme den intellektuelle Kultur en Plads
som selvstændigt medvirkende Aarsag ved det menneskelige
Livs Udvikling. Dersom Hjemmet, Kirken og Staten ikke kunne
bruge den selvstændige Videnskab, bliver det deres egen Sag;
saa vist som Sandhed er en uundværlig Livsbetingelse, vise de
derved, at de ere inde paa en fordærvelig Vej; --- men det er
meningsløst at bilde sig ind, at man støtter Videnskaben, naar
man ikke giver den fuldstændig Frihed. Specialiteter og Kuriositeter
kunne trives, --- Bastarder, avlede af fantastisk Mystik og
forstandig Tænkning, kunne tumle sig, --- gamle Lærdomme
kunne holdes oppe uden videnskabelig Frihed\footnote[1]{I Følge en i Aaret 1857 mellem Paven og Wiirtemberg sluttet Overenskomst sattes det katolsk-teologiske Fakultet i Tübingen under Biskoppens Ledelse og Opsigt. Biskoppen skulde give Professorerne Fuld-}; men de store
Skoler, fra de højeste til de laveste. Forældrenes og Statens Ret er underordnet
Kirken. Og i selve Skolen er (ifølge en Udtalelse af Inkvisitionskollegiet
af 30. Juni 1775) Religionsundervisningen Hovedsagen; alle
% --- p. 445 ---
\setcounter{footnote}{0}%
\opage{445}og dristige Tanker, som føre vor Erkendelse og derved igen
ogsaa Livet frem, vil man forgæves søge, og uden dem bliver
selv det hyggeligste Liv i Hjem, Kirke og Stat dog kun en
kummerlig Ting. Skal Skolen paa rette Maade udfylde sin
Plads i Samfundet, maa Lærerstanden anerkendes som en selvstændig
Stand, jevnbyrdig med Samfundets andre Organer og
uafhængig af kirkelige og politiske Autoriteter. Dens Opgave
er at virke for Folkedannelsen, og dens Forpligtelse til at tage
Hensyn til det Bestaaende i Familieliv, Kirke og Stat maa
være underordnet denne rent menneskelige Opgave. En saadan
selvstændig Stilling for Lærerstanden vil være vigtigere end
Skolelove og Skoleplaner. Ved de højeste Læreanstalter er
dette Maal i det Væsentlige naat; det gælder at gennemføre det
ogsaa for Almenskolen.

Til Videnskabens Frihed og Selvstændighed hører ikke blot,
at den faar Lov til at udvikle sig indenfor sine egne Omraader,
uden at hæmmes af ydre Autoriteter, men ogsaa, at den ikke
hindres i at øve den Indflydelse paa den almindelige Livs- og
Verdensanskuelse, som det ligger i dens Natur at øve. Det
magt til at holde teologiske Forelæsninger og kunde atter tage denne
Fuldmagt tilbage. Han skulde prøve deres Hefter og deres Forelæsningsbøger!
Det akademiske Senat i Tübingen nedsatte en Komité for at undersøge,
om det teologiske Fakultet under disse Omstændigheder endnu
kunde siges at være et Led af Universitetet. Denne Komité kom til det
Resultat, at Professorerne i den katolske Teologi fra nu af ikke kunde
betragtes som Repræsentanter for den frie Videnskab og derfor ikke
vælges til Medlemmer af det akademiske Senat. Biskop Ketteler (Freiheit,
Autoritåt und Kirche. p. 23) er meget forarget herover; men det er
ikke let at indse, hvorledes der kan være Tale om fri Videnskab under
stadig Censur af en Biskop! --- Forskellen er iøvrigt ikke stor, om det
er en levende Biskop eller en død Bog, man skal underordne sig. Hvor
der ikke er ubetinget Frihed for Tænkningen, kan der ikke være Videnskab.
En stadig Vanskelighed gør sig her gældende ved de teologiske
Fakulteters Stilling, idet Menigheder og Kirker fordre Kontrol med, hvad
de lære, altsaa bestemme de Resultater, til hvilke de skulle komme.
Paven f. Ex. vægrer sig ved at anerkende teologiske Fakulteter, som ikke
ere under en saadan Kontrol. I Frankrig bleve derfor de katolsk-teologiske
Fakulteter afskaffede paa Statsuniversiteterne i Aaret \textsc{Isbs}, da man
ikke fandt dem nyttige uden pavelig Anerkendelse.
% --- p. 446 ---
\setcounter{footnote}{0}%
\opage{446}videnskabelige Verdensbillede bestemmer nødvendigvis vor Livsanskuelse
og vor Tro. Ti det er i den virkelige Verden, at
vore Overbevisninger skulle staa deres Prøve, godtgøre deres
Kraft; og denne virkelige Verden kunne vi nu engang ikke
lære at kende paa anden Maade end ved Videnskabens Hjælp.
En Tro, der skyr Sandheden, er en blot Drøm. Vanskeligheden
ligger i, at de Ændringer, den videnskabelige Forskens Resultater
frembringe i en Livsanskuelse, i Regelen foregaa langsomt
og umærkeligt, og at det ikke altid kan afgøres med Sikkerhed,
hvilke de definitive Konsekvenser af et videnskabeligt Resultat
ere. Her stilles da store Krav til Sandhedskærlighed og Ærlighed.
Det sokratiske „Kend dig selv!“ finder her Anvendelse
som en Opfordring til intellektuel Selvprøvelse for at komme
til Afgørelse af, hvor Tyngdepunktet for vor Livsanskuelse ligger,
og i hvilken Grad de overleverede religiøse Forestillinger
ere blevne ændrede ved de Erkendelsesresultater, som ere tilegnede.
Her er Mulighed for et stort Selvbedrag og for mange
Illusioner. Hvad der her ene kan hjælpe, er en virkelig intellektuel
Kultur, en indøvet videnskabelig Tænkemaade, der
med rigtig Takt formaar at fastholde de væsentlige Synspunkter,
at sondre mellem Hypoteser og afgjorte Sandheder, at vove,
hvor der maa voves, og holde sin Dom tilbage, hvor der ikke
kan vides Noget. Den intellektuelle Kultur er en Kunst, som
langtfra findes hos alle Videnskabsmænd. Man kan bore sig
ned i en Specialitet og arbejde med Dygtighed der, og dog
være blottet for al Evne og Takt, hvor man staar overfor andre
Forhold. Enten kaster man sig saa, naar man kommer ud af
sin Hule, i Armene paa den første den bedste Tro, eller man
overdriver de paa det specielle Omraade vundne Resultater ved
at føre Alt tilbage til dem, eller man havner i en mere eller
mindre blaseret Tvivl. Man hævder Videnskabens Almagt, eller
man erklærer, at den, afset fra nogle faa specielle Omraader,
Intet kan lære os om Tilværelsen. Og der svinges maaske
mellem begge Paastande, saaledes at der en Stund tales, som
om alle Gaader vare løste, og naar dette er drevet til Overmaal,
Videnskaben pludselig erklæres bankerot. Overfor saadanne
Svingninger formaar sund intellektuel Kultur --- der be%
% --- p. 447 ---
\setcounter{footnote}{0}%
\opage{447}staar i en Forbindelse af kritisk Tænkning og personlig Kunst
--- ene at bevare Kontinuiteten i den aandelige Udvikling\footnote[1]{Smlgn. min Afhandling Filosofi som Kunst. (Mindre Arbejder. I.).\textsuperscript{1} Adolf Schmidt: Pariser Zustande. III. p. 391.}.
Den betinger Evnen til paa rette Maade at lede Tankens Resultater
over i det personlige og det sociale Liv.

3. Historien viser, at Overgreb fra Hjem, Kirke eller Stat
overfor Skolen ikke naa deres Maal. De bero dels paa en Overvurdering
af den intellektuelle Kulturs Virkninger, dels paa en
Miskendelse af, at direkte Paavirkning ofte endog fører til det
helt modsatte Resultat af det tilsigtede. --- Forstandsoplysning
bestemmer ikke strax Livets hele Retning. Interesser, Følelser
og Drifter ændres ikke strax ved de Kundskaber og Forestillinger,
som indpræntes. „Faren“ er altsaa ikke saa stor og
overvældende, som man mener. --- Og paa den anden Side kan
den ivrige Indskærpen af en vis Række Forestillinger netop
vække en Trang til at slaa ind paa Forestillingsrækker, der gaa
i stik modsat Retning. En saadan Kontrastvirkning faar ofte
den ene Slægt til at gaa i modsat Retning af den anden. Evner
og Drifter, som ikke fandt Næring ved de Forestillinger, hvormed
man under Opdragelsen fik sin Bevidsthed fyldt, trænge
sig da senere saa meget des stærkere og voldsommere frem.

I Frankrig have de skiftende Regeringer i det sidste Aarhundrede
alle søgt at benytte Skolen for deres Formaal, uden
at deres Varighed derfor er bleven større\footnotemark[1]. Det Eneste, man
opnaar, er at forsynde sig mod den intellektuelle Kulturs Tarv.
Grundbetingelsen for al højere Dannelse er, at den maa søges
for sin egen Skyld. ---

Skønt det er nødvendigt, at Staten støtter Videnskaben,
vilde det dog ikke være heldigt, om alle videnskabelige Forskere
i et Folk virkede i Statens Tjeneste. Ligesom hvor Staten
optræder som fordelende Magt paa den materielle Kulturs Omraade,
er det ogsaa her nødvendigt, at den har Konkurrenter i
private Bestræbelser. Naar den videnskabelige Forsken er sund
og frodig, voxer den ud af Livet og er ikke blot fremkaldt ved
den Haandsrækning, Staten kan yde.
% CHECK p. 447: note mark without a note

% --- p. 448 ---
\setcounter{footnote}{0}%
\opage{448}Staten vil dog i Regelen kunne vise større Upartiskhed end
private og kirkelige Institutioner. Netop fordi Staten omfatter
hele Folket, vil dens Blik være rettet mod de universelle Synspunkter
for Dannelsen, og de Fordomme, med hvilke enkelte
Kredse og Sekter træde op overfor Videnskaben, ville ikke saa
let komme til at raade i den. Privatskolen er afhængig af
„Publikum“, Statsskolen af Folket.

Naar Statsskolen ordnes og ledes i sand videnskabelig Aand,
vil den øve en forenende og forbindende Indflydelse. Ungdommen
fra alle Stænder og Klasser vil dér kunne mødes
og stifte Broderskab. Forudsætningen herfor er, at ingen særlig
politisk eller religiøs Retning kommer til at herske i Skolen.
Paa det Punkt, hvor de politiske og religiøse Forskelligheder
begynde, er det Hjemmets Sag at opdrage. Skolen kan kun
yde den Opdragelse, som videnskabelig Kultur bevirker. Det
kan have sin Vanskelighed at drage Grænserne mellem det
fælles Tankeliv, som har sit faste Grundlag i Videnskaben, og
de politisk-religiøse Anskuelser, der væsentligt bestemmes ved
Følelser og Traditioner. Men kun ved alvorlig Bestræbelse
for at fastslaa saadanne Grænser kan den intellektuelle Kultur
hævdes som selvstændigt Element i Menneskehedens Udvikling.

\grouphead{B. ÆSTETISK KULTUR.}

\chapterhead{XXX.}{KUNSTEN OG LIVET.}

1. Vi forstaa en Ting, naar vi se den paa dens bestemte
Plads i Tingenes Række, baaren af andre Ting og selv igen
bærende andre Ting. Videnskaben dvæler kun ved det Enkelte
og Individuelle for at finde Lovene for dets Sammenhæng med
% --- p. 449 ---
\setcounter{footnote}{0}%
\opage{449}den øvrige Verden. Naar vi derimod forholde os æstetisk til
Tingene, se vi dem hver for sig som et ejendommeligt og afsluttet
Hele, glæde os ved Billedet, blot fordi det giver os
noget Ejendommeligt og Karakteristisk. Hvor vi ikke selv kunne
finde eller danne ejendommelige og karakteristiske Billeder,
hjælper Kunsten os dermed.

Baade den intellektuelle og den æstetiske Følelse ere beslægtede
med de sympatiske Følelser, idet vor Lyst og Ulyst i
dem ikke bestemmes ved vor fysiske Selvopholdelsestrang, men
ved Hengivelsen til Noget, der ligger ud over os selv. Men
den æstetiske Følelse har mere Præget af Sympati end den
intellektuelle Følelse. Denne sidste bestemmes mere ved Forholdet
mellem Tingene end ved selve den enkelte Ting i dens
ejendommelige Væsen. Den æstetiske Følelse bestemmes ved,
at vi søge at gøre os til Et med Tingene, leve deres Liv, leve
selve Livet om i vor Fantasi. Der er et æstetisk Velbehag
allerede ved den frie Brug af Lemmerne, ved Legen, i hvilken
vi efterligne Sysler, som ellers foretages i alvorligt og praktisk
Øjemed. Ja, al Kunst kan kaldes en Leg, forsaavidt som den
fører Livet eller Dele af Livet frem for os i et Billede, i en
ideel Gengivelse og opfordrer os til at bevæges ved det, som
om det var et virkeligt Liv.

2. Etisk Betydning faar Udviklingen af den æstetiske Følelse
allerede derved, at den vænner til at betragte Tingene
uden egoistisk Bagtanke, vænner til at glemme os selv over
Synets Værdi. Der virker i det virkeligt Skønne en Magt,
som tiltvinger sig Anerkendelse, og ikke blot Anerkendelse'
men Kærlighed og Beundring. Baade hos den, der formaar at
frembringe skønne Værker, og hos den, der staar som æstetisk
Beskuer overfor Kunstens eller Naturens Værker, ytrer sig en
uinteresseret (d. v. s. uegoistisk) Trang, samtidigt med, at Anskuelsesevnen,
Opmærksomheden og Produktionsevnen sættes i
kraftig og harmonisk Virksomhed. Og den derved vakte Følelse
vil ikke blive indskrænket til Betragtningens eller Produktionens
Øjeblikke, men vil paavirke baade Livsanskuelse og
Livsførelse og her øve en harmonisk og idealiserende Indflydelse.
Hvor den rent moralske Dom og Tugt vilde være for
% --- p. 450 ---
\setcounter{footnote}{0}%
\opage{450}plumpt et Redskab --- i mange af Livets zarteste og individuelleste
Forhold --- der kan den æstetiske eller kunstneriske Takt ofte
føre over Vanskeligheder og afværge Katastrofer. Den æstetiske
Opdragelse af Mennesket, for hvilken \textsc{Schiller} tog Ordet\footnote[1]{Smlgn. Den ny ere Filosofis Historie\textsuperscript{2}. II. p. 128 f.\textsuperscript{1} Smlgn. Erwin Rohde: Psyche\textsuperscript{2} Il. p. 48 f.}, er
dog Mere end blot Middel for det moralske Liv; den fører til
en fri og alsidig Brug af Kræfterne istedetfor den Delthed, de
stridende Elementer i Menneskenaturen medføre, og istedetfor
den Ensidighed, som Arbejdets Deling medfører for det enkelte
arbejdende Menneskes Vedkommende. --- Med de æstetiske Goders
uinteresserede Karakter hænger det sammen, at de ere
fælles Goder, at de kunne deles af Mange, uden at deres Værdi
forringes. Det Fællesskab, de gøre muligt, har sin store etisksociale
Betydning. --- Men især er det af Betydning, at æstetisk
Lyst betinges ved Brugen af Kræfter, som ogsaa det virkelige
Liv lægger Beslag paa. Herpaa beror Alvoren i den
æstetiske Leg. Æstetisk Nydelse er Mere end blot passiv Underholdning.
Den opnaas kun, hvor vi føle os betagne af et
virkeligt Livsbillede. Den etiske Virkning er her ikke forskellig
fra den fuldkomne æstetiske Virkning. \textsc{Aristoteles} har
defineret Tragedien som en Efterligning af en betydningsfuld
Handling, en Efterligning, der vækker vor Frygt og Medlidenhed
og derved renser disse Følelser. Ved „Rensning“ forstod
han sandsynligvis, at Følelserne vækkes og udformes paa harmonisk
Vis ved Hjælp af den kunstneriske Fremstilling, og at
de derved befries fra det Vilde og Kaotiske. Ligesom Musiken
paa én Gang udløser og harmoniserer Følelserne, saaledes
formaar ogsaa den dramatiske Kunst det\footnotemark[1]. Hvad Aristoteles
sagde om Tragedien, kunne vi udstrække til al Kunst. I
Legen eller ved Billedet vækkes de samme Følelser som ved
de virkelige Forhold, der fremstilles; men de vækkes saaledes,
at det Pinlige og Egoistiske falder bort. De virkelige Oplevelser
komme over os, uden at vi altid formaa at finde Sammenhæng
i dem; Tilfældigheden synes at herske. De frembyde
desuden flere forskellige Sider og sætte derfor ogsaa flere 
% CHECK p. 450: note mark without a note
for%
% --- p. 451 ---
\setcounter{footnote}{0}%
\opage{451}skellige Følelser i Bevægelse. Vor Stemning ved Oplevelsen
er derfor ikke „ren“, d. v. s. forskellige, indbyrdes ikke samstemmende
Følelser gøre sig gældende. I det kunstneriske Billede
er derimod Alt anlagt saaledes, at der vækkes en Totalstemning,
som svarer til det karakteristiske Hovedtræk i Tingen
eller Begivenheden. Uden at betage Billedet dets konkret-individuelle
Karakter fremhæver den kunstneriske Behandling det
Væsentlige for os. Vi befris derved fra det Usammenhængende
og faa et samlet Indtryk, kunne derfor ogsaa langt lettere i vor
Følelse gøre os til Et med, hvad der fremstilles. Dermed hænger
det ogsaa sammen, at den kunstneriske Gengivelse lærer os
bedre at „forstaa“ Tingen. Den giver os vel ingen videnskabelig
Forklaring af den, men ved at vise os den i hele dens
Ejendommelighed tydeliggør Kunsten, hvilken Berettigelse denne
Ting har til at være. Hvad den saa ellers er, saa er den en
karakteristisk Del af Verden, eller rettere er selv en lille Verden.
Og denne lille Verden rykkes ligesom ud af Tiden ved
at fæstnes i et Billede; hvad der gik flygtigt forbi som et Mellemled
mellem Fortid og Fremtid, lever i Kunsten et evigt
Liv. Ved Kunstens Hjælp befris vi ikke blot fra Spredtheden
og Tilfældigheden, men ogsaa fra Forgængeligheden. --- Medens
det i Tragedien er Frygt og Medlidenhed, der „renses“, er det
i Komedien Magtfølelsen og Selvfølelsen. Det er ikke personlig
Haan, der faar Udslag i den Latter, den komiske Poesi forskaffer
os, men en Følelse af Befrielse ved at se det Endelige,
Selvmodsigende og Daarlige i dets hele Nøgenhed og dog som
Noget, der nu engang ogsaa hører med til Livet\footnote[1]{Smlgn. Psykologi VI E, 9 c.}.

3. Naar Talen er om, hvorvidt en æstetisk Fremstilling
stemmer med, hvad Etiken fordrer, ville Mange strax tænke
paa det Spørgsmaal, om stærkt sanselige Billeder og Skildringer
ere tilladelige eller ikke. Dette Spørgsmaal er der dog efter
den Maade, hvorpaa vi her have opfattet den æstetiske Kulturs
etiske Betydning, ingen Grund til at opkaste. Der kan, hvad
Indhold og Fremstillingsmaade angaar, ikke være nogen Strid
imellem, hvad Æstetiken virkeligt fordrer, og hvad der er etisk
% --- p. 452 ---
\setcounter{footnote}{0}%
\opage{452}berettiget. Alt, hvad der virkeligt har æstetisk Værdi, maa ogsaa
være etisk berettiget. Vi staa her overfor Livsværdier, der høre
til de Formaal, al Etik skal tjene. Det Skønne hører til det,
der gør Livet værdt at leve; det hører til det samvittighedsfrie
Omraade. (VIII, 1). --- En anden Sag er den pædagogiske Anvendelse
af Kunsten. Man kan ikke give Enhver ethvert Digterværk
i Haanden. Men det har Intet med den æstetiske Værdi
at gøre. Den, hvis sanselige Drifter sættes i stærk Bevægelse
ved en Skildring, kommer ikke til at forholde sig æstetisk til
den: hans Følelse „renses“ ikke, men ægges. Den unge Mand,
der i en af Lukians Fortællinger lod sig lukke inde om Natten
i den knidiske Afrodites Tempel for at omfavne Billedstøtten,
var ikke just bevæget af æstetisk Følelse.

Det følger af sig selv, at naar Kunsten er et ideelt Liv,
vil Kunstværdien svare til Livsværdien, og et Kunstværks Værdi
beror derfor ikke blot paa den Dygtighed eller det Geni, hvormed
Emnet er grebet og udført, men ogsaa paa Værdien af det
Liv, det fremstiller. Her ligger, hvad jeg har kaldet Kunstens
aabne Sted (Psykologi V B, 12), idet Livsvurderingen spiller
ind med. Paa dette Grænsepunkt vil der stedse føres en Strid.
Moralisterne ville her let komme til at tage for tungt paa
Sagen. De gøre let Kunstneren Uret ved at overse, at selve
den kunstneriske Behandling, det kunstneriske Synspunkt stiller
Kunstneren i et andet Forhold til et Emne end det rent praktiske.
--- Kunsten skal ikke anstille nogen etisk Vurdering, skal
ikke direkte moralisere. Men den ægte Kunst vil heller ikke
demoralisere ved paa ensidig Maade at rette Opmærksomheden
mod visse Sider af Livet. Kan Tilskueren forse sig paa et
Kunstværk, saa kan Kunstneren ogsaa forse sig paa Livet. Naturligvis
kan der netop være forskellige Anskuelser om, hvorledes
Livet er, og hvad Opmærksomheden især skal rettes paa.
Enhver betydelig Kunstner har her en Kamp at føre for at gøre
sit Blik paa Tingene gældende. En Kunstners Geni er uadskilleligt
fra hans Personlighed. Grænsen kan ikke drages bestemt
og er stedse unaturlig. Derfor kommer der en uvilkaarlig Vurdering
ind i hans Værk, selv hvor han --- af Overbevisning,
eller af „Fornemhed“, eller af Doktrinarisme --- bestræber sig
% --- p. 453 ---
\setcounter{footnote}{0}%
\opage{453}for at holde al Vurdering ude. Det er denne skjulte og uvilkaarlige
Vurdering (der slet ikke behøver at forudsætte nogen
„Absicht“), som med Rette kan gøres til Genstand for etisk
Kritik. Men den er ingenlunde let at opdage og at konstatere.
Den vil især fremtræde som en uberettiget Generaliseren. Kunsten
kan ifølge sin Natur kun give det Enkelte, Individuelle, Specielle;
ti kun det kan forenes til et levende Billede. Men om det
enkelte Billede giver et Exempel paa en almindelig Regel eller
blot fremstiller en Undtagelse, der skyldes særegne Forhold,
--- derom kan Kunsten som Kunst Intet sige. Derfor kan den ikke
sætte Problemer under Debat, uden forsaavidt den ved sin
Fremstilling kan godtgøre, at noget Saadant som det Fremstillede
dog ogsaa. gives i Verden, ligger indenfor Livets Muligheder.
En virkelig Problemstillen forudsætter et stort og rigt Stof og
megen Reflexion.

Meget i Modstanden mod den moderne æstetiske Realisme
skriver sig sikkert fra, at Mange i Kunsten blot søge Underholdning
og Hvile og derfor ikke ønske at se Tilværelsens bitre
og mørke Sider. Man vil slet ikke have sin Frygt og Medlidenhed
vækket. De gamle Tragedier var man vant til; man
havde fornøden Anvisning til at høre Skæbnens tunge Skridt i
dem; men i en moderne Tragedie som „Gengangere“ formaaede
man ikke at høre dem. Emnet var her for nær paa Livet. Og
dog har den moderne Realisme i sine ypperste Værker egentligt
ikke gjort Andet end at føre de store Synspunkter dybere
ned i det virkelige Liv og vise os Alvoren dér, hvor man var vant
til kun at se Hverdagsliv og Hverdagshistorier\footnote[1]{at de her have villet prøve, om Tragedien definitivt skulde være død. De ville prøve, „om de Smaas og de Fattiges Ulykker vilde tale ligesaa højt til Interessen, Følelsen og Medlidenheden som de Stores og de Riges“.}. Kunsten virker
her opdragende ved at faa os til at lukke Øjnene op, ved at
styrke vor Følelse af Livets Alvor og vække vor Sympati med
dets Lidelse. Den lærer os især i store, anskuelige Billeder
at se den dybe Sammenhæng i Menneskelivet. Uden direkte
Moraliseren kan Digteren blot derved, at han med sikker Haand

*) I Fortalen til „Germinie Lacerteux“ erklære Brødrene Goncourt,
% --- p. 454 ---
\setcounter{footnote}{0}%
% CHECK p. 454: notes paired with the marks in order (the layer misread a number)
\opage{454}peger paa de afgørende Aarsager til Karakterernes Udvikling
eller til Skæbnens Gang, give os en dybere Belæring og en
bedre Anvisning til at vurdere Handlinger og Livsordninger, end
der kan hentes fra nogen filosofisk Etik. Hvad man har sagt
om Goethes „Wahlverwandtschaften“, at det „ikke er skrevet
for Moralens Skyld, men har en Moral“\footnote[1]{Richard Meyer: Goethe. Berlin 1895. p. 286.}, det gælder om ethvert
betydeligt Digterværk. Og vi kunne føje til: det gælder
om det, jo mere det er gennemtrængt af Realismens og Determinismens
Aand, ti des bestemtere kan der peges paa det Punkt
i Begivenhedernes Række, hvorfra de skæbnesvangre Virkninger
udspringe. Paa en gribende Maade peger saaledes Fru Alving
i „Gengangere“ tilbage paa det Punkt, hvor hendes Mands
Ulykke blev til og Grunden til hans Forfald, med alle dets
Konsekvenser, blev lagt. Derved opdukker der et Lys midt i
det Mørke, i hvilket denne Tragedie fører os ind\footnote[2]{Gengangere (3. Akt): „Du kom før til at tale om Livsglæden: og da gik der ligesom et nyt Lys op for mig over alle Tingene i mit Liv\dots{} Du skulde have kendt din Far, da han var ganske ung Løjtnant. I ham var Livsglæden oppe!\dots{} Din stakkars Far fandt aldrig noget Afløb for den overmægtige Livsglæde, som var i ham. Jeg bragte heller ikke Søndagsvejr ind i hans Hjem“.}. Den gamle
græske Tragikers Ord, at de Dødelige maa vinde Kundskab
gennem Lidelse (nd\&si fiå\&og), bekræftes her. Det Tragiske er,
at Lidelsen kan være saa overvældende, at den vundne Kundskab
for den Lidende selv kun bliver et Lys, der falder tilbage paa
Fortiden. Men for den, der ved Kunstnerens Haand følger den
Lidendes Skæbne „i Frygt og Medlidenhed“, kan dette Lys
komme til at straale fremad. --- Digterværker af „idealistisk“
Art have ofte saare liden etisk Værdi, fordi de ikke anerkende
Livets faste Aarsagssammenhæng.

4. Saa meget Kunsten kan og skal være for Livet, saa
maa den dog ikke træde istedetfor Livet, og det virkelige Liv
maa ikke behandles blot som æstetisk Objekt. --- Dette vil
ikke saa let ske med den alvorlige Kunstner, der ikke forholder
sig nydende, men arbejdende, og for hvem Kunsten er en alvorlig
Livsopgave, en social Opgave, idet han, ligesom paa sin
% --- p. 455 ---
\setcounter{footnote}{0}%
\opage{455}Vis Videnskabsmanden, føler sig sat paa en Post for at iagttage
Livet og lære de Andre at se det. Han føler, hvad \textsc{Michel}
\textsc{Angelo} har udtalt, at „den ægte Kunst er ved sig selv from
og ædel formedelst den Aand, i hvilken den arbejder; ti Intet
gør Sjælen saa from og rén, som Bestræbelsen for at frembringe
noget Fuldendt“. Og han føler noksom den Modstand, som
Virkeligheden, det genstridige Stof gør imod, hvad der staar for
ham i hans Fantasi. Han gaar maaske til Grunde ved denne
Modstand. Langt nærmere ligger Faren for de overvejende receptive
Naturer og Dilettanterne. De komme let til at behandle
Livet som en æstetisk Leg. \textsc{Schiller} har sagt, at Mennesket
kun er ret Menneske, hvor det leger. Hans Mening dermed
var, at Fantasiens og Legens Verden er Menneskets eget Værk,
og at der kræves et frit Hjerte og en energisk Villie for at
unddrage sig Virkelighedens Tryk og holde denne ideelle Verden
fast. Men hans Ord ere blevne brugte som Motto for en æstetisk
Livsanskuelse, der forholder sig spøgende og ironisk overfor
ethvert praktisk Forhold, uden dog derfor altid at tage Kunsten
selv mere alvorligt. Arbejdet overlader man til Filistrene; Genierne
i egen Indbildning se fra deres ideelle Højde ned paa Verden
som paa en snurrig, dem uvedkommende Ting. --- Den realistiske
Retning i Kunsten kan ligesaa vel som Romantiken føre til, at
der leves i en indbildt Verden, og at man bliver fremmed for
Virkeligheden. Realisten bevæger sig ligesaa vel i Fantasien
som Idealisten, og Faren kan for ham maaske endog ligge nærmere
end for denne. Den overspændte Idealist vil i Regelen
have Bevidstheden om, at han lever i to Verdener, Drømmeverdenen
og den prosaiske Verden; han haaner denne sidste,
men kan derfor meget godt forstaa at tage den, som den er.
For Realisten, hvis Fantasi søger at mætte sig med Indtryk fra
den virkelige Verden, vil Fristelsen være større til at behandle
alle Livsforhold paa æstetisk Maade.

Denne Fare fremtræder i Historien til enhver Tid, hvor
Kunst og æstetisk Interesse staa i Forgrunden. I hvert Tilfælde
bliver det den Arv, som den næste Slægt modtager; denne
møder „med Dannelsens Fordringer, men uden Oprindelighed
% --- p. 456 ---
\setcounter{footnote}{0}%
% CHECK p. 456: notes paired with the marks in order (the layer misread a number)
\opage{456}og uden levende Skaberkraft“\footnote[1]{Med disse Ord karakteriserer Hettner (Italienische Studien. Zur Geschichte der Renaissance. p. 175) den Slægt, der fulgte efter Rafaels og Michel Angelos Tid. Men han finder Grunden til det bratte Fald i selve den italienske Renaissances Væsen. Den havde sprængt de middelalderlige Lænker, men formaaede hverken ad det etiske Sindelags eller den arbejdende Tænknings Vej at udforme et nyt Menneskeideal. Det var en Tidsalder uden faste sædelige Begreber og Forbilleder, et let Bytte for den kirkelige Reaktion.}. Selv en stor og betydningsfuld
Kunst (som den italienske i Renaissancetiden) kan staa uden
Sammenhæng med de virkelige, i Tiden kæmpende Magter og
uden Sammenhæng med det hele Folks Liv.

5. Kunsten skal give Livets Indhold Form og Klarhed,
udvide Blikket og Sympatien, pege paa den Vej, ad hvilken
Udviklingen skal gaa. Den store Kunstner er tillige Noget af
en Profet. Kunsten kan aldrig blive en blot Privatsag. Hele
Folket og hele Tiden skal i den lære sig selv at kende. Hvert
Folk og hver Tid maa derfor have sin egen Kunst og kunne
ikke ganske leve paa andre Folks og andre Tiders Kunst, saa
stor Betydning det end kan have at lære den at kende. Selv
hvor det Almenmenneskelige udtales og skildres, maa det dog
ske i den særlige Form, som svarer til det enkelte Folks Erfaringer.
Fordringen om en national og tidssvarende Kunst kan
begrundes paa to forskellige Maader. Kunstneren ser og arbejder
kun med sine bedste Kræfter, hvor han er besjælet af
sit eget Hjems og sin egen Tids Følelser og Forestillinger; kun
da kan han tilegne sig, hvad han ser, og gengive det. Og paa
den anden Side forstaar hvert Folk og hver Tid kun ret sig
selv; hvad der skal paavirke dem æstetisk, maa være Kød af
deres Kød og Blod af deres Blod. --- Naturligvis kan det ikke
nytte, at Kunsten er national og tidssvarende, naar den ikke
er virkelig Kunst. \textsc{Julius} \textsc{Lange} har træffende vist, hvor uheldige
Følger det kan have, naar de nationale Fordringer til
Kunsten faa Overhaand over Kunstens egne Fordringer. „Betingelsen
for en fremragende og ypperlig Kunst beror“, siger
han\footnote[2]{Sergei og Thorvaldsen. Studier i den nordiske Klassicismes Fremstilling af Mennesket. Kbhvn. 1886. p. 77.}, „mere paa, at Nationen lærer at gaa ind paa Kunstens
% --- p. 457 ---
\setcounter{footnote}{0}%
\opage{457}ejendommelige Opgaver, end paa Kravet til, at Kunsten skal
gaa ind paa Nationens. Nationen maa ikke blot stille sig i Forhold
til Kunsten som dens Herre; deraf drager den i Grunden
ingen Fordel; den maa først og fremmest gøre sig til dens
Lærling, lære af den at se Aanden i Form, Lys og Farve, lære
den umiddelbare Anskuelse af Mennesket og Naturen“. Der
gives et fælles Tanke- og Følelsesindhold, som ethvert Folk og
enhver Tid maa leve sig ind i og suge Næring af. Gennem
Homer, Dante, Shakespeare og Goethe leve vi den europæiske
Menneskeheds aandelige Liv om igen.

Hvad Kunstens Fremtid angaar, er der --- formedelst Sammenhængen
mellem Kunst og Liv --- kun da Grund til at frygte
for den, hvis Livet selv bliver ringere, --- hvis det føres i
mindre Stil og uden den Kamp og Spænding, store Opgaver og
Ideer føre med sig. Geniet, der paa en Gang er det store
Barn og den store Opdager, vil atter og atter aabne Livets
Dybder trods Fremtidens Filistre, som det har gjort det trods
Fortidens. Det gælder væsentligt om at bevare Sansen aaben
og Hjertet levende; Aandens Verden vil ikke lukkes. Og Haand
i Haand med Kunsten vil den alvorlige Livsførelse, der tager
Livet som den Kamp, det er, virke til at hævde Betingelserne
for den store Livspoesi, som er og bliver den eneste Form, i
hvilken vi formaa at sammenfatte Livets Indhold og Værdi. Her
gaa det etiske Opgør og den æstetiske Trang i Et.

\grouphead{C. RELIGIØS KULTUR.}

\chapterhead{XXXI.}{ETIK OG RELIGIØS FØLELSE.}

1. Etiken maa, som tidligere (II, 3) udtalt, bygge paa saa
faa Forudsætninger som muligt. Den maa ikke forlange nogen
Særstilling i Videnskaben og ikke forsøge paa at rokke ved de
% --- p. 458 ---
\setcounter{footnote}{0}%
\opage{458}Principer, Resultater og Hypoteser, som andre Videnskaber opstille.
Men saaledes som det religiøse Liv har udviklet sig i
Menneskeslægten, har det været knyttet til Antagelser og Dogmer,
der stedse paany have ført til Strid med Videnskaben. Etiken
kunde da synes at være henvist til at træde i Forhold til den
religiøse Kultur som til et rent historisk Fænomen, en fremmed
Magt, som den maa regne med, som den kan kritisere og
vurdere, men som den efter sin Natur og sine Forudsætninger
ikke har noget Slægtskab med. Dog vilde det være en forhastet
Slutning.

Ti for det Første have de historisk fremtraadte Religioner,
i hvert Tilfælde deres højere Former, hver til sin Tid været
etiske Magter. Etiske Forestillinger have udgjort en væsentlig
Del af de positive Religioners Indhold. Fra denne Side set maa
der dog altsaa være et vist Slægtskab mellem Etik og positiv
Religion. — For det Andet er det et Spørgsmaal, om der paa
det Standpunkt, hvorpaa Etiken stiller sig, ikke skulde være
Mulighed for en Følelse, der --- selvom man ikke vil kalde
den religiøs Følelse --- dog i sin psykologiske Natur er beslægtet
med den religiøse Følelse, saaledes som denne fremtræder
i de højere positive Religioner. Jeg vil i det nærmest
Følgende forsøge at beskrive en saadan Følelse, medens jeg opsætter
Undersøgelsen af den positive Religions etiske Betydning
til næste Kapitel.

2. Ved Livsfølelse forstaas i Psykologien den Følelse af
Lyst eller Ulyst, som svarer til det organiske Livs Gang i os.
Den hænger altsaa sammen med den Lethed og Kraft, hvormed
Aandedræt, Blodomløb og den hele Ernæringsvirksomhed foregaar
i os\footnote[1]{i Sammenhæng med de andre Hovedproblemer.}. Efterhaanden som Bevidstheden udvikler sig, bestemmes
Følelsen ikke blot ved vor organiske Tilstand, men

') Smlgn. Psykologi VI A, 3 a. --- Den følgende Beskrivelse af den
religiøse Følelse i Betydning af kosmisk Livsfølelse er en videre Udførelse
af, hvad kort er fremstillet i Psykologi VI C, 8 b. I min Religionsfilosofi
har jeg forsøgt en Drøftelse af det religiøse Problem baade fra det erkendelsesteoretiske,
det psykologiske og det etiske Synspunkt. Og i Den
menneskelige Tanke (p. 369—392) har jeg behandlet det religiøse Problem
% --- p. 459 ---
\setcounter{footnote}{0}%
\opage{459}ogsaa ved et større eller mindre Indbegreb af Forestillinger.
Vor Følelse knytter sig til Mangt og Meget, der ligger ud over
vor Organismes Grænser. I den intellektuelle og den æstetiske
Følelse er vor Lyst og Ulyst bestemt ved Erkendelsesvirksomheden
og ved de Billeder, Naturen eller Kunsten faar vor Fantasi
til at danne. Men naar vi ved Tankens Hjælp have besindet
os paa vor Stilling i Tilværelsen, have indsét, at vi med al vor
Higen, alle vore Planer og alle vore Idealer staa som et enkelt
Led i den store, uoverskuelige Række af Aarsager og Virkninger,
--- saa opstaar der en Følelse ved Livet, ikke blot ved det
Liv, som rører sig i vor egen Organisme, men ved det Liv,
der rører sig i det hele Univers, af hvilket vi ere Led. Vor
Livsfølelse udvides og bestemmes ved Livets og Udviklingens
Gang i Verden, saavidt vi kunne danne os nogen Forestilling
om den. Der opstaar en kosmisk Livsfølelse, som danner en
Analogi til den organiske Livsfølelse. Fra den organiske Livsfølelse
adskiller den kosmiske Livsfølelse sig ved sit Tankeindhold,
som udgøres af alle de Erfaringer, vi have vundet, og
alle de Forestillinger, vi have dannet os om Tingenes Gang.
Fra den intellektuelle og den æstetiske Følelse adskiller den
kosmiske Livsfølelse sig ved sin personlige og reale Karakter.
I Tankevirksomhed og Fantasi glemme vi os selv over, hvad
vi tænke og anskue; i organisk Livsfølelse have vi blot med
os selv som organiske Væsener at gøre; men i den kosmiske
Livsfølelse bestemmes vor Lyst eller Ulyst ved vor hele Personligheds
og vore højeste Livsværdiers Stilling indenfor Verdensudviklingen.

3. Den kosmiske Livsfølelse forudsætter en Verdensanskuelse.
Den behøver dog ikke noget stort Apparat af spekulative Hypoteser.
Hvor saadanne Hypoteser udvikles, ville de endog ofte
ved nærmere Undersøgelse vise sig snarere at være Følelsens
Virkninger end dens Aarsager. Den kan lade sig nøje med de
to Hovedtanker, der ovenfor (XXVIII, 5) fremhævedes som de
vigtigste for den intellektuelle Kultur, nemlig Tanken om Tilværelsen
som en lovordnet Sammenhæng og Tanken om Tilværelsen
som en stor Udviklingsproces.

Ethvert Led i den store Aarsagssammenhæng bæres af
% --- p. 460 ---
\setcounter{footnote}{0}%
\opage{460}andre Led og bærer igen andre. Den Kraft, hvormed det udfylder
sin Plads, kan det ikke have givet sig selv fra først af;
det har modtaget den, og det maa i hvert Øjeblik paany modtage
den for at kunne hævde sig. Selv i min mest energiske
Villiesakt, hvor min personlige Ejendommelighed lægger sig
stærkest for Dagen, og hvor min Selvfølelse synes at have
fastest Grund at bygge paa, selv dér forbruger jeg en Kapital,
jeg ikke selv fra først af har bragt til Veje. Jeg er afhængig
netop i min højeste Aktivitet, og desmere, jo mere Kraft denne
lægger Beslag paa. Men Afhængigheden optræder ikke blot
ved, at den Energi, jeg raader over, er mig given, men ogsaa
ved, at den Energi, jeg raader over, er begrænset. Min Skæbne
er indflettet i den store Udviklingsproces, og min Villie kan
kun for en Del gribe bestemmende ind i denne. Jeg omtumles
af Haab og Frygt under den store Bølgegang af Udvikling og
Opløsning, Opstaaen og Forgaaen, som Naturen frembyder.
Hvem har Ret, og paa hvis Side er Sejren, --- hos den fremadstræbende,
udviklende Tendens i Verden eller hos den hæmmende
og opløsende Tendens?

Forholdet mellem min Villie og min Skæbne bestemmer
altsaa min Følelse ved Livet. Denne Følelse kan være af rent
eller overvejende egoistisk Art. Men naar jeg ikke tænker paa
min Skæbne som dette enkelte Individ, --- naar min Sympati
med Andre og min etiske Følelse er vakt, faar den kosmiske
Livsfølelse en anden Karakter. Disse Følelser, som føre mig til
at opstille og virke for Idealer, der ligge langt udover min individuelle
Selvopholdelse, have dog udviklet sig efter bestemte
Naturlove; og at en saadan Udvikling har været mulig, giver
os et Vidnesbyrd om, at der virker værdifulde Kræfter i Tilværelsen.
(Smlgn. IV, 5). Naturen staar --- netop naar vi anlægge
en naturalistisk Forklaring paa Samvittighedens Opstaaen
og Udvikling --- for os som et Hjemsted for ideale Kræfter.
Hvad Udviklingen saa ellers har ført med sig, dette har den i
hvert Tilfælde ogsaa ført med sig! Der har udviklet sig en
Livstrang og en Livsdrift af anden Art end den blotte fysiske
Selvopholdelsestrang. Hvad jeg føler i mig, i min Samvittig%
% --- p. 461 ---
\setcounter{footnote}{0}%
\opage{461}hed, er ligesaa vel en Verdenskraft, som de Kræfter, der ytre
sig under de materielle Massers Vexelvirkning.

Men her opstaar nu igen en betydningsfuld Vanskelighed.
Vi kunne ikke forklare Tilværelsen ud fra den etiske Følelse.
Verden er ikke konstrueret efter de Principer, vor Samvittighed
stiller for os som de højeste. Hvad der er vort etiske Villieslivs
højeste Formaal, kunne vi ikke erklære for Verdensudviklingens
Formaal. Ikke blot er selve Tanken om et Verdensformaal
videnskabeligt set ligesaa uigennemførlig som Tanken om en
første Aarsag. Men de skrigende Disharmonier i Verden, den
Lidelse og Grusomhed, den Sum af Ulykke og Forbrydelse,
hvormed Udviklingen købes, og Fremskridtet --- forsaavidt der
kan paavises noget Fremskridt --- vindes, alt dette gør det
logisk og etisk umuligt at hævde et etisk Princip som Kilde
til Verdensudviklingen. Ethvert teologisk og filosofisk Forsøg
paa at overvinde denne Vanskelighed har vist sig frugtesløst.
Den ortodoxe Teologi har blot skudt Spørgsmaalet længere tilbage,
og den spekulative Filosofi har forflygtiget og bortforklaret
Vanskelighederne. Den eneste Maade, hvorpaa man kan slippe
fra disse Vanskeligheder, er at lade være at tænke paa dem,
og denne Udvej falder ikke lige let for alle Individer\footnote[1]{Flere teologiske Kritikere have været saa elskværdige at fortolke denne Sætning saaledes, at jeg selv hørte til dem, hvem det faldt let at lade være at tænke paa det omtalte Problem. De have ikke forstaaet, at Sætningen er ironisk ment. Jeg sigtede dels til dem, hvis Interesse overfor saadanne Spørgsmaal er sløvet, saa at de slet ikke beskæftige sig med dem, dels til de Teologer, der tilsyneladende beskæftige sig med dem, men hjælpe sig over Vanskelighederne paa saa let en Maade og med saa overfladiske Argumenter, at denne Beskæftigelse ikke kan kaldes en Tænkning. --- I min Religionsfilosofi'\textsuperscript{2}, p. 243---249 har jeg udviklet de Tanker, som for mig ere de væsentligste overfor det Problem, vi her staa ved.}. Ikke
af Hovmod, men netop af Samvittighed og af klar Indsigt i
vor Erkendelses Grænser maa vi lade dette Spørgsmaal staa
aabent. Filosofen ser ingen Grund til at tage Fornuften fangen
til Fordel for formentlige Løsninger, som for uhildet Kritik
indeholde uovervindelige Selvmodsigelser eller komme i Strid
% --- p. 462 ---
\setcounter{footnote}{0}%
\opage{462}med, hvad der ad Erfaringens eller Tænkningens Vej er vundet,
eller som maaske endog frembyde store etiske Betænkeligheder.
En ærlig Tvivl er bedre end en tankeløs Tro. Livet er fuldt
af store Modsætninger; men de blive først til Modsigelser, naar
man vil gøre ét Modsætningsled til det Hele og ud fra det forklare
Tilværelsen. Af et godt Verdensprincip kan det Onde i
Verden ikke forklares; af et ondt Verdensprincip, som den
absolute Pessimisme stiller op, kan det Gode i Verden ikke
forklares.

Men vor Medfølelse med Lidelsen i Verden og vor Trang
til at fastholde den ideale Betragtningsmaade behøver heller
ikke at rokkes ved, at vi ikke kende Verdensdramaets Afslutning,
eller en Gang vide, om det har nogen Afslutning. Det
Gode og Værdifulde i Tilværelsen er en kæmpende Magt, og
med enhver alvorlig Kamp er der forbundet Uvished og Spænding.
Den kosmiske Livsfølelse vil derfor svinge mellem Frygt
og Haab ogsaa hos dem, hos hvem den idethele har Frejdighedens
Præg. ---

En Følelse som den her beskrevne findes ganske sikkert.
Den betegner et Standpunkt, hvor der er gjort Alvor af, at
Religion er en Følelsessag. Den indeslutter intet Dogme, men
opstaar ganske simpelt ved den etiske Følelses Forhold til den
virkelige Verden. De, der føle Trang dertil, kunne forsøge ved
Hjælp af Spekulationer og Symboler at vinde en yderligere
Forstaaelse; der vil intet Betænkeligt være deri, naar der blot
ikke rokkes ved to Ting: ved Ideen om Tilværelsen som lovbunden
Sammenhæng og ved Etikens Uafhængighed af dogmatiske
Antagelser. Udover den videnskabelige Erkendelse kan
den religiøse Tanke kun hæve sig ved at opstille Fordringen
om en Afslutning af vor Verdensopfattelse, som den stadigt
fremskridende Erfaring umuliggør, og ved at tyde det Verdensbillede,
Erfaringen giver os, i Analogi med og som Udtryk for
Sjælelivets indre Tilstande. Men udover spekulative Hypoteser
kan man ad disse Veje ikke komme\footnote[1]{\textasciitilde{} „Analogi“). --- Den menneskelige Tanke p. 306---345.}. Og langt stærkere end

') Smlgn. Den nyere Filosofis Historie. (Registret under „Ide“ og
% --- p. 463 ---
\setcounter{footnote}{0}%
\opage{463}den teoretiske Interesse for at forstaa Verden vil under den
religiøse Følelses Indflydelse Trangen til Selvforstaaelse være.
Jo stærkere den ved Livets Gang og Vilkaar bestemte Følelse
rører sig, og jo mere dunkelt dens Udspring er, fordi saa
mange smaa og store Livserfaringer bestemme den, desmere
vil Mennesket søge at vinde saadanne Udtryk for den, at hans
indre Liv kan blive forstaaeligt og anskueligt for ham. Og ikke
blot for ham selv, men ogsaa for Andre. Ti han vil ikke
kunne slaa sig til Ro med, at hvad der rører sig hos ham,
skulde være aldeles særligt for ham. Gennem Tegn og Symboler
søger han da at kundgøre for Andre, hvorledes han er
stemt ved Livet, for at se, om der ikke hos dem skulde røre
sig noget Lignende. Gennem den uvilkaarlige Trang til Selvforstaaelse
og til aandeligt Fællesskab blive de religiøse Forestillinger
oprindeligt til. Og paa Grund af Erkendelsens Uafsluttelighed
kunne de ikke blive Andet end Symboler.

Den religiøse Følelse stiller ingen etiske Opgaver, som ikke
vilde stilles uden den. Den udtrykker den inderligste og højeste
Selvforstaaelse, Mennesket kan vinde med Hensyn til sin Stilling
i Verden og med Hensyn til den Skæbne, som vederfares
det, han lever for. Den „renser“ Sindet, ligesom den æstetiske
Følelse; men Rensningen er her mere indgribende, fordi det er
Menneskets egen virkelige Stilling i det virkelige Verdensliv,
som bestemmer Følelsen. Ved sin nøje Sammenhæng med den
etiske Følelse er den kosmiske Livsfølelse fri for Sentimentalitet
og Kvietisme. Vi bestige en Bjergtop, ikke for at bygge og
bo deroppe og unddrage os det virkelige Liv, men for at aande
den rene og stærke, maaske ogsaa skarpe Luft deroppe og for at
udvide og Klare Blikket, saa vi bedre kunne finde os til Rette
i de snevrere Forhold, til hvilke vore praktiske Opgaver knytte
os. Der er et Villieselement, som er nøje forbundet med den
religiøse Følelse, og som hindrer den fra at blive spekulativ
eller sentimental Sværmen. Den mellem Haab og Frygt svingende
Stemning, som Menneskets egen og de højeste Livsværdiers
skiftende Skæbne fremkalder, fører til den ejendommelige
Koncentration og dybe Selvbesindelse, hvori den religiøse
Følelse bestaar; men af denne samme Koncentration udspringer
% --- p. 464 ---
\setcounter{footnote}{0}%
\opage{464}Villien til at holde ud paa den gode Sags Side, til at arbejde
sammen med alle Magter i Tilværelsen, der tjene hine Livsværdier.
Den religiøse Tro, der snart mere har Taalmodets,
snart mere Frejdighedens Karakter, er en Villen, der fastholder
Tanken om Idealet trods al indre og ydre Hæmning og Modstand
og trods al Uvished og Tvivl. Der er her et indre Arbejdsomraade
af central Art, hvor der kan vises Udholdenhed og
Tapperhed ikke mindre end paa det ydre Livs Omraade. Fra
den rent etiske Side ere de Opgaver, som her ligge, dog alle
rede omtalte tidligere (se X, 3 —4).

4. Den beskrevne Følelse rører sig ikke lige stærkt hos
alle Individer. Nogle Naturer ere særligt modtagelige for den,
ligesom andre Naturer for intellektuel eller æstetisk Følelse.
Og den er af en saa sammensat Beskaffenhed, at det ikke kan
undre, om dens Elementer findes i højst afvigende indbyrdes
Forhold hos de enkelte Individer, saa at den fremtræder med
meget forskellig Klangfarve. --- Allerede Forholdet mellem kosmisk
og organisk Livsfølelse kan være højst forskelligt De
kunne ofte staa i afgjort Modsætning til hinanden. Temperamentet
raader ikke ubetinget for den Maade, hvorpaa vi føle
Verden. Den, der er melankolsk af Temperament, og hvis
organiske Livsfølelse altsaa overvejende har en mørk og tung
Karakter, kan derfor godt være Optimist i sin Betragtning af
Tingenes Gang i Naturen og Historien. Det Mørke i Verden
kan være indskrænket til den lille Plads, han selv udfylder, og
ved en Kontrastvirkning kan den øvrige Verden netop komme
til at staa i et lyst Skær for ham. Omvendt er der dem, hvis
lykkelige organiske Grundstemning afbrydes og forstyrres, naar
de kaste Blikket omkring sig i Verden. --- Nogle ville lægge
stor Vægt paa det Gaadefulde og Uudgrundelige i Tilværelsen,
Andre paa den lovmæssige Sammenhæng, Videnskaben efterhaanden
formaar at paavise i Tilværelsen. --- Hos Nogle vil
det være denne, hos Andre hin Klasse af Erfaringer, som især
bestemmer Følelsen, --- hos Nogle f. Ex. især Naturlivet, hos
Andre den menneskelige Historie. --- Hos Nogle (f. Ex. Spinoza)
har den kosmiske Livsfølelse væsentligt Resignationens Karakter,
hos Andre (f. Ex. Bruno og Fichte) har den Entusiasmens og
% --- p. 465 ---
\setcounter{footnote}{0}%
\opage{465}Virksomhedsdriftens Karakter. --- En afgørende Forskel er der
mellem Optimismen og Pessimismen. Livets Gang giver dem
begge delvis Ret, og de ere ikke uforenelige, uden naar de
hver for sig formuleres som spekulative Systemer. Hos Platon
og i den kristelige Teologi paa den ene Side, og hos Schopenhauer
paa den anden Side ere de to modsatte Former for kosmisk
Livsfølelse udviklede til diametralt modsatte Tankebygninger.
--- En betydelig Forskel vil lægge sig for Dagen i
den Originalitet og Selvstændighed, med hvilken de Enkelte
formaa at forme og udtrykke de Forestillinger, i hvilke deres
religiøse Følelse faar Luft. Hos Nogle ville disse Forestillinger
væsentligt afhænge af Familiens eller Nationens overleverede
religiøse Former; hos Andre vil der paa selvstændig Maade
genopstaa et lignende Følelsesliv som det, hvoraf den overleverede
Religionsform udsprang, kun at de ændrede kulturhistoriske
Forhold ville medføre nye Nuancer eller særlig Betoning
af enkelte Sider; hos Andre igen vil Følelseslivet i
Vexelvirkning med Livserfaringerne udfolde sig paa en saa
særegen Maade, at de ikke kunne finde Forstaaelse hos deres
Omgivelser og heller ikke kunne anerkende de gældende Symboler
som Udtryk for den Maade, paa hvilken Livet rører sig
i deres Indre. Mennesker af denne sidste Art maa da gaa
deres ensomme Veje, forsaavidt de ikke have den symbolskabende
Evne og den Dybde i Personligheden, der kan gøre
dem til religiøse Forbilleder og Ledere\footnote[1]{Smlgn. den udførligere Skildring af Forskelligheder i Type og Disposition Religionsfilosofi\textsuperscript{2}, p. 114---128; 258---282.}.

Men trods disse Forskelligheder er der dog her et fælles
aandeligt Omraade, hvor i de store Træk ogsaa en fælles Udvikling
kan være mulig, og hvor den Ene ikke behøver at staa
den Anden i Vejen. Hvor den beskrevne Følelse er ægte, bestemmes
den ved den virkelige Verden og dens Vilkaar, og
det er dog den samme Verden, som fremtræder for Alle, selvom
deres oprindelige Følelsesdispositioner (Temperamenter) og deres
Erfaringer ere forskellige. Naar de ret gøre sig fortrolige med,
hvorledes Verden er, maa de ogsaa tage selve den Omstændig%
% --- p. 466 ---
\setcounter{footnote}{0}%
\opage{466}hed med i Betragtning, at den kan og maa ses fra saa mange
forskellige Standpunkter. Verden kende vi jo dog kun gennem
de Billeder af den, der danne sig i hver Enkelts Bevidsthed;
og i disse Billeder er det ikke de samme Træk, som overalt
fremtræde tydeligst og mest levende; de ses heller ikke alle i
samme Belysning. Verden bliver herved netop rigere og mangfoldigere,
og Sympatien udvides, jo mere man forstaar at sætte
sig ind i Andres Maade at opfatte og føle Verden paa. I ethvert
inderligere Samliv mellem Mennesker maa der gøre sig et
Fællesskab eller dog Fortrolighed og Forstaaelse gældende i
denne Henseende. Den aandelige Hjælp, som det ene Individ
kan bringe det andet, bestaar især i at fremme Sandheden og
Klarheden af hans Følelse ved Livet. Dertil er det ikke nødvendigt,
at det ene Individ staar paa ganske samme Standpunkt
eller tilhører ganske samme religiøse Type som det andet; hvor
der er Sans for Personlighed og for personlig Sandhed, vil man
kunne hjælpe en Anden til fuld Klarhed og til fuldstændigt
Udtryk for sit Væsen, selvom Ens eget Væsen er af en ganske
anden Art. Nogle Individer egne sig særligt til at yde saadan
Sjælesorg, og øve den ofte uden at vide det. Men Ingen er
helt afskaaren derfra. Det almindelige Præstedømme, som
Protestantismen har proklameret, men aldrig gjort fuldt Alvor
af, vil kunne blive Sandhed, jo mere Følelseslivet faar Lov at
udvikle sig frit og ejendommeligt hos de enkelte Individer, og
jo mere samtidigt Sympatien og den psykologiske Sans voxe.
Læger og Præster ere de Nærmeste til her at øve en særlig
Indflydelse. Men de første have endnu for mange fysiologiske
Problemer at løse, til at de ret kunne udvikle deres psykologiske
Sans; de ere for tilbøjelige til at betragte Sjælelivet fra dets
ydre, materielle Side. Og de sidste maa efter deres Standpunkts
Medfør først og fremmest sikre sig, under hvilken Rubrik i
Dogmatiken de skulle henføre de sjælelige Fænomener, de have
for sig; de mangle det frie menneskelige Blik. ---

Selvom man vil kalde den kosmiske Livsfølelse for religiøs
Følelse, saa behøver den i hvert Tilfælde ingen Kirke at stifte
og ingen Kultus at grundlægge, ligesaa lidt som den behøver
at støtte sig til noget Dogme. Det Samfund,  den fører til, er
% --- p. 467 ---
\setcounter{footnote}{0}%
\opage{467}det frieste af alle; det ytrer sig kun i den gensidige Forstaaelse
og Hjælp. Dens Kirke er selve den store Natur, og dens Kultus
er Arbejde, Samliv med Mennesker og med Naturen, Liv i
Videnskab og i Kunst. Der hører i hvert Tilfælde, som vi ville
faa at se, ganske særegne Betingelser til, for at en fælles Symbolik
og en fælles Kreds af ydre Former skal danne sig.

\chapterhead{XXXII.}{POSITIVE RELIGIONERS SOCIAL-ETISKE BETYDNING.}

1. Der gives overordentligt mange Nuancer og Former
for religiøst Liv og religiøs Tro, og naar Frihed en Gang for
Alvor kommer til at herske paa det aandelige Omraade, vil der
gives end flere end nu. Historisk have hidtil de positive Religioner
haft den største etiske og sociale Betydning. To Ting
karakterisere positiv Religion: Kultus og Dogme. Religionen
er ikke fra først af blot Følelse eller et rent subjektivt Forhold
til Tilværelsens Magter; den er en Vej, ad hvilken Mennesket
mener at opnaa bestemte praktiske Formaal, at sørge for sin
legemlige og aandelige Frelse\footnote[1]{da Sprogbrugen ikke er af afgørende Betvdning.}. KuUushandlingen, det ældste
og mest stadige Element i en positiv Religion, har en magisk

I Som jeg min Religionsfilosofi* (p. 231 ---234) har søgt at vise,
er det Afgørende i Begrebet positiv Religion, at der bygges paa historisk
Overlevering, og at der er et Fællesskab i det religiøse Livs Former og
Udtryk. „Positiv“ Religion er at forstaa i Lighed med „positiv“ Ret\&
Kultus og Dogme ere altsaa kun Specielle, historisk betingede Udtryk for
positiv Religion; de ere afledede i Forhold til det Væsentlige i' dette
Begreb. Der kunde tænkes et religiøst Samfund med historisk Overlevering
og fælles Symbol uden Kultus eller Dogmer i disse Ords strenge
Betydning. --- jeg har dog ikke villet ændre Fremstillingen her i Etiken,
% --- p. 468 ---
\setcounter{footnote}{0}%
\opage{468}Karakter, idet Mennesket gennem den mener at gribe ind i en
overnaturlig Tingenes Orden eller selv at staa som Genstand for
Paavirkning fra en overnaturlig Tingenes Orden. I Kultushandlingen
mødes Mennesket umiddelbart med sin Guddom; derfor
staar den som det centrale Element i den positive Religion.
Fællesskab i Kultus er det oprindelige religiøse Fællesskab, og
et Fællesskab af stor Betydning, da det Forbindende er de
højeste og mest afgørende Erfaringer, som kendes. Dogmet er
en Antagelse, som bygger paa en overnaturlig Aabenbaring,
der afslører Tilværelsens Mysterium, en autentisk Forklaring
fra Tilværelsens Kilde om, hvad Meningen med den hele
Tilværelse er. I sin Kultus og i sine Dogmer har den positive
Religion altsaa ikke blot menneskelige Tanker om Tilværelsen
og menneskelige Følelser ved den, men Guddommens
egne Tanker og nærværende Virken. Saasnart Dogmet gøres
til et symbolsk Begreb og Kultushandlingen til en smuk Skik,
ere vi udenfor positiv Religion. Denne staar og falder med
Guddommens punktuelle Nærværelse paa bestemt Tid og Sted
og under en bestemt Form. Den positivt religiøse Bevidsthed
skelner ikke mellem symbolske Forestillinger og objektive Sandheder,
der afsløre Tilværelsens Inderste, heller ikke mellem
symbolske Skikke og Handlinger, der gribe ind i det Inderste
af Tilværelsen. De mægtige symbolske Former, som positive
Religioner have skabt, bleve netop kun mulige ved, at Menneskene
endnu ikke kendte den Forskel mellem Symbol og Virkelighed,
som er Reflexionens og Analysens Værk. Naar dette
Værk er øvet, kunne de store Billeder endnu beholde deres
Værdi som symbolske Udtryk for menneskelige Erfaringer og
Følelser, der stadigt paany gøres; men det er et Spørgsmaal,
om den symboldannende Virksomhed vilde være bleven udløst
fra først af med saa høj Grad af Energi, naar den ikke var
udsprungen af den faste Forvisning om at have med objektive
Realiteter at gøre. Den af Tvivl og Reflexion usvækkede, udelte
Tilstand har været en Betingelse for Frembringelsen af nogle
af de mægtigste Billeder og Former, Menneskeslægten raader
over. Det er en væsentlig Side ved det religiøse Problem, om
den store Produktionskraft kan bevares, naar Distinktionen
% --- p. 469 ---
\setcounter{footnote}{0}%
\opage{469}mellem Symbol og Virkelighed er traadt i Kraft. Men saa Meget
er vist, at der i Kultus og Dogmer er indvævet nogle af Menneskeslægtens
betydningsfuldeste Erfaringer. I Dyrkelsen af
fælles Guder føles Fællesskabet i Familien, Stammen, Folket
eller Menneskeheden. Guden krænkes, naar der tilføjes En af
hans Dyrkere en Krænkelse. I Guddommens Ophøjethed faar
Følelsen ved det Store og Vældige i Livet og i Naturen sit
Udtryk. Og at selv Guderne lide, og maaske lide med Villie,
helliger Lidelsen og gør det lettere at bære den. Stadigt og
uvilkaarligt formedes de Tanker og de Handlinger, der mentes
at have Gyldighed og Betydning for en anden Verden, paa
Grundlag af Paavirkninger, Erfaringer og Stemninger fra det
virkelige Liv. Det er ikke mindst derpaa, de positive Religioners
etiske Betydning beror.

De positive Religioners Udviklingshistorie viser os, at de
staa i et Vexelvirkningsforhold til den praktiske Etik (den positive
Moralitet se I, 2). De have øvet stor Indflydelse paa denne,
men ere ogsaa selv blevne paavirkede af den. De positive Religioners
Tilhængere kunne fra deres Standpunkt kun antage
det Første, ikke det Andet. I deres Religion er den absolute
Sandhed aabenbaret dem; af den skal Menneskelivet lære, men
den har Intet at lære af Menneskelivet. Den historiske Opfattelse
derimod, for hvilken Alt, hvad der fremtræder i Historien,
ogsaa selv er et Resultat af historisk Udvikling, fastholder, at
Religionernes etiske Indflydelse først bliver mulig ved, at de
optage i sig etiske Forestillinger, som have udviklet sig til en
vis Grad uafhængigt af religiøs Indflydelse. Mennesket maa fra
sin egen Erfaring kende de Egenskaber, som det tillægger Guddommen.
I det menneskelige Liv maa Kærlighed og Retfærdighed
have godtgjort deres Betydning, før Tanken om en kærlig og
retfærdig Guddom kan opstaa. Hvorledes skulde Mennesket
ellers kunne forbinde nogen Mening med disse Ord, naar de
bruges om Guddommen? Der finder en Potensering, en Idealisering
Sted: Egenskaberne tillægges Guddommen i en Grad og
Fylde, der overgaar al Forstand. Den positivt religiøse Følelse
er bestemt ved Forestillingen om et eller flere Væsener, der
ere langt ophøjede over menneskelig Natur og menneskelige
% --- p. 470 ---
\setcounter{footnote}{0}%
% CHECK p. 470: notes paired with the marks in order (the layer misread a number)
\opage{470}Vilkaar, hvis Evne langt overgaar, hvad menneskelig Forstand
kan omfatte, men som dog tænkes i Analogi med menneskelige
Væsener. Alle enkelte Egenskaber, som tillægges dem, maa
derfor undergaa en Udvidelse og Forhøjelse, men der maa ligge
Erfaringer til Grund, som kunne udvides og forhøjes. Gudeverdenen
udstyres med de Egenskaber, der staa for den menneskelige
Bevidsthed som de ypperste. Derved forstaa vi, at efterhaanden
som Menneskene forbedres, eller i hvert Tilfælde efterhaanden
som Sæderne mildnes, forbedres og mildnes ogsaa
Gudernes Karakter. Istedetfor de vilde og grusomme Krigsguder
træde efterhaanden Kærlighedens og Barmhjertighedens Guder.
Religionernes Historie viser os en fremskridende Humanisering
af Dogme og Kultus. Religion og Moral have neppe fælles
Udspring. De have sandsynligvis først paa et senere Udviklingstrin
forbundet sig med hinanden: „det religiøse Forhold er i
Tidens Løb blevet moraliseret“\footnote[1]{Chantepie de la Saussaye: Religionsgeschichte. Freiburg 1887. I. p. 35.}. --- Paa lavere Trin kan
Gudernes Hævntrang kun stilles ved Blod. Men de virkelige
Menneskeofre falde efterhaanden bort eller erstattes ved symbolske
Skikke, hvis oprindelige Betydning de Troende sjeldent
ret anskueliggøre sig for deres Fantasi. I den græske Religion
kan denne Udvikling særligt tydeligt paavises. I den jødiske
Religion spillede Sonofre en stor Rolle; men allerede Profeternes
Jehova tørster ikke efter Blod, og Kærlighed er ham bedre end
mange Ofre. Ved Paulus' Indflydelse gik alligevel Ideen om
Sonofret til den vrede Gud over i den kirkelige Kristendom,
og den fastholdes her ofte endnu i en Form, der ganske minder
om de primitive Forestillinger\footnote[2]{Et Præstemøde i Lund (September 1864) udtalte: „Ingen Kristendom uden en Forsoning --- og en Forsoning i Blod!“ og henviste udtrykkeligt til, at „Religionshistorien taler om en Forsoning i Blod gennem de blodige Ofre, der forekomme i saa mange Religioner“.}. — I Forestillingen om Helvede
og den evige Fordømmelse haves et andet Exempel. Selvom
denne Forestilling endnu fastholdes af den ortodoxe Kirkelære,
er det dog langtfra med den glødende Hævnlyst og den levende
% --- p. 471 ---
\setcounter{footnote}{0}%
\opage{471}Fantasi, som hos Oldtidens Kristne\footnote[1]{Tertullian og Cyprian glædede sig til at se deres Forfølgere pines i Helvedes Flammer, medens de selv sad frelste ved Guds Højre. Hos Augustinus og Thomas Aquinas er der allerede indtraadt en Mildnelse, idet Synet af de Fordømtes Kvaler egentligt kun er en Kontrast, der skal bibringe de Frelste en saameget des stærkere Følelse af den Naade, der er vederfaret dem. --- I vore Dage skulle nidkære Sjælesørgere undertiden trøste dem, hvis Nærmeste de antage fordømte, med, at al Erindring om disse Fordømte vil være udslettet hos de Frelste, saa at disses egen Salighed intet Skaar vil lide ved hines Pinsler. Mange ville dog mene, at der ved denne Antagelse ikke betegnes noget meget stort Fremskridt i Humanitet. --- Derimod ytrer et saadant sig i Schleiermachers Indvending mod den evige Fordømmelse: de Saliges Glæde vilde forspildes ved Tanken om Medmenneskers Pinsel!}. --- Grundlaget har forandret
sig; de religiøse Forestillinger, som have uddannet sig under
Indflydelse af tidligere Tiders Følelsesliv, holde sig vel endnu,
men der lægges en anden Betydning i dem, eller de drages ikke
mere frem for Bevidstheden i deres hele Skarphed og Tydelighed.
Sympatien, Menneskekærligheden er voxet, og derfor har
Fantasien ogsaa mistet sin Energi i denne Retning.

2. Det er ikke blot den etiske Udviklings Resultater, som
saaledes omsættes i Dogme og Kultus. Vi forstaa først de positive
Religioners Betydning, naar vi se dem som Fortætninger
af alle Sider af det aandelige Liv. Den intellektuelle og den
æstetiske Udviklings Indflydelse spores ikke mindre end den
etisk-sociales, naar vi sammenligne højere Former for Dogme og
Kultus med lavere. --- De højere Religioner have optaget mere
Fornufterkendelse i sig end de lavere. Forestillingen om Guddommen
bringes stedse til en vis Grad i Harmoni med den
voxende Erkendelse af Natursammenhængen. Udviklingen fra
Fetischisme til Polyteisme og derfra igen til Monoteisme er for
en Del betinget ved voxende Forstaaelse af Sammenhængen i
Naturen, der nøder til at begrænse de overnaturlige Kræfters
Indgriben. Den højere Religion er rationalistisk i Forhold til
den lavere. --- Ogsaa de æstetiske Evner tages i Beslag af Religionerne.
Fantasien søger at danne saa levende, anskuelige og
tiltalende Billeder som muligt af de Væsener, som ere Troens
% --- p. 472 ---
\setcounter{footnote}{0}%
\opage{472}Genstand. Den anstrenger sig for at forestille det Store og Ophøjede.

Den positive Religion er oprindeligt den eneste Form for
ideel Kultur. Det etiske, det intellektuelle og det æstetiske Liv
føres væsentligt i Religionens Form. Alle det aandelige Livs
Sider og Retninger fortættes og udformes i Dogme og Kultus.
Derfor vil positiv Religion heller ikke tilfredsstille nogen enkelt
Side af Menneskenaturen for sig alene, hverken Tanke, Fantasi,
Følelse eller Villie; men den vender sig mod dem alle og er i
sine klassiske Tider Alt for Mennesket i aandelig Henseende.
Der kendes da ingen særlig Videnskab, ingen særlig Kunst,
ingen særlig Etik, og derfor bliver Religionen selv heller ikke
nogen særlig Ting ved Siden af andre Bestræbelser. Religionens
kulturhistoriske Betydning beror paa dens mærkværdige fortættende
Evne overfor saa Meget af det, der staar eller stod
for Menneskene som sandt, skønt og godt. Med den samlede
Kraft, som vindes ved denne Fortætning, virker Religionen paa
Menneskene. Den nærer og opdrager dem med Elementer,
den selv har indsuget og omformet. --- Det er Religionshistoriens
Opgave at paavise de forskellige Fortætningsprocesser,
som her ere foregaaede, og de Elementer, som ere optagne til
forskellige Tider. En stor Religion bliver ikke udformet paa
én Gang eller af et enkelt Menneske. Mange Slægter og forskelligartede
Personligheder yde Bidrag til Kultus og Dogme
gennem uvilkaarlig Projiceren af deres Livserfaringer. Det religiøse
Geni lægger sig for Dagen i Evnen til at indlede en ny
Fortætning af det aandelige Livs Elementer til en given Tid;
og den religiøse Sans bestaar i Trangen til og Modtageligheden
for aandelig Næring i denne Form.

3. Under Kulturens fortsatte Udvikling indtræder der en
Arbejdets Deling paa det aandelige Omraade. Spredning træder
i Stedet for Fortætning, Analyse i Stedet for Syntese. Hver
enkelt Side og Retning af det aandelige Liv lægger nu Beslag
paa særlig Opmærksomhed og Energi. Der opstaar en særlig
Videnskab og en særlig Kunst, og Etiken søger sit selvstændige,
af Dogme og Kultus uafhængige Grundlag. Da opstaar
der Strid mellem Tro og Viden, mellem teologisk og filosofisk
% --- p. 473 ---
\setcounter{footnote}{0}%
\opage{473}Etik, mellem Kirke og Stat. Paa dette Udviklingstrin kommer
den positive Religion til at lide af en Vanskelighed, en Selvmodsigelse,
som fremkaldes ved det, der var dens Styrke paa
de tidligefe Trin: den skal udtrykke og bestemme hele det aandelige
Liv og kommer dog selv til at staa som en speciel Livsytring
ved Siden af Videnskab, Kunst, etisk Liv og Kulturvirksomhed.

En positiv Religions etiske Betydning vil dels bero paa
den bestemte Maade, paa hvilken den forbinder de forskellige
(intellektuelle, æstetiske, etiske) Bestanddele, dels særligt paa,
hvilke etiske Forestillinger den har optaget i sig. Ved dens
Opløsning kan der gaa en betydningsfuld etisk Kraft til Grunde.
Dels er der Naturer, som kun ere modtagelige for aandelig
Næring i fortættet Form, og som derfor nu føle sig splittede
og tomme. Dels er det ikke sagt, at de enkelte Elementer formaa
at virke ligesaa kraftigt udenfor den ejendommelige religiøse
Fortætning, som i denne. Ilten kan jo som Bestanddel af Kulsyren
være med til at fremkalde Virkninger, den ikke formaar
at fremkalde i fri Form. Ligesaa lidt som en positiv Religions
etiske Betydning følger af sig selv, ligesaa lidt følger det derfor
af sig selv, at dens Opløsning er et Fremskridt.

Bestanddelene af en positiv Religions Indhold behøve ikke
at falde bort ved dens Opløsning. De menneskelige Erfaringer,
som finde deres potenserede og idealiserede Udtryk i de religiøse
Forestillinger, kunne stedse gøres paany. Derfor danne
religiøse Motiver ikke nogen fuldstændig Modsætning til andre
Motiver, men kunne indbefatte disse i sig. Menneskekærlighed
er et etisk Motiv; Tro paa Gud som den, hvis Væsen er Kærlighed,
er et religiøst Motiv; men denne Tro forudsætter, at
man af egen Erfaring kender Kærligheden. Dette Forhold mellem
religiøse og etiske Motiver gør, at Sammenligning mellem
de forskellige Religioners etiske Betydning er saa vanskelig.
Det er ikke let at afgøre, med hvilken Grad af Selvstændighed
det enkelte Element virker, og hvor meget der beror paa den
Form og Sammenhæng, i hvilken Religionen lader det fremtræde.

Selvom de Kræfter, der i positiv Religion virke paa kon%
% --- p. 474 ---
\setcounter{footnote}{0}%
\opage{474}centreret Maade, meget vel skulde kunne udrette det Samme
ved at virke hver for sig, saa bliver dog den Ulempe tilbage,
at der mangler et samlet Udtryk for den ideelle Kultur. Der
vil stedse gøre sig en naturlig Trang gældende til at modtage
en samlet Paavirkning af Livet istedetfor stedse kun at have at
gøre med en enkelt Side af det. Arbejdsdelingen fører til Delthed
og Disharmoni, naar nye Fortætningsprocesser ikke kunne
foregaa. Det vil være af væsentlig Betydning for Menneskeslægtens
Fremtid, om saadanne Fortætninger kunne bringes i
Stand, uden at de behøve at antage Dogmets eller Kultushandlingens
Skikkelse. Dertil vilde udfordres en inderlig Vexelvirkning
af Videnskab og Kunst, af Teori og Praxis, af Erkendelse
og Følelse. Menneskelivets store Grunderfaringer og
Grundstemninger maatte finde deres Udtryk i Symboler, som
uden at slaas dogmatisk fast tale ligesaa meget til Følelsen
som til Forstanden, ligesaa meget til Villien som til Fantasien.
Om denne Opgave kan løses, maa Fremtiden vise. Den menneskelige
Bevidsthed er endnu ikke moden dertil. Vi leve foreløbigt
i Kritikens og Analysens Tid, og bære Mærkerne deraf
i vort hele aandelige Liv. Der maa under det overvældende
Indtryk af de religiøse Traditioner sættes saa stort Arbejde ind
paa aandelig Frigørelse, at der ikke let bliver Energi tilovers
til positiv Udformning af en ny Livsanskuelse. Og tillige maa
denne Frigørelseskamp føres af de Enkelte hver for sig. Harmen
over den Maade, paa hvilken den selvstændige religiøse
Udvikling kues og hæmmes, kommer da let til at spille en
større Rolle end Trangen til at lindre den aandelige Nød i
Slægten. En energisk religiøs Tænker, \textsc{Christoph} \textsc{Schrempf},
har --- i Anledning af den 50. Aarsdag for en königsbergsk
Frimenigheds Stiftelse --- udtalt sig herom\footnote[1]{Die Wahrheit. Halbmonatschrift zur Vertiefung in die Fragen und Aufgaben des Menschenlebens. V. Stuttgart 1896. p. 261.} paa følgende Maade:
„Sjælen i denne [den frikirkelige] Bevægelse var, at Pligten og
Retten til Selvstændighed paatrængte sig med Magt. Dette har
givet mig Noget at tænke paa. Ogsaa mig har Selvhævdelsen
bragt i Konflikt med Kirken, og hvor den religiøse Spænding
% --- p. 475 ---
\setcounter{footnote}{0}%
\opage{475}for Øjeblikket er stor, tør man formode, at den især foraarsages
af Uvillien over religiøs Underkuelse\dots{} Jeg ser deri et Tegn
paa, at vi endnu befinde os i de forberedende Stadier, ikke
allerede i den afgørende Bevægelse af den religiøse Bevidsthed.
Denne Bevægelse vil ikke udgaa fra en Stræben efter Selvstændighed,\dots{}
men fra en nyvakt Følelse af Broderskab\dots{}
Navnligt ville de ledende Aander ikke saa meget bevæges af
Bekymringen for deres eget religiøse Selv, som af Medlidenhed
med Folket, der nu igen er som Faar uden Hyrder. Naturligvis
vil denne Medlidenhed, hvis den er af rette Art, igen indskærpe
større Agtelse for ethvert Menneskes Selv. Den vil ikke finde
Menneskets sande Nød i Kalamiteter, heller ikke i den sociale
Ulighed, men i Faren forat miste sit Selv. Den vil altsaa se
sin Hovedopgave i Opdragelsen til Selvstændighed. Men den
religiøse Reforms egentlige Drivkraft vil dog ikke være Harmen
over religiøs Underkuelse, men --- som hos Jesus og Luther
--- Medynken“. Det er her træffende udtalt, at Trangen til
aandelig Selvhævdelse hos dem, der frigøre sig fra den religiøse
Tradition, nutildags let udvikles paa Bekostning af Evnen til
Hengivelse og Medfølelse. --- Endnu en Følge af religiøs Emancipation
har man med Rette henledet Opmærksomheden paa.
\textsc{Stanton} Coit, en af de mest energiske Forkæmpere for en
af Religion uafhængig Etik, siger i en Tale om „Farerne
ved Radikalisme i Religionen“: „Den blotte Bortkasten af
overleverede Lærdomme og Stræben efter egne Ideer giver
Mennesket en vis Tilskyndelse til stadig Forandring. Det maa
stedse have noget Nyt\dots{} Nybagte Fritænkere ere som Fugle,
der, naar de ere slupne ud af Buret, flyve rastløst videre fra
Gren til Gren, fra Dal til Bjerg og fra Bjerg til Sø, uden nogensinde
at stanse for at bygge en hjemlig Rede for deres Hjerte,
men stedse ile videre og videre, indtil tilsidst deres Vinger
lammes, og de falde til Jorden i en øde Ørken“\footnote[1]{von Gizycki. Leipzig 1890. p. 93 f.}. Ogsaa
denne Rastløshed, der gør det vanskeligt for Mange at opnaa
en ny Inderlighed, et nyt aandeligt Hjem efter Opgivelsen af

'i Stanton Coit: Die ethische Bewegung in der Religion. Uebers.
% --- p. 476 ---
\setcounter{footnote}{0}%
\opage{476}det gamle, lune Hjem i Traditionens Ly, er et Moment, der
hindrer positive Nydannelser paa Livsanskuelsernes Omraade.

4. Endnu en saare vigtig Side ved de positive Religioner
maa fremdrages. De Dogmer og Kultusformer, som den religiøse
Fortætningsproces danner, ere fælles for større eller mindre
Kredse af Mennesker. Positiv Religion og Kirke kunne
ikke adskilles. Det er ikke det isolerede Individ, hos hvem
den religiøse Proces gaar for sig; hvad den Enkelte udtaler,
er kun, hvad der mere eller mindre rører sig hos Mange, en
Tilslutning til eller en Videreførelse af, hvad der er givet i
Traditionen. Paa den positive Religions Omraade føle selv de
mest produktive Individer sig ikke som dem, der give, men
som dem, der modtage. Individerne holdes sammen om en
fælles Tradition; og de religiøse Stridigheder opstaa, naar der
bliver Spørgsmaal om, i hvilken Retning Traditionen skal føres
videre. Et positivt religiøst Samfund, en Kirke, opstaar ikke
ved fri Sammenslutning af Individer, der søge hverandre, men
derved, at det samme Ord lyder til Alle. --- I sin simpleste
Form bestaar dette Samfund af en enkelt Familie, indenfor
hvilken de Afdødes Aander dyrkes fra Slægt til Slægt. Arneilden
holdes vedlige til deres Ære, og deres Bud agtes som
højeste Lov. Allerede her er der en hellig Historie. I en saadan
Familiereligions Udøvelse deltage kun de, der høre til
Slægten; Familiefaderen, som leder den, staar som den højeste
levende Myndighed i religiøs som i andre Henseender. Endnu
den Dag i Dag er denne Art Kultus den mest udbredte paa
Jorden, idet den bestaar under, ved Siden af eller som Led af
de nationale Religioner. Kun i Kristendommen og Muhammedanismen
er denne Dyrkelse helt forsvunden\footnote[1]{Henry Maine: Early Law and Custom. p. 57 ff.}. Den nationale
Religion opstod ved, at hele Staten havde sine fælles Guder
og Helte og sin fælles Arne. Staten var ligesom Familien oprindeligt
baade et religiøst og et politisk Samfund. Der var
ingen Forskel paa Stat og Kirke, saa lidt som paa religiøse,
juridiske og etiske Love. Det var enhver Borgers Pligt at deltage
i den offentlige Kultus, men Fremmede vare udelukkede
% --- p. 477 ---
\setcounter{footnote}{0}%
\opage{477}fra den. Hver Stat havde sine Guder, ligesom hver Familie
havde sine. Man kæmpede for sine Fædres Guder, ligesom
for sine Fædres Land. --- Allerede her kan det religiøse Samfund
begynde at vise sig forskelligt fra det politiske, naar der
udvikler sig en særlig Præstestand, hvis Hverv det bliver at
bevare de gamle Traditioner og Skikke. Men afgørende træder
Forskellen først frem, naar Ideen om Kirken som universelt religiøst
Samfund opstaar. Familieforskelligheder og nationale
Skranker faa nu kun underordnet Betydning. Der skal ikke
mere være Jøde eller Hedning, Græker eller Barbar. Et rent
aandeligt Samfund skal grundlægges. Ogsaa her er dog det
religiøse Samfund, den aandelige Enhed betinget ved den fælles
Overlevering af en hellig Historie, fra hvilken de Enkelte hente
den højeste Lov for deres Liv og den faste Betryggelse for
deres Skæbne. For den europæiske Menneskeheds Vedkommende
gjorde Kristendommen denne Ide gældende. Den græske
Filosofi var vel ad sin egen Vej og belært af de historiske Erfaringer
efter Alexander den Stores Tid naat til Forestillingen
om alle Menneskers Lighed og fælles Natur og om almindelig
Menneskekærlighed. Og uden denne hele Forandring i Tænkemaaden
havde der neppe heller udviklet sig en Verdensreligion
af Jødedommen. Her har Udviklingen i den græske Verden
forberedt Kristendommen. Men her viser sig netop den religiøse
Fortætningsproces' Betydning: ti kun som Elementer i en positiv
Religion have Fællesskabet og Menneskekærligheden kunnet
skaffe sig en videregaaende Anerkendelse.

Saaledes har den positive Religion i stedse videre Omfang
godtgjort sin samfundsstiftende og samfundsbevarende Kraft.
Højdepunktet er naat i Ideen om Kirken som universelt Menneskesamfund.
Vi ville da anstille en nærmere Prøvelse af
denne Ide.

5. Det universelle Menneskesamfund, som den positive
Religion forkynder, skal ikke have udviklet sig paa naturlig
Maade; dets Stiftelse er en overnaturlig Akt, der maa gentages,
hver Gang et nyt Individ skal optages i det. Den Enkelte kan
ikke ved egen Bestræbelse faa Adgang til det. Udefra, som
historisk Overlevering kommer Beretningen om dets Stiftelse
% --- p. 478 ---
\setcounter{footnote}{0}%
\opage{478}til ham. Vel skal Samfundet være universelt, men den Tradition,
i hvilken man skal leve sig ind for at være Medlem af
det, strømmer dog gennem en snever Kanal. Den, der ikke
møder denne Strøm paa sin Vej, er fortabt. Alt kommer da
an paa, om Traditionen naar os, og om den er ægte. Herfor
maa vi have Garantier, og Troen paa Garantierne bliver derfor
nødvendigvis det Vigtigste. Positiv religiøs Tro bliver derfor
med logisk Nødvendighed Tro paa Kirken, den Kanal, gennem
hvilken den religiøse Tradition kommer til os. Denne Overgang
fra Troen paa det Garanterede til Troen paa Garantien
ses baade i Katolicismens og i Protestantismens nyere Historie.
Dogmet om Pavens Ufejlbarhed, som først proklameredes halvandet
Aartusinde efter de øvrige Dogmer, er her især karakteristisk.

Ligesom den positive Religion, hvad dens Indhold angaar,
paa videre fremskredne Kulturtrin kom til at lide af den Modsigelse,
at den skal være Alt og dog bliver en særlig Livsytring
ved Siden af andre, saaledes kommer den her, hvad
dens Grundlag angaar, til at lide under den Modsigelse, at den
som Betingelse for det universelle etiske Menneskesamfund opstiller
en Tro, der kan mangle, uden at derfor nogen af de
ædleste etiske Egenskaber behøver at mangle, og som indeholder
saare Meget, hvis Livsværdi det er vanskeligt at paavise.
Ogsaa her (ligesom ved den materielle Kultur, se XXVII, 1)
stønne vi under Vægten af et Apparat, som dog er tilvejebragt
for at gøre Livet lettere at leve! Den, hvis Fornuft ikke tillader
ham at gaa ind paa Dogmet og hylde Kirkens Autoritet, kan
derfor meget vel „hungre og tørste efter Retfærdighed“, og
være „sagtmodig og ydmyg af Hjertet“. Ja, det kan være, at
det netop er Hunger og Tørst efter Retfærdighed, der gør, at
man ikke kan gaa ind paa Dogmet. (Smlgn. XXXI, 3). Der
kan være Punkter i den dogmatiske Lære, som Samvittigheden
trods den redeligste Prøvelse ikke kan anerkende. Overfor
Sligt har Kirken kun det Svar, at det er Selvretfærdighed og
 Hovmod.  Den fordrer Lydighed som første Betingelse. (Smlgn.

X, 4). Den Modsigelse, som her lægger sig for Dagen, er en
% --- p. 479 ---
\setcounter{footnote}{0}%
\opage{479}Modsigelse mellem Tro og Kærlighed. Troen sætter Skranker,
Kærligheden ophæver dem. Troen adskiller, Kærligheden forener.
Denne Modsigelse strækker sig gennem hele Kirkens
Historie. Det var Menneskekærligheden, der sprængte de nationale
Skranker; den vil ogsaa formaa at sprænge de dogmatiske
Skranker. Dens Indflydelse har medvirket til Religionsfriheden.

6. Religionsfrihedens Princip ophæver alle ydre, borgerlige
og politiske Forskelle, som tidligere vare knyttede til Trosforskellighederne.
Det ophæver først og fremmest de Straffe, der
tidligere sattes for Overtrædelse af religiøse Pligter og for Vantro.
Langsomt er paa denne Maade en Række Forbrydelser,
der forhen sattes i første Linie, forsvunden fra de borgerlige
Lovboger\footnotemark[1]. Dog er det ikke at opfatte som et rent negativt
Princip. Det er ikke en Proklamation af Ligegyldigheden, men
en naturlig Følge af det udvidede Fællesskab og af Menneskekærligheden.
Det muliggør ikke blot, at Enhver uforstyrret af
Andre kan trække sig tilbage til sit eget Indre, men det hviler
paa Forudsætningen om et fælles menneskeligt Grundlag bag
alle Trosforskelligheder. Denne Betragtning er af Vigtighed; ti
kun ved at fremhæve den kan man hindre den store sociale
Kraft, som laa i de positive Religioner, i at gaa aldeles tabt.
Der ligger en stor Kraft i den Sammenslutning, som hviler paa
fælles Anskuelse om Livets største Problemer. Naar vi faa Lov
til hver at blive salige paa sin Facon, kan det let føre til, at
vi sejle hver sin Sø og tabe den gensidige Forstaaelse. Men
selvom Ligegyldigheden og Blaseretheden ogsaa finde deres Regning
ved Religionsfriheden, saa er det dog ikke den væsentlige
Side ved den.

I Religionsfriheden ligger for det Første den Forudsætning,
at Etik er uafhængig af Religion. Forskellighed i Trosbekendelse,
Modsætningen mellem Antagelse og Forkastelse af et Dogme
kan ikke være Et med Forskellen mellem Godt og Ondt. Paa
et Standpunkt, hvor Troen er det eneste Princip for alt Godt,
*) Durckheim: De la division du travail social. Paris 1893. p.
% CHECK p. 479: note mark without a note

% --- p. 480 ---
\setcounter{footnote}{0}%
\opage{480}172---177. Ikke-Tro for alt Ondt, er Religionsfrihed utilstedelig. Saa stor
Forskel man end vil gøre mellem det borgerlige Liv og Samvittighedens
indre Liv, saa kan denne Forskel dog aldrig blive
saa stor, at man kunde give Mennesker, der mangle de første
Betingelser for etisk Livsførelse, fuld Ret til aktiv Deltagelse i
det borgerlige og politiske Liv. Der er i det borgerlige og politiske
Liv baade Plads og Brug for de højeste etiske Egenskaber,
og Ingen kan anerkendes som selvstændigt Led af
sit Folk, naar Muligheden af at lægge saadanne Egenskaber
for Dagen hos ham er aldeles udelukket. Efter den strenge
Ortodoxi har den Ikke-Troende jo endog begaaet den største
Synd af alle, Moderen til al anden Synd: den ikke at ville tage
Fornuften fangen. Med en saadan Opfattelse er Religionsfrihedens
Anerkendelse aabenbart uforenelig.

For det Andet hviler Religionsfriheden paa den Forudsætning,
at hvert Individ har sin personlige Ejendommelighed, som
det er af Betydning at faa udviklet saa fyldigt og selvstændigt
som muligt. Hvad han saa tror eller ikke tror, saa skal han
behandles som selvstændig Personlighed, og han skal vinde Sandheden
ad sin ejendommelige Vej. Enhver personlig og ejendommelig
Livsytring har sin Værdi, naar den ikke træder hindrende
i Vejen for mere omfattende Hensyn. Selv ved de
simpleste aandelige Virksomheder gør der sig individuelle Forskelligheder
gældende, som man først i den nyeste Tid er bleven
ret opmærksom paa. Hvor meget mere maa da det højere
aandelige Liv forme sig forskelligt, og hvor forsigtig maa man
være med at drage Konsekvenser fra ydre Ligheder til indre!
Ord, Symbol og Handling kunne være de samme, og dog kan
der være den største Forskellighed i det Indre. Kun den, der
har Evne til at skelne mellem Skal og Kærne, opdager disse
Forskelligheder, denne Rigdom. Den, der stanser ved den positive
eller negative Trosbekendelse, kommer ikke længere end
til Overfladen. Blomsten, som voxer i det høje Græs, opdages
kun, naar man bøjer de lange Straa til Side. Det aandelige
Slægtskab viser sig at være et andet end det dogmatiske. Og
ikke blot kan der bag samme ydre Trosbekendelse ligge et højst
forskelligt indre Følelsesliv; men der kan ogsaa bag forskellig
% --- p. 481 ---
\setcounter{footnote}{0}%
\opage{481}Trosbekendelse ligge et nøje beslægtet Følelsesliv. Det kan være
vanskeligt at trænge ind dertil; vi staa endnu langt tilbage i
psykologisk Sans for det personlige Liv under alle Former, og
de dogmatiske Stridigheder vil endnu i lang Tid være en Hindring
for den fulde Udvikling af denne Sans. Men ved Religionsfriheden
er Vejen banet for dens Udvikling.

7. De fælles Interesser og det fælles Liv i det borgerlige
Samfund og i Menneskekærlighedens Tjeneste ville efterhaanden
føre baade til en fuldere Anerkendelse af Etikens Selvstændighed
og til bedre Forstaaelse af det personlige Liv trods de forskellige
Trosformer. Jo flere Omraader der kan arbejdes paa,
uden at Trosforskellighederne spille nogen Rolle, desmere vil en
saadan Anerkendelse og Forstaaelse udvikle sig, selvom det sker
halvt ubevidst. Her som andre Steder (smlgn. XIII, 4) gaar
fælles Virksomhed forud for Sympati, Praxis for Teori.

Kirken kan fra sit Standpunkt naturligvis ikke gaa ind paa
disse Konsekvenser af Religionsfrihedens Princip. Ligesom den
romerske Kirke vedblivende protesterer mod Ophævelsen af Kirkens
verdslige Magt og mod den Lære, at Kirketugten ikke maa
medføre timelige Straffe, saaledes maa enhver ortodox Kirke
fastholde Forskellen mellem Tro og Ikke-Tro som Livets største
Modsætning og betragte de „relative“ Dyder, der kunne findes
hos Ikke-Troende, som forsvindende i Forhold til hin absolute
Modsætning. Men Livet gaar sin Gang trods alle Dogmer og
udfolder praktisk de Konsekvenser, som dogmatiske Fordomme
hindre i at drage teoretisk.

Derimod er der Intet, der hindrer Fritænkeren i strax at
drage de omtalte Konsekvenser. Vel er Kritikens og Polemikens
Tid ikke forbi overfor de positive Religioner. Her er en Kamp,
som endnu længe stedse maa tages op paany. Men Fritænkeren,
som ikke kan hindres af dogmatiske Fordomme, bør føle,
at han har en langt større Forpligtelse end den Troende til at
vise Menneskekærlighed og Forstaaelse overfor anderledes Troende.
Kun dersom vor Tanke og vor Følelse virkeligt ere blevne udvidede,
har Frigørelsen fra den positive Religions Lærdomme
nogen etisk Betydning. Det ikke at tro Noget er saa omtrent
det Ringeste, der kan siges om et Menneske, og kan ikke i og
% --- p. 482 ---
\setcounter{footnote}{0}%
\opage{482}for sig give ham nogen Værdi, især da en negativ Trosbekendelse
kan antages paa en ligesaa uselvstændig og tankeløs Maade,
som en positiv kan antages. Hos den Fritænker, der tager Livet
alvorligt, er det ikke blot den „kolde Forstand“, der raader.
Det er den frie Samvittighed nok saa meget som den frie Forsken,
der besjæler ham. \textsc{George} \textsc{Eliot} har for sit Vedkommende
udtalt dette paa følgende Maade: „Jeg siger det nu, og
jeg siger det én Gang for alle, at jeg nu bestemmes i min Adfærd
af langt højere Hensyn og af en langt ædlere Forestilling
om Pligt, end det nogensinde var Tilfældet, medens jeg nærede
de evangeliske Anskuelser“\footnote[1]{Life of George Eliot. I, p. 148. (Tauchnitz Edition).}.

Til hvilken Nytte er al vor historiske Forsken, al vor Fordybelse
i forskellige Kulturperioder, naar der ikke derved udvikles
en større og grundigere Forstaaelse af det menneskelige
Liv, som føres omkring os, selvom det iklæder sig Former og
Symboler, der ikke stemme med vore Anskuelser? --- Hvis Fritænkeren
ikke overgaar den Troende i Sympati og i historiskpsykologisk
Forstaaelse af Aandslivet, vil den blotte negative
Kritik ikke være tilstrækkelig til at give ham nogen virkelig
Overlegenhed.

Om selve den religiøse Strid nogensinde vil ende, —om
der til alle Tider vil staa Troende og Fritænkere overfor hverandre,
derom kunne vi Intet vide, og det vedkommer heller ikke
Etiken. Saalænge Modsætningen bestaar, vil der finde en stadig
Vexelvirkning Sted mellem de to Retninger. Nogle Følelser og
Forestillinger begunstiges ved den ene, andre ved den anden
Retning; og dersom disse Følelser og Forestillinger ere af væsentlig
Betydning for det menneskelige Liv, maa de tilegnes
ogsaa af den Retning, der ikke selv har fostret dem. Hvor
Meget har f. Ex. vor Tids frie Forsken ikke lært af Romantiken
og den religiøse Reaktion i Aarhundredets Begyndelse? Og hvor
stor Indflydelse har ikke stadigt Erfaringens og Videnskabens
Indflydelse paa den religiøse Bevidsthed, ikke blot gennem deres
Resultater, men især gennem den Tænkemetode, de langsomt øve
Menneskene i at anvende? --- Paa denne Maade gaar Udviklin%
% --- p. 483 ---
\setcounter{footnote}{0}%
\opage{483}gen --- langsomt og gennem mange Svingninger --- fremad. Men
om der nogensinde naas en blivende Harmoni, kan først en
meget fjern Fremtid vise.

\chapterhead{XXXIII.}{STAT OG KIRKE.\footnote[1]{Kfr. f. Ex. Taines interessante Skildring af den katolske Kirkes Udvikling i Frankrig efter Revolutionen. Le régime nouveau. II p.}}

1. Kirken har ikke blot i Fortiden været, men er endnu
og vil foreløbigt vedblive at være en af de største Kulturmagter,
som virke i det menneskelige Liv. Det har tydeligt vist sig
ved den religiøse Bevægelse i det 19. Aarhundrede, idet Kirken
blomstrede frem til indre Liv og ydre Magt\footnotemark[1], skønt den kort
forud syntes at ligge under for ydre Fjender og at være forladt
af alle aandelige Kræfter. Formedelst sin store Indflydelse
og sin store Livskraft falder den ind under den sociale Etiks
Belysning. Aldeles afset fra den objektive Gyldighed af det
Grundlag, hvorpaa Kirken bygger, er det af afgørende Betydning,
at den kommer til --- indenfor de rette Grænser --- at
yde det Bidrag til Udviklingen, som den formaar at yde.

Kirken har et stort etisk Ansvar, som dens Repræsentanter
med en mærkelig Blindhed netop i den nyeste Tid ikke synes
at gøre sig klart. Jo mere Oplysning og Kundskaber bredes i
videre Kredse, desmere vil uundgaaeligt ogsaa Tvivl og Kritik
overfor Kirkelærens dogmatiske Paastande brede sig. De intellektuelle
Vanskeligheder ved Kirketroen ville føles af stedse Flere.
Men Kirken har netop i den nyeste Tid skærpet sit Standpunkt:
uden Tro ingen Etik! Kirken maa fra dette Standpunkt finde
det at være den rette Konsekvens, naar den, der forkaster Troen,

*) Smlgn. med dette Kapitel min Afhandling Stat og Kirke i Tidsskriftet
„Det ny Aarhundrede“ (1904).
% CHECK p. 483: note mark without a note

% --- p. 484 ---
\setcounter{footnote}{0}%
\opage{484}derpaa kaster sig ud i Bestialiteten. Apostelen Paulus har jo
allerede sagt de skæbnesvangre Ord: „Hvis Kristus ikke er
opstanden, da lader os æde og drikke, ti i Morgen skulle vi
dø“. Dersom denne Tankegang vandt Udbredelse --- og visse
Prædikanter gøre Alt, hvad de kunne, derfor --- vilde det se
uhyggeligt ud for Fremtiden. Kirken ser ikke, hvilket Hazardspil
den her indlader sig paa. I blind Tillid til, at dens Dogmer
ville sejre, sætter den Alt ind paa dette ene Kort. Hvilke Fortjenester
den saa ellers kan have indlagt sig af Menneskeheden,
saa vil det blive et svært Regnskab for den at opgøre, hvis
den ortodox-pietistiske Retning vedvarende kommer til at bestemme
dens Optræden overfor de store Masser, hvis vigtigste
aandelige Leder og Opdrager den endnu er. Dette Regnskab
vil blive den afkrævet i Menneskehedens Navn. Ti enhver Kirke,
lad den kalde sig selv med de stolteste Navne, er dog kun en
Sekt i Forhold til Menneskeheden. Kirken er til for Menneskenes
Skyld, og Menneskene ikke for Kirkens Skyld.

Det er Kirkens Tendens til at gøre etiske Spørgsmaal afhængige
af dogmatiske Principer, der især gør det nødvendigt
at sætte bestemte Grænser for dens Virksomhed. Dermed miskendes
ikke, at Kirken overfor dem, der slutte sig til den, øver
en stor opdragende Indflydelse. Hvor etiske Ideer ikke vilde
kunne virke ved egen Kraft, kunne de virke som Elementer i
Religionen, og hvor intellektuel og æstetisk Kultur ellers ikke
kunne naa hen, blive de mulige i Religionens Form. For Mange
af alle Lag i Samfundet staar Religionen endnu som Indbegrebet
af al aandelig Kultur. Det er ingenlunde saa, at det
kun skulde være „den store Mængde“, der ikke kan undvære
Religionen. Blandt Bønder og Arbejdere findes der ligesaa vel Kritikere
og Tvivlere som blandt de saakaldte Dannede, og Kritiken
og Tvivlen er dér ofte netop ægtest og alvorligst, fordi den er
fremvoxen af upaavirket Tænken og Erfaring. Der findes langt
Flere, end man i Regelen aner, i „den store Mængde“, hvis
Livsanskuelse formes paa selvstændig Maade. De „Dannedes“
Dannelse har ikke altid Selvstændighedens Præg; den bestaar
for Manges Vedkommende i Delagtighed i en vis intellektuel
Tradition, en Tradition, man ofte for tidligt  og med for lidet
% --- p. 485 ---
\setcounter{footnote}{0}%
\opage{485}Selvarbejde optager. Og selv blandt de virkeligt Dannede findes
der utvivlsomt stadigt saadanne, for hvem det er en aandelig
Nødvendighed, at al ideel Kultur, al Forstaaelse af Livet og al
Følelse ved Livet tilsidst fortætter sig i store farverige Symboler,
ved hvilke Forskellen mellem Billede og Virkelighed ikke
længere skal gælde. Det er dette, der gør det religiøse Problem
saa indviklet.

2. Forbindelsen mellem Kirke og Stat er ingen Tilfældighed.
Paa den ene Side maatte Kirken fordre at gennemtrænge hele
Samfundet, naar den skulde give Indbegrebet af al aandelig
Kultur og være det universelle Menneskesamfund. Den maatte
betragte Staten som sin Tjener. Paa den anden Side maatte
Staten --- saalænge de kirkelige Forestillinger havde et ubestridt
Herredømme over Sindene --- mene at have en religiøs Mission
og anse det for sin højeste Opgave at udbrede Religionen og
værne om den. Dog betragtede Staten sig ikke blot som Middel
for Kirken. Den betragtede ogsaa Kirken og Religionen som
sine Midler. Den ansaa religiøs Tro som en nødvendig Betingelse
for, at den kunde faa gode Borgere, og at der kunde
herske Sikkerhed og Fred i Landet. Man forbød Kætteri ikke
blot, fordi det var en Fornærmelse mod Gud, men ogsaa fordi
det var skadeligt mod Rigets Fred og Velfærd\footnote[1]{prudence and Ethics. London 1882). p. 160.}.

Saalænge positiv Religion var den eneste aandelige Kulturmagt,
havde Kirkens Herredømme over Staten eller dog i Staten
--- som Kirkestat eller Statskirke --- sin Berettigelse. Men da
Tiderne forandrede sig, begyndte Kirken selv at indse, at det
religiøse Liv stod sig bedst ved Frihed. Og da nu de moderne
Stater omfatte Mennesker med forskellige religiøse Standpunkter,
--- da mange Omraader, der tidligere kunde varetages af Kirken,
fordrede en selvstændig Behandling, --- saa er Religionsfrihed
mere og mere bleven anerkendt eller fordret baade fra Kirkens
og Statens Side. Forholdet til Kirke og Religion skal da aldeles
ingen Følge have for den Enkeltes Stilling i Staten, ikke
berøve ham nogen af hans Rettigheder eller befri ham fra nogen

') Smlgn. F. Pollock: The Theory of Persecution. (Essays in Juris%
% --- p. 486 ---
\setcounter{footnote}{0}%
% CHECK p. 486: notes paired with the marks in order (the layer misread a number)
\opage{486}af hans Pligter som Borger i Staten. Vi leve dog i denne Henseende
endnu i en Overgangstilstand. Det borgerlige Liv er
endnu ingenlunde frigjort for kirkelig Indblanding. Der staar
mange Former og Institutioner endnu tilbage fra den Tid, da
Staten betragtede enten sig som Religionens Tjener eller Religionen
som sit Middel. Et Par Exempler skal nævnes herpaa.

Ved Eden vilde Staten opnaa at skaffe sig den størst mulige
Garanti for, at et Individ talte sandt eller lovede ærligt.
En saadan Garanti fandt den ved at appellere til det, der maatte
være det Helligste for Individet. Eden er egentligt\footnote[1]{Post: Die Grundlagen des Rechts. p. 441 f.} en Levning
af den hos de Vilde og hos Middelalderens Barbarer saa
meget brugte „Guds Dom“: den, der sværger, nedkalder en
overnaturlig Straf over sig. Men Eden skulde ingen Trosbekendelse
være. Det faldt Ingen ind, at noget Individ kunde undlade
at tro paa en overnaturlig Straffemyndighed. Man henviste
trygt den Enkelte til Forudsætninger, som man antog, at Alle
delte. Dersom Eden var en Trosbekendelse, vilde dens Bestaaen
som Institution være en grov Krænkelse af Religionsfrihedens
Princip, saaledes som dette er slaaet fast i de fleste moderne
Staters Grundlove. Naar Edsinstitutionen bestaar efter Anerkendelsen
af Religionsfriheden, maa det være berettiget at opfatte
den som en Opfordring til at afgive sin Erklæring eller sit
Løfte efter den alvorligste og samvittighedsfuldeste Prøvelse og
Selvbesindelse, idet Enhver fæster Tanken paa det, der for ham
staar som det Helligste. Man aflægger da Eden i samme Aand
som den oprindeligt indstiftedes. Men saalænge Eden anvendes
i sin gamle teologiske Form, staar Døren aaben for juridisk
Chikane. Man paastaar, at der i Eden ligger en Trosbekendelse,
og faar derved et Middel til at udelukke dem, der ere
ærlige nok til at vedkende sig deres frie religiøse Overbevisning,
fra den fulde Nydelse af deres borgerlige Rettigheder\footnote[2]{Et engelsk Exempel paa saadant Misbrug af overleverede Edsformularer meddeles hos F. Pollock: The Oath of Allegiance. (Essays p. 187. 197). --- I den danske Rigsdag er det blevet udtalt af en højtstaaende Jurist, at hvis der skal ligge en forhøjet Garanti i Eden, maa denne være mere end en højtidelig Forsikring, maa være Anraabelse af en Magt, som kan „gribe ind i de menneskelige Vilkaar i dette eller det kommende Liv“, som idethele kan „gøre En Noget“. (Folketingets Møde den 20. December 1880). Der synes her at ligge et etisk Standpunkt til Grund af lignende Art som den unge Piges i Poul Møllers „En dansk Students Eventyr“. Hun var saa bange for Helvede, men beroligedes ved, at det ikke var Andet end --- den onde Samvittighed. Hvad kan den „gøre En“?}.
% --- p. 487 ---
\setcounter{footnote}{0}%
% CHECK p. 487: notes paired with the marks in order (the layer misread a number)
\opage{487}Man ser ikke, til hvilke Konsekvenser dette vil lede. Biskop
\textsc{Martensen}, der glæder sig meget over, at Staten, som ellers
„i saa mange Henseender har givet efter for den religionsløse
Humanitet“, ved at anvende Edsinstitutionen indrømmer, at den
ikke kan undvære den religiøse Garanti, gør med fuld Ret opmærksom
paa, at denne enkelte Garanti, Eden, bliver uden
Betydning, naar man ikke sørger bedre for Folkets Religiøsitet
idethele\footnote[1]{Individuel Etik. p. 278.}. Altsaa: Staten maa sørge for, at alle Borgere have
de fornødne dogmatiske Forudsætninger; disse maa derfor af al
Magt indskærpes hos Alle. Man ender her konsekvent med
Inkvisitionen. --- Med Religionsfrihedens Princip stemmer det
bedst, naar Formularen ændres saaledes, at den ingen teologiske
Udtryk indeholder. Man undgaar derved ogsaa den Fare, som
allerede \textsc{Kant} gjorde opmærksom paa, nemlig at bibringe Folk
den Tro, at Løgn ikke er saa slem, naar man blot ikke sværger
paa den.

Et andet Exempel frembyder Formen for Ægteskabets offentlige
Indgaaelse. Her er den eneste med Religionsfriheden fuldt
stemmende Ordning den, at der gives en fælles, rent borgerlig
Form for Indgaaelsen, medens det nu i mange Lande kun er
dem, der ikke ønske kirkelig Vielse, som blive borgerligt viede.
Det Naturlige --- og tillige det med den religiøse Aand stemmende
--- er dog, at de, der ønske kirkelig Vielse, kunne føje
den til den borgerlige, ikke at de, der ikke ønske den, skulle
skaffe sig fritagne for den. Ved at lade Valget staa aabent
mellem borgerlig og kirkelig Vielse, slaar man desuden\footnote[2]{Smlgn. E. Zeller: Staat und Kirche. Berlin 1873. p. 226.} fast,
at de ere Modsætninger, der udelukke hinanden, og nærer derved
de religiøse Fordomme. --- Man har ved Diskussionen om dette
% --- p. 488 ---
\setcounter{footnote}{0}%
\opage{488}Spørgsmaal altfor meget glemt, at det hverken er den kirkelige
eller den retslige Sanktion, der etisk stifter Ægteskabet. Ægteskabet
stiftes ved fri Pagt mellem Mand og Kvinde, og Kirkens
eller Statens Anerkendelse er af underordnet Betydning. (Smlgn.
XVII, 4).

Lignende Exempler kunne hentes fra Begravelsesskikkene,
fra Kirkens Herredømme i Folkeskolen o. s. v. „Men de anførte
Exempler ere tilstrækkelige til at vise Nødvendigheden af, at
der tilvejebringes en Ordning af det borgerlige Livs Forhold,
som er uberørt af teologiske Forudsætninger.

3. Det religiøse Liv, som Kirken værner om, vil blive
sandere og mere intensivt, jo mere den opgiver at herske i
den borgerlige Verden. Den historiske Udvikling i de sidste
hundrede Aar viser tydeligt nok, hvor dyb en Rod Kirken endnu
har i Folkene. Samtidigt med, at dens ydre Forrettigheder efterhaanden
falde bort, er dens indre Liv blevet kraftigere. Hvor
den helt eller delvis er skilt fra Staten, synes den at have
størst Indflydelse paa Sindene. Den Forbitrelse, som i det 18.
Aarhundrede vaktes mod Kirken, angik langt mere dens ydre
Herredømme end dens Dogmer og Kultus. Den katolske Kirkes
Historie efter Reformationen, Kirkens Stilling i Nordamerika,
den tyske Kirkes Stilling efter Civilstandsloven af 1875 viser,
at Kirken selv meget godt kan staa sig ved Religionsfrihedens
Gennemførelse. Man har fra kirkelig Side næret for stor Frygt
for, og fra mange Fritænkeres Side bygget for store Forhaabninger
paa Kirkens Adskillelse fra Staten.

Egentligt er det snarere Kirkens Udskillelse af Staten end
en Skilsmisse mellem Kirke og Stat, som i det Foregaaende er
fordret. Der ligger i det Foregaaende kun, at det borgerlige
Liv i Staten skal ordnes efter almindelige menneskelige, ikke
efter særlige religiøse Forudsætninger. Staten maa forudsætte,
at dens Borgere, selvom de ikke høre til nogen Kirke, have
de fornødne etiske Forudsætninger til at udøve de Rettigheder
og Pligter, der tilkomme dem. Men dermed er endnu Intet afgjort
om, hvorledes Kirkens fremtidige Stilling skal være. Fordi
den ikke skal raade i Staten og de almindelige borgerlige Forhold,
er det ikke sagt, at Staten slet Intet har at gøre med
% --- p. 489 ---
\setcounter{footnote}{0}%
\opage{489}den. Saalænge Kirken endnu lever i Folket, er den en Kulturmagt
af stor Betydning, som Staten ikke kan stille sig fuldstændigt
ligegyldig overfor.

Selvom der indenfor Kirken rejste sig en stærk religiøs
Bevægelse, en Begejstring som den første Menigheds, saa at
man gav Afkald paa enhver Støtte og Hjælp fra Statens Side,
selvom Kirken altsaa blev en ren Privatforening, vilde Staten
dog derfor ikke ganske kunne tabe den af Sigte. Staten maatte
føre en vis Kontrol med den for at forebygge, at der fra den
udgik Farer for den offentlige Fred og for det etiske Grundlag,
hvorpaa Staten hviler. Den maatte hindre fanatiske Overgreb
paa anderledes Troende og Undertrykkelse af Kirkens egne Tilhængere.
Den vilde her, ligesom overfor Familien, have den
Opgave at beskytte de Svage mod Overgreb, selv indenfor private
og meget inderlige Forhold. Staten repræsenterer det almindelige
etiske Princip overfor de specielle religiøs-etiske Principer,
som de forskellige Religionssamfund hylde. „Sædeligheden
staar som en almindeligere Lov over Dogmerne og de religiøse
Sætninger“\footnote[1]{Denne Udtalelse fremkom paa den danske grundlovgivende Rigsdag (i Mødet den 3. Maj 1849) fra den daværende Kultusminister (Madvig) under Debatten om den nuværende Grundlovs § 76.}. Saadanne religiøse Foreninger eller Ordener, der
forstyrre Freden mellem de forskellige Trosbekendelser og angribe
selve Statens Forudsætninger, vil Staten nødsages til at
have et vaagent Øje med, skønt den vil hævde sin Værdighed
bedst ved at overlade den frie Udvikling at føre ud over de
Misligheder, som her kunne fremkomme. --- Jeg ser her bort
fra, at Staten maa afgøre, om religiøse Privatforeninger ere at
anerkende som juridiske Personer, om de skulle have Ret til
at besidde fast Ejendom, om de skulle have Ret til at undervise
o. s. v. ---

I Lande, hvor Kirken har slaaet Rod i Folket og gennem
Aarhundreder har gennemsyret dets Liv, vil det være naturligt,
om Staten vedblivende yder den Understøttelse som et Samfund,
indenfor hvilket den største Del af Folket finder den vigtigste
Næring for sin aandelige Trang. Staten betragter altsaa Kirken
som historisk given Kulturmagt, og den Understøttelse, den yder
% --- p. 490 ---
\setcounter{footnote}{0}%
\opage{490}Kirken, er bygget paa de samme Motiver som den, der ydes
Videnskab og Kunst. Staten kan ikke direkte producere Kultur
af nogensomhelst Art. Dens Virksomhed i Kulturens Tjeneste
er stedse indirekte. Den kan hverken producere eller udrydde
Religion; men den kan yde materiel Støtte og retslige Former
for det religiøse Samfund. Om den foruden den i Folket historisk
overleverede Kirke ogsaa vil støtte andre Religionssamfund,
vil bero paa den Betydning, den tillægger disse. --- Som Betingelser
for den materielle og retslige Støtte, Kirken modtager
fra Staten, maa denne kræve Ret til at øve en mere indgaaende
Kontrol, end der kan øves overfor en ren Privatforening. Der
er især to Fordringer, som Staten i den aandelige Kulturs og
i Aandsfrihedens Interesse bestemt maa stille.

Det kan ikke være Staten ligegyldigt, hvilken Forfatning
Kirken har. Den maa arbejde hen til en Kirkeforfatning, der
kan tjene til at fremme Folkets hele Liv og Udvikling. En
Forfatning, hvorved der skabtes et Hierarki, der øvede et uindskrænket
aandeligt Herredømme over Kirkens Medlemmer, vil
Staten ikke kunne gaa ind paa. Den kan ikke begunstige Aandstyranni
og systematisk Indsnevring af Folkets aandelige Horizont.
Den vil saavidt muligt begunstige Lærefriheden og arbejde hen
til, at Folket kan samle sig om de religiøse Lærere, der vinde
dets frie Tilslutning. Den vil ikke i Længden kunne tillade
Kirken at fordre et bestemt Læresystem lagt til Grund af alle
religiøse Lærere trods den store Usandsynlighed for, at de alle
med fuld og selvstændig Overbevisning skulde have tilegnet
sig det. Den kirkelige Enhed og det kirkelige Fællesskab købes
for dyrt, naar de, der skulle lede den religiøse Udvikling, uvilkaarligt
eller vilkaarligt foranlediges til at afholde sig fra at
prøve det, de skulle lære, saa uhildet, som Sandhedskærligheden
fordrer\footnote[1]{Et oplysende Exempel fra den nyeste Tid frembyder en Kendelse af den würtembergske „Disciplinarret for evangeliske Gejstlige“ af 21. Febr. 1896, hvorved en Præst (Pastor Steudel) blev afsat, fordi han i „væsentlige“ Punkter afveg fra den augsburgske Konfession og Liturgien. Det blev i Dommen lagt denne Mand til Last, at „den ensidige Udformning og Pleje af hans etiske Betænkeligheder tildels havde sin Kilde i}. --- Staten vil tillige fordre, at de, der skulle være
% --- p. 491 ---
\setcounter{footnote}{0}%
\opage{491}Lærere i Kirken, faa en vis videnskabelig Dannelse. Det er af
ikke ringe Betydning, at de unge Mænd, der skulle træde i
Kirkens Tjeneste, gøre deres Studier ved en videnskabelig Højskole
sammen med Landets øvrige studerende Ungdom. De lære
Livet at kende fra flere Sider. Selve det kammeratlige Samliv
gør altfor stor Ensidighed umulig, og de faa en vis videnskabelig
Samvittighed foruden den religiøse Samvittighed.

Den Indflydelse, Staten paa denne Maade kan øve paa
Kirken, støttes ved dens Virken for Kulturen paa andre Omraader.
Jo bedre Folkeundervisningen er, og jo frodigere det
videnskabelige Liv blomstrer i Folket, des større Fordringer vil
der stilles til de Gejstlige, og des mindre vil man finde sig i
et Hierarkis Herredømme. Staten kan støtte de enkelte Kulturmagter,
de intellektuelle, æstetiske og religiøse; men det indbyrdes
Forhold, hvad enten det bliver harmonisk eller disharmonisk,
maa den overlade til Kræfternes frie Vexelvirkning. Den
kan skaffe Midler og materielt Grundlag til Førelsen af en
aandelig Kamp, men ikke selv gribe bestemmende ind i denne
Kamp. Til Folket høre Alle, Troende og Ikke-Troende, og Folkets
Liv udfolder sig under den stadige Brydning mellem de
forskellige aandelige Retninger.

4. Hvorlænge et saadant positivt Forhold mellem Stat og
Kirke vil vedblive at være naturligt, maa de praktiske Statsmænd
afgøre. Dette Spørgsmaals Afgørelse vil være forskelligt
fra Land til Land; den beror paa Folkets Karakter og paa dets
Traditioner. Der er dog flere forskellige Forhold, der mere og
mere ville virke i Retning af en Udskillelse af Kirken fra
en vis Overvurdering af sig selv, ved hvilken han førtes til at gøre en
Konflikt, som Hundreder ved Siden af ham udkæmpe i Stilhed paa en
eller anden Maade, til Genstand for et offentligt Martyrium“. Den høje
kirkelige Myndighed indrømmer altsaa her, at Hundreder af Præster befinde
sig i en alvorlig Konflikt med Kirkelæren, en Konflikt, ved hvilken
ikke blot intellektuelle, men ogsaa etiske Betænkeligheder spille en Rolle.
Det er meget usandsynligt, at dette skulde være Noget, der kun forekommer
i den würtembergske Kirke. --- Smlgn. om denne Sag Christoph
Schrempf: Ein kirchliches Ereigniss. (I det af Schrempf udgivne 
% CHECK p. 491: note whose mark was not found
\footnote[1]{skrift Die Wahrheit. VI. Stuttgart 1896).}Tids%
% --- p. 492 ---
\setcounter{footnote}{0}%
\opage{492}Staten som det Naturlige. Religionsfrihedens Gennemførelse, og
i Sammenhæng med den Gennemførelsen af rent menneskelige
Livsformer virke stadigt i Retning af at isolere Kirken. Indenfor
Kirken selv vil den strenge Betoning af Troslærdommene i
deres bogstavelige Opfattelse frembringe en stigende Modsætning
til Kulturlivet paa andre Omraader. Og paa den anden
Side vil Kirkens Bestræbelse for saavidt muligt at virke ved
frivillige Kræfter og frivillige Tilskud mere og mere stille den
uafhængigt, saa at den ikke behøver at fremme Troens (eller
de Troendes) Sag ved Midler, Staten opkræver fra de Ikke-Troende,
der af en eller anden Grund ikke have udmeldt sig
af Statskirken.

Spørgsmaalet er, om en fuldstændig Udskillelse af Kirken
vil være et Gode for Folkets Liv. Selvom den bliver naturlig
og nødvendig, er det jo ikke givet, at den betegner et Fremskridt.
Dersom Kirken faar sine egne Seminarier og sin egen
præstelige Forfatning, vil Fanatisme og Bornerthed brede sig.
Spaltningen i Folket i religiøs Henseende vil blive større, og
Religionen træde i den skarpeste Modsætning til den øvrige
Kultur. De største Farer kunne udspringe deraf for Freden og
Kulturen i Folket. Som Erfaring viser, kunne snevre, barbariske
og livsfjendtlige Lærdomme faa Indpas i Kirken trods den dæmpende
og humaniserende Indflydelse, Forbindelsen med Staten
medfører; men værre vil det sikkert blive, hvis Præsterne helt
afløses af Lægprædikanter, og hvis Sammenhængen med Folkets
Aandsliv paa andre Omraader afbrydes. Kun dersom Kirken
en Gang for alle skulde anerkende og optage saadanne Retninger,
vilde det kunne blive absolut betænkeligt for Staten at
vedligeholde Forbindelsen; 'saa vilde Kirken ikke længere kunne
øve nogen sand folkeopdragende Magt og vilde ikke længere
være en af den ideelle Kulturs Repræsentanter.

Fra Kirkens Side viser der sig i den nyeste Tid en iog
for sig naturlig Tilbøjelighed til at skaffe sig de Fordele, Religionsfriheden
medfører, og dog beholde de Fordele, Stillingen
som Statskirke medfører. Den fordrer fuld Frihed til at ordne
sine egne Sager, fuld Selvstyrelse, og dog samtidigt saa megen
Indflydelse som muligt paa det borgerlige Liv og saa megen
% --- p. 493 ---
\setcounter{footnote}{0}%
\opage{493}Understøttelse som muligt af Staten. Kirken fremtræder, som
man rigtigt har bemærket\footnote[1]{Zeller: Staat und Kirche. p. 132.}, nutildags ofte i Religionsfrihedens
Navn med Fordringer, som den i tidligere Tider udledte af sin
guddommelige Mission. Det er ikke blot Katolicismen, som paa
denne Maade benytter sig af Frihedens Navn\footnotemark[1]. Ogsaa den
protestantiske Statskirke har rejst Besværing over, at medens
Sekterne have Ret til selvstændigt at ordne deres Anliggender,
saa staar den under Statens Myndighed og er altsaa det eneste
religiøse Samfund i Landet, der ikke har Religionsfrihed!\footnotemark[1]. Det
vil med andre Ord sige: man vil fra kirkelig Side helst have
baade i Pose og i Sæk og indser ikke, at jo mere man faar i
Sækken, des mindre kan man faa i Posen. Et Enten---Eller er
idethele det, som nyere Teologer synes at have størst Vanskelighed
ved at forstaa. Kun hvis Kirken vil vende tilbage til
en rent Privatforenings Stilling, vil den kunne faa en saa godt
som fuldstændig Selvstyrelse. Naar den tager imod Understøttelse
fra Staten, forandres Forholdet.

Om en fuldstændig Skilsmisse mellem Kirke og Stat er et
Gode, vil bero paa, under hvilke Forhold den sker. Den er
ikke, som især mange Fritænkere have ment, et Tryllemiddel,
som bortrydder det religiøse Spørgsmaals Vanskeligheder. Den
overdrevne Betydning, man har tillagt Skilsmisse af Kirke og
Stat, hænger sammen med den overdrevne Betydning, man
idethele tillægger Staten. De religiøse Problemer skaffes ikke
ud af Verden, fordi Staten trækker sig tilbage fra Kirken, naar
denne endnu har Rod i Folket. Det er ikke Præsteskab og
andre Institutioner, der frembringe og opretholde Religionerne;
men det er Folkets Trang og Tilslutning, der gøre Præsteskabets
og Kirkens Bestaaen mulig). Sandsynligvis vil Baandet mellem
% CHECK p. 493: note whose mark was not found
\footnote[2]{'-') Ketteler: Freiheit, Autorität und Kirche. p. 159 ff.}
% CHECK p. 493: note whose mark was not found
\footnote[3]{Smlgn. Artiklen „Religionsfrihed'c i 2den Udgave af „Nordisk Konversationslexikon“, og Martensen: Social Etik. p. 417 ff.}
% CHECK p. 493: note whose mark was not found
\footnote[4]{Frederik den Store protesterede mod den Antagelse, at Fyrsterne havde forbundet sig med Præsterne for at udbytte Folket i Fællesskab: det var tværtimod Folkets Overtro, der gav Præsterne en Magt, ved hvilken ogsaa Regeringernes Hænder bleve bundne! (E. Zeller: Friedrich als Philosoph. p. 32).}
% CHECK p. 493: note mark without a note

% --- p. 494 ---
\setcounter{footnote}{0}%
% CHECK p. 494: notes paired with the marks in order (the layer misread a number)
\opage{494}Stat og Kirke blive løsere og løsere. Allerede nu arbejde frivillige
Kræfter ivrigt og i stort Antal i Kirkens Tjeneste. Selv
den engelske Statskirke, den stiveste af alle, bygger nu for en
meget stor Del paa „Voluntarismen“\footnote[1]{Arthur Elliot: The State and the Church. London 1882. p. 127 f.}. Og i de nordamerikanske
Fristater foreligger et mærkeligt Exempel paa, hvorledes Kirken
(eller Kirkerne) efter Adskillelsen fra Staten paa Frivillighedens
Grund har bevaret og forøget sin aandelige Indflydelse og sine
materielle Midler. „I Løbet af blot tre Slægtled har det amerikanske
demokratiske Samfund gennemført den fuldstændige Adskillelse
af Kirke og Stat, en Reform, som intet andet Folk
nogensinde har forsøgt“. „En stor Del af den moralske Energi,
som tre Slægtled have været i Stand til at spare fra Arbejdet
paa materielt Udkomme, er bleven anvendt til Gennemførelsen
af Frivillighedssystemet i Religionen\footnote[2]{Amerikaneren Charles W. Eliot, citeret hos James Bryce: The American Commonwealth. London 1888. III. p. 347 f.}“. Det nordamerikanske
Samfund er fra gammel Tid af gennemtrængt af Religionsfrihedens
Aand; fri, selvstændig Tænkemaade agtes idethele saa
højt, og den offentlige Mening værner i den Grad mod alle
Overgreb, at Statens humaniserende Indflydelse ikke savnes\footnote[3]{Dog hævder Moncure Conway, at ortodox Kristendom trods Alt stadigt er en national amerikansk Institution, fordi al kirkelig Ejendom er fri for Skat, --- fordi Søndagsfreden bestemmes ved Lov, --- og fordi der stadigt ansættes og lønnes nationale Kapellaner. (Autobiography. 1904. II. p. 291).}.
Her, som paa saa mange andre Punkter, afvige Forholdene i
de europæiske Lande, der i saa lange Tider have levet under
hierarkisk og aristokratisk Styrelse, i høj Grad fra de amerikanske.
Hvis Kirken ogsaa i Europa en Gang løsnes fuldstændigt
fra Staten, er det at haabe, at Religionsfriheden i den Grad vil
have gennemtrængt Folkets hele Tænkemaade, og at Folkets
Fordringer til sine religiøse Lærere maa være saa store, at de
Farer, som nu vilde medføres ved en saadan Løsning, ikke opstaa.

% --- p. 495 ---
\setcounter{footnote}{0}%

\opage{495}\grouphead{3. FILANTROPISK KULTUR.}

\chapterhead{XXXIV.}{FILANTROPIENS VÆSEN OG BETYDNING.}

1. Naar Kultur betyder Bearbejdelse og Videreførelse af
Noget, der er givet af Naturen, maa der ogsaa gives en særegen
filantropisk Kultur. Trangen til Afhjælpelse af Andres
Lidelse og til Forhøjelse af deres Glæde er ligesaa vel en
menneskelig Trang som den til at opholde Livet og udvikle Erkendelse
og Fantasi. Den filantropiske Kultur gaar, ligesom den
materielle og den ideelle Kultur, ud paa at udvikle og organisere
en væsentlig Side af den menneskelige Natur.

Den filantropiske Kulturs Opgave er at arbejde hen til, at
de Goder, den materielle og den ideelle Kultur frembringe,
blive saa Mange til Del som muligt. Baade paa den ideelle
og paa den materielle Kulturs Omraade kan der, som vi have
set, gøre sig en indsnevrende Tendens gældende, saa at Kulturgoderne
kun blive begrænsede Kredse til Del. I de filantropiske
Bestræbelser gør sig en expansiv, udadgaaende Tendens gældende,
som lægger Vægten paa den videst mulige Fordeling af
Goderne, ikke blot paa deres Produktion. Allerede i det Foregaaende,
ved de store sociale Problemer paa den materielle og
den ideelle Kulturs Omraade (se især XXV, 3; XXVIII, 5;
XXXII, 5), have vi truffet denne expansive Tendens. Det var
jo egentligt den, der først gjorde Kulturproblemerne til etiske
Problemer.

Den Trang, som fyldestgøres gennem den filantropiske
Kultur, er derfor ogsaa af mere umiddelbart og direkte etisksocial
Natur end de Fornødenheder, den materielle og den ideelle
Kultur fyldestgøre. Selve Sympatien kan jo knytte Individet
sammen med Andre, endog før det træder i Forbindelse med
dem gennem fælles Kulturarbejde. Men ogsaa indirekte kommer
den til at virke samfundsstiftende, og det er særligt denne
indirekte Indflydelse, som paa dette Sted skal omtales. Sympatien
% --- p. 496 ---
\setcounter{footnote}{0}%
\opage{496}stifter ikke blot et Samfund mellem Giver og Modtager, men
ogsaa mellem alle dem, der give i Fællesskab, mellem alle
dem, hvis Sympati fører dem til at arbejde i samme Retning
paa Afhjælpelse af den materielle og aandelige Nød i Verden.
Der bliver en særegen Klasse af Formaal og Opgaver, som
kunne samle Kræfterne og gennem fælles Virksomhed udvikle
Fællesskabet videre.

2. Ad to helt forskellige Veje kunde man dog søge at
begrunde den Paastand, at de filantropiske Bestræbelser ikke
med Rette danne nogen særegen Gruppe indenfor den sociale
Etik.

Fra den ene Side set kunde det synes, som om hele
Etiken egentligt var en Lære om filantropisk Kultur. Dersom
den uinteresserede og universelle Sympati er det psykologiske
Grundlag, som forudsættes ved den etiske Vurdering, saa synes
der ikke at kunne gives nogen Del af Etikens Indhold, som
ikke er Udtryk for denne Sympati og fyldestgør Trangen til at
lindre og hjælpe. Hvorfor da opstille et specielt Afsnit om filantropisk
Kultur? --- En Tankegang, som denne, vilde tilspidse
sig i den Sætning: „Intet er fornødent uden Menneskekærlighed!“
Den lyder smukt, men den er ikke sand. Ti netop hvor Menneskekærligheden
skal kunne virke og udrette Noget, bliver der
Brug for mange forskellige Kræfter og Evner, for at Velfærden
i Verden kan fremmes. Velfærd bestaar jo netop selv i, at
Evner og Drifter udvikles hos saa Mange som muligt og i saa
stor Fylde og Harmoni som muligt. Derfor beholde den materielle
og den ideelle Kultur, Familien og Staten deres selvstændige
Betydning. Man kan ikke udrydde alle Kræfter og
Drifter i Menneskenaturen med Undtagelse af selve Menneskekærligheden.
Man vilde derved udslette de personlige Forskelligheder
og gøre Menneskelivet fattigt. Menneskekærligheden selv
maatte i hvert Tilfælde fordre, at de skulde frembringes paany.
Den skal være en Kraft, der forædler Kulturarbejdet, men den
skal ikke træde i Stedet for dette eller gøre det til blot Middel.
Arbejde, Kunst og Videnskab have først og fremmest Værdi som
Udtryk for dybtliggende Interesser, og kun fordi de saaledes
have en selvstændig Værdi, blive de ogsaa af etisk Betydning.
% --- p. 497 ---
\setcounter{footnote}{0}%
\opage{497}Der gives Filantroper, som virke med stor Opofrelse og Energi,
men som ringeagte og forkaste alle Kulturbestræbelser, der ikke
direkte og umiddelbart „forbedre“ Menneskene og deres Kaar.
De opfatte Videnskab og Kunst som Udtryk for egoistisk Nydelse
og Forfængelighed. Om den irske Maadeholdsapostel Pater
Matthew fortælles, at han frygtede Kundskab, intellektuel og
moralsk Frihed næsten ligesaa meget som Brændevinen. Mange
Filantroper mangle ogsaa Sans for Betydningen af de uheldigere
stillede Klassers selvstændige og aktive Kamp for Tilværelsen.
Man vil gerne give Hjælp; men man forarges over, at den
Klasse, til hvilken de Nødlidende høre, organiserer sig for at
kunne sætte sine Krav igennem i Samfundet, f. Ex. ved Striker.
Den frejdige Menneskekærlighed vil derimod netop tro paa den
frie Livsudfoldelses Betydning og glæde sig over alle de forskellige
Former, i hvilke de menneskelige Kræfter bevæge sig,
selv hvor der er Ansvar og Fare forbundet med dem. ---

Paa den anden Side kunde man paastaa, at det var gavnligst,
om man slet ikke lod Sympatien ytre sig umiddelbart.
„Man forøger“, siger man da, „Slægtens og de Enkeltes Velfærd
ved Arbejde i Kulturens Tjeneste. Jo flere materielle og
ideelle Goder der frembringes, des større en Kapital har Menneskeheden
at raade over, og det er jo dog det, det kommer
an paa. Altsaa ingen direkte Indgriben, ingen umiddelbar Virksomhed
i filantropisk Retning, men energisk Virksomhed paa
den materielle og ideelle Kulturs Omraade!“ --- Denne Opfattelse
overser først og fremmest, at Frembringelsen af en Kapital
for Slægten ikke er nok; det kommer ogsaa an paa, hvorledes
denne Kapital er fordelt mellem Slægtens Medlemmer. Hvis
der gives dem, der gaa tomhændede bort ved Fordelingen, saa
kan umiddelbar Hjælp og Indgriben ikke undværes. Og det er
ikke blot Samfundsordningen, der stedse volder Nød og Ulykke
ved den ulige Fordeling af de frembragte Goder. Nød og
Ulykke opstaa ogsaa ved Naturens Gang, uafhængigt af menneskelige
Forhold og menneskelige Villier. Saa højt Kulturen
end kan tænkes at stige, saa vil den dog neppe nogensinde
kunne stoppe den Kilde til Nød og Ulykke, som ligger i Men%
% --- p. 498 ---
\setcounter{footnote}{0}%
\opage{498}neskenes Vilkaar som Naturvæsener. Det er endeligt ikke Alle,
der have Betingelserne for at arbejde med paa Kulturopgaverne.
Det er en Gunst og en Lykke at høre til Kulturens Arbejdere,
og det bliver en særlig Opgave at virke hen til, at denne Lykke
kan falde i saa Manges Lod som muligt.

3. De to nys kritiserede modstaaende Anskuelser hænge
nøje sammen med de to forskellige Sider, Kampen for Tilværelsen
frembyder.

Den førstes Blik er fængslet af de Lidelser, Tilværelseskampen
fører med sig. De Væsener, hvis Natur og Kraft ikke
svare til de Krav, Forholdene stille, gaa under efter først at
have gennemgaaet en længere eller kortere Lidelsestid. De
Svage nedtrampes i Slagets Hede, og de Mødige segne om,
før Kampen er endt. Naar Menneskekærligheden først vaagner,
vender den sig naturligt med stor Inderlighed og med Forkærlighed
mod disse Lidende, Forkuede og Nedtrykte. Den indskrænker
sig ikke til at opsøge dem og opofre sig for dem;
men den ser med Mistanke paa de Lykkelige, paa dem, der le,
og finder tilsidst Livets egentlige, sande Stemning i Sorgen og
Graaden. Den lader sig ikke nøje med at trøste dem for, hvad
de maa lide, men siger endog til dem, at de i Lidelsen have
fundet Livets egentlige Sandhed. Den hilder sig derved i en
forunderlig Modsigelse, da den jo stedse gaar ud paa at raade
Bod paa Lidelsen, og altsaa søger at bringe Mennesket ud af
den Tilstand, der efter dens Udsagn skulde være den normale.

Den anden Opfattelse dvæler hovedsageligt ved, at gennem
Kampen for Livet udvikles de nye Livsformer. Det er nu en
Gang saa, at Forholdene stille deres Krav, og at de Individer og
Racer, som ikke kunne fyldestgøre disse Krav, tidligere eller
senere maa gaa til Grunde. Selv den mest velmente Indgriben
forlænger kun Lidelserne for de til Undergang Dømte og hindrer
Slægten i at drage den Nytte af sine Erfaringer, som den
bør drage. Man afhjælper Øjeblikkets Nød paa dette enkelte
Sted, men ser ikke, at man derved ofte frembringer Ulykke og
Misforhold paa andre Punkter og forsinker eller stanser Udviklingen
for dem, der have Evne til at gaa frem. Ligesom en
Indianerflok paa sit Jagt- eller Krigstog ofte maa overlade Børn,
% --- p. 499 ---
\setcounter{footnote}{0}%
\opage{499}Svage og Gamle til deres Skæbne for ikke at hæmmes i sin
frie Bevægelse, saaledes maa Menneskeslægten under sin store
Fremmarsch ikke belemres med dem, der ikke kunne følge med.
Jo flere Kræfter der anvendes ved Ambulancen, des færre har
man til Raadighed i Kampen.

De to Betragtninger kunne meget godt forbindes med hinanden,
og ere i Virkeligheden forbundne i den moderne Udviklingslære.
Der fordres af ethvert levende Væsen, at det maa
kunne lempe sig efter Livsforholdene. Denne Tillempelse bringes
i Stand paa to Maader, dels indirekte, dels direkte. Den indirekte
Tillempelse sker ved, at de Individer og Racer, der ikke
kunne magte Forholdene, svækkes og forgaa. Den direkte Tillempelse
sker ved, at der vækkes og udvikles nye Evner, som
kunne magte Forholdene. Jo højere et Udviklingstrin et Individ
staar paa, des større en Rolle spiller denne direkte Tillempelse,
og den menneskelige Kultur er dens højeste Form. Al materiel
og ideel Kultur er en Form for den Virksomhed, ved hvilken
Mennesket gør sig til Herre over Livsforholdene. Men den
indirekte Tillempelse øver stedse sit Værk, og det er nu netop
den filantropiske Kulturs Opgave at hindre, at den antager den
brutale Form, som den har paa lavere Udviklingstrin. Intet
Steds kommer maaske det egentligt Menneskelige saa ejendommeligt
frem, som i Omsorgen for og Medlidenheden med
Livsspirer og Livsformer, der ellers vilde gaa til Grunde. Det
at mildne Tilværelseskampens Følger er en Opgave, Mennesket
stiller sig, saa snart det er kommet udover det allerraaeste
Trin; og det er en Opgave, ved hvilken det kun fortsætter,
hvad der allerede ytrer sig i de Instinkter, der lade Dyr arbejde
og opofre sig for deres Afkom. De Lidende høre med
til Slægten ligesaa vel som de Sunde og Kraftige. Blikket maa
ikke ensidigt vendes mod den ene eller den anden Side; derved
vilde vi faa en Dualisme i Slægten, der strider mod den universelle
Kærlighed. Kan det end siges, at Buddhismen og Kristendommen
have forset sig paa Lidelsen og Ulykken i Verden og
manglet Blik for den aktive Kulturvirksomheds Betydning, ---
og har der end været Medlidenhed og Menneskekærlighed 32' i
% --- p. 500 ---
\setcounter{footnote}{0}%
\opage{500}Verden før deres Tid\footnote[1]{Ortodoxe Teologer lade sig ikke nøje med at vise, at Hedningenes Dyder ere glimrende Laster, men frakende ogsaa Hedningene de Dyder, de virkeligt have. Saaledes siger Martensen: „1 Hedenskabet var Barmhjertigheden udslukt og ukendt“ (Individuel Etik. p. 304), og han udvikler, hvorledes „de mange Hospitaler og milde Stiftelser pludseligt fremstaa som noget for Hedenskabet Nyt og Ukendt, alt eftersom Kristendommen udbreder sig i Romerriget. Ti iblandt alle de rige Hedninge og iblandt alle de vise Hedninge have Ingen nogensinde tænkt paa Sligt til Hjælp for deres lidende Brødre, men kun overladt disse til deres Skæbne“. \{Soclal Etik. p. 160 f.). Dette er historisk urigtigt. At Barmhjertighedens Aand ikke var Hedningene ukendt, vil allerede enhver Læser af Odysseen vide: Fattige, Fremmede og Bønfaldende staa der under Zeus's Beskyttelse, som det fl ere Steder udtales med stor Inderlighed. I Aten var der et særligt Alter for Medlidenheden; Staten understøttede dem, der formedelst legemlig Svaghed ikke selv kunde skaffe sig deres Underhold, og velhavende Privatmænd ansaa det for deres Pligt at hjælpe de Fattige. Tukydid fortæller i sin Beskrivelse af Pesten i Athen, hvorledes Mange, uden at lade sig afskrække af Faren for Smitte eller af Sygdommens Uhygge, gik ind i Husene for at pleje de Syge. I Platons „Love“ (p. 729—730) lægges meget stor Vægt paa Pligterne mod Fremmede, Venneløse og Bønfaldende. Saadanne Bud ere fra den gammelariske Etik gaaede over i den græske (se Leist: Alt-Arisches Jus gentium). --- Heller ikke vare Hospitaler og milde Stiftelser noget aldeles Nyt og Ukendt i Hedenskabet. Navnlig opstod der i den romerske Kejsertid, væsentligt under stoisk Indflydelse, en hel Række Foranstaltninger og Stiftelser til Bedste for Slaver, Børn og Hjælpeløse. --- Smlgn. L. Schmidt: Die Ethik der alten Griechen. Il, p. 289 f. --- Schiern: Om Humaniteten i den ældre romerske Kejserlovgivning. (Historiske Studier. II).}, --- saa staar det som deres udødelige
Fortjeneste at have hævdet de Lidendes Menneskerettigheder.
De have vendt Menneskenes Blik mod Sider af Livet, fra hvilke
det har en naturlig Tilbøjelighed til at holde sig borte. De have
derved først ret lært Menneskene at kende Livet. Trods den
højt spændte Idealisme i disse Verdensreligioner have de derfor
alligevel en afgjort realistisk Karakter. Den Menneskekærlighed,
de forkyndte, fik vel ofte en passiv og en sentimental
Karakter og var ikke altid skikket til at forædle Kulturarbejdet;
men i sit egentlige Væsen var den dog en ny Form for aandelig
Kraft, Udslag af et indre Livs Fylde, der formaaede at bære
Mere end det, den Enkeltes personlige Skæbne paalagde. (Smlgn.
% --- p. 501 ---
\setcounter{footnote}{0}%
\opage{501}IX, 1). Den frygter ikke for, at Slægten skal blive ringere ved,
at de Lidende ikke overlades til deres Skæbne. Den vil fortsætte
Kampen for Livet saalænge som muligt; den kapitulerer
ikke let overfor Nøden, og den formaar at finde og benytte
Muligheder dér, hvor det ligegyldige Øje ikke bliver dem var.
Den frygter ikke for, at Udviklingen skal forsinkes gennem den
Hjælp, der ydes dem, der ikke kunne følge med. Maalet ligger
jo ikke blot i Fremtiden. Den enkelte Slægt er ikke blot Middel
eller Gennemgangsled for fremtidige Slægter, men dens Nød
og Lidelse vejer med som et selvstændigt Lod paa Vægtskaalen,
naar Livets Værdi skal bedømmes. Herpaa beror den filantropiske
Kulturs Berettigelse og Betydning.

4. Kærligheden er ofte blind, og det bliver da tilfældigt,
hvilken Gavn den stifter. Indsigt i de virkelige Forhold, først
og fremmest i Modtagerens Natur og Stilling, er en nødvendig
Betingelse for, at Kærligheden ikke skal medføre fordærvelige
Virkninger. (Smlgn. XII, 2). Kærlighed skal udvikle sig til at
blive Et med Retfærdighed. (III, 9; X, 4). Den hviler paa den
Overbevisning, at der hos den, der skal hjælpes, findes et ejendommeligt
Midtpunkt, en Kraft, der skal vækkes til selvstændig
Virksomhed. Den fastslaar derfor ingen Afstand mellem Giver
og Modtager; Gaven rækkes ikke ovenfra nedad; Modtagerens
Selvstændighed anerkendes. Derved er Menneskekærlighed forskellig
fra „Naade og Barmhjertighed“, som ere uetiske Begreber.
I Begrebet „Naade“ have vi et mærkeligt Exempel paa,
hvorledes helt modsatte Standpunkter netop ved Modsætningen
kunne komme til at bestemme hinanden. Naar man opfatter
Kærligheden som Naade, vil man stille den i den størst mulige
Modsætning til den Tjenstvillighed, der er Gengæld for ydede
Tjenester eller udspringer af, at man føler sin egen Magt som
begrænset. Naade vises af den, der har fuldkommen Magt over
Andre og staar saa højt hævet over dem, at han slet ikke kan
modtage Noget af dem. Dette Begreb er hentet fra Krigens
Verden; ti saaledes staar den sejrende Kriger overfor den afvæbnede
Fjende. Naade er Sejrherrens Afkald paa at bruge sin
Magt. Men etisk set kan der aldrig forekomme en Situation,
hvor Modtagerens selvstændige Ret skulde være aldeles forsvun%
% --- p. 502 ---
\setcounter{footnote}{0}%
\opage{502}den, og hvor Giveren staar som absolut Udgangspunkt. De
staa begge som Led af Slægten, underlagte en fælles Lov, der
giver Enhver af dem sin Ret\footnote[1]{--- Gneist: Das Selfgovernment in England. 3. Aufl. p. 260.}.

At Giveren aldrig er absolut Udgangspunkt, er let at vise.
Det kan endog siges til ham: „Hvad haver du, som du ikke
haver annammet?“ De gunstige Betingelser, under hvilke han
har begyndt Kampen for Livet og er naat saa vidt, han er kommen,
har han ikke selv frembragt. Han staar her i en Gæld
til Slægten, som han kun kan afbetale ved den Hjælp, han yder
Andre under deres Kamp for Livet. Han har maaske desuden,
bevidst eller ubevidst, været Aarsag til, at Betingelserne for
Andre ikke bleve saa gunstige, som de ellers vilde have været.
Eller de gunstige Betingelser, som ere komne ham til Gode,
hænge sammen med en Ordning af Samfundet, som ikke fuldt
tilfredsstiller Retfærdighedens Fordringer. \textsc{Kant} opkaster det
Spørgsmaal, om Evnen til at gøre vel ikke for Størstedelen er
et Resultat af Statens Uretfærdighed, der bevirker en Ulighed
i Velstand, som igen gør Velgørenhed nødvendig, --- og om
under saadanne Omstændigheder den Bistand, som den Rige yder
de Nødlidende, overhovedet fortjener Navn af Velgørenhed -). ---
Historien viser, at Fattigdommen især tager til, naar overleverede
Institutioner opløses uden at erstattes af andre, som kunne afgive
en lignende Støtte, som hine ydede. Saaledes udviklede
Fattigdommen sig i Middelalderen efter det personlige Slaveris
Ophør, senere ved Livegenskabets Ophør og ved Lavsvæsenets
Ophør, overhovedet altsaa ved Overgangen fra bundet Arbejde
til frit Arbejde\footnotemark[1]. Slaveri, Livegenskab og Lavstvang stemmede
ikke med Retfærdighedens Fordringer; men ligesaa lidt var der
fuld Retfærdighed i, at de faldt bort, uden at der dannedes nye
Former, der kunde afgive den Hjælp, som de havde ydet. (Smlgn.
XXIII, 4). Den Enkelte arver her den Skyld, Slægten har paa-

') Kun naar Begrebet Naade fortolkes som en højere Retfærdighed
(se Etiske Undersøgelser p. 63), bortfalde disse Betænkeligheder.

\footnotemark[1] Tugendlehre § 31.

 Smlgn. Lecky: History of European Morals. Il, p. 78---87.
% CHECK p. 502: note mark without a note

% --- p. 503 ---
\setcounter{footnote}{0}%
% CHECK p. 503: notes paired with the marks in order (the layer misread a number)
\opage{503}draget sig, arver den ved at fortsætte den og nyde Godt af
dens Følger. Den filantropiske Kultur staar fra denne Side set
som en Bestræbelse efter Genoprettelse af, hvad der er forsømt
og forskyldt. \textsc{Friedrich} \textsc{Nietzsche} har villet udlede den store
Elendighed og Nød, der findes i den europæiske Verden, af
„Slavemoralen“, der har holdt de Svage og Syge i Live paa de
Stærkes og Sundes Bekostning. I Modsætning til denne Paastand
er det vistnok med fuld Ret blevet hævdet, at hvis der
er sket en Forøgelse af de syge og ulykkelige Individers Tal,
saa er den for en stor Del netop at udlede af den Udbytning
af de Svage ved de Stærke, som Nietzsche i sin „Herremoral“
tager Ordet for. De „Herrer“, Nietzsche beundrer, opofre med
„god“ Samvittighed Andres Sundhed og Ære for deres egen
Nydelse eller Fordel\footnote[1]{F. Tønnies: Der Nietzsche-Kultus. Leipzig 1897. p. 109.}.

Hertil kan endnu føjes, at den individuelle Ejendomsret ---
som tidligere (XXVI, 15) paavist --- kun kan begrundes etisk,
naar Ejendommen opfattes som et Middel, der betros den Enkelte,
for at han ved det skal løse en social Opgave. Den saakaldte
Velgørenhed staar altsaa ikke som noget helt Nyt, som
en ganske særegen Dyd, men som en naturlig Følge, en Pligt.

5. Den store Betydning, som Buddhismen og Kristendommen
have haft og have i Filantropiens Historie, indskrænkes
ikke blot ved Blindheden af den Medlidenhed og Sympati, de
vække\footnote[2]{Man har endog kunnet paastaa, at Kirken ved sin ufornuftige Almissegivning har frembragt en stor Del af de Onder, den søgte at raade Bod paa. F. Raumer: Geschichte der Hohenstaufen. VI, p. 578. Lecky: History of European Morals. Il, p. 101.}, men ogsaa derved, at det for en meget stor Del var
religiøse Motiver, som laa til Grund netop for de mest storartede
Kærlighedsgerninger. Hvor Menneskekærligheden faar en
teologisk Motivering, er der et Mellemled mellem Giver og Modtager,
og disse staa ikke i et umiddelbart Forhold til hinanden.
Kærlighedsgerningen udspringer ikke umiddelbart af Medfølelsen,
men faar Karakteren af en Bodsgerning. Hvad man gav til de
Fattige, gav man ikke saa meget for deres Skyld, som for at
vinde sin egen Salighed. Almisse --- sagde man i Middelalderen
% --- p. 504 ---
\setcounter{footnote}{0}%
% CHECK p. 504: notes paired with the marks in order (the layer misread a number)
\opage{504}--- bærer hundredfoldig Frugt og udsletter Synderne, ligesom
Vand slukker lid\footnote[1]{Raumer anf. Skr. p. 575. --- Lecky anf. Skr. p. 85. 99.}. Grunden til, at den store Iver i Retning
af Godgørenhed kølnedes for en Tid ved Middelalderens Slutning,
var tildels den, at Troen paa de gode Gerningers Evne
til at skaffe Syndsforladelse faldt bort ved Reformationen. De
Fattige vare blot Midler eller Objekter, paa hvilke Bodfærdigheden
skulde øve sig: det Væsentlige blev, at man berøvede sig
selv Noget, ikke at de Andre fik Noget. Med mistænksomme
Øjne ser man derfor ofte fra dette Standpunkt paa Bestræbelser,
der gaa ud paa at gøre Almissegivning overflødig. Den
irske Politiker O'Connell, en ivrig Katolik, nærede af denne Grund
Betænkelighed ved den moderne Fattiglovgivning: Tiggeriet
ansaa han for noget i sig selv Helligt, som det var ugudeligt
at ophæve. „Den middelalderlige Opfattelse af Almissens Fortjenstlighed
maatte nødvendigvis føre med sig, at man saa paa
Tiggerne med et særligt mildt Blik. De betragtedes ingenlunde
som Verdens Skumpelskud; Tiggerlivet ansaas tværtimod som et
Gud særligt velbehageligt Liv. Naar Alt kom til Alt, saa kunde
det maaske være Tvivl underkastet, om det var Giverne, der
vare de egentligt Velgørende, og ikke snarere Tiggerne. Disse
gav jo Folk Anledning til at give Almisser og fremmede derved
Givernes evige Vel; for timelige Gaver fik man evige i
Stedet. Og man saa ikke med kritiske Øjne paa Tiggerne, om
de vare værdige Trængende eller ej; man øste ud i Flæng til
Værdige og Uværdige“\footnote[2]{M. H. Nielsen: Fattigplejen i Danmark før Reformationen. (Aarbog for dansk Kulturhistorie 1895). p. 85. F}.

Den teologiske Motivering har desuden den Mislighed, at
den naturligt kommer til at gøre Værdigheden til at modtage
Hjælp afhængig af den dogmatiske Tro. Filantropien bliver konfessionelt
begrænset. Der sættes derved en Begrænsning, som
er fremmed for Menneskekærlighedens Væsen og hindrer den i
frit at udfolde sig. Desuden fremmes Hykleriet, naar en bestemt
Konfession eller maaske endog Deltagelse i visse kirkelige
Skikke blive Betingelse eller dog begunstigende Omstændighed
for at nyde Godt af Velgerningen.

% --- p. 505 ---
\setcounter{footnote}{0}%
\opage{505}Derfor er Filantropiens Frigørelse fra de konfessionelle
Skranker af overordentligt stor etisk Betydning. Endnu synes
dog de fleste Mennesker at behøve teologiske Motiver, naar Velgørenhed
skal have nogen rigtig Art. Det er ikke blot middelalderlige
Katoliker, som forstaa at handle med Himlen, og som
Intet give bort uden fyldestgørende Anvisning paa Erstatning i
en anden Verden. Ogsaa hos Protestanter findes Spor nok af
saadan „hinsidig Verdslighed“ (other-worldliness). Hvad paa den
anden Side dem angaar, hvis Tro gennembryder de konfessionelle
Skranker, saa voxer ikke altid Evnen til Medfølelse og
Opofrelse med den kritiske Sans. Erkendelsen trives ofte paa
Følelsens Bekostning, og det kan se ud, som om teoretisk Sneverhed
var Betingelsen for praktisk Expansion. --- Menneskekærligheden,
som har overvundet saa mange Vanskeligheder, vil
dog ogsaa nok overvinde denne. Den er født i en Krybbe, men
den vil voxe, og den vil forsøge at gennemtrænge hele Verden.

\chapterhead{XXXV.}{FILANTROPIENS ORGANISATION.}

1. Allerede den Enkelte maa organisere sin Velgørenhed.
Han maa raade med de Midler, han kan og vil anvende til Hjælp
for Andre, saaledes at de paa bedste Maade komme til at opfylde
deres Hensigt. Han maa ikke lade sig lede af blot Lune,
saa at det bliver en Tilfældighed, at han giver den Ene og
nægter den Anden. Lunet vil naturligvis især raade, jo mindre
personligt et Forhold man stiller sig i til dem, der hjælpes, saaledes
overfor Tiggere. Naar menneskekærlige Mænd som James
Mill og Erkebiskop Whately roste sig af, at de aldrig havde
givet en Tigger Noget, var det et rigtigt Princip, da det blev
suppleret ved aktiv Hjælp til dem, hvis Trang man kendte, og
overfor hvem man derfor ogsaa kunde finde den bedste Maade
% --- p. 506 ---
\setcounter{footnote}{0}%
\opage{506}at hjælpe paa. Det er en let Sag at tjene en Dagløn paa 3 å
4 Kr. ved at gaa fra Dør til Dør i København; det læres af
Udsagn fra de mange Betlere, der fremstilles for Kriminalretten\footnote[1]{tive and delinquent classes. Boston 1901. p. 139.}.
Tiggeriet lønner sig ofte bedre end virkeligt Arbejde.
Men den Velgørenhed, der fremmer Sligt, er paa Afveje og
trænger til Organisation. Man har endog sagt, at den, der giver
i Blinde, Uvidenhed og Tankeløshed, bærer sig ad som den,
der affyrer en Kanon mod en Menneskemasse\footnotemark[1]. --- Den Enkeltes
Midler, Erfaring og Indsigt ere begrænsede. Hans Skærv
kan ofte gøre den bedste Virkning ved at lægges til Andres,
og der gør sig naturligt en vis Arbejdsdeling gældende, idet
nogle Individer særligt egne sig til som Undersøgere, Ledere
eller Fordelere at virke i filantropiske Foreningers Tjeneste. Der
kan her udvikle sig et frit Bureaukrati, en Art Æresposter, som
maaske en Gang ville være forbundne med ikke mindre Anseelse
end Statens Embeder. Dog vil det ikke være af det Gode,
om det personlige Forhold mellem Giver og Modtager rent faldt
bort. Det er den levende Kærne i den filantropiske Kultur, det
Frihedselement, som denne ligesaa lidt som nogen anden Art af
Kultur kan mangle. Den individuelle Velgørenhed og den kammeratlige
Velgørenhed ville altid være nødvendige Forudsætninger
for og Overgangsled til større Foreningers systematiske Virksomhed.
Den individuelle og den kammeratlige Filantropi ville
netop stræbe efter, at det systematiske Maskineri kun sættes i
Virksomhed, hvor det er aldeles nødvendigt. Det gamle Naboforhold,
som saa let taber sin Betydning under Kulturens Gang,
især i de store Byer, maa stedse staa som det Ideal, man søger
at fastholde eller genoprette under de forandrede Forhold.
\textsc{Bosanquet} har derfor med Rette sagt, at Velgørenhed bør
betyde en Tjeneste, den ene Nabo yder den anden, ikke blind
Almissegivning. Selve de filantropiske Foreninger bør virke med
saa stor Decentralisation, som kan forenes med det Overblik og

:) Københavns Understøttelsesforening 1874—1599. Ved Jul. Rüdinger.
p. 4.

-) Henderson: Introduction to the study of the dependent, 
% CHECK p. 506: note mark without a note
defec%
% --- p. 507 ---
\setcounter{footnote}{0}%
\opage{507}den Sammenslutning, som det netop er disse Foreningers Øjemed
at tilvejebringe. Bidraget til den filantropiske Forening
bliver ellers let en Art Skat, man halvt uvilligt paatager sig,
fordi det nu én Gang er Skik; og de, der modtage Hjælpen,
komme ogsaa let til at føle sig staaende overfor et Maskineri.
Kunsten vil være at forbinde den individualiserende Virksomhed
med tilstrækkelig Regelmæssighed og Ensartethed i den hele
Virksomhed. Uden Organisation af Filantropien arbejdes der i
Blinde og paa tilfældig Maade. De Fattige vænnes til at lyve
for at skaffe sig Hjælp fra flere Steder paa én Gang, og den
Hjælp, der kan komme fra det enkelte Sted, er ogsaa hyppigt
saa ringe, at den ikke udretter Andet end at gøre Modtageren
i Stand til at skaffe sig Arbejde ved at underbyde sine Kammerater.
Det viste sig for en Del Aar siden i London, at der
ikke manglede Pengemidler til Afhjælpelse af Nøden, heller
ikke i Regelen Personer og Institutioner, der vilde deltage i
Hjælpearbejdet, men at det Hele af Mangel paa Organisation,
ikke mindst af Mangel paa tilstrækkeligt Kendskab til Forholdene,
ofte gjorde mere Skade end Gavn\footnote[1]{B. Bosanquet: The principles and chief dangers of the administration of charity. (Journal of Ethics. III). --- Om filantropiske Organisationer se foruden den ovenfor anførte Beretning om Københavns Understøttelsesforening Beskrivelsen af de store engelske og amerikanske Organisationer hos Henderson: Introduction. Part. II, chap. 9: The charity organization society.}.

Den frie filantropiske Foreningsvirksomhed er endnu kun i
sin Begyndelse. Den har sikkert en stor Fremtid for sig; Opgaver
vil den i hvert Tilfælde ikke komme til at mangle.

2. Hverken den individuelle eller den organiserede Filantropi
vil dog kunne magte alle Opgaver eller yde Garanti mod
den værste Nød. Dertil kræves en fastere Organisation og større
Midler, som kun Staten raader over. I Fattigvæsenet som i
Skolevæsenet er det Statens Opgave at sørge for, at de mest
elementære Betingelser ere tilstede. Det er ikke blot, fordi
Staten derved indirekte sørger for Sikkerheden. Ligesaa lidt
som den understøtter den ideelle Kultur for Sikkerhedens Skyld,
ligesaa lidt er det af blot klog Beregning, at den organiserer
% --- p. 508 ---
\setcounter{footnote}{0}%
\opage{508}en Hjælp til de Fattige. Den træder til dér, hvor den yderste
Nød er paa Færde, selvom den offentlige Sikkerhed ikke trues.
Og ligesaa lidt som Enkeltmands Hjælp skal den Hjælp, Staten
yder, være en Naadessag, men betragtes som en etisk-social
Retfærdighedshandling. Ved at indrette et Fattigvæsen indrømmer
den, at den har Noget at oprette. I en ideal Stat vilde der
ingen Fattige være; at der gives dem, der mangle det til Livets
Ophold Nødvendige, skønt de have Evne og Villie til at arbejde,
vil stedse hænge sammen med Fejl i den sociale Organisation.
Deraf følger dog ikke, som \textsc{Fichte} mente, at Fattigunderstøttelsen
først og fremmest er Statens Sag og først i anden Række
Enkeltmænds Sag. Statens Indgriben maa her, som paa andre
Punkter, kun være den sidste Udvej, det nødvendige Supplement
til de frie Kræfters Virksomhed. Den frie Filantropi faar
netop den Opgave: saalænge som muligt at forebygge, at Staten
maa gribe ind med sit tunge Apparat.

Fra to helt forskellige Sider er Statens filantropiske Virksomhed
bleven angreben. --- Man har i Statsfattigvæsenet fundet
en slem Kommunisme, hvor Staten yder sin Gave uden at
have tilstrækkelig Indflydelse paa Modtageren. Naar det ikke
kan forbydes dem, der modtage Fattigunderstøttelse, at indgaa
Ægteskab, bliver Fattigdommen jo forøget ved Statshjælpen. Man
har derfor i Stedet for Fattigvæsenet forlangt et Statsassurancevæsen,
i hvilket Arbejderne skulde være tvungne til at deltage\footnote[1]{Saaledes de tyske Nationaløkonomer Wagner og Schåffle. (Schönbergs Handbuch der polit. Oekonomie. 1. Ausg. Il, p. 598. 615). --- Smlgn. ovenfor XXVI, 14. f}.
Men det er ikke let at forstaa, hvorledes man kan faa Folk til
at assurere deres Fremtid, naar de Intet have til Underhold i
Nutiden. Dette Forslag forudsætter altsaa, at den værste Nød allerede
paa en eller anden Maade er afhjulpen. Men hvorledes dette
Resultat skal kunne naas, er netop Vanskeligheden. --- Fra en
helt anden Side vil man have Staten holdt aldeles borte fra
dette Omraade. Ligesaa stor og smuk en Ting den private Velgørenhed
er, ligesaa fordærvelig en Ting anser man Statsvelgørenhed
for at være. Kun den frie Menneskekærlighed --- mener
% --- p. 509 ---
\setcounter{footnote}{0}%
\opage{509}man --- lader Velgørenheden blive udøvet med den rette Ansvarsfølelse,
saa at man ikke ved at understøtte Uværdige fører
Slægten til Tilbagegang istedetfor til Fremgang. Et Bureaukrati
kan ikke have den rette Ansvarsfølelse; det kommer ikke i det
personlige Forhold til de Lidende. Desuden er det farligt at
forøge Statens Magt; man skaber derved let Farer for den individuelle
Frihed. En velgørende Despot bliver dog ved at være
en Despot. Staten har blot den Opgave at sørge for Ret og
Sikkerhed; i Staten kan ikke, som i Familien, Sympatien umiddelbart
raade\footnote[1]{Smlgn. Herbert Spencer: The man versus the state (London 1884), hvor denne Tankegang er gennemført i konsekvent Form. I andet Bind af Principles of Ethics (1893) (Chap. 7: Relief of the poor), hvor Vanskelighederne og Betænkelighederne ved al Filantropi skarpt fremhæves, mener Spencer dog, at Statsunderstøttelse af Fattige har en retfærdig Basis, da den Ret, forrige Tiders Livegne havde til at bruge Jorden, berøvedes dem i Tidens Løb af Godsejerne: Fattigloven er paa en Maade en Anerkendelse af den urgamle Ret til Jorden. Men Statshjælp er og bliver for Spencer „en Art social Opiumsæden“, som lindrer i Øjeblikket, men forværrer senere. Der behøves en Afvænningsperiode for at komme over til \textasciitilde{},Selvhjælps og personlig Godgørenheds sunde Tilstand“.}. Hertil er at svare, at der næppe kan drages
saa stærke Grænser mellem Velgørenhed og Retfærdighed, som
den anførte Tankegang forudsætter. Den sande Velgørenhed er
selv en Retfærdighedsakt, og det er derfor rigtigt, at den danske
Grundlov (§ 84) anerkender, at den Trængende har Ret til
offentlig Understøttelse. Vel er Omsorg for Sikkerheden Statens
mest elementære Opgave; men det udelukker ikke, at Staten
kan træde hjælpende til med sin Organisation og sine Midler,
hvor vigtige humane Formaal ikke kunne naas blot ved de frie
Kræfters Virksomhed. Dersom vi paa egen Haand kunde sørge
for vor Sikkerhed, vilde vi helst undvære Statsmagten; Grunden
til, at vi kalde Staten til Hjælp, er paa alle Omraader væsentligt
den samme: de frie Kræfters Utilstrækkelighed, en Utilstrækkelighed,
der ikke mindst skyldes ufuldkommen Organisation.
Faren er ganske vist ogsaa overalt den samme: nemlig den,
at det System og den Magt, der stilles paa Benene, skal brede
sig for meget paa de egentlige Formaals Bekostning. Denne
% --- p. 510 ---
\setcounter{footnote}{0}%
\opage{510}Fare existerer dog ogsaa, skønt naturligvis i mindre Grad, ved
den organiserede frie Filantropi; den indtræder overalt, hvor
det umiddelbare Forhold mellem Giver og Modtager giver Plads
for et middelbart.

3. Hvor Staten virker, staar Tvang altid i Baggrunden, og
en systematisk Fremgangsmaade anvendes, som ikke altid lader
de individuelle, faktiske Forhold komme til deres Ret. I stor
Maalestok viste det sig i England ved Virkningen af den fra
Dronning Elisabeths Tid stammende ældre engelske Fattiglovgivning.
Dels ved Administrationens Magelighed, dels ved den
for stærke Centralisation og Tilbøjeligheden til at behandle alle
Tilfælde efter et abstrakt Skema, dels ved manglende Kontrol
fra Skatteydernes Side opstod der de fordærveligste Forhold. Da
Understøttelserne bleve givne uden nærmereUndersøgelse, vænnedes
de Fattige i Befolkningen til Ubetænksomhed og Ligegyldighed.
De tænkte ikke paa den Dag i Morgen; de selv og deres
Børn vilde jo i Nødstilfælde blive underholdte af det Offentlige,
maaske endog faa det bedre end Mange af dem, der arbejdede
haardt for at kunne udrede deres Andel af Fattigskatten. Ikke
blot Ligegyldighed og Dovenskab, ogsaa Utaknemlighed og
Egoisme avledes. Børn vægrede sig ved at hjælpe deres gamle
Forældre og overlod det til det Offentlige. Demoralisationen
blev saa stor, at Lordkansleren Brougham endog i det engelske
Overhus kunde udtale, at de engelske Fattiglove vare den vigtigste
Aarsag til Befolkningens sædelige Fordærvelse og Forbrydelsernes
Tiltagen. Det blev dernæst (fra Slutningen af det
18. Aarhundrede) almindeligt, at der af Kommunerne ydedes et
Tilskud til Arbejdslønnen, hvor denne var utilstrækkelig. Herved
forvandledes Fattigunderstøttelsen egentligt til en Understøttelse
af Arbejdsgiverne, som nu kunde presse Lønnen helt ned under
det, som var nødvendigt til Arbejdernes Underhold!\footnote[1]{Das Armenwesen. Frei nach Duchatel und Nav ille. Von einem deutschen Staatsbeamten. Weimar 1837. p. 502 ff. --- Gneist: Das Selfgovernment in England. 3. Aufl. p. 700 f.} Den
nyere Lovgivning (fra 1834 af) har søgt at raade Bod herpaa.
Men det har stadigt vist sig at være vanskeligt at afgøre, hvor
% --- p. 511 ---
\setcounter{footnote}{0}%
\opage{511}Trangen virkeligt er tilstede og hvor ikke. Man indførte i England
den saakaldte „Arbejdshusprøve“ (workhouse test), idet man
mente, at arbejdsdygtige Mennesker, der ikke vilde finde sig i
det Arbejde, det Offentlige paalagde dem som Betingelse for
Understøttelse, derved godtgjorde, at de ikke vare virkeligt trængende.
Men et saa mekanisk Middel blev kun nødvendigt, fordi
der ikke anvendtes tilstrækkelig Decentralisation\footnote[1]{Smlgn. Gneist anf. Skr. p. 740: „Naar man ved Vurderingen af Nøden og Fastsættelsen af de rette Midler til at afhjælpe den mangler den Forstaaelse, som en Nabo kan have, og som forudsætter personligt og stedligt Bekendtskab, som det kun kan findes i levende Kommunalforbindelser, saa kommer man til den skematiske workhouse test paa Humanitetens Bekostning og til sædelig Fornedrelse for de arbejdende Klasser. Det faktiske Spørgsmaal om Nødens Grund kan kun afgøres i en levende Kommunalforbindelse ved Kommunens Medlemmer; dersom noget Saadant mangler, saa kommer den forsørgende Statsmagt til den mekaniske workhouse test af samme Grund som den straffende Statsmagt en Gang kom til Pinebænken“. Smlgn. hermed Bosanquets ovenfor (1) omtalte Sætning: Charity for us means neighbourly service.}. Ide mindre
Kredse (Kommuner og endnu snevrere Naboforhold) er der
Mulighed for personligt Kendskab til alle de enkelte Individer, og
Staten bør derfor saa meget som muligt overlade Understøttelsens
Fordeling til saadanne Kredse. I denne Retning gaa ogsaa
de nyere Ordninger af denne Sag. Her er endnu mere end
paa andre Omraader Selvstyrelsen den eneste Ordning, der kan
føre til Maalet. Hvor de mindre Kredse have Selvstyrelse, kan
hvert enkelt Tilfælde blive behandlet efter dets individuelle
Beskaffenhed. I „Elberfeldersystemet“ -) er man gaaet saa vidt i
Decentralisation og Individualisering, at enhver Medhjælper kun
har med fire Familier at gøre. Deres Forhold kan han da
grundigt lære at kende, og det bliver tillige muligt for enhver
% CHECK p. 511: note whose mark was not found
\footnote[2]{Lon ing: Armenflege und Armenpolizei. (Schönbergs Handbuch. 1. Ausg. II), p. 588. 601 f. --- Ifølge Henderson (Introduction. p. 164) er Elberfeldersystemet tildels opstaaet med den af Chalmers i Skotland grundlagte Fattigpleje som Forbillede. „The friendly visitor“ er i Chalmers' System den gode Nabo, der bringer Hjælp og Raad i Nøden og sætter Livet i Gang, hvor det truer med at gaa istaa.}
% --- p. 512 ---
\setcounter{footnote}{0}%
\opage{512}Borger at overtage det kommunale Æreshverv som Medhjælper,
da hans Tid ikke tages saa stærkt i Beslag derved. Ved en
saadan Ordning nærmer Statsfilantropien sig mere og mere de
private filantropiske Organisationer. Den virker bedst, hvor den
samvirker med disse. Gennem en Række af Overgangsformer
afkortes den store Afstand mellem det personlige Forhold og
den upersonlige Statsindgriben.

% --- p. 513 ---
\setcounter{footnote}{0}%

\opage{513}\grouphead{C. STATEN.}

\chapterhead{XXXVI.}{FOLK OG STAT.}

1. Naar man vil stille Staten i skarp Modsætning til
Familien og det frie Kultursamfund, er det den tvingende Magt,
som tydeligst karakteriserer den. Men det er ogsaa kun, hvor
hin Modsætning skal fremstilles i sin yderste Skarphed, at' dette
Moment træder saa stærkt frem. Magten staar stedse i Baggrunden;
men hvor den er det eneste sammenholdende Baand,
er Staten sin Opløsning nær. Statens Væsen træder først frem
for os i sit hele Omfang, naar vi se, at den formaar at vække
Følelser og Motiver, der endog kunne besejre dem, som Familieforhold
og Kulturformaal vække: som naar den Enkelte
ofrer ikke blot sit eget Liv, men ogsaa Hensynet til sin Families
Tarv i Statens Tjeneste, og naar Kunstneren og Videnskabsmanden
underordne deres Fags Interesser under Fædrelandets
Tarv. Saadanne Følelser kan Staten kun vække, fordi
dens Bestaaen og dens Virksomhed ere nødvendige Betingelser
for Opnaaelsen af store almenmenneskelige Formaal. Dette har
allerede vist sig i det Foregaaende, idet baade Undersøgelsen
af Familien og af det frie Kultursamfund førte til at paapege
Opgaver, som kun Staten kunde løse. Forsaavidt have vi allerede
i det Foregaaende forudsat et bestemt Begreb om Staten,
som det nu gælder nærmere at udvikle saaledes, at baade Mod%
% --- p. 514 ---
\setcounter{footnote}{0}%
\opage{514}sætningen mellem Staten og de andre Samfundsformer og den
nøje Sammenhæng mellem dem kan finde sin Begrundelse. En
Undersøgelse af Forholdet mellem Folk og Stat vil tjene dette
Øjemed.

2. Ved Omtalen af Familien og det frie Kultursamfund
som selvstændige sociale Livsformer have vi set bort fra, at de
i Virkeligheden stedse forekomme i en bestemt historisk og
national Sammenhæng. Paa de mest primitive Udviklingstrin er
der ingen skarp Grænse at drage mellem Familie, Kultursamfund
og Stat. I Flokken eller Stammen, saaledes som den fremtræder
paa tidlige Udviklingstrin, hersker Blodets Baand, selvom
det ikke er det Eneste, der forbinder. Den primitive Flok er
tillige et Kultursamfund, forsaavidt der paa dette Trin idethele
er Tale om Kultur. Arbejdet øves i Fællesskab, og hvad man
besidder, ejes af Flokken i Fællesskab. Den ideelle Kultur er
repræsenteret af Troen paa Stammens Guder. Og Flokken er
endeligt en lille Stat, idet den samler sig til fælles Modstand
mod ydre Fjender, og idet Familiefaderen eller Høvdingen ogsaa
indadtil udøver Magt for at holde de Genstridige i Tømme. ---
Det gaar her som overalt: de skarpe og tydelige Forskelle fremkomme
først i Løbet af Udviklingen og gøre sig ikke gældende
fra først af.

En Flok er ikke et Folk, men bliver først dertil, naar der
udvikler sig en mere bevidst Følelse af Fællesskab, end der kan
være tilstede, hvor man ligesaa hurtigt splittes ad i smaa Grupper,
som man i Nødens Stund samles til gensidig Bistand. En
sammenhængende Erindring og Tradition er en Forudsætning
for et Folks Tilbliven. Gennem fælles Skæbne og fælles Virksomhed
opstaar enhver Fællesfølelse (smlgn. XIII, 3 o. fl. St.),
ogsaa Nationalfølelsen. Den fælles Skæbne og den fælles Virksomhed
betinges først og fremmest ved, at Alle bo og bygge i
det samme Land. Dettes Beskaffenhed betinger Arbejdets fælles
Natur saavel som de fælles fysiske Tilskikkelser, fælles Stemninger
overfor Naturen, fælles Frygt og fælles Glæde\footnote[1]{Om Betingelserne for Dannelse af et Folk se Spencer: Principles of Sociology. II. p. 265---287 (Political integration).}. Seiv et
% --- p. 515 ---
\setcounter{footnote}{0}%
\opage{515}Nomadefolk har dog sine Pladser, til hvilke det fra Tid til anden
vender tilbage, saa at der bliver Sammenhæng i dets Erindringer.
Har det ingen anden fast Besiddelse, saa har det dog
Fædrenes Grave, til hvilke det føler sig knyttet, og om hvilke det
kredser som om et Centrum\footnote[1]{der ellers lever som Nomade, sig et Stykke Jord for deri at begrave Sara.}. Det er ligegyldigt, hvad der fra
først af har ført Flokken sammen, --- om den er opstaaet ved
Nedstamning fra samme Familie, eller ved Sammenslutning af
flere omvankende Individer, eller ved at en stærkere Flok underlagde
sig en svagere. Hovedsagen er, at der kommer et saadant
Samliv i Gang, at et fælles Tanke- og Følelsesliv udvikler
sig under fælles Skæbne og fælles Virksomhed. Det fælles Sprog
er selv ofte først et Produkt af dette Samliv, og det er først
paa højere Udviklingstrin, at det kan komme til at spille en
afgørende Rolle. Sproget er ikke altid det Afgørende. Det
schweiziske Folk er et Exempel paa, at Nationalfølelsens og
Sprogets Grænser ikke behøve at falde sammen; Hovedsagen er,
at der udvikler sig en Fællesbevidsthed, og dertil behøves blot
en fælles Historie.

Gennem fælles Skæbne og fælles Virksomhed udvikle sig
de for en Menneskegruppe fælles Vaner og Sæder. Livet former
sig uvilkaarligt paa en vis bestemt Maade og i visse bestemte
Retninger. Disse uvilkaarligt dannede Former virke saa
igen bestemmende for Fremtiden, og det ikke blot paa ubevidst
Maade, men ogsaa gennem mere eller mindre bevidst Trang og
Følelse. Der føles en Trang til at foretage samme Handling
paa samme Vis under samme Forhold. Dyrkning af Jorden,
Jagt og Krig, Klædning og Vaaben, Ægteskabs Indgaaelse, Opdragelse,
Gudsdyrkelse, Hævn og Nydelser --- Alt ordnes og
øves ikke saa meget efter, hvad man ved Overvejelse finder
hensigtsmæssigt og fornuftigt, som efter den Maade, hvorpaa
Fædrene have ordnet og øvet det. Et Brud derpaa eller en

Smlgn. Herodot IV, 127, hvor Skyterne vægre sig ved at indlade
sig i Kamp med Perserne, da de ikke have Byer eller dyrket Land
at forsvare; kun da ville de kæmpe, hvis Perserne finde deres Fædres
Grave og ødelægge dem. --- I første Mosebog Kp. 23 køber Abraham,
% --- p. 516 ---
\setcounter{footnote}{0}%
\opage{516}Hindring deri kan vække en lignende Lidenskab, som et Dyr
føler, naar det hindres i at foretage en instinktmæssig Handling.
Der føles Forargelse over den, der vil indføre ny Sæd, og
en med Forargelse blandet Forundring er den Følelse, hvormed
man iagttager andre Menneskegruppers Sæder og Skikke. ---
Det er --- saaledes kunne vi ogsaa udtrykke os --- den fælles
positive Moralitet (se I, 2), som i Forbindelse med de fælles
Erindringer gør en Menneskegruppe til et Folk.

3. Til fuld og klar Bevidsthed kommer Nationalfølelsen
først ved Sammenstødet med og Forargelsen over det Fremmede.
Det er en psykologisk Lov, at der maa en mere eller mindre
stærk Kontrast til, for at en Bevidsthedstilstand skal faa sin bestemte,
udprægede Karakter. Historien viser da ogsaa, at Nationalfølelsen
kun har kunnet virke forenende ved at virke sondrende:
den har forenet ved den Bevidsthed om fælles Ejendommelighed,
som især vaktes ved fælles Følelse af Forargelse,
Foragt eller Fjendskab overfor det Fremmede.

Det er fra først af Mere end en æstetisk Antipati, der virker.
Alt Fremmed og Ukendt kan paa primitive Trin vække
Frygt eller Had, fordi det kan bringe Fare for Liv, Besiddelse
og Frihed. Man sætter sig i hvert Tilfælde for en Sikkerheds
Skyld i Forsvarstilstand. Men selvom Frygten bortfalder, vedligeholder
Antipatien sig. Fra den nationale Selvfølelses Standpunkt
fælder man en skarp Dom over det Fremmede. Det er
to etiske Systemer, som her i Kød og Blod staa overfor hinanden,
og en Mægling eller Forsoning kan kun ske gennem en
historisk Udvikling, der medfører Fællesskab i Skæbne og Arbejde.
Der hører længere Tids roligt Samliv til for at opdage,
at det Fremmede har sin ejendommelige Værdi. Det er en Udvidelse
af Sympatien og en Udvidelse af Forestillingerne, som
maa gaa for sig.

Sin hele Brynde kan Nationalfølelsen kun have, hvor den fremtræder
med Instinktets Styrke og Blindhed. Den har været en mægtig
historisk Kraft. Den har gjort Historien til en stor Krigshistorie,
men netop ved at volde Krigen har den virket sammenholdende og
koncentrerende. Krigen medfører de mest indgribende fælles Erfaringer
for hele Gruppen. Fælles Fare og Frygt, fælles Spænding og
% --- p. 517 ---
\setcounter{footnote}{0}%
\opage{517}Haab, fælles Nederlag og fælles Sejre, alt dette binder sammen
med det stærkeste Kit. Og dertil kommer, at paa den instinktive
Nationalfølelses Trin staar Alt paa ét Kort: det gælder i Kampen
alle Livets Goder paa én Gang; besejres man, maa man være beredt
paa at miste Frihed, Ejendom, maaske Livet selv; ja endog
Guderne mister man, idet de Sejrendes Guder fortrænge dem.
Det gælder at være eller ikke være. I \textsc{Aischylos}' Tragedie
Perserne raabes der fra græsk Side, da Salaminerslaget begynder:
„Hellenersønner, befrier Eders Fædreland, Eders Børn og
Hustruer, de fædrene Guders Boliger og Fædrenes Grave! Kampen
gælder nu Alt!“ --- Anderledes naar Udryddelseskrigenes
Tid er forbi. Selve Krigen føres med større Humanitet, og det
er ikke længere alle Livets Goder, der staa paa Spil. Der gives
et almenmenneskeligt Samfundsliv, som rækker ud over de nationale
Forskelligheder, og Modsætningen mellem vor egen Nationalitet
og den fremmede falder ikke længere sammen med
Modsætningen mellem Godt og Ondt. Vil da ikke Nationalfølelsen
engang have udspillet sin Rolle som mægtig historisk
Drivkraft, og vil ikke den etiske Betydning, den har haft, være
forbi? Den synes at maatte mangle tilstrækkelig Næring og
derfor heller ikke at ville kunne bevare tilstrækkelig Energi.

4. Nationalfølelsen har det fælles med enhver Følelse, der
nærmer sig Instinktet, at den ikke fremkaldes ved den Værdi,
som iog for sig kan tillægges det, hvortil den er knyttet. Vi
elske ikke vort Land og vort Folk, fordi vi anse dem for de
bedste, smukkeste, rigeste og bedst udstyrede; men vi elske
dem først og fremmest, fordi det er vort Land og vort Folk.
Det er en rent subjektiv, ikke en objektiv Vurdering, der anstilles
i Nationalfølelsen. Paa naiv Maade giver dette sig Udtryk
ide Superlativer, hvormed vi udstyre det, vi elske; disse
Superlativer ere Udtryk for Følelsens Inderlighed og Styrke,
men ikke for en objektiv Undersøgelses Resultat. Vi kunne
endog meget godt indrømme, at andre Lande og Folk besidde
en større objektiv Fortræffelighed; den Mangel have de dog
altid, at de ikke ere vort Land og vort Folk: „skønnere de er
maaske, ak, men det er ikke de“. Det gaar her som med Selvopholdelsen:
vi opholde Livet, først og fremmest fordi det er
% --- p. 518 ---
\setcounter{footnote}{0}%
\opage{518}vort Liv, ikke fordi vi have undersøgt det og fundet det at være
et udmærket Liv. Paa denne dens halvt ubevidste, instinktmæssige
Karakter beror det, at Nationalfølelsen kan bevares,
selvom Horizonten udvides og Antipatien overfor det Fremmede
falder bort. Med den skarpe Kontrast til det Fremmede falder
ikke nødvendigvis den inderlige Følelse ved det Hjemlige bort.
Og kun i det Hjemlige have vi et fast Tilhold og et Grundlag,
udfra hvilket vi kunne orientere os i Verden:).

Den fremskridende Kultur forøger vel Enheden og Fællesskabet
i Menneskeslægten, men den aabner ogsaa Blikket for
de ejendommelige Nuancer og Forskelligheder. Hvert Folk tager
Livet paa sin Maade, og derved bliver Livet rigere og mangfoldigere
for den, der ser paa det med sympatisk Forstaaelse.
Barbaren tramper det Fremmede ned, fordi han ikke forstaar
det, og er der Fare for de nationale Ejendommeligheder, saa
er det snarere, fordi der er for meget Barbari i Verden, end
fordi der er for megen Kultur. Den ene Nation har endnu ikke
ret indset Nødvendigheden af, at den anden skal existere. Den
fremskridende Kultur aabner baade Blikket for de forskellige
Folks Ejendommelighed og giver det enkelte Folk større Lejlighed
til at yde sit Bidrag til den fælles Udvikling.

Ligesom Familien og det frie Kultursamfund ere naturlige
Mellemled mellem Individet og Staten, saaledes ere Nationaliteterne
naturlige Mellemled mellem de enkelte Individer og hele
Slægten. Det Almenmenneskelige hindres derved ikke i sin Udvikling;
tværtimod bliver det derved rigere og mangfoldigere.

5. Staten er det organiserede Folk. Staten forudsætter en
saadan Organisation, at Folket kan optræde som Enhed udadtil,
og at ogsaa dets indre Anliggenders Ledelse og Ordning frembyder
en vis Enhed. Man kan ikke paapege noget bestemt Punkt,
hvor Staten bliver til. Allerede Familien og Flokken kunne
have en vis Organisation, saa de baade gøre Modstand mod ydre
% CHECK p. 518: note whose mark was not found
\footnote[1]{I min Afhandling Nationalitet og Kultur (Tilskueren 1895) har jeg beskrevet Nationalfølelsen som Genkendelsesfølelse og som Fællesskabsfølelse og tillige betonet, at ved ethvert dyberegaaende Arbejde maa de nationalt bestemte Elementer i Personligheden virke med.}
% --- p. 519 ---
\setcounter{footnote}{0}%
\opage{519}Angreb og ordne Forholdet mellem deres Medlemmer. En stor
Familie er allerede en lille Stat.

Ved den fastere Organisation, som Staten forudsætter, vinder
Folket den samme Fordel, som Dyreorganismen vinder ved
at have et Nervesystem. I Planten føres der et yppigt organisk
Liv; Stofskifte, Væxt og Forplantning gaa for sig med stor Energi.
Men kun, hvor et Nervesystem forbinder Organismens Dele, kan
denne optræde aktivt og med samlet Kraft overfor Omverdenen,
rette sin Virksomhed mod bestemte Punkter og bestemmes
ved Hensyn til, hvad der ligger langt udover de umiddelbare
Omgivelser. Uden Statsorganisation fører Folket et Planteliv.
Der kan udvikle sig en vis Kultur; materiel Virksomhed, Tanke- og
Følelsesliv kunne til en vis Grad udfolde sig; men det Hele
har Tilfældighedens og Spredthedens Præg. En videre Udvikling
er betinget ved Koncentration og fast Sammenhæng af Alt,
hvad der hører til Folket. I Kampen for Tilværelsen bliver en
saadan Koncentration af den største Betydning. Et Rige, som
er splidagtigt med sig selv, kan ikke bestaa. Krigen har derfor
til alle Tider været en Aarsag til fastere Organisation. Dog er
det kun paa de mest elementære Trin, at Koncentrationen overfor
ydre Fare er det herskende Hensyn. Den Organisation, som
først tilvejebringes for at møde ydre Modstand, anvendes naturligt
til at hjælpe og beskytte alle de Interesser og Formaal, som
Folkets Liv omfatter. Der gives to Typer af Statsliv, en, som
overvejende er bestemt ved Nødvendigheden af Magtudfoldelse
udadtil, og en anden, som væsentligt bestemmes ved Hensynet
til det indre Liv i Folket selv '). Saalænge en Stat maa kæmpe
for Tilværelsen overfor andre Stater, vil den første Type være
overvejende; saa snart fredeligere Forhold indtræde, vil den
anden Type udvikle sig. Det vil være let at se, hvilken af de
to Typer der vil være heldigst for en sund Udvikling af materiel
og ideel Kultur.

Allerede paa de primitive Udviklingstrin er der Noget, som
er stærkere end den blotte Magt og tager den i sin Tjeneste.
Det er de uskrevne Love, som indeholdes i Folkets Sædvaner.
% CHECK p. 519: note whose mark was not found
\footnote[1]{Spencer: Principles of Ethics. Il. p. 181---187.}
% --- p. 520 ---
\setcounter{footnote}{0}%
\opage{520}(Smlgn. 2). Ordningen af Folkets Forhold kan aldrig ske blot
gennem Magtbud. Selve Magthaverne ere i deres Indre mere
eller mindre betagne af Frygt eller Ærefrygt overfor det Overleverede.
Ingen Magthaver staar saaledes over sin Tid, at han
ikke har sine Forbilleder og Præcedenser, han følger. I saadanne
Forbilleder og Præcedenser have vi den primitive Form
for Retten.

Vi have saaledes forskellige Momenter i Statens Begreb.
Magten er Statens sidste Middel og Fundament; Retten er den
Norm, som bestemmer Magtens Anvendelse; og Kulturen er det
Formaal, som den af Retten ledede Magtanvendelse skal tjene] )-
Og ved alle disse Momenter bliver endnu Spørgsmaalet om Statens
etiske Betydning at opkaste.

\chapterhead{XXXVII.}{RET OG MORAL.}

1. Hvor Retten kun bestaar ide Sædvaner og Overleveringer,
der uvilkaarligt beherske Livet, gør Modsætningen mellem
Moral og Ret sig ikke gældende. Ved Ret forstaa vi paa vort
moderne Standpunkt Indbegrebet af de i bestemte Kundgørelser
udtalte Regler for Magtens Anvendelse, medens vi ved Moral
forstaa det, der kundgøres for os af Samvittigheden i vort Indre.
Men en saadan skarp Modsætning mellem et Ydre og et Indre,
mellem ydre Magt og indre Følelse er først Frugt af en lang
Udviklingsgang. Der er oprindeligt ingen Forskel mellem 
% CHECK p. 520: note whose mark was not found
\footnote[1]{Syndikalismen tager i sin Kritik af Staten kun Hensyn til Magtmomentet, ikke til Ret og Kultur. Og i Magten ser den kun en herskende Minoritets Middel. Selv Bakunins paa Kommunerne byggede Ordning afvises, fordi Kommunen er en politisk Institution. Den indlader sig dog ikke paa Spørgsmaalet om, hvorledes Staten vil falde sammen, og hvorledes den nye Ordning kan organiseres. (Poul Louis: Le Syndicalisme contre I'État). (1913).}Sæd%
% --- p. 521 ---
\setcounter{footnote}{0}%
\opage{521}vane, Tradition, Ret, Moral og Religion. Der gør sig en Trang
gældende til at behandle lignende Tilfælde paa lignende Maade.
At Noget én Gang er gaaet for sig paa en vis Maade, er Grund
nok til, at det næste Gang skal gaa for sig paa samme Maade.
\textsc{Tocqueville} har endog ment at iagttage, at der hos Pariserne
har udviklet sig en vis Sæd og Skik, efter hvilken de opføre
sig under og efter de Revolutioner, de gøre. ') Selv her har
altsaa Gentagelse og Vane virket! --- Som allerede bemærket.
er det ikke en rent ubevidst Vane, her raader. I Trangen til
at følge de én Gang banede Veje virker ogsaa Ærefrygt for
Fortiden og for de dunkle Magter, fra hvilke de raadende
Skikke og Overleveringer\footnote[1]{\textsuperscript{1} Souvenirs (Paris 1893) p. 105: Il sest formée parmi nous une espece de moralité particuliére ou désordre, et un code spécial pour les jours d'émeute. D'après ces lois exceptionnelles, le meurtre est toléré, la devastation est permise, mais le vol est sévérement défendu}') udledes. Disses første Udspring
skjuler sig i det Ubekendte; derved faa de en hemmelighedsfuld
Karakter. Og da de indeholde Summen af de Erfaringer,
Forfædrene have gjort i Livets forskellige Forhold, afgive de
ofte en bedre Rettesnor end Individets egen mere begrænsede
Indsigt.

Hvor disse uskrevne Love optegnes og bekendtgøres, træder
Statsmagten bestemtere frem som den, der hævder og gennemfører
Lovene, og disse faa nu Karakteren af udtrykkelige Villiestilkendegivelser,
af almindelige Regler for Handlen. Retten er
nu ikke længere Indbegrebet af Folkets uvilkaarlige Normer,
men er den herskende Magts Villie. I Modsætning til Statsmagten
staa de enkelte Individer med deres personlige Overbevisning,
og her kan det da komme til en Konflikt mellem Retten
og Moralen. De ere ikke længere forenede i Overleveringen,
% CHECK p. 521: note whose mark was not found
\footnote[2]{--- Vi mangle paa Dansk et Ord, som ganske udtrykker, hvad der ligger i det tyske «Sitte“. „Sitte“ er forskellig fra Vane ved, at den fremtræder som normerende for Bevidstheden, medens ved Vanen ingen tydelig Bevidsthed behøver at gøre sig gældende. Der er ved „Sitte“ et indre, subjektivt Moment, som kan mangle ved Vanen, der nærmer sig den rene Reflexbevægelse. --- Smlgn. den fortrinlige Udvikling af Begrebet „Sitte“ hos Jhering: Der Zweck im Recht. II, p. 19—41. Det defineres her som „verpflichtende Gewohnheit“.}
% --- p. 522 ---
\setcounter{footnote}{0}%
\opage{522}men hvad Retten byder eller tillader, forbyder maaske Moralen,
og hvad denne byder, forbyder maaske Retten. Der kan her
opstaa en Konflikt, som ofte faar en tragisk Udgang. Samfundet
maa hævde sin Retsorden, og Individet maa hævde sin Overbevisning,
maa „adlyde Gud mere end Menneskene“\footnote[1]{maa betragte Sagen paa denne Maade. Lovens Affatter synes at have villet tilsløre Muligheden af en virkelig Konflikt. Men denne Mulighed er stedse til Stede, saasnart der overhovedet kan gøres Forskel mellem Ret og Moral.}. Den
Enkelte har her ikke Andet at gøre end at appellere fra den
nuværende Retsorden til Fremtidens paa bedre Grundlag byggede
Retsorden. --- Men saadanne Konfliktstilfælde maa ikke
være bestemmende for Fastsættelsen af Forholdet mellem Ret
og Moral idethele. Ret og Moral have historisk udviklet sig af
en fælles Rod, og de blive kun i Undtagelsestilfælde helt fremmede
for hinanden.

2. Naar det primitive Udviklingstrin, hvor Forskellen mellem
Ret og Moral endnu ikke træder frem, forlades, og naar
Retten udvikles som et særligt Omraade for sig, saa er det især
fire Punkter, ved hvilke den viser sig at være forskellig fra
Moralen: den støtter sig til ydre Tvang; den angaar kun den
ydre Handling; den er mere universel; den er mere elementær.

a) En moralsk Handling maa være udsprungen af den Handlendes
egen indre Følelse og Villie. Men Retshandlingen kan
fremtvinges ved ydre Magt, ligegyldigt, om den Handlendes
Sind stemmer med den eller ikke. Muligheden af Tvang er
Rettens ydre Kendemærke. Retsordenen kundgør sig ved fysisk

') Smlgn. Platons Apologi. Kap. 17. --- Apostlenes Gerninger IV,
19; V, 29. --- Jurister holde ikke af at tænke sig den Mulighed, at den
Enkelte kan have etisk Ret overfor „Retten“. Selv A. S. Ørsted
siger: „Enhver er pligtig at underordne sin private Mening den, som
ligger til Grund for Lovene. Idetmindste maa han af enhver offentlig
Myndighed betragtes som pligtig dertil“. (Om de første Grundsætninger
for Straffelovgiv ningen. --- Eunomia II, p. 202). --- den danske borgerlige
Straffelov § 42 tales der om „den urigtige Formening, at den ved
Loven forbudne Handling efter Samvittighedens eller Religionens Bud
er tilladt eller endog befalet“, og det synes altsaa, som om det slaas fast,
at Uretten altid maa være paa den Enkeltes Side. Saa er Ørsted forsigtigere,
naar han blot hævder, at den offentlige Myndighed nødvendigvis
% --- p. 523 ---
\setcounter{footnote}{0}%
\opage{523}Magt, medens den moralske Verdensorden er en indre. Sanktionen
er ved Retten en ydre, ved Moralen en indre. (Smlgn.
IV, 5---6).

b) Hermed hænger det nøje sammen, at Retten blot fordrer
ydre Handling. Sindelaget og Villien bryder den sig ikke
om; dem kan den ikke ramme. For den rent etiske Vurdering
kan derimod den ydre Handling ofte blive af forsvindende Betydning
i Forhold til de Motiver, som have ført til den. Disse
Motiver kunne være højst betænkelige, skønt den ydre Handling
i og for sig ikke er fordærvelig; og de kunne være anerkendelsesværdige,
skønt den ydre Handling i og for sig er forkastelig.

c) Derfor kan Retten ogsaa kun angaa saadanne Forhold, i
hvilke det Samme med saa stor Tilnærmelse som muligt kan
fordres af Alle. De store individuelle Forskelligheder træde mere
frem i Sindelag og Motiver end ide ydre Handlinger. Mennesker
af samme Folk ligne hverandre mere i deres Færd end i deres
Tankegang. Handlingen er ikke saa individuel som Sindelaget.

d) Og hvad der saaledes kan og maa fordres af Alle, og
som kan gøres, ligegyldigt hvilke Motiverne ere, og derfor ogsaa
kan fremtvinges ved Magt, --- det kan kun være visse elementære
etiske Handlinger, det vil sige, Handlinger, som gaa ud
paa Opretholdelse af de allermest nødvendige Betingelser for
menneskeligt Samliv. Retten tenderer stedse --- eller bør stedse
tendere --- mod et Minimum. Det ligger allerede deri, at Tvangen
altid er et Onde, og derfor, i Kraft af Velfærdsprincipet,
bør indskrænkes til det mindst Mulige, det Uundværligste. Retten
skal ved sine strenge Bestemmelser fastsætte de Grænser,
ned under hvilke Menneskets Handlen ikke maa synke, men
over hvilke den kan hæve sig saa højt, den formaar.

3. Retsordenen er en af Menneskene dels uvilkaarligt dels
vilkaarligt tilvejebragt Naturorden. Skønt den er forskellig fra
den etiske Orden, virker den dog som Forberedelse og Grundlag
for denne. Den øver dels en indirekte dels en direkte Opdragelse.
Indirekte virker den ved at kue alle Forsøg paa at
sætte sig ud over de nødvendigste Betingelser for menneskeligt
Samliv. Derved hæmmes efterhaanden de Tendenser, som stride
% --- p. 524 ---
\setcounter{footnote}{0}%
\opage{524}mod disse Betingelser. Direkte opdrager den ved at afgive en
fast Ramme, indenfor hvilken de frie Kræfter trygt kunne røre
sig. Den betinger en Følelse af Sikkerhed og Stadighed, uden
hvilken enhver højere Bestræbelse bliver umulig.

Men Sammenhængen mellem Ret og Moral er langtfra udtømt
hermed. Den Opgave at hævde Retten ved Magten tilfalder
naturligt dem, der raade for Samfundets samlede koncentrerede
Magt. Den etiske Berettigelse og Nødvendighed af denne
Magtanvendelse beror paa, at det Heles Tarv staar over den
Enkeltes Lyster. Det er den samme Forudsætning, som gælder
for enhver Etik, der ikke hylder Øjebliksstandpunktet eller det
individualistiske Standpunkt. Det er en etisk Opgave at hævde
de Grundbetingelser, uden hvilke Samfundslivet ikke kan udvikle
sig. Der maa nødvendigvis sættes Skranker for de enkelte
Individers Trang til at brede sig i Tilværelsen, dersom et Samfundsliv
skal være muligt. --- Ikke blot ved sine Virkninger,
men ogsaa ved den Forudsætning, hvorpaa den bygger, faar den
ved Magt hævdede Ret altsaa en etisk Karakter.

Trods den stadige Mulighed for en Konflikt mellem Ret
og Moral, er Retten dog stedse etisk betinget og Retsordenen
en Del af den etiske Orden. Netop dette er det, som giver en
saadan Konflikt dens tragiske Karakter. Det er et Sammenstød
af etiske Magter, eller rettere af Magter, der hver for sig maa
gøre Fordring paa etisk Myndighed, --- ti naar der finder et
Sammenstød Sted, maa en af Parterne have mistet sin etiske
Karakter. Hvilken af dem der i det enkelte Tilfælde har forskertset
sin Myndighed, det er netop det store Spørgsmaal. Er
det min Samvittighed, som, naar den strider mod en gældende
Ret, blot er et Hjernespind, eller er det Lovbuddet, som, naar
den strider mod min Samvittigheds klare Udsagn, blot er en
vilkaarlig og uforsvarlig Fordring? Den sidste Afgørelse vil i
alle Tilfælde ligge i den Enkeltes Samvittighed; ti Samvittigheden
dømmer Alt, ogsaa sig selv (IV, 3). Vælger jeg at bøje
mig for den positive Ret, maa jeg ligesaa vel kunne forsvare
det for min Samvittighed, som hvis jeg vælger at bryde den
positive Ret.

Goos, hvis Udvikling af Forholdet mellem Ret og Moral
% --- p. 525 ---
\setcounter{footnote}{0}%
\opage{525}for mig har været meget lærerig, udtaler!), at den positive Ret
kun forspilder sin etiske Berettigelse, „naar et under Rettens
Navn raadende Despoti steg til en saadan Højde, at det fra
Sædelighedsideens Standpunkt kunde siges, at ingen Ret var
bedre end en saadan Ret“. Grænsen for Rettens Uret overfor
Moralen synes dog her at være afstukken for snevert. Ikke
blot hvor Despotiet, men ogsaa hvor Trægheden hersker i Rettens
Navn, kan en positiv Retsbestemmelse miste sin etiske Karakter.
En positiv Lov kan staa i Vejen for en Ordning, den
etiske Lov kræver, og det er ikke sagt, at den etiske Fordrings
Fyldestgørelse kan vente, til den gældende Ret ad den i Statens
Love foreskrevne Vej afløses af en bedre Ordning. Ved
enhver gavnlig Reformation og Revolution sker der et etisk
Gennembrud, som sprænger den positive Lovs Skranker. Under
sædvanlige Forhold kunne de etiske Fordringer paa stille og
jevn Vis faa Afløb over i den positive Rets Former; men der
kan være opsamlet en saadan etisk Spændkraft, at et pludseligt
Gennembrud bliver nødvendigt. Ikke blot altsaa, naar den positive
Ret staar meget lavt, men ogsaa naar den etiske Fordring
staar meget højt, vil Konflikten komme, og komme som en etisk
Nødvendighed.

4. Hvor bestemt Forskellen mellem Ret og Moral gør sig
gældende, naar et højere Udviklingstrin er naat, kan ses af, at
man netop fra et alvorligt etisk Standpunkt har villet hævde,
at Ret og Moral slet ikke have nogen fælles Rod, og at Retten
ikke har nogen etisk Karakter. \textsc{Kant} og \textsc{Fichte} opfattede Retten
som en rent ydre Ordning, der nødvendigvis maa være til
Stede, naar selvstændige Væsener skulle leve sammen. Den ene
Persons Frihed maa indskrænkes saaledes, at den anden Persons
ligesaa store Frihed bliver mulig. Men Retten fastsætter
kun en ydre Ordning, forbyder faktiske Overgreb paa Andres
Frihed. Retten fordrer blot Legalitet, ydre Overensstemmelse
med de for flere Personers Samliv gældende Regler. Etiken fordrer
derimod Moralitet, indre Tilslutning til den erkendte Pligt.
Ligesaa lidt som den gode Villie har Noget at gøre paa Rettens
% CHECK p. 525: note whose mark was not found
\footnote[1]{gjjj\textsuperscript{1} Almindelig Retslære I. p. 47.}
% --- p. 526 ---
\setcounter{footnote}{0}%
\opage{526}Omraade, ligesaa lidt har Legaliteten som saadan nogen etisk
Betydning.

\textsc{Fichte} har gjort det mest karakteristiske Forsøg paa at
begrunde Retslæren som en selvstændig, af Etiken aldeles uafhængig
Videnskab. Rettens Princip forudsætter blot den Beslutning
hos et Individ, at det vil leve sammen med andre Individer.
Hvorfor Individet fatter denne Beslutning, om det sker
af etiske eller af ikke-etiske Grunde, vedkommer ikke Retslæren.
Med den, der ikke vil leve sammen med andre Individer, har
Retslæren Intet at gøre; intet af dens Argumenter gælder overfor
ham. Men foreligger der en Beslutning af nævnte Indhold,
saa viser Retslæren, at den kun kan gennemføres, naar Individet
er villigt til at indskrænke sin Frihed. Har man sagt A,
maa man ogsaa sige B, og vil man ikke sige B, maa man ogsaa
lade være at sige A: den, der ikke vil indskrænke sin Frihed»
kan ikke leve sammen med andre Individer. Vi komme altsaa
--- siger Fichte --- til Retten, uden i mindste Maade at forudsætte
Moralen. --- Men dertil kommer endnu, at Moral og Ret
kunne stride mod hinanden. Den moralske Lov forbyder mig
at tage den Fattiges sidste Lam, naar han ikke kan betale sin
Gæld til mig; men den juridiske Lov tillader mig det. Hvorledes
kan en saadan Tilladelse have nogen etisk Begrundelse?
Af en fælles Forudsætning kan der ikke udspringe to hinanden
modstridende Slutninger\footnote[1]{J. G. Fichte: Grundlage des Naturrechts. Jena und Leipzig 1796. Einleitung § 2 og Erstes Hauptstiick § 4. --- Før Fichte havde Anselm Feuerbach udtalt en lignende Teori, dog ikke i saa gennemført en Form. (Smlgn. A. v. Feuerbach: Biographischer Nachlass, I. p. 23).}.

Hertil er for det Første at sige, at Retsordenen meget godt
kan have etisk Betydning, fordi den hævdes ved ydre Tvang»
og fordi Individet allerede af rent egoistiske Motiver ledes til
at anerkende den. Naar Retten hævder de mest elementære
etiske Fordringer, er det intet Under, at den slaar Noget fast,
som man ogsaa fra helt andre Standpunkter end det egentligt
etiske kan erkende for nødvendigt. Sikkerhed, Orden og fast
Organisation kunne godt være etiske Fornødenheder, fordi de
% --- p. 527 ---
\setcounter{footnote}{0}%
\opage{527}ogsaa ere egoistiske Fornødenheder. De faa etisk Værdi ved
det, der kan udfolde sig paa Grundlag af dem, og de udspringe
af Velfærdsprincipet ligesom alle etiske Fordringer.

Hvad dernæst Muligheden af en Konflikt angaar, er der
intet Underligt i, at en saadan kan opstaa, da de rent elementære
etiske Fordringer nødvendigvis maa være rummeligere og
lavere end de højere Fordringer, der ikke kunne hævdes ved
ydre Magt. Min Pligt sætter snevrere Grænser for min Handlefrihed
end min juridiske Ret. Retten er, som tidligere (2) bemærket,
mere universel end Pligten, og den etiske Berettigelse
(smlgn. VIII, 1) falder ikke sammen med den juridiske Berettigelse.
Bliver jeg staaende paa det juridiske Standpunkt og
skyder de strengere etiske Fordringer til Side, er det intet Under,
at en Konflikt opstaar. --- At den lavere Kreds af Rettigheder
maa existere, betinges ved, at der ellers vilde opstaa værre
Ulemper end de, der følge med en Konflikt mellem Ret og Moral.
Kreditor kan benytte sig af Loven til umenneskelig Adfærd
mod Debitor; men dersom der ikke her gjaldt haarde Bestemmelser,
vilde Letsindighed i at laane og fortære fostres, da
man kunde stole paa stedse at beholde Noget tilbage. En saadan
uetisk Adfærd fra Debitors Side vilde der være langt større
Fristelse til, end der er til umenneskelig Benyttelse af Kreditors
juridiske Ret. Naturligvis bør Bestemmelserne ikke være
haardere end strengt nødvendigt, og der er vel ikke Spørgsmaal
om, at der her som paa saa mange andre Punkter er store Fremskridt
i Humanitet endnu at gøre.

Vel gør Retten ikke Regning paa Mere end Legalitet, ydre
Overensstemmelse med de gældende Love. Men intet Retssamfund
kan i Virkeligheden bestaa, naar ikke andre end rent egoistiske
Motiver føre til Lydighed mod Retsordenen. Ska! Retsordenen
være mere end en ordnet Krigstilstand, maa den have
sit subjektive Grundlag i Mere end den blotte Egeninteresse\footnote[1]{Smlgn. Om Grundlaget for den humane Etik. p. ff. ff.}.
Og dette vil den faa, naar den svarer til sit Formaal, at værne
og frede om Livets Udvikling i de forskellige Retninger. Menneskene
ville da se op til Retsordenen som Betingelsen for de
% --- p. 528 ---
\setcounter{footnote}{0}%
\opage{528}største Goder. Den vil indgyde dem ikke blot Frygt og Afhængighedsfølelse,
men ogsaaTaknemmelighed og Beundring. De ville
i den se en Ramme, indenfor hvilken Alt, hvad de ønske og tilstræbe,
trygt kan udvikle sig, og selv naar de indse Nødvendigheden
af at kritisere og ændre den, ja maaske endog helt bryde
med det Spor, paa hvilket den er kommen ind, saa er det, fordi
de ville gøre den bedre i Stand til at udføre sit Hverv indenfor
den etiske Verden.

Det kan undertiden netop være etisk Pligt at gøre hensynsløs
Brug af den Ret, man besidder efter den gældende Orden,
for derved at hævde dennes Gyldighed og rette Forstaaelse
og fremme Ærbødigheden for de enkelte Individers Frihed og
Selvstændighed. Kampen for Retten er Kamp for et af Menneskelivets
vigtigste Goder, om end Jhering\footnote[1]{Kampen for Retten. Dansk Overs. p. 48 ff.} gaar for vidt, naar han
paastaar, at af de to Bud: „gør ikke Uret!“ og „taal ikke Uret!“
er det sidste det vigtigste. Det betegnede et stort etisk Fremskridt,
da den Sætning første Gang opstilledes, at det er bedre
at lide Uret end at gøre Uret (se XII, 2). Jherings Paastand
er kun rigtig, hvor man af Magelighed og Ligegyldighed finder
sig i Uretten, eller hvor man har mistet Blikket for Rettens praktiske
Betydning for det menneskelige Liv.

5. Der er to Hovedpunkter i Rettens Udviklingshistorie, som
ere af særlig Interesse i etisk Henseende. Det ene angaar den
Myndighed, som fastsætter og hævder Retten, det andet det Subjekt,
som besidder Retten, eller som skal fyldestgøre Rettens
Fordring.

a) Saalænge Retten endnu ikke udsondres fra Sæd og Skik,
men indeholdes i disse i nøje Forening med Moral og Religion,
kan den ikke formuleres i klare og bestemte Sætninger. Vejen
staar aaben for megen Vilkaarlighed og Tilfældighed. Det er
derfor et stort Fremskridt, naar den i Samfundet raadende Magt
vedkender sig bestemte Sætninger som Norm for sin Optræden.
Magten underlægger sig derved en højere Lov og kan nu vurderes
efter, hvorvidt den fyldestgør denne Lov. Den kan maales
med en Alen, den selv har anerkendt. Det er et vigtigt Vende%
% --- p. 529 ---
\setcounter{footnote}{0}%
\opage{529}punkt i Autoritetens Historie!). Magtens Inkonsekvenser ville i
Længden komme over dens eget Hoved. Selvom den oprindeligt
kun har taget Retten paa som en Slags Forklædning, saa har
den derved udæsket en farlig Fjende. Den har udæsket Kritiken,
og den vil ikke kunne stanse denne igen uden ved systematisk
at tilintetgøre enhver frittænkende Bevidsthed. At den
ikke formaar dette, viser sig ved, at den i Regelen --- selvom
det sker ved nok saa mange Fordrejninger og Sofisterier ---
søger at skaffe sig et Skin af retslig Begrundelse for sin Adfærd.
Naar Kritiken først er begyndt paa ét Punkt, gaar den snart
videre. Den indskrænker sig ikke til at prøve Overensstemmelsen
mellem den enkelte Magtanvendelse og den Ret, Magten
selv har erkendt som gældende; men den vender sig mod selve
den gældende Ret og fordrer dens Uligheder og Ulemper fjernede.
Allerede gennem den offentlige Kritik og Debat faa Folkets Tanker
og Følelser Indflydelse paa Rettens Fastsættelse og Anvendelse.
Magthaverne selv ville i Længden ikke kunne unddrage
sig denne Indflydelse; i deres egen Bevidsthed gøre disse Tanker
og Følelser sig til en vis Grad gældende, ligesom allerede
paa Sædvanerettens Standpunkt Magthaverne ogsaa føle Overleveringens
Magt (XXXVI, 5). --- Man har endog prist det som
den sande frie Forfatning, hvor den endelige Afgørelse af Sagerne
vel ligger hos den absolute Monark, men hvor der hersker
ubetinget Diskussionsfrihed, saa at der øves en Kontrol af
Alle overfor Alle. Monarken slaar da blot det fast, der som Frugt
af den alsidige Debat har udviklet sig i Folkets Sind\footnote[1]{Jomfru Marias ubesmittede Undfangelse, fordi den er slaaet fast af Paven; men den blev slaaet fast, fordi Katolikerne troede paa den!}. Men
her er der dog endnu for stor Afstand mellem Debatten og Afgørelsen
og ingen Garanti for, at denne virkeligt kommer i saa

*) Om Grundlaget for den humane Etik. p. 75—79.

\footnotemark[1] Smlgn. F. C. Sibbern: Dikaiosyne. Kbhvn. 1843. Se min Afhandling
om Sibbern i „Mindre Arbejder“ I. p. 99---103. --- Paa lignende Maade
skulde i Følge Kardinal Newman (Apologia pro vita sua. 1864. p. 271 ff.)
selve den ufejlbare Pave kun have den Opgave at „definere“, hvad Kirken
allerede i Forvejen tror paa: han skulde blot slaa Resultatet af Trosudviklingen
indenfor Kirken fast. Katolikerne, sagde Newman, tro ikke paa
% CHECK p. 529: note mark without a note
% --- p. 530 ---
\setcounter{footnote}{0}%
\opage{530}nøje Forbindelse med hin som muligt. Retten kommer derfor
endnu ikke i den inderlige Sammenhæng med Folkebevidstheden,
som er nødvendig for dens sikre Bestaaen. Det sker kun, hvor
Folket har politisk Frihed, og dermed direkte Indflydelse paa
Rettens Fastsættelse gennem Lovgivningen. Kun derved vil det
ogsaa være muligt, at Magthavernes Særinteresser mere og mere
kunne underordnes Samfundets fælles Interesser. --- Men den
levende Retsfølelse i Folket er --- hvilken Forfatning man saa
end tænker sig --- stedse det sidste Værn for Retsordenen\footnote[1]{Smlgn. Jhering: Der Zweck im Recht. I. 2. Aufl. p. 381: „Paa den nationale Retsfølelses moralske Magt beror i sidste Instans Rettens hele Sikkerhed. Ikke paa Forfatningen --- man udtænke sig den saa kunstig man vil; der kan ikke tænkes nogen, som faktisk berøvede Statsmagten Muligheden af at træde Loven under Fødder. Ikke paa de Eder, ved hvilke man mener at sikre den --- Erfaringen viser, hvor ofte den brydes. Ikke paa den Hellighedens og Ukrænkelighedens Nimbus, med hvilken Teorien beklæder Loven --- den imponerer ikke Vilkaarligheden“.},
ligesom den er den Kilde, af hvilken denne fra først af udvikler
sig. Hvor denne Kilde stanser, er Livet ude, selvom Statsmekanismens
Hjul endnu en Tid kunne blive ved at klapre.

b) Paa primitive Trin er det ikke de enkelte Individer, men
Slægterne, Familiegrupperne, Retsordenen har med at gøre. Individet
betragtes kun som Led af sin Familie og sin Stamme.
Indenfor disse Grupper hersker der Fællesskab i Ejendom, maaske
endog i Henseende til Kvinder og Børn. Ægteskabet er ikke
en Sag mellem to Individer, men mellem to Slægter, som derved
komme i Forbindelse med hinanden. Sker der et Mord
eller et Ran, rammes hele den Slægt, til hvilken Gerningsmanden
hører, af Gengældelsen. Der existerer ingen individuelle Forpligtelser
og ingen individuel Skyld. Først efterhaanden som
den instinktmæssige Sædvane udvikler sig til egentlig Retsfølelse,
foregaar der en Emancipation af de enkelte Individer, ved hvilken
de hver for sig blive til retslige Personer, Midtpunkter for
Rettigheder og Forpligtelser, Subjekter, der kunne paadrage sig
Skyld og bøde for den. Endnu findes der Spor af den primitive
Tankegang i vor Retsorden (smlgn. f. Ex. XXII, 1).

6. Et Mellemled mellem positiv Ret og etisk Overbevisning
% --- p. 531 ---
\setcounter{footnote}{0}%
\opage{531}dannes af den offentlige Mening. Den kan defineres som det
populære Udtryk for den positive Moralitet. Den er et ..moralsk
Politi“, der øver en mægtig Indflydelse ved sin Ros og sin
Dadel. Den trækker noget snevrere Grænser for det
Berettigede end Retten. Retten holder sig til den ydre Handling; den offentlige
Mening gaar videre og fælder en Dom om Karakteren. Der
kommer for dens Domstol Forhold i Betragtning, som den juridiske
Domstol ikke kan tage Hensyn til. Derfor frikender den
undertiden dem, der dømmes juridisk, skønt den endnu oftere
anklager dem, der ikke kunne stilles for en juridisk Domstol.

Den kan danne sig paa meget forskellig Maade og være
af meget forskellig Værdi. Ofte kan dens Oprindelse være vanskelig
at efterspore, fordi den er Udtryk for en lang Række stadige
Erfaringer, der hver for sig ikke have draget Opmærksomheden
til sig, men hvis Resultat først træder klart frem. I andre
Tilfælde kan det være pludselige Erfaringer, som gøres i Fællesskab,
og som strax sætte samme Stemning i Bevægelse hos
Alle. Paa begge Maader komme Meninger til at „ligge i Luffen
«, og paa begge Maader kunne de gøre sig gældende med
stor Styrke. Det Ubekendte i deres Udspring giver dem en vis
mystisk Karakter og paatrykker dem Stempel af det Selvfølgelige.
Ofte kan det være enkelte fremragende Mænd, der bestemme
den offentlige Mening ved deres Autoritet, især naar
de formaa at iklæde deres Tanker slaaende Ord. Men det kan
ogsaa være ukendte og ubetydelige Personers ihærdige Gentagelser,
som faa Meninger til „at ligge i Luften“. Der finder
en stadig Vexelvirkning Sted mellem den snevre Kreds af Mennesker,
der gøre nye Erfaringer og faa nye Ideer, og den store
Kreds, som overvejende modtager deres Meninger paa passiv
Maade. Ingen vil, hvor aktivt hans Tanke- og Følelsesliv ogsaa
er, kunne undgaa at blive bestemt ved Tidens og Folkets eller
Standens og Familiens Forudsætninger og Overleveringer; men
man kan stille sig i et mere eller mindre aktivt Forhold til
disse og have et mere eller mindre vaagent Øje for nye Muligheder.
I forskellige Lande vil Forholdet mellem dem, der aktivt
forme de gældende Meninger, dem, der passivt optage dem,
og dem, der --- --- ingen Mening have, være forskelligt. James
% --- p. 532 ---
\setcounter{footnote}{0}%
\opage{532}\textsc{Bryce}, der i sit Værk om den amerikanske Republik anstiller
en interessant Undersøgelse af den offentlige Menings Opstaaen
og Betydning, mener, at i de forenede Stater er Antallet af dem,
der høre til den første Klasse, meget ringe, medens det i England
er forholdsvis stort, men at til Gengæld den tredie Klasse
hist er forholdsvis langt mindre end i England — og vilde være
endnu langt mindre, hvis Indvandrere og Negere ikke satte Forholdstallet
ned. Det er derfor i de forenede Stater vanskeligere
end andre Steder at paavise den offentlige Menings successive
Udvikling; dens Opstaaen og Udvikling har dér endnu mere end
andre Steder Karakteren af en Væxt. Der er faa overlegne Intelligenser,
men der er en stor Mængde Mennesker, som ere i
Stand til at optage Ideerne, og disse spredes overordentligt hurtigt
i store Kredse\footnote[1]{Holtzendorff: Wesen und Werth der öffentlichen Meinung. Miinchen 1879. p. 93.}. --- Ved Udbredelsen virker ofte kun i ringe
Grad Eftertanke og Overbevisning med, men Efterligningstrangen
spiller en Hovedrolle, og ikke sjeldent bliver en offentlig
Mening til, blot fordi den allerede synes at være til!\footnotemark[1]

Den offentlige Mening kan ofte etisk set staa højere end
nogen af de Personer, der bidrage til den og hylde den. Dette
er i sig selv ikke underligere end, at Retsordenen ogsaa staar
højere end den virkelige Tilstand i Landet, hvilket allerede
Straffebestemmelsernes Nødvendighed viser, --- eller end, at
ogsaa den Enkeltes etiske Overbevisning kan have en Idealitet,
som hans virkelige Villen og Handlen er langt fra at besidde.
Men det hænger sikkert ogsaa tildels sammen med, at vi hver
for sig dadle Splinterne hos Andre og glemme Bjælkerne hos
os selv; der resulterer deraf i den offentlige Mening en Strenghed,
som de Enkelte hver for sig vel vogte sig for at anvende
i deres Selvkritik\footnotemark[1].

') James Bryce: The American Commonwealth. London 1885. III.
p. 10 —13; 99---105. --- Bryce fandt særligt i Nordamerika et mere kritisk
Blik paa Pressen hos Publikum idethele; det amerikanske Publikum er
køligere og forsigtigere end det europæiske og imponeres ikke saa let af
det hemmelighedsfulde „vi“. (p. 37 f.).
% CHECK p. 532: note whose mark was not found
\footnote[2]{Smlgn. Riimelin: Ueber den Zusammenhang der sittlichen und intellectuellen Bildung. (Reden und Aufsatze. Neue Folge. p. 24).}
% CHECK p. 532: note mark without a note

% --- p. 533 ---
\setcounter{footnote}{0}%
\opage{533}Ligesom Retsordenen saaledes bør ogsaa den offentlige Mening
indskrænkes til et nødvendigt Minimum. Den er mere nærgaaende
end Retten, idet den ogsaa bedømmer Karakteren; men
desuagtet raader den ikke over de grundige Metoder, der staa
til Retsordenens Tjeneste, naar det gælder om at belyse en
Handlings Natur. Karakteren er vanskeligere at bedømme end
den ydre Handling. Villiens Udviklingsgang vil altid være mere
eller mindre skjult for de Udenforstaaende. Den offentlige Mening
er alligevel ikke tilbøjelig til at erkende sine Skranker; den
tiltror sig let Alvidenhed og bliver intolerant. Den skærer Alle
over én Kam. Den beherskes ofte af Stands-, Race-, Parti- og
Religionsfordomme. Trods alt dette spiller den dog en vigtig
Rolle baade overfor Samfundsmagten og overfor de enkelte Individer.
Den er paa sin Vis, ligesom Retsordenen, en Art Naturorden,
som sætter Skranker og Betingelser. Den øver en Kontrol,
som ikke kan undværes. Men ligesom Retsordenen tilsidst
udspringer af Folkets Retsfølelse og kun bestaar ved Sammenhængen
med denne, saaledes er det ogsaa af største Betydning,
at den offentlige Menings Strøm faar Tilløb fra den alvorlige
Overbevisning hos de Enkelte, saa at den ikke bliver flygtig og
overfladisk. Og langt hyppigere, end der indtræder Konflikt mellem
Ret og Moral, vil der indtræde Konflikt mellem den offentlige
Mening og den Enkeltes Samvittighed.

Det var især Frygt for, at den offentlige Menings voxende
Magt skulde blive farlig for Karakterens Selvstændighed, der i
sin Tid førte \textsc{Stuart} \textsc{Mill} til at skrive sin Bog „Om Friheden“.
Men den Løsning, han gav af Spørgsmaalet, var, som jeg ovenfor
(VIII, 5) har søgt at vise, uheldig. Man kan ikke give Mill
Ret i, at Individets indre Selvudvikling kun vedkommer det selv
og ikke Andre. Der gives ikke den Karakteregenskab hos et
Menneske, som ikke kan faa Betydning for hans Forhold til
Andre og for hans Stilling i Samfundet. Og man vil heller ikke
kunne hindre de Andre i at anstille en Vurdering af det Karakterbillede,
der danner sig af ham i deres Bevidsthed. Det er
ikke selve Berettigelsen til at dømme, men det er Usikkerheden
i de Data, hvorefter der dømmes, som Vægten maa lægges paa,
naar den offentlige Mening skal vises tilbage indenfor sine rette
% --- p. 534 ---
\setcounter{footnote}{0}%
\opage{534}Grænser. Paa den anden Side har den Enkelte ofte selv Skylden
for den offentlige Miskendelse, der bliver ham til Del, naar
han ikke gør nok for, at hans Adfærd kan komme til at staa i
det rette Lys for Andre. Man har ikke Lov at lade sig gøre
til Martyr, naar man ikke har gjort alt Muligt for at stille sin
Sags virkelige Sammenhæng saa klart frem som muligt. Selv
Ironikeren Sokrates tog Bladet fra Munden, da det blev Alvor,
og da det gjaldt om at hindre hans Landsmænd i at begaa en
Forbrydelse.

\chapterhead{XXXVIII.}{STATENS ETISKE BETYDNING.}

1. Mest direkte og iøjnefaldende vilde Statens etiske Betydning
være, dersom de havde Ret, der opfatte Staten som en
umiddelbar Aabenbaring af det Godes Idé, som en etisk Autoritet,
der staar uafhængig af og hævet over de individuelle Samvittigheder.
Statsmagten vilde da være Individernes etiske Formynder,
som de havde at bøje sig for.

I sin mest extreme Form fremtræder denne Opfattelse, hvor
man paastaar, at Staten skal bygges paa et teologisk Grundlag.
Staten er her ikke blot en etisk Magt, men --- føjer man til ---
da al Etik beror paa religiøs Autoritet, maa Staten tilsidst bygge
paa religiøs Grund. Staten behøver --- hedder det --- „en øverste,
tilforladelig Maalestok“ for at kunne anstille en Vurdering af de
menneskelige Formaal, som det er dens Opgave at fremme. Udvortes
Retfærdighed kan ikke gennemføres uden indvortes Retfærdighed,
og denne forudsætter igen et religiøst Sindelag, som
kun kan vindes ved Kristendommen. Staten maa altsaa være
kristelig Stat. Dette forudsætter „ikke, at levende personlig Kristendom
skal findes hos Alle i Folket, men at Folket bøjer sig
under den kristelige Overleverings Autoritet“. Statens Kristelighed
viser sig først og fremmest i, at den støtter Kirken; der%
% --- p. 535 ---
\setcounter{footnote}{0}%
\opage{535}næst i. at den sørger for kristelig Sæd og Skik i Folket og for
kristelige Skoler, og at den idethele indpræger de kristelige Grundsætninger
i sine Love og Indretninger\footnote[1]{Martensen: Den sociale Etik. p. 125 ff. Martensens Begreb om ..kristelig Stat“ betyder dog et stort Afslag i Forhold til Stahls, der krævede udelukkende Beskyttelse for den kristelige Kirke og kristelig Trosbekendelse som Betingelse for alle Embeder og for Sæde i Repræsentationen. Og Stahls Lære er igen et Afslag i Forhold til, hvad den hellige Alliance krævede. (Smlgn. Bluntschli: Politik. p. 226---228). Den katolske Opfattelse gaar naturligvis videre endnu, idet den fordrer Kirkens Overhøjhed anerkendt i alle Spørgsmaal, der baade angaa Kirke og Stat. Paven protesterede mod den hellige Alliance, fordi Fyrsterne \_der vilde optræde som Jesu Stedfortrædere.}.

Der er noget Tvetydigt i denne Lære. Den siger os ikke
med rene Ord, hvor den tilforladelige Maalestok egentligt findes.
Kristelig Overlevering og kristelige Grundsætninger ere løse Begreber,
især naar det udtrykkeligt indrømmes, at man ikke forudsætter,
at Alle i Folket have levende, personlig Tro paa Kristendommen.
Ingen kan nægte, at Kristendommen har givet betydningsfulde
Bidrag til Udviklingen af den etiske Bevidsthed
i Menneskeslægten. Men disse Bidrag maa suppleres med Bidrag
fra andre Sider og ere hverken i Form eller Indhold tilstrækkelige
til at tjene som tilforladelig Maalestok. Naar man
fra teologisk Side roser sig af at have en tilforladelig Maalestok,
bliver der kun Mening deri, hvis man konsekvent fastholder Autoritetsprincipet.
Altsaa Kirketroens Autoritet (hvad enten denne
nu er Paven, Bibelen eller Trosbekendelsen) skal være Statens
Grundlag. Det bliver et Teokrati eller en Kirkestat, som følger
heraf. Staten bliver ikke et rent menneskeligt Samfund, men
afhængig af en overnaturlig Autoritet. Om noget Saadant kan
der paa det Standpunkt, hvorpaa vi her have stillet os, ikke
være Tale. Ligesom vi have hævdet Forskningens og Samvittighedens
Frihed, saaledes maa vi ogsaa hævde Statens Frihed og
Selvstændighed overfor det teologiske Autoritetsprincip. Forholdet
er endog, som vi have set i det Foregaaende (XXXIII, 3),
det omvendte af, hvad \textasciitilde{}den kristelige Stats“ Tilhængere tænke
sig: det er Staten, som tilsidst afgør, hvorvidt de forskellige
% --- p. 536 ---
\setcounter{footnote}{0}%
% CHECK p. 536: notes paired with the marks in order (the layer misread a number)
\opage{536}religiøse Konfessioner stemme med „Sædelighed og offentlig Orden“;
den kan ikke selv hente sit Sædelighedsbegreb fra nogen
af dem.

Fra Kristendommens eget Standpunkt er der kun da Mening
i Ideen om en kristelig Stat, naar man kan gaa ud fra,
at Alle i Folket ere Kristne. Men man sagde jo udtrykkeligt,
at man gik ikke ud derfra! Kristendommen skal altsaa raade
over dem, der ikke tro paa den! --- Hele denne Tale om kristelig
Stat er Uklarhed fra først til sidst.

Overfor højkirkelige Teologers Uklarheder er det velgørende
at se \textsc{Grundtvigs} sunde og klare Opfattelse af Statens Væsen.
Saa dybt han personlig var greben af Kristendommen, indsaa
han dog, at „hvor de borgerlige Forhold i Sandhed skal genfødes
og vinde Fasthed, dér maa man lade Kristendommen som
noget aldeles Frit og Uberegneligt aldeles ude af Regningen
\dots{} og holde sig til den Menneskenatur, Historien paa ethvert
givet Sted aabenbarer“\footnote[1]{Mands Minde. p. 127 smlgn. p. 569 f.}.

2. Selvom man ikke vil bygge Staten paa et teologisk
Grundlag, kan man dog mene, at Staten er Udtryk for en højere
Sædelighed, at dens Love lære den Enkelte, hvad der er godt
og rigtigt. Den Enkelte har vel sine subjektive Idealer, men i
Staten træder Idealet ham i Møde som en virkelig Magt. Den
antike Statsopfattelse, saaledes som den udvikledes af Platon og
Aristoteles, gik i denne Retning. I nyere Tid har \textsc{Hegel} især
gjort denne Opfattelse gældende i extrem Form. „Staten er den
sædelige Idés Virkelighed“. „Individet har kun Sandhed og
Sædelighed, forsaavidt det er Led af Staten“. Statens Magt er
„den sig som Villie virkeliggørende Fornuft“\footnote[2]{Hegel: Philosophie des Rechts. § 257--- 258. --- Smlgn. Den nyere Filosofis Historie'\textsuperscript{1} 11. p. 180—183.}.

Denne Opfattelse har det fælles med den foregaaende, at
den ophæver den enkelte Personligheds Selvstændighed og frie
Overbevisning. De enkelte Individer betragtes som blot afhængige
og modtagende i Forhold til Staten, ikke hvert for sig som
konstituerende Medlemmer af den. Man giver Staten et Slags
% --- p. 537 ---
\setcounter{footnote}{0}%
\opage{537}mystisk Liv udenfor eller over dens Borgere. Men Staten lever
kun i sine Borgere. (Smlgn. VIII, 4). Den aandelige og fysiske
Magt, som Staten raader over, er den, som dens Borgere formaa
at yde: de Tanker og Følelser, den Kraft, som de besidde. Den
Sædelighed, som er i Staten, maa være i dens Borgere. Eller
kunde maaske en sædelig Stat bestaa af usædelige Borgere? Det
er vel ligesaa umuligt, som at en Stat kan kaldes kristelig, naar
den ikke bestaar af troende Kristne.

Den omtalte Opfattelse passer bedst paa det Udviklingstrin,
hvor en Modsætning mellen Statsmagt og individuel Frihed
endnu ikke gør sig gældende, --- hvor Forskellen mellem Ret,
Moral, Sæd og Skik endnu ikke er opstaaet. De etiske Forestillinger
ses her ikke som frembragte af det enkelte Individ; de
ere Bestanddele af en ærværdig Overlevering, som de Enkelte
føle sig afhængige af og bøje sig for. Men hos de civiliserede
Nationer suppleres og kontrolleres denne instinktmæssige Tilslutning
til det Overleverede ved den offentlige Kritik, og de
enkelte Individer forholde sig ikke blot modtagende, men virke
aktivt tilbage paa Statslivet.

Ved at tillægge Staten direkte etisk Myndighed nødes man
i Virkeligheden tilsidst til at appellere til den fysiske Magt.
Statsmagten kan ikke lade det komme an paa, om Individet
føler sig overbevist om dens etiske Berettigelse. Den bruger da
sine skarpe Midler, og „den sædelige Idés Virkeliggørelse“ rykker
da frem med Politi og Kanoner. Den overspændte Etik
afslører sig som forklædt Fysik.

3. Man har nu netop fundet Statens Væsen i den blotte
Magt. Det er især Pessimister, som betone dette Moment i Statens
Begreb. Thomas \textsc{Hobbes} opfattede Staten som den ubetingede
Villie, overfor hvilken de Enkelte opgive deres individuelle
Villie, for at Fred og Sikkerhed kan herske. \textsc{Schopenhauer}
betegnede Staten som en Tvangsanstalt, hvis eneste Øjemed
er at beskytte de Enkelte mod hverandre indbyrdes og
hele Samfundet mod ydre Fjender. \textsc{Taine} finder som den egentlige
Kærne bag det hele Statsmaskineri „Gendarmen i Vaaben
mod den Vilde, den Røver og den Afsindige, som Enhver af
% --- p. 538 ---
\setcounter{footnote}{0}%
\opage{538}os skjuler i sit Hjertes Inderste, slumrende eller lænkede, men
endnu stedse levende“\footnote[1]{Hobbes: DeciveV,9. --- Schopenhauer: Die beiden Grundprobleme der Ethik. 2. Aufl. p. 217. --- Taine: L'ancien régime. p. 316.}.

Denne Opfattelse kan ikke blot beraabe sig paa, at Statens
sidste Middel stedse er Magtanvendelse, men ogsaa derpaa, at
de fleste Stater historisk ere grundlagte ved Erobring. Ved
Sværdet er Staten bleven til, og ved Sværdet bestaar den. Men
det har dog stedse vist sig, at den blotte Magt er for usikkert
et Baand. De ved Erobring grundlagte Riger vandt først Fasthed,
da der kunde danne sig fælles Skikke, Livsformer og Traditioner.
Gennem fælles Skæbne og Virksomhed voxe de sammenbragte
Menneskegrupper sig sammen til et Folk. (Smlgn.
XXXVI, 2). Ved den blotte Magt kan man samle en Flok;
men en virkelig Stat forudsætter, at der er et Folk. Hvor Statsmagt
og Folkeliv gaa i helt forskellige Retninger, er der en
ustadig Ligevægt tilstede, og et lille Stød kan være nok til at
kuldkaste det Hele.

Den blotte Magt kan hævde Sikkerheden og Enheden i Folket.
Ingen knyr, og Ingen forlader Geleddet. Men Sikkerhed
og Enhed ere kun Goder, 'forsaavidt som de ere Betingelser for
en fri og tryg Udfoldelse af Folkelivet. Man maa her vogte
sig for at gøre Midlet til Formaal. Jo mere man lader Sikkerhedsmidlerne
brede sig, jo flere Poster man opstiller for at forebygge
Farer og Ulykker, desmere hindrer man den frie Bevægelse,
som maaske ad sin egen Vej vilde kunne føre til at holde
Ulykkerne borte. En nervøs Stirren mod mulige Farer kan ødelægge
hele Livet i Folket, ligesom det enkelte Individ kan hindres
i at leve af Frygt for Døden. Der kan naturligvis være
Tider for Folket, hvor det særligt eller udelukkende bør være
optaget af at sørge for sin Sikkerhed og sin Enhed. Men det
er urigtigt at finde Statens Væsen udtømt i den Rolle, den i
saadanne Tider er nødt til at spille. Centralisationen og den
voldsomme Sammenslutning er en Kraftanstrengelse, der kan
være nødvendig og gavnlig, men som ved at blive stadig og
ubegrænset vilde kvæle al naturlig og selvstændig Udvikling.

% --- p. 539 ---
\setcounter{footnote}{0}%
\opage{539}Desuden kan og skal Magten ikke være blind. Den maa
ledes af Regler og Principer, og disse kunne ikke hentes fra
Magten selv, men maa udledes af de Formaal, Magten skal
værne om. Hvad Magten kan frembringe, Sikkerhed og Enhed,
er kun Middel for højere Formaal, som Magten ikke kan frembringe,
men som den frie Udvikling frembringer, og som Magten
først da kan træde i Forhold til. Af den blotte Magt kan
hverken Kulturen eller Retten udledes.

4. Baade naar man opfatter Staten som umiddelbar Aabenbaring
af det Godes Idé, og naar man opfatter den som den
over Alle og Alt raadende Magt, lægger man Hovedvægten paa
Enheden i Staten og gør den til Et og Alt. Overfor saadanne
Opfattelser gøres det med Rette gældende, at det virkeligt Levende
i Staten dog er de enkelte Individer. Enheden bestaar i
deres Fællesskab og Sammenslutning, men er kun en Form.
hvis Betydning beror paa det Indhold, hvoraf den fyldes. Det
kunde da endog synes, at man maa lægge de levende Bestanddele,
Elementerne, de enkelte Individer til Grund ved Betragtningen
af Statens Væsen og opfatte Statens Væsen som betinget
ved deres Sammenslutning. Enheden bliver da ikke oprindelig,
men afledet. Vi komme herved til den individualistiske
Opfattelse, ifølge hvilken Staten beror paa en Overenskomst
eller Kontrakt mellem de enkelte Individer.

Denne Teori har man ofte gjort Uret ved at antage, at dens
Mening nødvendigvis maatte være, at Staten historisk var opstaaet
ved en Kontrakt. Det er i hvert Tilfælde ikke alle individualistiske
Statsteoretikeres Mening\footnotemark[1]. Teorien bruger Begrebet
Kontrakt for at tydeliggøre, hvorledes Statssamfundet skulde
være ordnet, hvis det svarede til sit Ideal. Man bruger det
altsaa som en Maalestok ved Vurderingen af den virkelige Stat,
ikke som Middel til Forklaring af den virkelige Stats historiske
Opstaaen. Den fuldkomne Stat --- mener man --- vilde være
en saadan, hvor Alt var indrettet saaledes, som de enkelte Individer
ved fri Overenskomst kunde ønske det, --- saaledes altsaa,

) Smlgn. Den nyere Filosofis Historie. Registret under ..Naturret“
% CHECK p. 539: note mark without a note

% --- p. 540 ---
\setcounter{footnote}{0}%
\opage{540}at der var taget Hensyn til Alles Interesser, forsaavidt de kunde
bringes i Harmoni med hverandre.

Hvad man kan efterspore om den historiske Udvikling af
Forholdet mellem Stat og Individ, giver heller ikke den Forestilling,
at Staten historisk opstaar ved Sammenslutning af selvstændige
Individer. Afset fra Erobringer, som jo allerede forudsætte
Stater givne, synes Statsdannelsen at være gaaet for sig
ved Sammensmeltning af Flokke, Grupper eller enkelte Familier,
undertiden vel ogsaa ved simpel Væxt og Udvikling indenfor
en og samme Familie. Om enkelte selvstændige Individer
er der paa tidligere Trin egentligt ikke Tale; de betragtes kun
som Dele af Familien eller Stammen. Men under den historiske
Udvikling foregaar der en Opløsning eller „Søndergnidning“ af
de oprindelige Grupper. Individet emanciperes successivt. Frihedsprincipet
arbejder sig langsomt frem. (Smlgn. XXIII). Denne
Emancipationsproces træder tydeligere og tydeligere frem, naar
man gaar fra Øst til Vest: den er mindre fremtrædende i Asien,
hvor f. Ex. de indiske Landsbysamfund have holdt sig til den
Dag i Dag, end i Europa, og mindre fremtrædende i det østlige
Europa (de slaviske Lande) end i det vestlige (England og
Frankrig)\footnote[1]{Henry Maine: Ancient Law. p. 311. --- Early History of Institutions. p. 386 f.}. En gennemført individualistisk Opfattelse tænker
sig denne Emancipationsproces fuldendt og maaler de historisk
givne Statsformer efter den Grad, i hvilken de nærme sig til
dette Ideal.

Individualismens store Betydning ligger først og fremmest
i dens Fremhæven af Personlighedsprincipet, af hvert enkelt Individs
Værdi. Det er Statens, som ethvert andet Samfunds, store
Opgave at arbejde --- direkte og indirekte --- hen til Udviklingen
af et frit personligt Liv hos saa mange Individer som
muligt. Den fuldkomne Stat vilde være en saadan, der tilfredsstiller
det personlige Liv hos Alle, hver Enkelts Trang til Udvikling
af sine Evner og Drifter. Individualismen har derfor
grebet en væsentlig Side ved Statens Ideal.

Dernæst har Individualismen ogsaa Ret i, at de enkelte In%
% --- p. 541 ---
\setcounter{footnote}{0}%
\opage{541}dividers Overvejelser og bevidste Valg faktisk mere og mere
faa Indflydelse paa Statens Ordning og paa Statssagernes Afgørelse.
Hvor der raader politisk Frihed, har den Enkelte Adgang
til at lade sin Villie blive medbestemmende i de offentlige
Anliggender. Ikke blot ved Grundlæggelsen af nye Stater
eller Statsforbund\footnote[1]{p. 259. Diihring: Cursus der Philosophie. p. 268. Naar det her siges, at kun de Indgreb i Suveræniteten ere tilstedelige, som ere nødvendige for at hævde Alles lige Suverænitet, saa kan man aabenbart ikke udlede disse Indskrænkninger af Individets Suverænitet; hvorfor skulde den begrænse sig selv? Indskrænkningen forudsætter et Standpunkt, som paa én Gang omfatter alle Individer, og hvor Sondringen altsaa ikke gælder.}, eller ved Indførelsen af nye Forfatninger,
men ogsaa ved de enkelte Afgørelser under Statslivets sædvanlige
Gang kommer det mere og mere an paa at tilvejebringe
Overensstemmelse mellem de forskellige medvirkende individuelle
Villier. Ad denne Vej kan Alt, hvad der lever og udvikler sig
i de individuelle Personligheders Indre og i Samfundets mindre
Kredse, komme Statslivet til Gode. Individualismen peger altsaa
ikke blot hen paa det sidste Formaal for Statslivet, men ogsaa
paa et af de vigtigste Midler til dets Fremgang.

Men i sine extreme Former lægger Individualismen ensidigt
Vægten paa Sondringen mellem Individerne, betragter dem
hvert for sig som absolut selvstændige og opstiller Begrebet om
Individets Suverænitet. Et er, at ethvert enkelt Individ er et
ejendommeligt og selvstændigt Udgangspunkt, et Andet er at
tillægge det en absolut Betydning afset fra dets Forhold til
andre Individer. Hvis Ordet Suverænitet skal tages i streng
Forstand\footnotemark[1], vil ethvert Samfund kun kunne blive en Hob eller
en Dynge af Individer; der vil ikke kunne blive nogen indre

') Alexander Hamilton foreholdt under Forhandlingerne om den
nordamerikanske Forbundsforfatning sine Landsmænd, at „det synes at
have været forbeholdt dette Lands Befolkning ved deres Adfærd og Exempel
at afgøre det vigtige Spørgsmaal, om Samfund af Mennesker i Virkeligheden
have Evne til ved Overlæg og Valg at indrette en god Regering,
eller om de ere bestemte til, hvad deres politiske Konstitution angaar,
stedse at afhænge af Tilfældet og Magten“. The Federalist. 1787. Nr. 1.
(Ny Udgave, Boston 1882, p. 49).

-) Hvilketsjældentsker. Smlgn. Stuart Mill: Levned. Dansk Overs.
% CHECK p. 541: note mark without a note

% --- p. 542 ---
\setcounter{footnote}{0}%
\opage{542}Sammenhæng mellem dem. De gaa paa Akkord med hverandre,
men leve ikke et fælles Liv. Et virkeligt fælles Liv bliver kun
muligt, naar Individet fra først af betragter sig og betragtes
som et Led af Slægten og Samfundet. Suveræniteten --- den
absolute Selvstændighed og Magtfylde --- kan ikke ligge hos
nogen Enkelt eller hos alle Enkelte betragtede som hver for
sig isolerede. Den kan kun ligge hos Folket, til hvilket de Enkelte
høre ved Fællesskab i Skæbne og Virksomhed.

Konsekvent maa Individualismen gaa et Skridt videre. Ingen
Statsform kan uafbrudt lade de individuelle Villier raade. Selv
hvor den største politiske Frihed hersker, er det kun ved enkelte
særdeles vigtige Sager, at det hele Folk adspørges, saaledes
ved Folkeafstemningerne (Referenda) i Schweiz og nogle
andre Lande. I Mellemtiden mellem saadanne Afstemninger træde
de individuelle Villier tilbage, og Statens centrale Organer raade
for Tingenes Gang. Den strenge Individualisme maa betragte
en saadan Mellemtid som en Ufrihedsperiode. Det er derfor
konsekvent, naar \textsc{Rousseau} skriver: „Det engelske Folk tror,
at det er frit. Det tager højligt fejl. Det er kun frit, saalænge
der vælges Parlamentsmedlemmer; saasnart disse ere valgte, er
det Slave, er det Intet!“\footnote[1]{Contrat social. 111, 13.}. Rousseau forkastede derfor Repræsentationssystemet:
Suveræniteten kan ikke repræsenteres! Men
han oversaa, at heller ikke indenfor det enkelte Individ staar
den vælgende Villie isoleret fra de andre Elementer i Menneskets
Væsen. Den har kun Betydning ved at være Udslag af hele
Individets Natur og Karakter, og naar Individet tager Konsekvenserne
af sit Valg, bøjer det sig derfor heller ikke for nogen
fremmed Magt, men for sit eget Værk. Det er saaledes heller
ikke noget Skaar i Vælgernes Frihed, at de ere bundne ved
deres eget Valg. Uden en saadan Bundethed er der hverken
Karakter i det enkelte Individ eller Folkekarakter i den hele
Nation. ---

5. At saa forskellige Opfattelser af Staten som de nu omtalte
kunne gøre sig gældende, kommer først og fremmest af,
at Staten er et mangesidigt Væsen, og at dens Begreb inde%
% --- p. 543 ---
\setcounter{footnote}{0}%
\opage{543}holder Elementer, hvis indbyrdes Forhold det er vanskeligt at
bestemme. Men det hænger ogsaa sammen med, at Staten endnu
er i sin Vorden, idet de forskellige Elementer i dens Væsen
endnu ikke ere bragte til fast Sammenhæng indbyrdes, og at
under Udviklingens Gang snart et, snart et andet Element trænger
sig frem i Forgrunden. Magt, Ret og Kultur staa til forskellige
Tider i forskelligt Forhold til hverandre. Spørgsmaalet er nu,
om vi kunne angive et Normalforhold imellem dem, d. v. s. et
Forhold, som bedst svarer til Velfærdsprincipets Fordring.

Af stor Betydning er her den Erkendelse, at det Uvilkaarlige
stedse gaar forud for det Vilkaarlige. De Indretninger og
Ordninger, som Menneskene træffe med fuld Bevidsthed og Vilkaarlighed,
forudsætte stedse, at der uvilkaarligt har udviklet
sig et Indhold, som kan ordnes og anvendes paa en vis Maade.
Vi begynde i vore Indretninger og Ordninger aldrig absolut
forfra. Dette gælder psykologisk set for den vælgende Villies
Forhold til de andre Sider af Sjælelivet, og det gælder sociologisk
set for Statens Forhold til Folket. Staten skal ordne og
værne om Folkets Liv, men den frembringer det ikke fra først
af. --- Først og fremmest er det etiske Liv i sin højeste Form
og Udvikling uafhængigt af Staten. Det skyldes Følelser, som
udvikle sig under Individernes Samliv, men som ikke kunne
fremtvinges ved Magt. Staten kan indgyde Frygt ved sin blotte
Magt; men derved kommer den ikke langt. Ærefrygt indgyder
den først, naar Individet finder, at Staten ledes af de Principer,
det selv føler sig bundet til i sit Indre. Naar Staten fordrer
Underkastelse af etiske Grunde, appellerer den altsaa til en
Magt, den ikke selv har frembragt, men som den forudsætter,
er tilstede. De etiske Principer, som Staten slaar fast og anvender
i sin Forfatning og i sine Love, have først udviklet sig
i Folkets Bevidsthed. --- Staten frembringer heller ikke Familielivet.
Dette har sin Rod i Menneskets Natur; dets Kræfter og
dets Former ere kun i deres ydre Fremtræden og Kundgørelse,
ikke i deres Udspring og inderste Væsen underlagte Statens
Indgriben. --- Kulturen er heller ikke Statens Værk. Den udvikler
sig i det frie Samfund. Baade den materielle og den
aandelige Kulturs Oprindelse fører os stedse tilbage til enkelte
% --- p. 544 ---
\setcounter{footnote}{0}%
\opage{544}Individer, i hvis Bevidsthed de nyeTanker ere blevne til. Indenfor
den Enkeltes Bevidsthed have disse nye Tanker deres første
Kamp for Tilværelsen at bestaa. Naar de træde ud i Verden,
maa de først samle smaa Kredse eller Menigheder om sig,
indenfor hvilke de kunne finde deres videre Pleje og Udvikling, og
derved blive skikkede til at paavirke større Kredse. De betydningsfulde
Udviklingsprocesser i Historien foregaa gerne sporadisk,
ud fra spredte Punkter, ikke fra Centrum. De tage ofte deres
Begyndelse dér, hvorfra man mindst skulde vente det. Man
spørger da forundret: Kan noget Godt komme fra Galilæa? I
en afsides Krog af det store Romerrige fødtes en Kraft, der
skulde leve længere og virke videre end de Magter, man i
Rom betragtede som evige. Hvem kunde i det 12. Aarhundrede
tænke sig, at de Bevægelser, der ytrede sig i de franske „Kommuner“,
vare Indledning til en fuldstændig Omordning af hele
det offentlige og sociale Liv? I mørke Værksteder og i foragtede
Arbejderkredse blev et af det 19. Aarhundredes betydningsfuldeste
sociale Fænomener, de engelske Fagforeninger, til.
Hvad der saaledes har udviklet sig i det Skjulte, kan senere
komme Statslivet til Gode paa mange Maader, men dette kan
ikke selv fra først af frembringe det. --- Det Samme gælder
endog, som vi have set, for Rettens Vedkommende. Den udvikler
sig som en uvilkaarlig Form for Livet i Folket; Staten
kan kun fæstne, systematisere og hævde den.

Staten er ligesaa lidt produktiv som den vælgende Villie
er det hos det enkelte Menneske. Mulighederne, mellem hvilke
der vælges, maa stedse være givne forud i Menneskets Natur.
Valget har naturligvis sin store Betydning, selvom det ikke er
nogen skabende Akt. Hvad der vælges, var til forud, men gaar
nu over i en fastere og bestemtere Form.

6. Den Sætning, at Staten ikke er produktiv, gælder dog
kun med Hensyn til Livets Indhold. Sin Produktivitet udfolder
Staten dér, hvor det gælder om at ordne dette Indhold i Rettens
Form. Her stilles der til enhver Tid en Række betydningsfulde
Opgaver, ved hvis Behandling Statsmandens Genialitet skal vise
sig. Den store Statsmands Begavelse danner et Sidestykke til
den store Religionsstifters. Heller ikke dennes Betydning ligger
% --- p. 545 ---
\setcounter{footnote}{0}%
\opage{545}jo (smlgn. XXXII, 1—2) i, at han frembringer helt nye Tanker,
men deri, at han paa ejendommelig Maade „fortætter“ aandelige
Elementer og Kræfter, som i Forvejen røre sig. Den store
Statsmand er den, der med genialt og sympatisk Blik opdager
de nye Livsspirer i Folket og --- direkte eller indirekte --- yder
dem den Hjælp og Støtte, de behøve. Naar Egoisme og Ærgerrighed
ikke sløve hans Blik, ved han, at han ikke formaar
Andet end at fremhjælpe, hvad der uden ham er spiret frem.
Hans Produktivitet bestaar i at finde Form og Midler, hvorigennem
denne Hjælp kan ydes. --- Denne Opgave vil væsentligt
lettes, hvor der hersker politisk Frihed. Der kommer et
inderligere Forhold mellem Staten og det frie Samfund, mellem
Livets Form og Livets Indhold, naar Folket gennem sine Repræsentanter
deltager i Afgørelsen af, hvorledes Statsmagten
skal anvendes. Først ved den politiske Frihed bliver Staten
virkeligt det organiserede Folk, idet Kultursamfund og Stat nu
kunne komme i levende Vexelvirkningsforhold til hinanden.

\textsc{Gneist}, som paa en saa grundig og interessant Maade har
skildret Selvstyrelsens Udvikling og Organisation i England, betragter
hele den nyere sociale og politiske Bevægelse (for Englands
Vedkommende fra den første Parlamentsreform af) som
en stor Opløsningsproces. Samfundet --- mener han --- har
trængt sig ind i Staten og bemægtiget sig dens Magt for sine
Interesser. Hvorledes han opfatter Forholdet mellem det frie
Samfund og Staten, kan ses af følgende Udtalelse: „Hvad Striden
mellem Pligter og Drifter betyder i den Enkeltes Liv, det betyder
i Folkenes Liv den evige Kamp mellem Stat og Samfund“\footnote[1]{Das Selfgovernment in England. 3. Aufl. p. 1017.}.
Men selvom man, med Gneist, ved Samfundet særligt vil tænke
paa det økonomiske Samfund, i hvilket den materielle Kultur
plejes, saa er det altfor ensidig en Dom. Der kan paa den
materielle Kulturs Omraade udvikles Kræfter og Følelser af nok
saa megen Værdi som de, Statslivet fostrer. Den materielle
Kultur fremmer ikke blot Egennytten og den egoistiske Beregning;
den udvikler Villien, det praktiske Blik og Evnen til
Samvirken med Andre. Og det er intet Under, at de, som føle
% --- p. 546 ---
\setcounter{footnote}{0}%
% CHECK p. 546: notes paired with the marks in order (the layer misread a number)
\opage{546}disse Kræfter hos sig, ogsaa mene at have Ret til at tale et
Ord med i Landets offentlige Anliggender. Det er Betingelsen
for, at Staten kan vedblive at være Udtryk for, hvad der rører
sig i Folket. --- Men det er ensidigt her blot at tænke paa det
økonomiske og materielle Samfundsliv. Den Bevægelse, som har
ført til Valgrettens Udvidelse, hænger nøje sammen med den
intellektuelle og religiøse Udvikling, og om denne kan man
dog ikke sige, at den forholder sig til Statens Virksomhed som
Driften til Pligten\footnote[1]{I mærkelig Modsætning til Gneists gentagne Klager over „det gamle Samfunds Opløsning“ staar en enkelt Udtalelse af ham (p. 996 Note). „Forestillingen om en Opløsning af Samfundet“ --- siger han --- „beror kun paa, at tilvante Forestillinger ikke mere formaa at opdage de vante Skranker“. --- Naar det blot er de gamle Rubriker, man savner, synes der ikke at være saa megen Grund til Klage.}. --- Tilmed indrømmer Gneist, at den gamle
Tingenes Orden havde sine store Skyggesider: man tog ikke
Hensyn til Middelklassens Tarv; man drog ingen Omsorg for
Massernes sædelige og aandelige Løftelse eller for Statens Kulturøjemed\footnote[2]{'0 Anf.  Skr. p. 891. 938.}.
En Statsorden, mod hvilken saadanne Bebrejdelser
kunne rejses, fortjener dog sikkerligt at afløses af en anden
Ordning! Gneists store Beundring for den gammel-engelske Selvstyrelse
er ikke ret forligelig med denne skarpe Dom, der lader
Opløsningsprocessen komme som en Nemesis. ---

Det er ingen Undervurdering af Statens Betydning, naar
dens Produktivitet indskrænkes til at finde og hævde Former
for et Indhold, som er givet uden den. Men denne Opfattelse
giver os, hvis den er rigtig, en Advarsel imod at lægge for stor
Vægt paa det politiske Liv i dets Sondring fra Livets andre
Sider. I Statslivets Kampe glemmes det kun altfor ofte, at
Folkets Kræfter først og fremmest maa udvikle sig paa andre
Omraader, i Familielivet og i det frie Kultursamfund, og at de
kun, naar dette er sket og stadigt paany sker, paa frugtbar
Maade kunne anvendes paa det politiske Omraade. I Politiken
kommer man let til at sætte Formen over Indholdet.

7. I de faste Former, ved hvilke Folkets Liv ordnes, og
ved hvilke det bliver i Stand til at optræde med samlet Kraft
% --- p. 547 ---
\setcounter{footnote}{0}%
\opage{547}udåd til, lægger Folkets Enhed sig tydeligst for Dagen. De enkelte
Individer, Familierne og det frie Kultursamfund frembyde
Billedet af en Mangfoldighed af skiftende Fænomener. Staten
repræsenterer ikke blot Folkets Enhed, men ogsaa de vigtigste
Betingelser for dets varige, udover de enkelte Individers og
Slægtleds Tid rækkende Bestaaen. I Staten har Folket sit mest
koncentrerede Udtryk.

Ved denne Statens Natur er den almindelige Regel given
for, hvilke Opgaver den vil kunne løse. Naar Staten er Folkets
centraliserede Magt, maa det, den skal kunne fremhjælpe, være
saaledes beskaffent, at det kan ledes systematisk udfra ét Punkt,
og at Magtanvendelse kan bruges som sidste Middel til at holde
det oppe. Hvor disse to Betingelser mangle, kan Staten Intet
gøre. Men der kan ikke nævnes noget Omraade, hvor de nævnte
Betingelser ikke paa et eller andet Udviklingstrin og i en eller
anden Grad kunne være tilstede. --- Hvad og hvor Meget Staten
inddrager under sin Virksomhed, vil være forskelligt under forskellige
Forhold. Hvad en Tid ikke kan tænke sig, at Staten
skulde give sig af med, vil maaske til en anden Tid blive dens
Hovedopgave. Enten paatager den sig hele Virksomheden, som
ved Rets- og Forsvarsvæsen, eller den opretter Anstalter, med
hvilke den tillader Privatanstalter at konkurrere, som ved Undervisningsvæsenet,
eller den støtter Privatanstalter med materiel
Bistand, eller den indskrænker sig til at sikre de frie Bestræbelser
den Støtte, som den almindelige Retsorden yder.

Statens virkelige Styrke bestaar i den Tilslutning, Statsmagten
formaar at vinde hos hele Folket. Jo større Frihed,
Staten formaar at overlade Individ, Familie og Kultursamfund,'
uden at dog dens Enhed lider derunder, des stærkere er den.
Den personlige og politiske Friheds Udstrækning er derfor et
af de vigtigste Midler til at bedømme en Stats Udviklingstrin.

Den fulde Harmoni mellem Stat og Folk, Magt og Frihed,
Enhed og Mangfoldighed er et meget fjernt Maal. I
Virkeligheden er enhver Stat en Nødstat, svarer kun højst ufuldkomment
til sit Begreb. --- --- Paa den ene Side optræder Magten endnu
stedse paa mere eller mindre brutal Maade. At øve Magt er
forbundet med altfor stor Tilfredsstillelse for den menneske%
% --- p. 548 ---
\setcounter{footnote}{0}%
\opage{548}lige Egoisme, til at Misbrug ikke skulde indfinde sig. Dertil
komme Frygt og blind Selvopholdelsestrang, som kunne føre
Statsmagten til at brede sig paa Bekostning af de Formaal, den
skulde tjene. Der vil endnu behøves en lang Udviklingsgang,
inden vi komme saa vidt, at Tvang kun anvendes, hvor den i
Kraft af Velfærdsprincipet er aldeles nødvendig. --- Paa den
anden Side ytrer --- ikke sjeldent foranlediget ved Statsmagtens
Overgreb --- Friheden sig ofte som Vilkaarlighed, Kritiksyge,
Mistænksomhed og egoistisk Individualisme. --- Disse to Extremer
udæske gensidigt hinanden. Der maa snart gøres Front
til den ene Side, snart til den anden Side, for at Udviklingen
gennem rytmiske Svingninger langsomt kan føre os frem i Retning
af det angivne Maal.

\chapterhead{XXXIX.}{STATENS STRAFFENDE MYNDIGHED.}

1. Historisk har Straffen udviklet sig af Hævn- eller Gengældelsesinstinktet\footnote[1]{Forud for Gengældelsestrangen gaar dog efter nyere Undersøgelser Trangen til Renselse fra den Besmittelse, som Forbryderen, især Morderen, bragte over hele den sociale Gruppe, hvor hans Gerning var øvet. Gruppens Reaktion var da paa en Gang en Desinfektion og en Exorcisme. Overgangen til Straf bestod i, at Renselsen fuldbyrdedes gennem Forbryderens Landflygtighed eller Død. Se herom Frazer: Psyche's Task. A discourse concerning the influence of superstition on the growth of institutions. (1909). p. 79 f. 11913).}.
I sin simpleste Form ytrer dette Instinkt
sig ved en Bevægelse, som er rettet mod den Kant, fra hvilken
man har modtaget et stærkt og smerteligt Indtryk. Man
farer løs mod den virkelige eller formentlige Angriber. Dette
Instinkt er dels en særegen Form for Selvopholdelsesinstinktet,
forsaavidt som den reagerende Bevægelse fører til at skaffe
% --- p. 549 ---
\setcounter{footnote}{0}%
\opage{549}Aarsagen til Smerte bort, dels hænger det sammen med den
Trang til Bevægelse, til at faa Luft, som følger med enhver
ophidset Bevidsthedstilstand. Det virker blindt. Det kender
ingen andre Grænser end Udtømmelsen af sin Kraft. En forholdsvis
ringe Paavirkning kan fremkalde den voldsomste Tilbagevirkning,
naar Tilstanden i Forvejen er ophidset, ligesom
en lille Gnist kan fremkalde den største Explosion. Men Blindheden
viser sig ikke blot i, at der ofte ikke er noget rimeligt
Forhold mellem Paavirkning og Tilbagevirkning. Den viser sig
ogsaa i, at der ikke skelnes mellem den virkelige og den tilsyneladende
Aarsag til den smertelige Paavirkning, og deri, at
Tilbagevirkningens videre Følger ikke blive Genstand for Opmærksomhed.
Det Væsentlige er at faa Luft for den indre Uro
og Smerte; --- hvilke Objekter det gaar ud over, og hvor meget
det gaar ud over dem, spiller ingen Rolle.

Den Udvikling, som Hævninstinktet undergaar under Samfundslivets
og Kulturens Udvikling, bestaar dels i, at Genstanden
for den hævnende Tilbagevirkning individualiseres, saa at
Ingen rammes uden netop den, fra hvem Paavirkningen udgik,
--- dels i, at den hævnende Tilbagevirknings Styrke og Art bestemmes
i Forhold til Paavirkningens Styrke og Art og til den
Grad af Bevidsthed, med hvilken den udøvedes, --- dels i, at
der ogsaa tages Hensyn til de videre Følger, den hævnende
Tilbagevirkning faar for det Individ, som rammes af den, og
for Samfundet idethele. Alt dette forudsætter, at Instinktet
stanses i Farten, --- at der anbringes et Mellemrum mellem
Paavirkning og Tilbagevirkning, dette Mellemrum, der overhovedet
er saa betydningsfuldt i Bevidsthedens og særligt i
Villiens Udvikling\footnote[1]{Psykologi IV, 4. VII B, I—3.}, —og endeligt at Tilbagevirkningen ikke
bliver enkelte Individers eller Slægters, men Statens Sag. Først
ved denne Række af gennemgribende Forandringer bliver Hævnen
til Straf. Vi ville nærmere betragte nogle af disse Punkter.

2. Det enkelte Individ betragtes paa primitive Kulturtrin
ikke som isoleret, hverken hvad Rettigheder eller Pligter angaar.
Det er Familier og Slægter, ikke selvstændige Individer, der
% --- p. 550 ---
\setcounter{footnote}{0}%
\opage{550}staa overfor hverandre. Hævnen øves derfor af ethvert Medlem
af den Slægt, til hvilken det forulempede Individ hører,
mod ethvert Individ af den Slægt, til hvilken Forulemperen
hører. Blod skal flyde, naar Blod er flydt; det er Hovedsagen;
hvis Blod, det er det Underordnede. Hvor den gamle Familieorganisation
endnu raader, som i Albanien, raader ogsaa denne
Betragtningsmaade endnu. Albaneren modtager derfor ikke
den Fremmede, før han ved Forespørgsel har forvisset sig om,
at der ikke er Blod mellem deres Slægter. I Middelalderen
havde Gildebrødrene den Forpligtelse at hævne hverandre og
bøde for hverandre, som forhen Medlemmer af samme Slægt
havde. „Alle skulle bære det,“ siger en gammel Gildestatut,
„naar En forsér sig, og Alle lide det Samme“. Det skyldes den
fremskridende Frigørelse af Personlighederne, at efterhaanden
Forestillingen om Skyldens Begrænsning til det enkelte Individ
udvikler sig\footnote[1]{Post anf. Skr. p. 356 ff. --- Fustel de Coulanges: La cité antique. 4 éd. p. 108.}.

Ligesaa lidt som man oprindeligt skelnede mellem det handlende
Individ og hele hans Slægt, ligesaa lidt skelnede man
mellem den ydre Handling og den Villie, som lagde sig for
Dagen i den. Man holdt sig til Handlingens Virkning. Hævninstinktet
giver sig ikke Tid til at sætte sig ind i den Handlendes
Tilstand; det mangler desuden ogsaa Evne dertil. Man
holder sig til det Faktum, at der er udgaaet en smertende Virkning
fra ham, at han er et Væsen, der volder Smerte; Andet
bryder man sig ikke om at vide om ham. --- Man skelner ikke
mellem forsætlige og uforsætlige Forulempelser, mellem den
Skade, der voldes med velberaad Hu (dolus), og den, der skyldes
Ligegyldighed og Forsømmelighed (culpa). Det ufrivillige
og det villede Mord besmittede i lige høj Grad. Og man straffer
— endnu i Middelalderen og senere --- Dyr ligesaa vel
som Mennesker!\footnotemark[1].

At Hævnen avlede ny Hævn til Gengæld, saa at der ind-

*) Smlgn. herom Spencer: Political Institutions. p. 549 ff. --- Post:
Die Grundlagen des Rechts. p. 354 ff. --- L. Brentano: Die Arbeitergilden.
I, p. 2. --- R. Keyser: Efterladte Skrifter. Il, p. 359.
% CHECK p. 550: note mark without a note

% --- p. 551 ---
\setcounter{footnote}{0}%
\opage{551}lededes en endeløs Række af Aktioner og Reaktioner mellem
Slægterne, --- at det hele Samfund led under denne Strid, ---
at Raahed og Grusomhed næredes i Sindene, --- kom ikke i
Betragtning, hvor det ubetinget bydende Instinkt raadede.

3. Naar den Erkendelse opstaar, at de stridende Individer
og Slægter høre til ét Samfund, som lider ved deres Strid, og
naar de tillige selv føle Ulemperne ved de endeløse Hævnkampe,
kan det før omtalte Mellemrum indtræde. Der bliver
Mulighed for Mægling, og en Overvejelse anstilles, før der
skrides til Gengældelsesakten. De første Domstole vare Mæglingsdomstole.
Var den Skyldige udfunden, og kunde man ikke
bringe Forlig tilveje, blev han overladt til den Forulempedes
Hævn, enten uden videre eller paa visse Betingelser. Formaaede
den Forulempede ikke selv at hævne sig, hjalp Samfundet
ham dermed. At Statsmagtens Indgriben dog foreløbigt
ikke betragtedes som en Indgriben af en højeste Instans, ses
af, at den, der ikke var tilfreds med Dommen, i Frankrig (før
Ludvig den Helliges Retsreform) havde Ret til at udfordre
Dommeren til Tvekamp: det var en Appel fra Menneskenes Dom
til „Guds Dom gennem Sværdet“ (jugement de Dieu par I'épée).
--- Det næste Skridt er, at Samfundet selv tager hele Gengældelsen
i sin Haand. Det maa da raade over tilstrækkelig fysisk
Magt baade til at holde Hævninstinktet tilbage, og til at udføre
Dommen. Denne Magt falder naturligt sammen med den, der
opbydes mod ydre Fjender, saa at den retshaandhævende og
den forsvarende Myndighed komme til at hvile i samme Haand.
Det var vel ikke blot Hensigtsmæssighedsgrunde, som bevirkede
denne Overgang af Gengældelsesretten fra Individer og Slægter
til den højeste Samfundsmagt. Magthavernes egen Interesse,
deres Indflydelse og deres Indtægt forøgedes ved, at de greb
mæglende og dømmende ind i de enkelte Individers og Slæg-

Maine: Early Law and Custom. p. 170 ff. --- Spencer: Political
Institutions. p. 618f. --- Post: Die Grundlagen des Rechts. p. 401 ff.
--- Paa en interessant Maade fremtræder Overførelsen af Straffemyndigheden
til Statsmagten i Frankrigs Historie i Middelalderen, især i Ludvig
% CHECK p. 551: note whose mark was not found
\footnote[1]{den Helliges Retsreformer. Hen r i Martin: Histoire de France. 4. éd. IV p. 290—303.}
% --- p. 552 ---
\setcounter{footnote}{0}%
\opage{552}ters Mellemværender, og ved at Bøderne, et af de vigtigste
Gengældelsesmidler, gik i deres Kasse. Men som saa ofte i
Historien blev her det, der hovedsageligt skyldtes egoistiske
Motiver, Betingelse for et væsentligt Fremskridt i etisk-social
Henseende. Nu kunde en Strid for Retten med Grunde og Vidnesbyrd
afløse Striden med blanke Vaaben. Den hele Sag blev
nu betragtet fra et højere Stade. Det kom til at gælde om
noget Mere end blot Tilfredsstillelse af individuelle Instinkter.
Samfundet vil bevare sin Fred og sin Sikkerhed, en af dets
første Livsbetingelser, og for dette Hensyn maa baade den Forulempede
og den Forulempende, baade Hævner og Angriber
bøje sig. Der kommer nu Straffevedtægter og Straffelove, som
regulere Gengældelsen.

I Stedet for den blinde Hast, hvormed Hævnen fuldbyrdedes,
saalænge Instinktet raadede, træder nu en Afvejen af de Omstændigheder,
under hvilke Forulempelsen er øvet, og en Udmaalen
af den Grad af Lidelse, der paaføres Forulemperen.
Selvom denne Udmaalen endnu foretages paa vel udvortes Vis,
saa at man fordrer Øje for Øje og Tand for Tand, er det dog
et Fremskridt i Forhold til det Raseri, der ikke kjendte anden
Grænse end egen Udmattelse eller Ofrets Tilintetgørelse. Der
er dog nu Mulighed for, at man kan kaste Blikket længere ud
end blot til den øjeblikkelige Tilfredsstillelse af den kollektive
eller individuelle Hævntrang. Man kan overveje, om den Lidelse,
der tilføjes, virkeligt er nødvendig, og om den ikke, naar
den tilføjes i samme Maal, som Gengældelsestrangen fordrer
det, fører Misligheder med sig, som ere større end den Mislighed,
der ligger i, at Gengældelsestrangen maa hæmmes. Man
vænner sig til at betragte den hele Sag fra Samfundets Synspunkt
og at holde sig til, hvilken Betydning den Lidelse, der
tilføjes Forulemperen, har for hele Samfundet --- Forulemperen
selv indbefattet. Man opdager, at jo stærkere Statsmagten er,
des mildere kunne Straffene være, uden at Samfundets Fred og
Sikkerhed lider derunder, at paa den anden Side grusomme
Straffe vække og nære Vildhed og Raahed i Folkets Sind og
derved selv kunne blive farlige for Samfundsfreden, --- og at
det nok saa meget kommer an paa Straffens Uundgaaelighed,
% --- p. 553 ---
\setcounter{footnote}{0}%
\opage{553}som paa dens Strenghed\footnote[1]{.Man erfarede dette i det 18. Aarhundrede, da man i flere Lande søgte at virke ved stedse strengere Straffe. Om Bayern f. Ex., der særligt udmærkede sig i denne Henseende, siges der: „Forbrydelserne tiltog, jo flere Galger og Stejler der opstilledes ved Landevejene“. (A. v. Feuerbach: Biogr. Nachlass. I. p. 452). Sel ve den Grusomhed, hvormed der straffedes, bidrog sikkert til at øge Raaheden. --- Romilly motiverede sin Reform af de strenge engelske Straffelove med, at Forbrydelserne vare i Tiltagen.}. Man anerkender, at fordi et Menneske
er anklaget, er han dermed ikke berøvet sine Menneskerettigheder,
og man drager Omsorg for, at hans Sag kan blive
stillet i det gunstigst mulige Lys.

Det oprindelige Udgangspunkt i Hævntrangen er hermed
definitivt forladt. For Hævntrangen og for Gengældelseslæren
er den Lidelse, der tilføjes Forulemperen, Formaalet. Straffen
villes for dens egen Skyld, og Eftertanken omfatter ikke, hvad
der ligger udover det Øjeblik, da Gengældelsen har ramt sit
Offer. Ved Hæmning af den umiddelbare Hævn- og Gengældelsestrang,
kan virkelig Retfærdighed, der er Andet og Mere
end blot „Lige for Lige“, komme til at raade. Hovedvægten
lægges nu ikke paa det Instinkt, der kræver Tilfredsstillelse,
men paa, hvorledes Samfundets Tarv kræver, at Forulemperen
(som selv vedbliver at høre til Samfundet) skal behandles. Hvad
det primitive Instinkt aldeles saa bort fra, bliver nu det ledende
Synspunkt. Det vil vise sig, at det ogsaa kun derved er muligt
at give en etisk Begrundelse af Statens Straffemyndighed.

4. Naar det Spørgsmaal opkastes, hvorpaa Statens Straffemyndighed
begrunder sig, vil det i det Foregaaende udviklede
Forhold mellem Samfundet og Staten paa den ene Side og Individet
paa den anden Side være vejledende for vor Afgørelse.
(Smlgn. III, VIII og XXXVIII).

Det er Statens Opgave at hævde Retsordenen som en af
Grundbetingelserne for det menneskelige Samfunds Udvikling.
Uden Sikkerhed og Fred kan hverken det enkelte Individ, Familien
eller de frie Kultursamfund trives. Staten anvender derfor
sin tvingende Magt, naar et Individ vægrer sig ved at opfylde
de Pligter, som ifølge den gældende Retsorden paahvile
% --- p. 554 ---
\setcounter{footnote}{0}%
\opage{554}det. Og hvor Individet paa et eller andet Punkt bryder Retsordenen,
træder Statsmagten til for at opretholde denne. Dette
kan ikke ske blot ved at stanse og tilbagevise Angrebet. I
Bruddet paa Retsordenen kundgør der sig en Villie, som, hvis
det tillodes den at gaa konsekvent frem, som den har begyndt
i den udøvede Gerning, vilde sprænge Statssamfundet. Derfor
underkaster Staten Overtræderen en saadan Behandling, at Forholdet
mellem hans Villie og Retsordenen kan blive et harmonisk
og ikke, som i den udøvede Handling, et disharmonisk
Forhold. Den paatvinger ham en Opdragelse, ved hvilken han
kan blive i Stand til at holde sig indenfor de Grænser, Retsordenen
foreskriver. Han øves i Selvbeherskelse, og han vænnes
til Arbejde. Den optræder herved ikke som en helt fremmed
Magt overfor ham. Den maa intet Øjeblik glemme, at han
selv er et Led af Slægten og af Samfundet. Den behandler
ham ikke som blot Middel for sine Formaal, ofrer ham ikke
for Samfundets Skyld. Hvad det gælder om, er saavidt muligt
at faa Samfundets og Overtræderens Interesse til at falde sammen,
saa at social Sikkerhed og Orden hævdes igennem en
Opdragelse, der bliver til Gavn for selve den, der har krænket
dem.

Naar Staten som Retsordenens Haandhæver staar overfor
et Individ, som har brudt Retsordenen, saa er det egentligt to
etiske Systemer, der staa overfor hinanden. Med Rette har
Grotius i sit berømte Værk fremstillet dette Forhold som et
Krigsforhold, en Krig mellem den Enkelte og Samfundet, i Analogi
med Krigene mellem Enkelte indbyrdes og mellem Samfundene
indbyrdes\footnote[1]{Se Den nyere Filosofis Historie}. Princip staar mod Princip, og ingen Forhandling
er mulig, saalænge Individet fastholder de Forudsætninger,
der førte ham til Retsbruddet. Hvor et Sammenstød
mellem forskellige Principer og Systemer finder Sted, dér er,
som vi have set (III, 13), praktisk-psykologisk Paavirkning det
eneste Middel, der kan føre til en Forstaaelse. Det gælder om
Forandring af det psykologiske Grundlag, en saadan Forandring,
som al Opdragelse tilsigter. Før Opdragelsen er fuldendt, kan
% --- p. 555 ---
\setcounter{footnote}{0}%
\opage{555}Individet naturligvis ikke anerkende Berettigelsen af den Behandling,
det underkastes. Men Statens Berettigelse beror paa,
at den vil gøre ham skikket til at leve i Samfundet.

Den Opdragelse, Individet underkastes, er først og fremmest
en politisk eller juridisk Opdragelse til Legalitet. Individet
lærer, at det staar sig bedst ved at rette sig efter Loven i
sin Handlen, hvilke Tanker det saa ellers gør sig. Det er Selvopholdelsesdriften,
som der først og fremmest appelleres til.
Ved de elementære Fordringer, Retsordenen stiller, kommer det
jo, naar galt skal være, kun an paa den ydre Præstation. Men
der er ingen Grund til, at Staten ved sin Behandling af det
Individ, der har forbrudt sig, altid skal blive staaende ved den
ydre, juridiske Opdragelse. Denne bør være bestemmende for
Straffens Art og Varighed. Men derved udelukkes ikke, at man
kan forsøge at paavirke Individet i etisk Henseende for at udrydde
det Onde med Rode. Først naar Individets hele Sindelag
og Tænkemaade er udviklet saaledes, at det ikke blot af
egoistisk Interesse, men af indre Overbevisning bøjer sig for
Retsordenen, først da er denne fuldstændigt genoprettet ved
Bruddet, og mere end genoprettet, da den i Stedet for en Fjende
har vundet en Ven. Arbejde vil være det vigtigste Opdragelsesmiddel.
Meget ofte er Aarsagen til Overtrædelsen den, at Individet
søger Vinding gennem Bedrag og Overgreb istedetfor
gennem ærligt og redeligt Arbejde. Men en mere direkte Paavirkning
er ikke udelukket. Den etiske Opdragelse maa naturligvis
ikke øves saaledes, at der derved sker Brud paa Religionsfriheden.
Den maa i Form og Metode være forskellig fra
de forskellige Individer efter deres Standpunkt. Det vilde være
en Extrastraf, hvis Individet skulde nødes til at lade sig undervise
i en anden Konfession end dets egen. Hvis Individet er
konfessionsløst, maa Opdragelsen ledes af en Lærer eller en
anden til alvorlig Paavirkning egnet Personlighed. Hvor Bruddet
paa Retsordenen, hvad der saare ofte er Tilfældet, udspringer
af legemlig og aandelig Forkommenhed, vil der kunne udrettes
Noget ved, at Individet --- maaske for første Gang i sit
Liv --- kommer til at staa overfor Mennesker, som det føler,
at det kan have fuld Tillid til. Allerede dette vil, uden megen
% --- p. 556 ---
\setcounter{footnote}{0}%
\opage{556}Dogmatiseren og Moraliseren, kunne medføre vigtige etiske Forandringer
i Individets Natur. Vel hører denne etiske Opdragelse
mere ind under den Maade, hvorpaa Straffen fuldbyrdes, end
under selve Straffen. Den kan kun finde Sted, hvor Individet
viser sig modtageligt eller endog føler Trang til Paavirkning i
den omtalte Retning. Man har endog villet afvise Statens Ret
til at gribe ind i den Enkeltes etiske Udvikling, fordi intet
Menneske har Lov til at gøre sig til Dommer over Andres
indre Sindelag, og Staten ikke har Ret til at anvende Magt
for at forbedre den myndige Statsborger\footnote[1]{Fichte: Grandlage des Naturrechts. II, p. 114. --- Goos: Indledning til den danske Strafferet. p. 63. p. 8.\textsuperscript{1} Om de første Grundregler for Straffelovgivningen. (Eunomia II).}. Men fra den sociale
Etiks Synspunkt er der ikke Grund til at skelne saa skarpt
mellem Straf og Straffuldbyrdelse, som det af tekniske Grunde
maaske er nødvendigt i Retsvidenskaben. Den, som dømmes til
en Straf, dømmes jo ogsaa til Straffuldbyrdelsen og til Alt,
hvad den fører med sig, og i mange Tilfælde er den etiske
Forbedring enten en naturlig Fortsættelse af eller en nødvendig
Betingelse for den juridiske. Det vil vise sig umuligt at fastholde
en skarp Forskel mellem dem. Forholdet mellem det juridiske
og det etiske Element i selve den retsstridige Handling
vil jo ogsaa være forskelligt i hvert enkelt Tilfælde. Ved nogle
Retsbrud (som politiske Forbrydelser og Politiforseelser) kan det
etiske Element saa godt som mangle, eller der foreligger maaske
endog en af de Konflikter mellem Moral og Ret, som ikke
kunne undgaas. De største og mest indgribende Brud paa Retsordenen
behøve ikke ogsaa at være de i etisk Henseende mest
betænkelige. „Omstændighedernes Forvikling“, siger A. S. Ørsted\footnotemark[1],
„kan ofte føre et Menneske, der langtfra ikke er død
for det Gode og Ædle, til de farligste og skadeligste Forbrydelser,
medens en Forbrydelse af lidet betydelige Følger kan
forraade et i Bund og Grund fordærvet Sindelag. Hvo vil paastaa,
at Mord eller at Statsforbrydelser ubetinget røbe større
sædelig Fordærvelse end Bedrageri?“ Til dette Ørsteds Spørgsmaal
maa der efter flere Sagkyndiges Erfaringer svares, at 
% CHECK p. 556: note mark without a note
Mor%
% --- p. 557 ---
\setcounter{footnote}{0}%
\opage{557}dere netop meget ofte staa moralsk højere end andre Forbrydere.
(Se nedenfor i 6). --- Man kan derfor ikke begrunde Statens
Strafferet uden at fordre, at Straffen individualiseres saa
meget som muligt. Det er særligt det stedse forskellige Forhold
mellem det juridiske og det etiske Element, som gør denne
Fordring nødvendig. Uden Individualisering bliver det Individ,
der straffes, behandlet som blot Middel for Samfundets Trang
til at opretholde sin Retsorden. Men ved den individualiserende
Opdragelse, som vi fordre, bliver der intet Punkt, hvor den
Virksomhed, hvormed Retsordenen opretholdes, ikke tillige er
en pædagogisk Virksomhed overfor det Individ, der har overtraadt
den. Statens Straffevirksomhed faar des større etisk Begrundelse
og Berettigelse, jo mere den nærmer sig til dette
Ideal. Statens Ret til at underkaste Overtræderen en Opdragelse,
der gør ham skikket til at fyldestgøre de elementære Betingelser
for Samliv med Andre i et ordnet Samfund, hviler
væsentiigt paa samme Grundlag som dens Ret til at drage Omsorg
for Børns Opdragelse og Undervisning. Begge Steder hævder
Staten det Niveau, under hvilket dens Medlemmer ikke
maa synke, og søger at hjælpe dem op, som endnu ikke have
naat det, eller som ere sunkne ned derunder\footnotemark[1].

Den Begrundelse af Statens Straffemyndighed, der bygger
paa Straffens pædagogiske Virkning, har sin store Betydning ved
at opstille det Ideal, Straffevæsenets Ordning skal stræbe at
nærme sig til. Det er ingen Indvending mod den, at dette
Ideal ikke naas i Virkeligheden, ja, at de saakaldte Forbedringshuse
meget ofte ere Forværringsanstalter -). Dette viser kun,

) Den her givne Udvikling af Straffemyndighedens ideale Begrundelse
slutter sig nærmest til en enkelt Side af Fichtes Teori. Fichtes
Teori som Helhed er en Straffetruselsteori (ligesom A. Feuerbachs og A.
S. Ørsteds) og uholdbar, da den konsekvent maa gøre det Individ, der
straffes, til blot Middel for Samfundet. Men ved den nærmere Udvikling
af sin Teori kommer han ind paa Tanker, der kunne benyttes som
Grundlag for en helt anden Teori. løvrigt er den pædagogiske Straffeteori
af meget gammel Oprindelse. Den stammer fra Platon som en
nødvendig Konsekvens af denne store Tænkers Brud med den gamle
Gengældelseslære. (Smlgn. XII, 2).

\footnotemark[1] Garofalo: La criminologie. Paris \textsc{Isss}. p. 98; 129 ff.
% CHECK p. 557: note mark without a note

% --- p. 558 ---
\setcounter{footnote}{0}%
\opage{558}at der behøves et Ideal, som kan tjene til Maalestok og vise
den Retning, i hvilken der skal arbejdes. Og mod den Betragtning,
hin Begrundelse hviler paa, nemlig at det enkelte Individ
er og vedbliver at være Led af Samfundet, selvom det krænker
Samfundets Love, ligesaavel som syge og uarbejdsdygtige Individer
høre med til Samfundet, kan der ikke rejses Indvendinger.
Det Individ, som bryder Retsordenen, er selv blevet til og
har udviklet sig indenfor Samfundet. Dets Karakter og hele
Tænkemaade er for en stor Del bestemt ved den Aand, der
raadede i Samfundet, og ved de Forhold, Samfundet havde lagt
til Rette for det. I Retsbruddet kunne vi ofte se en simpel
Virkning af Misligheder og Brøst i Samfundet. Ved at dømme
Overtræderen dømmer Samfundet forsaavidt sig selv. Naar Samfundet
saaledes i Retsbruddet for en større eller mindre Del
maa se sit eget Værk, sin egen Skyld, forvandler Berettigelsen
til at opdrage Overtræderen sig til en Forpligtelse, og Nødvendigheden
af den individualiserende Behandlingsmaade bliver saa
meget mere indlysende. Kun en Opfattelse, der betragter Individet
som absolut Udgangspunkt for alle sine Handlinger, vilde
kunne hævde, at det ved at overtræde Retsordenen aldeles
skiller sig fra Samfundet, saa at dette ikke længere har nogen
Forpligtelse overfor det. Det er dog ikke blot indeterministiske
Anskuelser, der kunne komme til dette Resultat. Den saakaldte
italienske Skole i Kriminologien (\textsc{Lombroso} og hans Skole),
der forklarer Vaneforbryderne som atavistiske Fornyelser af en
Mennesketype, de civiliserede Samfund væsentligt ere komne
ud over, mener ligeledes at kunne hævde, at Samfundet ikke
har nogen Medskyld eller Forpligtelse overfor Forbryderne\footnote[1]{Se min Kritik af den italienske Skoles Lære i Etiske Undersøgelser. p. 73—81.}.
Det er dog klart, at det allerede er et Tegn paa Ufuldkommenhed
i Samfundet, at arvelige Dispositioner stadigt holde sig og
kun vente paa god Lejlighed til at gaa over i Virkelighed. Det
viser, at Samfundsaanden ikke virkeligt formaar at gennemtrænge
de enkelte Individers Natur, og at Opdragelse og sociale
Forhold ikke ere, som de bør være\footnotemark[1].

') Kfr. Garofalo: La criminologie. p. 55 ff., 117 ff.
% CHECK p. 558: note mark without a note

% --- p. 559 ---
\setcounter{footnote}{0}%
\opage{559}Jo mere man faar Øjet op for Samfundsforholdenes Indflydelse
paa Individernes Karakter og Handlemaade, desmindre
vil man vente store Virkninger af Straffesystemet alene. Dette
maa ses som et Led i hele den store Opdragelsesproces, Slægten
underkastes gennem Samfundets Institutioner. Som i den
sædvanlige Barneopdragelse vil ogsaa i den sociale Opdragelse
en indirekte og forebyggende Indflydelse mere og mere træde
i Stedet for den direkte Indgriben, ved hvilken man i Straffen
søger at raade Bod paa den Skade, der er sket. Allerede
\textsc{Thomas} \textsc{More} rettede i sin Utopia den Anklage mod Samfundet,
at det først frembragte Tyve og siden straffede dem. Som
i Medicinen vil ogsaa i Retslæren sikkert Hygiejne og Profylaktik
træde i Stedet for den Kureren, der først kan anvendes,
naar Ulykken er sket, og da ofte anvendes forgæves.

--- Den pædagogiske Straffeteori giver den mest ideale Begrundelse
af Statens Straffemyndighed. Men den er ikke tilstrækkelig.
Naar den ene og alene lægger Vægt paa, at det
enkelte Individ voxer ud af Samfundet og stedse vedbliver at
være Led af dette, overser den, at den, der overtræder Samfundets
Lov, optræder som Samfundets Fjende og afgiver et
Forbillede, som let kan finde Efterfølgelse. Forbrydelsen etablerer
en Krigstilstand mellem Individ og Samfund og stiller derved
en anden Opgave for Samfundet end den blot opdragende,
nemlig den at hævde Retsordenen og hindre, at det enkelte
Brud bliver et Præcedens. Derfor maa ogsaa Modsætningsforholdet
mellem Samfundet og det enkelte Individ betones. Denne
Betragtning lægges til Grund i Afskrækkelses- og Straffetruselsteorierne.
Her ses Straffens Formaal i at virke afskrækkende,
idet der statueres et Exempel, ligesom naar en Glente nagles
til Ladeporten, eller man opfatter Sagen saaledes, at Statsmagten
i sine Straffelove udtaler Trusler mod dem, der udøve visse
Handlinger, og at Straffen, som udgaar over den enkelte Forbryder,
staar som Vidnesbyrd om, at Truslerne ere alvorligt
mente. „Retten til at afskrække og true udledes af Statens Ret
til at være, af dens Betydning for menneskelig Velfærd.

Afskrækkelsesteorien, der især hyldes af Jurister og Stats%
% --- p. 560 ---
\setcounter{footnote}{0}%
\opage{560}mænd (blandt Filosoferne af Schopenhauer)\footnote[1]{Die Welt als Wille und Vorstellung. II. Kap. 47. --- Kfr. i den danske Literatur A. S. Ørsted: Om de første Grundregler for Straffelovgivningen. (Eunomia II). 1817. --- I den nyeste Tid har F. Tønnies taget Ordet for Straffetruselsteorien. (Probleme des Verbrechens und der Strafe --- i Tidsskriftet «Deutschland“. Oktbr.—Novbr. 1902).}, staar iet interessant
Forhold til Gengældelsesteorien og Opdragelsesteorien.
Den betragter ikke Straffen som Formaal, men som Middel.
Derved skiller den sig fra Gengældelseslæren og faar ligesom
den pædagogiske Teori en teleologisk Karakter. Men medens
den pædagogiske Teori betragter selve det Individ, der straffes,
som Formaal, betragter Afskrækkelsesteorien det som Middel,
idet det straffes for at indgyde Andre Skræk eller for at vise, at
Straffetruslerne virkeligt udføres. Formaalet for Straffen ligger
for den i Virkningen paa det øvrige Samfund, især i den Sikkerhed,
der opnaas. Den lægger ikke saa stor Vægt paa, hvorledes
Straffen virker paa dem, der straffes, som paa, hvorledes
den virker paa dem, der ikke straffes. Det kommer an paa, at
Haabet om at undgaa Straf slaas ned. Og det gælder især om
at hindre det første, farlige Skridt paa Forbrydelsens Bane.
Denne Teori ligner Gengældelsesteorien deri, at den ikke forudsætter
nogen Interesse for, hvorledes det gaar Individet under
og efter Straffens Fuldbyrdelse. Konsekvent vil den heller ikke
lægge nogen Vægt paa den Grad af Bevidsthed, med hvilken
Lovovertrædelsen er sket; det afgørende Synspunkt bliver Handlingens
Farlighed, der gør det nødvendigt gennem Afskrækkelsen
at sætte en Bom for alle Forsøg i lignende Retning. Hvad
dens erfaringsmæssige Virkninger angaar, har den det fælles
med Opdragelsesteorien, at den ikke fuldstændigt bekræftes ved
Erfaring. Skønt det er lettere at fremkalde Skræk end at bevirke
en virkelig Karakterforandring, saa viser Erfaringen dog,
at Forbrydelser meget ofte begaas af Individer, der flere Gange
ere straffede, og det kriminologiske Problem ligger netop i de
Tilbagefaldendes (Recidivisternes) store Antal. End ikke den
Straf, som især maatte formodes at virke afskrækkende paa
Andre, Dødsstraffen, har erfaringsmæssigt denne Virkning.

% --- p. 561 ---
\setcounter{footnote}{0}%
\opage{561}Men baade Opdragelses- og Afskrækkelsesteorien angive
Formaal, der nødvendigvis maa tilstræbes overfor Lovovertræderne.
I en fuldstændig Straffeteori maa begge være forenede.
Straffen vil da paa én Gang virke karakterændrende paa den,
der straffes, og staa som et Exempel paa, at Retsordenen ikke
kan krænkes. Det Individ, der straffes, vil da staa paa én Gang
som Formaal og som Middel. Til at gennemføre en saadan Opfattelse
udfordres ganske vist en Kunst og en Menneskekundskab,
som endnu ikke staa til Raadighed. Man er først i de
senere Aar begyndt at studere Fangerne og Fængselsvæsenet
paa videnskabelig Maade. Men den afgørende Maalestok for
Straffevæsenets Fuldkommenhed eller Ufuldkommenhed vil dog
være at hente fra den Grad, i hvilken det er lykkedes at forbinde
Opdragelse og Afskrækkelse. Denne Maalestok er kun en
speciel Anvendelse af den frie Personligheds Princip (VIII, 6),
og Straffeteorien hænger saaledes sammen med Etikens Grundtanker.
En filosofisk Straffeteori kan ikke give Andet og Mere
end en saadan med de etiske Grundtanker stemmende Maalestok
til Brug ved Vurderingen af det faktisk bestaaende Straffevæsen.
Erfaringen viser ogsaa, at det mere og mere bliver Principet
om den nøje Forbindelse af Opdragelse og Afskrækkelse, der
bestemmer Dommen over Straffevæsenet, som det fremtræder
paa givet Sted og til given Tid. Snart faar Opdragelsen, snart
Afskrækkelsen Overhaand som ledende Synspunkt\footnote[1]{Smlgn. Francis Hagerup: Et Blad af Straffens Historie. (Nordisk Tidsskrift, udg. af den Letterstedtske Forening. 1893).}. I Modsætning
til de tidligere herskende Gengældelses- og Afskrækkelsesteorier
gjorde det 18. Aarhundredes Filantroper den pædagogiske
Teori gældende. Man fordrede en human Behandling „af
dem, der skulde straffes, og stillede deres Forbedring som Straffens
Formaal. Selvom denne filantropiske Bevægelse ofte havde
en sentimental Karakter og ofte førte til Slaphed og ubetimelig
Lemfældighed i Straffevæsenet, er det utvivlsomt, at den gjorde
en god Gerning ved at hævde det straffede Individs Ret. Den
appellerede til et højere Retfærdighedsbegreb, end de tidligere
Synsmaader kendte, et Retfærdighedsbegreb, som det ganske
% --- p. 562 ---
\setcounter{footnote}{0}%
\opage{562}vist ikke lykkedes at fremstille i klare og bestemte Former. I
den nyeste Tid er der en Tendens til at fremdrage Afskrækkelsessynspunktet
igen; men da den filantropiske Bevægelse
havde ført til at sætte Frihedsstraffe i Stedet for de tidligere
saa hyppige legemlige Straffe, og da Frihedsstraffene, især de
kortvarige, have vist sig ikke at virke tilstrækkeligt afskrækkende,
staar man her overfor store praktiske og teoretiske Vanskeligheder.
Tillige have de kriminalantropologiske Undersøgelser
begyndt at yde et Indblik i Forbrydernes Natur og have
særligt godtgjort, at der findes store Forskelligheder mellem de
Individer, der overtræde Statens Love, saa at den Behandling,
der overgaar dem, ogsaa maa være forskellig og ikke blot kan
være bestemt ved selve Overtrædelsen som ydre Faktum.
Baade Gengældelses-, Afskrækkelses- og Forbedringsteorierne
oversaa de individuelle Forskelligheder og vare tilbøjelige til at
betragte alle Mennesker som ens. I den nyeste Tid har man
søgt at tage Hensyn til de individuelle Forskelligheder ved at
indføre betingede Straffedomme, d. v. s. saadanne, der kun udføres,
naar Individet endnu en Gang forser sig, og ubestemte
Straffedomme, d. v. s. saadanne, hvor Straffetidens Længde beror
paa, hvorledes Individets Karakter ytrer og udvikler sig
under Straffuldbyrdelsen. Man er tillige bleven mere og mere
klar over, at Forbrydelsen er et socialt Fænomen, der hænger
sammen med mange andre sociale Fænomener. Det kriminologiske
Problem staar ikke isoleret, men hænger nøje sammen
med de andre sociale Problemer. Man kommer derved ind paa
den allerede omtalte Opfattelse af Straffen som Middel til Bekæmpelse
af Forbrydelsen betragtet som sOcialt Onde. Denne
Opfattelse har den „internationale kriminalistiske Forening“ lagt
til Grund, og har tillige indskærpet, at Straffen ikke maa løsrives
af sin Forbindelse med de andre Midler til Bekæmpelse,
særligt til Forebyggelse af Forbrydelsen. Det ubestemte Udtryk
Bekæmpelse er altsaa sat i Stedet for de ældre Teoriers bestemtere
Udtryk Forbedring, Afskrækkelse, Gengældelse. Kriminalpolitik
træder i Stedet for Straffeteori. Det bliver fremtidig
Erfarings Sag at afgøre, hvor Meget der kan naas ad den forebyggende
Vej, især ved Forbedring af Opdragelsesvæsenet og
% --- p. 563 ---
\setcounter{footnote}{0}%
\opage{563}af de økonomiske Forhold, og hvor Meget der kun kan naas
gennem Tilføjelse af Straffelidelse. Man vil sikkert komme til
at stille større Fordringer til Samfundet, naar der skal tages
Hensyn til alle Variationer af Lovovertrædere, end da man lod
sig nøje med --- saa snart en Lovovertrædelse mentes at være
konstateret --- at sætte Vedkommende i Hullet eller paa Galejerne
eller at hugge Hovedet af ham. Stedse vil dog Maalestokken
ved Vurderingen af den Maade, paa hvilken Forbryderne
behandles, være den, i hvilken Grad det enkelte Individ
kommer til at staa som Formaal og ikke blot som Middel.
Under en eller anden Form vil det Problem, der kom frem
under Kampen mellem Gengældelses-, Afskrækkelses- og Forbedringsteorierne,
stedse fremtræde paany.

Naar man opfatter Straffen som et Middel ved Siden af
flere andre til Bekæmpelse af Forbrydelsen som socialt Fænomen,
kan det blive et Spørgsmaal, om man da endnu tager
Ordet Straf i samme Betydning som den, i hvilken det forekom
i de ældre Opfattelser. (Smlgn. ovenfor V, 5 f.). En skarpsindig
Tænker har benægtet dette. „Begrebet Straf“, siger F.
\textsc{Tonnies}, „maa forvandles til et højere Begreb, nemlig til Begrebet
om en hensigtsmæssig Behandling af Forbryderen.“ Til
en saadan hensigtsmæssig Behandling behøver i og for sig Lidelse
ikke at høre. Naar Individet er blevet tvunget til, saavidt
det er muligt, at oprette den Skade, det har voldet Andre, saa
maa dets øvrige Behandling konsekvent være af helbredende
og opdragende Art\footnotemark[1].

5. Efter den her udviklede Opfattelse straffes der ikke
blot, fordi en Overtrædelse har fundet Sted, men ogsaa for at
ingen Overtrædelse skal finde Sted (ikke blot: qvia peccatum

> F. Tönnies: Die Verhiitung des Verbrechens. (Deutsche Worte.
Wien 1891). --- Tønnies omtaler min i første Udgave af dette Værk fremstillede
Teori som en ren Forbedringsteori. Det var dog min Mening
(som jeg i de senere Udgaver har Søgt at faa klarere frem), at ogsaa Afskrækkelse
staar som nødvendigt Formaal, og at Maalestokken for et givet
Straffevæsens Fuldkommenhed ligger i den Grad, i hvilken de to 'Formaal
forbindes, navnligt forbindes saaledes, at Afskrækkelsen ikke breder
sig paa Opdragelsens Bekostning. (Se 1. Udg. p. 379).
% CHECK p. 563: note mark without a note

% --- p. 564 ---
\setcounter{footnote}{0}%
\opage{564}est, men ogsaa: ne peccetur). Det er Fremtiden, ikke Fortiden,
som gør Straffen nødvendig. At ingen Overtrædelse finder Sted,
vil jo sige, at Retsordenen staar ved Magt, og det er det, der
skal opnaas. Det er altsaa her et klart og bestemt Formaal,
som opnaas eller dog kan opnaas ved Straffen, --- om ikke ved
Straffen alene, saa ved den i Forbindelse med andre Midler
mod det sociale Onde, som Forbrydelsen er. Og dette Formaal
finder igen sin Begrundelse ved det almindelige Velfærdsprincip.

Ved den Antagelse, at Straffen maa tjene et Formaal og
ikke har sin Gyldighed „i sig selv“, staar den udviklede Opfattelse
i bestemt Modsætning til Gengældelsesteorien, ifølge
hvilken det skal være en Nedværdigelse baade af Straffemyndigheden
og af den Personlighed, der straffes, naar Straffen betragtes
som Middel. Og da Gengældelsesteorien endnu har saa
stor Indflydelse paa de almindelige Forestillinger om Forbrydelse
og Straf, at den maaske endog maa siges at være en af
de væsentligste Hindringer for en rationel Ordning af Straffevæsenet, maa vi noget nærmere undersøge denne Teori og
prøve dens Berettigelse. Gengældelsesteorien hævder det som
en evig, selvindlysende etisk Sandhed, at hvo som øver Ondt,
over hans Hoved skal Ondt komme. Men er det virkeligt selvindlysende?
Og er det en etisk Tanke? Hvorfor skal det Onde,
hvorfor skal Lidelsen i Verden forøges, fordi der en Gang er
øvet Ondt? Det er jo dog ligefrem at gøre Ondt værre, saalænge
man ikke paaviser et bestemt Gode, der opnaas ved det
Onde, man lader følge efter det en Gang udøvede. Men et
saadant Formaal vil man ikke kunne angive uden at beraabe
sig paa Velfærdsprincipet og altsaa komme ind paa «Nyttemoralen“.
Kant betragtede det som et „kategorisk Imperativ“, at et Menneske
skal behandles, som han selv behandler Andre. Men det
er Magtsprog uden Begrundelse. Derimod følger det med stor
Klarhed af Velfærdsprincipet, at Tilføjelse af Lidelse ikke er
tilstedelig, lige meget hvem den rettes imod, naar der ikke
derved opnaas et saameget des større Gode for den Enkelte
selv eller for Samfundet eller for begge. Ondt maa ikke gengældes
med Ondt, naar dette sidste Onde ikke er et forklædt
% --- p. 565 ---
\setcounter{footnote}{0}%
\opage{565}Gode. Og det følger ligeledes af Velfærdsprincipet, at Ingen
maa behandles som blot Middel --- heller ikke, eller allermindst,
som Middel for blinde Instinkters Tilfredsstillelse eller abstrakte
Maximers Anvendelse. Gengældelsen er enten et Udslag af
Hævninstinktet eller et abstrakt Fornuftprincip, der paastaas at
være selvindlysende.

,Men“ --- vil man sige --- „det var ogsaa en Fejl af Kant
at betragte hin Sætning som en ren Fornuftsætning. Den faar
først sin Betydning, naar den betragtes som Udtryk for Samfundets
etiske Følelse, der er oprørt over Forbrydelsen og først
tilfredsstilles ved Forbryderens Lidelse. Forbryderen maa udsone
sin Skyld; først derefter kan han igen gælde som Led af
Samfundet“. --- Hermed gør man ikke Sagen bedre. Ti med
hvilken Ret skal den Enkelte lide for at tilfredsstille de Andres
Følelse? Man maa godtgøre denne Følelses Værdi og Berettigelse,
før man kan fordre en Tilfredsstillelse af den, som gaar
ud over et Menneske. Menneskers Følelser kunne sættes i Bevægelse
ved mange forskellige Ting, og man maa altsaa godtgøre,
at denne bestemte Følelse har Ret til at faa sin Fordring
opfyldt. Hvis man forsøger at godtgøre dens Berettigelse, vil
man atter havne i Velfærdsprincipet. Kun da er den offentlige
Indignation og Oprørthed berettiget, naar den er motiveret ved
en Krænkelse af Samfundets Livsbetingelser og fordrer, hvad
der er fornødent til at hævde disse. Det forholder sig med den
Lidelse, Staten paafører i Straffen, som med den, der føles i
Angeren. Heller ikke denne var jo (se V, 3) „i sig selv“ af
etisk Værdi, men kun ved at være et nødvendigt psykologisk
Gennemgangsled.

Dersom den oprørte Følelse virkeligt er af sympatisk Natur
og ikke en Opblussen af det primitive Hævninstinkt, vil den
vel fordre en Soning, en bestemt Lidelse, der skal komme over
Overtræderen for hans Gerning; men den vil ikke stanse derved
som det sidste Maal. Hvis den gør det, virker kun Trangen til
at faa Luft uden Hensyn til Følgerne, denne Trang, som saa
ofte forvandler de sympatiske Følelser til Egoisme\footnote[1]{Psykologi VI C, 7.}. Den sande
% --- p. 566 ---
\setcounter{footnote}{0}%
\opage{566}Sympati er det ikke blot om at gøre at faa Luft; den vil forøge
Lystfølelsen i Verden og forminske Lidelsen; det vilde da
være en Selvmodsigelse, om den paa noget Punkt fandt sit
Maal naat ved Tilføjelsen af Lidelse. --- Ogsaa Overtræderen
selv ønsker ofte Soning. Denne Trang indeholder intet Mystisk
og heller intet Bevis for Gengældelseslæren. Det er psykologisk
forstaaeligt, at den, hos hvem Angeren begynder at virke, føler
en Lettelse ved den fysiske Lidelse. At underkaste sig denne
er en Udladning, er en Vej til at faa Luft for den stærke
aandelige Spænding. Ved at underkaste sig Lidelsen viser han
dernæst sin alvorlige Villie til at genoprette det normale Forhold
til Retsordenen; ved at lide den Straf, denne foreskriver,
fyldestgør han jo allerede en af dens Fordringer. Det forholder
sig hermed som med den smertelige Selverkendelse, hvilken
det Individ, der virkeligt har brudt med en forkastelig Villiesretning,
villigt gennemgaar. (Se V, 3). En Morder vil, som
\textsc{Fichte} har bemærket, kunne føle det retfærdigt, at han bliver
dræbt, selvom det sker ved en anden Forbryders Haand eller
ved en Magt, der ikke véd Noget om hans Forbrydelse.

Bag den „Fornuft“ og den „Følelse“, man beraaber sig
paa til Fordel for Gengældelsestankens selvindlysende Karakter,
ligger det gamle, i Menneskenaturen dybt indgroede Hævninstinkt.
Naar vi gaa tilbage til det, forstaa vi godt, at Begrundelsen
mangler. Ti Instinktet ledes ikke af Tanken om noget Formaal,
men bestaar blot i Trangen til at foretage visse Handlinger. Og
da Handlingerne i dette Tilfælde have været og tildels endnu
ere af stor Betydning for Individets og Slægtens Selvopholdelse,
forstaa vi, at Instinktet holder sig og er vanskeligt at overvinde.
Det var derfor saa mærkeligt et Fremskridt i Menneskehedens
etiske Udvikling, da man begyndte at tvivle om det
Selvindlysende i den Sætning, at Ondt skal gengældes med Ondt
(XII, 2). --- Højst forunderligt er det, at selv \textsc{Kant} kunde mene
i Gengældelsesprincipet at have en evig etisk Sandhed, og at
Fichte og A. S. \textsc{Ørsted}, der forkastede dette Princip i Retslæren,
kunde holde paa det i Etiken!\footnote[1]{Fichte: Grundlage des Naturrechts. II, p. 128. --- A. S. Ørsted: Om de første Grundregler for Straffelovgiv ningen. (Eunomia II). p. 4. ---}.

% --- p. 567 ---
\setcounter{footnote}{0}%
\opage{567}Den Energi, som lægger sig for Dagen i Hævn- og Gengældelsesinstinktet,
behøver ikke at gaa tabt, fordi Instinktet
maa undergaa en væsentlig Metamorfose. Der behøves stadigt
stærke og' levende Følelser af Indignation og Harme for at hævde
den krænkede Retsorden og træde i Skranken for den. Det er
kun det Formaalsløse i disse Følelser, vi tale imod. Fordi Tilføjelse
af Lidelse ikke er det sidste Maal, kan der meget vel
være Plads og Brug for en levende Retfærdighedsfølelse. Ved
„Nemesis“ forstod Aristoteles den Følelse, vi have, naar vi se
Lykken falde i Uværdiges Lod\footnotemark[1]. En saadan Harme over Misforhold
mellem indre og ydre Lykke er en vigtig etisk Kraft.
Men denne Harme behøver ikke at stanse ved Ønsket om Lidelse
for dem, der have øvet Ondt; den svækkes ikke ved, at
denne Lidelse betragtes som Middel til en Karakterforandring
hos Udøveren.

Hvor Gengældelsestrangen ikke ligefrem virker i Selvopholdelsens
Tjeneste, beror dens Tilfredsstillelse paa en Illusion.
Hvad opnaar man egentligt ved Gengældelsen? Det Onde, som
er sket, kan ikke gøres ugjort, og hvorledes kan det være en
Lindring for os, at Gerningsmanden kommer til at lide? Illusionen
er her en lignende som ved Misundelse og Skadefrohed.
Man har opnaat, at man har faaet Luft; men Intetsomhelst i
I en interessant Afhandling Vergeltung und Zurechnung (i 5. og 6. Bind
af „Vierteljahrsschrift für wissenschaftliche Philosophie“) har E. La as
søgt at hævde Gengældelsens etiske Betydning og at opfatte Straffen som
„eticeret Hævn“. Men jeg kan ikke finde, at han fremfører Noget, som
besvarer den ovenfor fremførte Tvivl, og det forekommer mig at være
uberettiget at bruge Ordet Hævn, naar enhver egoistisk Trang skal tænkes
borte. (Se især 5. Bind p. 152 f.). I lignende Retning som Laas udtale
sig to indbyrdes saa forskellige Tænkere som Lotze (Grundziige der
praktischen Philosophie § 64—65 ) og Dühring (Cursus der Philosophie.
p. 226). Ifølge Westermarck Origin and development of moral ideas.
I. p. 122 ere Hævn og Straf ikke Moder og Datter, men stamme fra en
fælles Moder, Følelsen af socialt Nag (public resentment). Maaske stammer
dette Nag igen fra Følelsen af Besmittelse, som Frazer lægger til Grund.
--- Deraf at Straffen har sit historiske Udspring af Slægts- og Privathævnen,
følger ikke, at dens etiske Berettigelse grunder sig paa Hævninstinktet.
(Smlgn. I, 4; XIII, 4).

\footnotemark[1] Eth. Nic. II, 7.
% CHECK p. 567: note mark without a note

% --- p. 568 ---
\setcounter{footnote}{0}%
\opage{568}Tingenes Orden er forandret, uden forsaavidt som Lidelsens
Sum i Verden er forøget. Denne Illusion ved Hævnen har
\textsc{Bjørnson} skildret i „Bergljot“:

\begin{quote}
Hævn? --- Hvem nævner Hævn?\\
Kan Hævnen vække mine Døde\\
eller dække over mig for Kulden?\\
Findes i den et tilstængt Enkesæde\\
eller Trøst for en barnløs Mor?
\end{quote}

6. Det kunde synes, som om Gengældelsesprincipet egnede
sig fortræffeligt til at anvendes ved Straffens Udmaaling. Det
synes at være en klar og rimelig Tanke, at man skal behandles
ganske, som man behandler Andre. Man finder det selvindlysende,
at den, der har dræbt en Anden, selv skal dræbes.
Men man vil allerede (paa vort nuværende Trin af etisk Udvikling)
blive noget betænkelig ved at slaa Øjet ud paa den,
der har slaaet Øjet ud paa en Anden. Og nu den, der har
begaaet Voldtægt, hvorledes skal han straffes efter Gengældelsesprincipet?
Kant mente ved Kastration. Men det er allerede
en stærk Afvigelse fra Principet Lige for Lige. Dette kan ikke
gennemføres, og hvor det synes at kunne anvendes, er det
snarere en vis Symbolik, der træder i Kraft, end en virkelig,
fornuftig Forbindelse mellem Gerning og Straf. --- Vil man
konsekvent gennemføre Gengældelsesprincipet, maa man gaa
ligesaa vidt som nogle vilde Folkeslag. En Basutoneger f. Ex.,
hvis Søn var bleven saaret i Hovedet ved et Slag med en Stok,
vilde have fat i Gerningsmanden for at slaa ham med den
samme Stok paa samme Sted af Hovedet og staaende paa samme
Plet Jord, hvorpaa Gerningsmanden havde staaet. Først saa følte
han sig fuldstændigt fyldestgjort.

Gengældelseslæren kan ikke begrunde, at Handlinger, der
udspringe af Forsæt (dolus), straffes strengere end de, der udspringe
af Uagtsomhed (culpa). Det er Handlingen, som skal
gengældes; hvorledes kan Gengældelsen blive forskellig, eftersom
Handlingen er foretagen med Villie eller af Uagtsomhed?
Det vilde dog være absurd, om man vilde udsætte den, der
ved Uagtsomhed havde voldet en Andens Død, for en lignende
% --- p. 569 ---
\setcounter{footnote}{0}%
\opage{569}Fare som den, han havde bragt den Anden i! Man maatte saa
sørge for, at den ogsaa førte Døden med sig. Gengældelsen
vilde dog ikke være fuldstændig, da Døden netop ikke vilde
blive voldet ved Uagtsomhed\footnote[1]{Smlgn. A. S. Ørsteds træffende Kritik af Gengældelse som Strafferettens Maalestok: Om de første Grundregler Eunomia II p. 12---24.}. Men kan man ikke tage Hensyn
til Forskellen mellem forsætlig og uagtsom Handling, saa bliver
det klart, at Gengældelsen i Virkeligheden er Udslag af et blindt
Instinkt, der lader os slaa til uden at undersøge Forholdene.
Det er en underlig Skæbne for en saa højt spændt etisk Anskuelse
som Gengældelseslæren, der gerne med Foragt ser ned
paa „Nyttemoralen“ og „Humaniteten“.

Gengældelsesteorien vil endeligt ikke kunne begrunde, hvad
de fleste nyere Straffelove fastsætte om Strafskyldens Forældelse.
Det er efter den Begrundelse af Straffemyndigheden, som ovenfor
er given, selvfølgeligt, at Retsbrud, der ligge langt tilbage i
Tiden og først nu opdages, ikke straffes. Det Motiv, der sætter
Straffemagten i Bevægelse, mangler nemlig her. Ligesom Etiken
ser fremad og kun ser tilbage for at kunne se desbedre fremad
(V, 3 f.), saaledes har Staten ikke med Fortiden, men med Nutiden
og Fremtiden at gøre; den langt tilbage i Tiden udøvede
Handling er hverken noget sikkert Kendemærke paa, hvorledes
Gerningsmandens Karakter nu er beskaffen, ejheller indeholder
den nogen Fare for den nu bestaaende Retsorden. Den praktiske
Interesse er falden bort. Men for Gengældelseslæren maa
der her være et Hul i Verdensordenen, og ingen nok saa lang
Tids Forløb kan begrunde, at Straffen ikke tilføjes. Hvorfor
skulde Blodet høre op at raabe?

--- I Følge den Begrundelse af Strafferetten, som ovenfor er
given, er der et tvesidigt Formaal, som skal naas ved Straffen:
Retsordenens Genoprettelse og Forandring af Gerningsmandens
Karakter. Ved ingen Straffemaade maa disse to Hensyn skilles
ad. Af de Straffe, som nutildags anvendes, er der to, som
stride herimod, nemlig Dødsstraffen og den livsvarige Fængselsstraf.
De kunne derfor ikke retfærdiggøres, og vi finde da ogsaa,
at de anvendes sjeldnere og sjeldnere. Kun en Femtedel af
% --- p. 570 ---
\setcounter{footnote}{0}%
% CHECK p. 570: notes paired with the marks in order (the layer misread a number)
\opage{570}de Dødsstraffe, der idømmes i Europa, blive fuldbyrdede. Disse
Straffe ere et Udtryk for, at vi endnu staa paa et barbarisk
Udviklingstrin. Vi leve under en Krigstilstand, der gør, at Afskrækkelsesmomentet
kommer til at spille en større Rolle, end
det etisk set kan forsvares. Men Grunden er ogsaa særligt
den, at vi endnu ere saa overordentligt langt tilbage i Straffepsykologi
og Straffepædagogik. Derfor formaa vi endnu ikke at
indvirke saaledes paa Overtræderens Karakter, at Straffen kan
afsluttes uden at ende med Døden. At lade den, der skal opdrages,
„sidde over“ paa Livstid, er en underlig pædagogisk
Metode, og ikke mindre underligt er det at tage Livet af ham
for at faa Bugt med hans Karakter. Det er let nok at erklære
et Menneske for uforbederligt; men hvorfra henter man Beviset
for, at alle Midler ere udtømte? Vi ere endnu ikke saa fremskredne
i psykologisk Indsigt og pædagogisk Praxis, at vi tør
udstede en saadan Fredløserklæring. Adskillige Erfaringer foreligge,
der vise, at dødsdømte, men derpaa benaadede Mordere
have ført et skyldfrit Liv senere hen. I nogle Tilfælde var det
maaske, fordi det slette Exempel og den ydre Lejlighed holdtes
borte; men i andre Tilfælde var det sikkert, fordi en Karakterforandring
var indtraadt\footnote[1]{Holtzendorff: Das Verbrechen des Mordes. Berlin 1875. p. 177 f. --- Prosper Despine: Psychologie naturelle. Paris 1868. Il. p. 260.}.

Efter den Undersøgelse, som tidligere (V) er anstillet om
„Villiens Frihed“, er der ingen Grund til her at gaa ind paa
dette Spørgsmaal. Naar Etiken kan forliges med Determinismen,
vil Retslæren ogsaa kunne det. Men derfor vil det alligevel ---
netop efter den her givne Begrundelse af Strafferetten --- være
af stor Vigtighed, at de subjektive Villiesforhold hos dem, hvem
Straffen skal tilføjes, blive Genstand for Opmærksomhed. ---
Det kommer for det Første an paa, hvorvidt Handlingen virkeligt
er villet, hvorvidt den er Frugt af Individets klare Overvejelse
og bevidste Beslutning, eller kun har været Genstand
for Forsæt, ikke for egentlig Beslutning\footnote[2]{Psykologi VII B, 1.}, eller kun ved Mangel
paa Opmærksomhed og Eftertanke er bleven voldet af Individet,
% --- p. 571 ---
\setcounter{footnote}{0}%
\opage{571}eller endeligt er bleven udført i en Tilstand af usædvanlig Ophidselse.
I disse forskellige Tilfælde er der en meget forskellig
Grad af Disharmoni mellem Individets Villie og Retsordenen,
og Straffen bør nødvendigvis variere derefter. Dog har man
ofte lagt for stor Vægt paa den Grad af Bevidsthed og Overlæg,
med hvilken Handlingen er bleven udført. En lang og bevidst
Overvejelse kan tyde paa, at der har været store Hindringer i
Individets Sind at overvinde, inden den forbryderiske Beslutning
kunde blive til. I saadanne Tilfælde er da Overvejelse
en Formildelsesgrund. En Karakter, der kun efter Overvindelse
af indre Modstand kan foretage en Handling, staar over den
Karakter, der slet ikke føler nogen Betænkelighed ved en saadan
Handling. --- For det Andet skal Straffen jo paavirke Individets
Villie som et Motiv, der har vist sig at være fornødent. Individet
skal opdrages. Men det kommer da an paa, om Individet
er normalt, saa at de Paavirkninger, man bringer det under,
kunne faa nogen Indflydelse. Er Individet afgjort sindssygt, vil
den opdragende Indflydelse i sædvanlig Form være unyttig eller
endog skadelig; det maa da henvises til en Sindssygeanstalt.
Hvis Individet vel ikke er sindssygt i sædvanlig Forstand, men
(selvom den Handling, der bringer det for Retten, er forholdsvis
ubetydelig) aabenbarer stærke og farlige Tendenser i sin Natur,
som det sandsynligvis ikke vil kunne overvinde ved den Evne
til Selvbeherskelse, det raader over, maa det henvises til et
„Fængselsas\textbackslash{}i“. en Mellemform mellem Fængsel og Sindssygeanstalt\footnote[1]{Smlgn. A. Forel: Du traitement des causes pathologiques du crime (Bulletin de I'union international de droit pénal. Vol. X, p. 390.}.
Og jo mere saadanne Individers Natur studeres, des
flere saadanne Mellemformer vil der kræves.

% --- p. 572 ---
\setcounter{footnote}{0}%

\opage{572}\chapterhead{XL.}{STATENS FORFATNING.}

1. Da Franskmændene i Aaret 1789 vare i Begreb med
at lave sig en helt ny Forfatning, skrev en amerikansk Statsmand,
\textsc{Morris}, der opholdt sig i Paris, at man vilde have en
amerikansk Konstitution uden at betænke, at man ikke havde
amerikanske Borgere, og at der idetmindste vilde behøves en
Menneskealder, for at man kunde blive vant til den nye Forfatning
og virkeligt leve under den. Samtidigt forudsagde Wilhelm
v. Humboldt, at den nye Forfatning ikke vilde bestaa,
fordi den var et Værk af Fornuften. Fornuften --- lærte han
--- kan vel ordne et givet Stof, men ikke selv frembringe det.
Den kan lede og vække ved at opstille Forbilleder, men de
levende Kræfter udvikle sig langsomt gennem historiske Erfaringerl).
--- Hvad saaledes en praktisk Statsmand og en filosofisk
Tænker vare enige om, har Historien tilstrækkeligt bekræftet.
Statsformer og Statsforfatninger kunne ikke laves paa én Gang
og udefra eller ovenfra paatrykkes Folket. Der skal siden Begyndelsen
af det 19. Aarhundrede være lavet 350 nye Konstitutioner\footnote[1]{Taine: La Révolution. I, p. 158. 183. --- W. v. Humboldt: Ideen iiber Staatsverfassung, durch die neue französische Constitution veranlasst. (1791). (Werke I. p. 301 ff.).}.
Allerede Tallet viser, at saare mange af dem maa
være forsvundne igen, fordi de ikke have kunnet slaa Rod.
Forfatninger maa voxe frem. Det er for os her ingen ny Sandhed.
Statsformen eller Statsforfatningen er jo kun den Del af Retsordenen,
som angaar Statens Styrelse, og det gælder for al
Ret, at kun naar den fremgaar af eller fører til sædvanemæssig
Virksomhed, faar den fast Bestaaen som Form for et virkeligt
Liv. Alt er ikke gjort med Formulering af Principer og Skemaer.
Som en Sammenligning mellem de kontinentale Folkeslags og
det engelske og nordamerikanske Folks Historie viser, udvikler
% CHECK p. 572: note whose mark was not found
\footnote[2]{Maine: Popular government. London 1885p. 174}
% --- p. 573 ---
\setcounter{footnote}{0}%
\opage{573}Forfatningslivet sig sundere og kraftigere, hvor Forfatningen bliver
til gennem en lang Række smaa Udveje eller Forbedringer, der
anvendes overfor enkelte givne Tilfælde, end naar en Forfatning
konstrueres og indføres paa én Gang. At den nordamerikanske
Unionsforfatning var enkelte Mænds bevidste Arbejde
og indførtes paa én Gang, strider ikke herimod. Ti det Geniale
ved denne Forfatning var netop den Maade, paa hvilken den
benyttede de i tidligere engelske Kolonistaters Forhold og Institutioner
liggende Muligheder\footnote[1]{Smlgn. min Af handling Alexander Hamilton og den nordamerikanske Unionsforfatning. .Mindre Arbeider.}. Denne Forfatning har da ogsaa
haft den for en „indført“ Forfatning enestaaende Lykke, at dens
Hundredaarsjubilæum har kunnet fejres under gode Udsigter
for fortsat Bestaaen. --- Fornuften, den bevidste Overvejelse,
mister, som ogsaa Humboldt antyder, derfor ikke sin Betydning.
Den virker paa mange Maader, ordnende, ledende og
vækkende, tilbage paa det helt eller halvt Ubevidste. Den skærper
Opmærksomheden for de væsentlige Betingelser, ad varer mod
de Farer, som true, og opdager nye Udviklingsmuligheder.
(Smlgn. III, 19---20).

Heraf følger, at det ikke kan være Etikens Opgave at konstruere
en ideal Statsforfatning. Det nytter ikke at forfærdige
Rammen, naar man ikke véd, om der er Noget at udfylde den
med. Hvad der er det Afgørende, er de Muligheder og Udgangspunkter,
Folket efter sin Natur og sin historiske Udvikling
frembyder. Længe førend man kan begynde at diskutere om
Forfatninger, har der faktisk udviklet sig en indre Forfatning,
som lægger sig for Dagen i Folkets Sæd og Skik, i dets Karakter
og dets Kultur, og til den maa den ydre, officielle Ordning
af Forholdet mellem Folkets forskellige Elementer slutte
sig, hvis den skal have Bestaaen og Betydning. Rammen er
naturligvis ikke ligegyldig. Er den for snever, trykker den Indholdet
og hindrer dets Udfoldelse; er den for vid, ordnes Indholdet
ikke i fast Form. Og en daarlig Ramme kan ikke blot
være en, som vilkaarligt sættes paa udefra; den kan ogsaa
uvilkaarligt have udviklet sig. Men man kan kun raade Bod
% --- p. 574 ---
\setcounter{footnote}{0}%
\opage{574}paa saadanne Misligheder ved at støtte sig til Muligheder og
Udgangspunkter, som ere givne i Folket, saaledes som dette nu
engang er. Al sund Udvikling gaar indefra udad, nedefra opad.

Den Sætning, at Forfatninger ikke vilkaarligt kunne vælges
og indføres, stiller os dog ikke aldeles magtesløse overfor Statslivet.
Vi kunne maaske ikke udrette Noget overfor Udviklingens
færdige Resultater, og vi kunne ikke paa én Gang sætte Noget
i Værk, som fordrer lang Tids Udvikling for at trives. Men vi
kunne indvirke paa de Aarsager, der bestemme den fremtidige
Udvikling. Vi kunne lære af Erfaringen og efter moden Overvejelse
lægge Forholdene bedre til Rette, end-de hidtil have
ligget. Det er kun for de tidligste Perioder af social og politisk
Udvikling, at det gælder, at de ydre Forhold bestemme Samfundets
og Statens Ordning, uden at bevidst Indsigt og Valg
spille nogensomhelst Rolle. Hvor vi finde hensigtsmæssige Ordninger
bestaaende, have vi ikke altid Ret til at antage, at deres
Opstaaen skyldes klog Forudseenhed og Beregning. Men deraf
følger dog ikke, at man (med \textsc{Spencer}) kan opstille som almindelig
Regel, at det er de givne Betingelser og ikke de bevidste
Hensigter, som bestemme Statsformen\footnote[1]{Conditions and not intentions determine. Spencer: Political Institutions. p. 395.}. Naar en højere
Grad af aandelig Udvikling er naat, opstaar der jo dog netop
Indsigt i de givne Betingelser, og denne Indsigt vil da kunne
blive vejledende for vor Handlen. Naar historisk Erfaring, social
og politisk Indsigt har udviklet sig, vilde det stride mod
al Psykologi, om de ikke ogsaa til en vis Grad bestemte vor
praktiske Adfærd. Ingen af vore Forestillinger er aldeles uden
Indflydelse paa vor Følelse og Villie. Den Spencerske Sætning
betegner Højdepunktet af Reaktion mod det 18. Aarhundredes
entusiastiske, men naive Tillid til Fornuften. Reaktionen begynder
med, at man ikke længere tror paa, at det er „Ideer“,
der regere Verden: ikke Tanker, men Følelser og Lidenskaber
bestemme de menneskelige Handlinger, og Tankerne ere selv
kun Sideresultater af Følelserne. Det næste Skridt er, at disse
Følelser og Lidenskaber igen finde deres Forklaring ved de ydre
% --- p. 575 ---
\setcounter{footnote}{0}%
\opage{575}Forhold, under hvilke Menneskene leve og virke. Tanken og
den bevidste Hensigt betragtes kun som Produkter af en hel
Aarsagsrække, hvis første Led søges i de ydre Livsforhold. Men
naar Tanken én Gang er frembragt, kan den dog virke tilbage
paa de Betingelser, den skylder sin Tilblivelse. Hvad der er
Virkning, er selv igen Aarsag. Det overser den Spencerske
Teori. --- Naar Indsigten er tilstrækkeligt grundig, vil den ogsaa
medføre en Forstaaelse af, hvor overordentligt sammensat et
menneskeligt Samfundsliv er, og hvor mangesidige og forgrenede
Virkninger derfor enhver Indgriben paa et enkelt Punkt kan
have. Den vil derfor føre til Varsomhed og til den Overbevisning,
at indirekte Indgriben snarere naar Maalet end direkte.
Den vil søge at tjene Livet, ikke at hovmesterere det.

2. Hvad der stedse bringer et irrationelt Element ind selv
i den fornuftigst udtænkte Statsordning, er den Rolle, Magten
nødvendigvis maa spille. Al Magtanvendelse er nødvendig For
at bevare Sikkerhed og Fred udadtil og indadtil, er et Vidnesbyrd
om, at der endnu virker Drifter og Lidenskaber i Menneskenaturen,
som ikke ere tillempede efter et paa en fast
Retsorden bygget Samfunds Fordringer. Det er forsaavidt, for at
bruge Taines L'dtryk, nødvendigt at have Gendarmen paa
Vagt imod den Vilde. Men Raadigheden over Magt og Udøvelsen
af Magt tiltaler ogsaa Menneskenaturens Drifter og Lidenskaber.
Al Egoisme er egentligt en Slags Magtfølelse. Den brutale Side
af Menneskets Væsen kan ligesaa vel ytre sig i raa Anvendelse
af Magten, som i Trods mod den Orden, Magten skal hævde.
Gendarmen har selv Noget af den Vilde i sig. Menneskene
maa maaske regeres med stramme Tøjler; men de, der holde
Tøjlerne, ere ogsaa Mennesker.

„Dersom Mennesker vare Engle“, sagde \textsc{Alexander} \textsc{Hamilton},
,vilde ingen Regering behøves. Dersom det var Engle,
der regerede over Mennesker, vilde man ingen Kontrol behøve
at føre med Regeringen“ *). Det Problem, der skal løses i
Statens Forfatning, er netop dette: der behøves en til Magt
støttet Kontrol for at styre de Drifter i Menneskenaturen, der
% CHECK p. 575: note whose mark was not found
\footnote[1]{The Federalist. (1787). Boston 1882. p. 399.}
% --- p. 576 ---
\setcounter{footnote}{0}%
\opage{576}true med at sprænge Retsordenen; men der behøves dernæst
ogsaa en Kontrol med dem, som udøve hin Kontrol; og hvorledes
kan denne sidste Kontrol faa den fornødne Myndighed
uden at svække den Magt, som de, der øve hin første Kontrol,
nødvendigvis maa raade over?

3. Paa lavere Trin af politisk Udvikling gaar man i Virkeligheden
ofte ud fra, at det er Engle, der regere. Man ser
med blind Tillid hen til de Regerende som Indbegrebet af al
Visdom. Saaledes i den teokratiske Stat, og i enhver Stat, hvor
Regeringen staar som Folkets Formynder. Her synes da kun
den første Art af Kontrol at findes, ikke den anden. Denne vil
dog ved nærmere Betragtning ogsaa kunne paavises. Den bestaar
i selve Folkets Tro paa Regeringens Visdom og Dygtighed.
Kun saa langt som denne Tro naar, vil Regeringen, selv den
mest patriarkalske, kunne gaa trygt. Dertil kommer Nationalfølelsen
og hvilke andre Følelser Regeringen maatte kunne vække.
Den blotte Magt vil ikke i Længden være tilstrækkelig, hvor
den indre Tilslutning mangler. De Regerende blive, især hvor
det gælder aktiv og samlet Optræden, i Virkeligheden stedse
afhængige af Folkets Stemning og Tænkemaade, hvilken de i
Regelen selv for en stor Del have. Derhos kan Regeringen
ikke umiddelbart besørge Alt; den maa have Mellemmænd, anvende
et bureaukratisk System, som den kun ufuldkomment
kan kontrollere, og af hvilket den let bliver afhængig\footnote[1]{De danske Retsfilosofer i Midten af det 18. Aarhundrede opfattede endog Enevældet som en fri Forfatning, fordi det støttede sig til den offentlige Mening og Embedsstanden. Se E. Holm: Om det Syn paa Kongemagt, Folk og borgerlig Frihed, der udviklede sig i den dansknorske Stat i Midten af 18. Aarhundrede. (Universitetsprogram). 1883. p. 66- 85. 95.}.

Den offentlige Mening, der allerede under Formynderregimentet
spiller en vigtig Rolle, vil, hvor Folkets Udvikling er
sund og kraftig, gaa over til at at blive en offentlig Villie, der
kræver aktiv Meddelagtighed i Sagernes Afgørelse (Smlgn. XXXVII,
5 og XXXVIII, 5). Forudsætningen herfor er, at Folket ved
Virksomhed indenfor det frie Samfund og i de mindre Kredse
er blevet sig sin Kraft bevidst. Den politiske Frihed i det mo%
% --- p. 577 ---
\setcounter{footnote}{0}%
\opage{577}derne Europa er et Resultat af Udviklingsprocesser, som ere
foregaaede paa den materielle og den ideelle Kulturs Omraade.
Under den frodige Fremblomstren af Handel, Industri og Agerbrug
udfoldede sig en Intelligens og en Handledygtighed, der
ikke kunde holdes indenfor saa snevre Grænser, som de gamle
Statsformer fastsatte. Gennem den frie Udvikling af det religiøse
Liv --- især hos Dissenters --- steg Selvfølelsen og Frihedsfølelsen
i en Grad, der gjorde det umuligt at udelukke Folket
fra Indflydelse paa Afgørelsen af Sager, der angik dets eget Ve
og Vel. De engelske Sekter have faaet verdenshistorisk Betydning
ved deres Indflydelse paa det engelsk-amerikanske Statsliv.
Den nyere Videnskab hidførte ikke blot ved sine Resultater et
helt nyt Blik paa Naturen og Historien, men fremkaldte især
ved sin Metode og ved den Lyst til Forsken og Kritik, den
vakte, en fri og dristig Betragtningsmaade, som ogsaa maatte
strække sig til det politiske Omraade. Man kunde ikke længere
finde sig i at blive kontrolleret uden at kontrollere. Den politiske
Frihed er et Udtryk for, at det menneskelige Liv maa
ledes af de samme Love paa alle Omraader; den, der fra Roden
af vil skaffe Trangen til politisk Frihed ud af Verden, maa
stanse alt selvstændigt Arbejde paa den materielle og ideelle
Kulturs Omraade.

Dette er ikke saaledes at forstaa, som om Folket Først kan
opdrages til politisk Frihed og derefter blive delagtig i denne.
Saalænge det ikke har Friheden, kan det ikke heller erhverve
de Evner og Følelser, som den aktive Deltagelse i Statslivet
forudsætter. Atter her gælder det aristoteliske Princip, ti kun
ved at have politisk Frihed kan et Folk blive politisk frit. Friheden
er ogsaa her baade Middel og Maal. (Smlgn. XXIII, 2).
Tvangsregimente og Formynderskab kunne ikke opdrage til Frihed.
Kun gennem Øvelse og Vane udvikles Evnerne, og Øvelse i
at bruge Friheden kan man ikke faa, naar man ikke har Friheden.
--- Denne Betragtning ophæver ganske vist ikke al
Vanskelighed; den flytter denne længere tilbage, idet det nu
maa gælde om, naar et Folk er kommet saa vidt i Udvikling,
at politisk Frihed kan blive et Middel for dets videre Udvikling.
Men der er vundet ikke Lidet ved denne ændrede Problem%
% --- p. 578 ---
\setcounter{footnote}{0}%
% CHECK p. 578: notes paired with the marks in order (the layer misread a number)
\opage{578}stilling, da det lettere kan afgøres, naar et Middel er anvendeligt,
end naar en fuldkommen Færdighed er til Stede.

Det er heller ikke saaledes at forstaa, som om den politiske
Frihed blot var en Ret, Folket erhvervede sig ved sin
Dygtighed paa andre Omraader. Dens Besiddelse og Udøvelse
kan blive en Pligt ikke mindre end en Ret. Staten bliver først
virkeligt til, naar saa Mange som muligt aktivt deltage i Statslivet.
Staten har jo ingen Bestaaen udenfor eller over de enkelte
Individer, men bestaar i dem, i deres Sind og i deres
Villie. Jo flere medvirkende Villier, des fastere Grundlag for
Statsbygningen. Den fri og levende Tilslutning a f de enkelte
Individer er Statens sande Styrke. Den politiske Frihed knytter
de enkelte Individer til Staten gennem deres Selvvirksomhed i
Statsøjemed, og det er Individernes Pligt at vise sig som aktive
Medlemmer af Staten. Store Statsmænd have derfor betragtet
den almindelige Valgret som en nødvendig Betingelse for, at
Staten kunde opbygges paa saa dyb og fast en Grund som
muligt\footnote[1]{Alexander Hamilton: The Federalist. (1787). Boston 1882. p. LIX. Smlgn. mine Mindre Arbejder. Første Række. p. 128. --- Bismarck: Smlgn. Fiirst Bismarck als Redner. Vollständige Sammlung der parlamentarischen Reden Bismarcks. III. p. 194 ff. For Bismarck havde den almindelige Valgret politisk Betydning i flere Henseender: dels som national tysk Tradition fra 1849, dels som Modvægt mod de dynastiske Interesser i Enkeltstaterne, dels som Modsætning til det paa Klassevalg og indirekte Valg grundede preussiske Deputeretkammer, dels endeligt som Skræmmemiddel overfor Udlandets reaktionære Regeringer. Smlgn. Sy bel: Die Begriindung des deutschen Reichs. Il. p. 451. 540. V. p. 379.}. Den politiske Friheds Udvidelse er ingen Opløsning
af Staten, men netop selve Statslivets Fuldkommengørelse.

Den politiske Frihed er et af de vigtigste Opdragelsesmidler
for et Folk\footnote[2]{Smlgn. James Bryce: The American Commonwealth. London 1888. II, p. 316: „Almindelig Stemmeret, saaledes som den existerer i de forenede Stater, er ikke blot et stort Sikkerhedsmiddel i Nutiden, men er maaske den vældigste opdragende Kraft, Menneskemasserne nogensinde have været underkastede“. --- Se ogsaa H. F. Feilberg: Dansk Bondeliv. København 1889. p. 130. 137. T}. Den opfordrer den Enkelte til at skaffe sig al
den Oplysning og Kundskab, han formaar, for at kunne udføre
% --- p. 579 ---
\setcounter{footnote}{0}%
\opage{579}sit Hverv som aktivt Medlem af Statssamfundet. Men især lærer
han at kaste Blikket ud over den snevre Horisont, hans Livsgerning
indskrænker ham til, og han føler sit Menneskeværd
som den, der har sin Skærv til Folkets Liv at yde. --- Foruden
at være en naturlig Frugt af Virksomheden paa Kulturens
Omraade, er den politiske Frihed altsaa ogsaa et Middel til at
befæste Staten ved at knytte Individerne nærmere til den og
en Spore til energisk Udvikling af Personligheden og dens Evner.

Selvom man ikke er nogen Tilhænger af Formynderregimente
eller Bureaukrati, kan man godt indrømme, at den politiske
Frihed har sine store Skyggesider. Den kan misbruges, hvad
jo netop det Bedste altid er udsat for. Men her drejer det sig
ikke om en historisk Skildring af den; det gælder derimod om
at vise dens Nødvendighed og den Værdi, der fra et etisk
Synspunkt maa tillægges den. Vi maa underkaste os Frihedens
Farer og Ulemper for at blive delagtige i de uundværlige Goder,
den er Betingelsen for.

4. Den politiske Frihed have vi foreløbigt defineret som
Ret til aktiv Deltagelse i Statslivet. Men hvad forstaar man
nærmere herved? Vi kunne jo ikke alle umiddelbart regere.
Sligt kunde tilnærmelsesvis lade sig gøre i Oldtidens smaa Stater,
hvor Stat og By faldt sammen; men det gik ogsaa kun dér,
saalænge de faldt sammen. For større Riger med mange Byer
og Kommuner passer det ikke. I nyere Tid sker derfor de enkelte
Individers Deltagelse i Statslivet gennem Valg af Repræsentanter.
Valget er den Villiesakt, ved hvilken Enhver --- foruden
gennem den Deltagelse i den offentlige Debat, som heller ikke
Formynderregimentet udelukker --- kan gribe ind i Statslivets
Gang. Hvor Intet kan blive Lov, --- hvor ingen Skat kan tages,
--- og hvor ingen afgørende Bestemmelse angaaende Statens
indre og ydre Skæbne kan fattes uden Indvilligelse af Folkets
valgte Repræsentanter, dér bestaar en fri Forfatning, hvad enten
den er nedskreven paa Papir eller ikke; og dér kan herske et
Tillidsforhold mellem Regering og Folk af en dybere og fastere
Art end det, der kan være til Stede under en Formynderregering,
hvor Folket passivt følger, fordi det maaske vel faar Lov at
ytre sin Mening, men ikke at have nogen Villie.

% --- p. 580 ---
\setcounter{footnote}{0}%
\opage{580}„Men“, kunde man spørge, „er det dog ikke det Vigtigste,
at de klogeste og bedste Mænd regere, hvad enten de nyde
Folkets Tillid eller ikke?“ --- Hertil maa svares, at det netop
er det store Spørgsmaal, hvorledes man skal faa afgjort, hvilke
Mænd der er de klogeste og bedste. Den aandelige Højde er
ikke saa let at maale som den legemlige. Den store Strid i
den menneskelige Verden staar jo netop om Maalestokken for
Klogskab og Godhed, og den politiske Strid er kun en Forgrening
af denne store Strid. Men saa Meget er vist, at de
virkeligt klogeste og bedste Statsmænd (man kan være klog og
god uden at være Statsmand) ville formaa at udfinde ikke blot,
hvad der tjener til Folkets Tarv, men ogsaa hvorledes det kan
virkeliggøres paa bedste Maade. De ville indse, at kun det
tjener til Folkets Tarv, som stemmer med dets Natur og dets
Kræfter, og de ville forstaa at vække og lede disse Kræfter.
De betragte ikke Folket som et dødt Stof, en Masse, der passivt
skal tage mod den Form, man vil give det. De ville ikke
blot arbejde for Folket, men ogsaa med Folket og ville formaa
at vinde Folkets Tillid, i hvilken de ville se en af de allervigtigste
Betingelser for at kunne udrette Noget.

5. Selv den frieste Forfatning kan dog ikke være ganske
simpel og usammensat, hvis den skal holde sig. Der er to indbyrdes
modsatte Farer, som den er udsat for, og som maa forebygges
ved særegne Indretninger.

Den ene Fare bestaar i, at de uvilkaarlige Tilskyndelser»
de Fordomme og Lidenskaber, der raade hos den store Mængde,
umiddelbart komme til at slaa over i Handling uden at være
prøvede og lutrede gennem rolig og alsidig Overvejelse. Staten
kommer da til at føre en Instinkt- eller Driftpolitik, ikke en
Villiespolitik. Atter her (smlgn. f. Ex. XXXIX, 1) viser sig Vigtigheden
af, at der tilvejebringes et Mellemrum mellem Tilskyndelse
og Handling. Folket maa, ligesom det enkelte Individ,
kunne vise Selvbeherskelse, maa kunne hæmme uvilkaarlige
Tilskyndelser, for at Overvejelse kan finde Sted. En saadan
Selvbeherskelse er en væsentlig Egenskab ved enhver Karakter,
og at Folket har Karakter, skal netop vise sig ved, at ikke det
enkelte Øjebliks Indskydelser og Tilskyndelser, men stadige og
% --- p. 581 ---
\setcounter{footnote}{0}%
\opage{581}Fast grundede Tanker og Følelser bestemme dets Færd. (Smlgn.
III, 4---7 og XI, 8—9). Fordi Intet kan blive gyldigt i Staten
uden Folkets Villie, derfor behøve ikke Folkets øjeblikkelige Tilskyndelser
strax at træde over i Virkeligheden. En Tidsafstand
mellem Drift og Beslutning er nødvendig. Folket maa gennem
sin Forfatning beskyttes mod sine egne Overilelser. Den Erfaring
og den Sagkundskab, der findes indenfor Folket, maa faa Lejlighed
til at gøre sig gældende, før Afgørelsen sker. Hvorledes
dette Maal kan naas i det enkelte Folks Forfatning, vil for
Størstedelen bero paa Folkets Historie. Som et enkelt Exempel
skal blot nævnes, at de nordamerikanske Fristater ved den Maade,
hvorpaa Forholdet mellem udøvende og lovgivende Magt og
mellem de to Kamre er ordnet, og ved den uafhængige Stilling,
der er givet den højeste Domstol, have udviklet en Forfatning,
som --- paa Grund af de Garantier, den i den omtalte Henseende
frembyder --- har vundet Anerkendelse seh' hos dem,
der se med meget stor Betænkelighed paa „folkeligt Regimentes“
Udvikling\footnote[1]{karakteristisk, at den Forfatter, der med størst Sagkundskab og Talent i de senere Aar har udtalt sig mod Demokratiet som Regeringsform, ikke er langt fra at betragte den nordamerikanske Forfatning som et Ideal.}. --- Det maa vel mærkes, at hvad der trænger til
at hæmmes ved Besindelsens og Overvejelsens Hjælp, ikke blot
kan være Utaalmodighed og ustyrlig Trang til Forandring, men
ogsaa konservativ Inerti. Der kan være Fordomme, Egoisme
og Sneverhed i den folkelige Betragtningsmaade, som først efterhaanden
vige for klarere Indsigt og friere Tanker. \textbackslash{}Vigtige og
velbegrundede Reformer kunne støde paa Modstand hos den
store Mængde, der i sin Antipati mod de nye Ideer maaske
endog, i hvert Tilfælde for en Tid, foretrækker et frihedsfjendsk
Regimente. Folkeafstemningernes Historie i Frankrig afgiver oplysende
Exempler herpaa. Jesuiterne i Luzern Fandt i Fyrrerne
deres Regning ved at appellere til en almindelig Folkeafstemning.
Den ved den schweiziske Forbundsforfatning af 1874 indførte
Appel til Folket (Referendum) har vist sig at være en
meget konservativ Institution. Efter de schweiziske Forbundslove
skal der ske Folkeafstemning om en Lov, naar 30,000

*) Henry Maine: Popular Government. London 1885 --- Det er
% --- p. 582 ---
\setcounter{footnote}{0}%
% CHECK p. 582: notes paired with the marks in order (the layer misread a number)
\opage{582}Vælgere fordre det. I Tidsrummet fra 1874 til 1893 anvendtes
Referendum ved 19 Love (af 196). Af disse 19 Love godkendte
Folket kun 6, men forkastede 13. I Kanton Ziirich, hvor enhver
Lov skal underkastes Folkeafstemning, godkendtes i Tidsrummet
fra 1869 til 1893 kun 97 Love af 128 Love, som
Kantonalraadet havde antaget. Gennemgaaende forkastes i de
Kantoner, hvor tvungent Referendum bestaar, en Fjerdedel eller
Mere af de af den lovgivende Magt vedtagne Love\footnote[1]{Lawrence Lowell: The Referendum and Initiative. (Journal of Ethics VI). p. 52. --- Referendum fandtes tidligere i enkelte Kantoner, ligesom det allerede i det 18. Aarhundrede fandtes i nogle af de nordamerikanske Fristater og senere i Australien. Smlgn. Stewardson i „Journal of Ethics“ XIII. p. 134.}.

En anden, helt modsat Fare er den, at der under en paa
politisk Frihed grundet Forfatning skal blive overvejet, forhandlet
og talt for meget og for længe, saa at der ikke villes og handles
med den Energi, som er nødvendig for Statens Bestaaen og
Fremgang. Mellemrummet, der har vist sig at være saa nødvendigt,
breder sig altsaa her for meget\footnote[2]{Ogsaa her have vi en psykologisk Parallel. Smlgn. Psykologi VIII B, 3.}. Denne Fare vil dog,
ligesom den foregaaende, forminskes, jo mere der raader et
Tillidsforhold mellem Folket og dets Regering. Kun saalænge
Mistanke og Frygt for Overgreb fra Magthavernes Side gør sig
gældende, vil man søge at indskrænke den udøvende Magt saa
meget som muligt. Hvor derimod et Tillidsforhold udvikler sig,
vil der ogsaa være Redebonhed til at stille saa megen Magt til
Regeringens Raadighed, som er nødvendig til at løse de foreliggende
Opgaver. Uden et saadant Tillidsforhold kan ingen fri
Stat være en stærk Stat, og kan egentligt slet ingen Stat være
en stærk Stat. --- Men det er, som allerede tidligere (XXXVIII,
7 Slutning) bemærket, netop her, det viser sig, at enhver Stat
er en Nødstat. Det rette Forhold mellem Magt og Frihed vil
først gennem mange Svingninger og Kampe, --- Kampe, i hvilke
der kan være Skyld paa begge Sider, --- tilnærmelsesvis kunne
bringes til Veje.

6. Det er, hvor politisk Frihed raader, Majoriteten af Væl%
% --- p. 583 ---
\setcounter{footnote}{0}%
\opage{583}gerne i Landet, der kommer til at bestemme Sagernes Afgørelse.
Dette kan ikke være anderledes. Naar der skal træffes en Afgørelse,
er det rimeligere, at Minoriteten retter sig efter Majoriteten,
end omvendt. Majoritetens Villie er den største bevægende
Kraft, og kun den egner sig til at repræsentere hele
Samfundets Villie, naar absolut Enstemmighed ikke er tilstede.
Den Fordring, der rettes til Minoriteten om at bøje sig for
Majoritetens Villie, har til sin Forudsætning, at ogsaa Minoriteten
vil, at Samfundet fremdeles skal bestaa; af denne forudsatte
Villie følger med Nødvendighed Anerkendelsen af den stærkeste
Villie i Samfundet som den afgørende\footnote[1]{Smlgn. John Locke: Of Civil Government. 11. 96.}. Det Eneste, Minoriteten
med Billighed kan fordre, er at faa Lejlighed til frit at
udtale sine Anskuelser og fremføre sine Grunde. Minoriteten maa
have Lov at forsøge paa at overbevise Majoriteten og derved
selv at blive Majoritet. Derfor er det af stor Betydning, at
Majoriteterne blive repræsenterede, og Tanken om Forholdstalvalgmaaden,
som viser den mest praktiske Vej hertil, vil gribe
mægtigt ind i den fremtidige politiske Udvikling. En virkelig
Repræsentation af Folket vil kun derved bringes i Stand.

Alt Stort og Godt er først udgaaet fra Enkeltmænd og har
fundet sin første Pleje i mindre Kredse. (Se XXXVIII, 5. Men
deri ligger ingen Indvending mod Majoritetens Berettigelse til'
at bestemme den endelige Afgørelse. Ingen Enkeltmand og ingen
Minoritet kan fordre umiddelbart og øjeblikkeligt Herredømme
for sine Ideer. De kunne det saa meget mindre, som jo ikke
engang Majoriteten har Ret til øjeblikkeligt at sætte sin Villie
igennem. Et Minoritetsherredømme vilde jo være et værre Tyranni
end et Majoritetsherredømme. Ti det er jo tilfældigt,
hvem Minoriteten kommer til at bestaa af, saa man har ikke
Lov at antage, at det altid skulde være de Bedste og Klogeste.
Og Minoriteten kan ligesaavel være radikal som konservativ.
Men det vil i alle Tilfælde være flere, der maa bøje deres Villie
og deres Fornuft, naar Minoriteten regerer. --- Naturligvis kan
Minoriteten have Retten paa sin Side. Den Art tragisk Konflikt
kan ligesaa lidt undgaas her, som den f. Ex. kan undgaas i
% --- p. 584 ---
\setcounter{footnote}{0}%
\opage{584}Forholdet mellem Ret og Moral (XXXVII, 3). Minoritetens Ret
kan ikke gaa videre end til frit at udvikle sine Tanker og
gøre dem gældende før Afgørelsen, samt til at kritisere selve
Afgørelsen.

Det er en hyppig Anklage mod den demokratiske Forfatning,
der lader Majoritetens Villie raade, at den fører til Majoritetstyranni,
og at den er ugunstig for originale Menneskers
og Ideers Opstaaen. Allerede i det mærkelige Skrift, der spillede
saa stor en Rolle ved den amerikanske Unionsforfatnings
Tilblivelse, har \textsc{Alexander} \textsc{Hamilton} drøftet, hvorledes denne
Fare kunde undgaas. Han saa en Skranke for Majoritetstyranniet
i den store Mængde af Interesser og Retninger, som en
Stat, der ikke er altfor lille, vil indbefatte. Ligesom den religiøse
Frihed sikres ved Sekternes Mangfoldighed, saaledes vil
den politiske Frihed sikres ved Interessernes Mangfoldighed.
Herfra hentede Hamilton netop et Argument til Fordel for Unionens
Gennemførelse, idet han mente, at Friheden vilde blive
mindre i de nordamerikanske Stater, naar enhver af dem vilde
danne et lille, afsluttet Hele for sig\footnote[1]{The Federalist. (1787). Nr. 51. (Boston 1882. p. 401).}. Da \textsc{Tocqueville} et
halvt Aarhundrede senere besøgte Amerika og derefter skrev
sin berømte Bog om Demokratiet i Amerika, fandt han i Majoritetstyranniet
en Hovedfejl ved det amerikanske Folk og dets
Regeringsform. „Jeg kender“, sagde han, „intet Land, hvor der,
i det Hele taget, hersker mindre aandelig Uafhængighed og
sand Diskussionsfrihed end i Amerika. Majoriteten drager en
frygtindjagende Kreds omkring Tanken“. Den nyeste Undersøgelse
af amerikanske Tilstande er dog kommen til et andet Resultat.
\textsc{James} \textsc{Bryce} er tilbøjelig til at tro, at Tocqueville her, som
paa andre Punkter, har været noget for hastig i sine Domme.
Han fandt i hvert Tilfælde intet Majoritetstyranni, hverken i
politisk eller i religiøs Henseende. Ligesaa lidt er han enig
med Tocqueville i den Paastand, der atter og atter er bleven
gentagen senere, at Demokratiet skulde være en ugunstig Jordbund
for originale Mennesker og originale Ideer. Han fandt,
at Amerikanerne netop havde Tendens til Heltedyrkelse, og at
% --- p. 585 ---
\setcounter{footnote}{0}%
\opage{585}den, der har ydet noget Fremragende paa et af det aandelige
Livs Omraader, agtes højere i Amerika end i de fleste europæiske
Lande. Overhovedet havde efter Bryces Erfaringer de
højere Klasser i Amerika ikke tabt i Dannelse ved den demokratiske
Udvikling, medens de lavere Klasser havde vundet i
Uafhængighed. Tocqueville synes paa doktrinær Vis at have
tillagt Regeringsformen en altfor stor Betydning for Folkets Liv\footnote[1]{James Bryce: The American Commonwealth. London 1888. I. p. 14; III. P. 134—143; 169; 534—549.}.

7. Den politiske Frihed medfører en ejendommelig Art af
Samfundsdannelse, der er ligesaa fri som Samfundsdannelse
paa Kulturens Omraade og dog er paa det Nøjeste knyttet til
Statslivet og griber ind i dette. Den fører naturligt til Sammenslutning
mellem dem, der stemme overens i de store politiske
Spørgsmaal, altsaa til Dannelse af politiske Partier. At der danner
sig forskellige Partier, har dels sin Grund deri, at de store
politiske Spørgsmaal kunne ses fra forskellige Sider, dels deri,
at de enkelte Individers Tankegang og Sindelag ere forskellige
og ved fremskridende Civilisation sikkert ville blive mere og
mere forskellige. Partivæsenet er nødvendigt overalt, hvor der
hersker Frihed, og der er neppe Nogen, som deklamerer mod
det, uden at han jo har et eller andet Parti, han bevidst eller
ubevidst undtager. De Ulemper, det medfører, kunne kun afhjælpes
ved bedre og mere energisk Brug af Friheden, altsaa
af selve den Aarsag, som har kaldt det til Live.

Det er af den største Betydning, hvilke Motiver og Synspunkter
der ere ledende ved Partidannelsen. Historien viser os
snart Racemodsætninger, snart Religionsmodsætninger, snart Klasse- og
Standsmodsætninger som bestemmende ved politiske Partiers
Dannelse. Men de medføre hver for sig store Ulemper.

Ved Racemodsætningen ligger Faren for en Deling af Staten
klart for Dagen, og ofte arbejdes der endog direkte mod dette
Maal, saa at Partimodsætningen kun er den første Antydning
af den fuldstændige Sønderdeling. Fra 1830 til 1864 bestemtes
den vigtigste Partimodsætning indenfor den daværende danske
Stat mere og mere paa denne Maade, og Resultatet blev ogsaa
en virkelig Søndersplittelse.

% --- p. 586 ---
\setcounter{footnote}{0}%
\opage{586}Hvor Religionsforskelligheder danne Grundlaget for politisk
Partideling, er det til Skade baade for Religionen, hvis Omraade
er det indvortes Menneske, og hvis Inderlighed krænkes ved,
at den bruges som Vaaben og Fane i det ydre Livs Kampe,
--- og for Staten, hvis Forhold ikke kunne ordnes efter religiøse
--- Synspunkter. Jo mere Kirken udskilles fra Staten, desmere maa
dette Grundlag for Partidannelse falde bort. Fra et Standpunkt,
som vil have, at kirkelige Ideer skulle beherske Verden, er
det konsekvent, naar man\footnote[1]{Ketteler: Freiheit, Autoritåt und Kirche. p. 99.} mener, at de politiske Partiers Forskellighed
tilsidst maa bestemmes ved Antagelsen eller Forkastelsen
af det Overnaturlige. Og i Lande, hvor Befolkningen
er delt i to store kirkelige Samfund, der hvert for sig med
Skinsyge vogter paa sin Stilling og sine Interesser, dér vil religiøs
Begrundelse af Partidelingen ikke let kunne holdes borte.
Den kan, naar den drives til det Yderste, ligesaa vel som Racemodsætningen
føre til en Spaltning af Staten. Men den er da
idetmindste naturlig og ærlig. I Lande derimod, hvor der hersker
Religionsfrihed, og hvor Mennesker af samme religiøse Tro
kunne høre til forskellige politiske Partier og Mennesker af
forskellig religiøs Tro til samme politiske Parti, --- dér kan
Beraabelse paa Religionen i Partikampen kun have til Formaal
at nedsætte og mistænkeliggøre Modstanderne.

Det kan ikke være Andet, end at Klasse- og Standsmodsætninger
maa øve stor Indflydelse paa Partidelingen. Bag ethvert
stort politisk Spørgsmaal har der hidtil staaet et socialt
Spørgsmaal. Den politiske Kamp er ofte kun Formen, under
hvilken sociale Modsætninger lægge sig for Dagen. Det følger
allerede deraf, at Fremskridtet i Historien for en stor Del bestaar
i, at stedse nye Samfundslag træde op for at faa politisk
Frihed og Ret til at tale med. Saadant var det Fremskridt, der
skete i Frankrig 1789, i England 1832, 1867 og 1884, og hos
os 1848. Det store Maal, den historiske Udvikling stiler mod,
er at gøre Ende paa al død og passiv Masse indenfor Samfundet
og at give saa Mange som muligt Adgang til fri Brug
af deres Kræfter paa det materielle og aandelige Omraade, og
% --- p. 587 ---
\setcounter{footnote}{0}%
\opage{587}politisk Frihed er dels en Følge af, dels en Aarsag til denne
Kræfternes frie Brug. --- Men hvor Stands- og Klassemodsætningen
ensidigt er bestemmende, kommer der ligesaa vel snevre
og uvedkommende Hensyn ind, som naar Racer og Sekter bestemme
Partidelingen. Særinteresserne komme til at spille den
største Rolle paa Bekostning af Samfundets fælles Vel, og Folket
spaltes i forskellige Lag, der tilsidst ikke mere forstaa hverandre.
Især træder Klasseegoismen tydeligt frem, hvor et Samfundslag,
som en Tid har været det begunstigede og det herskende,
gør Alt for at holde det næste nede, skønt Forholdenes
Natur og Vigtigheden af en sund og fredelig Udvikling synes
at indskærpe Nødvendigheden af at komme de nye Lag i Møde
paa en saadan Maade, at deres politiske Opdragelse kan ske
uden for store Vanskeligheder.

Enhver Samfundsklasse har ligesom ethvert enkelt Individ
sit bestemte Kald indenfor Samfundet som Helhed. Den ved
Klassemodsætningen bestemte Partideling maa derfor, hvis en
sund Udvikling skal finde Sted, underordnes og efterhaanden
gøre Plads for en Partideling, der udspringer af selve Grundbetingelserne
\textbackslash{}for hele Samfundets Bestaaen \textbackslash{} og Udvikling. Idealet
for Partidannelsen vilde være, at hvert Parti samlede sig om
en af disse Grundbetingelser. Det drejer sig her især om to
Principer, Orden og Fremskridt: Orden som Grundlag for Livets
rolige Gang og til Betryggelse for, hvad der er opnaaet, og
Fremskridt, fordi Forholdene ændre sig, og den, der ikke gaar
frem, gaar tilbage. Til denne Modsætning i Samfundets Grundbetingelser
svarer en Hovedforskellighed i de enkelte Individers
Begavelse og Karakter. Nogle ere anlagte til at bevare, Andre
til at erhverve; Nogles Sind er rettet med Pietet tilbage mod
det gode Gamle, Andres ser med Forhaabning fremad mod det
gode Nye. Den ideale Partideling vilde være tilstede, hvor
Individer af alle forskellige Samfundsklasser forenede sig om
det ene eller det andet af disse to store Principer. Det er ikke
noget godt Tegn, naar man strax af en Mands sociale Stilling
kan slutte sig til hans politiske Anskuelse. De to store naturbegrundede
politiske Partier bør være repræsenterede i alle
Samfundsklasser fra den højeste til den laveste. De mindre
% --- p. 588 ---
\setcounter{footnote}{0}%
\opage{588}Grupper ville da kunne finde deres Plads paa den ene eller
den anden Side. I Englands og Nordamerikas politiske Partiforhold
er det hidtil gaaet paa denne Maade, og derpaa beror
ikke mindst den sunde Karakter, den politiske Udvikling gennemgaaende
har haft i disse Lande.

Ved et saadant Grundlag for Partidelingen ville Partierne
naturligt betragte hinanden indbyrdes som Supplementer. Det
ene Parti vil føle, at det trænger til det andets Kritik. Og Kritiken
vil væsentligt gaa ud paa, at Modpartiet ikke svarer til
sit Navn og sin Idé. Det konservative Parti bliver ikke kritiseret
fra et helt fremmed Synspunkt, naar man gør det opmærksomt
paa, at dets Opgave er at bevare, hvad der har
Værdi, og at sørge for, at Kontinuiteten i Folkets Liv ikke
brydes, men ikke blindt og stædigt at modsætte sig enhver
Forandring. En overlegen konservativ Statsmand vil netop føre
sit Parti til at foretage Reformer, da de saa kunne gennemføres
med Pietet mod det Overleverede og uden altfor bratte Overgange.
Fremskridtspartiet svigter sin egen Idé, naar det i sin
Iver for at forandre det Overleverede krænker de varige og
dybereliggende Betingelser for Fremskridt, som bestaa i Folkekarakterens
Alvor og Dygtighed. Rastløs Hast og blot negativ
Kritik ere ikke gode Betingelser for Karakterudvikling. Karakteren
udvikles gennem det stadige Arbejde; Reformarbejdets overordentlige
Betydning maa ikke gøre blind for, at uden Kontinuitet
i Livsforholdene kunne de stærke og dybe Følelser ikke udvikle
sig, som ere Drivkraften ved al Handlen, hvad enten denne
gaar ud paa at reformere eller at konservere.

Modsætningen mellem det konservative Parti og Fremskridtspartiet
forbinder sig paa forskellig Maade med to andre Modsætningsforhold,
hvorved højst forskellige Kombinationer blive
mulige. Modsætningen mellem Centralisation og Decentralisation
falder ikke uden videre sammen med Modsætningen mellem
Orden og Fremskridt, og ligesaa lidt er dette Tilfældet med
Modsætningen mellem Magtudfoldelse og Kulturarbejde. Til hver
enkelt Tid bestemmes Partiforholdene i et Land ved det forskellige
Forhold mellem alle tre Modsætningspar. ---

Partidannelsen er en Organisation, uden hvilken det poli%
% --- p. 589 ---
\setcounter{footnote}{0}%
\opage{589}tiske Liv ikke kan have Sammenhæng og Klarhed. De forskellige
Ønsker og Planer, som bevæge sig i Folkets Sind, føres
ved den tilbage til nogle faa Hovedtanker. Som ethvert Samfund
har' ogsaa Partiet sin Egoisme; det kan sætte, og sætter
meget ofte, sine specielle Interesser over det Offentliges Interesser.
Men det vil, som sagt, i saadanne Tilfælde komme i Modstrid
med sin egen Idé og vil tidligere eller senere komme til at
mærke Følgerne heraf. Retsfølelsen og den offentlige Mening i
Folket vil desuden altid bevare en vis Uafhængighed af Partivæsenet
og øve sin Indflydelse paa dettes Udvikling. I Nordamerika
have professionelle Politikere skaffet sig en stor Indflydelse
indenfor de enkelte Partier og ofte misbrugt denne paa
den skammeligste Maade, --- Noget, hvis Grund maa søges dels
i de altfor talrige Valg, dels i det politiske Raastof, Indvandringen
stadigt bringer, og som let falder i Professionalisternes Hænder.
Men Folkets sunde Forstand og den offentlige Menings Magt
sætter bestemte Grænser for Partivæsenets Misbrug og formaar
at indgyde egoistiske Partigængere den fornødne Respekt\footnote[1]{James Bryce. III, p. 151 ff, 333 ff.}. Et
Parti kan ligesaa vel trænge til Reform, som hele Staten kan
det. Og Reformens Udspring er i begge Tilfælde det samme:
fra Livets frie Udvikling i de mindre Kredse komme nye Problemer
og Opgaver, til hvilke de bestaaende Partier nødes at
tage Stilling. Partierne ere Midler for Folket og skulle ligesaa
lidt som Staten hovmesterere dette. De ere i sig selv ligesaa
lidt realt produktive som Staten. De ere at ligne med Ringe,
der dannes, naar en Sten er kastet i Vandet; det egentligt
Produktive er Stenkastet; Ringene ere de Former, under hvilke
dets Virkning forplanter sig hen over Vandfladen.

Forholdet mellem den Enkelte og det Parti, han har sluttet
sig til, eller hvis Repræsentant han er, har Karakteren af et
etisk Forhold. Der kan her indtræde Konflikter af samme pinlige
og tragiske Karakter, som overalt, hvor man maa bryde
med en mindre Helhed for den større Helheds Skyld. I den
berømte Tale, i hvilken Sir Robert \textsc{Peel} tog Afsked med det
politiske Liv, sagde han bl. a.: \textasciitilde{}„Idet jeg nedlægger Magten,
% --- p. 590 ---
\setcounter{footnote}{0}%
\opage{590}vil jeg efterlade et Navn, som jeg frygter vil blive alvorligt
dadlet af mange Mænd, der uden nogen personlig Interesse,
blot med Henblik paa det offentlige Vel, bittert ville beklage
Bruddet af Partibaandene, i den Overbevisning, at Troskab mod
Partioverenskomster og Opretholdelsen af store Partier ere mægtige
og væsentlige Midler for Statsstyrelsen“. Han erkendte
altsaa, at han ved at gennemføre Kornlovenes Afskaffelse mod
sit Partis (Torypartiets) Villie havde brudt et Baand af i og for
sig etisk Natur. Men han fastholdt, at et højere etisk Hensyn
nødte ham dertil: Nødvendigheden af at skaffe en Uretfærdighed
bort, som fordyrede og forbitrede den fattige Arbejder hans daglige
Brød\footnote[1]{Abridged Edition. London 1877. p. 367.}. I denne Udtalelse er egentligt hele Partivæsenets
Etik indeholdt.

8. Selvom man tænker sig Partiorganisationen nok saa
fuldkommen, og Forholdet mellem Regering og Folk nok saa
harmonisk, ville Staten og dens Forfatning alligevel svæve i
Luften, naar der ikke gives et fastere Grundlag for det politiske
Liv end det hidtil omtalte. Det er ikke nok, at det, som rører
sig hos de enkelte Individer, i Familien og i de frie Kultursamfund
gennem den politiske Frihed og gennem Partiernes
Kamp kan faa Indflydelse paa Statslivets Udvikling. Der er her
endnu for stort et Spring fra det individuelle Liv og det frie
Samfundsliv paa den ene Side og Staten paa den anden Side.
Dette Mellemrum kan kun udfyldes af Selvstyrelsen, d. v. s.
derved, at saa mange som muligt af de offentlige Anliggender
besørges ved Borgernes egen Virksomhed og ikke blot ved et
af Regeringen indsat Bureaukrati. Selvstyrelsen danner et Mellemled
mellem Virksomheden i det frie Kultursamfund og den
egentlige Statsvirksomhed. I Kultursamfundet er det de fælles
Formaal, Vægten lægges paa; ved Selvstyrelsen er det den lokale
Sammenhøren, som gør den fælles Virksomhed naturlig.
Og ved denne af lokal Sammenhøren (Naboforhold) betingede
Samvirksomhed adskiller Selvstyrelsen sig fra det egentlige Statsliv,
som forener den over Landets hele Omraade spredte Befolk-

*) Molesworth: History of England from the year 1830—1874.
% --- p. 591 ---
\setcounter{footnote}{0}%
\opage{591}ning. De smaa lokale Kredse (Kommuner, Amter) have hver
for sig Interesser og Opgaver, hvis Tilfredsstillelse fordrer Organisation,
men som dog ikke kunne opfattes og behandles efter
deres fulde Ejendommelighed fra den fjernere staaende Centralregerings
Standpunkt. Det er fælles materielle, ideelle og filantropiske
Opgaver, der skulle løses, og ved hvis Løsning de smaa
lokale Kredse have den Fordel, at de lettere kunne indsamle
nøjagtige Erfaringer, --- at alle Individer saa omtrent kunne
kende hverandre, --- og at Maskineriet kan indskrænkes til det
mindst Mulige. Sammenhængen med Staten kan dog ikke undværes.
Dels maa Staten øve en Kontrol for at hindre Overgreb
og Vilkaarlighed overfor enkelte Individer i Kredsen, dels maa
Erfaringerne fra de andre Lokalkredse gøres frugtbringende. Der
er ligesaa lidt noget nødvendigt fjendtligt Forhold mellem Selvstyrelse
og Statsstyrelse som mellem politisk Frihed og Statsorden.
Ligesom Staten netop styrkes ved, at alle Borgere have
den Ret og Pligt at deltage i Afgørelsen af Statssagerne, saaledes
styrkes den ogsaa, jo mere selvstændigt og kraftigt Livet
rører sig i de lokale Kredse. Vi faa da ikke et bureaukratisk
Statsmaskineri, tilspidset i Regeringen, paa den ene Side og alle
de enkelte, isolerede Individer paa den anden Side. En saadan
Dualisme havde det enevældige Monarki ført til i mange europæiske
Lande\footnote[1]{Smlgn. Tocqueville: L'ancien régime et la révolution. 7 éd. p. 100. --- Gneist: Das Selfgovernment in England. 3 Aufl. p. 978. --- Wallace: Rusland. Dansk Overs. I, p. 168---174. --- E. Holm: Danmarks og Norges indre Historie under Enevælden fra 1660 til 1720. II. p. 407. --- Det franske Nationalkonvent var fjendtligt stemt mod l'esprit de localité hos Departementer og Kommuner. Ved Lov af 24. August 1793 ophævedes Kommunernes juridiske Existens, idet deres Ejendom og deres Gæld overgik til Staten. Derved mistede den politiske Frihed sine Støttepunkter. Arthur Desjardins: De la liberté politique dans l'état moderne. p. 12. --- Nationalkonventet traadte altsaa i de enevældige Herskeres Spor.}. Og det kunde ikke være anderledes. Alle
Arter Enevælde opstille et Enten---Eller: enten Stat eller Individ.
Jo mere Magt der lægges over paa den ene Side, desmindre
bliver der tilbage paa den anden Side. Fordrer man en absolut
Koncentration af al offentlig Virksomhed, saa maa den indivi%
% --- p. 592 ---
\setcounter{footnote}{0}%
\opage{592}duelle og borgerlige Selvvirksomhed forsvinde. Centralmagten
ser da med Mistanke paa enhver selvstændig Bevægelse, og
med en lignende Mistanke vil saa igen Befolkningen naturligt
komme til at betragte Regeringens Foranstaltninger, hvad enten
de gaa i reformerende eller konserverende Retning. Gennem
Selvstyrelsen vænnes derimod Individerne til at interessere sig
for og arbejde for fælles og offentlige Anliggender. Ved Virksomheden
i den mindre Kreds næres Sansen for den større
Kreds. Selvstyrelse er derfor et nødvendigt Grundlag for et
frit politisk Liv. --- I denne som i mange andre Henseender
er England et Forbillede. Den engelske Forfatning hviler, som
især \textsc{Gneist} har paavist, paa det faste Grundlag, som Selvstyrelsen
i de mindre, lokale Kredse afgiver. Den engelske Forfatning
er fremgaaet af Forvaltningen. Der er derved paa én
Gang tilvejebragt en god Skole for Deltagelsen i det politiske
Liv og tillige et Omraade, hvor fælles Virksomhed kan finde
Sted uden nødvendigvis at berøres af de store politiske Modsætninger.

Kun ved at hvile paa et saadant Grundlag bliver Staten
virkeligt, hvad den efter sin Definition skulde være: det organiserede
Folk.

\chapterhead{XLI.}{SLUTNING.}

1. Den etiske Verden bestaar af en Række mindre Verdener,
som hver for sig har sit Midtpunkt. Hvert enkelt Individ
er en lille Verden for sig, selvom en større Verdens Kræfter
gribe ind i den. Paa lignende Maade forholder det sig med
Familie, Kultursamfund og Stat. De omfattes tilsidst alle af det
store Humanitetsrige, der strækker sig ligesaa langt som Menneskeslægten,
og som endog paa enkelte Punkter udvider sig
% --- p. 593 ---
\setcounter{footnote}{0}%
\opage{593}til at omfatte alle Væsener med Evne til at føle Lyst og Ulyst.
(Smlgn. XII, 3). Flere Gange (f. Ex. III, 17; XII, 1; XIV, 5;
XXXVI, 4) er det paavist, at det vilde være en Misforstaaelse
at mene, at man bedst virkede for den store Verdens Tarv ved
at overspringe de mindre Verdener. Der er et omvendt Forhold
mellem Følelsens Styrke og Omfang, og det er derfor af afgørende
Betydning, at Interessen koncentreres i de mindre Kredse
og voxer sig stærk dér, før den gaar op i de større Kredse.
Troskab i det Smaa er Betingelsen for at kunne udrette Noget
i det Store. I den etiske Verden hersker en Kontinuitetens Lov
ligesom i den fysiske Verden; man kan ikke springe Noget over.
Ideen om det almindelige Humanitetsrige har ogsaa foreløbigt
kun den Betydning, at den forbyder os at opstille vilkaarlige
Skranker og erklærer, at ethvert menneskeligt Væsen har et
Midtpunkt i sig selv, som forhindrer, at det nogensinde kan
betragtes som blot Middel. Humanitetsriget ligger ikke udenfor
og udover Familie, Kultursamfund og Stat, men afgiver netop
--- som det stadigt har vist sig i det Foregaaende --- Principerne
for den etiske Ordning af disse Kredse. Vi ere ikke
først Familiemedlemmer, Kulturarbejdere og Statsborgere og dernæst
Mennesker; men vi skulle netop leve som Mennesker og
behandle hverandre som Mennesker i alle Forhold indenfor Familie,
Kultursamfund og Stat. Ikke blot paa Omfanget af vore
Bestræbelser, men ogsaa paa Beskaffenheden faar da Ideen om
det almindelige Menneskesamfund Betydning og lagdes derfor
ogsaa paa alle Punkter til Grund i den foregaaende Fremstilling.
Humanitetsriget kan være, hvor blot To eller Tre ere til
Stede.

Men derved udelukkes ikke, at der kan blive Tale om en
særegen Organisation af det almindelige Menneskesamfund som
et Samfund, der omfatter alle Nationer og Stater. Det vil være
en naturlig Fortsættelse af Udviklingen indenfor det enkelte
Folk. Ligesom Folket faar sin Organisation i Staten, skulde der
saa ikke kunne tænkes en Organisation, der forbandt alle Stater
indbyrdes til en stor Helhed?

De forskellige Stater leve endnu i Naturtilstanden i Forhold
til hverandre. Det er Magt, der staar mod Magt, og den
% --- p. 594 ---
\setcounter{footnote}{0}%
\opage{594}raadende Trang er Trangen til at brede sig paa Andres Bekostning.
Ideen om den evige Fred staar vel som et fjærnt Maal;
men hver enkelt Stat vil helst gaa ind til denne evige Fred
saa stor og rig som muligt. Og det er ikke engang sagt, at
man indrømmer, at dette Maal nogensinde kan naas. En stor
tysk Feltherre har endog kaldt Krigen et nødvendigt Element i
den guddommelige Verdensorden. Det er en Arv, der fra en
tidligere Tids tyske Filosofer nu synes at være gaaet over til
Generalerne, at tillægge sig et fortroligt Kendskab til den guddommelige
Verdensorden. Fra et naturligere Synspunkt staar
Krigen som en Rest af den Brutalitet, der i Regelen er betegnende
for lavere menneskelige Udviklingstrin. Deraf følger nu
ingenlunde med Nødvendighed, at Krigen nogensinde vil ophøre;
ti det er jo slet ikke sagt, at vi naa den højeste Grad af Udvikling,
vi kunne tænke os. Men i Etiken have vi heller ikke
med Fremtidsfantasier at gøre. Derimod er det en etisk Overvejelse,
hvorvidt der er Mulighed for at virkeliggøre Humanitetsriget
og Freden og trænge Brutaliteten tilbage i højere Grad
end hidtil. Derved henpeges fra en ny Side paa en hel Række
af Bestræbelser, som tildels allerede ere omtalte fra andre
Synspunkter.

2. En saadan Krigstilstand, som nu raader mellem Staterne
indbyrdes, raadede tidligere indenfor de enkelte Stater mellem
de forskellige Individer, de forskellige Slægter og Klasser. Den
moderne Stat blev først til ved Stansning af denne indre Krig.
En af de første Former, under hvilke det lykkedes at trænge
Angrebslysten og Hævntrangen tilbage, var Mæglingen. De
første Domstole vare Mæglingsdomstole. Hvad der er naaet i
det Mindre, vil man ogsaa med Grund prøve at gennemføre i
det Større. Mæglingsdomstole ere jo allerede ikke saa sjeldne
mellem Stater, selvom de Spørgsmaal, der henvises til dem,
ikke ere de egentligt brændende og afgørende; dem forbeholder
man sig helst til Afgørelse ved Magt --- naar man mener at
besidde Magten.

Større Forhaabning giver en anden Betragtning. --- Selve
Krigene føres mere og mere humant. For en Menneskealder
siden lykkedes det Genferen \textsc{Henry} \textsc{Dunant}'s filantropiske Be%
% --- p. 595 ---
\setcounter{footnote}{0}%
% CHECK p. 595: notes paired with the marks in order (the layer misread a number)
\opage{595}stræbelser at bevirke, at de Saarede og hele det militære eller
frivillige Sanitetspersonale, hvis Opgave er at pleje dem, erklæredes
for neutraliserede midt under Kampen. Efter at Dunant
med stor Opofrelse havde forberedt Sagen og forhandlet med
forskellige Regeringer om det, indbød det schweiziske Forbundsraad
til en Kongres i Genéve i August 1864, hvor den saakaldte
Genferkonvention kom i Stand. Fra den Tid af betegner
det røde Kors, Genferkonventionens Symbol, et fredlyst Sted,
selvom Kampen raser rundt om\footnote[1]{Rudolf Müller: Henry Dunant, der Begründer des Roten Kreuzes und der Genfer Konvention. Stuttgart 1896.}. —Og uagtet de store Magter
staa paa en stadig Krigsfod overfor hverandre, saa er deres
Samkvem dog langt mere betinget ved Fredens Interesser end
ved Krigens. Medens Folkeretten paa \textsc{Grotius}'s Tid væsentligt
kun bestod i Regler for Krigsførelse og Fredsslutning, som det
ogsaa angives i Titelen paa hans Værk\footnote[2]{Om Retten i Krig og Fred (De jure belli et pacis). Værket udkom 1625. I Begyndelsen af første Kapitel siger Grotius, at det er Krigen, han vil omhandle, men at han ogsaa kommer til at tale om Freden,* da Krigen udspringer af den og igen fører til den.}, saa er det nu især
Fredens Forhold og Opgaver, som afgive Stof for de folkeretlige
Bestemmelser, som overholdes af de forskellige Stater\footnote[3]{„Den moderne internationale Ret bestaar for Størstedelen af Regler, hvis Existens skyldes fysiske og sociale Betingelser, der ikke vare til for to Aarhundreder siden. Paa Grund af Letheden ved at rejse, paa Grund af Post- og Telegrafforbindelsen saavel som paa Grund af udvidede moralske Begreber og mere oplyste Handelsgrundsætninger er Nationernes indbyrdes Samkvem og Samkvemmet mellem de forskellige Nationers Borgere langt betydeligere og vigtigere i Fredstid end i Krigstid“. Sheldon A mos: The science of Law. 2. ed. London 1574. p. 341.}.
Det er Kulturens (den materielle, ideelle og filantropiske Kulturs)
Opgaver, som bringe Folkeslagene i Berøring, og saa hyppige
Krigene i nogle Perioder kunne være, saa tiltage de fredelige
Berøringer dog i den Grad i Hyppighed og Mangfoldighed, at
de langt have Overvægten. Vi have i en anden Sammenhæng
omtalt Handelens store samfundsstiftende Betydning. Den ideelle
Kultur følger i den materielles Fodspor, og gaar undertiden
endog forud for den. Baade de egoistiske og de ideelle Interesser
% --- p. 596 ---
\setcounter{footnote}{0}%
\opage{596}virke saaledes hen til at forbinde Mennesker af forskellige
Nationer. Der bestaar faktisk et Menneskesamfund, for hvilket
de nationale Grænser ingen afgørende Betydning har. --- I sit
mærkelige Skrift Zum ewigen Frieden (1795) lagde \textsc{Kant} især
Vægt paa, at de egoistiske Interesser og Handelsaanden maatte
hidføre en Fredstilstand og et universelt Statsforbund, skønt det
naturligvis var hans Overbevisning, at dette først faar sin rette
Begrundelse, naar dets Anerkendelse udspringer af et etisk Sindelag.
Desværre viser det sig, at Handelsinteresser ikke altid
forene, men ofte nok saa meget bringe i Strid. Europæiske og
asiatiske Magter strides om Markeder til Afsætning af deres
uhyre Produktion. De jage efter Kunder, og de komme i Strid
om Jagtgrundene, --- ligesom Indianerne, der følge Bøflernes
Spor, komme i Strid om Bøffeldistrikterne.

Den fredelige Udvikling vil begunstiges, jo fuldkomnere
det indenfor den enkelte Stat lykkes at trænge det brutale Magtelement
tilbage, jo mere altsaa den politiske Frihed virkeligt
kommer til at gælde. Magtanvendelse udadtil og indadtil hænge
nøje sammen.

Etisk set blev det, allerede da Menneskekærlighedens Idé
først kom frem (III, 8; 13), slaaet fast, at der gives et almindeligt
Menneskerige. Dersom nu Menneskekærligheden voxer i
Slægten, vil den ogsaa føre til dette Riges praktiske Anerkendelse
og indskrænke Krigstilstanden til det mindst Mulige. ---
Det er ikke blot mellem civiliserede Folk indbyrdes, at denne
praktiske Anerkendelse endnu i høj Grad mangler. Det gælder
ikke mindre, hvad Forholdet mellem de europæiske Folkeslag
og de andre Verdensdeles Beboere angaar. Europæerne ere
overalt optraadte som Jordens fødte Herrer og have uden videre
betragtet de andre Racer som retsløse, deres Jord som herreløs.
Ligesom Thomas \textsc{More's} Utopianere mene de, at det er retfærdigt
at bekrige et Folk, der ikke udnytter sit eget Lands
Hjælpekilder, naar man selv behøver disse til sin Bestaaen.
Dette Princip fører tit og ofte i Praxis til, at de vilde Folk betragtes
som retsløse. Men Samvittigheden har her begyndt at
røre sig. Ikke blot have Filosofer for længe siden protesteret
% --- p. 597 ---
\setcounter{footnote}{0}%
\opage{597}mod denne Adfærd\footnote[1]{Kant: Zum ewigen Frieden. Königsberg 1795. p. 42.}, men der har dannet sig „Samfund til Beskyttelse
for de uciviliserede Urfolk“\textbackslash{} og den store Statsmand,
der ledede Grundlæggelsen af Frankrigs Kolonialrige, hævdede
energisk den højere Races Pligter overfor de lavere\footnotemark[1].

Selvom den evige Fred aldrig skulde komme til at herske,
er der altsaa baade Mulighed for og etisk Nødvendighed af en
Fremgang i Retning af det Ideal, den angiver. De Bestræbelser,
som kunne føre i denne Retning, ere for Størstedelen blot Fortsættelse
af Bestræbelser, som allerede fremtræde indenfor de
mindre Kredse.

3. For næsten hundrede Aar siden opkastede Kant det
Spørgsmaal, hvorvidt Menneskeslægten var i stadig Fremadskriden
i moralsk Henseende. Erfaringen, mente han, var usikker og
tvetydig i denne Henseende, medmindre der kunde paavises en
Kendsgerning, som afgjort vidnede om et moralsk Anlæg, en
Tendens til det Gode i den menneskelige Natur. Paa en saadan
Kendsgerning kunde Troen paa Fremskridtet støtte sig. En
saadan Kendsgerning fandt han i den almindelige Begejstring,
det franske Folks Forsøg paa at grundlægge en med Fornuften
og Retfærdigheden bedre stemmende Samfundstilstand vakte rundt
om i Europa. „Den Revolution hos en aandfuld Nation“, siger
han, „som vi i vore Dage have set gaa for sig, kan lykkes
eller strande, --- den kan i den Grad være opfyldt af Elendighed
og Rædselsgerninger, at et veltænkende Menneske, selvom han,
naar han for anden Gang udførte den, kunde haabe at have
Held med sig, dog ikke vilde gøre Experimentet paa saadan
Bekostning, --- men denne Revolution finder alligevel i alle Tilskueres
Gemyt (naar de ikke ere indviklede med i Spillet) en
Deltagelse, der grænser til Entusiasme, og hvis Ytring selv er
% CHECK p. 597: note whose mark was not found
\footnote[2]{L. Felix: Der Einfluss der Sitten und Gebråuche auf die Entwickelung des Eigenthums. Leipzig 1886. p. 395. --- løvrigt oplyses det her, at ogsaa Kineserne undertrykke det himmelske Riges Urindvaanere og betragte dem som retsløse.}
% CHECK p. 597: note whose mark was not found
\footnote[3]{A. Rambaud: Jules Ferry. Paris 1903. p. 317. 516 f. (Ferry hævdede bl. A., at man ikke havde Ret til at straffe de Vilde for „Traktatbrud“, saalænge man ikke havde lært dem, hvad en Traktat en.}
% CHECK p. 597: note mark without a note

% --- p. 598 ---
\setcounter{footnote}{0}%
\opage{598}forbunden med Fare, saa at den kun kan stamme fra et moralsk
Anlæg i Menneskeslægten“. „Et saadant Fænomen i den
menneskelige Historie glemmes ikke igen, fordi det har afsløret
et Anlæg og en Evne til Fremgang i den menneskelige Natur,
som ingen Politiker havde kunnet udspekulere af Tingenes tidligere
Gang“\footnote[1]{Streit der Fakultåten. 1798. (Vermischte Schriften. Halle 1799. III) p. 440 f. 445.}.

Hvad Kant søgte, er der til alle Tider Grund til at søge.
Og han har sikkert Ret i, at det Væsentlige er, om der er en
indre Kraft og Tendens at bygge paa. Kun hvis vor Erfaring
giver os lignende Kendsgerninger, som den, Kant henviste til,
er der Grund til at dele hans Haab. Hele den Fremstilling af
Etiken, som her er given, hviler netop paa den Forudsætning,
at der i Menneskenaturen gives en Evne til universel og uinteresseret
Sympati, og den Vurdering af menneskelige Handlinger
og Livsformer, som det her er forsøgt at gennemføre, hviler
paa dette Grundlag.

Men vi vilde dog sikkert i vore Dage fordre Mere, end
det 18. Aarhundredes fremragende Aander fordrede. Det var
en storartet Idealisme, der prægede Slutningen af hint Aarhundrede.
En Tro paa den menneskelige Aands Kræfter, en Tro,
der var saameget desmere berettiget, som nogle af den menneskelige
Aands betydningsfuldeste Værker stamme fra denne
Tid, var under forskellige Former raadende hos Tidens Ordførere.
Intet Under, at man nærede den Overbevisning, at naar blot
alle Skranker og Hindringer for Menneskeaandens fri Udfoldelse
fjernedes, var der ingen Grænser for Fremskridtet. Til ingen
Tid er der maaske gaaet en renere og ædlere Begejstring
gennem Verden.

Man har oplevet Meget siden den Tid. Vi ere blevne
meget klogere og maa se til, at vi ikke blive gammelkloge. Vi
have ikke længere den faste Tro paa Menneskeaandens indre
Kraft til at bygge og ordne Livets Forhold efter sine Idealer.
Vi have faaet Blik for de mange indviklede Betingelser, som
her gøre sig gældende. Vi have lært adskillig Historie, ikke
blot den, der er gaaet for sig i det sidste Aarhundrede, og som
% --- p. 599 ---
\setcounter{footnote}{0}%
\opage{599}ligger os saa nær; ogsaa tidligere Tider staa nu for os i et
andet Lys. Vi staa mere eller mindre tvivlende overfor Reformatorer
og Samfundsforbedrere. De historiske Magters Sejghed og
Udholdenhed, den dybt liggende Aare, som bringer Fortidens
Blod til endnu at pulsere i Nutidens Organer, de faste Baand,
ved hvilke det Nuværende, selv hos dem, der mene at repræsentere
det nyeste Nye, hænger sammen med det Gamle, Glemte
eller endog Foragtede, --- alt dette have vi lært at aabne vort
Blik for. Vi beundre vel den mægtigt frembrydende Strøm,
men give os dernæst til at søge efter de mange smaa
Bække, hvis Sammenflyden den skylder sin Kraft. Den bedste, den redeligste
og den mest energiske Villie formaar Intet, naar Betingelserne
ikke ere der. Og ofte synes vi, at Betingelserne for et
Fremskridt ere saa mange, at man maa mistvivle, om de virkeligt
kunne komme til Stede. Historien faar for os derved ofte
en tragisk Karakter, idet vi saa ofte have set de bedste Kræfter
strande eller føre til Ulykker, fordi Virkeligheden gjorde for
stærk en Modstand, og de smaa Bække ikke kunde forenes til
en stor Strøm.

Dog behøver den historiske Betragtningsmaade, som hos
os er traadt i Stedet for det 18. Aarhundredes mere abstrakte
Tankeretning, ikke at kvæle den idealistiske Tro, saa streng en
Prøve denne end kan sættes paa. Vi maa blot gøre os fortrolige
med, at alt Stort tager Tid, og at Begyndelsen til det
Store ofte er uanselig. Det er en Tanke, som ogsaa i den
foregaaende Fremstilling gentagne Gange er fremdragen. Og i
Forbindelse dermed staar en anden Ting: det Uvilkaarlige gaar
forud for det Vilkaarlige; vi maa leve og gøre Erfaringer, før
vi kunne forme de Tanker, der kunne virke tilbage paa Livet;
det er Livet selv, som frembringer de Former, som vi kunne
fæstne og udvikle og bruge til at føre Livet videre. Ved Hjælp
af den inderlige Sammenhæng mellem det Smaa og det Store,
som den moderne Udviklingslære har paavist, er det os muligt
at forbinde vor Realisme med vore Forgængeres Idealisme, idet
vi hylde den Regel: „kun at begejstres for det Store, men at
være trofaste ogsaa i det Smaa“.

% The Register (pp. 601-606, PDF 621-626): to be set.

\end{document}
